%!TEX spellcheck = English
%%%%%%%%%%%%%%%%%%%%%%%%%%%%%%%%%%%%%%%%%%%%%%%%%%%%%%%%%%%%%%%%%%%%%%%%%%%%%%%%%
%%  MODELE D'EXPOSE.  Olivier Bournez.
%%
%%  Revu apres la preparation du tutoriel CiE / MCU / CCA 2026, qui a fait
%%  apparaitre une demi-douzaine de pieges couteux. Chacun est corrige ici et
%%  documente sur place, pour ne pas les repayer.
%%
%%  PAQUETS UBUNTU NECESSAIRES
%%    texlive-fonts-extra   (bbm, dsfont)      texlive-latex-extra  (thmtools)
%%    texlive-science       texlive-pictures   texlive-plain-generic (mathdots)
%%    texlive-bibtex-extra  (alphaurl.bst)     texlive-lang-french  lmodern
%%
%%  COMPILATION
%%    pdflatex, bibtex, puis pdflatex deux fois.
%%
%%  FICHIERS ATTENDUS A COTE
%%    un sous-repertoire FIG/, et \FIGCOMMONS pointant sur les
%%    logos. La bibliographie est prise dans BIBDESKDIR/@@reference-biblio.bib,
%%    accessible par copie locale ou par BIBINPUTS.
%%%%%%%%%%%%%%%%%%%%%%%%%%%%%%%%%%%%%%%%%%%%%%%%%%%%%%%%%%%%%%%%%%%%%%%%%%%%%%%%%

\providecommand\ifnotSLIDE[1]{}

\documentclass[compress,dvipsnames]{beamer}

%%%%%%%%%%%%%%%%%%%%%%%%%%%%%%%%%%%%%%%%%%%%%%%%%%%%%%%%%%%%%%%%%%%%%%%%%%%%%%%%%
%%  PIEGE 1.  Charger macros.tex dans un expose demande deux precautions.
%%
%%  (a) \providecommand\ifnotSLIDE[1]{} AVANT \documentclass, ci-dessus. C'est
%%      lui qui dit a macros.tex de ne PAS installer ses environnements de
%%      theoreme, que olivier.sty redefinit ensuite. Sans cette ligne, dix
%%      collisions sur theorem, lemma, corollary, definition, example et leurs
%%      compteurs, impossibles a neutraliser proprement ensuite.
%%
%%  (b) Cinq commandes restent definies des deux cotes. On les rend
%%      redefinissables juste avant olivier: \let ... \relax ne detruit rien,
%%      \newcommand accepte les commandes \relax, et ce sont les versions de
%%      expose.sty qui gagnent, ce qu'on veut en mode expose.
%%
%%  NE PAS charger macro-Bournez.sty. Il duplique 31 definitions de macros.tex,
%%  dont PTIME, NP, tu, powerset et les dix-huit briques GPAC, et fait passer le
%%  compte de 5 collisions a 66 sans rien apporter. Les briques de dessin sont
%%  deja dans macros.tex.
%%%%%%%%%%%%%%%%%%%%%%%%%%%%%%%%%%%%%%%%%%%%%%%%%%%%%%%%%%%%%%%%%%%%%%%%%%%%%%%%%

\providecommand\nolncs[1]{#1}	

%% Chemins par défaut (surchargables avant 
\providecommand\nolncs[1]{#1}	

%% Chemins par défaut (surchargables avant 
\providecommand\nolncs[1]{#1}	

%% Chemins par défaut (surchargables avant \input{macros} via \providecommand,
%% ou après via \renewcommand) :
\providecommand\BIBFILES{@@reference-biblio}
\providecommand\FIGCOMMONS{/Users/bournez/lib/LaTeX/Perso/FIG-COMMONS}

\newcommand\affiliationofficielleolivier{LIX, CNRS, École polytechnique, Institut Polytechnique de Paris, Palaiseau, France}
\newcommand\affiliationofficielleolivierslides{LIX, CNRS, École polytechnique, \\ Institut Polytechnique de Paris, Palaiseau, France}


%%%%%%%%%%%%%%%%%%%%%%%%
%
%              begin. Switch modes compatibilités
%
%%%%%%%%%%%%%%%%%%%%%%%

%% Apxproof mode
\providecommand\apxproofmode[1]{}
	% sinon: \newcommand\apxproof[1]{#1}
	
\providecommand\apxproofmodesubsection[1]{}
	% sinon: \newcommand\apxproofmodesubsection[1]{#1}
% utiliser: \appendixapxproofsubsection au lieu de \appendix dans ce cas. 
		

%% Lipics mode
\providecommand\lipicsmode{}
	% sinon: \newcommand\lipicsmode[1]{#1}
		% en vrai sert qu'à définir \nolipics
		
%%%%%%%%%%%%%%%%%%%%%%%%
%
%              end. Switch modes compatibilités
%
%%%%%%%%%%%%%%%%%%%%%%%


%% begin gestion des switchs

% en fait \nopilcs est la seule commande utilisée

\newcommand\nolipics[1]{#1}
\lipicsmode{\renewcommand\nolipics[1]{}}

%% Décorations (fourier + cadres colorés mdframed autour de definition,
%% theorem, exercice, example, remark, ...) : actives par défaut hors slides.
%% Pour les désactiver dans un document : \newcommand\NODECORATION[1]{}
%% avant \input{macros}. Retirées automatiquement en mode lipics.
\providecommand\NODECORATION[1]{\ifnotSLIDE{#1}}
\lipicsmode{\renewcommand\NODECORATION[1]{}}

\providecommand\PASSIAPXPROOF[1]{#1}
\apxproofmode{\renewcommand\PASSIAPXPROOF[1]{}}



%% end gestion des switchs

\def\MACROSPOINTEXOUVERT{}


\providecommand\ifnotTDEXAM[1]{#1}
\providecommand\ifnotSLIDE[1]{#1}


  
% pour preuves en appendix
\providecommand\PREUVESENAPPENDIX[1]{}



%%%% STYLES: APRES VOLS


\NODECORATION{
\usepackage{fourier}
}

%% Repli si fourier (et donc fourier-orns, qui définit \danger) n'est pas
%% chargé ; doit rester APRÈS le bloc ci-dessus, sinon fourier-orns échoue
%% avec « \danger already defined ».
\providecommand\danger{}

%%%% Box around environment
\usepackage[framemethod=tikz]{mdframed}

\newcommand\ENVCOLOR[1]{#1}
\newcommand\BEGINTEMPENVCOLOR[1]{\let\envcolorwas\ENVCOLOR\renewcommand\ENVCOLOR[1]{#1}}
\newcommand\ENDTEMPENVCOLOR{\let\ENVCOLOR\envcolorwas}


\NODECORATION{
\surroundwithmdframed[hidealllines=true,backgroundcolor=\ENVCOLOR{blue!10},roundcorner=8pt]{exercice}
\surroundwithmdframed[hidealllines=true,backgroundcolor=\ENVCOLOR{blue!10},roundcorner=8pt]{exercise}
\surroundwithmdframed[hidealllines=true,backgroundcolor=\ENVCOLOR{blue!10},roundcorner=8pt]{exercicee}
\surroundwithmdframed[hidealllines=true,backgroundcolor=\ENVCOLOR{blue!10},roundcorner=8pt]{exerciced}
\surroundwithmdframed[hidealllines=true,backgroundcolor=\ENVCOLOR{blue!10},roundcorner=8pt]{exercicec}
\surroundwithmdframed[hidealllines=true,backgroundcolor=\ENVCOLOR{blue!10},roundcorner=8pt]{exercicedc}
\surroundwithmdframed[hidealllines=true,backgroundcolor=\ENVCOLOR{orange!20},roundcorner=8pt]{example}
\surroundwithmdframed[hidealllines=true,backgroundcolor=\ENVCOLOR{orange!20},roundcorner=8pt]{exemple}
\surroundwithmdframed[hidealllines=true,backgroundcolor=\ENVCOLOR{orange!10},roundcorner=8pt]{remark}
\ifnotSLIDE{\surroundwithmdframed[backgroundcolor=\ENVCOLOR{green!15},roundcorner=8pt]{definition}}
\surroundwithmdframed[hidealllines=true,backgroundcolor=\ENVCOLOR{gray!35}]{proposition}
\ifnotSLIDE{\surroundwithmdframed[backgroundcolor=\ENVCOLOR{gray!35},roundcorner=8pt]{theorem}}
\surroundwithmdframed[hidealllines=true,backgroundcolor=\ENVCOLOR{gray!35}]{lemma}
\surroundwithmdframed[hidealllines=true,backgroundcolor=\ENVCOLOR{gray!35}]{corollary}


\surroundwithmdframed[backgroundcolor=red!35,roundcorner=8pt]{theoremaprouver}
\surroundwithmdframed[backgroundcolor=red!35,roundcorner=8pt]{conjectureaprouver}
}

%%% tcolorbox

\usepackage{tcolorbox}
\tcbuselibrary{skins}


  %%. Commande magique
%% \tcolorboxenvironment{⟨name⟩}{⟨options⟩}: attach options de tcolorbox à environnement
%% alternative: je cree des boites moi meme: voici des exemples


\newtcolorbox{maboiteconfiance}[2][]{colback=red!5!white, colframe=red!75!black,fonttitle=\bfseries, colbacktitle=red!85!black,enhanced,
attach boxed title to top center={yshift=-2mm},
  title={#2},#1}
  
\newtcolorbox{maboiteresume}[2][]{colback=red!5!white, colframe=red!75!black,fonttitle=\bfseries, colbacktitle=red!85!black,enhanced,
attach boxed title to top center={yshift=-2mm},
  title={#2},#1}
  \newtcolorbox{maboiteresumeperso}[2][]{colback=red!5!white, colframe=red!75!black,fonttitle=\bfseries, colbacktitle=red!85!black,enhanced,
attach boxed title to top center={yshift=-2mm},
  title={#2},#1}
  \newtcolorbox{maboitecommentaire}[2][]{colback=red!5!white, colframe=red!75!black,fonttitle=\bfseries, colbacktitle=red!85!black,enhanced,
attach boxed title to top center={yshift=-2mm},
  title={#2},#1}
\newtcolorbox{maboitetruc}[2][]{colback=red!5!white, colframe=red!75!black,fonttitle=\bfseries, colbacktitle=red!85!black,enhanced,
attach boxed title to top left={yshift=-2mm},
  title={#2},#1}

%% Boîte des mises en évidence CARE. Couleurs pensées pour la
%% protanopie : l'axe bleu-jaune reste discriminant (jaune = \IMPORTANT
%% d'Olivier, bleu = CARE), et la distinction ne repose pas que sur la
%% couleur (cadre + bandeau étiqueté, lisibles même en niveaux de gris).
\newtcolorbox{maboiteimportantcare}[2][]{colback=cyan!10!white, colframe=blue!60!black,fonttitle=\bfseries, colbacktitle=blue!70!black,enhanced,
attach boxed title to top center={yshift=-2mm},
  title={#2},#1}
  


%%%% FIN STYLES



%%%%%%%%%%%%%%%%%%%%%%%%%%%%%%%%%%%%%%%%%%%%%%%%%
%%%%%%%%%%%%%%%%%%%%%%%%%%%%%%%%%%%%%%%%%%%%%%%%%
%
%
%                      Version ENSEIGNEMENT 2020
%
%
%%%%%%%%%%%%%%%%%%%%%%%%%%%%%%%%%%%%%%%%%%%%%%%%%
%%%%%%%%%%%%%%%%%%%%%%%%%%%%%%%%%%%%%%%%%%%%%%%%%

%english, french
\ifdefined\CESTDUFRANCAIS
	\providecommand\LANGUE[2]{#2}
	\usepackage[french]{babel}
\else
	\providecommand\LANGUE[2]{#1}
\fi
\providecommand\LANGUEinv[2]{\LANGUE{#2}{#1}}

%%%%%%%%%%%%%%%%%%%%%%%%%%%%%%%%%%%%%%%%%%%%%%%%%
%
%
%                      Packages
%
%
%%%%%%%%%%%%%%%%%%%%%%%%%%%%%%%%%%%%%%%%%%%%%%%%%


%%% SURLIGNEUR
\providecommand\ifnotmodeDIFFUSIONFINALE[1]{
%% 2021: Comprend pas mais assure
\ifdefined\SAFEMODE
\else
#1
\fi
%#1
}
\providecommand\ifmodeDIFFUSIONFINALE[1]{
\ifdefined\SAFEMODE
#1
\else
\fi
}
\newcommand\SURLIGNE[1]{
\ifnotmodeDIFFUSIONFINALE{
\BEGINTEMPENVCOLOR{blue!30}
{                        \colorbox{blue!30}{

\noindent \begin{minipage}{\dimexpr\textwidth-2\fboxsep\relax}
#1
\end{minipage}
}
}
\ENDTEMPENVCOLOR}
\ifmodeDIFFUSIONFINALE{
\noindent \begin{minipage}{\textwidth}
#1
\end{minipage}
}
}


\newenvironment{confiance}{\begin{maboiteresume}[colback=brown!50]{Résumé ici}}{\end{maboiteresume}}
\newcommand\CONFIANCE[2][Confiance ici]{
\begin{maboiteconfiance}[colback=yellow!50]{#1}
\textbf{ATTENTION : #2}
\end{maboiteconfiance}
}

\newenvironment{resume}{\begin{maboiteresume}[colback=brown!50]{Résumé ici}}{\end{maboiteresume}}
\newcommand\RESUME[2][Résumé ici]{
\begin{maboiteresume}[colback=brown!50]{#1}
#2
\end{maboiteresume}
}

\newenvironment{resumeperso}{\begin{maboiteresumeperso}[colback=brown!50]{Résumé perso ici}}{\end{maboiteresumeperso}}
\newcommand\RESUMEPERSO[2][Résumé perso  ici]{
\begin{maboiteresumeperso}[colback=brown!50]{#1}
#2
\end{maboiteresumeperso}
%
%
%\todo[inline, color=brown]{résumé: #1}
%\colorbox{brown}{
%\noindent \begin{minipage}{\textwidth}
%{
%\bf 
%#2
%}
%\end{minipage}
%}
}

\newcommand\DANSLAMARGE[1]{
\ifnotmodeDIFFUSIONFINALE{
\marginpar{\textcolor{brown}{#1}}}
}

\newenvironment{commentaireperso}{\begin{maboitecommentaire}[colback=brown!30]{Commentaire perso} Commentaire perso:}{\end{maboitecommentaire}}
\newcommand\COMMENTAIREPERSO[2][Commentaire perso]{
\begin{maboitecommentaire}[colback=brown!30]{#1}
#2
\end{maboitecommentaire}
}


\newcommand\SURSURLIGNE[2][truc surligné ici]{
%\todo[inline, color=orange]{surligné: #1}
\ifnotmodeDIFFUSIONFINALE{
\BEGINTEMPENVCOLOR{blue!60}
	                 \colorbox{blue!60}{
\noindent \begin{minipage}{\dimexpr\textwidth-2\fboxsep\relax}
{
\bf 
#2
}
\end{minipage}
}
\ENDTEMPENVCOLOR
}
\ifmodeDIFFUSIONFINALE{#2}
}

\newcommand\IMPORTANT[1]{
\ifnotmodeDIFFUSIONFINALE{
\BEGINTEMPENVCOLOR{yellow!100}
		        \colorbox{yellow!100}{

\noindent \begin{minipage}{\dimexpr\textwidth-2\fboxsep\relax}
{ 
#1
}
\end{minipage}
}
\ENDTEMPENVCOLOR
}
\ifmodeDIFFUSIONFINALE{#1}
}

%% Pendant de \IMPORTANT, réservé aux mises en évidence CARE ;
%% \IMPORTANT reste réservé aux annotations d'Olivier. Comme \IMPORTANT,
%% redevient du texte simple en mode DIFFUSIONFINALE.
\newcommand\IMPORTANTCARE[2][Important (CARE)]{
\ifnotmodeDIFFUSIONFINALE{
\begin{maboiteimportantcare}{#1}
#2
\end{maboiteimportantcare}
}
\ifmodeDIFFUSIONFINALE{#2}
}


\newcommand\QUESTION[1]{
\ifnotmodeDIFFUSIONFINALE{
\BEGINTEMPENVCOLOR{oceanfoam!100}
		        \colorbox{oceanfoam!100}{

\noindent \begin{minipage}{\dimexpr\textwidth-2\fboxsep\relax}
{ 
#1
}
\end{minipage}
}
\ENDTEMPENVCOLOR
}
\ifmodeDIFFUSIONFINALE{#1}
}




\newcommand\SOUSLIGNE[1]{
\colorbox{lightgray}{

\noindent \begin{minipage}{\dimexpr\textwidth-2\fboxsep\relax}
{ \color{gray}
#1
}
\end{minipage}
}
}



%% Gestion des cross-references cross-referencing
\usepackage{xr-hyper}
\usepackage{hyperref}

%%
\usepackage{versions}
\usepackage{amsmath,amssymb,mathrsfs,amsfonts}
%\usepackage{mathabx}
\usepackage{listings}
\usepackage{url}
\usepackage{versions}
\usepackage{color}
\usepackage[colorinlistoftodos]{todonotes}
%\usepackage{todonotes}
\usepackage{proof}
\usepackage{xstring}

%% Index
\usepackage{makeidx}
\makeindex
\providecommand\DOINDEXPRINT{}
\providecommand\SIGNALEMOTNOUV[1]{}
\ifdefined\INDEXPRINT
	\ifnotSAFEMODE{
%	\usepackage{showidx}
%	\let\oldindex\index
%	\renewcommand{\index}[1]{\oldindex{#1}\marginpar{LA: \tiny#1}}
	\renewcommand\DOINDEXPRINT{
	      	\let\oldindex\index
		\renewcommand{\index}[1]{\oldindex{##1}\marginpar{IDXLA: \tiny##1}}
		\renewcommand\SIGNALEMOTNOUV[1]{\marginpar{MNV: \small\textbf{##1}}}
	}
	}
\fi

% === gestion des accents


%soit iso latin1
%\usepackage[cyr]{aeguill}
%soit utf8
\usepackage[utf8]{inputenc}
\usepackage[T1]{fontenc}


%%%  Pour citation \Cref (cref, ref)	

\usepackage{cleveref}


%%%%%%%%%%%%%%%%%%%%%%%%%%%%%%%%%%%%%%%%%%%%%%%%%
%
%
%                      Gestion des accents spéciaux.
%
%
%%%%%%%%%%%%%%%%%%%%%%%%%%%%%%%%%%%%%%%%%%%%%%%%%

\input{macros-accents}
%% Tolérant tant que macros-markdown.tex n'est pas dans le dépôt :
\InputIfFileExists{macros-markdown}{}{\typeout{*** macros-markdown.tex introuvable: ignore ***}}


%\usepackage{PDFdraftcopy}


%-- Pour état des chapitres


%\newcommand\ETAT{\draftstring{BROUILLON}}
%
%\newcommand\brouillon{\renewcommand\ETAT{\draftstring{BROUILLON}}}
%
%\newcommand\final{\renewcommand\ETAT{\draftstring{\rien}}}
%
%\newcommand\travail{\renewcommand\ETAT{\draftstring{CHANTIER}}}
%
%\newcommand\bazar{\renewcommand\ETAT{\draftstring{BAZAR TOTAL}}}
%
%\final



%%%%%%%%%%%%%%%%%%%%%%%%%%%%%%%%%%%%%%%%%%%%%%%%%
%
%
%                      TRADUCTION EN COURS
%
%%%%%%%%%%%%%%%%%%%%%%%%%%%%%%%%%%%%%%%%%%%%%%%%%


% === volé Amaury



\usepackage{stmaryrd}
\usepackage{ifthen}
\usepackage{mathtools}
\usepackage{dsfont}
\PASSIAPXPROOF{\usepackage{thmtools}}
\usepackage{longtable}
\usepackage{mathdots}
\usepackage{booktabs}


%%%%%%%%%%%%%%%%%%%%%%%%%%%%%%%%%%%%%%%%%%%%%%%%%
%
%
%                      TIKZ STUFFS
%
%
%%%%%%%%%%%%%%%%%%%%%%%%%%%%%%%%%%%%%%%%%%%%%%%%%

%%% volé pour dessin des GPAC: PAS SUR QUE BESOIN de TOUT

\usepackage{tikz}
\usetikzlibrary{calc}
\usepgflibrary{arrows}
\usetikzlibrary{fit}
\usetikzlibrary{patterns}
\usetikzlibrary{decorations.pathreplacing}
\usetikzlibrary{shapes.multipart}
\usetikzlibrary{shapes.geometric}
\usetikzlibrary{shapes.misc}
\usetikzlibrary{positioning}
\usetikzlibrary{graphs}
\usetikzlibrary{topaths}
\usetikzlibrary{arrows.meta}
\usetikzlibrary{decorations.pathreplacing}
\usetikzlibrary{decorations.markings}
\usetikzlibrary{decorations.pathmorphing}
%
 \usetikzlibrary{calc}
 \usetikzlibrary{automata}
\usetikzlibrary{chains}
\usetikzlibrary{scopes}
 \usetikzlibrary{positioning}
 \usetikzlibrary{through,backgrounds}
\usepackage{circuitikz}
% % pour accolades




%%%%%%%%%%%%%%%%%%%%%%%%%%%%%%%%%%%%%%%%%%%%%%%%%
%
%
%                      MACROS locales
%
%
%%%%%%%%%%%%%%%%%%%%%%%%%%%%%%%%%%%%%%%%%%%%%%%%%


%====================
%                      Moi
%=====================


\newcommand\myaddress{ 
{\sc LIX, CNRS, Ecole Polytechnique, Institut Polytechnique de Paris, }\\ 91128 Palaiseau Cedex, FRANCE \\
{\small Olivier.Bournez@lix.polytechnique.fr}}



%====================
%                      un peu de trucs sales
%=====================

%=  pour style lipics et ses trucs

\providecommand\ifnotPOLYCHAPTERMODE[1]{#1}
\providecommand\ifPOLYCHAPTERMODE[1]{}

\ifnotPOLYCHAPTERMODE{
\ifdefined\thechapter
\providecommand\spnewtheorem[4]{\newtheorem{#1}{#2}[chapter]}
\providecommand\spnewtheoremi[5]{\newtheorem{#1}[#5]{#2}}
\else
\providecommand\spnewtheorem[4]{\newtheorem{#1}{#2}}
\providecommand\spnewtheoremi[5]{\newtheorem{#1}[#5]{#2}}
\fi
}
\ifPOLYCHAPTERMODE{
\providecommand\spnewtheorem[4]{\newtheorem{#1}{#2}}
\providecommand\spnewtheoremi[5]{\newtheorem{#1}[#5]{#2}}
}
\providecommand\nolipics[1]{#1}




%====================
%                      un peu de trucs franchement sales
%=====================


% on l'inclus pas ici.
%\input{macros-moins-propres}

%%%%%%%%%%%%%%%%%%%%%%%%%%%%%%%%%%%%%%%%%%%%%%%%%
%
%
%                      Page de Garde Cours X
%
%
%%%%%%%%%%%%%%%%%%%%%%%%%%%%%%%%%%%%%%%%%%%%%%%%%


\includeversion{metdate}

\newcommand\XPAGEDEGARDE[2]{
\thispagestyle{empty}
\title{{#1} \\[2cm] %\\[3cm] 
{\large #2}
}
\author{
 Olivier Bournez% \inst{1}
\\[0.5cm] 
bournez@lix.polytechnique.fr%{\myaddress}
}
%\date{
%\vspace{2cm}
%\includegraphics[height=2.5cm]{/Users/bournez/lib/LaTeX/Perso/FIG-COMMONS/logo_X_new}
%}
\begin{metdate}
\date{
\vspace{2cm}
Version \LANGUE{of}{du} \today \\[1cm] 
\includegraphics[height=3.5cm]{/Users/bournez/lib/LaTeX/Perso/FIG-COMMONS/logo_X_new}
}
\end{metdate}
\maketitle
}


%%%%%%%%%%%%%%%%%%%%%%%%%%%%%%%%%%%%%%%%%%%%%%%%%
%
%
%                      MACROS UTILES
%
%
%%%%%%%%%%%%%%%%%%%%%%%%%%%%%%%%%%%%%%%%%%%%%%%%%

%%%                      Macros
 
%-------------
%--- vectors  ---
%-------------

 
\newcommand\vectorl[1]{{\mathbf#1}}
\newcommand\tu[1]{\vectorl{#1}}

\newcommand\vp{\vectorl{p}}
\newcommand\va{\vectorl{a}}
\newcommand\vb{\vectorl{b}}
\newcommand\vc{\vectorl{c}}
\newcommand\vf{\vectorl{f}}
\newcommand\vn{\vectorl{n}}
\newcommand\vx{\vectorl{x}}
\newcommand\vX{\vectorl{X}}
\newcommand\vy{\vectorl{y}}
\newcommand\vz{\vectorl{z}}
\newcommand\ve{\vectorl{e}}
\newcommand\vv{\vectorl{v}}
\newcommand\vr{\vectorl{r}}
\newcommand\vu{\vectorl{u}}
\newcommand\vA{\vectorl{A}}
\newcommand\vB{\vectorl{B}}
\newcommand\vC{\vectorl{C}}
\newcommand\vD{\vectorl{D}}

%-------------
%--- overline ---
%-------------


\newcommand\ox{\overline{x}}
\newcommand\oy{\overline{y}}
\newcommand\oc{\overline{c}}
\newcommand\oa{\overline{a}}
\newcommand\ow{\overline{w}}
\newcommand\od{\overline{d}}
\newcommand\ob{\overline{b}}
\newcommand\oP{\overline{P}}
\newcommand\oee{\overline{e}}
\newcommand\ot{\overline{t}}
\newcommand\ou{\overline{u}}
\newcommand\ov{\overline{v}}
\newcommand\oz{\overline{z}}
\newcommand\am{\overline{a}}
\newcommand\om{\overline{m}}
\newcommand\oj{\overline{j}}

%----------------
%--- ensembles ---
%----------------

\newcommand\Q{\mathbb{Q}}
\newcommand\R{\mathbb{R}}
\newcommand\Z{\mathbb{Z}}
\providecommand\C{\mathbb{C}}
\newcommand\N{\mathbb{N}}
\newcommand\K{\mathbb{K}}
\newcommand\F{\mathbb{F}}

\newcommand\dyadic{\mathbb{D}}
\newcommand\Zd{{\Z_2}}
\newcommand\Zp{{\Z_p}}
%
\newcommand\RP{\mathbb{R}^{>0}}
\newcommand\QP{\mathbb{Q}^{>0}}


%----------------
%--- lettres cal ---
%----------------

\newcommand\mK{\mathcal{K}}
\newcommand\mD{\mathcal{D}}
\newcommand\mA{\mathcal{A}}
\newcommand\mL{\mathcal{L}}
\newcommand\mX{\mathcal{X}}
\newcommand\mG{\mathcal{G}}
\newcommand\mE{\mathcal{E}}
\newcommand\mH{\mathcal{H}}
\newcommand\mR{\mathcal{R}}
\newcommand\mF{\mathcal{F}}
\newcommand\mJ{\mathcal{J}}
\newcommand\mO{\mathcal{O}}
\newcommand\mP{\mathcal{P}}
\newcommand\mN{\mathcal{N}}
\newcommand\mS{\mathcal{S}}
\newcommand\mM{\mathcal{M}}
\newcommand\mC{\mathcal{C}}
\newcommand\mV{\mathcal{V}}
\newcommand\mI{\mathcal{I}}
\newcommand\mT{\mathcal{T}}

\newcommand\fC{\mathcal{C}}

%----------------
%--- lettres frak ---
%----------------


\newcommand\kM{\mathfrak{M}}
\newcommand\kN{\mathfrak{N}}
\newcommand\kU{\mathfrak{U}}
\newcommand\kg{\mathfrak{g}}



%--------------------------
%--- Façon noter noms classes ---
%---------------------------

\newcommand{\cp}[1]{\operatorname{#1}}


%--------------------------
%--- Calculabilité  ---
%---------------------------

% -- classes
\LANGUE{
\newcommand\RE{\cp{CE}}
}
{\newcommand\RE{\cp{RE}}
}

\newcommand\Rec{\cp{D}}
\newcommand\PR{\cp{PR}}
\newcommand\Gregn{\mE_n}


%--------------------------
%--- Réductions  ---
%---------------------------

\newcommand\lem{\le_{m}}
\newcommand\lelog{\le_{L}}
\newcommand\equivlog{\equiv_{L}}



%--------------------------
%--- Complexité  ---
%---------------------------

% COMPLEXITE

% -- P
\newcommand{\Ptime}{\cp{P}}      % P
\newcommand\PTIME{\Ptime}
\newcommand\PTIMEs[1]{{\PTIME}_{#1}}
\newcommand\PTIMEsp[1]{{\PTIME}_{#1}^0}
\renewcommand\P{\PTIME}
\newcommand{\FPtime}{\cp{FPTIME}}

% -- NP
\newcommand\NP{\cp{NP}}
\newcommand{\NPtime}{\NP}
\newcommand\NPs[1]{\cp{NP}_{#1}}
\newcommand{\FNP}{\cp{FNP}}

\newcommand\NDP{\cp{NDP}}
\newcommand\coNDP{\cp{coNDP}}
\newcommand\NDPs[1]{\cp{NDP}_{#1}}

% -- coNP
\newcommand\coNP{\cp{coNP}}

% -- PSPACE
\newcommand\PSPACE{\cp{ PSPACE}}
\newcommand\NPSPACE{\cp{ NPSPACE}}
\newcommand{\Pspace}{\PSPACE}
\newcommand{\FPspace}{\cp{FPSPACE}}


% -- LOGSPACE
\newcommand\LSPACE{\cp{LOGSPACE}}
\newcommand\NLSPACE{\cp{NLOGSPACE}}
\newcommand\coNLSPACE{\cp{coNLOGSPACE}}
\newcommand\coNL{\cp{coNL}}

% -- Autres
\newcommand\EXPTIME{\cp{EXPTIME}}
\newcommand\NEXPTIME{\cp{NEXPTIME}}
\newcommand\EXPSPACE{\cp{EXPSPACE}}
\newcommand\PH{\cp{PH}}
\newcommand\NC{\cp{NC}}
\newcommand\AC{\cp{AC}}
\newcommand\PPAD{\cp{\mathbf{PPAD}}}
\newcommand\PLS{\cp{\mathbf{PLS}}}

% -- Reverse math
\newcommand{\WKL}{\ensuremath{\docheck{\mathsf{WKL}_0}}}
\newcommand{\ACA}{\ensuremath{\docheck{\mathsf{ACA}_0}}}
\newcommand{\ATR}{\ensuremath{\docheck{\mathsf{ATR}_0}}}
\newcommand{\RCAO}{\ensuremath{\docheck{\mathsf{RCA}_0}}}
\newcommand{\PiCA}{\ensuremath{\docheck{\Pi^1_1\!\text{-}\mathsf{CA}_0}}}


% -- circuits
\newcommand\CIRCUIT[2]{\cp{CIRCUIT(#1,#2)}}
\newcommand\UCIRCUIT[2]{\cp{UCIRCUIT(#1,#2)}}

% -- Classes proba
\newcommand\PP{\cp{PP}}
\newcommand\RPP{\cp{RP}}
\newcommand\ZPP{\cp{ZPP}}
\newcommand\coRP{\cp{coRP}}
\newcommand\BPP{{\cp{BPP}}}

% -- IP
\newcommand\IP{\cp{IP}}
\newcommand\PCP{\cp{PCP}}

% -- non-uniforme
\newcommand\nuP{\mathbb{P}}
\newcommand\nuPs[1]{\mathbb{P}_{#1}}
\newcommand\nuPsp[1]{\mathbb{P}_{#1}^0}
\newcommand\nuNP{\mathbb{N}\mathbb{P}}
%\newcommand\poly{/{poly}}
\newcommand\Ppoly{\PTIME_{\poly}}
\newcommand\NPpoly{\NP_{\poly}}

% -- conditions d'uniformité
\newcommand\Luniforme{\cp{L}}
\newcommand\BCuniforme{\cp{BC}}


% Classes de circuits
\newcommand\TC{\cp{TC}}
\newcommand\LFP{\cp{LFP}}
\newcommand\FAC{\cp{FAC}}
\newcommand\FACC{\cp{FACC}}
\newcommand\FTC{\cp{FTC}}
\newcommand\FNC{\cp{FNC}}

% Logiques temorelles
\newcommand\CTL{\cp{CTL}}
\newcommand\LTL{\cp{LTL}}

% -- mesure temps/espace/taille

\newcommand\DTIME{\cp{DTIME}}
\newcommand\TIME[1]{\cp{TIME(#1)}}
\newcommand\TIMEs[2]{\cp{TIME_{#2}(#1)}}
\newcommand\TIMEsp[2]{\cp{TIME^0_{#2}(#1)}}
\newcommand\NTIME[1]{\cp{NTIME(#1)}}
\newcommand\NTIMEs[2]{\cp{NTIME_{#2}(#1)}}
\newcommand\NTIMEsp[2]{\cp{NTIME_{#2}^0(#1)}}
\newcommand\DSPACE{\cp{DSPACE}}
\newcommand\SPACE[1]{\cp{SPACE(#1)}}
\newcommand\SPACEs[2]{\cp{SPACE_{#2}(#1)}}
\newcommand\NSPACE[1]{\cp{NSPACE(#1)}}
\newcommand\NSPACEs[2]{\cp{NSPACE_{#2}(#1)}}
\newcommand\coNSPACE[1]{\cp{coNSPACE(#1)}}
\newcommand\TIMEON[2]{\cp{TIME(#1,#2)}}
\newcommand\SPACEON[2]{\cp{SPACE(#1,#2)}}
\newcommand\SIZE[1]{\cp{ size(#1)}}
\newcommand\ATIME{\cp{ATIME}}
\newcommand\ASPACE{\cp{ASPACE}}


%% autre:
\newcommand\LH{\cp{LH}}
\newcommand\FO{\cp{FO}}
\newcommand\SO{\cp{SO}}
\newcommand\ESOhorn{\cp{ESO_{horn}}}
\newcommand\SOhorn{\cp{SO_{horn}}}
\newcommand\ACzero{\cp{AC_0}}
\newcommand{\LINSPACE}{\cp{LINSPACE}}
\newcommand{\mFLINSPACE}{\mF_{\cp{LINSPACE}}}
\newcommand{\LOGSPACE}{\cp{LOGSPACE}}
\newcommand\ALOGTIME{\cp{ALOGTIME}}
\newcommand\RF{\cp{RF}}


\newcommand\BOOL[1]{\cp{{BOOLEAN(#1)}}}
\newcommand\DET{\cp{ DET}}

% % -- dirty


\newcommand{\cPspace}{\cp{\sharp PSPACE}}
\newcommand{\cAPtime}{\cp{\sharp APTIME}}
%\newcommand{\FPspace}{\cp{FPSPACE}}
%\newcommand{\FPspaceN}{\cp{FPSPACE}^{+}}
%\newcommand{\FPspaceN}{\cp{FPSPACE}}



%-------------------------
%--- Problèmes de décision  ---
%--------------------------

\newcommand\decisionname[1]{\cp{#1}}


\newcommand\Luniv{\decisionname{L_{univ}}}
\newcommand\HALTINGPROBLEM{\decisionname{HALTING-PROBLEM}}
\newcommand\SAT{\decisionname{SAT}}
\newcommand\MAXtSAT{\decisionname{MAX3SAT}}
\newcommand\REACH{\decisionname{REACH}}
\newcommand\CIRCUITVALUE{\decisionname{CIRCUITVALUE}}
\newcommand\MONOTONECIRCUITVALUE{\decisionname{MONOTONECIRCUITVALUE}}
\newcommand\THEOREMS{\decisionname{THEOREMS}}
\newcommand\QCIRCUITSAT{\decisionname{QCIRCUITSAT}}
\newcommand\QSAT{\decisionname{QSAT}}
\newcommand\QBF{\decisionname{QBF}}
\newcommand\SATVALUE{\decisionname{SATVALUE}}
\newcommand\SUCCINTSAT{\decisionname{SUCCINTSAT}}
\newcommand\SUCCINTSATVALUE{\decisionname{SUCCINTSATVALUE}}
\newcommand\GEOGRAPHY{\decisionname{GEOGRAPHY}}
\newcommand\PARITY{\decisionname{PARITY}}
\newcommand\MAJ{\decisionname{MAJ}}
\newcommand\CVP{\CIRCUITVALUE}
\newcommand\BOOLEANPROGRAMVALUE{BOOLEAN~PROGRAM~VALUE}

\newcommand\CLIQUE{\cp{CLIQUE}}
\LANGUEinv{
\newcommand\RECOUVREMENTDESOMMETS{\cp{RECOUVREMENT~DE~SOMMETS}}
\newcommand\RS{\RECOUVREMENTDESOMMETS}
}{
\newcommand\RECOUVREMENTDESOMMETSeng{\cp{VERTEX~COVER}}
\newcommand\RSeng{\RECOUVREMENTDESOMMETSeng}
}

\LANGUEinv{\newcommand\COUPUREMAXIMALE{\cp{COUPURE~MAXIMALE}}}
{\newcommand\COUPUREMAXIMALEeng{\cp{MAXIMAL~CUT}}}
\newcommand\CIRCUITHAMILTONIEN{\CHAMILTONIEN}
\LANGUEinv{\newcommand\VOYAGEURDECOMMERCE{\VCOMMERCE}}
{\newcommand\VOYAGEURDECOMMERCEeng{\VCOMMERCEeng}}

\LANGUEinv{\newcommand\CIRCUITLEPLUSLONG{\cp{CIRCUIT~LE~PLUS~LONG}}}
{\newcommand\CIRCUITLEPLUSLONGeng{\cp{LONGEST~CIRCUIT}}}
\LANGUEinv{\newcommand\SOMMEDESOUSENSEMBLE{\SUBSETSUM}}{
\newcommand\SOMMEDESOUSENSEMBLEeng{\SUBSETSUMeng}}
\LANGUEinv{
\newcommand\SACADOS{\cp{SAC~A~DOS}}}
{\newcommand\SACADOSeng{\cp{KNAPSACK}}}

%français
\LANGUEinv{
\newcommand\CHAMILTONIEN{\cp{CIRCUIT~HAMILTONIEN}}}
{\newcommand\CHAMILTONIENeng{\cp{HAMILTONIAN~CIRCUIT}}}
\LANGUEinv{
\newcommand\VCOMMERCE{\cp{VOYAGEUR~DE~COMMERCE}}}
{\newcommand\VCOMMERCEeng{\cp{TRAVELING~SALESMAN}}}
\LANGUEinv{
\newcommand\kCOLORABILITE{\cp{k-COLORABILITE}}}
{\newcommand\kCOLORABILITEeng{\cp{k-COLORABILITY}}}
\LANGUEinv{
\newcommand\COLORIAGE{\cp{BON~COLORIAGE}}}
{\newcommand\COLORIAGEeng{\cp{GOOD~COLORING}}}
\LANGUEinv{
\newcommand\tCOLORABILITE{\cp{3-COLORABILITE}}}
{\newcommand\tCOLORABILITEeng{\cp{3-COLORABILITY}}}
\LANGUEinv{
\newcommand\SUBSETSUM{\cp{SOMME~DE~SOUS~ENSEMBLE}}}
{\newcommand\SUBSETSUMeng{\cp{SUBSET~SUM}}}
\newcommand\PARTITION{\cp{PARTITION}}
\LANGUEinv{
\newcommand\NOMBREPREMIER{\decisionname{NOMBRE~ PREMIER}}}
{\newcommand\NOMBREPREMIEReng{\decisionname{PRIME~NUMBER}}}
\LANGUEinv{
\newcommand\CODAGE{\decisionname{CODAGE}}}
{\newcommand\CODAGEeng{\decisionname{ENCODING}}}
\newcommand\kSAT{\cp{k-SAT}}
\newcommand\tSAT{\cp{3-SAT}}
\newcommand\NAESAT{\cp{NAESAT}}
\newcommand\NAEtSAT{\cp{NAE3SAT}}
\newcommand\NAEqSAT{\cp{NAE4SAT}}
\newcommand\STABLE{\cp{\LANGUE{INDEPENDANT~SET}{STABLE}}}

%\newcommand\CM{\cp{COUPURE MAXIMALE}}
\LANGUEinv{
\newcommand\POSTpb{\cp{Problème~de~ la~correspondance~de~Post}}}
{\newcommand\POSTpbeng{\cp{Post~correspondence~problem}}}

\LANGUEinv{
\newcommand\PARITE{\cp{PARITE}}}
{\newcommand\PARITEeng{\cp{PARITY}}}


%-------------------------
%--- Modèles parallèles  ---
%--------------------------


% PARALLELISME
\newcommand\RAM{\cp{RAM}}
%
\newcommand\PRAM{\cp{PRAM}}
\newcommand\CRCW{\cp{CRCW}}
\newcommand\CREW{\cp{CREW}}
\newcommand\EREW{\cp{EREW}}


%------------------------
%--- codages  ---
%------------------------


\newcommand\codage[1]{\langle #1 \rangle}
\newcommand\tuple[1]{\codage{#1}}
\newcommand\paire[2]{\codage{#1,#2}}
\newcommand\triplet[3]{\codage{#1,#2,#3}}
\newcommand\quadruplet[4]{\codage{#1,#2,#3,#4}}
\newcommand\cinquplet[5]{\codage{#1,#2,#3,#4,#5}}


%------------------------
%--- descriptive set theory  ---
%------------------------

\newcommand\baire{\mathcal{N}}
\newcommand\peano{\overline{\mathcal{N}}} %% A MOI
\newcommand\cantor{\mathcal{C}}
%
\newcommand\bSigma{\mathbf{\Sigma}}
\newcommand\bPi{\mathbf{\Pi}}
\newcommand\bDelta{\mathbf{\Delta}}
%


%%%%%% ABOUT NOMS QUI CHANGES

\newcommand\WF{\cp{WF}}
\newcommand\WO{\WF}
\newcommand\LO{\cp{LO}}
\newcommand\DLO{\cp{DLO}}
% 
\newcommand\rk{\cp{rank}}


%% Moins propre:

\newcommand\I{\mathbb{I}}
\renewcommand\H{\mathbb{H}}
%
\newcommand\leKB{<_{KB}}
%
\newcommand\finiomega{{<\omega}}
\newcommand\compl{^\frown}
\newcommand\prefix{ ~ prefix~  } 
\newcommand\X{{X}}
\newcommand\Y{\mathcal{Y}}
\newcommand\subsetstrict{\subset}

\renewcommand\complement[1]{#1^{c}}


%------------------------
%--- schémas de récursion ---
%------------------------

%%                      Source papier MFCS 2019


% COMPLEXITY

%% MFCS differe de Clote:
%% Version à taille restreinte

\newcommand\co{\cp{co}}

\newcommand\mFPTIME{\mF_{PTIME}}
\newcommand{\mFPSPACE}{\mF_{SPACE}}
\newcommand{\mFPH}{\mF_{PH}}

\newcommand\SigmaP{\Sigma^P}
\newcommand\DeltaP{\Delta^{P}}
\newcommand\PiP{\Pi^{P}}
\newcommand\coSigmaP{\co\SigmaP}
\newcommand\coPiP{\mF{\co\PiP}}

%% autres classes:


%% Schemas pour complexité implicite

\newcommand\COMP{\cp{ COMP}}
\newcommand\CRN{\cp{ CRN}}
\newcommand\BRN{\cp{ BRN}}
\newcommand\SBRN{\cp{ SBRN}}
\newcommand\BR{\cp{ BR}}
\newcommand\kBR{k-\cp{ BR}}
\newcommand \WBPR{\cp{ WBPR}}
\newcommand \BMIN{\cp{ BMIN}}
\newcommand \BSUM{\cp{ BSUM}}
\newcommand\SBSUM{\cp{ SBSUM}}
\newcommand \BPROD{\cp{ BPROD}}
\newcommand \SBPROD{\cp{ SBPROD}}
\newcommand \SCOMP{\cp{ SCOMP}}
\newcommand \SR{\cp{ SR}}
\newcommand \SRN{\cp{ SRN}}

%% devrait etre défini
\newcommand\WBRN{\cp{ WBRN}}



%-- godel
\newcommand\INTDIV{\cp{INTDIV}}
\newcommand\DIV{\cp{DIV}}
\newcommand\EVEN{\cp{EVEN}}
\newcommand\ODD{\cp{ODD}}
\newcommand\PRIME{\cp{PRIME}}
\newcommand\POWER{\cp{POWER}}
\newcommand\BIT{\cp{BIT}}
\newcommand\negHP{\overline{\cp{HP}}}
\newcommand\HP{\cp{HP}}
\newcommand\negLuniv{\overline{\Luniv}}
\newcommand\VALCOMP{\cp{VALCOMP}}
\newcommand\mWINDOW{\cp{LEGAL}}
\newcommand\mmWINDOW{\cp{WINDOW}}
\newcommand\LENGTH{\cp{LENGTH}}
\newcommand\DIGIT{\cp{DIGIT}}
\newcommand\MATCH{\cp{MATCH}}
\newcommand\MOVE{\cp{MOVE}}
\newcommand\START{\cp{START}}
\newcommand\CELL{\cp{CELL}}
\newcommand\HALT{\cp{HALT}}
\newcommand\SUBST{\cp{SUBST}}
\newcommand\PROOF{\cp{PROOF}}
\newcommand\PROVABLE{\cp{PROVABLE}}
\newcommand\CONSIST{\cp{CONSIST}}



%---------------------
%--- calculus (AP) ---
%---------------------

% \jacobian{f}{x}
\newcommand{\jacobian}[1]{J_{#1}}
%\newcommand{\grad}[1]{\overrightarrow{\operatorname{grad}}~{#1}}
\newcommand{\grad}[1]{\nabla{#1}}
\newcommand{\scalarprod}[2]{{#1}\cdot{#2}}
\newcommand{\bigO}[1]{\mathcal{O}\left(#1\right)}
\newcommand\smallO{o}
\newcommand{\softO}[1]{\tilde{\mathcal{O}}\left(#1\right)}
\newcommand{\taylor}[3]{{T_{#1}^{#2}#3}}
\newcommand{\transpose}[1]{{#1}^T}
\newcommand{\pastsup}[2]{{\sup}_{#1}#2}

\DeclareMathOperator\arctanh{arctanh}



%---------------
%--- norms (AP) ---
%---------------


\newcommand{\inorm}[2]{\left\lVert{#1}\right\rVert_{#2}}
\newcommand{\twonorm}[1]{\inorm{#1}{2}}
\newcommand{\infnorm}[1]{\inorm{#1}{}}
\newcommand{\normsup}[1]{\infnorm{#1}{}}
%
\newcommand\hausdorff{d_H}
%
\newcommand\norm[1]{\infnorm{#1}}


%----------------
%--- topology ---
%----------------

% \openball{pt}{radius}
 \newcommand{\openball}[2]{B_{#2}(#1)}
\newcommand{\closedball}[2]{\ovl{B}_{#2}(#1)}
\newcommand{\closure}[1]{\ovl{#1}}
% \newcommand{\interior}[1]{\mathring{#1}}
% \newcommand{\border}[1]{\partial{#1}}




%-----------------
%--- functions (AP) ---
%-----------------
 
\newcommand{\lambdafun}[2]{(#1)\mapsto{#2}}
\newcommand{\lambdafunex}[2]{#1\mapsto{#2}}
\newcommand{\argmin}{\operatornamewithlimits{argmin}}
\newcommand{\argmax}{\operatornamewithlimits{argmax}}


 
%-----------------------
%--- Machines de Turing  ---
%-----------------------

\newcommand\blanc{\vectorl{B}}



%-----------------
%--- Logique ---
%-----------------
 
\newcommand\sem[1]{{[\semspace[}#1{]\semspace]}}
\newcommand\sems[1]{{[\semspace[}#1{]\semspace]}}
\newcommand\semss[1]{\sem{#1}_{\sigma'}}
\newcommand\semspace{\kern -.25ex}
\newcommand\true{true}
\newcommand\false{false}
%extension élémentaire:
\newcommand\elem{\preceq}
\newcommand\godel[1]{\lceil #1 \rceil}

%-----------------
%--- Model theory ---
%-----------------

\newcommand\hull[1]{\langle #1 \rangle}

%-----------------
%--- Ordinaux ---
%-----------------

\newcommand\omegaunCK{\omega_1^{CK}}
\newcommand\omegaunck{\omegaunCK}

\DeclareMathOperator{\cof}{cof}

%-----------------
%--- Surréels (QG) ---
%-----------------

\newcommand{\On}{\textbf{On}}
\newcommand{\Nobf}{\textbf{No}}
\newcommand{\Nobb}{\ensuremath{\mathbbmss{No}}}
\newcommand{\Nolt}[1]{\ensuremath{\Nobf_{< #1}}}
\newcommand{\Noleq}[1]{\ensuremath{\Nobf_{\leq #1}}}
\newcommand{\Noltbb}[1]{\ensuremath{\Nobb_{< #1}}}
\providecommand\No{\Nobf}
\newcommand\Noalpha{\Nolt{\alpha}}
\newcommand\Nolambda{\Nolt{\lambda}}

\DeclareMathOperator{\length}{length}
\DeclareMathOperator{\supp}{supp} % operateur support


%----------------------
%---Maths  de base (AP)  ---
%----------------------

\DeclareMathOperator{\dom}{dom} % operateur support

\DeclareMathOperator{\size}{size}

\newcommand{\powerset}[1]{\mathcal{P}(#1)}
\newcommand{\card}[1]{{\##1}}
\newcommand\priv{ - }
\newcommand\prive{-}
\renewcommand{\restriction}{\mathord{\upharpoonright}}
\newcommand\restrict[1]{_{\restriction #1}}

% \providecommand\mod{}
%% UNDEFINE:
\let\mod\relax
\let\div\relax
\DeclareMathOperator{\mod}{mod} 
\DeclareMathOperator{\div}{div} 

% partieentier
\newcommand\entieres[1]{\lceil #1 \rceil}
\newcommand\entierei[1]{\lfloor #1 \rfloor}


%----------------------
%---Symbole danger ---
%----------------------


%\providecommand*{\TakeFourierOrnament}[1]{{%
%\fontencoding{U}\fontfamily{futs}\selectfont\char#1}}
%\providecommand*{\danger}{{\Huge \TakeFourierOrnament{66}}}



%%%%%%%%%%%%%%%%%%%%%%%%%%%%%%%%%%%%%%%%%%%%%%%%%
%
%
%                      ASMeries
%
%
%%%%%%%%%%%%%%%%%%%%%%%%%%%%%%%%%%%%%%%%%%%%%%%%%



% CIRCUITS
%-- particuliers
\newcommand\REPLACE{\cp{REPLACE}}
\newcommand\EQUAL{\cp{EQUAL}}
\newcommand\EVALUE{\cp{EVALUE}}
\newcommand\TVALUE{\cp{TVALUE}}
\newcommand\UPDATE{\cp{UPDATE}}
\newcommand\ASSIGN{\cp{ASSIGN}}
\newcommand\PAR{\cp{PAR}}
\newcommand\DO{\cp{DO}}


% ASM
\newcommand\emplacement[2]{(#1,#2)}
\newcommand\contenu[2]{\sem{(#1,#2)}}
\newcommand\affecte[3]{(#1,#2)\leftarrow#3}
\newcommand\ctl{ctl\_state}


%%%%%%%%%%%%%%%%%%%%%%%%%%%%%%%%%%%%%%%%%%%%%%%%%
%
%
%                      GPACqueries
%
%
%%%%%%%%%%%%%%%%%%%%%%%%%%%%%%%%%%%%%%%%%%%%%%%%%



%------------------------
%--- generable zoo (GPAC) ---
%------------------------

\newcommand{\myop}[1]{\operatorname{#1}}
\newcommand{\abs}{\myop{abs}}
\newcommand{\sg}{\myop{sg}}
\newcommand{\ip}[1]{\myop{ip}_{#1}}
\newcommand{\rnd}{\myop{rnd}}
\newcommand{\lil}{\myop{lil}}
\newcommand{\lxh}{\myop{lxh}}
\newcommand{\hxl}{\myop{hxl}}
\newcommand{\plil}{\myop{plil}}
\newcommand{\reach}{\myop{reach}}
%\newcommand{\mreach}{\myop{mreach}}
\newcommand{\watchdog}{\myop{watchdog}}
\newcommand{\monitor}{\myop{monitor}}
%\newcommand{\norm}{\myop{norm}}
\newcommand{\mx}{\myop{mx}}
\newcommand{\mn}{\myop{mn}}
\newcommand{\sample}{\myop{sample}}
\newcommand{\select}{\myop{select}}
\newcommand{\ssine}[1]{\sin_{#1}}
\newcommand{\clamp}{\myop{clamp}}

\newcommand{\fnsg}[3]{tanh((#1)*(#2)*(#3))}
\newcommand{\fnipone}[3]{(1+\fnsg{(#1)-1}{#2}{#3})/2.}
\newcommand{\fnipinf}[1]{#1-sin(2*pi*(#1))/2./pi}%
\newcommand{\fnlil}[5]{%
\fnipone{(#3)-(#1)+1-((#2)-(#1))/8.}{#4+log(1+4./(#2-#1)*(#5)*(#5))+log(1+4)}{8./(#2-#1)}%
*\fnipone{#2-#3+1-(#2-#1)/8.}{#4+log(1+4/(#2-#1)*(#5)*(#5))+log(1+4)}{8./(#2-#1)}%
*4./(#2-#1)*(#5)}
\newcommand{\fnlxh}[5]{\fnipone{(#3)-(#1+#2)/2.+1}{#4+log(1+(#5)*(#5))}{2./(#2-#1)}*(#5)}
\newcommand{\fnhxl}[5]{\fnipone{(#1+#2)/2.-(#3)+1}{#4+log(1+(#5)*(#5))}{2./(#2-#1)}*(#5)}
\newcommand{\fnplilf}[4]{sin(((#4)-((#1)+(#2))/2.+(#3)/4.)*2.*pi/(#3))}
\newcommand{\fnplil}[6]{(#6)*\fnlxh{\fnplilf{#1}{#2}{#3}{#1}}%
{\fnplilf{#1}{#2}{#3}{#1+((#2)-(#1))/4.}}%
{\fnplilf{#1}{#2}{#3}{#4}}%
{#5+2.+log(1+(#6)*(#6))}%
{1./4.+2./((#2)-(#1))}}
\newcommand{\fnssine}[2]{sin(#2)*(1-cos(#2)/cos(#1))}



%--------------------
%--- complexity GPAC ---
%--------------------


%\ifthenelse{\equal{#1}{}}{Empty.}{Nonempty.}
\newcommand{\myclass}[1]{\operatorname{#1}}
% \newcommand{\gcomp}[2]{\ensuremath{\myclass{GCOMP}(#1,#2)}}
% \newcommand{\ggenerable}[1]{\textcolor{red}{\ensuremath{#1-}\text{generable}}}
 \newcommand{\gval}[2][]{\ensuremath{\myclass{GVAL}_{#1}\ifthenelse{\equal{#2}{}}{}{[#2]}}}
 \newcommand{\gpval}[1][]{\ensuremath{\myclass{GPVAL}_{#1}}}
 \newcommand{\gc}[3][]{\ensuremath{\myclass{ATSC}(#2,#3)}}
 \newcommand{\gpc}[1][]{\ensuremath{\myclass{ATSP}}}
\newcommand{\cgc}[1][]{\ensuremath{\myclass{ATS}}}
\newcommand{\gexpc}[1][]{\ensuremath{\myclass{AEXP}}}
\newcommand{\gwc}[2]{\ensuremath{\myclass{AW}(#1,#2)}}
\newcommand{\gpwc}{\ensuremath{\myclass{AWP}}}
\newcommand{\cgwc}{\ensuremath{\myclass{AW}}}
\newcommand{\grc}[3]{\ensuremath{\myclass{AR}(#1,#2,#3)}}
\newcommand{\gprc}{\ensuremath{\myclass{ARP}}}
\newcommand{\gsc}[3]{\ensuremath{\myclass{AS}(#1,#2,#3)}}
\newcommand{\gpsc}{\ensuremath{\myclass{ASP}}}
\newcommand{\goc}[3]{\ensuremath{\myclass{AO}(#1,#2,#3)}}
\newcommand{\gpoc}{\ensuremath{\myclass{AOP}}}
\newcommand{\cgoc}{\ensuremath{\myclass{AO}}}
\newcommand{\guc}[4]{\ensuremath{\myclass{AX}(#1,#2,#3,#4)}}
\newcommand{\gpuc}{\ensuremath{\myclass{AXP}}}
\newcommand{\glc}[2][]{\ensuremath{\myclass{AL}(#2)}}
\newcommand{\gplc}[1][]{\ensuremath{\myclass{ALP}}}
\newcommand{\cglc}[1][]{\ensuremath{\myclass{AL}}}

%\newcommand{\intinterv}[2]{\{#1,\ldots,#2\}}
\newcommand{\intinterv}[2]{\llbracket#1,#2\rrbracket}

\newcommand{\Rgen}{\R_{G}}
\newcommand{\Rpoly}{\R_{P}}

%%%%%%%%%%%%%%%%%%%%%%%%%%%%%%%%%%%%%%%%%%%%%%%%%
%
%
%                      Mes codages obscures
%
%
%%%%%%%%%%%%%%%%%%%%%%%%%%%%%%%%%%%%%%%%%%%%%%%%%




%%% Codages

% \newcommand\gammaoc{\gamma_{[0,1]}}
% \newcommand\gammaonc{\gamma_{\N}}
\newcommand\gammaTMc{\gamma^{TM}_{[0,1]}}
\newcommand\gammaTMcatan{\gamma^{TM}_{(0,1)^2}}
\newcommand\gammaTMnc{\gamma^{TM}_{\N^3}}
\newcommand\gammaTMncd{\gamma^{TM}_{\N^2}}
% \newcommand\gammaTMe{\gamma^{TM}_{*}}
% \newcommand\gammas{\gamma_{sab}}
% \newcommand\gammaord{\gamma_{cantor}}


\newcommand\enccompact{\nu_{X}}
 \newcommand\encinfinite{\nu_\N}
 \newcommand\enc{\nu}
 
%%%%%%%%%%%%%%%%%%%%%%%%%%%%%%%%%%%%%%%%%%%%%%%%%
%
%
%                      DESSINS ET IMAGES
%
%
%%%%%%%%%%%%%%%%%%%%%%%%%%%%%%%%%%%%%%%%%%%%%%%%%


\newcommand\IMAGE[2]{ \begin{center} 
    \includegraphics[#1]{#2}
  \end{center}
  }


  
%%%% POUR FAIRE DESSIN.
\newenvironment{dessinpp}[1]{
\begin{tikzpicture}[baseline=(current bounding
      box.west),#1]
}
{\end{tikzpicture}}




\newenvironment{dessinp}[1]{
\begin{center}\begin{tikzpicture}[baseline=(current bounding
      box.north),#1]
}
{\end{tikzpicture}
\end{center}}


\newenvironment{dessin}{
\begin{center}\begin{tikzpicture}[baseline=(current bounding
      box.north)]
}
{\end{tikzpicture}
\end{center}}



\newenvironment{dessind}{
\begin{tikzpicture}[baseline=(current bounding
      box.north)]
}
{\end{tikzpicture}
}



%%%%%%%%%%%%%%%%%%%%%%%%%%%%%%%%%%%%%%%%%%%%%%%%%
%
%
%                      POUR DESSINS GPAC
%
%
%%%%%%%%%%%%%%%%%%%%%%%%%%%%%%%%%%%%%%%%%%%%%%%%%


%%%% MACROS DE BASE.


\newcommand\constant[2]
{
node (#1)  [shape=circle,draw] {#2};
\draw (#1.east) -- +(0.3,0) coordinate (#1-o) ; 
\draw (#1.west) -- +(-0.3,0) coordinate (#1-i) ; 
}

\newcommand\basicproduct[2]
{
node (#1) {#2};
\draw (#1) +(-0.4,0.4) -- +(0.4,0.4) -- +(0.8,0) -- +(0.4,-0.4) -- +(-0.4,-0.4) --
+(-0.4,0.4);
\draw (#1) ++(0.8,0) -- ++(+0.3,0) coordinate (#1-o) ;
\draw (#1) ++ (-0.4,0.2) -- +(-0.3,0) coordinate (#1-i1) ;
\draw (#1) ++(-0.4,-0.2) -- +(-0.3,0) coordinate (#1-i2) ;
}

\newcommand\product[1]{\basicproduct{#1}{$+\Pi$}}
\newcommand\negativeproduct[1]{\basicproduct{#1}{$-\Pi$}}


%%BEGIN INTERNAL
\newcommand\basicsummer[1]{
coordinate (#1) {};
\draw (#1) +(-0.5,0.6) -- +(0.5,0) -- +(-0.5,-0.6) -- +(-0.5,0.6);
\draw (#1) ++(0.5,0) -- +(+0.3,0) coordinate (#1-o) ;
}

\newcommand\basicintegrator[1]{
coordinate (#1) {};
\draw (#1) +(-0.5,0.6) -- +(0.5,0) -- +(-0.5,-0.6) -- +(-0.5,0.6);
\draw (#1) ++(0.5,0) -- +(+0.3,0) coordinate (#1-o) ;
\draw (#1) +(-0.5,-0.6) -| +(-0.8,0.6) -- +(-0.5,0.6) ;
 \draw (#1) ++(-0.65,0.6) -- +(0,+0.3) coordinate (#1-init) ;
}

\newcommand\splitfour[2]{
\draw (#2) ++ (#1,0.4) -- +(-0.3,0) coordinate (#2-i1) ;
\draw (#2) ++ (#1,0.15) -- +(-0.3,0) coordinate (#2-i2) ;
\draw (#2) ++ (#1,-0.15) -- +(-0.3,0) coordinate (#2-i3) ;
\draw (#2) ++ (#1,-0.4) -- +(-0.3,0) coordinate (#2-i4) ;
}

\newcommand\splitthree[2]{
\draw  (#2) ++(#1,0.3) -- +(-0.3,0) coordinate (#2-i1) ;
\draw  (#2) ++(#1,0) -- +(-0.3,0) coordinate (#2-i2) ;
\draw  (#2) ++(#1,-0.3) -- +(-0.3,0) coordinate (#2-i3) ;
}

\newcommand\splittwo[2]{
\draw  (#2) ++(#1,0.2) -- +(-0.3,0) coordinate (#2-i1) ;
\draw  (#2) ++(#1,-0.2) -- +(-0.3,0) coordinate (#2-i2) ;
}

\newcommand\splitone[2]{
\draw (#2) ++(#1,0) -- +(-0.3,0) coordinate (#2-i1) ;
}

%% END INTERNAL

\newcommand\summerfour[1]{
\basicsummer{#1}
\splitfour{-0.5}{#1}
}

\newcommand\summerthree[1]{
\basicsummer{#1}
\splitthree{-0.5}{#1}
}

\newcommand\summertwo[1]{
\basicsummer{#1}
\splittwo{-0.5}{#1}
}

\newcommand\summerone[1]{
\basicsummer{#1}
\splitone{-0.5}{#1}
}


\newcommand\integratorfour[1]{
\basicintegrator{#1}
\splitfour{-0.8}{#1}
<}


\newcommand\integratorthree[1]{
\basicintegrator{#1}
\splitthree{-0.8}{#1}
}


\newcommand\integratortwo[1]{
\basicintegrator{#1}
\splittwo{-0.8}{#1}
}


\newcommand\integratorone[1]{
\basicintegrator{#1}
\splitone{-0.8}{#1}
}

%%%%% FIN MACROS




%%%%%%%%%%%%%%%%%%%%%%%%%%%%%%%%%%%%%%%%%%%%%%%%%
%
%
%                      Trucs d'Amaury
%
%
%%%%%%%%%%%%%%%%%%%%%%%%%%%%%%%%%%%%%%%%%%%%%%%%%


%----------------
%--- ref (AP) ---
%----------------

% \newcommand{\defref}[1]{Definition~\ref{#1}}
\newcommand{\lemref}[1]{Lemma~\ref{#1}}
\newcommand{\thref}[1]{Theorem~\ref{#1}}
\newcommand{\amaurycorref}[1]{Corollary~\ref{#1}}
 \newcommand{\figref}[1]{Figure~\ref{#1}}
% \newcommand{\figrefmult}[1]{Figures~{#1}}
% \newcommand{\secref}[1]{Section~\ref{#1}}
% %\renewcommand{\algref}[1]{Algorithm~\ref{#1}}
 \newcommand{\propref}[1]{Proposition~\ref{#1}}
% \newcommand{\exref}[1]{Example~\ref{#1}}
\newcommand{\remref}[1]{Remark~\ref{#1}}


 %-------------------
%--- polynomials ---
%-------------------

\newcommand{\sigmap}[1]{{\Sigma{#1}}}
\newcommand{\degp}[1]{{\operatorname{deg}(#1)}}
\newcommand{\poly}{\operatorname{poly}}


 %----------------------
%--- misc functions ---
%----------------------

\newcommand{\sgn}[1]{\operatorname{sgn}(#1)}
\newcommand{\intp}{\operatorname{int}}
 \newcommand{\fracp}{\operatorname{frac}}
\newcommand{\fiter}[2]{#1^{[#2]}}
\newcommand{\frestrict}[2]{#1\restriction_{#2}}
\newcommand{\indicator}[1]{\mathds{1}_{#1}}
\newcommand{\idfun}{\operatorname{id}}
% \newcommand{\floor}[1]{\left\lfloor#1\right\rfloor}
 \newcommand{\ceil}[1]{\left\lceil#1\right\rceil}
\newcommand{\round}[1]{\left\lfloor#1\right\rceil}
\DeclareMathOperator\erf{erf}


%------------
%--- sets ---
%------------

\newcommand{\A}{\mathbb{A}}
%\newcommand{\D}{\mathbb{D}}
\newcommand{\Ns}{\mathbb{N}^*}
%\newcommand{\C}{\mathbb{C}}
%\newcommand{\K}{\mathbb{K}}
% \MAT{height}{[width]}{[field]}
\newcommand{\MAT}[3]{M_{#1\ifthenelse{\equal{#2}{}}{}{,#2}}\ifthenelse{\equal{#3}{}}{}{\left(#3\right)}}

%\newcommand{\Rcomp}{\R_{c}}
\newcommand{\Rp}{\R_{+}}
\newcommand{\Rs}{\R^{*}}
\newcommand{\Rps}{\R_{+}^{*}}

%\newcommand{\Analytic}{C^\omega}
%\newcommand{\PTIME}{\ensuremath{\operatorname{P}}}
%\newcommand{\EXPTIME}{\ensuremath{\operatorname{EXPTIME}}}
%\newcommand{\NP}{\ensuremath{\operatorname{NP}}}
%\newcommand{\NPTIME}{\NP}
\newcommand{\FP}{\ensuremath{\operatorname{FP}}}
%\newcommand{\PSPACE}{\ensuremath{\operatorname{PSPACE}}}
% \newcommand{\FPSPACE}{\ensuremath{\operatorname{FPSPACE}}}



%%%%%%%%%%%%%%%%%%%%%%%%%%%%%%%%%%%%%%%%%%%%%%%%%
%
%
%                      Trucs de Quentin
%
%
%%%%%%%%%%%%%%%%%%%%%%%%%%%%%%%%%%%%%%%%%%%%%%%%%


\newcommand{\fct}[5]{\ensuremath{#1 : \left\{\begin{array}{c c c}
#2 & \rightarrow & #3\\
#4 & \mapsto & #5
                                             \end{array}\right.}}

%Nouvelle fraction, plus jolie
\newcommand{\f}[2]{	
	\mathchoice%
		{\dfrac{#1}{#2}}
    	{\dfrac{#1}{#2}}
		{\frac{#1}{#2}}
		{\frac{#1}{#2}}
}


%%Ensembles et combinatoires

%Ensemble avec propriété à drotie \{x | blabla\}
\newcommand{\enstq}[2]{\left\{ \left.#1\ \vphantom{#2}\right|\ #2  \right\}}
%meme chose mais avec des crochets
\newcommand{\crotq}[2]{\left[ \left.#1\ \vphantom{#2}\right|\ #2  \right]}
%même chose mais avec des brackets
\newcommand{\bratq}[2]{\left\langle \left.#1\ \vphantom{#2}\right|\ #2  
  \right\rangle}

\newcommand{\intn}[2]{\ensuremath{
		\left\llbracket\, #1\ ;\ #2\,\right\rrbracket
              }}


%%%%%%%%%%%%%%%%%%%%%%%%%%%%%%%%%%%%%%%%%%%%%%%%%
%
%
%               ENVIRONNEMENTS THEOREMS & CO
%
%
%%%%%%%%%%%%%%%%%%%%%%%%%%%%%%%%%%%%%%%%%%%%%%%%%


%----------------
%--- pas dans lipics ---
%-----------------


\ifnotSLIDE{
% commente car bug Fevrier 2026: \ifnotTDEXAM{\spnewtheorem{solution}{Solution}{\bfseries}{\rmfamily}}
\spnewtheorem{openproblem}{Open Problem}{\bfseries}{\rmfamily}
\spnewtheorem{remarkolivier}{Commentaire d'Olivier: }{\bfseries}{\rmfamily}
\spnewtheorem{postulate}{Postulate}{\bfseries}{\rmfamily}
}

%----------------
%--- dans lipics ---
%-----------------
%
\newenvironment{exercice}[1]{\begin{exercise} \label{exo:#1-stt}}{\end{exercise}}
\newenvironment{exerciced}[1]{\begin{exercisee}\label{exo:#1-stt}}{\end{exercisee}}
\newenvironment{exercicec}[1]{\begin{exercise} \label{exo:#1-stt} (\LANGUE{solution on}{corrigé} page \pageref{exo:#1-cor})
}{\end{exercise}}
\newenvironment{exercicedc}[1]{\begin{exercisee} \label{exo:#1-stt} (\LANGUE{solution on}{corrigé} page \pageref{exo:#1-cor})
}{\end{exercisee}}

\newenvironment{slideproposition}{{\bf Proposition}}{}

\ifnotSLIDE{
\nolipics{
\nolncs{\spnewtheorem{theorem}{\LANGUE{Theorem}{Théorème}}{\bfseries}{\rmfamily}}
\nolncs{\spnewtheoremi{definition}{\LANGUE{Definition}{Définition}}{\bfseries}{\rmfamily}{theorem}}
\nolncs{\spnewtheoremi{lemma}{\LANGUE{Lemma}{Lemme}}{\bfseries}{\rmfamily}{theorem}}
\nolncs{\spnewtheoremi{corollary}{\LANGUE{Corollary}{Corollaire}}{\bfseries}{\rmfamily}{theorem}}
\nolncs{\spnewtheoremi{proposition}{Proposition}{\bfseries}{\rmfamily}{theorem}}
\nolncs{\spnewtheoremi{example}{\LANGUE{Example}{Exemple}}{\bfseries}{\rmfamily}{theorem}}
%\spnewtheorem{sketch}{{\bf \LANGUE{Sketch of proof}{Démonstration (principe)}}}{\bfseries}{\rmfamily}
\nolncs{\spnewtheoremi{remark}{\LANGUE{Remark}{Remarque}}{\bfseries}{\rmfamily}{theorem}}
\nolncs{\spnewtheorem{exercise}{\LANGUE{Exercise}{Exercice}}{\bfseries}{\rmfamily}}
\spnewtheorem{exercisee}{\LANGUE{*Exercise}{*Exercice}}{\bfseries}{\rmfamily}
\spnewtheorem{summary}{\LANGUE{Summary}{Résumé}}{\bfseries}{\rmfamily}
 %\spnewtheorem{exerciceenglish}{Exercise}{\bfseries}{\rmfamily}
 %
 \newenvironment{sketch}{{\bf \LANGUE{Sketch of proof}{Démonstration (principe)}}:}{\hfill $\Box$ }
\nolncs{\spnewtheorem{conjecture}{\LANGUE{Conjecture}{Conjecture}}{\bfseries}{\rmfamily}}
 %%
 \spnewtheorem{theoremaprouver}{\LANGUE{Theorem TO BE PROVED}{Théorème (A PROUVER)}}{\bfseries}{\rmfamily}
  \spnewtheorem{conjectureaprouver}{\LANGUE{Conjecture TO BE VERIFIED}{Conjecture (A VERIFIER)}}{\bfseries}{\rmfamily}
 }}
            
 \makeatletter
 \ifnotSLIDE{
  \ifnotTDEXAM{
\@ifundefined{proof}{
	\newenvironment{proof}{{\bf \LANGUE{Proof}{Démonstration}}:}{\hfill $\Box$ 
	}{}
	}}}
\makeatother

 
%%%%%%%%%%%%%%%%%%%%%%%%%%%%%%%%%%%%%%%%%%%%%%%%%
%
%
%                      Macros dont pas fier du tout
%
%
%%%%%%%%%%%%%%%%%%%%%%%%%%%%%%%%%%%%%%%%%%%%%%%%%



\newcommand\equations[1]{
\begin{array}{lll}
#1
\end{array}
}



\newcommand\systeme[1]{
\left\{
\begin{array}{lll} 
#1
\end{array}
\right.
}



%------------------
%--- shorthands (AP)---
%------------------
\newcommand{\mtt}[1]{\mathtt{#1}}
\newcommand{\ovl}[1]{\overline{#1}}
\newcommand{\panic}[1]{\textcolor{red}{#1}}
%\NewEnviron{epanic}{{\color{red}\BODY}}
\newcommand{\idest}{\emph{i.e.\thinspace}}


%------------------
%--- mes trucs  ---
%------------------


\newcommand\implique{\Rightarrow}
\newcommand\ssi{\Leftrightarrow}
\newcommand\ac[1]{\{#1\}}
\newcommand\zero{\vectorl{0}}
\newcommand\un{\vectorl{1}}

\newcommand\technique[1]{{\scriptsize #1}}


%% Preuve and Co

\newcommand\sensdirect{ Direction $\Rightarrow$. }
\newcommand\sensindirect{ Direction $\Leftarrow$. }




\newcommand\donnee{\mbox{<donnée>}}



% pour aller vite

\newcommand\gM{\mathfrak{M}}
\newcommand\gMM{(M,f_1,\cdots,f_u,r_1,\cdots,r_v)}
%
\newcommand\mygM{\mK}
\newcommand\mygMM{\left ( \K, \{op_i\}_{i\in I}, r_1 \ldots,
r_\ell, \zero,\un \right )}
\newcommand\myM{\K} 

\newcommand\gMMM{(M,f_1,\cdots,f_u,c_1,\cdots,c_k,r_1,\cdots,r_v)}
\newcommand\gMMMM{(f_1,\cdots,f_u,r_1,\cdots,r_v)}
\newcommand\gMME{(M,f_1,\cdots,f_u,f'_1,\cdots,f'_\ell,r_1,\cdots,r_v)}
\newcommand\gMMS{(f_1,\cdots,f_u,f'_1,\cdots,f'_\ell,r_1,\cdots,r_v)}

\newcommand\bool{\mathfrak{B}}
\newcommand\gbool{(\{0,1\},\zero,\un,=)}

\newcommand\osg{\overline{sg}}
\newcommand\moins{\mathrel{\dot{-}}}


%%% TRUC DE BARWISE


\newcommand\UM{\kU_{\kM}}
\newcommand\VV{\mathbb{V}}
% Hereditarily finite
\newcommand\IHF{\mathbb{I}HF}
\newcommand\IHP{\IHYP}
\newcommand\IHYP{\mathbb{I}HYP}

\providecommand\citep[1]{\cite{#1}}

%% Cofinali


% probas
% défini dans amsmath
\renewcommand\Pr{\cp{Pr}}
\newcommand\Var{\cp{Var}}


% -- mesures
%\newcommand\card{card}
\newcommand\taille{\LANGUE{size}{taille}}
\newcommand\entree{e}
\newcommand\profondeur{\LANGUE{depth}{profondeur}}


%%%%%%%%%%%%%%%%%%%%%%%%%%%%%%%%%%%%%%%%%%%%%%%%%%%%%%%%%%%%%%%%%%%%%%%%%%%%
%                 MACHINES DE TURING (graphique)
%%%%%%%%%%%%%%%%%%%%%%%%%%%%%%%%%%%%%%%%%%%%%%%%%%%%%%%%%%%%%%%%%%%%%%%%%%%%



% Configuration de machine de Turing
\newcommand\configTuring[3]{#1\textbf{#2}#3}


%%% DESSIN CROIX PORU AIDER  A DESSINER

\newcommand\croix{  +(-1,-1) -- +(1,1) +(-1,1) --
  +(1,-1) }

%% DESSIN ACCOLADE

\newcommand\ACCOLADE[3]{
%#1: coordonnées de départ
%#2: coordonnées arrivée
%#3: texte
\draw[decorate,decoration={brace}]  (#1) -- (#2)
 node[midway,above=0.1cm] {#3};
}

%%%


\def\lcase{3.5 ex}
\def\hcase{3.5ex}
\tikzstyle{case} = []
%[draw,minimum width = \lcase,minimum height = 3.6ex,baseline]
\newcommand\dcase{ +(-0.5*\lcase,-0.5*\hcase) rectangle +(0.5*\lcase,0.5*\hcase)}
%\newcommand\lettre[1]{\node [name=l,case,on chain=1] {#1};}


\newcommand\RUBANMACHINEDETURING[6]{
%#1: texte
%#2: position de la tête
%#3: nb de blancs à gauche
%#4: nb de cases totales
%#5: texte au dessus du ruban
%#6: position du coin
    \draw (#6) +(\lcase,4ex ) node {#5};
    \draw (#6) +(-0.6,0)  node  {\ldots};  
    \draw (#6) node[coordinate] (t) {};
    \foreach \x in {1,2,...,#3} {
         \draw (t) node [case]   {$\blanc$} \dcase;
         \draw (t) ++(\lcase,0) node [coordinate] (t) {};
       };
   \StrLen{#1}[\Resultat] % ATTENTION IL FAUT CETTE SYNTAXE POUR FAIRE
                         % EQUIVALENT D UN EDEF
  
\foreach \x in {1,...,\Resultat} 
             { 
        \draw (t) node [case]   {\StrChar{#1}{\x}} \dcase;
        \draw (t) ++(\lcase,0) node [coordinate] (t) {};
             }
\pgfmathtruncatemacro{\z}{#4 - \Resultat - #3}

  \foreach \x in {1,2,..., \z} {
        \draw (t) node [case]   {$\blanc$} \dcase;
        \draw (t) ++(\lcase,0) node [coordinate] (t) {};
}
%% Convention: graphiquement, la case numéro n se trouve à (n*\lcase,0)
%% en comptant les case à partir de $0$

   \draw  (t) +(0.1,0) node [name=r] {\ldots}; 
}

\newcommand\RUBANMACHINEDETURINGtexte[6]{
%#1: texte
%#2: position de la tête
%#3: nb de blancs à gauche
%#4: nb de cases totales (équivalent)
%#5: texte au dessus du ruban
%#6: position du coin
    \draw (#6) +(\lcase,4ex ) node {#5};
    \draw (#6) +(-0.6,0)  node  {\ldots};  
    \draw (#6) node[coordinate] (t) {};
   \foreach \x in {1,2,...,#3} {
        \draw (t) node [case]   {$\blanc$} \dcase;
        \draw (t) ++(\lcase,0) node [coordinate] (t) {};
      };
   \draw (#6) ++(-2*\lcase,0.5*\lcase)-- +(\lcase*#4,0);
   \draw (#6) ++(-2*\lcase,-0.5*\lcase) -- +(\lcase*#4,0);
   
    \draw (t) ++(-0.5*\lcase,0) node [case,anchor=west]  (te) {#1};
    \draw (te.east)  node [anchor=west,coordinate] (t) {};

  \draw (t) ++(0.2*\lcase,0) node [coordinate] (t) {};
   \foreach \x in {1,2,...,#3} {
        \draw (t) node [case]   {$\blanc$} \dcase;
        \draw (t) ++(\lcase,0) node [coordinate] (t) {};
      };

    \draw  (t) +(0.1,0) node [anchor=west] {\ldots}; 
}


\newcommand\RUBANMACHINEDETURINGTETE[6]{
%
% dessine la tete
%
% output: (tete)
%
\RUBANMACHINEDETURING{#1}{#2}{#3}{#4}{#5}{#6}
% TETE
     \draw (#6) + ( #2 * \lcase + #3  *\lcase,- 0.5*\lcase )
     node[coordinate] (k) {};
     \draw[->] (k) +(0,-0.5) -- (k);
     \draw (k) +(0,-0.5) node[coordinate] (tete) {};
}


\newcommand\RUBANMACHINEDETURINGTETEtexte[6]{
%
% dessine la tete
%
% output: (tete)
%
\RUBANMACHINEDETURINGtexte{#1}{#2}{#3}{#4}{#5}{#6}
% TETE
     \draw (#6) + ( #2 * \lcase + #3  *\lcase,- 0.5*\lcase )
     node[coordinate] (k) {};
     \draw[->] (k) +(0,-0.5) -- (k);
     \draw (k) +(0,-0.5) node[coordinate] (tete) {};
}


\newcommand\MACHINEDETURING[6]{
%#1: texte
%#2: position de la tête
%#3: état
%#4: nb de blancs à gauche
%#5: nb de cases totales
%#6: texte en haut du ruban
%\begin{tikzpicture}[node distance=-0.15mm]

\RUBANMACHINEDETURINGTETE{#1}{#2}{#4}{#5}{#6}{0,0}
%
%% ETAT:
     \draw (tete) node [name=etat,draw,circle,anchor=north]  {#3};
     \draw (etat) -- (tete);
%\end{tikzpicture}
}

\newcommand\MACHINEDETURINGtexte[6]{
%#1: texte
%#2: position de la tête
%#3: état
%#4: nb de blancs à gauche
%#5: nb de cases totales
%#6: texte en haut du ruban
%\begin{tikzpicture}[node distance=-0.15mm]

\RUBANMACHINEDETURINGTETEtexte{#1}{#2}{#4}{#5}{#6}{0,0}
%
%% ETAT:
     \draw (tete) node [name=etat,draw,circle,anchor=north]  {#3};
     \draw (etat) -- (tete);
%\end{tikzpicture}
}

\newcommand\MACHINEDETURINGTROISRUBANS[9]{
%#1: texte1
%#2: position de la tête 1
%#3: texte2
%#4: position de la tête 2
%#5: texte3
%#6: position de la tête 3
%#7: état
%#8: nb de blancs à gauche 1
%#9: nb de cases totales
%\begin{tikzpicture}[node distance=-0.15mm]

\RUBANMACHINEDETURINGTETE{#1}{#2}{#8}{#9}{\LANGUE{tape}{ruban} $1$}{0,0}
\draw (tete) node[coordinate] (tete1) {};
\draw (tete1) +(8,0) node[coordinate] (tete1b) {};
\draw (tete1) -| (tete1b) ;

\RUBANMACHINEDETURINGTETE{#3}{#4}{#8}{#9}{\LANGUE{tape}{ruban} $2$}{0,-2}
\draw (tete) node[coordinate] (tete2) {};
\draw (tete2) +(-8.5,0) node[coordinate] (tete2b) {};
\draw (tete2) -| (tete2b) ;

\RUBANMACHINEDETURINGTETE{#5}{#6}{#8}{#9}{\LANGUE{tape}{ruban} $3$}{0,-4}
\draw (tete) node[coordinate] (tete3) {};

%% ETAT:
     \draw (tete3) node [name=etat,draw,circle,anchor=north]  {#7};
     \draw (etat) -- (tete3);
    \draw  (etat) -|  (tete2b);
\draw  (etat) -|  (tete1b);
%\end{tikzpicture}
}

\newcommand\MACHINEDETURINGTROISRUBANSENGLISH[9]{
%#1: texte1
%#2: position de la tête 1
%#3: texte2
%#4: position de la tête 2
%#5: texte3
%#6: position de la tête 3
%#7: état
%#8: nb de blancs à gauche 1
%#9: nb de cases totales
%\begin{tikzpicture}[node distance=-0.15mm]

\RUBANMACHINEDETURINGTETE{#1}{#2}{#8}{#9}{tape $1$}{0,0}
\draw (tete) node[coordinate] (tete1) {};
\draw (tete1) +(8,0) node[coordinate] (tete1b) {};
\draw (tete1) -| (tete1b) ;

\RUBANMACHINEDETURINGTETE{#3}{#4}{#8}{#9}{tape $2$}{0,-2}
\draw (tete) node[coordinate] (tete2) {};
\draw (tete2) +(-8.5,0) node[coordinate] (tete2b) {};
\draw (tete2) -| (tete2b) ;

\RUBANMACHINEDETURINGTETE{#5}{#6}{#8}{#9}{tape $3$}{0,-4}
\draw (tete) node[coordinate] (tete3) {};

%% ETAT:
     \draw (tete3) node [name=etat,draw,circle,anchor=north]  {#7};
     \draw (etat) -- (tete3);
    \draw  (etat) -|  (tete2b);
\draw  (etat) -|  (tete1b);
%\end{tikzpicture}
}



\newcommand\MACHINEDETURINGTROISRUBANStexte[9]{
%#1: texte1
%#2: position de la tête 1
%#3: texte2
%#4: position de la tête 2
%#5: texte3
%#6: position de la tête 3
%#7: état
%#8: nb de blancs à gauche 1
%#9: nb de cases totales
% \begin{tikzpicture}[node distance=-0.15mm]

\RUBANMACHINEDETURINGTETEtexte{#1}{#2}{#8}{#9}{\LANGUE{tape}{ruban} $1$}{0,0}
\draw (tete) node[coordinate] (tete1) {};
\draw (tete1) +(8,0) node[coordinate] (tete1b) {};
\draw (tete1) -| (tete1b) ;

\RUBANMACHINEDETURINGTETEtexte{#3}{#4}{#8}{#9}{\LANGUE{tape}{ruban} $2$}{0,-2}
\draw (tete) node[coordinate] (tete2) {};
\draw (tete2) +(-8.5,0) node[coordinate] (tete2b) {};
\draw (tete2) -| (tete2b) ;

\RUBANMACHINEDETURINGTETEtexte{#5}{#6}{#8}{#9}{\LANGUE{tape}{ruban} $3$}{0,-4}
\draw (tete) node[coordinate] (tete3) {};

%% ETAT:
     \draw (tete3) node [name=etat,draw,circle,anchor=north]  {#7};
     \draw (etat) -- (tete3);
    \draw  (etat) -|  (tete2b);
\draw  (etat) -|  (tete1b);
%\end{tikzpicture}
}


\newcommand\MACHINEDETURINGDEUXRUBANStexte[9]{
%#1: texte1  
%#2: position de la tête 1 
%#3: texte2 inutilisée
%#4: position de la tête 2 inutilisée
%#5: texte3
%#6: position de la tête 3
%#7: état
%#8: nb de blancs à gauche 1
%#9: nb de cases totales
% \begin{tikzpicture}[node distance=-0.15mm]

\RUBANMACHINEDETURINGTETEtexte{#1}{#2}{#8}{#9}{\LANGUE{tape}{ruban} $1$}{0,0}
\draw (tete) node[coordinate] (tete1) {};
\draw (tete1) +(8,0) node[coordinate] (tete1b) {};
\draw (tete1) -| (tete1b) ;

% \RUBANMACHINEDETURINGTETEtexte{#3}{#4}{#8}{#9}{\LANGUE{tape}{ruban} $2$}{0,-1.5}
% \draw (tete) node[coordinate] (tete2) {};
% \draw (tete2) +(-5.5,0) node[coordinate] (tete2b) {};
% \draw (tete2) -| (tete2b) ;

\RUBANMACHINEDETURINGTETEtexte{#5}{#6}{#8}{#9}{\LANGUE{tape}{ruban} $2$}{0,-1.5}
\draw (tete) node[coordinate] (tete3) {};

%% ETAT:
     \draw (tete3) node [name=etat,draw,circle,anchor=north]  {#7};
     \draw (etat) -- (tete3);
%   \draw  (etat) -|  (tete2b);
\draw  (etat) -|  (tete1b);
%\end{tikzpicture}
}


\def\argdix{2}
\def\argonze{25}
\newcommand\MACHINEDETURINGQUATRERUBANStextep[9]{
%#1: texte1
%#2: position de la tête 1
%#3: texte2
%#4: position de la tête 2
%#5: texte3
%#6: position de la tête 3
%#7: texte4
%#8: position de la tête 3
%#9: état
%#10: nb de blancs à gauche 1
%#11: nb de cases totales
% \begin{tikzpicture}[node distance=-0.15mm]

%% BUG sur #10 devenu \argdix
%% BUG SUR #11 devenu \argonze

\RUBANMACHINEDETURINGTETEtexte{#1}{#2}{\argdix}{\argonze}{\LANGUE{tape}{ruban} $1$}{0,0}
% \draw (tete) node[coordinate] (tete1) {};
% \draw (tete1) +(8,0) node[coordinate] (tete1b) {};
% \draw (tete1) -| (tete1b) ;

\RUBANMACHINEDETURINGTETEtexte{#3}{#4}{\argdix}{\argonze}{\LANGUE{tape}{ruban} $2$}{0,-1.5}
% \draw (tete) node[coordinate] (tete2) {};
% \draw (tete2) +(-5.5,0) node[coordinate] (tete2b) {};
% \draw (tete2) -| (tete2b) ;

\RUBANMACHINEDETURINGTETEtexte{#5}{#6}{\argdix}{\argonze}{\LANGUE{tape}{ruban} $3$}{0,-3.2}
\draw (tete) node[coordinate] (tete3) {};

\RUBANMACHINEDETURINGTETEtexte{#7}{#8}{\argdix}{\argonze}{\LANGUE{tape}{ruban} $4$}{0,-4.9}
\draw (tete) node[coordinate] (tete4) {};

%% ETAT:
     \draw (tete4) node [name=etat,draw,circle,anchor=north]  {#9};
%      \draw (etat) -- (tete3);
%     \draw  (etat) -|  (tete2b);
% \draw  (etat) -|  (tete1b);
%\end{tikzpicture}
}


%%%%%%%%%%%%%%%%%%%%%%%%%%%%%%%%%%%%%%%%%%%%%%%%%%%%%%%%%%%%%%%%%%%%%%%%%%%%
%                 MACHINES DE TURING (automate)
%%%%%%%%%%%%%%%%%%%%%%%%%%%%%%%%%%%%%%%%%%%%%%%%%%%%%%%%%%%%%%%%%%%%%%%%%%%%



\newcommand\dec{1ex}
\newcommand\myACCOLADE[4]{
\ACCOLADE{#1*\lcase+#4*\lcase-0.5*\lcase,0.5*\lcase+\dec}{#1*\lcase+#4*\lcase-0.5*\lcase+#3*\lcase,0.5*\lcase+\dec}{#2}
}


\newcommand\M{\Sigma}


\newcommand\sommet[2]{\node[state] (#1) #2 {$#1$};}
\newcommand\sommetinitial[2]{\node[state,initial] (#1) #2 {$#1$};}
\newcommand\sommetfinal[2]{\node[state,accepting] (#1) #2 {$#1$};}
\newcommand\gtransition[5]{\draw (#1) edge node {#2/#4 $#5$} (#3);}
\newcommand\gtransitionba[5]{\draw (#1) [loop above] edge node {#2/#4
    $#5$} (#3);}
\newcommand\gtransitionbb[5]{\draw (#1) [loop below] edge node {#2/#4
    $#5$} (#3);}
\newcommand\gtransitionbleft[5]{\draw (#1) [loop left] edge node {#2/#4
    $#5$} (#3);}
\newcommand\gtransitionbright[5]{\draw (#1) [loop right] edge node {#2/#4
    $#5$} (#3);}
\newcommand\gtransitionbl[5]{\draw (#1) [bend left] edge node {#2/#4
    $#5$} (#3);}
\newcommand\gtransitionbr[5]{\draw (#1) [bend right] edge node {#2/#4
    $#5$} (#3);}
\newcommand\gtransitionbadouble[9]{\draw (#1) [loop above] edge node {
\begin{tabular}{l}
#2/#4 $#5$ \\
  #7/#8 $#9$ \\
\end{tabular}
} (#3);}

\newcommand\SCALE{}

\newcommand\dessinmtp[2]{
\begin{center}
\begin{tikzpicture}[->,>=stealth',shorten >=1pt,auto,node distance=2.3cm,semithick,#2] 
#1
\end{tikzpicture}
\end{center}
}

\newcommand\dessinmt[1]{\dessinmtp{#1}{scale=1}}


\newcommand\ANIMATIONMT[1]{
\begin{overprint}
\begin{dessinp}{scale=0.52}
\foreach \av/\q/\b  [count=\xi]
in  
{#1}
{ 
\only<\xi| handout 0>{ 
 \StrLen{\av}[\Resintm] % ATTENTION IL FAUT CETTE SYNTAXE POUR FAIRE
                         % EQUIVALENT D UN EDEF
 \MACHINEDETURING{\av \b}{\Resintm}{$\q$}{5}{20}{}
}}
 \end{dessinp}
\end{overprint}
}

\newcommand\ANIMATIONMTd[1]{
\begin{overprint}
\begin{dessinp}{scale=0.52}
\foreach \av/\q/\b  [count=\xi]
in  
{#1}
{ 
\only<\xi| handout 0>{ 
 \StrLen{\av}[\Resintm] % ATTENTION IL FAUT CETTE SYNTAXE POUR FAIRE
                         % EQUIVALENT D UN EDEF
 \hspace{-1cm}\MACHINEDETURING{\av \b}{\Resintm}{$\q$}{5}{33}{}
}}
 \end{dessinp}
\end{overprint}
}

\def\MAXDET{20}

\newcommand\DIAGRAMMEESPACETEMPS[1]{
\draw node[coordinate] (tt) {};
\foreach \av/\q/\b  in  
{#1}
{ 

\RUBANMACHINEDETURING{\av{$\q$}\b}{0}{4}{\MAXDET}{}{tt};
\draw (tt) +(0,-\hcase) node[coordinate] (tt) {};
}
%\foreach \x in {0,...,19} {\draw (tt) + (\x*\lcase,0) node {$\vdots$};}
}



\newcommand\progun{
/q_0/0011, % initial
X/q_1/011,X0/q_1/11,X/q_2/0Y1,/q_2/X0Y1,
X/q_0/0Y1,XX/q_1/Y1,XXY/q_1/1, XX/q_2/YY, X/q_2/XYY,
XX/q_0/YY,XXY/q_3/Y,XXYY/q_3/B,XXYYB/q_4/B
}

\newcommand\DIAGRAMMEESPACETEMPSun{
\DIAGRAMMEESPACETEMPS
%% RECOPIE progun
{/q_0/0011, % initial
X/q_1/011,X0/q_1/11,X/q_2/0Y1,/q_2/X0Y1,
X/q_0/0Y1,XX/q_1/Y1,XXY/q_1/1, XX/q_2/YY, X/q_2/XYY,
XX/q_0/YY,XXY/q_3/Y,XXYY/q_3/B,XXYYB/q_4/B}
%% FIN RECOPIE
}

\newcommand\progdeux{
/q_0/0010, X/q_1/010, X0/q_1/10, X/q_2/0Y0,/q_2/X0Y0,
X/q_0/0Y0,XX/q_1/Y0,XXY/q_1/0,XXY0/q_1/%B,
}


\newcommand\DIAGRAMMEESPACETEMPSdeux{
\DIAGRAMMEESPACETEMPS
{
%% RECOPIE progdeux
/q_0/0010, X/q_1/010, X0/q_1/10, X/q_2/0Y0,/q_2/X0Y0,
X/q_0/0Y0,XX/q_1/Y0,XXY/q_1/0,XXY0/q_1/%B,
%% FIN RECOPIE
}
}

\newcommand\DIAGRAMMEESPACETEMPSdeuxvariante{
\def\memolcase{\lcase}
\def\memoMAXDET{20}
\def\MAXDET{15}
\def\lcase{6 ex}
\DIAGRAMMEESPACETEMPS
{
%% RECOPIE progdeux
/q_00/010, X/q_10/10, X0/q_11/0, X/q_20/Y0,/q_2X/0Y0,
X/q_00/Y0,XX/q_1Y/0,XXY/q_10/,XXY0/q_1B/%B,
%% FIN RECOPIE
}
\def\lcase{\memolcase}
\def\MAXDET{\memoMAXDET}
}


\newcommand\ANIMATIONun{
\ANIMATIONMT
%% RECOPIE progun
{/q_0/0011, % initial
X/q_1/011,X0/q_1/11,X/q_2/0Y1,/q_2/X0Y1,
X/q_0/0Y1,XX/q_1/Y1,XXY/q_1/1, XX/q_2/YY, X/q_2/XYY,
XX/q_0/YY,XXY/q_3/Y,XXYY/q_3/B,XXYYB/q_4/B}
%% FIN RECOPIE
}



\newcommand\DESSINALGO[3]{
% argument 1= ou
% argument 2= texte 
% argument 3= nom du noeud ex m
\draw node[draw,#1,minimum size=1cm] (#3) {#2};

% pour fitting box
\draw (#3.west) +(-0.2cm,0) node (#3c2) {};
\draw (#3.north) +(0,+0.2cm) node (#3c3) {};
\draw (#3.south) +(0,-0.2cm) node (#3c4) {};
}

\newcommand\DESSINALGOa[3]{
\DESSINALGO{#1}{#2}{#3}
\draw (#3.east) +(-0.1cm,0.2cm) node(#3a) {};
\draw[->] (#3a) -- ++(0.5cm,0) node (#3aa) {};
\draw (#3aa) node[anchor=west] (#3aaa) {\LANGUE{Accept}{Accepte}};
\draw (#3aaa.east) node(#3aaaa) {};
}

\newcommand\DESSINALGOe[4]{
\DESSINALGO{#1}{#2}{#3}
\draw (#3.east) +(-0.1cm,0.2cm) node(#3a) {};
\draw[->] (#3a) -- ++(0.5cm,0) node (#3aa) {};
\draw (#3aa) node[anchor=west] (#3aaa) {#4};
\draw (#3aaa.east) node(#3aaaa) {};
}


\newcommand\DESSINALGOar[3]{
\DESSINALGOa{#1}{#2}{#3}
\draw (#3.east) +(-0.1cm,-0.2cm) node(#3r) {};
\draw[->] (#3r) -- ++(0.5cm,0) node (#3rr) {};
\draw (#3rr) node[anchor=west] (#3rrr) {\LANGUE{Reject}{Refuse}};
\draw (#3rrr.east) node(#3rrrr) {};
}


\newcommand\WINDOWarg[7] {
%#1: position
%#2..7: contenu du tableau
\draw (#1) node[case]  {#2} \dcase;
\draw (#1) ++ (\lcase,0) node[case] {#3} \dcase;
\draw (#1) ++ (2*\lcase,0) node[case] {#4} \dcase;
\draw (#1) ++ (0,-\lcase) node[case] {#5} \dcase;
\draw (#1) ++ (\lcase,-\lcase) node[case] {#6} \dcase;
\draw (#1) ++ (2*\lcase,-\lcase) node[case] {#7} \dcase;
}

\newcommand\WINDOW[3]{
%#1: positioin
%#2: contenu
%#3: étiquette
\WINDOWarg{#1}
{\StrChar{#2}{1}}{\StrChar{#2}{2}}{\StrChar{#2}{3}}
{\StrChar{#2}{4}}{\StrChar{#2}{5}}{\StrChar{#2}{6}}
\draw (#1) +(-\lcase*1.2,-\lcase*0.5) node {#3};
}





%%%%%%%%%%%%%%%%%%%%%%%%%%%%%%%%%%%%%%%%%%%%%%%%%%%%%%%%%%%%%%%%%%%%%%%%%%%%
%                 LISTINGS
%%%%%%%%%%%%%%%%%%%%%%%%%%%%%%%%%%%%%%%%%%%%%%%%%%%%%%%%%%%%%%%%%%%%%%%%%%%%


% LISTINGS
\RequirePackage{listings}
 \lstset{language=Java,
 mathescape=true,
 showspaces=false,
 showstringspaces=false,
 frame=single,
morekeywords={recursive,then,div,function,integer,read,out,else,par,endpar,seq,endseq,and,not,read,write,repeat,until,undef,let,to,endfor,endif,endwhile,guess,in,parallel},
% basicstyle=\ttfamily\small,
% basewidth={1.1ex},
% keywordstyle=\bf\color{blue},
% %stringstyle=\bf\ttfamily\color{green},
% commentstyle=\bfseries\color{magenta},
 literate={:=}{{$\gets$ }}1 {<=}{{$\leq$}}2  {>=}{{$\geq$}}2
 {!=}{{$\neq$}}1 
 }
% \lstset{escapechar=\%}
\let\lst\lstinline
\def\beginlisting{\begin{lstlisting}}
\def\endlisting{\end{lstlisting}}  




\newenvironment{algo}[1]{
  \lstset{escapechar=\@}
  \def\input{$#1$}
  \def\out{out}
  \def\BEGIN{\textbf{read} \input }
 \def\BEGINN{\textbf{repeat} }
  \def\END{\textbf{until out!=undef}}
  \def\ENDD{\textbf{write \out}}
  \def\ENDT{\textbf{until $\ctl$ = $q_a$ or $\ctl$ = $q_r$}}
  \def\ENDR{\textbf{until $\ctl$ = $q_a$}}
  \def\ENDP{\textbf{until $\ctl$ = $q_a$}}
}{}
\newenvironment{algon}[1]{
  \lstset{numbers=left}
  \begin{algo}{#1}
}{\end{algo}}
\newcommand\FSM[3]{FSM(#1,\lst{#2},#3)}
\newcommand\FSMi[3]{FSM(#1,{#2},#3)}




%%%%%%%%%%%%%%%%%%%%%%%%%%%%%%%%%%%%%%%%%%%%%%%%%
%
%
%                      TRADUCTION EN COURS
%
%%%%%%%%%%%%%%%%%%%%%%%%%%%%%%%%%%%%%%%%%%%%%%%%%

\newcommand\ENTREPOT[1]{\NDLR{ENTREPOT ICI: #1}}
\newcommand\scite[2][]{\cite[#1]{#2}}
\newcommand\nospellcheck[1]{{#1}}
\newcommand\nd{nd}
\newcommand\NL{
 
}

\newenvironment{FRANCAIS}{
\small
\color{orange}
}
{\normalsize
\color{black}
}


\newenvironment{SEMIFRANCAIS}{
\small
\color{violet}
}
{\normalsize
\color{black}
}

\newcommand\QFRANCAIS[1]{
\begin{FRANCAIS}
#1
\end{FRANCAIS}
}


%%%%%%%%%%%%%%%%%%%%%%%%%%%%%%%%%%%%%%%%%%%%%%%%%
%
%
%                      Wikipedia
%
%%%%%%%%%%%%%%%%%%%%%%%%%%%%%%%%%%%%%%%%%%%%%%%%%

\providecommand\tightlist{}
\providecommand\Complex{\C}
\newenvironment{myalgie}{\begin{array}{lll}}{\end{array}}

\newcommand\nl{\ \\ }




%%%%%%%%%%%%%%%%%%%%%%%%%%%%%%%%%%%%%%%%%%%%%%%%%
%
%
%                      Pour citations précises
%
%%%%%%%%%%%%%%%%%%%%%%%%%%%%%%%%%%%%%%%%%%%%%%%%%



\newcommand\citeprecis[2]{{\cite[#1]{#2}}}
\newcommand\citeprecisinv[2]{\citeprecis{#2}{#1}}

\usepackage{csquotes}



%%%%%%%%%%%%%%%%%%%%%%%%%%%%%%%%%%%%%%%%%%%%%%%%%
%
%
%                      Travail sur couleurs
%
%   toutes les réponses pointent dans le même sens, et le schéma correspond exactement 
%. au profil perceptif d’un protanope** (déficit des cônes L / longueurs d’onde longues)
%  MAIS SEVERE
%
%%%%%%%%%%%%%%%%%%%%%%%%%%%%%%%%%%%%%%%%%%%%%%%%%


% ================================
% GROUPE 1 — Verts sombres
% ================================
\definecolor{deepgreen}{RGB}{0,160,40}
\definecolor{forestgreen}{RGB}{0,150,70}
\definecolor{mossgreen}{RGB}{0,140,90}
\definecolor{greenwave}{RGB}{0,150,90}
\definecolor{forestmist}{RGB}{0,160,90}
\definecolor{darkteal}{RGB}{0,130,70}

% ================================
% GROUPE 2 — Teal / verts bleutés
% ================================
\definecolor{seagreen}{RGB}{0,130,100}
\definecolor{freshteal}{RGB}{0,140,130}
\definecolor{greenteal}{RGB}{0,140,120}
\definecolor{blueteal}{RGB}{0,150,120}      % ancien "teal" (doublon corrigé)
\definecolor{lagoon}{RGB}{0,180,120}
\definecolor{greenmist}{RGB}{0,170,110}

% ================================
% GROUPE 3 — Turquoises / Aqua
% ================================
\definecolor{springgreen}{RGB}{0,190,100}
\definecolor{mintgreen}{RGB}{10,200,150}
\definecolor{oceanfoam}{RGB}{0,200,150}
\definecolor{aquagreen}{RGB}{30,210,170}
\definecolor{coastalaqua}{RGB}{40,210,180}
\definecolor{softaqua}{RGB}{40,210,190}

% ================================
% GROUPE 4 — Cyan / Aqua clairs
% ================================
\definecolor{emeraldteal}{RGB}{0,180,140}
\definecolor{deepcyan}{RGB}{0,180,160}
\definecolor{lightcyan}{RGB}{0,200,160}
\definecolor{aquamarine}{RGB}{0,200,160}
\definecolor{lightaqua}{RGB}{60,230,200}
\definecolor{mistblue}{RGB}{80,240,210}

% ================================
% GROUPE 5 — Bleus clairs
% ================================
\definecolor{paleaqua}{RGB}{80,200,190}
\definecolor{glacier}{RGB}{90,150,200}
\definecolor{skyblueA}{RGB}{60,150,220}      % ancien "skyblue"
\definecolor{lightsky}{RGB}{90,170,230}
\definecolor{iceblue}{RGB}{120,190,240}
\definecolor{frost}{RGB}{150,210,250}

% ================================
% GROUPE 6 — Bleus moyens
% ================================
\definecolor{steelblue}{RGB}{30,100,160}
\definecolor{cadetblue}{RGB}{70,120,180}
\definecolor{coolblue}{RGB}{110,170,220}
\definecolor{morningblue}{RGB}{140,200,240}
\definecolor{blueiceA}{RGB}{40,150,210}       % ancien "blueice" (doublon corrigé)
\definecolor{blueair}{RGB}{60,170,225}

% ================================
% GROUPE 7 — Azurs / Bleu saturé
% ================================
\definecolor{tealblue}{RGB}{0,110,130}
\definecolor{deepbluecyan}{RGB}{0,110,170}
\definecolor{azur}{RGB}{0,70,180}
\definecolor{brightazure}{RGB}{20,130,200}
\definecolor{lightazure}{RGB}{80,180,235}
\definecolor{skylight}{RGB}{100,200,245}

% ================================
% GROUPE 8 — Bleus sombres / violets
% ================================
\definecolor{nightblue}{RGB}{40,0,110}
\definecolor{slateblue}{RGB}{20,40,120}
\definecolor{violetA}{RGB}{60,50,160}         % ancien "violet" (distinct de violetB ci-dessous)
\definecolor{blueviolet}{RGB}{80,60,180}
\definecolor{purpleA}{RGB}{100,70,200}        % ancien "purple"
\definecolor{lightpurple}{RGB}{130,90,210}

% ================================
% GROUPE 9 — Violets saturés / indigos
% ================================
\definecolor{lilac}{RGB}{160,110,220}
\definecolor{indigoA}{RGB}{50,0,130}          % ancien "indigo"
\definecolor{deepindigo}{RGB}{70,10,150}
\definecolor{royalpurple}{RGB}{90,20,170}

% ================================
% GROUPE 10 — Fermeture cercle perceptif
% ================================
\definecolor{coolmint}{RGB}{0,130,70}
\definecolor{deepsea}{RGB}{0,130,70}

% ================================
% PALETTE DEUTÉRANOPE — noms uniques
% ================================
\definecolor{crimson}{RGB}{200,30,60}
\definecolor{brickred}{RGB}{180,40,40}
\definecolor{coral}{RGB}{220,90,70}
\definecolor{salmon}{RGB}{240,130,120}

\definecolor{magenta}{RGB}{200,40,150}
\definecolor{fuchsia}{RGB}{220,60,170}
\definecolor{violetB}{RGB}{140,70,190}        % ancien "violet"
\definecolor{indigoB}{RGB}{90,50,180}         % ancien "indigo"

\definecolor{royalblueA}{RGB}{50,80,200}      % ancien "royalblue"
\definecolor{azure}{RGB}{30,120,230}
\definecolor{skyblueB}{RGB}{80,160,230}       % second "skyblue"
\definecolor{tealB}{RGB}{0,150,170}           % second "teal"

\definecolor{marine}{RGB}{0,80,140}




%% Alternative

% === Générateur automatique fill/text pour protanopes ===
% Usage : \ProtaColorPair{mistblue}
% Suppose que mistblue est déjà défini par \definecolor{mistblue}{RGB}{...}

\newcommand{\ProtaColorPair}[1]{%
  % fill = couleur originale
  \colorlet{#1fill}{#1}%
  % text = version assombrie idéalement lisible
  \colorlet{#1text}{#1!80!black}%
}

%Utilisation
%\ProtaColorPair{mistblue}

%\textcolor{mistbluefill}{Couleur claire} \quad
%\textcolor{mistbluetext}{Couleur pour texte}

	\foreach \i/\name in {
	  1/deepgreen,
	  2/greenwave,
	  3/lagoon,
	  4/oceanfoam,
	  5/softaqua,
	  6/mistblue,
 	 7/lightsky,
	  8/blueiceA,   % <-- corrigé
	  9/deepcyan,
	  10/azur,
	  11/blueviolet,
	  12/royalpurple}
	  {\ProtaColorPair{\name}}
	  


\newcommand\ROUEPROTANOPEDOUZE{
	% Roue chromatique protanope — 12 couleurs
	\begin{tikzpicture}[scale=3]

	\foreach \i/\name/\nametext in {
	  1/deepgreen/deepgreen,
	  2/greenwave/greenwave,
	  3/lagoon/lagoontext,
	  4/oceanfoam/oceanfoam,
	  5/softaqua/softaquatext,
	  6/mistblue/mistbluetext,
 	 7/lightsky/lightskytext,
	  8/blueiceA/blueiceA,   % <-- corrigé
	  9/deepcyan/blueiceA,
	  10/azur/azur,
	  11/blueviolet/blueviolet,
	  12/royalpurple/royalpurple}
	{
	  % secteur coloré
 	 \filldraw[fill=\name, draw=black, line width=0.2pt]
	    ({90-30*\i}:1)
	    -- ({90-30*(\i-1)}:1)
	    -- (0,0) -- cycle;

 	 % label dans la bonne couleur
	  \node at ({90-30*(\i-0.5)}:1.25)
	    {\textcolor{\nametext}{\small \nametext}};
	}

	\end{tikzpicture}
}


\newcommand\ROUEPROTANOPEVINGT{
\begin{tikzpicture}[scale=3]

\foreach \i/\name/\nametext in {
  1/deepgreen/deepgreen,
  2/greenwave/greenwave,
  3/lagoon/lagoon,
  4/oceanfoam/oceanfoam,
  5/softaqua/softaquatext,
  6/mistblue/mistbluetext,
  7/lightsky/lightsky,
  8/blueiceA/blueiceA,      % corrigé
  9/deepcyan/deepcyan,
  10/azur/azur,
  11/brightazure/brightazure,
  12/skyblueA/skyblueA,     % corrigé
  13/coolblue/coolblue,
  14/glacier/glacier,
  15/marine/marine,
  16/violetA/violetA,      % corrigé
  17/blueviolet/blueviolet,
  18/purpleA/purpleA,      % corrigé
  19/royalpurple/royalpurple,
  20/indigoA/indigoA}      % corrigé
{
  % Secteur coloré
  \filldraw[fill=\name, draw=black, line width=0.2pt]
    ({90-18*\i}:1)
    -- ({90-18*(\i-1)}:1)
    -- (0,0) -- cycle;

  % Étiquette
  \node at ({90-18*(\i-0.5)}:1.25)
    {\textcolor{\nametext}{\small \nametext}};
}

\end{tikzpicture}
}

\newcommand\ROUEDEUTERANOPEDOUZE{
% Roue chromatique deutéranope — 12 couleurs
\begin{tikzpicture}[scale=3]

\foreach \i/\name in {
  1/crimson,
  2/brickred,
  3/coral,
  4/salmon,
  5/magenta,
  6/fuchsia,
  7/violetB,       % corrigé
  8/indigoB,       % corrigé
  9/royalblueA,    % corrigé
  10/azure,
  11/skyblueB,     % corrigé
  12/tealB}        % corrigé
{
  % Secteur coloré
  \filldraw[fill=\name, draw=black, line width=0.2pt]
    ({90-30*\i}:1)
    -- ({90-30*(\i-1)}:1)
    -- (0,0) -- cycle;

  % Label dans la couleur
  \node at ({90-30*(\i-0.5)}:1.25)
    {\textcolor{\name}{\small \name}};
}

\end{tikzpicture}
}



% === Macro intelligente pour foncer une couleur de façon lisible pour protanope ===
% Usage : \protadarken{mistblue}{40}  -> crée mistblue@dark
% Le 40 signifie : 40% de mélange avec du noir

\newcommand{\DarkenForText}[2]{%
  % #1 = nom de la couleur
  % #2 = pourcentage (40 conseillé si tu ne sais pas)
  %
  % On crée une nouvelle couleur #1@dark
  \colorlet{#1text}{black!#2!#1}%
}





%%%%

% Groupe 1
\DarkenForText{deepgreen}{20}
\DarkenForText{forestgreen}{20}
\DarkenForText{mossgreen}{20}
\DarkenForText{greenwave}{20}
\DarkenForText{forestmist}{20}
\DarkenForText{darkteal}{20}

% Groupe 2
\DarkenForText{seagreen}{20}
\DarkenForText{freshteal}{20}
\DarkenForText{greenteal}{20}
\DarkenForText{blueteal}{20}
\DarkenForText{lagoon}{20}
\DarkenForText{greenmist}{20}

% Groupe 3
\DarkenForText{springgreen}{20}
\DarkenForText{mintgreen}{20}
\DarkenForText{oceanfoam}{20}
\DarkenForText{aquagreen}{20}
\DarkenForText{coastalaqua}{20}
\DarkenForText{softaqua}{20}

% Groupe 4
\DarkenForText{emeraldteal}{20}
\DarkenForText{deepcyan}{20}
\DarkenForText{lightcyan}{20}
\DarkenForText{aquamarine}{20}
\DarkenForText{lightaqua}{20}
\DarkenForText{mistblue}{30}

% Groupe 5
\DarkenForText{paleaqua}{20}
\DarkenForText{glacier}{20}
\DarkenForText{skyblueA}{20}
\DarkenForText{lightsky}{20}
\DarkenForText{iceblue}{20}
\DarkenForText{frost}{20}

% Groupe 6
\DarkenForText{steelblue}{20}
\DarkenForText{cadetblue}{20}
\DarkenForText{coolblue}{20}
\DarkenForText{morningblue}{20}
\DarkenForText{blueiceA}{20}
\DarkenForText{blueair}{20}

% Groupe 7
\DarkenForText{tealblue}{20}
\DarkenForText{deepbluecyan}{20}
\DarkenForText{azur}{20}
\DarkenForText{brightazure}{20}
\DarkenForText{lightazure}{20}
\DarkenForText{skylight}{20}

% Groupe 8
\DarkenForText{nightblue}{20}
\DarkenForText{slateblue}{20}
\DarkenForText{violetA}{20}
\DarkenForText{blueviolet}{20}
\DarkenForText{purpleA}{20}
\DarkenForText{lightpurple}{20}

% Groupe 9
\DarkenForText{lilac}{20}
\DarkenForText{indigoA}{20}
\DarkenForText{deepindigo}{20}
\DarkenForText{royalpurple}{20}

% Groupe 10
\DarkenForText{coolmint}{20}
\DarkenForText{deepsea}{20}

% Palette deutéranope
\DarkenForText{crimson}{20}
\DarkenForText{brickred}{20}
\DarkenForText{coral}{20}
\DarkenForText{salmon}{20}

\DarkenForText{magenta}{20}
\DarkenForText{fuchsia}{20}
\DarkenForText{violetB}{20}
\DarkenForText{indigoB}{20}

\DarkenForText{royalblueA}{20}
\DarkenForText{azure}{20}
\DarkenForText{skyblueB}{20}
\DarkenForText{tealB}{20}

% Couleur supplémentaire
\DarkenForText{marine}{20}

%%%%%%%%%%%%%%%%%%%%%
%
%.   apx proof
%
%%%%%%%%%%%%%%%%%%%%%



% do
\newcommand\PREUVESENAPPENDIXCOMPATIBLE{
	\newenvironment{appendixproof}{{\bf \LANGUE{Proof}{Démonstration}}:}{\hfill $\Box$}
	\newtheorem{theoremrep}[theorem]{Theorem}
	\newtheorem{lemmarep}[theorem]{Lemma}
	\newtheorem{propositionrep}[theorem]{Proposition}
	\newtheorem{corollaryrep}[theorem]{Corollary}
	\newtheorem{summaryrep}[theorem]{Summary}	
}

\providecommand\apxparameter{}

%% proof inline
%\providecommand\apxparameter{appendix=inline}
%% proof in appendix
%\providecommand\apxparameter{}


\apxproofmode{
\usepackage[\apxparameter]{apxproof}
\theoremstyle{plain}

\newtheoremrep{theorem}{Theorem}
%\newtheoremrep{theoremconjecture}[theorem]{Theorem-To-Be-Proved=Conjecture}
%\newtheoremrep{lemmaconjecture}[theorem]{Conjecture-To-Be-Proved=Conjecture}
\newtheoremrep{claim}[theorem]{Theorem}
\newtheoremrep{proposition}[theorem]{Proposition}
\newtheoremrep{lemma}[theorem]{Lemma}
\newtheoremrep{corollary}[theorem]{Corollary}
%%bug mars 26: \newtheoremrep{summary}[theorem]{Summary}
% Du coup:
\newtheoremrep{resume}[theorem]{Summary}
%bug mars 26: \newtheoremrep{openproblem}[theorem]{Open Problem}
%\newtheoremrep{claim}[theorem]{Claim}
}


%\apxproofmodesubsection{
%%apx proof, pour que compteur correct
%\makeatletter
%\AtBeginDocument{%
%\let\axp@section\subsection
%  \let\axp@oldsection\subsection
%}
%\makeatother
%\renewcommand\appendixprelim{\section{Missing proofs}}
%}


\makeatletter
\apxproofmodesubsection{
\let\axp@section\subsection
%\let\axp@oldsection\subsection
\renewcommand\appendixprelim{\section{Proofs from main documents}}
}
\makeatother


\makeatletter
\newcommand\appendixapxproofsubsection{
\appendix
\let\apx@orig@appendix\appendix
\renewcommand{\appendix}{}
}
\makeatother


%%% Pour checker/vérifier que bon


\newcommand\docheck[1]{\textcolor{check}{#1}}

\usepackage{xcolor}
%\definecolor{chose}{RGB}{0,110,50}
%\definecolor{forestgreen}{RGB}{0,150,70}
%\definecolor{ctruc}{RGB}{164,164,164}
%\definecolor{oceanfoam}{RGB}{0,200,150}
\definecolor{check}{RGB}{160,160,0}




 via \providecommand,
%% ou après via \renewcommand) :
\providecommand\BIBFILES{@@reference-biblio}
\providecommand\FIGCOMMONS{/Users/bournez/lib/LaTeX/Perso/FIG-COMMONS}

\newcommand\affiliationofficielleolivier{LIX, CNRS, École polytechnique, Institut Polytechnique de Paris, Palaiseau, France}
\newcommand\affiliationofficielleolivierslides{LIX, CNRS, École polytechnique, \\ Institut Polytechnique de Paris, Palaiseau, France}


%%%%%%%%%%%%%%%%%%%%%%%%
%
%              begin. Switch modes compatibilités
%
%%%%%%%%%%%%%%%%%%%%%%%

%% Apxproof mode
\providecommand\apxproofmode[1]{}
	% sinon: \newcommand\apxproof[1]{#1}
	
\providecommand\apxproofmodesubsection[1]{}
	% sinon: \newcommand\apxproofmodesubsection[1]{#1}
% utiliser: \appendixapxproofsubsection au lieu de \appendix dans ce cas. 
		

%% Lipics mode
\providecommand\lipicsmode{}
	% sinon: \newcommand\lipicsmode[1]{#1}
		% en vrai sert qu'à définir \nolipics
		
%%%%%%%%%%%%%%%%%%%%%%%%
%
%              end. Switch modes compatibilités
%
%%%%%%%%%%%%%%%%%%%%%%%


%% begin gestion des switchs

% en fait \nopilcs est la seule commande utilisée

\newcommand\nolipics[1]{#1}
\lipicsmode{\renewcommand\nolipics[1]{}}

%% Décorations (fourier + cadres colorés mdframed autour de definition,
%% theorem, exercice, example, remark, ...) : actives par défaut hors slides.
%% Pour les désactiver dans un document : \newcommand\NODECORATION[1]{}
%% avant 
\providecommand\nolncs[1]{#1}	

%% Chemins par défaut (surchargables avant \input{macros} via \providecommand,
%% ou après via \renewcommand) :
\providecommand\BIBFILES{@@reference-biblio}
\providecommand\FIGCOMMONS{/Users/bournez/lib/LaTeX/Perso/FIG-COMMONS}

\newcommand\affiliationofficielleolivier{LIX, CNRS, École polytechnique, Institut Polytechnique de Paris, Palaiseau, France}
\newcommand\affiliationofficielleolivierslides{LIX, CNRS, École polytechnique, \\ Institut Polytechnique de Paris, Palaiseau, France}


%%%%%%%%%%%%%%%%%%%%%%%%
%
%              begin. Switch modes compatibilités
%
%%%%%%%%%%%%%%%%%%%%%%%

%% Apxproof mode
\providecommand\apxproofmode[1]{}
	% sinon: \newcommand\apxproof[1]{#1}
	
\providecommand\apxproofmodesubsection[1]{}
	% sinon: \newcommand\apxproofmodesubsection[1]{#1}
% utiliser: \appendixapxproofsubsection au lieu de \appendix dans ce cas. 
		

%% Lipics mode
\providecommand\lipicsmode{}
	% sinon: \newcommand\lipicsmode[1]{#1}
		% en vrai sert qu'à définir \nolipics
		
%%%%%%%%%%%%%%%%%%%%%%%%
%
%              end. Switch modes compatibilités
%
%%%%%%%%%%%%%%%%%%%%%%%


%% begin gestion des switchs

% en fait \nopilcs est la seule commande utilisée

\newcommand\nolipics[1]{#1}
\lipicsmode{\renewcommand\nolipics[1]{}}

%% Décorations (fourier + cadres colorés mdframed autour de definition,
%% theorem, exercice, example, remark, ...) : actives par défaut hors slides.
%% Pour les désactiver dans un document : \newcommand\NODECORATION[1]{}
%% avant \input{macros}. Retirées automatiquement en mode lipics.
\providecommand\NODECORATION[1]{\ifnotSLIDE{#1}}
\lipicsmode{\renewcommand\NODECORATION[1]{}}

\providecommand\PASSIAPXPROOF[1]{#1}
\apxproofmode{\renewcommand\PASSIAPXPROOF[1]{}}



%% end gestion des switchs

\def\MACROSPOINTEXOUVERT{}


\providecommand\ifnotTDEXAM[1]{#1}
\providecommand\ifnotSLIDE[1]{#1}


  
% pour preuves en appendix
\providecommand\PREUVESENAPPENDIX[1]{}



%%%% STYLES: APRES VOLS


\NODECORATION{
\usepackage{fourier}
}

%% Repli si fourier (et donc fourier-orns, qui définit \danger) n'est pas
%% chargé ; doit rester APRÈS le bloc ci-dessus, sinon fourier-orns échoue
%% avec « \danger already defined ».
\providecommand\danger{}

%%%% Box around environment
\usepackage[framemethod=tikz]{mdframed}

\newcommand\ENVCOLOR[1]{#1}
\newcommand\BEGINTEMPENVCOLOR[1]{\let\envcolorwas\ENVCOLOR\renewcommand\ENVCOLOR[1]{#1}}
\newcommand\ENDTEMPENVCOLOR{\let\ENVCOLOR\envcolorwas}


\NODECORATION{
\surroundwithmdframed[hidealllines=true,backgroundcolor=\ENVCOLOR{blue!10},roundcorner=8pt]{exercice}
\surroundwithmdframed[hidealllines=true,backgroundcolor=\ENVCOLOR{blue!10},roundcorner=8pt]{exercise}
\surroundwithmdframed[hidealllines=true,backgroundcolor=\ENVCOLOR{blue!10},roundcorner=8pt]{exercicee}
\surroundwithmdframed[hidealllines=true,backgroundcolor=\ENVCOLOR{blue!10},roundcorner=8pt]{exerciced}
\surroundwithmdframed[hidealllines=true,backgroundcolor=\ENVCOLOR{blue!10},roundcorner=8pt]{exercicec}
\surroundwithmdframed[hidealllines=true,backgroundcolor=\ENVCOLOR{blue!10},roundcorner=8pt]{exercicedc}
\surroundwithmdframed[hidealllines=true,backgroundcolor=\ENVCOLOR{orange!20},roundcorner=8pt]{example}
\surroundwithmdframed[hidealllines=true,backgroundcolor=\ENVCOLOR{orange!20},roundcorner=8pt]{exemple}
\surroundwithmdframed[hidealllines=true,backgroundcolor=\ENVCOLOR{orange!10},roundcorner=8pt]{remark}
\ifnotSLIDE{\surroundwithmdframed[backgroundcolor=\ENVCOLOR{green!15},roundcorner=8pt]{definition}}
\surroundwithmdframed[hidealllines=true,backgroundcolor=\ENVCOLOR{gray!35}]{proposition}
\ifnotSLIDE{\surroundwithmdframed[backgroundcolor=\ENVCOLOR{gray!35},roundcorner=8pt]{theorem}}
\surroundwithmdframed[hidealllines=true,backgroundcolor=\ENVCOLOR{gray!35}]{lemma}
\surroundwithmdframed[hidealllines=true,backgroundcolor=\ENVCOLOR{gray!35}]{corollary}


\surroundwithmdframed[backgroundcolor=red!35,roundcorner=8pt]{theoremaprouver}
\surroundwithmdframed[backgroundcolor=red!35,roundcorner=8pt]{conjectureaprouver}
}

%%% tcolorbox

\usepackage{tcolorbox}
\tcbuselibrary{skins}


  %%. Commande magique
%% \tcolorboxenvironment{⟨name⟩}{⟨options⟩}: attach options de tcolorbox à environnement
%% alternative: je cree des boites moi meme: voici des exemples


\newtcolorbox{maboiteconfiance}[2][]{colback=red!5!white, colframe=red!75!black,fonttitle=\bfseries, colbacktitle=red!85!black,enhanced,
attach boxed title to top center={yshift=-2mm},
  title={#2},#1}
  
\newtcolorbox{maboiteresume}[2][]{colback=red!5!white, colframe=red!75!black,fonttitle=\bfseries, colbacktitle=red!85!black,enhanced,
attach boxed title to top center={yshift=-2mm},
  title={#2},#1}
  \newtcolorbox{maboiteresumeperso}[2][]{colback=red!5!white, colframe=red!75!black,fonttitle=\bfseries, colbacktitle=red!85!black,enhanced,
attach boxed title to top center={yshift=-2mm},
  title={#2},#1}
  \newtcolorbox{maboitecommentaire}[2][]{colback=red!5!white, colframe=red!75!black,fonttitle=\bfseries, colbacktitle=red!85!black,enhanced,
attach boxed title to top center={yshift=-2mm},
  title={#2},#1}
\newtcolorbox{maboitetruc}[2][]{colback=red!5!white, colframe=red!75!black,fonttitle=\bfseries, colbacktitle=red!85!black,enhanced,
attach boxed title to top left={yshift=-2mm},
  title={#2},#1}

%% Boîte des mises en évidence CARE. Couleurs pensées pour la
%% protanopie : l'axe bleu-jaune reste discriminant (jaune = \IMPORTANT
%% d'Olivier, bleu = CARE), et la distinction ne repose pas que sur la
%% couleur (cadre + bandeau étiqueté, lisibles même en niveaux de gris).
\newtcolorbox{maboiteimportantcare}[2][]{colback=cyan!10!white, colframe=blue!60!black,fonttitle=\bfseries, colbacktitle=blue!70!black,enhanced,
attach boxed title to top center={yshift=-2mm},
  title={#2},#1}
  


%%%% FIN STYLES



%%%%%%%%%%%%%%%%%%%%%%%%%%%%%%%%%%%%%%%%%%%%%%%%%
%%%%%%%%%%%%%%%%%%%%%%%%%%%%%%%%%%%%%%%%%%%%%%%%%
%
%
%                      Version ENSEIGNEMENT 2020
%
%
%%%%%%%%%%%%%%%%%%%%%%%%%%%%%%%%%%%%%%%%%%%%%%%%%
%%%%%%%%%%%%%%%%%%%%%%%%%%%%%%%%%%%%%%%%%%%%%%%%%

%english, french
\ifdefined\CESTDUFRANCAIS
	\providecommand\LANGUE[2]{#2}
	\usepackage[french]{babel}
\else
	\providecommand\LANGUE[2]{#1}
\fi
\providecommand\LANGUEinv[2]{\LANGUE{#2}{#1}}

%%%%%%%%%%%%%%%%%%%%%%%%%%%%%%%%%%%%%%%%%%%%%%%%%
%
%
%                      Packages
%
%
%%%%%%%%%%%%%%%%%%%%%%%%%%%%%%%%%%%%%%%%%%%%%%%%%


%%% SURLIGNEUR
\providecommand\ifnotmodeDIFFUSIONFINALE[1]{
%% 2021: Comprend pas mais assure
\ifdefined\SAFEMODE
\else
#1
\fi
%#1
}
\providecommand\ifmodeDIFFUSIONFINALE[1]{
\ifdefined\SAFEMODE
#1
\else
\fi
}
\newcommand\SURLIGNE[1]{
\ifnotmodeDIFFUSIONFINALE{
\BEGINTEMPENVCOLOR{blue!30}
{                        \colorbox{blue!30}{

\noindent \begin{minipage}{\dimexpr\textwidth-2\fboxsep\relax}
#1
\end{minipage}
}
}
\ENDTEMPENVCOLOR}
\ifmodeDIFFUSIONFINALE{
\noindent \begin{minipage}{\textwidth}
#1
\end{minipage}
}
}


\newenvironment{confiance}{\begin{maboiteresume}[colback=brown!50]{Résumé ici}}{\end{maboiteresume}}
\newcommand\CONFIANCE[2][Confiance ici]{
\begin{maboiteconfiance}[colback=yellow!50]{#1}
\textbf{ATTENTION : #2}
\end{maboiteconfiance}
}

\newenvironment{resume}{\begin{maboiteresume}[colback=brown!50]{Résumé ici}}{\end{maboiteresume}}
\newcommand\RESUME[2][Résumé ici]{
\begin{maboiteresume}[colback=brown!50]{#1}
#2
\end{maboiteresume}
}

\newenvironment{resumeperso}{\begin{maboiteresumeperso}[colback=brown!50]{Résumé perso ici}}{\end{maboiteresumeperso}}
\newcommand\RESUMEPERSO[2][Résumé perso  ici]{
\begin{maboiteresumeperso}[colback=brown!50]{#1}
#2
\end{maboiteresumeperso}
%
%
%\todo[inline, color=brown]{résumé: #1}
%\colorbox{brown}{
%\noindent \begin{minipage}{\textwidth}
%{
%\bf 
%#2
%}
%\end{minipage}
%}
}

\newcommand\DANSLAMARGE[1]{
\ifnotmodeDIFFUSIONFINALE{
\marginpar{\textcolor{brown}{#1}}}
}

\newenvironment{commentaireperso}{\begin{maboitecommentaire}[colback=brown!30]{Commentaire perso} Commentaire perso:}{\end{maboitecommentaire}}
\newcommand\COMMENTAIREPERSO[2][Commentaire perso]{
\begin{maboitecommentaire}[colback=brown!30]{#1}
#2
\end{maboitecommentaire}
}


\newcommand\SURSURLIGNE[2][truc surligné ici]{
%\todo[inline, color=orange]{surligné: #1}
\ifnotmodeDIFFUSIONFINALE{
\BEGINTEMPENVCOLOR{blue!60}
	                 \colorbox{blue!60}{
\noindent \begin{minipage}{\dimexpr\textwidth-2\fboxsep\relax}
{
\bf 
#2
}
\end{minipage}
}
\ENDTEMPENVCOLOR
}
\ifmodeDIFFUSIONFINALE{#2}
}

\newcommand\IMPORTANT[1]{
\ifnotmodeDIFFUSIONFINALE{
\BEGINTEMPENVCOLOR{yellow!100}
		        \colorbox{yellow!100}{

\noindent \begin{minipage}{\dimexpr\textwidth-2\fboxsep\relax}
{ 
#1
}
\end{minipage}
}
\ENDTEMPENVCOLOR
}
\ifmodeDIFFUSIONFINALE{#1}
}

%% Pendant de \IMPORTANT, réservé aux mises en évidence CARE ;
%% \IMPORTANT reste réservé aux annotations d'Olivier. Comme \IMPORTANT,
%% redevient du texte simple en mode DIFFUSIONFINALE.
\newcommand\IMPORTANTCARE[2][Important (CARE)]{
\ifnotmodeDIFFUSIONFINALE{
\begin{maboiteimportantcare}{#1}
#2
\end{maboiteimportantcare}
}
\ifmodeDIFFUSIONFINALE{#2}
}


\newcommand\QUESTION[1]{
\ifnotmodeDIFFUSIONFINALE{
\BEGINTEMPENVCOLOR{oceanfoam!100}
		        \colorbox{oceanfoam!100}{

\noindent \begin{minipage}{\dimexpr\textwidth-2\fboxsep\relax}
{ 
#1
}
\end{minipage}
}
\ENDTEMPENVCOLOR
}
\ifmodeDIFFUSIONFINALE{#1}
}




\newcommand\SOUSLIGNE[1]{
\colorbox{lightgray}{

\noindent \begin{minipage}{\dimexpr\textwidth-2\fboxsep\relax}
{ \color{gray}
#1
}
\end{minipage}
}
}



%% Gestion des cross-references cross-referencing
\usepackage{xr-hyper}
\usepackage{hyperref}

%%
\usepackage{versions}
\usepackage{amsmath,amssymb,mathrsfs,amsfonts}
%\usepackage{mathabx}
\usepackage{listings}
\usepackage{url}
\usepackage{versions}
\usepackage{color}
\usepackage[colorinlistoftodos]{todonotes}
%\usepackage{todonotes}
\usepackage{proof}
\usepackage{xstring}

%% Index
\usepackage{makeidx}
\makeindex
\providecommand\DOINDEXPRINT{}
\providecommand\SIGNALEMOTNOUV[1]{}
\ifdefined\INDEXPRINT
	\ifnotSAFEMODE{
%	\usepackage{showidx}
%	\let\oldindex\index
%	\renewcommand{\index}[1]{\oldindex{#1}\marginpar{LA: \tiny#1}}
	\renewcommand\DOINDEXPRINT{
	      	\let\oldindex\index
		\renewcommand{\index}[1]{\oldindex{##1}\marginpar{IDXLA: \tiny##1}}
		\renewcommand\SIGNALEMOTNOUV[1]{\marginpar{MNV: \small\textbf{##1}}}
	}
	}
\fi

% === gestion des accents


%soit iso latin1
%\usepackage[cyr]{aeguill}
%soit utf8
\usepackage[utf8]{inputenc}
\usepackage[T1]{fontenc}


%%%  Pour citation \Cref (cref, ref)	

\usepackage{cleveref}


%%%%%%%%%%%%%%%%%%%%%%%%%%%%%%%%%%%%%%%%%%%%%%%%%
%
%
%                      Gestion des accents spéciaux.
%
%
%%%%%%%%%%%%%%%%%%%%%%%%%%%%%%%%%%%%%%%%%%%%%%%%%

\input{macros-accents}
%% Tolérant tant que macros-markdown.tex n'est pas dans le dépôt :
\InputIfFileExists{macros-markdown}{}{\typeout{*** macros-markdown.tex introuvable: ignore ***}}


%\usepackage{PDFdraftcopy}


%-- Pour état des chapitres


%\newcommand\ETAT{\draftstring{BROUILLON}}
%
%\newcommand\brouillon{\renewcommand\ETAT{\draftstring{BROUILLON}}}
%
%\newcommand\final{\renewcommand\ETAT{\draftstring{\rien}}}
%
%\newcommand\travail{\renewcommand\ETAT{\draftstring{CHANTIER}}}
%
%\newcommand\bazar{\renewcommand\ETAT{\draftstring{BAZAR TOTAL}}}
%
%\final



%%%%%%%%%%%%%%%%%%%%%%%%%%%%%%%%%%%%%%%%%%%%%%%%%
%
%
%                      TRADUCTION EN COURS
%
%%%%%%%%%%%%%%%%%%%%%%%%%%%%%%%%%%%%%%%%%%%%%%%%%


% === volé Amaury



\usepackage{stmaryrd}
\usepackage{ifthen}
\usepackage{mathtools}
\usepackage{dsfont}
\PASSIAPXPROOF{\usepackage{thmtools}}
\usepackage{longtable}
\usepackage{mathdots}
\usepackage{booktabs}


%%%%%%%%%%%%%%%%%%%%%%%%%%%%%%%%%%%%%%%%%%%%%%%%%
%
%
%                      TIKZ STUFFS
%
%
%%%%%%%%%%%%%%%%%%%%%%%%%%%%%%%%%%%%%%%%%%%%%%%%%

%%% volé pour dessin des GPAC: PAS SUR QUE BESOIN de TOUT

\usepackage{tikz}
\usetikzlibrary{calc}
\usepgflibrary{arrows}
\usetikzlibrary{fit}
\usetikzlibrary{patterns}
\usetikzlibrary{decorations.pathreplacing}
\usetikzlibrary{shapes.multipart}
\usetikzlibrary{shapes.geometric}
\usetikzlibrary{shapes.misc}
\usetikzlibrary{positioning}
\usetikzlibrary{graphs}
\usetikzlibrary{topaths}
\usetikzlibrary{arrows.meta}
\usetikzlibrary{decorations.pathreplacing}
\usetikzlibrary{decorations.markings}
\usetikzlibrary{decorations.pathmorphing}
%
 \usetikzlibrary{calc}
 \usetikzlibrary{automata}
\usetikzlibrary{chains}
\usetikzlibrary{scopes}
 \usetikzlibrary{positioning}
 \usetikzlibrary{through,backgrounds}
\usepackage{circuitikz}
% % pour accolades




%%%%%%%%%%%%%%%%%%%%%%%%%%%%%%%%%%%%%%%%%%%%%%%%%
%
%
%                      MACROS locales
%
%
%%%%%%%%%%%%%%%%%%%%%%%%%%%%%%%%%%%%%%%%%%%%%%%%%


%====================
%                      Moi
%=====================


\newcommand\myaddress{ 
{\sc LIX, CNRS, Ecole Polytechnique, Institut Polytechnique de Paris, }\\ 91128 Palaiseau Cedex, FRANCE \\
{\small Olivier.Bournez@lix.polytechnique.fr}}



%====================
%                      un peu de trucs sales
%=====================

%=  pour style lipics et ses trucs

\providecommand\ifnotPOLYCHAPTERMODE[1]{#1}
\providecommand\ifPOLYCHAPTERMODE[1]{}

\ifnotPOLYCHAPTERMODE{
\ifdefined\thechapter
\providecommand\spnewtheorem[4]{\newtheorem{#1}{#2}[chapter]}
\providecommand\spnewtheoremi[5]{\newtheorem{#1}[#5]{#2}}
\else
\providecommand\spnewtheorem[4]{\newtheorem{#1}{#2}}
\providecommand\spnewtheoremi[5]{\newtheorem{#1}[#5]{#2}}
\fi
}
\ifPOLYCHAPTERMODE{
\providecommand\spnewtheorem[4]{\newtheorem{#1}{#2}}
\providecommand\spnewtheoremi[5]{\newtheorem{#1}[#5]{#2}}
}
\providecommand\nolipics[1]{#1}




%====================
%                      un peu de trucs franchement sales
%=====================


% on l'inclus pas ici.
%\input{macros-moins-propres}

%%%%%%%%%%%%%%%%%%%%%%%%%%%%%%%%%%%%%%%%%%%%%%%%%
%
%
%                      Page de Garde Cours X
%
%
%%%%%%%%%%%%%%%%%%%%%%%%%%%%%%%%%%%%%%%%%%%%%%%%%


\includeversion{metdate}

\newcommand\XPAGEDEGARDE[2]{
\thispagestyle{empty}
\title{{#1} \\[2cm] %\\[3cm] 
{\large #2}
}
\author{
 Olivier Bournez% \inst{1}
\\[0.5cm] 
bournez@lix.polytechnique.fr%{\myaddress}
}
%\date{
%\vspace{2cm}
%\includegraphics[height=2.5cm]{/Users/bournez/lib/LaTeX/Perso/FIG-COMMONS/logo_X_new}
%}
\begin{metdate}
\date{
\vspace{2cm}
Version \LANGUE{of}{du} \today \\[1cm] 
\includegraphics[height=3.5cm]{/Users/bournez/lib/LaTeX/Perso/FIG-COMMONS/logo_X_new}
}
\end{metdate}
\maketitle
}


%%%%%%%%%%%%%%%%%%%%%%%%%%%%%%%%%%%%%%%%%%%%%%%%%
%
%
%                      MACROS UTILES
%
%
%%%%%%%%%%%%%%%%%%%%%%%%%%%%%%%%%%%%%%%%%%%%%%%%%

%%%                      Macros
 
%-------------
%--- vectors  ---
%-------------

 
\newcommand\vectorl[1]{{\mathbf#1}}
\newcommand\tu[1]{\vectorl{#1}}

\newcommand\vp{\vectorl{p}}
\newcommand\va{\vectorl{a}}
\newcommand\vb{\vectorl{b}}
\newcommand\vc{\vectorl{c}}
\newcommand\vf{\vectorl{f}}
\newcommand\vn{\vectorl{n}}
\newcommand\vx{\vectorl{x}}
\newcommand\vX{\vectorl{X}}
\newcommand\vy{\vectorl{y}}
\newcommand\vz{\vectorl{z}}
\newcommand\ve{\vectorl{e}}
\newcommand\vv{\vectorl{v}}
\newcommand\vr{\vectorl{r}}
\newcommand\vu{\vectorl{u}}
\newcommand\vA{\vectorl{A}}
\newcommand\vB{\vectorl{B}}
\newcommand\vC{\vectorl{C}}
\newcommand\vD{\vectorl{D}}

%-------------
%--- overline ---
%-------------


\newcommand\ox{\overline{x}}
\newcommand\oy{\overline{y}}
\newcommand\oc{\overline{c}}
\newcommand\oa{\overline{a}}
\newcommand\ow{\overline{w}}
\newcommand\od{\overline{d}}
\newcommand\ob{\overline{b}}
\newcommand\oP{\overline{P}}
\newcommand\oee{\overline{e}}
\newcommand\ot{\overline{t}}
\newcommand\ou{\overline{u}}
\newcommand\ov{\overline{v}}
\newcommand\oz{\overline{z}}
\newcommand\am{\overline{a}}
\newcommand\om{\overline{m}}
\newcommand\oj{\overline{j}}

%----------------
%--- ensembles ---
%----------------

\newcommand\Q{\mathbb{Q}}
\newcommand\R{\mathbb{R}}
\newcommand\Z{\mathbb{Z}}
\providecommand\C{\mathbb{C}}
\newcommand\N{\mathbb{N}}
\newcommand\K{\mathbb{K}}
\newcommand\F{\mathbb{F}}

\newcommand\dyadic{\mathbb{D}}
\newcommand\Zd{{\Z_2}}
\newcommand\Zp{{\Z_p}}
%
\newcommand\RP{\mathbb{R}^{>0}}
\newcommand\QP{\mathbb{Q}^{>0}}


%----------------
%--- lettres cal ---
%----------------

\newcommand\mK{\mathcal{K}}
\newcommand\mD{\mathcal{D}}
\newcommand\mA{\mathcal{A}}
\newcommand\mL{\mathcal{L}}
\newcommand\mX{\mathcal{X}}
\newcommand\mG{\mathcal{G}}
\newcommand\mE{\mathcal{E}}
\newcommand\mH{\mathcal{H}}
\newcommand\mR{\mathcal{R}}
\newcommand\mF{\mathcal{F}}
\newcommand\mJ{\mathcal{J}}
\newcommand\mO{\mathcal{O}}
\newcommand\mP{\mathcal{P}}
\newcommand\mN{\mathcal{N}}
\newcommand\mS{\mathcal{S}}
\newcommand\mM{\mathcal{M}}
\newcommand\mC{\mathcal{C}}
\newcommand\mV{\mathcal{V}}
\newcommand\mI{\mathcal{I}}
\newcommand\mT{\mathcal{T}}

\newcommand\fC{\mathcal{C}}

%----------------
%--- lettres frak ---
%----------------


\newcommand\kM{\mathfrak{M}}
\newcommand\kN{\mathfrak{N}}
\newcommand\kU{\mathfrak{U}}
\newcommand\kg{\mathfrak{g}}



%--------------------------
%--- Façon noter noms classes ---
%---------------------------

\newcommand{\cp}[1]{\operatorname{#1}}


%--------------------------
%--- Calculabilité  ---
%---------------------------

% -- classes
\LANGUE{
\newcommand\RE{\cp{CE}}
}
{\newcommand\RE{\cp{RE}}
}

\newcommand\Rec{\cp{D}}
\newcommand\PR{\cp{PR}}
\newcommand\Gregn{\mE_n}


%--------------------------
%--- Réductions  ---
%---------------------------

\newcommand\lem{\le_{m}}
\newcommand\lelog{\le_{L}}
\newcommand\equivlog{\equiv_{L}}



%--------------------------
%--- Complexité  ---
%---------------------------

% COMPLEXITE

% -- P
\newcommand{\Ptime}{\cp{P}}      % P
\newcommand\PTIME{\Ptime}
\newcommand\PTIMEs[1]{{\PTIME}_{#1}}
\newcommand\PTIMEsp[1]{{\PTIME}_{#1}^0}
\renewcommand\P{\PTIME}
\newcommand{\FPtime}{\cp{FPTIME}}

% -- NP
\newcommand\NP{\cp{NP}}
\newcommand{\NPtime}{\NP}
\newcommand\NPs[1]{\cp{NP}_{#1}}
\newcommand{\FNP}{\cp{FNP}}

\newcommand\NDP{\cp{NDP}}
\newcommand\coNDP{\cp{coNDP}}
\newcommand\NDPs[1]{\cp{NDP}_{#1}}

% -- coNP
\newcommand\coNP{\cp{coNP}}

% -- PSPACE
\newcommand\PSPACE{\cp{ PSPACE}}
\newcommand\NPSPACE{\cp{ NPSPACE}}
\newcommand{\Pspace}{\PSPACE}
\newcommand{\FPspace}{\cp{FPSPACE}}


% -- LOGSPACE
\newcommand\LSPACE{\cp{LOGSPACE}}
\newcommand\NLSPACE{\cp{NLOGSPACE}}
\newcommand\coNLSPACE{\cp{coNLOGSPACE}}
\newcommand\coNL{\cp{coNL}}

% -- Autres
\newcommand\EXPTIME{\cp{EXPTIME}}
\newcommand\NEXPTIME{\cp{NEXPTIME}}
\newcommand\EXPSPACE{\cp{EXPSPACE}}
\newcommand\PH{\cp{PH}}
\newcommand\NC{\cp{NC}}
\newcommand\AC{\cp{AC}}
\newcommand\PPAD{\cp{\mathbf{PPAD}}}
\newcommand\PLS{\cp{\mathbf{PLS}}}

% -- Reverse math
\newcommand{\WKL}{\ensuremath{\docheck{\mathsf{WKL}_0}}}
\newcommand{\ACA}{\ensuremath{\docheck{\mathsf{ACA}_0}}}
\newcommand{\ATR}{\ensuremath{\docheck{\mathsf{ATR}_0}}}
\newcommand{\RCAO}{\ensuremath{\docheck{\mathsf{RCA}_0}}}
\newcommand{\PiCA}{\ensuremath{\docheck{\Pi^1_1\!\text{-}\mathsf{CA}_0}}}


% -- circuits
\newcommand\CIRCUIT[2]{\cp{CIRCUIT(#1,#2)}}
\newcommand\UCIRCUIT[2]{\cp{UCIRCUIT(#1,#2)}}

% -- Classes proba
\newcommand\PP{\cp{PP}}
\newcommand\RPP{\cp{RP}}
\newcommand\ZPP{\cp{ZPP}}
\newcommand\coRP{\cp{coRP}}
\newcommand\BPP{{\cp{BPP}}}

% -- IP
\newcommand\IP{\cp{IP}}
\newcommand\PCP{\cp{PCP}}

% -- non-uniforme
\newcommand\nuP{\mathbb{P}}
\newcommand\nuPs[1]{\mathbb{P}_{#1}}
\newcommand\nuPsp[1]{\mathbb{P}_{#1}^0}
\newcommand\nuNP{\mathbb{N}\mathbb{P}}
%\newcommand\poly{/{poly}}
\newcommand\Ppoly{\PTIME_{\poly}}
\newcommand\NPpoly{\NP_{\poly}}

% -- conditions d'uniformité
\newcommand\Luniforme{\cp{L}}
\newcommand\BCuniforme{\cp{BC}}


% Classes de circuits
\newcommand\TC{\cp{TC}}
\newcommand\LFP{\cp{LFP}}
\newcommand\FAC{\cp{FAC}}
\newcommand\FACC{\cp{FACC}}
\newcommand\FTC{\cp{FTC}}
\newcommand\FNC{\cp{FNC}}

% Logiques temorelles
\newcommand\CTL{\cp{CTL}}
\newcommand\LTL{\cp{LTL}}

% -- mesure temps/espace/taille

\newcommand\DTIME{\cp{DTIME}}
\newcommand\TIME[1]{\cp{TIME(#1)}}
\newcommand\TIMEs[2]{\cp{TIME_{#2}(#1)}}
\newcommand\TIMEsp[2]{\cp{TIME^0_{#2}(#1)}}
\newcommand\NTIME[1]{\cp{NTIME(#1)}}
\newcommand\NTIMEs[2]{\cp{NTIME_{#2}(#1)}}
\newcommand\NTIMEsp[2]{\cp{NTIME_{#2}^0(#1)}}
\newcommand\DSPACE{\cp{DSPACE}}
\newcommand\SPACE[1]{\cp{SPACE(#1)}}
\newcommand\SPACEs[2]{\cp{SPACE_{#2}(#1)}}
\newcommand\NSPACE[1]{\cp{NSPACE(#1)}}
\newcommand\NSPACEs[2]{\cp{NSPACE_{#2}(#1)}}
\newcommand\coNSPACE[1]{\cp{coNSPACE(#1)}}
\newcommand\TIMEON[2]{\cp{TIME(#1,#2)}}
\newcommand\SPACEON[2]{\cp{SPACE(#1,#2)}}
\newcommand\SIZE[1]{\cp{ size(#1)}}
\newcommand\ATIME{\cp{ATIME}}
\newcommand\ASPACE{\cp{ASPACE}}


%% autre:
\newcommand\LH{\cp{LH}}
\newcommand\FO{\cp{FO}}
\newcommand\SO{\cp{SO}}
\newcommand\ESOhorn{\cp{ESO_{horn}}}
\newcommand\SOhorn{\cp{SO_{horn}}}
\newcommand\ACzero{\cp{AC_0}}
\newcommand{\LINSPACE}{\cp{LINSPACE}}
\newcommand{\mFLINSPACE}{\mF_{\cp{LINSPACE}}}
\newcommand{\LOGSPACE}{\cp{LOGSPACE}}
\newcommand\ALOGTIME{\cp{ALOGTIME}}
\newcommand\RF{\cp{RF}}


\newcommand\BOOL[1]{\cp{{BOOLEAN(#1)}}}
\newcommand\DET{\cp{ DET}}

% % -- dirty


\newcommand{\cPspace}{\cp{\sharp PSPACE}}
\newcommand{\cAPtime}{\cp{\sharp APTIME}}
%\newcommand{\FPspace}{\cp{FPSPACE}}
%\newcommand{\FPspaceN}{\cp{FPSPACE}^{+}}
%\newcommand{\FPspaceN}{\cp{FPSPACE}}



%-------------------------
%--- Problèmes de décision  ---
%--------------------------

\newcommand\decisionname[1]{\cp{#1}}


\newcommand\Luniv{\decisionname{L_{univ}}}
\newcommand\HALTINGPROBLEM{\decisionname{HALTING-PROBLEM}}
\newcommand\SAT{\decisionname{SAT}}
\newcommand\MAXtSAT{\decisionname{MAX3SAT}}
\newcommand\REACH{\decisionname{REACH}}
\newcommand\CIRCUITVALUE{\decisionname{CIRCUITVALUE}}
\newcommand\MONOTONECIRCUITVALUE{\decisionname{MONOTONECIRCUITVALUE}}
\newcommand\THEOREMS{\decisionname{THEOREMS}}
\newcommand\QCIRCUITSAT{\decisionname{QCIRCUITSAT}}
\newcommand\QSAT{\decisionname{QSAT}}
\newcommand\QBF{\decisionname{QBF}}
\newcommand\SATVALUE{\decisionname{SATVALUE}}
\newcommand\SUCCINTSAT{\decisionname{SUCCINTSAT}}
\newcommand\SUCCINTSATVALUE{\decisionname{SUCCINTSATVALUE}}
\newcommand\GEOGRAPHY{\decisionname{GEOGRAPHY}}
\newcommand\PARITY{\decisionname{PARITY}}
\newcommand\MAJ{\decisionname{MAJ}}
\newcommand\CVP{\CIRCUITVALUE}
\newcommand\BOOLEANPROGRAMVALUE{BOOLEAN~PROGRAM~VALUE}

\newcommand\CLIQUE{\cp{CLIQUE}}
\LANGUEinv{
\newcommand\RECOUVREMENTDESOMMETS{\cp{RECOUVREMENT~DE~SOMMETS}}
\newcommand\RS{\RECOUVREMENTDESOMMETS}
}{
\newcommand\RECOUVREMENTDESOMMETSeng{\cp{VERTEX~COVER}}
\newcommand\RSeng{\RECOUVREMENTDESOMMETSeng}
}

\LANGUEinv{\newcommand\COUPUREMAXIMALE{\cp{COUPURE~MAXIMALE}}}
{\newcommand\COUPUREMAXIMALEeng{\cp{MAXIMAL~CUT}}}
\newcommand\CIRCUITHAMILTONIEN{\CHAMILTONIEN}
\LANGUEinv{\newcommand\VOYAGEURDECOMMERCE{\VCOMMERCE}}
{\newcommand\VOYAGEURDECOMMERCEeng{\VCOMMERCEeng}}

\LANGUEinv{\newcommand\CIRCUITLEPLUSLONG{\cp{CIRCUIT~LE~PLUS~LONG}}}
{\newcommand\CIRCUITLEPLUSLONGeng{\cp{LONGEST~CIRCUIT}}}
\LANGUEinv{\newcommand\SOMMEDESOUSENSEMBLE{\SUBSETSUM}}{
\newcommand\SOMMEDESOUSENSEMBLEeng{\SUBSETSUMeng}}
\LANGUEinv{
\newcommand\SACADOS{\cp{SAC~A~DOS}}}
{\newcommand\SACADOSeng{\cp{KNAPSACK}}}

%français
\LANGUEinv{
\newcommand\CHAMILTONIEN{\cp{CIRCUIT~HAMILTONIEN}}}
{\newcommand\CHAMILTONIENeng{\cp{HAMILTONIAN~CIRCUIT}}}
\LANGUEinv{
\newcommand\VCOMMERCE{\cp{VOYAGEUR~DE~COMMERCE}}}
{\newcommand\VCOMMERCEeng{\cp{TRAVELING~SALESMAN}}}
\LANGUEinv{
\newcommand\kCOLORABILITE{\cp{k-COLORABILITE}}}
{\newcommand\kCOLORABILITEeng{\cp{k-COLORABILITY}}}
\LANGUEinv{
\newcommand\COLORIAGE{\cp{BON~COLORIAGE}}}
{\newcommand\COLORIAGEeng{\cp{GOOD~COLORING}}}
\LANGUEinv{
\newcommand\tCOLORABILITE{\cp{3-COLORABILITE}}}
{\newcommand\tCOLORABILITEeng{\cp{3-COLORABILITY}}}
\LANGUEinv{
\newcommand\SUBSETSUM{\cp{SOMME~DE~SOUS~ENSEMBLE}}}
{\newcommand\SUBSETSUMeng{\cp{SUBSET~SUM}}}
\newcommand\PARTITION{\cp{PARTITION}}
\LANGUEinv{
\newcommand\NOMBREPREMIER{\decisionname{NOMBRE~ PREMIER}}}
{\newcommand\NOMBREPREMIEReng{\decisionname{PRIME~NUMBER}}}
\LANGUEinv{
\newcommand\CODAGE{\decisionname{CODAGE}}}
{\newcommand\CODAGEeng{\decisionname{ENCODING}}}
\newcommand\kSAT{\cp{k-SAT}}
\newcommand\tSAT{\cp{3-SAT}}
\newcommand\NAESAT{\cp{NAESAT}}
\newcommand\NAEtSAT{\cp{NAE3SAT}}
\newcommand\NAEqSAT{\cp{NAE4SAT}}
\newcommand\STABLE{\cp{\LANGUE{INDEPENDANT~SET}{STABLE}}}

%\newcommand\CM{\cp{COUPURE MAXIMALE}}
\LANGUEinv{
\newcommand\POSTpb{\cp{Problème~de~ la~correspondance~de~Post}}}
{\newcommand\POSTpbeng{\cp{Post~correspondence~problem}}}

\LANGUEinv{
\newcommand\PARITE{\cp{PARITE}}}
{\newcommand\PARITEeng{\cp{PARITY}}}


%-------------------------
%--- Modèles parallèles  ---
%--------------------------


% PARALLELISME
\newcommand\RAM{\cp{RAM}}
%
\newcommand\PRAM{\cp{PRAM}}
\newcommand\CRCW{\cp{CRCW}}
\newcommand\CREW{\cp{CREW}}
\newcommand\EREW{\cp{EREW}}


%------------------------
%--- codages  ---
%------------------------


\newcommand\codage[1]{\langle #1 \rangle}
\newcommand\tuple[1]{\codage{#1}}
\newcommand\paire[2]{\codage{#1,#2}}
\newcommand\triplet[3]{\codage{#1,#2,#3}}
\newcommand\quadruplet[4]{\codage{#1,#2,#3,#4}}
\newcommand\cinquplet[5]{\codage{#1,#2,#3,#4,#5}}


%------------------------
%--- descriptive set theory  ---
%------------------------

\newcommand\baire{\mathcal{N}}
\newcommand\peano{\overline{\mathcal{N}}} %% A MOI
\newcommand\cantor{\mathcal{C}}
%
\newcommand\bSigma{\mathbf{\Sigma}}
\newcommand\bPi{\mathbf{\Pi}}
\newcommand\bDelta{\mathbf{\Delta}}
%


%%%%%% ABOUT NOMS QUI CHANGES

\newcommand\WF{\cp{WF}}
\newcommand\WO{\WF}
\newcommand\LO{\cp{LO}}
\newcommand\DLO{\cp{DLO}}
% 
\newcommand\rk{\cp{rank}}


%% Moins propre:

\newcommand\I{\mathbb{I}}
\renewcommand\H{\mathbb{H}}
%
\newcommand\leKB{<_{KB}}
%
\newcommand\finiomega{{<\omega}}
\newcommand\compl{^\frown}
\newcommand\prefix{ ~ prefix~  } 
\newcommand\X{{X}}
\newcommand\Y{\mathcal{Y}}
\newcommand\subsetstrict{\subset}

\renewcommand\complement[1]{#1^{c}}


%------------------------
%--- schémas de récursion ---
%------------------------

%%                      Source papier MFCS 2019


% COMPLEXITY

%% MFCS differe de Clote:
%% Version à taille restreinte

\newcommand\co{\cp{co}}

\newcommand\mFPTIME{\mF_{PTIME}}
\newcommand{\mFPSPACE}{\mF_{SPACE}}
\newcommand{\mFPH}{\mF_{PH}}

\newcommand\SigmaP{\Sigma^P}
\newcommand\DeltaP{\Delta^{P}}
\newcommand\PiP{\Pi^{P}}
\newcommand\coSigmaP{\co\SigmaP}
\newcommand\coPiP{\mF{\co\PiP}}

%% autres classes:


%% Schemas pour complexité implicite

\newcommand\COMP{\cp{ COMP}}
\newcommand\CRN{\cp{ CRN}}
\newcommand\BRN{\cp{ BRN}}
\newcommand\SBRN{\cp{ SBRN}}
\newcommand\BR{\cp{ BR}}
\newcommand\kBR{k-\cp{ BR}}
\newcommand \WBPR{\cp{ WBPR}}
\newcommand \BMIN{\cp{ BMIN}}
\newcommand \BSUM{\cp{ BSUM}}
\newcommand\SBSUM{\cp{ SBSUM}}
\newcommand \BPROD{\cp{ BPROD}}
\newcommand \SBPROD{\cp{ SBPROD}}
\newcommand \SCOMP{\cp{ SCOMP}}
\newcommand \SR{\cp{ SR}}
\newcommand \SRN{\cp{ SRN}}

%% devrait etre défini
\newcommand\WBRN{\cp{ WBRN}}



%-- godel
\newcommand\INTDIV{\cp{INTDIV}}
\newcommand\DIV{\cp{DIV}}
\newcommand\EVEN{\cp{EVEN}}
\newcommand\ODD{\cp{ODD}}
\newcommand\PRIME{\cp{PRIME}}
\newcommand\POWER{\cp{POWER}}
\newcommand\BIT{\cp{BIT}}
\newcommand\negHP{\overline{\cp{HP}}}
\newcommand\HP{\cp{HP}}
\newcommand\negLuniv{\overline{\Luniv}}
\newcommand\VALCOMP{\cp{VALCOMP}}
\newcommand\mWINDOW{\cp{LEGAL}}
\newcommand\mmWINDOW{\cp{WINDOW}}
\newcommand\LENGTH{\cp{LENGTH}}
\newcommand\DIGIT{\cp{DIGIT}}
\newcommand\MATCH{\cp{MATCH}}
\newcommand\MOVE{\cp{MOVE}}
\newcommand\START{\cp{START}}
\newcommand\CELL{\cp{CELL}}
\newcommand\HALT{\cp{HALT}}
\newcommand\SUBST{\cp{SUBST}}
\newcommand\PROOF{\cp{PROOF}}
\newcommand\PROVABLE{\cp{PROVABLE}}
\newcommand\CONSIST{\cp{CONSIST}}



%---------------------
%--- calculus (AP) ---
%---------------------

% \jacobian{f}{x}
\newcommand{\jacobian}[1]{J_{#1}}
%\newcommand{\grad}[1]{\overrightarrow{\operatorname{grad}}~{#1}}
\newcommand{\grad}[1]{\nabla{#1}}
\newcommand{\scalarprod}[2]{{#1}\cdot{#2}}
\newcommand{\bigO}[1]{\mathcal{O}\left(#1\right)}
\newcommand\smallO{o}
\newcommand{\softO}[1]{\tilde{\mathcal{O}}\left(#1\right)}
\newcommand{\taylor}[3]{{T_{#1}^{#2}#3}}
\newcommand{\transpose}[1]{{#1}^T}
\newcommand{\pastsup}[2]{{\sup}_{#1}#2}

\DeclareMathOperator\arctanh{arctanh}



%---------------
%--- norms (AP) ---
%---------------


\newcommand{\inorm}[2]{\left\lVert{#1}\right\rVert_{#2}}
\newcommand{\twonorm}[1]{\inorm{#1}{2}}
\newcommand{\infnorm}[1]{\inorm{#1}{}}
\newcommand{\normsup}[1]{\infnorm{#1}{}}
%
\newcommand\hausdorff{d_H}
%
\newcommand\norm[1]{\infnorm{#1}}


%----------------
%--- topology ---
%----------------

% \openball{pt}{radius}
 \newcommand{\openball}[2]{B_{#2}(#1)}
\newcommand{\closedball}[2]{\ovl{B}_{#2}(#1)}
\newcommand{\closure}[1]{\ovl{#1}}
% \newcommand{\interior}[1]{\mathring{#1}}
% \newcommand{\border}[1]{\partial{#1}}




%-----------------
%--- functions (AP) ---
%-----------------
 
\newcommand{\lambdafun}[2]{(#1)\mapsto{#2}}
\newcommand{\lambdafunex}[2]{#1\mapsto{#2}}
\newcommand{\argmin}{\operatornamewithlimits{argmin}}
\newcommand{\argmax}{\operatornamewithlimits{argmax}}


 
%-----------------------
%--- Machines de Turing  ---
%-----------------------

\newcommand\blanc{\vectorl{B}}



%-----------------
%--- Logique ---
%-----------------
 
\newcommand\sem[1]{{[\semspace[}#1{]\semspace]}}
\newcommand\sems[1]{{[\semspace[}#1{]\semspace]}}
\newcommand\semss[1]{\sem{#1}_{\sigma'}}
\newcommand\semspace{\kern -.25ex}
\newcommand\true{true}
\newcommand\false{false}
%extension élémentaire:
\newcommand\elem{\preceq}
\newcommand\godel[1]{\lceil #1 \rceil}

%-----------------
%--- Model theory ---
%-----------------

\newcommand\hull[1]{\langle #1 \rangle}

%-----------------
%--- Ordinaux ---
%-----------------

\newcommand\omegaunCK{\omega_1^{CK}}
\newcommand\omegaunck{\omegaunCK}

\DeclareMathOperator{\cof}{cof}

%-----------------
%--- Surréels (QG) ---
%-----------------

\newcommand{\On}{\textbf{On}}
\newcommand{\Nobf}{\textbf{No}}
\newcommand{\Nobb}{\ensuremath{\mathbbmss{No}}}
\newcommand{\Nolt}[1]{\ensuremath{\Nobf_{< #1}}}
\newcommand{\Noleq}[1]{\ensuremath{\Nobf_{\leq #1}}}
\newcommand{\Noltbb}[1]{\ensuremath{\Nobb_{< #1}}}
\providecommand\No{\Nobf}
\newcommand\Noalpha{\Nolt{\alpha}}
\newcommand\Nolambda{\Nolt{\lambda}}

\DeclareMathOperator{\length}{length}
\DeclareMathOperator{\supp}{supp} % operateur support


%----------------------
%---Maths  de base (AP)  ---
%----------------------

\DeclareMathOperator{\dom}{dom} % operateur support

\DeclareMathOperator{\size}{size}

\newcommand{\powerset}[1]{\mathcal{P}(#1)}
\newcommand{\card}[1]{{\##1}}
\newcommand\priv{ - }
\newcommand\prive{-}
\renewcommand{\restriction}{\mathord{\upharpoonright}}
\newcommand\restrict[1]{_{\restriction #1}}

% \providecommand\mod{}
%% UNDEFINE:
\let\mod\relax
\let\div\relax
\DeclareMathOperator{\mod}{mod} 
\DeclareMathOperator{\div}{div} 

% partieentier
\newcommand\entieres[1]{\lceil #1 \rceil}
\newcommand\entierei[1]{\lfloor #1 \rfloor}


%----------------------
%---Symbole danger ---
%----------------------


%\providecommand*{\TakeFourierOrnament}[1]{{%
%\fontencoding{U}\fontfamily{futs}\selectfont\char#1}}
%\providecommand*{\danger}{{\Huge \TakeFourierOrnament{66}}}



%%%%%%%%%%%%%%%%%%%%%%%%%%%%%%%%%%%%%%%%%%%%%%%%%
%
%
%                      ASMeries
%
%
%%%%%%%%%%%%%%%%%%%%%%%%%%%%%%%%%%%%%%%%%%%%%%%%%



% CIRCUITS
%-- particuliers
\newcommand\REPLACE{\cp{REPLACE}}
\newcommand\EQUAL{\cp{EQUAL}}
\newcommand\EVALUE{\cp{EVALUE}}
\newcommand\TVALUE{\cp{TVALUE}}
\newcommand\UPDATE{\cp{UPDATE}}
\newcommand\ASSIGN{\cp{ASSIGN}}
\newcommand\PAR{\cp{PAR}}
\newcommand\DO{\cp{DO}}


% ASM
\newcommand\emplacement[2]{(#1,#2)}
\newcommand\contenu[2]{\sem{(#1,#2)}}
\newcommand\affecte[3]{(#1,#2)\leftarrow#3}
\newcommand\ctl{ctl\_state}


%%%%%%%%%%%%%%%%%%%%%%%%%%%%%%%%%%%%%%%%%%%%%%%%%
%
%
%                      GPACqueries
%
%
%%%%%%%%%%%%%%%%%%%%%%%%%%%%%%%%%%%%%%%%%%%%%%%%%



%------------------------
%--- generable zoo (GPAC) ---
%------------------------

\newcommand{\myop}[1]{\operatorname{#1}}
\newcommand{\abs}{\myop{abs}}
\newcommand{\sg}{\myop{sg}}
\newcommand{\ip}[1]{\myop{ip}_{#1}}
\newcommand{\rnd}{\myop{rnd}}
\newcommand{\lil}{\myop{lil}}
\newcommand{\lxh}{\myop{lxh}}
\newcommand{\hxl}{\myop{hxl}}
\newcommand{\plil}{\myop{plil}}
\newcommand{\reach}{\myop{reach}}
%\newcommand{\mreach}{\myop{mreach}}
\newcommand{\watchdog}{\myop{watchdog}}
\newcommand{\monitor}{\myop{monitor}}
%\newcommand{\norm}{\myop{norm}}
\newcommand{\mx}{\myop{mx}}
\newcommand{\mn}{\myop{mn}}
\newcommand{\sample}{\myop{sample}}
\newcommand{\select}{\myop{select}}
\newcommand{\ssine}[1]{\sin_{#1}}
\newcommand{\clamp}{\myop{clamp}}

\newcommand{\fnsg}[3]{tanh((#1)*(#2)*(#3))}
\newcommand{\fnipone}[3]{(1+\fnsg{(#1)-1}{#2}{#3})/2.}
\newcommand{\fnipinf}[1]{#1-sin(2*pi*(#1))/2./pi}%
\newcommand{\fnlil}[5]{%
\fnipone{(#3)-(#1)+1-((#2)-(#1))/8.}{#4+log(1+4./(#2-#1)*(#5)*(#5))+log(1+4)}{8./(#2-#1)}%
*\fnipone{#2-#3+1-(#2-#1)/8.}{#4+log(1+4/(#2-#1)*(#5)*(#5))+log(1+4)}{8./(#2-#1)}%
*4./(#2-#1)*(#5)}
\newcommand{\fnlxh}[5]{\fnipone{(#3)-(#1+#2)/2.+1}{#4+log(1+(#5)*(#5))}{2./(#2-#1)}*(#5)}
\newcommand{\fnhxl}[5]{\fnipone{(#1+#2)/2.-(#3)+1}{#4+log(1+(#5)*(#5))}{2./(#2-#1)}*(#5)}
\newcommand{\fnplilf}[4]{sin(((#4)-((#1)+(#2))/2.+(#3)/4.)*2.*pi/(#3))}
\newcommand{\fnplil}[6]{(#6)*\fnlxh{\fnplilf{#1}{#2}{#3}{#1}}%
{\fnplilf{#1}{#2}{#3}{#1+((#2)-(#1))/4.}}%
{\fnplilf{#1}{#2}{#3}{#4}}%
{#5+2.+log(1+(#6)*(#6))}%
{1./4.+2./((#2)-(#1))}}
\newcommand{\fnssine}[2]{sin(#2)*(1-cos(#2)/cos(#1))}



%--------------------
%--- complexity GPAC ---
%--------------------


%\ifthenelse{\equal{#1}{}}{Empty.}{Nonempty.}
\newcommand{\myclass}[1]{\operatorname{#1}}
% \newcommand{\gcomp}[2]{\ensuremath{\myclass{GCOMP}(#1,#2)}}
% \newcommand{\ggenerable}[1]{\textcolor{red}{\ensuremath{#1-}\text{generable}}}
 \newcommand{\gval}[2][]{\ensuremath{\myclass{GVAL}_{#1}\ifthenelse{\equal{#2}{}}{}{[#2]}}}
 \newcommand{\gpval}[1][]{\ensuremath{\myclass{GPVAL}_{#1}}}
 \newcommand{\gc}[3][]{\ensuremath{\myclass{ATSC}(#2,#3)}}
 \newcommand{\gpc}[1][]{\ensuremath{\myclass{ATSP}}}
\newcommand{\cgc}[1][]{\ensuremath{\myclass{ATS}}}
\newcommand{\gexpc}[1][]{\ensuremath{\myclass{AEXP}}}
\newcommand{\gwc}[2]{\ensuremath{\myclass{AW}(#1,#2)}}
\newcommand{\gpwc}{\ensuremath{\myclass{AWP}}}
\newcommand{\cgwc}{\ensuremath{\myclass{AW}}}
\newcommand{\grc}[3]{\ensuremath{\myclass{AR}(#1,#2,#3)}}
\newcommand{\gprc}{\ensuremath{\myclass{ARP}}}
\newcommand{\gsc}[3]{\ensuremath{\myclass{AS}(#1,#2,#3)}}
\newcommand{\gpsc}{\ensuremath{\myclass{ASP}}}
\newcommand{\goc}[3]{\ensuremath{\myclass{AO}(#1,#2,#3)}}
\newcommand{\gpoc}{\ensuremath{\myclass{AOP}}}
\newcommand{\cgoc}{\ensuremath{\myclass{AO}}}
\newcommand{\guc}[4]{\ensuremath{\myclass{AX}(#1,#2,#3,#4)}}
\newcommand{\gpuc}{\ensuremath{\myclass{AXP}}}
\newcommand{\glc}[2][]{\ensuremath{\myclass{AL}(#2)}}
\newcommand{\gplc}[1][]{\ensuremath{\myclass{ALP}}}
\newcommand{\cglc}[1][]{\ensuremath{\myclass{AL}}}

%\newcommand{\intinterv}[2]{\{#1,\ldots,#2\}}
\newcommand{\intinterv}[2]{\llbracket#1,#2\rrbracket}

\newcommand{\Rgen}{\R_{G}}
\newcommand{\Rpoly}{\R_{P}}

%%%%%%%%%%%%%%%%%%%%%%%%%%%%%%%%%%%%%%%%%%%%%%%%%
%
%
%                      Mes codages obscures
%
%
%%%%%%%%%%%%%%%%%%%%%%%%%%%%%%%%%%%%%%%%%%%%%%%%%




%%% Codages

% \newcommand\gammaoc{\gamma_{[0,1]}}
% \newcommand\gammaonc{\gamma_{\N}}
\newcommand\gammaTMc{\gamma^{TM}_{[0,1]}}
\newcommand\gammaTMcatan{\gamma^{TM}_{(0,1)^2}}
\newcommand\gammaTMnc{\gamma^{TM}_{\N^3}}
\newcommand\gammaTMncd{\gamma^{TM}_{\N^2}}
% \newcommand\gammaTMe{\gamma^{TM}_{*}}
% \newcommand\gammas{\gamma_{sab}}
% \newcommand\gammaord{\gamma_{cantor}}


\newcommand\enccompact{\nu_{X}}
 \newcommand\encinfinite{\nu_\N}
 \newcommand\enc{\nu}
 
%%%%%%%%%%%%%%%%%%%%%%%%%%%%%%%%%%%%%%%%%%%%%%%%%
%
%
%                      DESSINS ET IMAGES
%
%
%%%%%%%%%%%%%%%%%%%%%%%%%%%%%%%%%%%%%%%%%%%%%%%%%


\newcommand\IMAGE[2]{ \begin{center} 
    \includegraphics[#1]{#2}
  \end{center}
  }


  
%%%% POUR FAIRE DESSIN.
\newenvironment{dessinpp}[1]{
\begin{tikzpicture}[baseline=(current bounding
      box.west),#1]
}
{\end{tikzpicture}}




\newenvironment{dessinp}[1]{
\begin{center}\begin{tikzpicture}[baseline=(current bounding
      box.north),#1]
}
{\end{tikzpicture}
\end{center}}


\newenvironment{dessin}{
\begin{center}\begin{tikzpicture}[baseline=(current bounding
      box.north)]
}
{\end{tikzpicture}
\end{center}}



\newenvironment{dessind}{
\begin{tikzpicture}[baseline=(current bounding
      box.north)]
}
{\end{tikzpicture}
}



%%%%%%%%%%%%%%%%%%%%%%%%%%%%%%%%%%%%%%%%%%%%%%%%%
%
%
%                      POUR DESSINS GPAC
%
%
%%%%%%%%%%%%%%%%%%%%%%%%%%%%%%%%%%%%%%%%%%%%%%%%%


%%%% MACROS DE BASE.


\newcommand\constant[2]
{
node (#1)  [shape=circle,draw] {#2};
\draw (#1.east) -- +(0.3,0) coordinate (#1-o) ; 
\draw (#1.west) -- +(-0.3,0) coordinate (#1-i) ; 
}

\newcommand\basicproduct[2]
{
node (#1) {#2};
\draw (#1) +(-0.4,0.4) -- +(0.4,0.4) -- +(0.8,0) -- +(0.4,-0.4) -- +(-0.4,-0.4) --
+(-0.4,0.4);
\draw (#1) ++(0.8,0) -- ++(+0.3,0) coordinate (#1-o) ;
\draw (#1) ++ (-0.4,0.2) -- +(-0.3,0) coordinate (#1-i1) ;
\draw (#1) ++(-0.4,-0.2) -- +(-0.3,0) coordinate (#1-i2) ;
}

\newcommand\product[1]{\basicproduct{#1}{$+\Pi$}}
\newcommand\negativeproduct[1]{\basicproduct{#1}{$-\Pi$}}


%%BEGIN INTERNAL
\newcommand\basicsummer[1]{
coordinate (#1) {};
\draw (#1) +(-0.5,0.6) -- +(0.5,0) -- +(-0.5,-0.6) -- +(-0.5,0.6);
\draw (#1) ++(0.5,0) -- +(+0.3,0) coordinate (#1-o) ;
}

\newcommand\basicintegrator[1]{
coordinate (#1) {};
\draw (#1) +(-0.5,0.6) -- +(0.5,0) -- +(-0.5,-0.6) -- +(-0.5,0.6);
\draw (#1) ++(0.5,0) -- +(+0.3,0) coordinate (#1-o) ;
\draw (#1) +(-0.5,-0.6) -| +(-0.8,0.6) -- +(-0.5,0.6) ;
 \draw (#1) ++(-0.65,0.6) -- +(0,+0.3) coordinate (#1-init) ;
}

\newcommand\splitfour[2]{
\draw (#2) ++ (#1,0.4) -- +(-0.3,0) coordinate (#2-i1) ;
\draw (#2) ++ (#1,0.15) -- +(-0.3,0) coordinate (#2-i2) ;
\draw (#2) ++ (#1,-0.15) -- +(-0.3,0) coordinate (#2-i3) ;
\draw (#2) ++ (#1,-0.4) -- +(-0.3,0) coordinate (#2-i4) ;
}

\newcommand\splitthree[2]{
\draw  (#2) ++(#1,0.3) -- +(-0.3,0) coordinate (#2-i1) ;
\draw  (#2) ++(#1,0) -- +(-0.3,0) coordinate (#2-i2) ;
\draw  (#2) ++(#1,-0.3) -- +(-0.3,0) coordinate (#2-i3) ;
}

\newcommand\splittwo[2]{
\draw  (#2) ++(#1,0.2) -- +(-0.3,0) coordinate (#2-i1) ;
\draw  (#2) ++(#1,-0.2) -- +(-0.3,0) coordinate (#2-i2) ;
}

\newcommand\splitone[2]{
\draw (#2) ++(#1,0) -- +(-0.3,0) coordinate (#2-i1) ;
}

%% END INTERNAL

\newcommand\summerfour[1]{
\basicsummer{#1}
\splitfour{-0.5}{#1}
}

\newcommand\summerthree[1]{
\basicsummer{#1}
\splitthree{-0.5}{#1}
}

\newcommand\summertwo[1]{
\basicsummer{#1}
\splittwo{-0.5}{#1}
}

\newcommand\summerone[1]{
\basicsummer{#1}
\splitone{-0.5}{#1}
}


\newcommand\integratorfour[1]{
\basicintegrator{#1}
\splitfour{-0.8}{#1}
<}


\newcommand\integratorthree[1]{
\basicintegrator{#1}
\splitthree{-0.8}{#1}
}


\newcommand\integratortwo[1]{
\basicintegrator{#1}
\splittwo{-0.8}{#1}
}


\newcommand\integratorone[1]{
\basicintegrator{#1}
\splitone{-0.8}{#1}
}

%%%%% FIN MACROS




%%%%%%%%%%%%%%%%%%%%%%%%%%%%%%%%%%%%%%%%%%%%%%%%%
%
%
%                      Trucs d'Amaury
%
%
%%%%%%%%%%%%%%%%%%%%%%%%%%%%%%%%%%%%%%%%%%%%%%%%%


%----------------
%--- ref (AP) ---
%----------------

% \newcommand{\defref}[1]{Definition~\ref{#1}}
\newcommand{\lemref}[1]{Lemma~\ref{#1}}
\newcommand{\thref}[1]{Theorem~\ref{#1}}
\newcommand{\amaurycorref}[1]{Corollary~\ref{#1}}
 \newcommand{\figref}[1]{Figure~\ref{#1}}
% \newcommand{\figrefmult}[1]{Figures~{#1}}
% \newcommand{\secref}[1]{Section~\ref{#1}}
% %\renewcommand{\algref}[1]{Algorithm~\ref{#1}}
 \newcommand{\propref}[1]{Proposition~\ref{#1}}
% \newcommand{\exref}[1]{Example~\ref{#1}}
\newcommand{\remref}[1]{Remark~\ref{#1}}


 %-------------------
%--- polynomials ---
%-------------------

\newcommand{\sigmap}[1]{{\Sigma{#1}}}
\newcommand{\degp}[1]{{\operatorname{deg}(#1)}}
\newcommand{\poly}{\operatorname{poly}}


 %----------------------
%--- misc functions ---
%----------------------

\newcommand{\sgn}[1]{\operatorname{sgn}(#1)}
\newcommand{\intp}{\operatorname{int}}
 \newcommand{\fracp}{\operatorname{frac}}
\newcommand{\fiter}[2]{#1^{[#2]}}
\newcommand{\frestrict}[2]{#1\restriction_{#2}}
\newcommand{\indicator}[1]{\mathds{1}_{#1}}
\newcommand{\idfun}{\operatorname{id}}
% \newcommand{\floor}[1]{\left\lfloor#1\right\rfloor}
 \newcommand{\ceil}[1]{\left\lceil#1\right\rceil}
\newcommand{\round}[1]{\left\lfloor#1\right\rceil}
\DeclareMathOperator\erf{erf}


%------------
%--- sets ---
%------------

\newcommand{\A}{\mathbb{A}}
%\newcommand{\D}{\mathbb{D}}
\newcommand{\Ns}{\mathbb{N}^*}
%\newcommand{\C}{\mathbb{C}}
%\newcommand{\K}{\mathbb{K}}
% \MAT{height}{[width]}{[field]}
\newcommand{\MAT}[3]{M_{#1\ifthenelse{\equal{#2}{}}{}{,#2}}\ifthenelse{\equal{#3}{}}{}{\left(#3\right)}}

%\newcommand{\Rcomp}{\R_{c}}
\newcommand{\Rp}{\R_{+}}
\newcommand{\Rs}{\R^{*}}
\newcommand{\Rps}{\R_{+}^{*}}

%\newcommand{\Analytic}{C^\omega}
%\newcommand{\PTIME}{\ensuremath{\operatorname{P}}}
%\newcommand{\EXPTIME}{\ensuremath{\operatorname{EXPTIME}}}
%\newcommand{\NP}{\ensuremath{\operatorname{NP}}}
%\newcommand{\NPTIME}{\NP}
\newcommand{\FP}{\ensuremath{\operatorname{FP}}}
%\newcommand{\PSPACE}{\ensuremath{\operatorname{PSPACE}}}
% \newcommand{\FPSPACE}{\ensuremath{\operatorname{FPSPACE}}}



%%%%%%%%%%%%%%%%%%%%%%%%%%%%%%%%%%%%%%%%%%%%%%%%%
%
%
%                      Trucs de Quentin
%
%
%%%%%%%%%%%%%%%%%%%%%%%%%%%%%%%%%%%%%%%%%%%%%%%%%


\newcommand{\fct}[5]{\ensuremath{#1 : \left\{\begin{array}{c c c}
#2 & \rightarrow & #3\\
#4 & \mapsto & #5
                                             \end{array}\right.}}

%Nouvelle fraction, plus jolie
\newcommand{\f}[2]{	
	\mathchoice%
		{\dfrac{#1}{#2}}
    	{\dfrac{#1}{#2}}
		{\frac{#1}{#2}}
		{\frac{#1}{#2}}
}


%%Ensembles et combinatoires

%Ensemble avec propriété à drotie \{x | blabla\}
\newcommand{\enstq}[2]{\left\{ \left.#1\ \vphantom{#2}\right|\ #2  \right\}}
%meme chose mais avec des crochets
\newcommand{\crotq}[2]{\left[ \left.#1\ \vphantom{#2}\right|\ #2  \right]}
%même chose mais avec des brackets
\newcommand{\bratq}[2]{\left\langle \left.#1\ \vphantom{#2}\right|\ #2  
  \right\rangle}

\newcommand{\intn}[2]{\ensuremath{
		\left\llbracket\, #1\ ;\ #2\,\right\rrbracket
              }}


%%%%%%%%%%%%%%%%%%%%%%%%%%%%%%%%%%%%%%%%%%%%%%%%%
%
%
%               ENVIRONNEMENTS THEOREMS & CO
%
%
%%%%%%%%%%%%%%%%%%%%%%%%%%%%%%%%%%%%%%%%%%%%%%%%%


%----------------
%--- pas dans lipics ---
%-----------------


\ifnotSLIDE{
% commente car bug Fevrier 2026: \ifnotTDEXAM{\spnewtheorem{solution}{Solution}{\bfseries}{\rmfamily}}
\spnewtheorem{openproblem}{Open Problem}{\bfseries}{\rmfamily}
\spnewtheorem{remarkolivier}{Commentaire d'Olivier: }{\bfseries}{\rmfamily}
\spnewtheorem{postulate}{Postulate}{\bfseries}{\rmfamily}
}

%----------------
%--- dans lipics ---
%-----------------
%
\newenvironment{exercice}[1]{\begin{exercise} \label{exo:#1-stt}}{\end{exercise}}
\newenvironment{exerciced}[1]{\begin{exercisee}\label{exo:#1-stt}}{\end{exercisee}}
\newenvironment{exercicec}[1]{\begin{exercise} \label{exo:#1-stt} (\LANGUE{solution on}{corrigé} page \pageref{exo:#1-cor})
}{\end{exercise}}
\newenvironment{exercicedc}[1]{\begin{exercisee} \label{exo:#1-stt} (\LANGUE{solution on}{corrigé} page \pageref{exo:#1-cor})
}{\end{exercisee}}

\newenvironment{slideproposition}{{\bf Proposition}}{}

\ifnotSLIDE{
\nolipics{
\nolncs{\spnewtheorem{theorem}{\LANGUE{Theorem}{Théorème}}{\bfseries}{\rmfamily}}
\nolncs{\spnewtheoremi{definition}{\LANGUE{Definition}{Définition}}{\bfseries}{\rmfamily}{theorem}}
\nolncs{\spnewtheoremi{lemma}{\LANGUE{Lemma}{Lemme}}{\bfseries}{\rmfamily}{theorem}}
\nolncs{\spnewtheoremi{corollary}{\LANGUE{Corollary}{Corollaire}}{\bfseries}{\rmfamily}{theorem}}
\nolncs{\spnewtheoremi{proposition}{Proposition}{\bfseries}{\rmfamily}{theorem}}
\nolncs{\spnewtheoremi{example}{\LANGUE{Example}{Exemple}}{\bfseries}{\rmfamily}{theorem}}
%\spnewtheorem{sketch}{{\bf \LANGUE{Sketch of proof}{Démonstration (principe)}}}{\bfseries}{\rmfamily}
\nolncs{\spnewtheoremi{remark}{\LANGUE{Remark}{Remarque}}{\bfseries}{\rmfamily}{theorem}}
\nolncs{\spnewtheorem{exercise}{\LANGUE{Exercise}{Exercice}}{\bfseries}{\rmfamily}}
\spnewtheorem{exercisee}{\LANGUE{*Exercise}{*Exercice}}{\bfseries}{\rmfamily}
\spnewtheorem{summary}{\LANGUE{Summary}{Résumé}}{\bfseries}{\rmfamily}
 %\spnewtheorem{exerciceenglish}{Exercise}{\bfseries}{\rmfamily}
 %
 \newenvironment{sketch}{{\bf \LANGUE{Sketch of proof}{Démonstration (principe)}}:}{\hfill $\Box$ }
\nolncs{\spnewtheorem{conjecture}{\LANGUE{Conjecture}{Conjecture}}{\bfseries}{\rmfamily}}
 %%
 \spnewtheorem{theoremaprouver}{\LANGUE{Theorem TO BE PROVED}{Théorème (A PROUVER)}}{\bfseries}{\rmfamily}
  \spnewtheorem{conjectureaprouver}{\LANGUE{Conjecture TO BE VERIFIED}{Conjecture (A VERIFIER)}}{\bfseries}{\rmfamily}
 }}
            
 \makeatletter
 \ifnotSLIDE{
  \ifnotTDEXAM{
\@ifundefined{proof}{
	\newenvironment{proof}{{\bf \LANGUE{Proof}{Démonstration}}:}{\hfill $\Box$ 
	}{}
	}}}
\makeatother

 
%%%%%%%%%%%%%%%%%%%%%%%%%%%%%%%%%%%%%%%%%%%%%%%%%
%
%
%                      Macros dont pas fier du tout
%
%
%%%%%%%%%%%%%%%%%%%%%%%%%%%%%%%%%%%%%%%%%%%%%%%%%



\newcommand\equations[1]{
\begin{array}{lll}
#1
\end{array}
}



\newcommand\systeme[1]{
\left\{
\begin{array}{lll} 
#1
\end{array}
\right.
}



%------------------
%--- shorthands (AP)---
%------------------
\newcommand{\mtt}[1]{\mathtt{#1}}
\newcommand{\ovl}[1]{\overline{#1}}
\newcommand{\panic}[1]{\textcolor{red}{#1}}
%\NewEnviron{epanic}{{\color{red}\BODY}}
\newcommand{\idest}{\emph{i.e.\thinspace}}


%------------------
%--- mes trucs  ---
%------------------


\newcommand\implique{\Rightarrow}
\newcommand\ssi{\Leftrightarrow}
\newcommand\ac[1]{\{#1\}}
\newcommand\zero{\vectorl{0}}
\newcommand\un{\vectorl{1}}

\newcommand\technique[1]{{\scriptsize #1}}


%% Preuve and Co

\newcommand\sensdirect{ Direction $\Rightarrow$. }
\newcommand\sensindirect{ Direction $\Leftarrow$. }




\newcommand\donnee{\mbox{<donnée>}}



% pour aller vite

\newcommand\gM{\mathfrak{M}}
\newcommand\gMM{(M,f_1,\cdots,f_u,r_1,\cdots,r_v)}
%
\newcommand\mygM{\mK}
\newcommand\mygMM{\left ( \K, \{op_i\}_{i\in I}, r_1 \ldots,
r_\ell, \zero,\un \right )}
\newcommand\myM{\K} 

\newcommand\gMMM{(M,f_1,\cdots,f_u,c_1,\cdots,c_k,r_1,\cdots,r_v)}
\newcommand\gMMMM{(f_1,\cdots,f_u,r_1,\cdots,r_v)}
\newcommand\gMME{(M,f_1,\cdots,f_u,f'_1,\cdots,f'_\ell,r_1,\cdots,r_v)}
\newcommand\gMMS{(f_1,\cdots,f_u,f'_1,\cdots,f'_\ell,r_1,\cdots,r_v)}

\newcommand\bool{\mathfrak{B}}
\newcommand\gbool{(\{0,1\},\zero,\un,=)}

\newcommand\osg{\overline{sg}}
\newcommand\moins{\mathrel{\dot{-}}}


%%% TRUC DE BARWISE


\newcommand\UM{\kU_{\kM}}
\newcommand\VV{\mathbb{V}}
% Hereditarily finite
\newcommand\IHF{\mathbb{I}HF}
\newcommand\IHP{\IHYP}
\newcommand\IHYP{\mathbb{I}HYP}

\providecommand\citep[1]{\cite{#1}}

%% Cofinali


% probas
% défini dans amsmath
\renewcommand\Pr{\cp{Pr}}
\newcommand\Var{\cp{Var}}


% -- mesures
%\newcommand\card{card}
\newcommand\taille{\LANGUE{size}{taille}}
\newcommand\entree{e}
\newcommand\profondeur{\LANGUE{depth}{profondeur}}


%%%%%%%%%%%%%%%%%%%%%%%%%%%%%%%%%%%%%%%%%%%%%%%%%%%%%%%%%%%%%%%%%%%%%%%%%%%%
%                 MACHINES DE TURING (graphique)
%%%%%%%%%%%%%%%%%%%%%%%%%%%%%%%%%%%%%%%%%%%%%%%%%%%%%%%%%%%%%%%%%%%%%%%%%%%%



% Configuration de machine de Turing
\newcommand\configTuring[3]{#1\textbf{#2}#3}


%%% DESSIN CROIX PORU AIDER  A DESSINER

\newcommand\croix{  +(-1,-1) -- +(1,1) +(-1,1) --
  +(1,-1) }

%% DESSIN ACCOLADE

\newcommand\ACCOLADE[3]{
%#1: coordonnées de départ
%#2: coordonnées arrivée
%#3: texte
\draw[decorate,decoration={brace}]  (#1) -- (#2)
 node[midway,above=0.1cm] {#3};
}

%%%


\def\lcase{3.5 ex}
\def\hcase{3.5ex}
\tikzstyle{case} = []
%[draw,minimum width = \lcase,minimum height = 3.6ex,baseline]
\newcommand\dcase{ +(-0.5*\lcase,-0.5*\hcase) rectangle +(0.5*\lcase,0.5*\hcase)}
%\newcommand\lettre[1]{\node [name=l,case,on chain=1] {#1};}


\newcommand\RUBANMACHINEDETURING[6]{
%#1: texte
%#2: position de la tête
%#3: nb de blancs à gauche
%#4: nb de cases totales
%#5: texte au dessus du ruban
%#6: position du coin
    \draw (#6) +(\lcase,4ex ) node {#5};
    \draw (#6) +(-0.6,0)  node  {\ldots};  
    \draw (#6) node[coordinate] (t) {};
    \foreach \x in {1,2,...,#3} {
         \draw (t) node [case]   {$\blanc$} \dcase;
         \draw (t) ++(\lcase,0) node [coordinate] (t) {};
       };
   \StrLen{#1}[\Resultat] % ATTENTION IL FAUT CETTE SYNTAXE POUR FAIRE
                         % EQUIVALENT D UN EDEF
  
\foreach \x in {1,...,\Resultat} 
             { 
        \draw (t) node [case]   {\StrChar{#1}{\x}} \dcase;
        \draw (t) ++(\lcase,0) node [coordinate] (t) {};
             }
\pgfmathtruncatemacro{\z}{#4 - \Resultat - #3}

  \foreach \x in {1,2,..., \z} {
        \draw (t) node [case]   {$\blanc$} \dcase;
        \draw (t) ++(\lcase,0) node [coordinate] (t) {};
}
%% Convention: graphiquement, la case numéro n se trouve à (n*\lcase,0)
%% en comptant les case à partir de $0$

   \draw  (t) +(0.1,0) node [name=r] {\ldots}; 
}

\newcommand\RUBANMACHINEDETURINGtexte[6]{
%#1: texte
%#2: position de la tête
%#3: nb de blancs à gauche
%#4: nb de cases totales (équivalent)
%#5: texte au dessus du ruban
%#6: position du coin
    \draw (#6) +(\lcase,4ex ) node {#5};
    \draw (#6) +(-0.6,0)  node  {\ldots};  
    \draw (#6) node[coordinate] (t) {};
   \foreach \x in {1,2,...,#3} {
        \draw (t) node [case]   {$\blanc$} \dcase;
        \draw (t) ++(\lcase,0) node [coordinate] (t) {};
      };
   \draw (#6) ++(-2*\lcase,0.5*\lcase)-- +(\lcase*#4,0);
   \draw (#6) ++(-2*\lcase,-0.5*\lcase) -- +(\lcase*#4,0);
   
    \draw (t) ++(-0.5*\lcase,0) node [case,anchor=west]  (te) {#1};
    \draw (te.east)  node [anchor=west,coordinate] (t) {};

  \draw (t) ++(0.2*\lcase,0) node [coordinate] (t) {};
   \foreach \x in {1,2,...,#3} {
        \draw (t) node [case]   {$\blanc$} \dcase;
        \draw (t) ++(\lcase,0) node [coordinate] (t) {};
      };

    \draw  (t) +(0.1,0) node [anchor=west] {\ldots}; 
}


\newcommand\RUBANMACHINEDETURINGTETE[6]{
%
% dessine la tete
%
% output: (tete)
%
\RUBANMACHINEDETURING{#1}{#2}{#3}{#4}{#5}{#6}
% TETE
     \draw (#6) + ( #2 * \lcase + #3  *\lcase,- 0.5*\lcase )
     node[coordinate] (k) {};
     \draw[->] (k) +(0,-0.5) -- (k);
     \draw (k) +(0,-0.5) node[coordinate] (tete) {};
}


\newcommand\RUBANMACHINEDETURINGTETEtexte[6]{
%
% dessine la tete
%
% output: (tete)
%
\RUBANMACHINEDETURINGtexte{#1}{#2}{#3}{#4}{#5}{#6}
% TETE
     \draw (#6) + ( #2 * \lcase + #3  *\lcase,- 0.5*\lcase )
     node[coordinate] (k) {};
     \draw[->] (k) +(0,-0.5) -- (k);
     \draw (k) +(0,-0.5) node[coordinate] (tete) {};
}


\newcommand\MACHINEDETURING[6]{
%#1: texte
%#2: position de la tête
%#3: état
%#4: nb de blancs à gauche
%#5: nb de cases totales
%#6: texte en haut du ruban
%\begin{tikzpicture}[node distance=-0.15mm]

\RUBANMACHINEDETURINGTETE{#1}{#2}{#4}{#5}{#6}{0,0}
%
%% ETAT:
     \draw (tete) node [name=etat,draw,circle,anchor=north]  {#3};
     \draw (etat) -- (tete);
%\end{tikzpicture}
}

\newcommand\MACHINEDETURINGtexte[6]{
%#1: texte
%#2: position de la tête
%#3: état
%#4: nb de blancs à gauche
%#5: nb de cases totales
%#6: texte en haut du ruban
%\begin{tikzpicture}[node distance=-0.15mm]

\RUBANMACHINEDETURINGTETEtexte{#1}{#2}{#4}{#5}{#6}{0,0}
%
%% ETAT:
     \draw (tete) node [name=etat,draw,circle,anchor=north]  {#3};
     \draw (etat) -- (tete);
%\end{tikzpicture}
}

\newcommand\MACHINEDETURINGTROISRUBANS[9]{
%#1: texte1
%#2: position de la tête 1
%#3: texte2
%#4: position de la tête 2
%#5: texte3
%#6: position de la tête 3
%#7: état
%#8: nb de blancs à gauche 1
%#9: nb de cases totales
%\begin{tikzpicture}[node distance=-0.15mm]

\RUBANMACHINEDETURINGTETE{#1}{#2}{#8}{#9}{\LANGUE{tape}{ruban} $1$}{0,0}
\draw (tete) node[coordinate] (tete1) {};
\draw (tete1) +(8,0) node[coordinate] (tete1b) {};
\draw (tete1) -| (tete1b) ;

\RUBANMACHINEDETURINGTETE{#3}{#4}{#8}{#9}{\LANGUE{tape}{ruban} $2$}{0,-2}
\draw (tete) node[coordinate] (tete2) {};
\draw (tete2) +(-8.5,0) node[coordinate] (tete2b) {};
\draw (tete2) -| (tete2b) ;

\RUBANMACHINEDETURINGTETE{#5}{#6}{#8}{#9}{\LANGUE{tape}{ruban} $3$}{0,-4}
\draw (tete) node[coordinate] (tete3) {};

%% ETAT:
     \draw (tete3) node [name=etat,draw,circle,anchor=north]  {#7};
     \draw (etat) -- (tete3);
    \draw  (etat) -|  (tete2b);
\draw  (etat) -|  (tete1b);
%\end{tikzpicture}
}

\newcommand\MACHINEDETURINGTROISRUBANSENGLISH[9]{
%#1: texte1
%#2: position de la tête 1
%#3: texte2
%#4: position de la tête 2
%#5: texte3
%#6: position de la tête 3
%#7: état
%#8: nb de blancs à gauche 1
%#9: nb de cases totales
%\begin{tikzpicture}[node distance=-0.15mm]

\RUBANMACHINEDETURINGTETE{#1}{#2}{#8}{#9}{tape $1$}{0,0}
\draw (tete) node[coordinate] (tete1) {};
\draw (tete1) +(8,0) node[coordinate] (tete1b) {};
\draw (tete1) -| (tete1b) ;

\RUBANMACHINEDETURINGTETE{#3}{#4}{#8}{#9}{tape $2$}{0,-2}
\draw (tete) node[coordinate] (tete2) {};
\draw (tete2) +(-8.5,0) node[coordinate] (tete2b) {};
\draw (tete2) -| (tete2b) ;

\RUBANMACHINEDETURINGTETE{#5}{#6}{#8}{#9}{tape $3$}{0,-4}
\draw (tete) node[coordinate] (tete3) {};

%% ETAT:
     \draw (tete3) node [name=etat,draw,circle,anchor=north]  {#7};
     \draw (etat) -- (tete3);
    \draw  (etat) -|  (tete2b);
\draw  (etat) -|  (tete1b);
%\end{tikzpicture}
}



\newcommand\MACHINEDETURINGTROISRUBANStexte[9]{
%#1: texte1
%#2: position de la tête 1
%#3: texte2
%#4: position de la tête 2
%#5: texte3
%#6: position de la tête 3
%#7: état
%#8: nb de blancs à gauche 1
%#9: nb de cases totales
% \begin{tikzpicture}[node distance=-0.15mm]

\RUBANMACHINEDETURINGTETEtexte{#1}{#2}{#8}{#9}{\LANGUE{tape}{ruban} $1$}{0,0}
\draw (tete) node[coordinate] (tete1) {};
\draw (tete1) +(8,0) node[coordinate] (tete1b) {};
\draw (tete1) -| (tete1b) ;

\RUBANMACHINEDETURINGTETEtexte{#3}{#4}{#8}{#9}{\LANGUE{tape}{ruban} $2$}{0,-2}
\draw (tete) node[coordinate] (tete2) {};
\draw (tete2) +(-8.5,0) node[coordinate] (tete2b) {};
\draw (tete2) -| (tete2b) ;

\RUBANMACHINEDETURINGTETEtexte{#5}{#6}{#8}{#9}{\LANGUE{tape}{ruban} $3$}{0,-4}
\draw (tete) node[coordinate] (tete3) {};

%% ETAT:
     \draw (tete3) node [name=etat,draw,circle,anchor=north]  {#7};
     \draw (etat) -- (tete3);
    \draw  (etat) -|  (tete2b);
\draw  (etat) -|  (tete1b);
%\end{tikzpicture}
}


\newcommand\MACHINEDETURINGDEUXRUBANStexte[9]{
%#1: texte1  
%#2: position de la tête 1 
%#3: texte2 inutilisée
%#4: position de la tête 2 inutilisée
%#5: texte3
%#6: position de la tête 3
%#7: état
%#8: nb de blancs à gauche 1
%#9: nb de cases totales
% \begin{tikzpicture}[node distance=-0.15mm]

\RUBANMACHINEDETURINGTETEtexte{#1}{#2}{#8}{#9}{\LANGUE{tape}{ruban} $1$}{0,0}
\draw (tete) node[coordinate] (tete1) {};
\draw (tete1) +(8,0) node[coordinate] (tete1b) {};
\draw (tete1) -| (tete1b) ;

% \RUBANMACHINEDETURINGTETEtexte{#3}{#4}{#8}{#9}{\LANGUE{tape}{ruban} $2$}{0,-1.5}
% \draw (tete) node[coordinate] (tete2) {};
% \draw (tete2) +(-5.5,0) node[coordinate] (tete2b) {};
% \draw (tete2) -| (tete2b) ;

\RUBANMACHINEDETURINGTETEtexte{#5}{#6}{#8}{#9}{\LANGUE{tape}{ruban} $2$}{0,-1.5}
\draw (tete) node[coordinate] (tete3) {};

%% ETAT:
     \draw (tete3) node [name=etat,draw,circle,anchor=north]  {#7};
     \draw (etat) -- (tete3);
%   \draw  (etat) -|  (tete2b);
\draw  (etat) -|  (tete1b);
%\end{tikzpicture}
}


\def\argdix{2}
\def\argonze{25}
\newcommand\MACHINEDETURINGQUATRERUBANStextep[9]{
%#1: texte1
%#2: position de la tête 1
%#3: texte2
%#4: position de la tête 2
%#5: texte3
%#6: position de la tête 3
%#7: texte4
%#8: position de la tête 3
%#9: état
%#10: nb de blancs à gauche 1
%#11: nb de cases totales
% \begin{tikzpicture}[node distance=-0.15mm]

%% BUG sur #10 devenu \argdix
%% BUG SUR #11 devenu \argonze

\RUBANMACHINEDETURINGTETEtexte{#1}{#2}{\argdix}{\argonze}{\LANGUE{tape}{ruban} $1$}{0,0}
% \draw (tete) node[coordinate] (tete1) {};
% \draw (tete1) +(8,0) node[coordinate] (tete1b) {};
% \draw (tete1) -| (tete1b) ;

\RUBANMACHINEDETURINGTETEtexte{#3}{#4}{\argdix}{\argonze}{\LANGUE{tape}{ruban} $2$}{0,-1.5}
% \draw (tete) node[coordinate] (tete2) {};
% \draw (tete2) +(-5.5,0) node[coordinate] (tete2b) {};
% \draw (tete2) -| (tete2b) ;

\RUBANMACHINEDETURINGTETEtexte{#5}{#6}{\argdix}{\argonze}{\LANGUE{tape}{ruban} $3$}{0,-3.2}
\draw (tete) node[coordinate] (tete3) {};

\RUBANMACHINEDETURINGTETEtexte{#7}{#8}{\argdix}{\argonze}{\LANGUE{tape}{ruban} $4$}{0,-4.9}
\draw (tete) node[coordinate] (tete4) {};

%% ETAT:
     \draw (tete4) node [name=etat,draw,circle,anchor=north]  {#9};
%      \draw (etat) -- (tete3);
%     \draw  (etat) -|  (tete2b);
% \draw  (etat) -|  (tete1b);
%\end{tikzpicture}
}


%%%%%%%%%%%%%%%%%%%%%%%%%%%%%%%%%%%%%%%%%%%%%%%%%%%%%%%%%%%%%%%%%%%%%%%%%%%%
%                 MACHINES DE TURING (automate)
%%%%%%%%%%%%%%%%%%%%%%%%%%%%%%%%%%%%%%%%%%%%%%%%%%%%%%%%%%%%%%%%%%%%%%%%%%%%



\newcommand\dec{1ex}
\newcommand\myACCOLADE[4]{
\ACCOLADE{#1*\lcase+#4*\lcase-0.5*\lcase,0.5*\lcase+\dec}{#1*\lcase+#4*\lcase-0.5*\lcase+#3*\lcase,0.5*\lcase+\dec}{#2}
}


\newcommand\M{\Sigma}


\newcommand\sommet[2]{\node[state] (#1) #2 {$#1$};}
\newcommand\sommetinitial[2]{\node[state,initial] (#1) #2 {$#1$};}
\newcommand\sommetfinal[2]{\node[state,accepting] (#1) #2 {$#1$};}
\newcommand\gtransition[5]{\draw (#1) edge node {#2/#4 $#5$} (#3);}
\newcommand\gtransitionba[5]{\draw (#1) [loop above] edge node {#2/#4
    $#5$} (#3);}
\newcommand\gtransitionbb[5]{\draw (#1) [loop below] edge node {#2/#4
    $#5$} (#3);}
\newcommand\gtransitionbleft[5]{\draw (#1) [loop left] edge node {#2/#4
    $#5$} (#3);}
\newcommand\gtransitionbright[5]{\draw (#1) [loop right] edge node {#2/#4
    $#5$} (#3);}
\newcommand\gtransitionbl[5]{\draw (#1) [bend left] edge node {#2/#4
    $#5$} (#3);}
\newcommand\gtransitionbr[5]{\draw (#1) [bend right] edge node {#2/#4
    $#5$} (#3);}
\newcommand\gtransitionbadouble[9]{\draw (#1) [loop above] edge node {
\begin{tabular}{l}
#2/#4 $#5$ \\
  #7/#8 $#9$ \\
\end{tabular}
} (#3);}

\newcommand\SCALE{}

\newcommand\dessinmtp[2]{
\begin{center}
\begin{tikzpicture}[->,>=stealth',shorten >=1pt,auto,node distance=2.3cm,semithick,#2] 
#1
\end{tikzpicture}
\end{center}
}

\newcommand\dessinmt[1]{\dessinmtp{#1}{scale=1}}


\newcommand\ANIMATIONMT[1]{
\begin{overprint}
\begin{dessinp}{scale=0.52}
\foreach \av/\q/\b  [count=\xi]
in  
{#1}
{ 
\only<\xi| handout 0>{ 
 \StrLen{\av}[\Resintm] % ATTENTION IL FAUT CETTE SYNTAXE POUR FAIRE
                         % EQUIVALENT D UN EDEF
 \MACHINEDETURING{\av \b}{\Resintm}{$\q$}{5}{20}{}
}}
 \end{dessinp}
\end{overprint}
}

\newcommand\ANIMATIONMTd[1]{
\begin{overprint}
\begin{dessinp}{scale=0.52}
\foreach \av/\q/\b  [count=\xi]
in  
{#1}
{ 
\only<\xi| handout 0>{ 
 \StrLen{\av}[\Resintm] % ATTENTION IL FAUT CETTE SYNTAXE POUR FAIRE
                         % EQUIVALENT D UN EDEF
 \hspace{-1cm}\MACHINEDETURING{\av \b}{\Resintm}{$\q$}{5}{33}{}
}}
 \end{dessinp}
\end{overprint}
}

\def\MAXDET{20}

\newcommand\DIAGRAMMEESPACETEMPS[1]{
\draw node[coordinate] (tt) {};
\foreach \av/\q/\b  in  
{#1}
{ 

\RUBANMACHINEDETURING{\av{$\q$}\b}{0}{4}{\MAXDET}{}{tt};
\draw (tt) +(0,-\hcase) node[coordinate] (tt) {};
}
%\foreach \x in {0,...,19} {\draw (tt) + (\x*\lcase,0) node {$\vdots$};}
}



\newcommand\progun{
/q_0/0011, % initial
X/q_1/011,X0/q_1/11,X/q_2/0Y1,/q_2/X0Y1,
X/q_0/0Y1,XX/q_1/Y1,XXY/q_1/1, XX/q_2/YY, X/q_2/XYY,
XX/q_0/YY,XXY/q_3/Y,XXYY/q_3/B,XXYYB/q_4/B
}

\newcommand\DIAGRAMMEESPACETEMPSun{
\DIAGRAMMEESPACETEMPS
%% RECOPIE progun
{/q_0/0011, % initial
X/q_1/011,X0/q_1/11,X/q_2/0Y1,/q_2/X0Y1,
X/q_0/0Y1,XX/q_1/Y1,XXY/q_1/1, XX/q_2/YY, X/q_2/XYY,
XX/q_0/YY,XXY/q_3/Y,XXYY/q_3/B,XXYYB/q_4/B}
%% FIN RECOPIE
}

\newcommand\progdeux{
/q_0/0010, X/q_1/010, X0/q_1/10, X/q_2/0Y0,/q_2/X0Y0,
X/q_0/0Y0,XX/q_1/Y0,XXY/q_1/0,XXY0/q_1/%B,
}


\newcommand\DIAGRAMMEESPACETEMPSdeux{
\DIAGRAMMEESPACETEMPS
{
%% RECOPIE progdeux
/q_0/0010, X/q_1/010, X0/q_1/10, X/q_2/0Y0,/q_2/X0Y0,
X/q_0/0Y0,XX/q_1/Y0,XXY/q_1/0,XXY0/q_1/%B,
%% FIN RECOPIE
}
}

\newcommand\DIAGRAMMEESPACETEMPSdeuxvariante{
\def\memolcase{\lcase}
\def\memoMAXDET{20}
\def\MAXDET{15}
\def\lcase{6 ex}
\DIAGRAMMEESPACETEMPS
{
%% RECOPIE progdeux
/q_00/010, X/q_10/10, X0/q_11/0, X/q_20/Y0,/q_2X/0Y0,
X/q_00/Y0,XX/q_1Y/0,XXY/q_10/,XXY0/q_1B/%B,
%% FIN RECOPIE
}
\def\lcase{\memolcase}
\def\MAXDET{\memoMAXDET}
}


\newcommand\ANIMATIONun{
\ANIMATIONMT
%% RECOPIE progun
{/q_0/0011, % initial
X/q_1/011,X0/q_1/11,X/q_2/0Y1,/q_2/X0Y1,
X/q_0/0Y1,XX/q_1/Y1,XXY/q_1/1, XX/q_2/YY, X/q_2/XYY,
XX/q_0/YY,XXY/q_3/Y,XXYY/q_3/B,XXYYB/q_4/B}
%% FIN RECOPIE
}



\newcommand\DESSINALGO[3]{
% argument 1= ou
% argument 2= texte 
% argument 3= nom du noeud ex m
\draw node[draw,#1,minimum size=1cm] (#3) {#2};

% pour fitting box
\draw (#3.west) +(-0.2cm,0) node (#3c2) {};
\draw (#3.north) +(0,+0.2cm) node (#3c3) {};
\draw (#3.south) +(0,-0.2cm) node (#3c4) {};
}

\newcommand\DESSINALGOa[3]{
\DESSINALGO{#1}{#2}{#3}
\draw (#3.east) +(-0.1cm,0.2cm) node(#3a) {};
\draw[->] (#3a) -- ++(0.5cm,0) node (#3aa) {};
\draw (#3aa) node[anchor=west] (#3aaa) {\LANGUE{Accept}{Accepte}};
\draw (#3aaa.east) node(#3aaaa) {};
}

\newcommand\DESSINALGOe[4]{
\DESSINALGO{#1}{#2}{#3}
\draw (#3.east) +(-0.1cm,0.2cm) node(#3a) {};
\draw[->] (#3a) -- ++(0.5cm,0) node (#3aa) {};
\draw (#3aa) node[anchor=west] (#3aaa) {#4};
\draw (#3aaa.east) node(#3aaaa) {};
}


\newcommand\DESSINALGOar[3]{
\DESSINALGOa{#1}{#2}{#3}
\draw (#3.east) +(-0.1cm,-0.2cm) node(#3r) {};
\draw[->] (#3r) -- ++(0.5cm,0) node (#3rr) {};
\draw (#3rr) node[anchor=west] (#3rrr) {\LANGUE{Reject}{Refuse}};
\draw (#3rrr.east) node(#3rrrr) {};
}


\newcommand\WINDOWarg[7] {
%#1: position
%#2..7: contenu du tableau
\draw (#1) node[case]  {#2} \dcase;
\draw (#1) ++ (\lcase,0) node[case] {#3} \dcase;
\draw (#1) ++ (2*\lcase,0) node[case] {#4} \dcase;
\draw (#1) ++ (0,-\lcase) node[case] {#5} \dcase;
\draw (#1) ++ (\lcase,-\lcase) node[case] {#6} \dcase;
\draw (#1) ++ (2*\lcase,-\lcase) node[case] {#7} \dcase;
}

\newcommand\WINDOW[3]{
%#1: positioin
%#2: contenu
%#3: étiquette
\WINDOWarg{#1}
{\StrChar{#2}{1}}{\StrChar{#2}{2}}{\StrChar{#2}{3}}
{\StrChar{#2}{4}}{\StrChar{#2}{5}}{\StrChar{#2}{6}}
\draw (#1) +(-\lcase*1.2,-\lcase*0.5) node {#3};
}





%%%%%%%%%%%%%%%%%%%%%%%%%%%%%%%%%%%%%%%%%%%%%%%%%%%%%%%%%%%%%%%%%%%%%%%%%%%%
%                 LISTINGS
%%%%%%%%%%%%%%%%%%%%%%%%%%%%%%%%%%%%%%%%%%%%%%%%%%%%%%%%%%%%%%%%%%%%%%%%%%%%


% LISTINGS
\RequirePackage{listings}
 \lstset{language=Java,
 mathescape=true,
 showspaces=false,
 showstringspaces=false,
 frame=single,
morekeywords={recursive,then,div,function,integer,read,out,else,par,endpar,seq,endseq,and,not,read,write,repeat,until,undef,let,to,endfor,endif,endwhile,guess,in,parallel},
% basicstyle=\ttfamily\small,
% basewidth={1.1ex},
% keywordstyle=\bf\color{blue},
% %stringstyle=\bf\ttfamily\color{green},
% commentstyle=\bfseries\color{magenta},
 literate={:=}{{$\gets$ }}1 {<=}{{$\leq$}}2  {>=}{{$\geq$}}2
 {!=}{{$\neq$}}1 
 }
% \lstset{escapechar=\%}
\let\lst\lstinline
\def\beginlisting{\begin{lstlisting}}
\def\endlisting{\end{lstlisting}}  




\newenvironment{algo}[1]{
  \lstset{escapechar=\@}
  \def\input{$#1$}
  \def\out{out}
  \def\BEGIN{\textbf{read} \input }
 \def\BEGINN{\textbf{repeat} }
  \def\END{\textbf{until out!=undef}}
  \def\ENDD{\textbf{write \out}}
  \def\ENDT{\textbf{until $\ctl$ = $q_a$ or $\ctl$ = $q_r$}}
  \def\ENDR{\textbf{until $\ctl$ = $q_a$}}
  \def\ENDP{\textbf{until $\ctl$ = $q_a$}}
}{}
\newenvironment{algon}[1]{
  \lstset{numbers=left}
  \begin{algo}{#1}
}{\end{algo}}
\newcommand\FSM[3]{FSM(#1,\lst{#2},#3)}
\newcommand\FSMi[3]{FSM(#1,{#2},#3)}




%%%%%%%%%%%%%%%%%%%%%%%%%%%%%%%%%%%%%%%%%%%%%%%%%
%
%
%                      TRADUCTION EN COURS
%
%%%%%%%%%%%%%%%%%%%%%%%%%%%%%%%%%%%%%%%%%%%%%%%%%

\newcommand\ENTREPOT[1]{\NDLR{ENTREPOT ICI: #1}}
\newcommand\scite[2][]{\cite[#1]{#2}}
\newcommand\nospellcheck[1]{{#1}}
\newcommand\nd{nd}
\newcommand\NL{
 
}

\newenvironment{FRANCAIS}{
\small
\color{orange}
}
{\normalsize
\color{black}
}


\newenvironment{SEMIFRANCAIS}{
\small
\color{violet}
}
{\normalsize
\color{black}
}

\newcommand\QFRANCAIS[1]{
\begin{FRANCAIS}
#1
\end{FRANCAIS}
}


%%%%%%%%%%%%%%%%%%%%%%%%%%%%%%%%%%%%%%%%%%%%%%%%%
%
%
%                      Wikipedia
%
%%%%%%%%%%%%%%%%%%%%%%%%%%%%%%%%%%%%%%%%%%%%%%%%%

\providecommand\tightlist{}
\providecommand\Complex{\C}
\newenvironment{myalgie}{\begin{array}{lll}}{\end{array}}

\newcommand\nl{\ \\ }




%%%%%%%%%%%%%%%%%%%%%%%%%%%%%%%%%%%%%%%%%%%%%%%%%
%
%
%                      Pour citations précises
%
%%%%%%%%%%%%%%%%%%%%%%%%%%%%%%%%%%%%%%%%%%%%%%%%%



\newcommand\citeprecis[2]{{\cite[#1]{#2}}}
\newcommand\citeprecisinv[2]{\citeprecis{#2}{#1}}

\usepackage{csquotes}



%%%%%%%%%%%%%%%%%%%%%%%%%%%%%%%%%%%%%%%%%%%%%%%%%
%
%
%                      Travail sur couleurs
%
%   toutes les réponses pointent dans le même sens, et le schéma correspond exactement 
%. au profil perceptif d’un protanope** (déficit des cônes L / longueurs d’onde longues)
%  MAIS SEVERE
%
%%%%%%%%%%%%%%%%%%%%%%%%%%%%%%%%%%%%%%%%%%%%%%%%%


% ================================
% GROUPE 1 — Verts sombres
% ================================
\definecolor{deepgreen}{RGB}{0,160,40}
\definecolor{forestgreen}{RGB}{0,150,70}
\definecolor{mossgreen}{RGB}{0,140,90}
\definecolor{greenwave}{RGB}{0,150,90}
\definecolor{forestmist}{RGB}{0,160,90}
\definecolor{darkteal}{RGB}{0,130,70}

% ================================
% GROUPE 2 — Teal / verts bleutés
% ================================
\definecolor{seagreen}{RGB}{0,130,100}
\definecolor{freshteal}{RGB}{0,140,130}
\definecolor{greenteal}{RGB}{0,140,120}
\definecolor{blueteal}{RGB}{0,150,120}      % ancien "teal" (doublon corrigé)
\definecolor{lagoon}{RGB}{0,180,120}
\definecolor{greenmist}{RGB}{0,170,110}

% ================================
% GROUPE 3 — Turquoises / Aqua
% ================================
\definecolor{springgreen}{RGB}{0,190,100}
\definecolor{mintgreen}{RGB}{10,200,150}
\definecolor{oceanfoam}{RGB}{0,200,150}
\definecolor{aquagreen}{RGB}{30,210,170}
\definecolor{coastalaqua}{RGB}{40,210,180}
\definecolor{softaqua}{RGB}{40,210,190}

% ================================
% GROUPE 4 — Cyan / Aqua clairs
% ================================
\definecolor{emeraldteal}{RGB}{0,180,140}
\definecolor{deepcyan}{RGB}{0,180,160}
\definecolor{lightcyan}{RGB}{0,200,160}
\definecolor{aquamarine}{RGB}{0,200,160}
\definecolor{lightaqua}{RGB}{60,230,200}
\definecolor{mistblue}{RGB}{80,240,210}

% ================================
% GROUPE 5 — Bleus clairs
% ================================
\definecolor{paleaqua}{RGB}{80,200,190}
\definecolor{glacier}{RGB}{90,150,200}
\definecolor{skyblueA}{RGB}{60,150,220}      % ancien "skyblue"
\definecolor{lightsky}{RGB}{90,170,230}
\definecolor{iceblue}{RGB}{120,190,240}
\definecolor{frost}{RGB}{150,210,250}

% ================================
% GROUPE 6 — Bleus moyens
% ================================
\definecolor{steelblue}{RGB}{30,100,160}
\definecolor{cadetblue}{RGB}{70,120,180}
\definecolor{coolblue}{RGB}{110,170,220}
\definecolor{morningblue}{RGB}{140,200,240}
\definecolor{blueiceA}{RGB}{40,150,210}       % ancien "blueice" (doublon corrigé)
\definecolor{blueair}{RGB}{60,170,225}

% ================================
% GROUPE 7 — Azurs / Bleu saturé
% ================================
\definecolor{tealblue}{RGB}{0,110,130}
\definecolor{deepbluecyan}{RGB}{0,110,170}
\definecolor{azur}{RGB}{0,70,180}
\definecolor{brightazure}{RGB}{20,130,200}
\definecolor{lightazure}{RGB}{80,180,235}
\definecolor{skylight}{RGB}{100,200,245}

% ================================
% GROUPE 8 — Bleus sombres / violets
% ================================
\definecolor{nightblue}{RGB}{40,0,110}
\definecolor{slateblue}{RGB}{20,40,120}
\definecolor{violetA}{RGB}{60,50,160}         % ancien "violet" (distinct de violetB ci-dessous)
\definecolor{blueviolet}{RGB}{80,60,180}
\definecolor{purpleA}{RGB}{100,70,200}        % ancien "purple"
\definecolor{lightpurple}{RGB}{130,90,210}

% ================================
% GROUPE 9 — Violets saturés / indigos
% ================================
\definecolor{lilac}{RGB}{160,110,220}
\definecolor{indigoA}{RGB}{50,0,130}          % ancien "indigo"
\definecolor{deepindigo}{RGB}{70,10,150}
\definecolor{royalpurple}{RGB}{90,20,170}

% ================================
% GROUPE 10 — Fermeture cercle perceptif
% ================================
\definecolor{coolmint}{RGB}{0,130,70}
\definecolor{deepsea}{RGB}{0,130,70}

% ================================
% PALETTE DEUTÉRANOPE — noms uniques
% ================================
\definecolor{crimson}{RGB}{200,30,60}
\definecolor{brickred}{RGB}{180,40,40}
\definecolor{coral}{RGB}{220,90,70}
\definecolor{salmon}{RGB}{240,130,120}

\definecolor{magenta}{RGB}{200,40,150}
\definecolor{fuchsia}{RGB}{220,60,170}
\definecolor{violetB}{RGB}{140,70,190}        % ancien "violet"
\definecolor{indigoB}{RGB}{90,50,180}         % ancien "indigo"

\definecolor{royalblueA}{RGB}{50,80,200}      % ancien "royalblue"
\definecolor{azure}{RGB}{30,120,230}
\definecolor{skyblueB}{RGB}{80,160,230}       % second "skyblue"
\definecolor{tealB}{RGB}{0,150,170}           % second "teal"

\definecolor{marine}{RGB}{0,80,140}




%% Alternative

% === Générateur automatique fill/text pour protanopes ===
% Usage : \ProtaColorPair{mistblue}
% Suppose que mistblue est déjà défini par \definecolor{mistblue}{RGB}{...}

\newcommand{\ProtaColorPair}[1]{%
  % fill = couleur originale
  \colorlet{#1fill}{#1}%
  % text = version assombrie idéalement lisible
  \colorlet{#1text}{#1!80!black}%
}

%Utilisation
%\ProtaColorPair{mistblue}

%\textcolor{mistbluefill}{Couleur claire} \quad
%\textcolor{mistbluetext}{Couleur pour texte}

	\foreach \i/\name in {
	  1/deepgreen,
	  2/greenwave,
	  3/lagoon,
	  4/oceanfoam,
	  5/softaqua,
	  6/mistblue,
 	 7/lightsky,
	  8/blueiceA,   % <-- corrigé
	  9/deepcyan,
	  10/azur,
	  11/blueviolet,
	  12/royalpurple}
	  {\ProtaColorPair{\name}}
	  


\newcommand\ROUEPROTANOPEDOUZE{
	% Roue chromatique protanope — 12 couleurs
	\begin{tikzpicture}[scale=3]

	\foreach \i/\name/\nametext in {
	  1/deepgreen/deepgreen,
	  2/greenwave/greenwave,
	  3/lagoon/lagoontext,
	  4/oceanfoam/oceanfoam,
	  5/softaqua/softaquatext,
	  6/mistblue/mistbluetext,
 	 7/lightsky/lightskytext,
	  8/blueiceA/blueiceA,   % <-- corrigé
	  9/deepcyan/blueiceA,
	  10/azur/azur,
	  11/blueviolet/blueviolet,
	  12/royalpurple/royalpurple}
	{
	  % secteur coloré
 	 \filldraw[fill=\name, draw=black, line width=0.2pt]
	    ({90-30*\i}:1)
	    -- ({90-30*(\i-1)}:1)
	    -- (0,0) -- cycle;

 	 % label dans la bonne couleur
	  \node at ({90-30*(\i-0.5)}:1.25)
	    {\textcolor{\nametext}{\small \nametext}};
	}

	\end{tikzpicture}
}


\newcommand\ROUEPROTANOPEVINGT{
\begin{tikzpicture}[scale=3]

\foreach \i/\name/\nametext in {
  1/deepgreen/deepgreen,
  2/greenwave/greenwave,
  3/lagoon/lagoon,
  4/oceanfoam/oceanfoam,
  5/softaqua/softaquatext,
  6/mistblue/mistbluetext,
  7/lightsky/lightsky,
  8/blueiceA/blueiceA,      % corrigé
  9/deepcyan/deepcyan,
  10/azur/azur,
  11/brightazure/brightazure,
  12/skyblueA/skyblueA,     % corrigé
  13/coolblue/coolblue,
  14/glacier/glacier,
  15/marine/marine,
  16/violetA/violetA,      % corrigé
  17/blueviolet/blueviolet,
  18/purpleA/purpleA,      % corrigé
  19/royalpurple/royalpurple,
  20/indigoA/indigoA}      % corrigé
{
  % Secteur coloré
  \filldraw[fill=\name, draw=black, line width=0.2pt]
    ({90-18*\i}:1)
    -- ({90-18*(\i-1)}:1)
    -- (0,0) -- cycle;

  % Étiquette
  \node at ({90-18*(\i-0.5)}:1.25)
    {\textcolor{\nametext}{\small \nametext}};
}

\end{tikzpicture}
}

\newcommand\ROUEDEUTERANOPEDOUZE{
% Roue chromatique deutéranope — 12 couleurs
\begin{tikzpicture}[scale=3]

\foreach \i/\name in {
  1/crimson,
  2/brickred,
  3/coral,
  4/salmon,
  5/magenta,
  6/fuchsia,
  7/violetB,       % corrigé
  8/indigoB,       % corrigé
  9/royalblueA,    % corrigé
  10/azure,
  11/skyblueB,     % corrigé
  12/tealB}        % corrigé
{
  % Secteur coloré
  \filldraw[fill=\name, draw=black, line width=0.2pt]
    ({90-30*\i}:1)
    -- ({90-30*(\i-1)}:1)
    -- (0,0) -- cycle;

  % Label dans la couleur
  \node at ({90-30*(\i-0.5)}:1.25)
    {\textcolor{\name}{\small \name}};
}

\end{tikzpicture}
}



% === Macro intelligente pour foncer une couleur de façon lisible pour protanope ===
% Usage : \protadarken{mistblue}{40}  -> crée mistblue@dark
% Le 40 signifie : 40% de mélange avec du noir

\newcommand{\DarkenForText}[2]{%
  % #1 = nom de la couleur
  % #2 = pourcentage (40 conseillé si tu ne sais pas)
  %
  % On crée une nouvelle couleur #1@dark
  \colorlet{#1text}{black!#2!#1}%
}





%%%%

% Groupe 1
\DarkenForText{deepgreen}{20}
\DarkenForText{forestgreen}{20}
\DarkenForText{mossgreen}{20}
\DarkenForText{greenwave}{20}
\DarkenForText{forestmist}{20}
\DarkenForText{darkteal}{20}

% Groupe 2
\DarkenForText{seagreen}{20}
\DarkenForText{freshteal}{20}
\DarkenForText{greenteal}{20}
\DarkenForText{blueteal}{20}
\DarkenForText{lagoon}{20}
\DarkenForText{greenmist}{20}

% Groupe 3
\DarkenForText{springgreen}{20}
\DarkenForText{mintgreen}{20}
\DarkenForText{oceanfoam}{20}
\DarkenForText{aquagreen}{20}
\DarkenForText{coastalaqua}{20}
\DarkenForText{softaqua}{20}

% Groupe 4
\DarkenForText{emeraldteal}{20}
\DarkenForText{deepcyan}{20}
\DarkenForText{lightcyan}{20}
\DarkenForText{aquamarine}{20}
\DarkenForText{lightaqua}{20}
\DarkenForText{mistblue}{30}

% Groupe 5
\DarkenForText{paleaqua}{20}
\DarkenForText{glacier}{20}
\DarkenForText{skyblueA}{20}
\DarkenForText{lightsky}{20}
\DarkenForText{iceblue}{20}
\DarkenForText{frost}{20}

% Groupe 6
\DarkenForText{steelblue}{20}
\DarkenForText{cadetblue}{20}
\DarkenForText{coolblue}{20}
\DarkenForText{morningblue}{20}
\DarkenForText{blueiceA}{20}
\DarkenForText{blueair}{20}

% Groupe 7
\DarkenForText{tealblue}{20}
\DarkenForText{deepbluecyan}{20}
\DarkenForText{azur}{20}
\DarkenForText{brightazure}{20}
\DarkenForText{lightazure}{20}
\DarkenForText{skylight}{20}

% Groupe 8
\DarkenForText{nightblue}{20}
\DarkenForText{slateblue}{20}
\DarkenForText{violetA}{20}
\DarkenForText{blueviolet}{20}
\DarkenForText{purpleA}{20}
\DarkenForText{lightpurple}{20}

% Groupe 9
\DarkenForText{lilac}{20}
\DarkenForText{indigoA}{20}
\DarkenForText{deepindigo}{20}
\DarkenForText{royalpurple}{20}

% Groupe 10
\DarkenForText{coolmint}{20}
\DarkenForText{deepsea}{20}

% Palette deutéranope
\DarkenForText{crimson}{20}
\DarkenForText{brickred}{20}
\DarkenForText{coral}{20}
\DarkenForText{salmon}{20}

\DarkenForText{magenta}{20}
\DarkenForText{fuchsia}{20}
\DarkenForText{violetB}{20}
\DarkenForText{indigoB}{20}

\DarkenForText{royalblueA}{20}
\DarkenForText{azure}{20}
\DarkenForText{skyblueB}{20}
\DarkenForText{tealB}{20}

% Couleur supplémentaire
\DarkenForText{marine}{20}

%%%%%%%%%%%%%%%%%%%%%
%
%.   apx proof
%
%%%%%%%%%%%%%%%%%%%%%



% do
\newcommand\PREUVESENAPPENDIXCOMPATIBLE{
	\newenvironment{appendixproof}{{\bf \LANGUE{Proof}{Démonstration}}:}{\hfill $\Box$}
	\newtheorem{theoremrep}[theorem]{Theorem}
	\newtheorem{lemmarep}[theorem]{Lemma}
	\newtheorem{propositionrep}[theorem]{Proposition}
	\newtheorem{corollaryrep}[theorem]{Corollary}
	\newtheorem{summaryrep}[theorem]{Summary}	
}

\providecommand\apxparameter{}

%% proof inline
%\providecommand\apxparameter{appendix=inline}
%% proof in appendix
%\providecommand\apxparameter{}


\apxproofmode{
\usepackage[\apxparameter]{apxproof}
\theoremstyle{plain}

\newtheoremrep{theorem}{Theorem}
%\newtheoremrep{theoremconjecture}[theorem]{Theorem-To-Be-Proved=Conjecture}
%\newtheoremrep{lemmaconjecture}[theorem]{Conjecture-To-Be-Proved=Conjecture}
\newtheoremrep{claim}[theorem]{Theorem}
\newtheoremrep{proposition}[theorem]{Proposition}
\newtheoremrep{lemma}[theorem]{Lemma}
\newtheoremrep{corollary}[theorem]{Corollary}
%%bug mars 26: \newtheoremrep{summary}[theorem]{Summary}
% Du coup:
\newtheoremrep{resume}[theorem]{Summary}
%bug mars 26: \newtheoremrep{openproblem}[theorem]{Open Problem}
%\newtheoremrep{claim}[theorem]{Claim}
}


%\apxproofmodesubsection{
%%apx proof, pour que compteur correct
%\makeatletter
%\AtBeginDocument{%
%\let\axp@section\subsection
%  \let\axp@oldsection\subsection
%}
%\makeatother
%\renewcommand\appendixprelim{\section{Missing proofs}}
%}


\makeatletter
\apxproofmodesubsection{
\let\axp@section\subsection
%\let\axp@oldsection\subsection
\renewcommand\appendixprelim{\section{Proofs from main documents}}
}
\makeatother


\makeatletter
\newcommand\appendixapxproofsubsection{
\appendix
\let\apx@orig@appendix\appendix
\renewcommand{\appendix}{}
}
\makeatother


%%% Pour checker/vérifier que bon


\newcommand\docheck[1]{\textcolor{check}{#1}}

\usepackage{xcolor}
%\definecolor{chose}{RGB}{0,110,50}
%\definecolor{forestgreen}{RGB}{0,150,70}
%\definecolor{ctruc}{RGB}{164,164,164}
%\definecolor{oceanfoam}{RGB}{0,200,150}
\definecolor{check}{RGB}{160,160,0}




. Retirées automatiquement en mode lipics.
\providecommand\NODECORATION[1]{\ifnotSLIDE{#1}}
\lipicsmode{\renewcommand\NODECORATION[1]{}}

\providecommand\PASSIAPXPROOF[1]{#1}
\apxproofmode{\renewcommand\PASSIAPXPROOF[1]{}}



%% end gestion des switchs

\def\MACROSPOINTEXOUVERT{}


\providecommand\ifnotTDEXAM[1]{#1}
\providecommand\ifnotSLIDE[1]{#1}


  
% pour preuves en appendix
\providecommand\PREUVESENAPPENDIX[1]{}



%%%% STYLES: APRES VOLS


\NODECORATION{
\usepackage{fourier}
}

%% Repli si fourier (et donc fourier-orns, qui définit \danger) n'est pas
%% chargé ; doit rester APRÈS le bloc ci-dessus, sinon fourier-orns échoue
%% avec « \danger already defined ».
\providecommand\danger{}

%%%% Box around environment
\usepackage[framemethod=tikz]{mdframed}

\newcommand\ENVCOLOR[1]{#1}
\newcommand\BEGINTEMPENVCOLOR[1]{\let\envcolorwas\ENVCOLOR\renewcommand\ENVCOLOR[1]{#1}}
\newcommand\ENDTEMPENVCOLOR{\let\ENVCOLOR\envcolorwas}


\NODECORATION{
\surroundwithmdframed[hidealllines=true,backgroundcolor=\ENVCOLOR{blue!10},roundcorner=8pt]{exercice}
\surroundwithmdframed[hidealllines=true,backgroundcolor=\ENVCOLOR{blue!10},roundcorner=8pt]{exercise}
\surroundwithmdframed[hidealllines=true,backgroundcolor=\ENVCOLOR{blue!10},roundcorner=8pt]{exercicee}
\surroundwithmdframed[hidealllines=true,backgroundcolor=\ENVCOLOR{blue!10},roundcorner=8pt]{exerciced}
\surroundwithmdframed[hidealllines=true,backgroundcolor=\ENVCOLOR{blue!10},roundcorner=8pt]{exercicec}
\surroundwithmdframed[hidealllines=true,backgroundcolor=\ENVCOLOR{blue!10},roundcorner=8pt]{exercicedc}
\surroundwithmdframed[hidealllines=true,backgroundcolor=\ENVCOLOR{orange!20},roundcorner=8pt]{example}
\surroundwithmdframed[hidealllines=true,backgroundcolor=\ENVCOLOR{orange!20},roundcorner=8pt]{exemple}
\surroundwithmdframed[hidealllines=true,backgroundcolor=\ENVCOLOR{orange!10},roundcorner=8pt]{remark}
\ifnotSLIDE{\surroundwithmdframed[backgroundcolor=\ENVCOLOR{green!15},roundcorner=8pt]{definition}}
\surroundwithmdframed[hidealllines=true,backgroundcolor=\ENVCOLOR{gray!35}]{proposition}
\ifnotSLIDE{\surroundwithmdframed[backgroundcolor=\ENVCOLOR{gray!35},roundcorner=8pt]{theorem}}
\surroundwithmdframed[hidealllines=true,backgroundcolor=\ENVCOLOR{gray!35}]{lemma}
\surroundwithmdframed[hidealllines=true,backgroundcolor=\ENVCOLOR{gray!35}]{corollary}


\surroundwithmdframed[backgroundcolor=red!35,roundcorner=8pt]{theoremaprouver}
\surroundwithmdframed[backgroundcolor=red!35,roundcorner=8pt]{conjectureaprouver}
}

%%% tcolorbox

\usepackage{tcolorbox}
\tcbuselibrary{skins}


  %%. Commande magique
%% \tcolorboxenvironment{⟨name⟩}{⟨options⟩}: attach options de tcolorbox à environnement
%% alternative: je cree des boites moi meme: voici des exemples


\newtcolorbox{maboiteconfiance}[2][]{colback=red!5!white, colframe=red!75!black,fonttitle=\bfseries, colbacktitle=red!85!black,enhanced,
attach boxed title to top center={yshift=-2mm},
  title={#2},#1}
  
\newtcolorbox{maboiteresume}[2][]{colback=red!5!white, colframe=red!75!black,fonttitle=\bfseries, colbacktitle=red!85!black,enhanced,
attach boxed title to top center={yshift=-2mm},
  title={#2},#1}
  \newtcolorbox{maboiteresumeperso}[2][]{colback=red!5!white, colframe=red!75!black,fonttitle=\bfseries, colbacktitle=red!85!black,enhanced,
attach boxed title to top center={yshift=-2mm},
  title={#2},#1}
  \newtcolorbox{maboitecommentaire}[2][]{colback=red!5!white, colframe=red!75!black,fonttitle=\bfseries, colbacktitle=red!85!black,enhanced,
attach boxed title to top center={yshift=-2mm},
  title={#2},#1}
\newtcolorbox{maboitetruc}[2][]{colback=red!5!white, colframe=red!75!black,fonttitle=\bfseries, colbacktitle=red!85!black,enhanced,
attach boxed title to top left={yshift=-2mm},
  title={#2},#1}

%% Boîte des mises en évidence CARE. Couleurs pensées pour la
%% protanopie : l'axe bleu-jaune reste discriminant (jaune = \IMPORTANT
%% d'Olivier, bleu = CARE), et la distinction ne repose pas que sur la
%% couleur (cadre + bandeau étiqueté, lisibles même en niveaux de gris).
\newtcolorbox{maboiteimportantcare}[2][]{colback=cyan!10!white, colframe=blue!60!black,fonttitle=\bfseries, colbacktitle=blue!70!black,enhanced,
attach boxed title to top center={yshift=-2mm},
  title={#2},#1}
  


%%%% FIN STYLES



%%%%%%%%%%%%%%%%%%%%%%%%%%%%%%%%%%%%%%%%%%%%%%%%%
%%%%%%%%%%%%%%%%%%%%%%%%%%%%%%%%%%%%%%%%%%%%%%%%%
%
%
%                      Version ENSEIGNEMENT 2020
%
%
%%%%%%%%%%%%%%%%%%%%%%%%%%%%%%%%%%%%%%%%%%%%%%%%%
%%%%%%%%%%%%%%%%%%%%%%%%%%%%%%%%%%%%%%%%%%%%%%%%%

%english, french
\ifdefined\CESTDUFRANCAIS
	\providecommand\LANGUE[2]{#2}
	\usepackage[french]{babel}
\else
	\providecommand\LANGUE[2]{#1}
\fi
\providecommand\LANGUEinv[2]{\LANGUE{#2}{#1}}

%%%%%%%%%%%%%%%%%%%%%%%%%%%%%%%%%%%%%%%%%%%%%%%%%
%
%
%                      Packages
%
%
%%%%%%%%%%%%%%%%%%%%%%%%%%%%%%%%%%%%%%%%%%%%%%%%%


%%% SURLIGNEUR
\providecommand\ifnotmodeDIFFUSIONFINALE[1]{
%% 2021: Comprend pas mais assure
\ifdefined\SAFEMODE
\else
#1
\fi
%#1
}
\providecommand\ifmodeDIFFUSIONFINALE[1]{
\ifdefined\SAFEMODE
#1
\else
\fi
}
\newcommand\SURLIGNE[1]{
\ifnotmodeDIFFUSIONFINALE{
\BEGINTEMPENVCOLOR{blue!30}
{                        \colorbox{blue!30}{

\noindent \begin{minipage}{\dimexpr\textwidth-2\fboxsep\relax}
#1
\end{minipage}
}
}
\ENDTEMPENVCOLOR}
\ifmodeDIFFUSIONFINALE{
\noindent \begin{minipage}{\textwidth}
#1
\end{minipage}
}
}


\newenvironment{confiance}{\begin{maboiteresume}[colback=brown!50]{Résumé ici}}{\end{maboiteresume}}
\newcommand\CONFIANCE[2][Confiance ici]{
\begin{maboiteconfiance}[colback=yellow!50]{#1}
\textbf{ATTENTION : #2}
\end{maboiteconfiance}
}

\newenvironment{resume}{\begin{maboiteresume}[colback=brown!50]{Résumé ici}}{\end{maboiteresume}}
\newcommand\RESUME[2][Résumé ici]{
\begin{maboiteresume}[colback=brown!50]{#1}
#2
\end{maboiteresume}
}

\newenvironment{resumeperso}{\begin{maboiteresumeperso}[colback=brown!50]{Résumé perso ici}}{\end{maboiteresumeperso}}
\newcommand\RESUMEPERSO[2][Résumé perso  ici]{
\begin{maboiteresumeperso}[colback=brown!50]{#1}
#2
\end{maboiteresumeperso}
%
%
%\todo[inline, color=brown]{résumé: #1}
%\colorbox{brown}{
%\noindent \begin{minipage}{\textwidth}
%{
%\bf 
%#2
%}
%\end{minipage}
%}
}

\newcommand\DANSLAMARGE[1]{
\ifnotmodeDIFFUSIONFINALE{
\marginpar{\textcolor{brown}{#1}}}
}

\newenvironment{commentaireperso}{\begin{maboitecommentaire}[colback=brown!30]{Commentaire perso} Commentaire perso:}{\end{maboitecommentaire}}
\newcommand\COMMENTAIREPERSO[2][Commentaire perso]{
\begin{maboitecommentaire}[colback=brown!30]{#1}
#2
\end{maboitecommentaire}
}


\newcommand\SURSURLIGNE[2][truc surligné ici]{
%\todo[inline, color=orange]{surligné: #1}
\ifnotmodeDIFFUSIONFINALE{
\BEGINTEMPENVCOLOR{blue!60}
	                 \colorbox{blue!60}{
\noindent \begin{minipage}{\dimexpr\textwidth-2\fboxsep\relax}
{
\bf 
#2
}
\end{minipage}
}
\ENDTEMPENVCOLOR
}
\ifmodeDIFFUSIONFINALE{#2}
}

\newcommand\IMPORTANT[1]{
\ifnotmodeDIFFUSIONFINALE{
\BEGINTEMPENVCOLOR{yellow!100}
		        \colorbox{yellow!100}{

\noindent \begin{minipage}{\dimexpr\textwidth-2\fboxsep\relax}
{ 
#1
}
\end{minipage}
}
\ENDTEMPENVCOLOR
}
\ifmodeDIFFUSIONFINALE{#1}
}

%% Pendant de \IMPORTANT, réservé aux mises en évidence CARE ;
%% \IMPORTANT reste réservé aux annotations d'Olivier. Comme \IMPORTANT,
%% redevient du texte simple en mode DIFFUSIONFINALE.
\newcommand\IMPORTANTCARE[2][Important (CARE)]{
\ifnotmodeDIFFUSIONFINALE{
\begin{maboiteimportantcare}{#1}
#2
\end{maboiteimportantcare}
}
\ifmodeDIFFUSIONFINALE{#2}
}


\newcommand\QUESTION[1]{
\ifnotmodeDIFFUSIONFINALE{
\BEGINTEMPENVCOLOR{oceanfoam!100}
		        \colorbox{oceanfoam!100}{

\noindent \begin{minipage}{\dimexpr\textwidth-2\fboxsep\relax}
{ 
#1
}
\end{minipage}
}
\ENDTEMPENVCOLOR
}
\ifmodeDIFFUSIONFINALE{#1}
}




\newcommand\SOUSLIGNE[1]{
\colorbox{lightgray}{

\noindent \begin{minipage}{\dimexpr\textwidth-2\fboxsep\relax}
{ \color{gray}
#1
}
\end{minipage}
}
}



%% Gestion des cross-references cross-referencing
\usepackage{xr-hyper}
\usepackage{hyperref}

%%
\usepackage{versions}
\usepackage{amsmath,amssymb,mathrsfs,amsfonts}
%\usepackage{mathabx}
\usepackage{listings}
\usepackage{url}
\usepackage{versions}
\usepackage{color}
\usepackage[colorinlistoftodos]{todonotes}
%\usepackage{todonotes}
\usepackage{proof}
\usepackage{xstring}

%% Index
\usepackage{makeidx}
\makeindex
\providecommand\DOINDEXPRINT{}
\providecommand\SIGNALEMOTNOUV[1]{}
\ifdefined\INDEXPRINT
	\ifnotSAFEMODE{
%	\usepackage{showidx}
%	\let\oldindex\index
%	\renewcommand{\index}[1]{\oldindex{#1}\marginpar{LA: \tiny#1}}
	\renewcommand\DOINDEXPRINT{
	      	\let\oldindex\index
		\renewcommand{\index}[1]{\oldindex{##1}\marginpar{IDXLA: \tiny##1}}
		\renewcommand\SIGNALEMOTNOUV[1]{\marginpar{MNV: \small\textbf{##1}}}
	}
	}
\fi

% === gestion des accents


%soit iso latin1
%\usepackage[cyr]{aeguill}
%soit utf8
\usepackage[utf8]{inputenc}
\usepackage[T1]{fontenc}


%%%  Pour citation \Cref (cref, ref)	

\usepackage{cleveref}


%%%%%%%%%%%%%%%%%%%%%%%%%%%%%%%%%%%%%%%%%%%%%%%%%
%
%
%                      Gestion des accents spéciaux.
%
%
%%%%%%%%%%%%%%%%%%%%%%%%%%%%%%%%%%%%%%%%%%%%%%%%%


%% MOI A LA MAIN:
\usepackage{bbm}

\DeclareUnicodeCharacter{2212}{-}
\DeclareUnicodeCharacter{2297}{\ensuremath{\otimes}}       
\DeclareUnicodeCharacter{2A33}{\ensuremath{\smashtimes}}
\DeclareUnicodeCharacter{1D56F}{\ensuremath{\mathfrak{D}}}  
\DeclareUnicodeCharacter{1D40E}{\ensuremath{\mathbf{O}}}  
\DeclareUnicodeCharacter{1D427}{\ensuremath{\mathbf{n}}}  
\DeclareUnicodeCharacter{3008}{\ensuremath{\langle}} % 〈 
\DeclareUnicodeCharacter{3009}{\ensuremath{\rangle}} %  〉

% (possible de voir: http://tug.ctan.org/info/symbols/comprehensive/symbols-letter.pdf)

%% SOURCE:  https://github.com/agda/agda/blob/master/doc/user-manual/unicode-symbols-sphinx.tex.txt

%%%%%%%%%%%%%%%%%%%%%%%%%%%%%%%%%%%%%%%%%%%%%%%%%%%%%%%%%%%%%%%%%%%%%%%%%%%%%%
% Please keep the below list ordered by code points!

\DeclareUnicodeCharacter{00AC}{\ensuremath{\neg}}                     % Symbol ¬
\DeclareUnicodeCharacter{00B1}{\ensuremath{\pm}}                      % Symbol ±
\DeclareUnicodeCharacter{00B2}{\ensuremath{^2}}                       % Symbol ²
\DeclareUnicodeCharacter{00B3}{\ensuremath{^3}}                       % Symbol ³
\DeclareUnicodeCharacter{00B7}{\ensuremath{\cdot}}                    % Symbol ·
\DeclareUnicodeCharacter{00B9}{\ensuremath{^1}}                       % Symbol ¹
\DeclareUnicodeCharacter{00D7}{\ensuremath{\times}}                   % Symbol ×
\DeclareUnicodeCharacter{02B0}{\ensuremath{^h}}                       % Symbol ʰ
\DeclareUnicodeCharacter{02B3}{\ensuremath{^r}}                       % Symbol ʳ
\DeclareUnicodeCharacter{02E2}{\ensuremath{^s}}                       % Symbol ˢ
\DeclareUnicodeCharacter{0302}{\ensuremath{\hat{\phantom{x}}}}        % Symbol ̂
\DeclareUnicodeCharacter{0332}{\ensuremath{\underline{\phantom{x}}}}  % Symbol ̲
\DeclareUnicodeCharacter{0308}{\ensuremath{\ddot{\phantom{x}}}}       % Symbol ̈
\DeclareUnicodeCharacter{0393}{\ensuremath{\Gamma}}                   % Symbol Γ
\DeclareUnicodeCharacter{0394}{\ensuremath{\Delta}}                   % Symbol Δ
\DeclareUnicodeCharacter{0398}{\ensuremath{\Theta}}                   % Symbol Θ
\DeclareUnicodeCharacter{039B}{\ensuremath{\Lambda}}                  % Symbol Λ
\DeclareUnicodeCharacter{039E}{\ensuremath{\Xi}}                      % Symbol Ξ
\DeclareUnicodeCharacter{03A0}{\ensuremath{\Pi}}                      % Symbol Π
\DeclareUnicodeCharacter{03A3}{\ensuremath{\Sigma}}                   % Symbol Σ


% There is not `\Tau` command in LaTeX, so we use `\mathrm{T}`. See
% http://tex.stackexchange.com/questions/1368/why-can-i-only-use-some-capital-greek-letters-inside-my-equations.
\DeclareUnicodeCharacter{03A4}{\ensuremath{\mathrm{T}}}  % Symbol Τ

\DeclareUnicodeCharacter{03A5}{\ensuremath{\Upsilon}}   % Symbol Υ
\DeclareUnicodeCharacter{03A6}{\ensuremath{\Phi}}       % Symbol Φ
\DeclareUnicodeCharacter{03A8}{\ensuremath{\Psi}}       % Symbol Ψ
\DeclareUnicodeCharacter{03A9}{\ensuremath{\Omega}}     % Symbol Ω
\DeclareUnicodeCharacter{03B1}{\ensuremath{\alpha}}     % Symbol α
\DeclareUnicodeCharacter{03B2}{\ensuremath{\beta}}      % Symbol β
\DeclareUnicodeCharacter{03B3}{\ensuremath{\gamma}}     % Symbol γ
\DeclareUnicodeCharacter{03B4}{\ensuremath{\delta}}     % Symbol δ
\DeclareUnicodeCharacter{03B5}{\ensuremath{\epsilon}}   % Symbol ε
\DeclareUnicodeCharacter{03B6}{\ensuremath{\zeta}}      % Symbol ζ
\DeclareUnicodeCharacter{03B7}{\ensuremath{\eta}}       % Symbol η
\DeclareUnicodeCharacter{03B8}{\ensuremath{\theta}}     % Symbol θ
\DeclareUnicodeCharacter{03B9}{\ensuremath{\iota}}      % Symbol ι
\DeclareUnicodeCharacter{03BA}{\ensuremath{\kappa}}     % Symbol κ
\DeclareUnicodeCharacter{03BB}{\ensuremath{\lambda}}    % Symbol λ
\DeclareUnicodeCharacter{03BC}{\ensuremath{\mu}}        % Symbol μ
\DeclareUnicodeCharacter{03BD}{\ensuremath{\nu}}        % Symbol ν
\DeclareUnicodeCharacter{03BE}{\ensuremath{\xi}}        % Symbol ξ
\DeclareUnicodeCharacter{03C0}{\ensuremath{\pi}}        % Symbol π
\DeclareUnicodeCharacter{03C1}{\ensuremath{\rho}}       % Symbol ρ
\DeclareUnicodeCharacter{03C2}{\ensuremath{\varsigma}}  % Symbol ς
\DeclareUnicodeCharacter{03C3}{\ensuremath{\sigma}}     % Symbol σ
\DeclareUnicodeCharacter{03C4}{\ensuremath{\tau}}       % Symbol τ
\DeclareUnicodeCharacter{03C5}{\ensuremath{\upsilon}}   % Symbol υ
\DeclareUnicodeCharacter{03C6}{\ensuremath{\varphi}}    % Symbol φ
\DeclareUnicodeCharacter{03C7}{\ensuremath{\chi}}       % Symbol χ
\DeclareUnicodeCharacter{03C8}{\ensuremath{\psi}}       % Symbol ψ
\DeclareUnicodeCharacter{03C9}{\ensuremath{\omega}}     % Symbol ω
\DeclareUnicodeCharacter{03D1}{\ensuremath{\vartheta}}  % Symbol ϑ

\DeclareUnicodeCharacter{0432}{\ensuremath{\mathrm{PDF\;TODO}}}  % Symbol в
\DeclareUnicodeCharacter{0435}{\ensuremath{\mathrm{PDF\;TODO}}}  % Symbol е
\DeclareUnicodeCharacter{0438}{\ensuremath{\mathrm{PDF\;TODO}}}  % Symbol и
\DeclareUnicodeCharacter{043C}{\ensuremath{\mathrm{PDF\;TODO}}}  % Symbol м
\DeclareUnicodeCharacter{0440}{\ensuremath{\mathrm{PDF\;TODO}}}  % Symbol р
\DeclareUnicodeCharacter{0442}{\ensuremath{\mathrm{PDF\;TODO}}}  % Symbol т
\DeclareUnicodeCharacter{041F}{\ensuremath{\mathrm{PDF\;TODO}}}  % Symbol П
\DeclareUnicodeCharacter{0BA8}{\ensuremath{\mathrm{PDF\;TODO}}}  % Symbol ந
\DeclareUnicodeCharacter{0BBF}{\ensuremath{\mathrm{PDF\;TODO}}}  % Symbol  ி
\DeclareUnicodeCharacter{1100}{\ensuremath{\mathrm{PDF\;TODO}}}  % Symbol ᄀ
\DeclareUnicodeCharacter{11F9}{\ensuremath{\mathrm{PDF\;TODO}}}  % Symbol ᇹ

\DeclareUnicodeCharacter{1D48}{\ensuremath{^d}}  % Symbol ᵈ
\DeclareUnicodeCharacter{1D50}{\ensuremath{^m}}  % Symbol ᵐ
\DeclareUnicodeCharacter{1D56}{\ensuremath{^p}}  % Symbol ᵖ
\DeclareUnicodeCharacter{1D62}{\ensuremath{_i}}  % Symbol ᵢ
\DeclareUnicodeCharacter{1D63}{\ensuremath{_r}}  % Symbol ᵣ
\DeclareUnicodeCharacter{1D64}{\ensuremath{_u}}  % Symbol ᵤ
\DeclareUnicodeCharacter{1D9C}{\ensuremath{^c}}  % Symbol ᶜ

% The command `\text` requires `\usepackage{amsmath}` and the command
% `\textasciiacute` requires `\usepackage{textcomp}`.

\DeclareUnicodeCharacter{2032}{\ensuremath{\text{\textasciiacute}}}  % Symbol ′

\DeclareUnicodeCharacter{2070}{\ensuremath{^0}}  % Symbol ⁰
\DeclareUnicodeCharacter{2074}{\ensuremath{^4}}  % Symbol ⁴
\DeclareUnicodeCharacter{2075}{\ensuremath{^5}}  % Symbol ⁵
\DeclareUnicodeCharacter{2076}{\ensuremath{^6}}  % Symbol ⁶
\DeclareUnicodeCharacter{2077}{\ensuremath{^7}}  % Symbol ⁷
\DeclareUnicodeCharacter{2078}{\ensuremath{^8}}  % Symbol ⁸
\DeclareUnicodeCharacter{2079}{\ensuremath{^9}}  % Symbol ⁹
\DeclareUnicodeCharacter{207A}{\ensuremath{^+}}  % Symbol ⁺
\DeclareUnicodeCharacter{207B}{\ensuremath{^-}}  % Symbol ⁻
\DeclareUnicodeCharacter{207F}{\ensuremath{^n}}  % Symbol ⁿ
\DeclareUnicodeCharacter{2080}{\ensuremath{_0}}  % Symbol ₀
\DeclareUnicodeCharacter{2081}{\ensuremath{_1}}  % Symbol ₁
\DeclareUnicodeCharacter{2082}{\ensuremath{_2}}  % Symbol ₂
\DeclareUnicodeCharacter{2083}{\ensuremath{_3}}  % Symbol ₃
\DeclareUnicodeCharacter{2084}{\ensuremath{_4}}  % Symbol ₄
\DeclareUnicodeCharacter{2085}{\ensuremath{_5}}  % Symbol ₅
\DeclareUnicodeCharacter{2086}{\ensuremath{_6}}  % Symbol ₆
\DeclareUnicodeCharacter{2087}{\ensuremath{_7}}  % Symbol ₇
\DeclareUnicodeCharacter{2088}{\ensuremath{_8}}  % Symbol ₈
\DeclareUnicodeCharacter{2089}{\ensuremath{_9}}  % Symbol ₉
\DeclareUnicodeCharacter{208A}{\ensuremath{_+}}  % Symbol ₊
\DeclareUnicodeCharacter{208B}{\ensuremath{_-}}  % Symbol ₋
\DeclareUnicodeCharacter{2091}{\ensuremath{_e}}  % Symbol ₑ
\DeclareUnicodeCharacter{2095}{\ensuremath{_h}}  % Symbol ₕ
\DeclareUnicodeCharacter{2096}{\ensuremath{_k}}  % Symbol ₖ
\DeclareUnicodeCharacter{2097}{\ensuremath{_l}}  % Symbol ₗ
\DeclareUnicodeCharacter{2098}{\ensuremath{_m}}  % Symbol ₘ
\DeclareUnicodeCharacter{2099}{\ensuremath{_n}}  % Symbol ₙ
\DeclareUnicodeCharacter{209A}{\ensuremath{_p}}  % Symbol ₚ
\DeclareUnicodeCharacter{209B}{\ensuremath{_s}}  % Symbol ₛ
\DeclareUnicodeCharacter{209C}{\ensuremath{_t}}  % Symbol ₜ

\DeclareUnicodeCharacter{2113}{\ensuremath{\mathpzc{l}}} % Symbol ℓ

% The command `\mathbbm` requires `\usepackage{bbm}`.
\DeclareUnicodeCharacter{2115}{\ensuremath{\mathbbm{N}}}  % Symbol ℕ
\DeclareUnicodeCharacter{211A}{\ensuremath{\mathbbm{Q}}}  % Symbol ℚ
\DeclareUnicodeCharacter{211D}{\ensuremath{\mathbbm{R}}}  % Symbol ℝ
\DeclareUnicodeCharacter{2124}{\ensuremath{\mathbbm{Z}}}  % Symbol ℤ

\DeclareUnicodeCharacter{2135}{\ensuremath{\aleph}}           % Symbol ℵ
\DeclareUnicodeCharacter{2190}{\ensuremath{\leftarrow}}       % Symbol ←
\DeclareUnicodeCharacter{2192}{\ensuremath{\rightarrow}}      % Symbol →
\DeclareUnicodeCharacter{2194}{\ensuremath{\leftrightarrow}}  % Symbol ↔

% The command `\twoheadrightarrow` requires `\usepackage{amssymb}`.
\DeclareUnicodeCharacter{21A0}{\ensuremath{\twoheadrightarrow}}  % Symbol ↠

\DeclareUnicodeCharacter{21A6}{\ensuremath{\mapsto}}  % Symbol ↦
\DeclareUnicodeCharacter{21A6}{\ensuremath{\mapsto}}  % Symbol ↦

% The command `\restriction` requires `\usepackage{amssymb}`.
\DeclareUnicodeCharacter{21BE}{\ensuremath{\restriction}}  % Symbol ↾

\DeclareUnicodeCharacter{21D0}{\ensuremath{\Leftarrow}}       % Symbol ⇐
\DeclareUnicodeCharacter{21D2}{\ensuremath{\Rightarrow}}      % Symbol ⇒
\DeclareUnicodeCharacter{21D4}{\ensuremath{\Leftrightarrow}}  % Symbol ⇔

\DeclareUnicodeCharacter{2200}{\ensuremath{\forall}}      % Symbol ∀
\DeclareUnicodeCharacter{2203}{\ensuremath{\exists}}      % Symbol ∃
\DeclareUnicodeCharacter{2204}{\ensuremath{\not\exists}}  % Symbol ∃
\DeclareUnicodeCharacter{2205}{\ensuremath{\emptyset}}    % Symbol ∅
\DeclareUnicodeCharacter{2208}{\ensuremath{\in}}          % Symbol ∈
\DeclareUnicodeCharacter{2209}{\ensuremath{\not\in}}      % Symbol ∉
\DeclareUnicodeCharacter{220B}{\ensuremath{\ni}}          % Symbol ∋
\DeclareUnicodeCharacter{220C}{\ensuremath{\not\ni}}      % Symbol ∌

% The command `\blacksquare` requires `\usepackage{amssymb}`.
\DeclareUnicodeCharacter{220E}{{\tiny \ensuremath{\blacksquare}}} % Symbol ∎

\DeclareUnicodeCharacter{220F}{{\tiny \ensuremath{\prod}}}  % Symbol ∏
\DeclareUnicodeCharacter{2211}{\ensuremath{\sum}}           % Symbol ∑
\DeclareUnicodeCharacter{2218}{\ensuremath{\circ}}          % Symbol ∘
\DeclareUnicodeCharacter{2219}{\ensuremath{\bullet}}        % Symbol ∙
\DeclareUnicodeCharacter{221E}{\ensuremath{\infty}}         % Symbol ∞
\DeclareUnicodeCharacter{2223}{\ensuremath{\mid}}           % Symbol ∣
\DeclareUnicodeCharacter{2225}{\ensuremath{\parallel}}      % Symbol ∥
\DeclareUnicodeCharacter{2227}{\ensuremath{\wedge}}         % Symbol ∧
\DeclareUnicodeCharacter{2228}{\ensuremath{\vee}}           % Symbol ∨
\DeclareUnicodeCharacter{2229}{\ensuremath{\cap}}           % Symbol ∩
\DeclareUnicodeCharacter{222A}{\ensuremath{\cup}}           % Symbol ∪
\DeclareUnicodeCharacter{222B}{\ensuremath{\int}}           % Symbol ∫
\DeclareUnicodeCharacter{2234}{\ensuremath{\therefore}}     % Symbol ∴

% The command `\dblcolon` requires `\usepackage{mathtools}`.
\DeclareUnicodeCharacter{2237}{\ensuremath{\dblcolon}} % Symbol ∷

\newcommand{\dotminus}{\mathop{\mbox{$-^{\hspace{-.5em}\cdot}\,$}}}
\DeclareUnicodeCharacter{2238}{\ensuremath{\dotminus}}  % Symbol ∸

\DeclareUnicodeCharacter{223C}{\ensuremath{\sim}}       % Symbol ∼
\DeclareUnicodeCharacter{2243}{\ensuremath{\simeq}}     % Symbol ≃
\DeclareUnicodeCharacter{2245}{\ensuremath{\cong}}      % Symbol ≅
\DeclareUnicodeCharacter{2248}{\ensuremath{\approx}}    % Symbol ≈
\DeclareUnicodeCharacter{224D}{\ensuremath{\asymp}}       % Symbol ≍ : OLIVIER
\DeclareUnicodeCharacter{2250}{\ensuremath{\doteq}}     % Symbol ≐

% The command `\coloneqq` requires `\usepackage{mathtools}`.
\DeclareUnicodeCharacter{2254}{\ensuremath{\coloneqq}}  % Symbol ≔

\newcommand{\eqq}{\stackrel{\text{\tiny ?}}{=}}
\DeclareUnicodeCharacter{225F}{\ensuremath{\eqq}}  % Symbol ≟

\DeclareUnicodeCharacter{2260}{\ensuremath{\neq}}         % Symbol ≠
\DeclareUnicodeCharacter{2261}{\ensuremath{\equiv}}       % Symbol ≡
\DeclareUnicodeCharacter{2262}{\ensuremath{\not\equiv}}   % Symbol ≢
\DeclareUnicodeCharacter{2264}{\ensuremath{\le}}          % Symbol ≤
\DeclareUnicodeCharacter{2265}{\ensuremath{\ge}}          % Symbol ≥
\DeclareUnicodeCharacter{226E}{\ensuremath{\nless}}       % Symbol ≮
\DeclareUnicodeCharacter{226F}{\ensuremath{\ngtr}}        % Symbol ≯
\DeclareUnicodeCharacter{2270}{\ensuremath{\nleq}}        % Symbol ≰
\DeclareUnicodeCharacter{2271}{\ensuremath{\ngeq}}        % Symbol ≱
\DeclareUnicodeCharacter{227A}{\ensuremath{\prec}}        % Symbol ≺
\DeclareUnicodeCharacter{227C}{\ensuremath{\preceq}}      % Symbol ≼
\DeclareUnicodeCharacter{2282}{\ensuremath{\subset}}      % Symbol ⊂
\DeclareUnicodeCharacter{2284}{\ensuremath{\not\subset}} % Symbol ⊄
\DeclareUnicodeCharacter{2283}{\ensuremath{\supset}}      % Symbol ⊃
\DeclareUnicodeCharacter{2286}{\ensuremath{\subseteq}}    % Symbol ⊆
\DeclareUnicodeCharacter{228E}{\ensuremath{\uplus}}       % Symbol ⊎
\DeclareUnicodeCharacter{2291}{\ensuremath{\sqsubseteq}}  % Symbol ⊑
\DeclareUnicodeCharacter{2293}{\ensuremath{\sqcap}}       % Symbol ⊓
\DeclareUnicodeCharacter{2294}{\ensuremath{\sqcup}}       % Symbol ⊔
\DeclareUnicodeCharacter{2295}{\ensuremath{\oplus}}       % Symbol ⊕
\DeclareUnicodeCharacter{22A2}{\ensuremath{\vdash}}       % Symbol ⊢
\DeclareUnicodeCharacter{22A4}{\ensuremath{\top}}         % Symbol ⊤
\DeclareUnicodeCharacter{22A5}{\ensuremath{\bot}}         % Symbol ⊥
\DeclareUnicodeCharacter{22C0}{\ensuremath{\bigwedge}}    % Symbol ⋀
\DeclareUnicodeCharacter{22C2}{\ensuremath{\bigcap}}      % Symbol ⋂
\DeclareUnicodeCharacter{22C3}{\ensuremath{\bigcup}}      % Symbol ⋃
\DeclareUnicodeCharacter{22EE}{\ensuremath{\vdots}}       % Symbol ⋮
\DeclareUnicodeCharacter{22EF}{\ensuremath{\cdots}}       % Symbol ⋯

% The command `\iddots` requires `\usepackage{mathdots}`.
\DeclareUnicodeCharacter{22F0}{\ensuremath{\iddots}} % Symbol ⋰

\DeclareUnicodeCharacter{230A}{\ensuremath{\lfloor}}  % Symbol ⌊
\DeclareUnicodeCharacter{230B}{\ensuremath{\rfloor}}  % Symbol ⌋
\DeclareUnicodeCharacter{2308}{\ensuremath{\lceil}}   % Symbol ⌈
\DeclareUnicodeCharacter{2309}{\ensuremath{\rceil}}   % Symbol ⌉
\DeclareUnicodeCharacter{2329}{\ensuremath{\langle}}  % Symbol 〈
\DeclareUnicodeCharacter{232A}{\ensuremath{\rangle}}  % Symbol 〉

% The command `\Return` requires `\usepackage{keystroke}`.
\DeclareUnicodeCharacter{23CE}{\ensuremath{\Return}}  % Symbol ⏎

% The command `\text` requires `\usepackage{amsmath}` and the command
% `\ding` requires `\usepackage{pifont}`.
\DeclareUnicodeCharacter{2460}{\ensuremath{\text{\ding{172}}}} % Symbol ①
\DeclareUnicodeCharacter{2461}{\ensuremath{\text{\ding{173}}}} % Symbol ②
\DeclareUnicodeCharacter{2462}{\ensuremath{\text{\ding{174}}}} % Symbol ③
\DeclareUnicodeCharacter{2463}{\ensuremath{\text{\ding{175}}}} % Symbol ④
\DeclareUnicodeCharacter{2464}{\ensuremath{\text{\ding{176}}}} % Symbol ⑤
\DeclareUnicodeCharacter{2465}{\ensuremath{\text{\ding{177}}}} % Symbol ⑥
\DeclareUnicodeCharacter{2466}{\ensuremath{\text{\ding{178}}}} % Symbol ⑦
\DeclareUnicodeCharacter{2467}{\ensuremath{\text{\ding{179}}}} % Symbol ⑧
\DeclareUnicodeCharacter{2468}{\ensuremath{\text{\ding{180}}}} % Symbol ⑨

\DeclareUnicodeCharacter{25C5}{\ensuremath{\triangleleft}}  % Symbol ◅

\DeclareUnicodeCharacter{2660}{\ensuremath{\spadesuit}} % Symbol ♠
\DeclareUnicodeCharacter{266D}{\ensuremath{\flat}}      % Symbol ♭
\DeclareUnicodeCharacter{266F}{\ensuremath{\sharp}}     % Symbol ♯

\DeclareUnicodeCharacter{2713}{\ensuremath{\checkmark}}  % Symbol ✓

% The commands `\llbracket` and `\rrbracket` require
% `\usepackage{stmaryrd}`.
\DeclareUnicodeCharacter{27E6}{\ensuremath{\llbracket}}  % Symbol ⟦
\DeclareUnicodeCharacter{27E7}{\ensuremath{\rrbracket}}  % Symbol ⟧

\DeclareUnicodeCharacter{27E8}{\ensuremath{\langle}}          % Symbol ⟨
\DeclareUnicodeCharacter{27E9}{\ensuremath{\rangle}}          % Symbol 〉
\DeclareUnicodeCharacter{27F6}{\ensuremath{\longrightarrow}}  % Symbol % ⟶

\DeclareUnicodeCharacter{2983}{\ensuremath{\mathrm{PDF\;TODO}}}  % Symbol ⦃
\DeclareUnicodeCharacter{2984}{\ensuremath{\mathrm{PDF\;TODO}}}  % Symbol ⦄
\DeclareUnicodeCharacter{2985}{\ensuremath{\mathrm{PDF\;TODO}}}  % Symbol ⦅
\DeclareUnicodeCharacter{2986}{\ensuremath{\mathrm{PDF\;TODO}}}  % Symbol ⦆

% The commands `\llparenthesis` and `\rrparenthesis` require
% `\usepackage{stmaryrd}`.
\DeclareUnicodeCharacter{2987}{\ensuremath{\llparenthesis}}  % Symbol ⦇
\DeclareUnicodeCharacter{2988}{\ensuremath{\rrparenthesis}}  % Symbol ⦈

\DeclareUnicodeCharacter{29F5}{\ensuremath{\setminus}}   % Symbol ⧵

\DeclareUnicodeCharacter{2A06}{\ensuremath{\bigcup}}  % Symbol ⋃
\DeclareUnicodeCharacter{2C7C}{\ensuremath{_j}}       % Symbol ⱼ

\DeclareUnicodeCharacter{D7B0}{\ensuremath{\mathrm{PDF\;TODO}}} % Symbol ힰ

% The command `\mathbbm{1}` requires `\usepackage{bbm}`.
\DeclareUnicodeCharacter{1D7D9}{\ensuremath{\mathbbm{1}}}  % Symbol 𝟙


%%%% MOI


\DeclareUnicodeCharacter{22AC}{\ensuremath{\neg\vdash}}
\DeclareUnicodeCharacter{279C}{\ensuremath{\to}}
\DeclareUnicodeCharacter{1D442}{ }
\DeclareUnicodeCharacter{2000}{ }
\DeclareUnicodeCharacter{1D440}{M}
\DeclareUnicodeCharacter{1D442}{O}
\DeclareUnicodeCharacter{20E3}{2}
\DeclareUnicodeCharacter{1D45A}{m}
\DeclareUnicodeCharacter{1D45B}{n}

\usepackage{pifont}
\DeclareUnicodeCharacter{2714}{\ding{52}}
\DeclareUnicodeCharacter{2705}{\ding{52}}
\DeclareUnicodeCharacter{2717}{\ding{55}}
\DeclareUnicodeCharacter{274C}{\ding{53}}


%% Tolérant tant que macros-markdown.tex n'est pas dans le dépôt :
\InputIfFileExists{macros-markdown}{}{\typeout{*** macros-markdown.tex introuvable: ignore ***}}


%\usepackage{PDFdraftcopy}


%-- Pour état des chapitres


%\newcommand\ETAT{\draftstring{BROUILLON}}
%
%\newcommand\brouillon{\renewcommand\ETAT{\draftstring{BROUILLON}}}
%
%\newcommand\final{\renewcommand\ETAT{\draftstring{\rien}}}
%
%\newcommand\travail{\renewcommand\ETAT{\draftstring{CHANTIER}}}
%
%\newcommand\bazar{\renewcommand\ETAT{\draftstring{BAZAR TOTAL}}}
%
%\final



%%%%%%%%%%%%%%%%%%%%%%%%%%%%%%%%%%%%%%%%%%%%%%%%%
%
%
%                      TRADUCTION EN COURS
%
%%%%%%%%%%%%%%%%%%%%%%%%%%%%%%%%%%%%%%%%%%%%%%%%%


% === volé Amaury



\usepackage{stmaryrd}
\usepackage{ifthen}
\usepackage{mathtools}
\usepackage{dsfont}
\PASSIAPXPROOF{\usepackage{thmtools}}
\usepackage{longtable}
\usepackage{mathdots}
\usepackage{booktabs}


%%%%%%%%%%%%%%%%%%%%%%%%%%%%%%%%%%%%%%%%%%%%%%%%%
%
%
%                      TIKZ STUFFS
%
%
%%%%%%%%%%%%%%%%%%%%%%%%%%%%%%%%%%%%%%%%%%%%%%%%%

%%% volé pour dessin des GPAC: PAS SUR QUE BESOIN de TOUT

\usepackage{tikz}
\usetikzlibrary{calc}
\usepgflibrary{arrows}
\usetikzlibrary{fit}
\usetikzlibrary{patterns}
\usetikzlibrary{decorations.pathreplacing}
\usetikzlibrary{shapes.multipart}
\usetikzlibrary{shapes.geometric}
\usetikzlibrary{shapes.misc}
\usetikzlibrary{positioning}
\usetikzlibrary{graphs}
\usetikzlibrary{topaths}
\usetikzlibrary{arrows.meta}
\usetikzlibrary{decorations.pathreplacing}
\usetikzlibrary{decorations.markings}
\usetikzlibrary{decorations.pathmorphing}
%
 \usetikzlibrary{calc}
 \usetikzlibrary{automata}
\usetikzlibrary{chains}
\usetikzlibrary{scopes}
 \usetikzlibrary{positioning}
 \usetikzlibrary{through,backgrounds}
\usepackage{circuitikz}
% % pour accolades




%%%%%%%%%%%%%%%%%%%%%%%%%%%%%%%%%%%%%%%%%%%%%%%%%
%
%
%                      MACROS locales
%
%
%%%%%%%%%%%%%%%%%%%%%%%%%%%%%%%%%%%%%%%%%%%%%%%%%


%====================
%                      Moi
%=====================


\newcommand\myaddress{ 
{\sc LIX, CNRS, Ecole Polytechnique, Institut Polytechnique de Paris, }\\ 91128 Palaiseau Cedex, FRANCE \\
{\small Olivier.Bournez@lix.polytechnique.fr}}



%====================
%                      un peu de trucs sales
%=====================

%=  pour style lipics et ses trucs

\providecommand\ifnotPOLYCHAPTERMODE[1]{#1}
\providecommand\ifPOLYCHAPTERMODE[1]{}

\ifnotPOLYCHAPTERMODE{
\ifdefined\thechapter
\providecommand\spnewtheorem[4]{\newtheorem{#1}{#2}[chapter]}
\providecommand\spnewtheoremi[5]{\newtheorem{#1}[#5]{#2}}
\else
\providecommand\spnewtheorem[4]{\newtheorem{#1}{#2}}
\providecommand\spnewtheoremi[5]{\newtheorem{#1}[#5]{#2}}
\fi
}
\ifPOLYCHAPTERMODE{
\providecommand\spnewtheorem[4]{\newtheorem{#1}{#2}}
\providecommand\spnewtheoremi[5]{\newtheorem{#1}[#5]{#2}}
}
\providecommand\nolipics[1]{#1}




%====================
%                      un peu de trucs franchement sales
%=====================


% on l'inclus pas ici.
%
%% defaults

\includeversion{NOTLFMI2020}
\includeversion{NOTCSE304}
\excludeversion{NOTDIFFUSIONFINALE}

%% Booleans (dont teste si existe ou pas, pour construire des macros)
% \SAFEMODE
% -> \DIFFUSIONFINALE (via entête-cours)
% -> \PAGEDEGARDE (via entête-cours)
% \POLYMODE
% \POLYCHAPTERMODE

%%%%%%%%%%%%%%%%%%%%%%%%%%%%%%%%%%%%%%%%%%%%%%%%%
%
%
%                      Hack pour DIFFUSION FINALE/ LFMI / Fin provisoire
%
%
%%%%%%%%%%%%%%%%%%%%%%%%%%%%%%%%%%%%%%%%%%%%%%%%%

%\def\DIFFUSIONFINALE{}

\providecommand\ifnotmodeDIFFUSIONFINALE[1]{
\ifdefined\DIFFUSIONFINALE
\else
#1
\fi
}

\providecommand\ifmodeDIFFUSIONFINALE[1]{
\ifdefined\DIFFUSIONFINALE
%% 2021: Comprend pas mais assure
\ifdefined\SAFEMODE
\else
#1
\fi
%#1
\fi
}



\ifnotmodeDIFFUSIONFINALE{\includeversion{NOTDIFFUSIONFINALE}}
\ifnotmodeDIFFUSIONFINALE{\newcommand\FINPROVISOIRELFMI[1]{}}

\providecommand\FINPROVISOIRELFMI{
  \section{Bibliographic notes}

  The current text is highly based (sometimes copied) from:
  
  \bibliographystyle{alpha}
  \bibliography{/Users/bournez/bibliographie/Bib-Files/bournez-avantbib2DOI,/Users/bournez/bibliographie/Bib-Files/perso}
  \end{document}
  }

%\def\POLYCHAPTERMODE{}

\providecommand\ifnotPOLYCHAPTERMODE[1]{
\ifdefined\POLYCHAPTERMODE
\else
#1
\fi
}

\providecommand\ifmodePOLYCHAPTERMODE[1]{
\ifdefined\POLYCHAPTERMODE
#1
\fi
}


%%%%%%%%%%%%%%%%%%%%%%%%%%%%%%%%%%%%%%%%%%%%%%%%%
%
%
%                      POUR INDICATIONS: ETAT
%
%
%%%%%%%%%%%%%%%%%%%%%%%%%%%%%%%%%%%%%%%%%%%%%%%%%


\newcommand\ABSTRACTINTERNAL[1]{
  \montruc[color=black!30]{
  \hrule
  \hrule
  \danger   ABSTRACT:\\[0.5cm] #1

  \vspace{0.5cm}
  \hrule
  \hrule
  }
  }

\newcommand\ABSTRACT[1]{
  \newpage
  \montruc[color=black!20]{
  \hrule
  \hrule
  \danger   ABSTRACT:\\[0.5cm] #1

  \vspace{0.5cm}
  \hrule
  \hrule
  }
  }

  
\newcommand\PARTIEFINALISEE{
  \newpage
  \montruc[color=black!15]{
  \hrule
  \hrule
  \danger   FINALISE A PARTIR ICI
  \hrule
  \hrule
  }
  }
\newcommand\PARTIESEMIFINALISEE{
  \newpage
  \montruc[color=black!15]{
  \hrule
  \hrule
  \danger   1/2 FINALISE A PARTIR ICI
  \hrule
  \hrule
  }
  }


\newcommand\PARTIENONFINALISEE{
  \newpage
  \montruc[color=black!30]{
  \hrule
  \hrule
  \danger   NON FINALISE A PARTIR ICI
  \hrule
  \hrule
  }
  }


%%%%%%%%%%%%%%%%%%%%%%%%%%%%%%%%%%%%%%%%%%%%%%%%%
%
%
%                      Affichage titre du cours: générique
%
%
%%%%%%%%%%%%%%%%%%%%%%%%%%%%%%%%%%%%%%%%%%%%%%%%%


  \newcommand\POLYWARNING[1]{
    
    \chapter*{Preliminaries}

    
\begin{center} 
  These are notes for the course: \\[1cm]
  \textbf{#1}
\end{center}

\bigskip
\fbox{\fbox{
\begin{minipage}{10cm}
\begin{center}
These notes change from days to days.  \\
Please do  not distribute. \\
\end{center}
\end{minipage}
}}

\bigskip
\begin{center}
  \textbf{
    Version of \today.
    }
\end{center}

 % This document is not stable.

\vfill
\begin{center}
Any comment (even about typos, orthography) is welcome: \\
send an email to bournez@lix.polytechnique.fr
\end{center}
}


\newcommand\MYCOURSEON[2]{ 
	%% External self-references
	\ifnotPOLYMODE{
		\ifdefined\MAINAUXFILE
			\externaldocument{\MAINAUXFILE}
		\fi
		}
	
	\ifdefined\PASLAPREMIEREFOIS
		\newpage
	\else
		\makeindex

     		\begin{document}
	\fi
	\gdef\PASLAPREMIEREFOIS{}

	\ifdefined\XPOLY
		\XPAGEDEGARDE{\TITRECOURS\ifPOLYCHAPTERMODE{\\[0.5cm] \Large \bf \LANGUE{Chapter:}{Chapitre:}  #2}}{\SUBTITRECOURS}
		\ifPOLYCHAPTERMODE{
		\renewcommand\thesection{\arabic{section}}
		\chapter*{#2}
		}
	\else
	
      \title{ \ifnotmodeDIFFUSIONFINALE{Course  #1 -}
      \ifPOLYCHAPTERMODE{\LANGUE{Chapter:}{Chapitre:} }
      #2}
      \author{
        Olivier Bournez% \inst{1}
        \\[0.5cm] %{\myaddress}
      }
      \date{Version of \today}
      % \institute{\myaddress}
      




      \ifnotPOLYCHAPTERMODE{
      	\maketitle
		}
      \ifPOLYCHAPTERMODE{\ifdefined\DOPAGEDEGARDE
      			\maketitle
		\fi
		\renewcommand\thesection{\arabic{section}}
%		\renewcommand\thesubsection{\arabic{section}.\ara}
      		\chapter*{#2}}

      % \begin{abstract}
      % \end{abstract}
      \fi
      


      \DOINDEXPRINT
      
      %% 2021: (meme si je comprends pas)
      \ifmodeDIFFUSIONFINALE{
		\excludeversion{recherche}
	}


    }



%%%%%%%%%%%%%%%%%%%%%%%%%%%%%%%%%%%%%%%%%%%%%%%%%
%
%
%                      Affichage titre du cours: variantes
%
%
%%%%%%%%%%%%%%%%%%%%%%%%%%%%%%%%%%%%%%%%%%%%%%%%%

    
\newcommand\MYCOURSE[1]{
  \MYCOURSEON{\csname numerocours#1\endcsname}{\csname nomcours#1\endcsname}
  
    \ifnotmodeDIFFUSIONFINALE{}%\listoftodos}

  }
  
  \newcommand\MYCOURSEnew[1]{
  \MYCOURSEON{???}{#1}
  
  \ifnotmodeDIFFUSIONFINALE{}%\listoftodos}

  }

%   ---------------------------
%   Pour compatilité chapter
%  ----------------------------

\newcommand\PASPREMIEREPAGE[1]{}
\providecommand\mychapter[1]{% \PASPREMIEREPAGE{\newpage} 
  \renewcommand\PASPREMIEREPAGE[1]{##1}  \hrule\hrule \section*{#1}
  \hrule\hrule \  \\[1cm]}
\providecommand\chapter[1]{\PASPREMIEREPAGE{\newpage} 
  \renewcommand\PASPREMIEREPAGE[1]{##1}  \hrule\hrule \section*{#1}
  \hrule\hrule \  \\[1cm]}

%   ---------------------------
%   if not sub
%  ----------------------------


\providecommand\IFNOTSUB[1]{#1}
\IFNOTSUB{
  
%\documentclass{article}
\usepackage{versions}
\includeversion{pasweb}
%\input{macros}

}


 
%%%%%%%%%%%%%%%%%%%%%%%%%%%%%%%%%%%%%%%%%%%%%%%%%
%
%
%               Déclaration de cours
%
%  (l'argument 1 = numéro sert pas)
%
%%%%%%%%%%%%%%%%%%%%%%%%%%%%%%%%%%%%%%%%%%%%%%%%%


\newcommand\COURS[3]{
  \expandafter\newcommand\csname numerocours#2\endcsname{#1}
  \expandafter\newcommand\csname nomcours#2\endcsname{#3}
}



\newcommand\NANCY[3]{
  \expandafter\newcommand\csname numerocours#2\endcsname{#1}
  \expandafter\newcommand\csname nomcours#2\endcsname{#3}
}


\newcommand\INFquatrecentdouze[3]{
  \expandafter\newcommand\csname numerocours#2\endcsname{#1}
  \expandafter\newcommand\csname nomcours#2\endcsname{#3}
}

\newcommand\INFcinqsixun[3]{
  \expandafter\newcommand\csname numerocours#2\endcsname{#1}
  \expandafter\newcommand\csname nomcours#2\endcsname{#3}
}



\newcommand\MPRI[3]{
  \expandafter\newcommand\csname numerocours#2\endcsname{#1}
  \expandafter\newcommand\csname nomcours#2\endcsname{#3}
}

\newcommand\COURSEINCLUDE[2][]{
  \renewcommand\IFNOTSUB[1]{}
  \renewcommand\MYCOURSE[1]{\chapter{\ifnotmodeDIFFUSIONFINALE{#1} \csname nomcours##1\endcsname}
    \ifnotmodeDIFFUSIONFINALE{\mychapter{Fichier: #2}}
  }
   \renewcommand\MYCOURSEnew[1]{\chapter{\ifnotmodeDIFFUSIONFINALE{#1} ##1 }
    \ifnotmodeDIFFUSIONFINALE{\mychapter{Fichier: #2}}
  }
%  \renewcommand\FINSUBMPRI{}
  \input{#2}
}

\newcommand\STRESSCOURSEINCLUDE[1]{
  \COURSEINCLUDE[!STRESS!]{#1}
  }

  \newcommand\subCOURSEINCLUDE[1] {
    \NDLR{BEGIN:  J'inclus/ Je répète ICI #1 }
    \input{INCLUDE/#1}
    \NDLR{END:  J'inclus/ Je répète ICI #1 }
}

%%%%%%%%%%%%%%%%%%%%%%%%%%%%%%%%%%%%%%%%%%%%%%%%%
%
%
%                      POUR INDICATIONS: MESSAGES
%
%
%%%%%%%%%%%%%%%%%%%%%%%%%%%%%%%%%%%%%%%%%%%%%%%%%


% ancienne version
% \newcommand\montruc[1]{\begin{center} {\fbox{\fbox{{#1}
%         }}} \end{center}}




\newcommand\montruc[2][]{
  \ifnotmodeDIFFUSIONFINALE{
  \begin{maboitetruc}[#1]{\danger} #2
  \end{maboitetruc}
%%  \todo[inline,color=green!20,caption={2do}, ssss#1]{\begin{minipage}{\textwidth-4pt}
%\todo[inline,color=green!5%,caption={2do}
%]{\begin{minipage}{\textwidth-4pt}
%    #2\end{minipage}}
}
}

\newcommand\montrucdanger[3][]{
  \ifnotmodeDIFFUSIONFINALE{
                    \montruc[#1]{\danger #2: #3}
                    }
}

\newcommand\montrucdangermargin[3][]{
  \ifnotmodeDIFFUSIONFINALE{
  \montruc[#1]{\danger #2: #3}
\marginpar{\danger #2: #3}
}
 }



%%%%%%%%%%%%%%%%%%%%%%%%%%%%%%%%%%%%%%%%%%%%%%%%%
%
%
%                      GESTION DES TRUCS MARQUES VERSIONS
%
%
%%%%%%%%%%%%%%%%%%%%%%%%%%%%%%%%%%%%%%%%%%%%%%%%%

\makeatletter
\renewcommand\beginmarkversion{\montrucdanger[colback=yellow!10]{=== begin NOT}{}\@Vs@sffbox{\@currenvir$>$}}
 \renewcommand\endmarkversion{\@Vs@sffbox{$<$\@currenvir}\montrucdanger[colback=yellow!10]{end=== NOT}{}}
\makeatother 


%%%%%%%%%%%%%%%%%%%%%%%%%%%%%%%%%%%%%%%%%%%%%%%%%
%
%
%                      VARIANTES DE MESSAGES
%
%
%%%%%%%%%%%%%%%%%%%%%%%%%%%%%%%%%%%%%%%%%%%%%%%%%

 \newcommand\NDLR[1]{\montrucdanger[colback=blue!20]{NDLR}{#1}}
 \newcommand\NLDR[1]{\NDLR{#1}}
 \newcommand\VOIRAUSSI[1]{\montrucdanger[colback=yellow!20]{VOIR AUSSI}{#1}}
 \newcommand\TODO[1]{\montrucdanger[colback=green!20]{TODO}{#1}}
 \newcommand\ATTENTION[1]{\montruc[colback=brown!20]{\danger ATTENTION#1}}
 \newcommand\COMPRENDPAS[1]{\montrucdanger[colback=yellow!20]{JE
     COMPRENDS PAS CA}{#1}}
  \newcommand\POSSIBLE[1]{\montrucdanger[colback=yellow!10]{
      Serait possible}{#1}}
 \newcommand\ENLEVE[1]{\montrucdanger[colback=gray!20]{
      ENLEVE}{#1}}
 
 
  \newcommand\DEPLACE[1]{
    \vspace{0.5cm}
    \hrule
    \hrule
      \vspace{0.5cm}
    \montrucdangermargin[colback=black!40]{DEPLACE/MISSING CONCEPT
    }{#1}
      \vspace{0.5cm}
    \hrule
    \hrule
      \vspace{0.5cm}
  }

\newcommand\BESOINDE[1]{\montrucdangermargin[colback=pink!20]{ATTTENTION
     BESOIN DE}{#1}}
\newcommand\MISSINGCONCEPT[1]{\montrucdangermargin[colback=pink!20]{MISSING CONCEPT
     }{#1}}
\newcommand\SAUTE[1]{\montrucdangermargin[colback=pink!20]{SAUTE
  }{#1}}
\newcommand\SHOULDBEHERE[1]{\BEGINTEMPENVCOLOR{blue!30}\montrucdanger[colback=pink!10]{
SHOULD
    BE HERE
  }{    
    \vspace{1cm}
    
    \centerline{$\vdots$}
    
   \centerline{ #1}
    
      \centerline{$\vdots$}
          
    \vspace{1cm}
    }\ENDTEMPENVCOLOR}
\newcommand\COULDBEHERE[1]{\montrucdangermargin[colback=pink!20]{COULD
    BE HERE
  }{#1}}



%%%%%%%%%%%%%%%%%%%%%%%%%%%%%%%%%%%%%%%%%%%%%%%%%
%
%
%                      POUR INDICATIONS: SOURCES
%
%
%%%%%%%%%%%%%%%%%%%%%%%%%%%%%%%%%%%%%%%%%%%%%%%%%


\newcommand\SOURCEURL[1]{
  \ifnotmodeDIFFUSIONFINALE{
    \marginpar{ \textcolor{magenta}{Source: \url{#1}} }
    }
}

\newcommand\SOURCETXT[1]{
  \ifnotmodeDIFFUSIONFINALE{
    \marginpar{ \textcolor{magenta}{Source: {#1}} }
    }
}

\newcommand\SOURCE[1]{
  \nocite{#1}
  \ifnotmodeDIFFUSIONFINALE{
    \marginpar{ \textcolor{magenta}{Source: \cite{#1}} }
    }
}


%--------------------------------
%--- Raccourci Sources particulières  ---
%--------------------------------

\newcommand\SOURCEJECH{\SOURCE{jech2003set}}
\newcommand\SOURCEKECHRIS{\SOURCE{kechris2012classical}}
\newcommand\SOURCENOTESSETTHEORY{\SOURCE{moschovakis2006notes}}
\newcommand\SOURCEKOZENCT{\SOURCE{LivreKozenTC}}
\newcommand\SOURCEMARKER{\SOURCE{marker2002descriptive}}
\newcommand\SOURCEMOSCHO{\SOURCE{moschovakis2009descriptive}}



%%%%%%%%%%%%%%%%%%%%%%%%%%%%%%%%%%%%%%%%%%%%%%%%%
%
%
%                      POUR INDICATIONS: DELIRES/COMMENTAIRES
%
%
%%%%%%%%%%%%%%%%%%%%%%%%%%%%%%%%%%%%%%%%%%%%%%%%%


%%%%%% ABOUT PLAN

%\newcommand\PLANBUT[1]{\montruc{\begin{minipage}{0.8\textwidth} #1\end{minipage}}}


%%%%%% ABOUT RECHERCHE


\ifnotmodeDIFFUSIONFINALE{
\newenvironment{recherche}
{}{}
%{\color{blue} BEGIN-RECHERCHE:
%%\ifmodeDIFFUSIONFINALE{diffusion finale}
%%\ifnotmodeDIFFUSIONFINALE{not diffusion finale}
%%	\ifdefined\SAFEMODE
%%		SAFEMODE
%%	\else
%%		NOTSAFEMODE
%%	\fi
%  \todo[inline, color=blue]{recherche ici} }{
%  END-RECHERCHE \color{black} }
\tcolorboxenvironment{recherche}{colback=blue!5!white, colframe=red!75!black,fonttitle=\bfseries, colbacktitle=blue!85!black,enhanced,
attach boxed title to top center={yshift=-2mm},
  title={FAIRE RECHERCHE ICI}}
}


%%%%%% ABOUT DELIRE


%\ifnotmodeDIFFUSIONFINALE{
%\newenvironment{delire}{\color{violet} BEGIN-DELIRE: }{
%  END-DELIRE \color{black} }
%}


%%%%%% ABOUT PERSO


%\ifnotmodeDIFFUSIONFINALE{
%\newenvironment{perso}{\color{brown}
%  BEGIN-PERSO: }{
%  END-PERSO\color{black} }
%}



%%%%%% ABOUT COMMENTAIRE

\ifnotmodeDIFFUSIONFINALE{
\newenvironment{commentaire}{\begin{maboitecommentaire}[colback=brown!30]{Commentaire} Commentaire :}{\end{maboitecommentaire}}
\newcommand\COMMENTAIRE[2][Commentaire ]{
\begin{maboitecommentaire}[colback=brown!30]{#1}
Commentaire : #2
\end{maboitecommentaire}
}
}

%%%%%% ABOUT HORSUJET


%\ifnotmodeDIFFUSIONFINALE{
%\newenvironment{probablementhorsujet}{\color{gray}
%  BEGIN-PROBABLEMENT HORS SUJET: }{
%  END-PROBABLEMENT HORS SUJET \color{black} }
%}

%%%%%% ABOUT BONUS


\ifnotmodeDIFFUSIONFINALE{
\newenvironment{pasbesoinpourlinstant}{NOT NEEDEED FOR
  THIS SECTION: }{ }
    \tcolorboxenvironment{pasbesoinpourlinstant}{colback=green!5!white, colframe=gray!75!black,fonttitle=\bfseries, colbacktitle=green!85!black,enhanced,
attach boxed title to top left={yshift=-2mm},
  title={PAS BESOIN POUR L'INSTANT}}
}

%%%%%% ABOUT BONUS


\ifnotmodeDIFFUSIONFINALE{
\newenvironment{bonus}{BONUS: }{
  }
  \tcolorboxenvironment{bonus}{colback=gray!5!white, colframe=gray!75!black,fonttitle=\bfseries, colbacktitle=gray!85!black,enhanced,
attach boxed title to top left={yshift=-2mm},
  title={BONUS}}
}

%%%%%% ABOUT BONUS


\ifnotmodeDIFFUSIONFINALE{
\newcommand\ATTENTIONREPETE[1]{
  \ATTENTION{Repete (attention modif): #1}
  \begin{remarkolivier}
Repete (attention modif): #1
\end{remarkolivier}
}
}


%%%%%%%%%%%%%%%%%%%%%%%%%%%%%%%%%%%%%%%%%%%%%%%%%
%
%
%                      MACROS Papiers Paulin
%
%
%%%%%%%%%%%%%%%%%%%%%%%%%%%%%%%%%%%%%%%%%%%%%%%%%


\newcommand{\x}{\times}
\newcommand{\s}{\sigma}
\newcommand\sep{.}
\newcommand\bzero{\zero}
\newcommand\bun{\un}
\renewcommand{\a}{\alpha}
\def\hd{{\sf hd}}
\def\tl{{\sf tl}}
\def\cons{{\sf cons}}
\def\Cons{{\sf Cons}}
\def\Pr{{\sf Pr}}
\def\Op{{\sf Op}}
\def\Rel{{\sf Rel}}
\def\Equal{{\sf Equal}}
\def\Eval{{\sf Eval}}
\def\Gate{{\sf Gate}}
\def\Selection{{\sf Selection}}
\def\Result{{\sf Result}}
\def\Circuit{{\sf Circuit}}
\def\next{{\sf next}}
\def\comp{{\sf comp}}
\def\finalcomp{{\sf finalcomp}}
\def\Pick{{\sf Pick}}
\def\Hd{{\sf Hd}}
\def\Tl{{\sf Tl}}
\def\RevHd{{\sf RevHd}}
\def\Node{{\sf Node}}
\def\Length{{\sf Length}}
\def\Tape{{\sf Tape}}
\def\Head{{\sf Head}}
\def\Time{{\sf Time}}
\def\unpad{{\sf unpad}}
\def\concat{{\sf concat}}
\def\rg{{\sf right}}
\def\lf{{\sf left}}
\def\st{{\sf state}}
\def\tope{{\sf top}}


%%%%%%%%%%%%%%%%%%%%%%%%%%%%%%%%%%%%%%%%%%%%%%%%%
%
%
%                      MACROS POLY INF561
%
%
%%%%%%%%%%%%%%%%%%%%%%%%%%%%%%%%%%%%%%%%%%%%%%%%%

\newcommand\buglst[1]{#1}

%\newcommand\mydefref[1]{\expandafter\newcommand\csname #1\endcsname{\ref{#1}}}
%\mydefref{refthrisc}
%\mydefref{refthcinqhuit}
%\mydefref{refthcinqsix}
%\mydefref{refthfondamcaract}
%\mydefref{refdefcircuit}
%%% sinon, fait à la main
%\newcommand\refchapalgo{\ref{chapalgo}}
%\newcommand\refchapmodeles{\ref{chapmodeles}}
%\newcommand\refchappreliminaires{\ref{chappreliminaires}}
%\newcommand\refchapcircuits{\ref{chapcircuits}}
%\newcommand \refpropfondamen{\ref{propfondamen}}
%\newcommand\refthhierspace{ref{thhierspace}}
%\newcommand\refchapcomplexitetemps{\ref{chapcomplexitetemps}}
%\newcommand\refdefverifun{\ref{defverifun}}
%\newcommand\refdefverif{\ref{defverif}}
%\newcommand\refthmachinevs{ref{thmachinevs}}
%\newcommand\reflemnotationconfig{\ref{lemnotationconfig}}
%\newcommand\reflemfondam{\ref{lemfondam}}
%\newcommand\refthtroisat{\ref{thtroisat}}
%\newcommand\refthfondam{\ref{thfondam}}
%\newcommand\refthsansram{\ref{thsansram}}
%\newcommand\refdefsram{\ref{defsram}}
%!TEX encoding = UTF-8 Unicode

%% TRUCS DU POLY INF561

\ifnotSLIDE{
\LANGUE{
\spnewtheoremi{exemple}{\NDLR{ECRIT EXEMPLE AU LIEU DE
  EXAMPLE}}{\bfseries}{\rmfamily}{theorem} 
\spnewtheoremi{remarque}{\NDLR{ECRIT REMARQUE AU LIEU DE REMARK}}{\bfseries}{\rmfamily}{theorem} 
}
{
\spnewtheoremi{exemple}{Exemple}{\bfseries}{\rmfamily}{theorem} 
\spnewtheoremi{remarque}{Remarque}{\bfseries}{\rmfamily}{theorem} 
}
}
            
            \newcommand\PERSOSOURCEPOSSIBLE[1]{\PERSOSOURCE{#1}\PERSOCOMMENTAIRE{#1}}
\newcommand\PERSOPOSSIBLE[1]{\POSSIBLE{#1}}
% \newcommand\PERSOPOSSIBLEbug[1]{\NDLR{\textbf{-- possible --
%       enleve car pb compilation}}}
\newcommand\PERSOETAIT[1]{\NDLR{ETAIT: #1}}
%   \begin{minipage}{10cm}
% \NDLR{\textbf{\tiny  -- était --      #1 -- était
%     --
%   \end{minipage}
% }}}

\providecommand\nl{}

\newcommand\PERSOSOURCE[1] { \SOURCETXT{#1}}
\newcommand\PERSOCOMMENTAIRE[1]{\NDLR{#1}}
\newcommand\PERSOCOMMENTAIREd[2]{\NDLR{#1 \textbf{#2}}}


\newenvironment{dessinpb}
{%\shorthandoff{:}
%\selectlanguage{english}
\begin{center}\begin{tikzpicture}[baseline=(current bounding
      box.north)]
%babelbug with tikz
}
{\end{tikzpicture}
%\selectlanguage{frenchb}
\end{center}}



\newcommand\PERSOMANQUE[1]{\NDLR{\textbf{-- manque --
      #1 --manque--}}}

\newcommand\PERSOENLEVE[1]{\ENLEVE{#1}}






%%%%%%%%%%%%%%%%%%%%%%%%%%%%%%%%%%%%%%%%%%%%%%%%%
%
%
%                      MACROS POLY INF412
%
%
%%%%%%%%%%%%%%%%%%%%%%%%%%%%%%%%%%%%%%%%%%%%%%%%%

%\mydefref{refEEXTtruc}
%\mydefref{refEEXTchappreliminaires}
%\mydefref{refEEXTsecarbrebinaire}
%\mydefref{refEEXTsectermes}
%\mydefref{refEEXTchappredicats}
%\mydefref{refEXTchappredicats}
%\mydefref{refEXTchappreliminaires}
%\mydefref{refEEXTdefTuring}
%\mydefref{refEEXTpropnondet}
%\mydefref{refEEXTchapcompletude}
%\mydefref{refEEXTdefconfigalternative}
%\mydefref{refEEXTchapmachinesdeTuring}
%\mydefref{refEEXTpbdedecision}
%\mydefref{refEEXTchapcalc}
%\mydefref{refEEXTincompletude}
%\mydefref{refEEXTthhierspace}
%\mydefref{refEEXTchapcomplexitetemps}

\tikzstyle{nd} = [shape=circle,draw]


\newcommand\mmod{mod}

  %%% VALUATIONS
\newcommand\mmV{v}
\newcommand\mmVb{\overline{\mmV}}
\newcommand\mmVe{\mmVb}

%%% Corrections

\newcommand\exoref[1]{\ref{exo:#1-stt}}

\newenvironment{correction}[1]
{\par\medskip\noindent \label{exo:#1-cor} \hbox{\bf Exercice \ref{exo:#1-stt} 
  \   (page \pageref{exo:#1-stt}). }}
{\par\medskip}

% PROBLEMES DE DECISION

\newcommand\DECISIONint[3]{
%Le problème de décision #1
\begin{itemize}
\item[\textbf{Donnée}:] #2
\item[\textbf{Réponse}:] #3
\end{itemize}
}

\newcommand\INDECIDABLE[5]{
\DECISION{#1}{#2}{#3}

\begin{#4} \label{#5}
Le problème $#1$ % ATTENTION ENLEVE $#1$ en 2020. REMIS 2023
est indécidable.
\end{#4}
}


\newcommand\INDECIDABLEenglish[5]{
\DECISIONenglish{#1}{#2}{#3}

\begin{#4} \label{#5}
The problem $#1$
is undecidable.
\end{#4}
}

\newcommand\INDECIDABLEq[5]{
\DECISION{#1}{#2}{#3}

\begin{#4} \label{#5}
Le problème $#1$
est ?
\end{#4}
}

\newcommand\INDECIDABLEqenglish[5]{
\DECISION{#1}{#2}{#3}

\begin{#4} \label{#5}
The problem $#1$
is ?
\end{#4}
}



\newcommand\DECISIONintenglish[3]{
%Le problème de décision #1
\begin{itemize}
\item[\textbf{Input}:] #2
\item[\textbf{Answer}:] #3
\end{itemize}
}

\newcommand\DECISIONi[4]{
\item $#1$
%Le problème de décision #1
\DECISIONint{#1}{#2}{#3}
}

\newcommand\DECISIONienglish[4]{
\item $#1$
%Le problème de décision #1
\DECISIONintenglish{#1}{#2}{#3}
}


\newcommand\DECISION[3]{
\begin{definition}[$#1$]\ 

%Le problème de décision #1
\DECISIONint{#1}{#2}{#3}
\end{definition}
}


\newcommand\DECISIONenglish[3]{
\begin{definition}[$#1$]\ 
%
%Le problème de décision #1
\DECISIONintenglish{#1}{#2}{#3}
\end{definition}
}



\newcommand\NPC[4]{ 
\DECISION{#1}{#2}{#3}

\begin{#4}
Le problème $#1$
est $\NP$-complet. \LANGUE{\myindex{NP-completeness}}{\myindex{NP-completude@$\NP$-complétude}}
\end{#4}
}


\newcommand\NPCenglish[4]{ 
\DECISIONenglish{#1}{#2}{#3}
\begin{#4}
The problem $#1$
is $\NP$-complete. \LANGUE{\myindex{NP-completeness}}{\myindex{NP-completude@$\NP$-complétude}}
\end{#4}
}


\newcommand\NPCp[5]{
\DECISION{#1}{#2}{#3}

\begin{#4}[#5]
Le problème $#1$
est $\NP$-complet.
\end{#4}
}


\newcommand\NPCpenglish[5]{
\DECISIONenglish{#1}{#2}{#3}

\begin{#4}[#5]
The problem $#1$
is $\NP$-complete.
\end{#4}
}




\newcommand\BRUNOCOMMENTAIRE[1]{\NDLR{Bruno commentaire: #1}}

\ifnotSLIDE{
%\spnewtheorem{exercicec}{Exercice}{\bfseries}{\rmfamily}
%\spnewtheorem{exerciced}{Exercice}{\bfseries}{\rmfamily}
%\spnewtheorem{exercicedc}{Exercice}{\bfseries}{\rmfamily}
%\spnewtheorem{exercicecenglish}{Exercice}{\bfseries}{\rmfamily}
%\spnewtheorem{exercicedenglish}{Exercice}{\bfseries}{\rmfamily}
%\spnewtheorem{exercicedcenglish}{Exercice}{\bfseries}{\rmfamily}
\spnewtheorem{notation}{Notation}{\bfseries}{\rmfamily}
}

%%% INDEX
 \newcommand\motnouv[1]{\emph{#1}\index{#1}}%\SIGNALEMOTNOUV{#1}}%\index{#1}}
\newcommand\INTmotnouv[1]{\emph{#1}}
\newcommand\myindex[1]{\index{#1}}
%\marginpar{-#1=}

\newcommand\synonyme[2]{\kindex{#1|see{#2}}}
\newcommand\indexM[1]{\index{$#1$}}
\newcommand\idx[1]{#1\index{#1}}
\newcommand\idxm[1]{#1\index{$#1$}}
\newcommand\idxM[1]{$#1$\index{$#1$}}
%\newcommand\idxmBUG[1]{#1}
\newcommand\idxmm[2]{#1\index{#2@$#1$}}
\newcommand\idxi[2]{#1\index{#2}}
%\newcommand\motnouv[1]{\INTmotnouv{#1}\index{#1}}
 \newcommand\motnouvi[2]{\INTmotnouv{#1}\index{#2}} % indexage manuel
 \newcommand\motnouva[2]{\INTmotnouv{#1}\index{#2@#1}} % indexage: 2ième
%                                 % argumet = version sans accent
\newcommand\motnouvm[1]{\INTmotnouv{#1}} %indexage manuel fait ailleurs


\newenvironment{PERSOVERBATIM}{}{}
 \newcommand{\joNP}[3]{{ Problème \textbf{ { \motnouv{##1}}:}
            \begin{itemize}
            \item  \textbf{ Instance:} {##2} 
            \item   \textbf{ Question:} {##3}
            \end{itemize}
          }}

%%%%%%%%%%%%%%%%%%%%%%%%%%%%%%%%%%%%%%%%%%%%%%%%%
%
%
%                      POUR COMPATIBILITE AVEC SLIDES & AUTRE
%
%
%%%%%%%%%%%%%%%%%%%%%%%%%%%%%%%%%%%%%%%%%%%%%%%%%


\newcommand\defin[1]{\textbf{#1}}


%%%   COMPATIBILITE AVEC SLIDES

\newcommand\soustitre[1]{\textbf{#1}}


%% COMPATIBILITE AVEC MA THESE

\providecommand\motnouv[1]{\emph{#1}}
\providecommand\motnouvi[2]{\emph{#1}}

\newcommand\papiernorm[1]{\myop{norm}}

%%% OLIVIER
\newcommand\olivier[1]{ \ATTENTION{Olivier: #1}}


%%%%%%%%%%%%%%%%%%%%%%%%%%%%%%%%%%%%%%%%%%%%%%%%%
%
%
%                      Page de Garde Cours X
%
%
%%%%%%%%%%%%%%%%%%%%%%%%%%%%%%%%%%%%%%%%%%%%%%%%%


\includeversion{metdate}

\newcommand\XPAGEDEGARDE[2]{
\thispagestyle{empty}
\title{{#1} \\[2cm] %\\[3cm] 
{\large #2}
}
\author{
 Olivier Bournez% \inst{1}
\\[0.5cm] 
bournez@lix.polytechnique.fr%{\myaddress}
}
%\date{
%\vspace{2cm}
%\includegraphics[height=2.5cm]{/Users/bournez/lib/LaTeX/Perso/FIG-COMMONS/logo_X_new}
%}
\begin{metdate}
\date{
\vspace{2cm}
Version \LANGUE{of}{du} \today \\[1cm] 
\includegraphics[height=3.5cm]{/Users/bournez/lib/LaTeX/Perso/FIG-COMMONS/logo_X_new}
}
\end{metdate}
\maketitle
}


%%%%%%%%%%%%%%%%%%%%%%%%%%%%%%%%%%%%%%%%%%%%%%%%%
%
%
%                      MACROS UTILES
%
%
%%%%%%%%%%%%%%%%%%%%%%%%%%%%%%%%%%%%%%%%%%%%%%%%%

%%%                      Macros
 
%-------------
%--- vectors  ---
%-------------

 
\newcommand\vectorl[1]{{\mathbf#1}}
\newcommand\tu[1]{\vectorl{#1}}

\newcommand\vp{\vectorl{p}}
\newcommand\va{\vectorl{a}}
\newcommand\vb{\vectorl{b}}
\newcommand\vc{\vectorl{c}}
\newcommand\vf{\vectorl{f}}
\newcommand\vn{\vectorl{n}}
\newcommand\vx{\vectorl{x}}
\newcommand\vX{\vectorl{X}}
\newcommand\vy{\vectorl{y}}
\newcommand\vz{\vectorl{z}}
\newcommand\ve{\vectorl{e}}
\newcommand\vv{\vectorl{v}}
\newcommand\vr{\vectorl{r}}
\newcommand\vu{\vectorl{u}}
\newcommand\vA{\vectorl{A}}
\newcommand\vB{\vectorl{B}}
\newcommand\vC{\vectorl{C}}
\newcommand\vD{\vectorl{D}}

%-------------
%--- overline ---
%-------------


\newcommand\ox{\overline{x}}
\newcommand\oy{\overline{y}}
\newcommand\oc{\overline{c}}
\newcommand\oa{\overline{a}}
\newcommand\ow{\overline{w}}
\newcommand\od{\overline{d}}
\newcommand\ob{\overline{b}}
\newcommand\oP{\overline{P}}
\newcommand\oee{\overline{e}}
\newcommand\ot{\overline{t}}
\newcommand\ou{\overline{u}}
\newcommand\ov{\overline{v}}
\newcommand\oz{\overline{z}}
\newcommand\am{\overline{a}}
\newcommand\om{\overline{m}}
\newcommand\oj{\overline{j}}

%----------------
%--- ensembles ---
%----------------

\newcommand\Q{\mathbb{Q}}
\newcommand\R{\mathbb{R}}
\newcommand\Z{\mathbb{Z}}
\providecommand\C{\mathbb{C}}
\newcommand\N{\mathbb{N}}
\newcommand\K{\mathbb{K}}
\newcommand\F{\mathbb{F}}

\newcommand\dyadic{\mathbb{D}}
\newcommand\Zd{{\Z_2}}
\newcommand\Zp{{\Z_p}}
%
\newcommand\RP{\mathbb{R}^{>0}}
\newcommand\QP{\mathbb{Q}^{>0}}


%----------------
%--- lettres cal ---
%----------------

\newcommand\mK{\mathcal{K}}
\newcommand\mD{\mathcal{D}}
\newcommand\mA{\mathcal{A}}
\newcommand\mL{\mathcal{L}}
\newcommand\mX{\mathcal{X}}
\newcommand\mG{\mathcal{G}}
\newcommand\mE{\mathcal{E}}
\newcommand\mH{\mathcal{H}}
\newcommand\mR{\mathcal{R}}
\newcommand\mF{\mathcal{F}}
\newcommand\mJ{\mathcal{J}}
\newcommand\mO{\mathcal{O}}
\newcommand\mP{\mathcal{P}}
\newcommand\mN{\mathcal{N}}
\newcommand\mS{\mathcal{S}}
\newcommand\mM{\mathcal{M}}
\newcommand\mC{\mathcal{C}}
\newcommand\mV{\mathcal{V}}
\newcommand\mI{\mathcal{I}}
\newcommand\mT{\mathcal{T}}

\newcommand\fC{\mathcal{C}}

%----------------
%--- lettres frak ---
%----------------


\newcommand\kM{\mathfrak{M}}
\newcommand\kN{\mathfrak{N}}
\newcommand\kU{\mathfrak{U}}
\newcommand\kg{\mathfrak{g}}



%--------------------------
%--- Façon noter noms classes ---
%---------------------------

\newcommand{\cp}[1]{\operatorname{#1}}


%--------------------------
%--- Calculabilité  ---
%---------------------------

% -- classes
\LANGUE{
\newcommand\RE{\cp{CE}}
}
{\newcommand\RE{\cp{RE}}
}

\newcommand\Rec{\cp{D}}
\newcommand\PR{\cp{PR}}
\newcommand\Gregn{\mE_n}


%--------------------------
%--- Réductions  ---
%---------------------------

\newcommand\lem{\le_{m}}
\newcommand\lelog{\le_{L}}
\newcommand\equivlog{\equiv_{L}}



%--------------------------
%--- Complexité  ---
%---------------------------

% COMPLEXITE

% -- P
\newcommand{\Ptime}{\cp{P}}      % P
\newcommand\PTIME{\Ptime}
\newcommand\PTIMEs[1]{{\PTIME}_{#1}}
\newcommand\PTIMEsp[1]{{\PTIME}_{#1}^0}
\renewcommand\P{\PTIME}
\newcommand{\FPtime}{\cp{FPTIME}}

% -- NP
\newcommand\NP{\cp{NP}}
\newcommand{\NPtime}{\NP}
\newcommand\NPs[1]{\cp{NP}_{#1}}
\newcommand{\FNP}{\cp{FNP}}

\newcommand\NDP{\cp{NDP}}
\newcommand\coNDP{\cp{coNDP}}
\newcommand\NDPs[1]{\cp{NDP}_{#1}}

% -- coNP
\newcommand\coNP{\cp{coNP}}

% -- PSPACE
\newcommand\PSPACE{\cp{ PSPACE}}
\newcommand\NPSPACE{\cp{ NPSPACE}}
\newcommand{\Pspace}{\PSPACE}
\newcommand{\FPspace}{\cp{FPSPACE}}


% -- LOGSPACE
\newcommand\LSPACE{\cp{LOGSPACE}}
\newcommand\NLSPACE{\cp{NLOGSPACE}}
\newcommand\coNLSPACE{\cp{coNLOGSPACE}}
\newcommand\coNL{\cp{coNL}}

% -- Autres
\newcommand\EXPTIME{\cp{EXPTIME}}
\newcommand\NEXPTIME{\cp{NEXPTIME}}
\newcommand\EXPSPACE{\cp{EXPSPACE}}
\newcommand\PH{\cp{PH}}
\newcommand\NC{\cp{NC}}
\newcommand\AC{\cp{AC}}
\newcommand\PPAD{\cp{\mathbf{PPAD}}}
\newcommand\PLS{\cp{\mathbf{PLS}}}

% -- Reverse math
\newcommand{\WKL}{\ensuremath{\docheck{\mathsf{WKL}_0}}}
\newcommand{\ACA}{\ensuremath{\docheck{\mathsf{ACA}_0}}}
\newcommand{\ATR}{\ensuremath{\docheck{\mathsf{ATR}_0}}}
\newcommand{\RCAO}{\ensuremath{\docheck{\mathsf{RCA}_0}}}
\newcommand{\PiCA}{\ensuremath{\docheck{\Pi^1_1\!\text{-}\mathsf{CA}_0}}}


% -- circuits
\newcommand\CIRCUIT[2]{\cp{CIRCUIT(#1,#2)}}
\newcommand\UCIRCUIT[2]{\cp{UCIRCUIT(#1,#2)}}

% -- Classes proba
\newcommand\PP{\cp{PP}}
\newcommand\RPP{\cp{RP}}
\newcommand\ZPP{\cp{ZPP}}
\newcommand\coRP{\cp{coRP}}
\newcommand\BPP{{\cp{BPP}}}

% -- IP
\newcommand\IP{\cp{IP}}
\newcommand\PCP{\cp{PCP}}

% -- non-uniforme
\newcommand\nuP{\mathbb{P}}
\newcommand\nuPs[1]{\mathbb{P}_{#1}}
\newcommand\nuPsp[1]{\mathbb{P}_{#1}^0}
\newcommand\nuNP{\mathbb{N}\mathbb{P}}
%\newcommand\poly{/{poly}}
\newcommand\Ppoly{\PTIME_{\poly}}
\newcommand\NPpoly{\NP_{\poly}}

% -- conditions d'uniformité
\newcommand\Luniforme{\cp{L}}
\newcommand\BCuniforme{\cp{BC}}


% Classes de circuits
\newcommand\TC{\cp{TC}}
\newcommand\LFP{\cp{LFP}}
\newcommand\FAC{\cp{FAC}}
\newcommand\FACC{\cp{FACC}}
\newcommand\FTC{\cp{FTC}}
\newcommand\FNC{\cp{FNC}}

% Logiques temorelles
\newcommand\CTL{\cp{CTL}}
\newcommand\LTL{\cp{LTL}}

% -- mesure temps/espace/taille

\newcommand\DTIME{\cp{DTIME}}
\newcommand\TIME[1]{\cp{TIME(#1)}}
\newcommand\TIMEs[2]{\cp{TIME_{#2}(#1)}}
\newcommand\TIMEsp[2]{\cp{TIME^0_{#2}(#1)}}
\newcommand\NTIME[1]{\cp{NTIME(#1)}}
\newcommand\NTIMEs[2]{\cp{NTIME_{#2}(#1)}}
\newcommand\NTIMEsp[2]{\cp{NTIME_{#2}^0(#1)}}
\newcommand\DSPACE{\cp{DSPACE}}
\newcommand\SPACE[1]{\cp{SPACE(#1)}}
\newcommand\SPACEs[2]{\cp{SPACE_{#2}(#1)}}
\newcommand\NSPACE[1]{\cp{NSPACE(#1)}}
\newcommand\NSPACEs[2]{\cp{NSPACE_{#2}(#1)}}
\newcommand\coNSPACE[1]{\cp{coNSPACE(#1)}}
\newcommand\TIMEON[2]{\cp{TIME(#1,#2)}}
\newcommand\SPACEON[2]{\cp{SPACE(#1,#2)}}
\newcommand\SIZE[1]{\cp{ size(#1)}}
\newcommand\ATIME{\cp{ATIME}}
\newcommand\ASPACE{\cp{ASPACE}}


%% autre:
\newcommand\LH{\cp{LH}}
\newcommand\FO{\cp{FO}}
\newcommand\SO{\cp{SO}}
\newcommand\ESOhorn{\cp{ESO_{horn}}}
\newcommand\SOhorn{\cp{SO_{horn}}}
\newcommand\ACzero{\cp{AC_0}}
\newcommand{\LINSPACE}{\cp{LINSPACE}}
\newcommand{\mFLINSPACE}{\mF_{\cp{LINSPACE}}}
\newcommand{\LOGSPACE}{\cp{LOGSPACE}}
\newcommand\ALOGTIME{\cp{ALOGTIME}}
\newcommand\RF{\cp{RF}}


\newcommand\BOOL[1]{\cp{{BOOLEAN(#1)}}}
\newcommand\DET{\cp{ DET}}

% % -- dirty


\newcommand{\cPspace}{\cp{\sharp PSPACE}}
\newcommand{\cAPtime}{\cp{\sharp APTIME}}
%\newcommand{\FPspace}{\cp{FPSPACE}}
%\newcommand{\FPspaceN}{\cp{FPSPACE}^{+}}
%\newcommand{\FPspaceN}{\cp{FPSPACE}}



%-------------------------
%--- Problèmes de décision  ---
%--------------------------

\newcommand\decisionname[1]{\cp{#1}}


\newcommand\Luniv{\decisionname{L_{univ}}}
\newcommand\HALTINGPROBLEM{\decisionname{HALTING-PROBLEM}}
\newcommand\SAT{\decisionname{SAT}}
\newcommand\MAXtSAT{\decisionname{MAX3SAT}}
\newcommand\REACH{\decisionname{REACH}}
\newcommand\CIRCUITVALUE{\decisionname{CIRCUITVALUE}}
\newcommand\MONOTONECIRCUITVALUE{\decisionname{MONOTONECIRCUITVALUE}}
\newcommand\THEOREMS{\decisionname{THEOREMS}}
\newcommand\QCIRCUITSAT{\decisionname{QCIRCUITSAT}}
\newcommand\QSAT{\decisionname{QSAT}}
\newcommand\QBF{\decisionname{QBF}}
\newcommand\SATVALUE{\decisionname{SATVALUE}}
\newcommand\SUCCINTSAT{\decisionname{SUCCINTSAT}}
\newcommand\SUCCINTSATVALUE{\decisionname{SUCCINTSATVALUE}}
\newcommand\GEOGRAPHY{\decisionname{GEOGRAPHY}}
\newcommand\PARITY{\decisionname{PARITY}}
\newcommand\MAJ{\decisionname{MAJ}}
\newcommand\CVP{\CIRCUITVALUE}
\newcommand\BOOLEANPROGRAMVALUE{BOOLEAN~PROGRAM~VALUE}

\newcommand\CLIQUE{\cp{CLIQUE}}
\LANGUEinv{
\newcommand\RECOUVREMENTDESOMMETS{\cp{RECOUVREMENT~DE~SOMMETS}}
\newcommand\RS{\RECOUVREMENTDESOMMETS}
}{
\newcommand\RECOUVREMENTDESOMMETSeng{\cp{VERTEX~COVER}}
\newcommand\RSeng{\RECOUVREMENTDESOMMETSeng}
}

\LANGUEinv{\newcommand\COUPUREMAXIMALE{\cp{COUPURE~MAXIMALE}}}
{\newcommand\COUPUREMAXIMALEeng{\cp{MAXIMAL~CUT}}}
\newcommand\CIRCUITHAMILTONIEN{\CHAMILTONIEN}
\LANGUEinv{\newcommand\VOYAGEURDECOMMERCE{\VCOMMERCE}}
{\newcommand\VOYAGEURDECOMMERCEeng{\VCOMMERCEeng}}

\LANGUEinv{\newcommand\CIRCUITLEPLUSLONG{\cp{CIRCUIT~LE~PLUS~LONG}}}
{\newcommand\CIRCUITLEPLUSLONGeng{\cp{LONGEST~CIRCUIT}}}
\LANGUEinv{\newcommand\SOMMEDESOUSENSEMBLE{\SUBSETSUM}}{
\newcommand\SOMMEDESOUSENSEMBLEeng{\SUBSETSUMeng}}
\LANGUEinv{
\newcommand\SACADOS{\cp{SAC~A~DOS}}}
{\newcommand\SACADOSeng{\cp{KNAPSACK}}}

%français
\LANGUEinv{
\newcommand\CHAMILTONIEN{\cp{CIRCUIT~HAMILTONIEN}}}
{\newcommand\CHAMILTONIENeng{\cp{HAMILTONIAN~CIRCUIT}}}
\LANGUEinv{
\newcommand\VCOMMERCE{\cp{VOYAGEUR~DE~COMMERCE}}}
{\newcommand\VCOMMERCEeng{\cp{TRAVELING~SALESMAN}}}
\LANGUEinv{
\newcommand\kCOLORABILITE{\cp{k-COLORABILITE}}}
{\newcommand\kCOLORABILITEeng{\cp{k-COLORABILITY}}}
\LANGUEinv{
\newcommand\COLORIAGE{\cp{BON~COLORIAGE}}}
{\newcommand\COLORIAGEeng{\cp{GOOD~COLORING}}}
\LANGUEinv{
\newcommand\tCOLORABILITE{\cp{3-COLORABILITE}}}
{\newcommand\tCOLORABILITEeng{\cp{3-COLORABILITY}}}
\LANGUEinv{
\newcommand\SUBSETSUM{\cp{SOMME~DE~SOUS~ENSEMBLE}}}
{\newcommand\SUBSETSUMeng{\cp{SUBSET~SUM}}}
\newcommand\PARTITION{\cp{PARTITION}}
\LANGUEinv{
\newcommand\NOMBREPREMIER{\decisionname{NOMBRE~ PREMIER}}}
{\newcommand\NOMBREPREMIEReng{\decisionname{PRIME~NUMBER}}}
\LANGUEinv{
\newcommand\CODAGE{\decisionname{CODAGE}}}
{\newcommand\CODAGEeng{\decisionname{ENCODING}}}
\newcommand\kSAT{\cp{k-SAT}}
\newcommand\tSAT{\cp{3-SAT}}
\newcommand\NAESAT{\cp{NAESAT}}
\newcommand\NAEtSAT{\cp{NAE3SAT}}
\newcommand\NAEqSAT{\cp{NAE4SAT}}
\newcommand\STABLE{\cp{\LANGUE{INDEPENDANT~SET}{STABLE}}}

%\newcommand\CM{\cp{COUPURE MAXIMALE}}
\LANGUEinv{
\newcommand\POSTpb{\cp{Problème~de~ la~correspondance~de~Post}}}
{\newcommand\POSTpbeng{\cp{Post~correspondence~problem}}}

\LANGUEinv{
\newcommand\PARITE{\cp{PARITE}}}
{\newcommand\PARITEeng{\cp{PARITY}}}


%-------------------------
%--- Modèles parallèles  ---
%--------------------------


% PARALLELISME
\newcommand\RAM{\cp{RAM}}
%
\newcommand\PRAM{\cp{PRAM}}
\newcommand\CRCW{\cp{CRCW}}
\newcommand\CREW{\cp{CREW}}
\newcommand\EREW{\cp{EREW}}


%------------------------
%--- codages  ---
%------------------------


\newcommand\codage[1]{\langle #1 \rangle}
\newcommand\tuple[1]{\codage{#1}}
\newcommand\paire[2]{\codage{#1,#2}}
\newcommand\triplet[3]{\codage{#1,#2,#3}}
\newcommand\quadruplet[4]{\codage{#1,#2,#3,#4}}
\newcommand\cinquplet[5]{\codage{#1,#2,#3,#4,#5}}


%------------------------
%--- descriptive set theory  ---
%------------------------

\newcommand\baire{\mathcal{N}}
\newcommand\peano{\overline{\mathcal{N}}} %% A MOI
\newcommand\cantor{\mathcal{C}}
%
\newcommand\bSigma{\mathbf{\Sigma}}
\newcommand\bPi{\mathbf{\Pi}}
\newcommand\bDelta{\mathbf{\Delta}}
%


%%%%%% ABOUT NOMS QUI CHANGES

\newcommand\WF{\cp{WF}}
\newcommand\WO{\WF}
\newcommand\LO{\cp{LO}}
\newcommand\DLO{\cp{DLO}}
% 
\newcommand\rk{\cp{rank}}


%% Moins propre:

\newcommand\I{\mathbb{I}}
\renewcommand\H{\mathbb{H}}
%
\newcommand\leKB{<_{KB}}
%
\newcommand\finiomega{{<\omega}}
\newcommand\compl{^\frown}
\newcommand\prefix{ ~ prefix~  } 
\newcommand\X{{X}}
\newcommand\Y{\mathcal{Y}}
\newcommand\subsetstrict{\subset}

\renewcommand\complement[1]{#1^{c}}


%------------------------
%--- schémas de récursion ---
%------------------------

%%                      Source papier MFCS 2019


% COMPLEXITY

%% MFCS differe de Clote:
%% Version à taille restreinte

\newcommand\co{\cp{co}}

\newcommand\mFPTIME{\mF_{PTIME}}
\newcommand{\mFPSPACE}{\mF_{SPACE}}
\newcommand{\mFPH}{\mF_{PH}}

\newcommand\SigmaP{\Sigma^P}
\newcommand\DeltaP{\Delta^{P}}
\newcommand\PiP{\Pi^{P}}
\newcommand\coSigmaP{\co\SigmaP}
\newcommand\coPiP{\mF{\co\PiP}}

%% autres classes:


%% Schemas pour complexité implicite

\newcommand\COMP{\cp{ COMP}}
\newcommand\CRN{\cp{ CRN}}
\newcommand\BRN{\cp{ BRN}}
\newcommand\SBRN{\cp{ SBRN}}
\newcommand\BR{\cp{ BR}}
\newcommand\kBR{k-\cp{ BR}}
\newcommand \WBPR{\cp{ WBPR}}
\newcommand \BMIN{\cp{ BMIN}}
\newcommand \BSUM{\cp{ BSUM}}
\newcommand\SBSUM{\cp{ SBSUM}}
\newcommand \BPROD{\cp{ BPROD}}
\newcommand \SBPROD{\cp{ SBPROD}}
\newcommand \SCOMP{\cp{ SCOMP}}
\newcommand \SR{\cp{ SR}}
\newcommand \SRN{\cp{ SRN}}

%% devrait etre défini
\newcommand\WBRN{\cp{ WBRN}}



%-- godel
\newcommand\INTDIV{\cp{INTDIV}}
\newcommand\DIV{\cp{DIV}}
\newcommand\EVEN{\cp{EVEN}}
\newcommand\ODD{\cp{ODD}}
\newcommand\PRIME{\cp{PRIME}}
\newcommand\POWER{\cp{POWER}}
\newcommand\BIT{\cp{BIT}}
\newcommand\negHP{\overline{\cp{HP}}}
\newcommand\HP{\cp{HP}}
\newcommand\negLuniv{\overline{\Luniv}}
\newcommand\VALCOMP{\cp{VALCOMP}}
\newcommand\mWINDOW{\cp{LEGAL}}
\newcommand\mmWINDOW{\cp{WINDOW}}
\newcommand\LENGTH{\cp{LENGTH}}
\newcommand\DIGIT{\cp{DIGIT}}
\newcommand\MATCH{\cp{MATCH}}
\newcommand\MOVE{\cp{MOVE}}
\newcommand\START{\cp{START}}
\newcommand\CELL{\cp{CELL}}
\newcommand\HALT{\cp{HALT}}
\newcommand\SUBST{\cp{SUBST}}
\newcommand\PROOF{\cp{PROOF}}
\newcommand\PROVABLE{\cp{PROVABLE}}
\newcommand\CONSIST{\cp{CONSIST}}



%---------------------
%--- calculus (AP) ---
%---------------------

% \jacobian{f}{x}
\newcommand{\jacobian}[1]{J_{#1}}
%\newcommand{\grad}[1]{\overrightarrow{\operatorname{grad}}~{#1}}
\newcommand{\grad}[1]{\nabla{#1}}
\newcommand{\scalarprod}[2]{{#1}\cdot{#2}}
\newcommand{\bigO}[1]{\mathcal{O}\left(#1\right)}
\newcommand\smallO{o}
\newcommand{\softO}[1]{\tilde{\mathcal{O}}\left(#1\right)}
\newcommand{\taylor}[3]{{T_{#1}^{#2}#3}}
\newcommand{\transpose}[1]{{#1}^T}
\newcommand{\pastsup}[2]{{\sup}_{#1}#2}

\DeclareMathOperator\arctanh{arctanh}



%---------------
%--- norms (AP) ---
%---------------


\newcommand{\inorm}[2]{\left\lVert{#1}\right\rVert_{#2}}
\newcommand{\twonorm}[1]{\inorm{#1}{2}}
\newcommand{\infnorm}[1]{\inorm{#1}{}}
\newcommand{\normsup}[1]{\infnorm{#1}{}}
%
\newcommand\hausdorff{d_H}
%
\newcommand\norm[1]{\infnorm{#1}}


%----------------
%--- topology ---
%----------------

% \openball{pt}{radius}
 \newcommand{\openball}[2]{B_{#2}(#1)}
\newcommand{\closedball}[2]{\ovl{B}_{#2}(#1)}
\newcommand{\closure}[1]{\ovl{#1}}
% \newcommand{\interior}[1]{\mathring{#1}}
% \newcommand{\border}[1]{\partial{#1}}




%-----------------
%--- functions (AP) ---
%-----------------
 
\newcommand{\lambdafun}[2]{(#1)\mapsto{#2}}
\newcommand{\lambdafunex}[2]{#1\mapsto{#2}}
\newcommand{\argmin}{\operatornamewithlimits{argmin}}
\newcommand{\argmax}{\operatornamewithlimits{argmax}}


 
%-----------------------
%--- Machines de Turing  ---
%-----------------------

\newcommand\blanc{\vectorl{B}}



%-----------------
%--- Logique ---
%-----------------
 
\newcommand\sem[1]{{[\semspace[}#1{]\semspace]}}
\newcommand\sems[1]{{[\semspace[}#1{]\semspace]}}
\newcommand\semss[1]{\sem{#1}_{\sigma'}}
\newcommand\semspace{\kern -.25ex}
\newcommand\true{true}
\newcommand\false{false}
%extension élémentaire:
\newcommand\elem{\preceq}
\newcommand\godel[1]{\lceil #1 \rceil}

%-----------------
%--- Model theory ---
%-----------------

\newcommand\hull[1]{\langle #1 \rangle}

%-----------------
%--- Ordinaux ---
%-----------------

\newcommand\omegaunCK{\omega_1^{CK}}
\newcommand\omegaunck{\omegaunCK}

\DeclareMathOperator{\cof}{cof}

%-----------------
%--- Surréels (QG) ---
%-----------------

\newcommand{\On}{\textbf{On}}
\newcommand{\Nobf}{\textbf{No}}
\newcommand{\Nobb}{\ensuremath{\mathbbmss{No}}}
\newcommand{\Nolt}[1]{\ensuremath{\Nobf_{< #1}}}
\newcommand{\Noleq}[1]{\ensuremath{\Nobf_{\leq #1}}}
\newcommand{\Noltbb}[1]{\ensuremath{\Nobb_{< #1}}}
\providecommand\No{\Nobf}
\newcommand\Noalpha{\Nolt{\alpha}}
\newcommand\Nolambda{\Nolt{\lambda}}

\DeclareMathOperator{\length}{length}
\DeclareMathOperator{\supp}{supp} % operateur support


%----------------------
%---Maths  de base (AP)  ---
%----------------------

\DeclareMathOperator{\dom}{dom} % operateur support

\DeclareMathOperator{\size}{size}

\newcommand{\powerset}[1]{\mathcal{P}(#1)}
\newcommand{\card}[1]{{\##1}}
\newcommand\priv{ - }
\newcommand\prive{-}
\renewcommand{\restriction}{\mathord{\upharpoonright}}
\newcommand\restrict[1]{_{\restriction #1}}

% \providecommand\mod{}
%% UNDEFINE:
\let\mod\relax
\let\div\relax
\DeclareMathOperator{\mod}{mod} 
\DeclareMathOperator{\div}{div} 

% partieentier
\newcommand\entieres[1]{\lceil #1 \rceil}
\newcommand\entierei[1]{\lfloor #1 \rfloor}


%----------------------
%---Symbole danger ---
%----------------------


%\providecommand*{\TakeFourierOrnament}[1]{{%
%\fontencoding{U}\fontfamily{futs}\selectfont\char#1}}
%\providecommand*{\danger}{{\Huge \TakeFourierOrnament{66}}}



%%%%%%%%%%%%%%%%%%%%%%%%%%%%%%%%%%%%%%%%%%%%%%%%%
%
%
%                      ASMeries
%
%
%%%%%%%%%%%%%%%%%%%%%%%%%%%%%%%%%%%%%%%%%%%%%%%%%



% CIRCUITS
%-- particuliers
\newcommand\REPLACE{\cp{REPLACE}}
\newcommand\EQUAL{\cp{EQUAL}}
\newcommand\EVALUE{\cp{EVALUE}}
\newcommand\TVALUE{\cp{TVALUE}}
\newcommand\UPDATE{\cp{UPDATE}}
\newcommand\ASSIGN{\cp{ASSIGN}}
\newcommand\PAR{\cp{PAR}}
\newcommand\DO{\cp{DO}}


% ASM
\newcommand\emplacement[2]{(#1,#2)}
\newcommand\contenu[2]{\sem{(#1,#2)}}
\newcommand\affecte[3]{(#1,#2)\leftarrow#3}
\newcommand\ctl{ctl\_state}


%%%%%%%%%%%%%%%%%%%%%%%%%%%%%%%%%%%%%%%%%%%%%%%%%
%
%
%                      GPACqueries
%
%
%%%%%%%%%%%%%%%%%%%%%%%%%%%%%%%%%%%%%%%%%%%%%%%%%



%------------------------
%--- generable zoo (GPAC) ---
%------------------------

\newcommand{\myop}[1]{\operatorname{#1}}
\newcommand{\abs}{\myop{abs}}
\newcommand{\sg}{\myop{sg}}
\newcommand{\ip}[1]{\myop{ip}_{#1}}
\newcommand{\rnd}{\myop{rnd}}
\newcommand{\lil}{\myop{lil}}
\newcommand{\lxh}{\myop{lxh}}
\newcommand{\hxl}{\myop{hxl}}
\newcommand{\plil}{\myop{plil}}
\newcommand{\reach}{\myop{reach}}
%\newcommand{\mreach}{\myop{mreach}}
\newcommand{\watchdog}{\myop{watchdog}}
\newcommand{\monitor}{\myop{monitor}}
%\newcommand{\norm}{\myop{norm}}
\newcommand{\mx}{\myop{mx}}
\newcommand{\mn}{\myop{mn}}
\newcommand{\sample}{\myop{sample}}
\newcommand{\select}{\myop{select}}
\newcommand{\ssine}[1]{\sin_{#1}}
\newcommand{\clamp}{\myop{clamp}}

\newcommand{\fnsg}[3]{tanh((#1)*(#2)*(#3))}
\newcommand{\fnipone}[3]{(1+\fnsg{(#1)-1}{#2}{#3})/2.}
\newcommand{\fnipinf}[1]{#1-sin(2*pi*(#1))/2./pi}%
\newcommand{\fnlil}[5]{%
\fnipone{(#3)-(#1)+1-((#2)-(#1))/8.}{#4+log(1+4./(#2-#1)*(#5)*(#5))+log(1+4)}{8./(#2-#1)}%
*\fnipone{#2-#3+1-(#2-#1)/8.}{#4+log(1+4/(#2-#1)*(#5)*(#5))+log(1+4)}{8./(#2-#1)}%
*4./(#2-#1)*(#5)}
\newcommand{\fnlxh}[5]{\fnipone{(#3)-(#1+#2)/2.+1}{#4+log(1+(#5)*(#5))}{2./(#2-#1)}*(#5)}
\newcommand{\fnhxl}[5]{\fnipone{(#1+#2)/2.-(#3)+1}{#4+log(1+(#5)*(#5))}{2./(#2-#1)}*(#5)}
\newcommand{\fnplilf}[4]{sin(((#4)-((#1)+(#2))/2.+(#3)/4.)*2.*pi/(#3))}
\newcommand{\fnplil}[6]{(#6)*\fnlxh{\fnplilf{#1}{#2}{#3}{#1}}%
{\fnplilf{#1}{#2}{#3}{#1+((#2)-(#1))/4.}}%
{\fnplilf{#1}{#2}{#3}{#4}}%
{#5+2.+log(1+(#6)*(#6))}%
{1./4.+2./((#2)-(#1))}}
\newcommand{\fnssine}[2]{sin(#2)*(1-cos(#2)/cos(#1))}



%--------------------
%--- complexity GPAC ---
%--------------------


%\ifthenelse{\equal{#1}{}}{Empty.}{Nonempty.}
\newcommand{\myclass}[1]{\operatorname{#1}}
% \newcommand{\gcomp}[2]{\ensuremath{\myclass{GCOMP}(#1,#2)}}
% \newcommand{\ggenerable}[1]{\textcolor{red}{\ensuremath{#1-}\text{generable}}}
 \newcommand{\gval}[2][]{\ensuremath{\myclass{GVAL}_{#1}\ifthenelse{\equal{#2}{}}{}{[#2]}}}
 \newcommand{\gpval}[1][]{\ensuremath{\myclass{GPVAL}_{#1}}}
 \newcommand{\gc}[3][]{\ensuremath{\myclass{ATSC}(#2,#3)}}
 \newcommand{\gpc}[1][]{\ensuremath{\myclass{ATSP}}}
\newcommand{\cgc}[1][]{\ensuremath{\myclass{ATS}}}
\newcommand{\gexpc}[1][]{\ensuremath{\myclass{AEXP}}}
\newcommand{\gwc}[2]{\ensuremath{\myclass{AW}(#1,#2)}}
\newcommand{\gpwc}{\ensuremath{\myclass{AWP}}}
\newcommand{\cgwc}{\ensuremath{\myclass{AW}}}
\newcommand{\grc}[3]{\ensuremath{\myclass{AR}(#1,#2,#3)}}
\newcommand{\gprc}{\ensuremath{\myclass{ARP}}}
\newcommand{\gsc}[3]{\ensuremath{\myclass{AS}(#1,#2,#3)}}
\newcommand{\gpsc}{\ensuremath{\myclass{ASP}}}
\newcommand{\goc}[3]{\ensuremath{\myclass{AO}(#1,#2,#3)}}
\newcommand{\gpoc}{\ensuremath{\myclass{AOP}}}
\newcommand{\cgoc}{\ensuremath{\myclass{AO}}}
\newcommand{\guc}[4]{\ensuremath{\myclass{AX}(#1,#2,#3,#4)}}
\newcommand{\gpuc}{\ensuremath{\myclass{AXP}}}
\newcommand{\glc}[2][]{\ensuremath{\myclass{AL}(#2)}}
\newcommand{\gplc}[1][]{\ensuremath{\myclass{ALP}}}
\newcommand{\cglc}[1][]{\ensuremath{\myclass{AL}}}

%\newcommand{\intinterv}[2]{\{#1,\ldots,#2\}}
\newcommand{\intinterv}[2]{\llbracket#1,#2\rrbracket}

\newcommand{\Rgen}{\R_{G}}
\newcommand{\Rpoly}{\R_{P}}

%%%%%%%%%%%%%%%%%%%%%%%%%%%%%%%%%%%%%%%%%%%%%%%%%
%
%
%                      Mes codages obscures
%
%
%%%%%%%%%%%%%%%%%%%%%%%%%%%%%%%%%%%%%%%%%%%%%%%%%




%%% Codages

% \newcommand\gammaoc{\gamma_{[0,1]}}
% \newcommand\gammaonc{\gamma_{\N}}
\newcommand\gammaTMc{\gamma^{TM}_{[0,1]}}
\newcommand\gammaTMcatan{\gamma^{TM}_{(0,1)^2}}
\newcommand\gammaTMnc{\gamma^{TM}_{\N^3}}
\newcommand\gammaTMncd{\gamma^{TM}_{\N^2}}
% \newcommand\gammaTMe{\gamma^{TM}_{*}}
% \newcommand\gammas{\gamma_{sab}}
% \newcommand\gammaord{\gamma_{cantor}}


\newcommand\enccompact{\nu_{X}}
 \newcommand\encinfinite{\nu_\N}
 \newcommand\enc{\nu}
 
%%%%%%%%%%%%%%%%%%%%%%%%%%%%%%%%%%%%%%%%%%%%%%%%%
%
%
%                      DESSINS ET IMAGES
%
%
%%%%%%%%%%%%%%%%%%%%%%%%%%%%%%%%%%%%%%%%%%%%%%%%%


\newcommand\IMAGE[2]{ \begin{center} 
    \includegraphics[#1]{#2}
  \end{center}
  }


  
%%%% POUR FAIRE DESSIN.
\newenvironment{dessinpp}[1]{
\begin{tikzpicture}[baseline=(current bounding
      box.west),#1]
}
{\end{tikzpicture}}




\newenvironment{dessinp}[1]{
\begin{center}\begin{tikzpicture}[baseline=(current bounding
      box.north),#1]
}
{\end{tikzpicture}
\end{center}}


\newenvironment{dessin}{
\begin{center}\begin{tikzpicture}[baseline=(current bounding
      box.north)]
}
{\end{tikzpicture}
\end{center}}



\newenvironment{dessind}{
\begin{tikzpicture}[baseline=(current bounding
      box.north)]
}
{\end{tikzpicture}
}



%%%%%%%%%%%%%%%%%%%%%%%%%%%%%%%%%%%%%%%%%%%%%%%%%
%
%
%                      POUR DESSINS GPAC
%
%
%%%%%%%%%%%%%%%%%%%%%%%%%%%%%%%%%%%%%%%%%%%%%%%%%


%%%% MACROS DE BASE.


\newcommand\constant[2]
{
node (#1)  [shape=circle,draw] {#2};
\draw (#1.east) -- +(0.3,0) coordinate (#1-o) ; 
\draw (#1.west) -- +(-0.3,0) coordinate (#1-i) ; 
}

\newcommand\basicproduct[2]
{
node (#1) {#2};
\draw (#1) +(-0.4,0.4) -- +(0.4,0.4) -- +(0.8,0) -- +(0.4,-0.4) -- +(-0.4,-0.4) --
+(-0.4,0.4);
\draw (#1) ++(0.8,0) -- ++(+0.3,0) coordinate (#1-o) ;
\draw (#1) ++ (-0.4,0.2) -- +(-0.3,0) coordinate (#1-i1) ;
\draw (#1) ++(-0.4,-0.2) -- +(-0.3,0) coordinate (#1-i2) ;
}

\newcommand\product[1]{\basicproduct{#1}{$+\Pi$}}
\newcommand\negativeproduct[1]{\basicproduct{#1}{$-\Pi$}}


%%BEGIN INTERNAL
\newcommand\basicsummer[1]{
coordinate (#1) {};
\draw (#1) +(-0.5,0.6) -- +(0.5,0) -- +(-0.5,-0.6) -- +(-0.5,0.6);
\draw (#1) ++(0.5,0) -- +(+0.3,0) coordinate (#1-o) ;
}

\newcommand\basicintegrator[1]{
coordinate (#1) {};
\draw (#1) +(-0.5,0.6) -- +(0.5,0) -- +(-0.5,-0.6) -- +(-0.5,0.6);
\draw (#1) ++(0.5,0) -- +(+0.3,0) coordinate (#1-o) ;
\draw (#1) +(-0.5,-0.6) -| +(-0.8,0.6) -- +(-0.5,0.6) ;
 \draw (#1) ++(-0.65,0.6) -- +(0,+0.3) coordinate (#1-init) ;
}

\newcommand\splitfour[2]{
\draw (#2) ++ (#1,0.4) -- +(-0.3,0) coordinate (#2-i1) ;
\draw (#2) ++ (#1,0.15) -- +(-0.3,0) coordinate (#2-i2) ;
\draw (#2) ++ (#1,-0.15) -- +(-0.3,0) coordinate (#2-i3) ;
\draw (#2) ++ (#1,-0.4) -- +(-0.3,0) coordinate (#2-i4) ;
}

\newcommand\splitthree[2]{
\draw  (#2) ++(#1,0.3) -- +(-0.3,0) coordinate (#2-i1) ;
\draw  (#2) ++(#1,0) -- +(-0.3,0) coordinate (#2-i2) ;
\draw  (#2) ++(#1,-0.3) -- +(-0.3,0) coordinate (#2-i3) ;
}

\newcommand\splittwo[2]{
\draw  (#2) ++(#1,0.2) -- +(-0.3,0) coordinate (#2-i1) ;
\draw  (#2) ++(#1,-0.2) -- +(-0.3,0) coordinate (#2-i2) ;
}

\newcommand\splitone[2]{
\draw (#2) ++(#1,0) -- +(-0.3,0) coordinate (#2-i1) ;
}

%% END INTERNAL

\newcommand\summerfour[1]{
\basicsummer{#1}
\splitfour{-0.5}{#1}
}

\newcommand\summerthree[1]{
\basicsummer{#1}
\splitthree{-0.5}{#1}
}

\newcommand\summertwo[1]{
\basicsummer{#1}
\splittwo{-0.5}{#1}
}

\newcommand\summerone[1]{
\basicsummer{#1}
\splitone{-0.5}{#1}
}


\newcommand\integratorfour[1]{
\basicintegrator{#1}
\splitfour{-0.8}{#1}
<}


\newcommand\integratorthree[1]{
\basicintegrator{#1}
\splitthree{-0.8}{#1}
}


\newcommand\integratortwo[1]{
\basicintegrator{#1}
\splittwo{-0.8}{#1}
}


\newcommand\integratorone[1]{
\basicintegrator{#1}
\splitone{-0.8}{#1}
}

%%%%% FIN MACROS




%%%%%%%%%%%%%%%%%%%%%%%%%%%%%%%%%%%%%%%%%%%%%%%%%
%
%
%                      Trucs d'Amaury
%
%
%%%%%%%%%%%%%%%%%%%%%%%%%%%%%%%%%%%%%%%%%%%%%%%%%


%----------------
%--- ref (AP) ---
%----------------

% \newcommand{\defref}[1]{Definition~\ref{#1}}
\newcommand{\lemref}[1]{Lemma~\ref{#1}}
\newcommand{\thref}[1]{Theorem~\ref{#1}}
\newcommand{\amaurycorref}[1]{Corollary~\ref{#1}}
 \newcommand{\figref}[1]{Figure~\ref{#1}}
% \newcommand{\figrefmult}[1]{Figures~{#1}}
% \newcommand{\secref}[1]{Section~\ref{#1}}
% %\renewcommand{\algref}[1]{Algorithm~\ref{#1}}
 \newcommand{\propref}[1]{Proposition~\ref{#1}}
% \newcommand{\exref}[1]{Example~\ref{#1}}
\newcommand{\remref}[1]{Remark~\ref{#1}}


 %-------------------
%--- polynomials ---
%-------------------

\newcommand{\sigmap}[1]{{\Sigma{#1}}}
\newcommand{\degp}[1]{{\operatorname{deg}(#1)}}
\newcommand{\poly}{\operatorname{poly}}


 %----------------------
%--- misc functions ---
%----------------------

\newcommand{\sgn}[1]{\operatorname{sgn}(#1)}
\newcommand{\intp}{\operatorname{int}}
 \newcommand{\fracp}{\operatorname{frac}}
\newcommand{\fiter}[2]{#1^{[#2]}}
\newcommand{\frestrict}[2]{#1\restriction_{#2}}
\newcommand{\indicator}[1]{\mathds{1}_{#1}}
\newcommand{\idfun}{\operatorname{id}}
% \newcommand{\floor}[1]{\left\lfloor#1\right\rfloor}
 \newcommand{\ceil}[1]{\left\lceil#1\right\rceil}
\newcommand{\round}[1]{\left\lfloor#1\right\rceil}
\DeclareMathOperator\erf{erf}


%------------
%--- sets ---
%------------

\newcommand{\A}{\mathbb{A}}
%\newcommand{\D}{\mathbb{D}}
\newcommand{\Ns}{\mathbb{N}^*}
%\newcommand{\C}{\mathbb{C}}
%\newcommand{\K}{\mathbb{K}}
% \MAT{height}{[width]}{[field]}
\newcommand{\MAT}[3]{M_{#1\ifthenelse{\equal{#2}{}}{}{,#2}}\ifthenelse{\equal{#3}{}}{}{\left(#3\right)}}

%\newcommand{\Rcomp}{\R_{c}}
\newcommand{\Rp}{\R_{+}}
\newcommand{\Rs}{\R^{*}}
\newcommand{\Rps}{\R_{+}^{*}}

%\newcommand{\Analytic}{C^\omega}
%\newcommand{\PTIME}{\ensuremath{\operatorname{P}}}
%\newcommand{\EXPTIME}{\ensuremath{\operatorname{EXPTIME}}}
%\newcommand{\NP}{\ensuremath{\operatorname{NP}}}
%\newcommand{\NPTIME}{\NP}
\newcommand{\FP}{\ensuremath{\operatorname{FP}}}
%\newcommand{\PSPACE}{\ensuremath{\operatorname{PSPACE}}}
% \newcommand{\FPSPACE}{\ensuremath{\operatorname{FPSPACE}}}



%%%%%%%%%%%%%%%%%%%%%%%%%%%%%%%%%%%%%%%%%%%%%%%%%
%
%
%                      Trucs de Quentin
%
%
%%%%%%%%%%%%%%%%%%%%%%%%%%%%%%%%%%%%%%%%%%%%%%%%%


\newcommand{\fct}[5]{\ensuremath{#1 : \left\{\begin{array}{c c c}
#2 & \rightarrow & #3\\
#4 & \mapsto & #5
                                             \end{array}\right.}}

%Nouvelle fraction, plus jolie
\newcommand{\f}[2]{	
	\mathchoice%
		{\dfrac{#1}{#2}}
    	{\dfrac{#1}{#2}}
		{\frac{#1}{#2}}
		{\frac{#1}{#2}}
}


%%Ensembles et combinatoires

%Ensemble avec propriété à drotie \{x | blabla\}
\newcommand{\enstq}[2]{\left\{ \left.#1\ \vphantom{#2}\right|\ #2  \right\}}
%meme chose mais avec des crochets
\newcommand{\crotq}[2]{\left[ \left.#1\ \vphantom{#2}\right|\ #2  \right]}
%même chose mais avec des brackets
\newcommand{\bratq}[2]{\left\langle \left.#1\ \vphantom{#2}\right|\ #2  
  \right\rangle}

\newcommand{\intn}[2]{\ensuremath{
		\left\llbracket\, #1\ ;\ #2\,\right\rrbracket
              }}


%%%%%%%%%%%%%%%%%%%%%%%%%%%%%%%%%%%%%%%%%%%%%%%%%
%
%
%               ENVIRONNEMENTS THEOREMS & CO
%
%
%%%%%%%%%%%%%%%%%%%%%%%%%%%%%%%%%%%%%%%%%%%%%%%%%


%----------------
%--- pas dans lipics ---
%-----------------


\ifnotSLIDE{
% commente car bug Fevrier 2026: \ifnotTDEXAM{\spnewtheorem{solution}{Solution}{\bfseries}{\rmfamily}}
\spnewtheorem{openproblem}{Open Problem}{\bfseries}{\rmfamily}
\spnewtheorem{remarkolivier}{Commentaire d'Olivier: }{\bfseries}{\rmfamily}
\spnewtheorem{postulate}{Postulate}{\bfseries}{\rmfamily}
}

%----------------
%--- dans lipics ---
%-----------------
%
\newenvironment{exercice}[1]{\begin{exercise} \label{exo:#1-stt}}{\end{exercise}}
\newenvironment{exerciced}[1]{\begin{exercisee}\label{exo:#1-stt}}{\end{exercisee}}
\newenvironment{exercicec}[1]{\begin{exercise} \label{exo:#1-stt} (\LANGUE{solution on}{corrigé} page \pageref{exo:#1-cor})
}{\end{exercise}}
\newenvironment{exercicedc}[1]{\begin{exercisee} \label{exo:#1-stt} (\LANGUE{solution on}{corrigé} page \pageref{exo:#1-cor})
}{\end{exercisee}}

\newenvironment{slideproposition}{{\bf Proposition}}{}

\ifnotSLIDE{
\nolipics{
\nolncs{\spnewtheorem{theorem}{\LANGUE{Theorem}{Théorème}}{\bfseries}{\rmfamily}}
\nolncs{\spnewtheoremi{definition}{\LANGUE{Definition}{Définition}}{\bfseries}{\rmfamily}{theorem}}
\nolncs{\spnewtheoremi{lemma}{\LANGUE{Lemma}{Lemme}}{\bfseries}{\rmfamily}{theorem}}
\nolncs{\spnewtheoremi{corollary}{\LANGUE{Corollary}{Corollaire}}{\bfseries}{\rmfamily}{theorem}}
\nolncs{\spnewtheoremi{proposition}{Proposition}{\bfseries}{\rmfamily}{theorem}}
\nolncs{\spnewtheoremi{example}{\LANGUE{Example}{Exemple}}{\bfseries}{\rmfamily}{theorem}}
%\spnewtheorem{sketch}{{\bf \LANGUE{Sketch of proof}{Démonstration (principe)}}}{\bfseries}{\rmfamily}
\nolncs{\spnewtheoremi{remark}{\LANGUE{Remark}{Remarque}}{\bfseries}{\rmfamily}{theorem}}
\nolncs{\spnewtheorem{exercise}{\LANGUE{Exercise}{Exercice}}{\bfseries}{\rmfamily}}
\spnewtheorem{exercisee}{\LANGUE{*Exercise}{*Exercice}}{\bfseries}{\rmfamily}
\spnewtheorem{summary}{\LANGUE{Summary}{Résumé}}{\bfseries}{\rmfamily}
 %\spnewtheorem{exerciceenglish}{Exercise}{\bfseries}{\rmfamily}
 %
 \newenvironment{sketch}{{\bf \LANGUE{Sketch of proof}{Démonstration (principe)}}:}{\hfill $\Box$ }
\nolncs{\spnewtheorem{conjecture}{\LANGUE{Conjecture}{Conjecture}}{\bfseries}{\rmfamily}}
 %%
 \spnewtheorem{theoremaprouver}{\LANGUE{Theorem TO BE PROVED}{Théorème (A PROUVER)}}{\bfseries}{\rmfamily}
  \spnewtheorem{conjectureaprouver}{\LANGUE{Conjecture TO BE VERIFIED}{Conjecture (A VERIFIER)}}{\bfseries}{\rmfamily}
 }}
            
 \makeatletter
 \ifnotSLIDE{
  \ifnotTDEXAM{
\@ifundefined{proof}{
	\newenvironment{proof}{{\bf \LANGUE{Proof}{Démonstration}}:}{\hfill $\Box$ 
	}{}
	}}}
\makeatother

 
%%%%%%%%%%%%%%%%%%%%%%%%%%%%%%%%%%%%%%%%%%%%%%%%%
%
%
%                      Macros dont pas fier du tout
%
%
%%%%%%%%%%%%%%%%%%%%%%%%%%%%%%%%%%%%%%%%%%%%%%%%%



\newcommand\equations[1]{
\begin{array}{lll}
#1
\end{array}
}



\newcommand\systeme[1]{
\left\{
\begin{array}{lll} 
#1
\end{array}
\right.
}



%------------------
%--- shorthands (AP)---
%------------------
\newcommand{\mtt}[1]{\mathtt{#1}}
\newcommand{\ovl}[1]{\overline{#1}}
\newcommand{\panic}[1]{\textcolor{red}{#1}}
%\NewEnviron{epanic}{{\color{red}\BODY}}
\newcommand{\idest}{\emph{i.e.\thinspace}}


%------------------
%--- mes trucs  ---
%------------------


\newcommand\implique{\Rightarrow}
\newcommand\ssi{\Leftrightarrow}
\newcommand\ac[1]{\{#1\}}
\newcommand\zero{\vectorl{0}}
\newcommand\un{\vectorl{1}}

\newcommand\technique[1]{{\scriptsize #1}}


%% Preuve and Co

\newcommand\sensdirect{ Direction $\Rightarrow$. }
\newcommand\sensindirect{ Direction $\Leftarrow$. }




\newcommand\donnee{\mbox{<donnée>}}



% pour aller vite

\newcommand\gM{\mathfrak{M}}
\newcommand\gMM{(M,f_1,\cdots,f_u,r_1,\cdots,r_v)}
%
\newcommand\mygM{\mK}
\newcommand\mygMM{\left ( \K, \{op_i\}_{i\in I}, r_1 \ldots,
r_\ell, \zero,\un \right )}
\newcommand\myM{\K} 

\newcommand\gMMM{(M,f_1,\cdots,f_u,c_1,\cdots,c_k,r_1,\cdots,r_v)}
\newcommand\gMMMM{(f_1,\cdots,f_u,r_1,\cdots,r_v)}
\newcommand\gMME{(M,f_1,\cdots,f_u,f'_1,\cdots,f'_\ell,r_1,\cdots,r_v)}
\newcommand\gMMS{(f_1,\cdots,f_u,f'_1,\cdots,f'_\ell,r_1,\cdots,r_v)}

\newcommand\bool{\mathfrak{B}}
\newcommand\gbool{(\{0,1\},\zero,\un,=)}

\newcommand\osg{\overline{sg}}
\newcommand\moins{\mathrel{\dot{-}}}


%%% TRUC DE BARWISE


\newcommand\UM{\kU_{\kM}}
\newcommand\VV{\mathbb{V}}
% Hereditarily finite
\newcommand\IHF{\mathbb{I}HF}
\newcommand\IHP{\IHYP}
\newcommand\IHYP{\mathbb{I}HYP}

\providecommand\citep[1]{\cite{#1}}

%% Cofinali


% probas
% défini dans amsmath
\renewcommand\Pr{\cp{Pr}}
\newcommand\Var{\cp{Var}}


% -- mesures
%\newcommand\card{card}
\newcommand\taille{\LANGUE{size}{taille}}
\newcommand\entree{e}
\newcommand\profondeur{\LANGUE{depth}{profondeur}}


%%%%%%%%%%%%%%%%%%%%%%%%%%%%%%%%%%%%%%%%%%%%%%%%%%%%%%%%%%%%%%%%%%%%%%%%%%%%
%                 MACHINES DE TURING (graphique)
%%%%%%%%%%%%%%%%%%%%%%%%%%%%%%%%%%%%%%%%%%%%%%%%%%%%%%%%%%%%%%%%%%%%%%%%%%%%



% Configuration de machine de Turing
\newcommand\configTuring[3]{#1\textbf{#2}#3}


%%% DESSIN CROIX PORU AIDER  A DESSINER

\newcommand\croix{  +(-1,-1) -- +(1,1) +(-1,1) --
  +(1,-1) }

%% DESSIN ACCOLADE

\newcommand\ACCOLADE[3]{
%#1: coordonnées de départ
%#2: coordonnées arrivée
%#3: texte
\draw[decorate,decoration={brace}]  (#1) -- (#2)
 node[midway,above=0.1cm] {#3};
}

%%%


\def\lcase{3.5 ex}
\def\hcase{3.5ex}
\tikzstyle{case} = []
%[draw,minimum width = \lcase,minimum height = 3.6ex,baseline]
\newcommand\dcase{ +(-0.5*\lcase,-0.5*\hcase) rectangle +(0.5*\lcase,0.5*\hcase)}
%\newcommand\lettre[1]{\node [name=l,case,on chain=1] {#1};}


\newcommand\RUBANMACHINEDETURING[6]{
%#1: texte
%#2: position de la tête
%#3: nb de blancs à gauche
%#4: nb de cases totales
%#5: texte au dessus du ruban
%#6: position du coin
    \draw (#6) +(\lcase,4ex ) node {#5};
    \draw (#6) +(-0.6,0)  node  {\ldots};  
    \draw (#6) node[coordinate] (t) {};
    \foreach \x in {1,2,...,#3} {
         \draw (t) node [case]   {$\blanc$} \dcase;
         \draw (t) ++(\lcase,0) node [coordinate] (t) {};
       };
   \StrLen{#1}[\Resultat] % ATTENTION IL FAUT CETTE SYNTAXE POUR FAIRE
                         % EQUIVALENT D UN EDEF
  
\foreach \x in {1,...,\Resultat} 
             { 
        \draw (t) node [case]   {\StrChar{#1}{\x}} \dcase;
        \draw (t) ++(\lcase,0) node [coordinate] (t) {};
             }
\pgfmathtruncatemacro{\z}{#4 - \Resultat - #3}

  \foreach \x in {1,2,..., \z} {
        \draw (t) node [case]   {$\blanc$} \dcase;
        \draw (t) ++(\lcase,0) node [coordinate] (t) {};
}
%% Convention: graphiquement, la case numéro n se trouve à (n*\lcase,0)
%% en comptant les case à partir de $0$

   \draw  (t) +(0.1,0) node [name=r] {\ldots}; 
}

\newcommand\RUBANMACHINEDETURINGtexte[6]{
%#1: texte
%#2: position de la tête
%#3: nb de blancs à gauche
%#4: nb de cases totales (équivalent)
%#5: texte au dessus du ruban
%#6: position du coin
    \draw (#6) +(\lcase,4ex ) node {#5};
    \draw (#6) +(-0.6,0)  node  {\ldots};  
    \draw (#6) node[coordinate] (t) {};
   \foreach \x in {1,2,...,#3} {
        \draw (t) node [case]   {$\blanc$} \dcase;
        \draw (t) ++(\lcase,0) node [coordinate] (t) {};
      };
   \draw (#6) ++(-2*\lcase,0.5*\lcase)-- +(\lcase*#4,0);
   \draw (#6) ++(-2*\lcase,-0.5*\lcase) -- +(\lcase*#4,0);
   
    \draw (t) ++(-0.5*\lcase,0) node [case,anchor=west]  (te) {#1};
    \draw (te.east)  node [anchor=west,coordinate] (t) {};

  \draw (t) ++(0.2*\lcase,0) node [coordinate] (t) {};
   \foreach \x in {1,2,...,#3} {
        \draw (t) node [case]   {$\blanc$} \dcase;
        \draw (t) ++(\lcase,0) node [coordinate] (t) {};
      };

    \draw  (t) +(0.1,0) node [anchor=west] {\ldots}; 
}


\newcommand\RUBANMACHINEDETURINGTETE[6]{
%
% dessine la tete
%
% output: (tete)
%
\RUBANMACHINEDETURING{#1}{#2}{#3}{#4}{#5}{#6}
% TETE
     \draw (#6) + ( #2 * \lcase + #3  *\lcase,- 0.5*\lcase )
     node[coordinate] (k) {};
     \draw[->] (k) +(0,-0.5) -- (k);
     \draw (k) +(0,-0.5) node[coordinate] (tete) {};
}


\newcommand\RUBANMACHINEDETURINGTETEtexte[6]{
%
% dessine la tete
%
% output: (tete)
%
\RUBANMACHINEDETURINGtexte{#1}{#2}{#3}{#4}{#5}{#6}
% TETE
     \draw (#6) + ( #2 * \lcase + #3  *\lcase,- 0.5*\lcase )
     node[coordinate] (k) {};
     \draw[->] (k) +(0,-0.5) -- (k);
     \draw (k) +(0,-0.5) node[coordinate] (tete) {};
}


\newcommand\MACHINEDETURING[6]{
%#1: texte
%#2: position de la tête
%#3: état
%#4: nb de blancs à gauche
%#5: nb de cases totales
%#6: texte en haut du ruban
%\begin{tikzpicture}[node distance=-0.15mm]

\RUBANMACHINEDETURINGTETE{#1}{#2}{#4}{#5}{#6}{0,0}
%
%% ETAT:
     \draw (tete) node [name=etat,draw,circle,anchor=north]  {#3};
     \draw (etat) -- (tete);
%\end{tikzpicture}
}

\newcommand\MACHINEDETURINGtexte[6]{
%#1: texte
%#2: position de la tête
%#3: état
%#4: nb de blancs à gauche
%#5: nb de cases totales
%#6: texte en haut du ruban
%\begin{tikzpicture}[node distance=-0.15mm]

\RUBANMACHINEDETURINGTETEtexte{#1}{#2}{#4}{#5}{#6}{0,0}
%
%% ETAT:
     \draw (tete) node [name=etat,draw,circle,anchor=north]  {#3};
     \draw (etat) -- (tete);
%\end{tikzpicture}
}

\newcommand\MACHINEDETURINGTROISRUBANS[9]{
%#1: texte1
%#2: position de la tête 1
%#3: texte2
%#4: position de la tête 2
%#5: texte3
%#6: position de la tête 3
%#7: état
%#8: nb de blancs à gauche 1
%#9: nb de cases totales
%\begin{tikzpicture}[node distance=-0.15mm]

\RUBANMACHINEDETURINGTETE{#1}{#2}{#8}{#9}{\LANGUE{tape}{ruban} $1$}{0,0}
\draw (tete) node[coordinate] (tete1) {};
\draw (tete1) +(8,0) node[coordinate] (tete1b) {};
\draw (tete1) -| (tete1b) ;

\RUBANMACHINEDETURINGTETE{#3}{#4}{#8}{#9}{\LANGUE{tape}{ruban} $2$}{0,-2}
\draw (tete) node[coordinate] (tete2) {};
\draw (tete2) +(-8.5,0) node[coordinate] (tete2b) {};
\draw (tete2) -| (tete2b) ;

\RUBANMACHINEDETURINGTETE{#5}{#6}{#8}{#9}{\LANGUE{tape}{ruban} $3$}{0,-4}
\draw (tete) node[coordinate] (tete3) {};

%% ETAT:
     \draw (tete3) node [name=etat,draw,circle,anchor=north]  {#7};
     \draw (etat) -- (tete3);
    \draw  (etat) -|  (tete2b);
\draw  (etat) -|  (tete1b);
%\end{tikzpicture}
}

\newcommand\MACHINEDETURINGTROISRUBANSENGLISH[9]{
%#1: texte1
%#2: position de la tête 1
%#3: texte2
%#4: position de la tête 2
%#5: texte3
%#6: position de la tête 3
%#7: état
%#8: nb de blancs à gauche 1
%#9: nb de cases totales
%\begin{tikzpicture}[node distance=-0.15mm]

\RUBANMACHINEDETURINGTETE{#1}{#2}{#8}{#9}{tape $1$}{0,0}
\draw (tete) node[coordinate] (tete1) {};
\draw (tete1) +(8,0) node[coordinate] (tete1b) {};
\draw (tete1) -| (tete1b) ;

\RUBANMACHINEDETURINGTETE{#3}{#4}{#8}{#9}{tape $2$}{0,-2}
\draw (tete) node[coordinate] (tete2) {};
\draw (tete2) +(-8.5,0) node[coordinate] (tete2b) {};
\draw (tete2) -| (tete2b) ;

\RUBANMACHINEDETURINGTETE{#5}{#6}{#8}{#9}{tape $3$}{0,-4}
\draw (tete) node[coordinate] (tete3) {};

%% ETAT:
     \draw (tete3) node [name=etat,draw,circle,anchor=north]  {#7};
     \draw (etat) -- (tete3);
    \draw  (etat) -|  (tete2b);
\draw  (etat) -|  (tete1b);
%\end{tikzpicture}
}



\newcommand\MACHINEDETURINGTROISRUBANStexte[9]{
%#1: texte1
%#2: position de la tête 1
%#3: texte2
%#4: position de la tête 2
%#5: texte3
%#6: position de la tête 3
%#7: état
%#8: nb de blancs à gauche 1
%#9: nb de cases totales
% \begin{tikzpicture}[node distance=-0.15mm]

\RUBANMACHINEDETURINGTETEtexte{#1}{#2}{#8}{#9}{\LANGUE{tape}{ruban} $1$}{0,0}
\draw (tete) node[coordinate] (tete1) {};
\draw (tete1) +(8,0) node[coordinate] (tete1b) {};
\draw (tete1) -| (tete1b) ;

\RUBANMACHINEDETURINGTETEtexte{#3}{#4}{#8}{#9}{\LANGUE{tape}{ruban} $2$}{0,-2}
\draw (tete) node[coordinate] (tete2) {};
\draw (tete2) +(-8.5,0) node[coordinate] (tete2b) {};
\draw (tete2) -| (tete2b) ;

\RUBANMACHINEDETURINGTETEtexte{#5}{#6}{#8}{#9}{\LANGUE{tape}{ruban} $3$}{0,-4}
\draw (tete) node[coordinate] (tete3) {};

%% ETAT:
     \draw (tete3) node [name=etat,draw,circle,anchor=north]  {#7};
     \draw (etat) -- (tete3);
    \draw  (etat) -|  (tete2b);
\draw  (etat) -|  (tete1b);
%\end{tikzpicture}
}


\newcommand\MACHINEDETURINGDEUXRUBANStexte[9]{
%#1: texte1  
%#2: position de la tête 1 
%#3: texte2 inutilisée
%#4: position de la tête 2 inutilisée
%#5: texte3
%#6: position de la tête 3
%#7: état
%#8: nb de blancs à gauche 1
%#9: nb de cases totales
% \begin{tikzpicture}[node distance=-0.15mm]

\RUBANMACHINEDETURINGTETEtexte{#1}{#2}{#8}{#9}{\LANGUE{tape}{ruban} $1$}{0,0}
\draw (tete) node[coordinate] (tete1) {};
\draw (tete1) +(8,0) node[coordinate] (tete1b) {};
\draw (tete1) -| (tete1b) ;

% \RUBANMACHINEDETURINGTETEtexte{#3}{#4}{#8}{#9}{\LANGUE{tape}{ruban} $2$}{0,-1.5}
% \draw (tete) node[coordinate] (tete2) {};
% \draw (tete2) +(-5.5,0) node[coordinate] (tete2b) {};
% \draw (tete2) -| (tete2b) ;

\RUBANMACHINEDETURINGTETEtexte{#5}{#6}{#8}{#9}{\LANGUE{tape}{ruban} $2$}{0,-1.5}
\draw (tete) node[coordinate] (tete3) {};

%% ETAT:
     \draw (tete3) node [name=etat,draw,circle,anchor=north]  {#7};
     \draw (etat) -- (tete3);
%   \draw  (etat) -|  (tete2b);
\draw  (etat) -|  (tete1b);
%\end{tikzpicture}
}


\def\argdix{2}
\def\argonze{25}
\newcommand\MACHINEDETURINGQUATRERUBANStextep[9]{
%#1: texte1
%#2: position de la tête 1
%#3: texte2
%#4: position de la tête 2
%#5: texte3
%#6: position de la tête 3
%#7: texte4
%#8: position de la tête 3
%#9: état
%#10: nb de blancs à gauche 1
%#11: nb de cases totales
% \begin{tikzpicture}[node distance=-0.15mm]

%% BUG sur #10 devenu \argdix
%% BUG SUR #11 devenu \argonze

\RUBANMACHINEDETURINGTETEtexte{#1}{#2}{\argdix}{\argonze}{\LANGUE{tape}{ruban} $1$}{0,0}
% \draw (tete) node[coordinate] (tete1) {};
% \draw (tete1) +(8,0) node[coordinate] (tete1b) {};
% \draw (tete1) -| (tete1b) ;

\RUBANMACHINEDETURINGTETEtexte{#3}{#4}{\argdix}{\argonze}{\LANGUE{tape}{ruban} $2$}{0,-1.5}
% \draw (tete) node[coordinate] (tete2) {};
% \draw (tete2) +(-5.5,0) node[coordinate] (tete2b) {};
% \draw (tete2) -| (tete2b) ;

\RUBANMACHINEDETURINGTETEtexte{#5}{#6}{\argdix}{\argonze}{\LANGUE{tape}{ruban} $3$}{0,-3.2}
\draw (tete) node[coordinate] (tete3) {};

\RUBANMACHINEDETURINGTETEtexte{#7}{#8}{\argdix}{\argonze}{\LANGUE{tape}{ruban} $4$}{0,-4.9}
\draw (tete) node[coordinate] (tete4) {};

%% ETAT:
     \draw (tete4) node [name=etat,draw,circle,anchor=north]  {#9};
%      \draw (etat) -- (tete3);
%     \draw  (etat) -|  (tete2b);
% \draw  (etat) -|  (tete1b);
%\end{tikzpicture}
}


%%%%%%%%%%%%%%%%%%%%%%%%%%%%%%%%%%%%%%%%%%%%%%%%%%%%%%%%%%%%%%%%%%%%%%%%%%%%
%                 MACHINES DE TURING (automate)
%%%%%%%%%%%%%%%%%%%%%%%%%%%%%%%%%%%%%%%%%%%%%%%%%%%%%%%%%%%%%%%%%%%%%%%%%%%%



\newcommand\dec{1ex}
\newcommand\myACCOLADE[4]{
\ACCOLADE{#1*\lcase+#4*\lcase-0.5*\lcase,0.5*\lcase+\dec}{#1*\lcase+#4*\lcase-0.5*\lcase+#3*\lcase,0.5*\lcase+\dec}{#2}
}


\newcommand\M{\Sigma}


\newcommand\sommet[2]{\node[state] (#1) #2 {$#1$};}
\newcommand\sommetinitial[2]{\node[state,initial] (#1) #2 {$#1$};}
\newcommand\sommetfinal[2]{\node[state,accepting] (#1) #2 {$#1$};}
\newcommand\gtransition[5]{\draw (#1) edge node {#2/#4 $#5$} (#3);}
\newcommand\gtransitionba[5]{\draw (#1) [loop above] edge node {#2/#4
    $#5$} (#3);}
\newcommand\gtransitionbb[5]{\draw (#1) [loop below] edge node {#2/#4
    $#5$} (#3);}
\newcommand\gtransitionbleft[5]{\draw (#1) [loop left] edge node {#2/#4
    $#5$} (#3);}
\newcommand\gtransitionbright[5]{\draw (#1) [loop right] edge node {#2/#4
    $#5$} (#3);}
\newcommand\gtransitionbl[5]{\draw (#1) [bend left] edge node {#2/#4
    $#5$} (#3);}
\newcommand\gtransitionbr[5]{\draw (#1) [bend right] edge node {#2/#4
    $#5$} (#3);}
\newcommand\gtransitionbadouble[9]{\draw (#1) [loop above] edge node {
\begin{tabular}{l}
#2/#4 $#5$ \\
  #7/#8 $#9$ \\
\end{tabular}
} (#3);}

\newcommand\SCALE{}

\newcommand\dessinmtp[2]{
\begin{center}
\begin{tikzpicture}[->,>=stealth',shorten >=1pt,auto,node distance=2.3cm,semithick,#2] 
#1
\end{tikzpicture}
\end{center}
}

\newcommand\dessinmt[1]{\dessinmtp{#1}{scale=1}}


\newcommand\ANIMATIONMT[1]{
\begin{overprint}
\begin{dessinp}{scale=0.52}
\foreach \av/\q/\b  [count=\xi]
in  
{#1}
{ 
\only<\xi| handout 0>{ 
 \StrLen{\av}[\Resintm] % ATTENTION IL FAUT CETTE SYNTAXE POUR FAIRE
                         % EQUIVALENT D UN EDEF
 \MACHINEDETURING{\av \b}{\Resintm}{$\q$}{5}{20}{}
}}
 \end{dessinp}
\end{overprint}
}

\newcommand\ANIMATIONMTd[1]{
\begin{overprint}
\begin{dessinp}{scale=0.52}
\foreach \av/\q/\b  [count=\xi]
in  
{#1}
{ 
\only<\xi| handout 0>{ 
 \StrLen{\av}[\Resintm] % ATTENTION IL FAUT CETTE SYNTAXE POUR FAIRE
                         % EQUIVALENT D UN EDEF
 \hspace{-1cm}\MACHINEDETURING{\av \b}{\Resintm}{$\q$}{5}{33}{}
}}
 \end{dessinp}
\end{overprint}
}

\def\MAXDET{20}

\newcommand\DIAGRAMMEESPACETEMPS[1]{
\draw node[coordinate] (tt) {};
\foreach \av/\q/\b  in  
{#1}
{ 

\RUBANMACHINEDETURING{\av{$\q$}\b}{0}{4}{\MAXDET}{}{tt};
\draw (tt) +(0,-\hcase) node[coordinate] (tt) {};
}
%\foreach \x in {0,...,19} {\draw (tt) + (\x*\lcase,0) node {$\vdots$};}
}



\newcommand\progun{
/q_0/0011, % initial
X/q_1/011,X0/q_1/11,X/q_2/0Y1,/q_2/X0Y1,
X/q_0/0Y1,XX/q_1/Y1,XXY/q_1/1, XX/q_2/YY, X/q_2/XYY,
XX/q_0/YY,XXY/q_3/Y,XXYY/q_3/B,XXYYB/q_4/B
}

\newcommand\DIAGRAMMEESPACETEMPSun{
\DIAGRAMMEESPACETEMPS
%% RECOPIE progun
{/q_0/0011, % initial
X/q_1/011,X0/q_1/11,X/q_2/0Y1,/q_2/X0Y1,
X/q_0/0Y1,XX/q_1/Y1,XXY/q_1/1, XX/q_2/YY, X/q_2/XYY,
XX/q_0/YY,XXY/q_3/Y,XXYY/q_3/B,XXYYB/q_4/B}
%% FIN RECOPIE
}

\newcommand\progdeux{
/q_0/0010, X/q_1/010, X0/q_1/10, X/q_2/0Y0,/q_2/X0Y0,
X/q_0/0Y0,XX/q_1/Y0,XXY/q_1/0,XXY0/q_1/%B,
}


\newcommand\DIAGRAMMEESPACETEMPSdeux{
\DIAGRAMMEESPACETEMPS
{
%% RECOPIE progdeux
/q_0/0010, X/q_1/010, X0/q_1/10, X/q_2/0Y0,/q_2/X0Y0,
X/q_0/0Y0,XX/q_1/Y0,XXY/q_1/0,XXY0/q_1/%B,
%% FIN RECOPIE
}
}

\newcommand\DIAGRAMMEESPACETEMPSdeuxvariante{
\def\memolcase{\lcase}
\def\memoMAXDET{20}
\def\MAXDET{15}
\def\lcase{6 ex}
\DIAGRAMMEESPACETEMPS
{
%% RECOPIE progdeux
/q_00/010, X/q_10/10, X0/q_11/0, X/q_20/Y0,/q_2X/0Y0,
X/q_00/Y0,XX/q_1Y/0,XXY/q_10/,XXY0/q_1B/%B,
%% FIN RECOPIE
}
\def\lcase{\memolcase}
\def\MAXDET{\memoMAXDET}
}


\newcommand\ANIMATIONun{
\ANIMATIONMT
%% RECOPIE progun
{/q_0/0011, % initial
X/q_1/011,X0/q_1/11,X/q_2/0Y1,/q_2/X0Y1,
X/q_0/0Y1,XX/q_1/Y1,XXY/q_1/1, XX/q_2/YY, X/q_2/XYY,
XX/q_0/YY,XXY/q_3/Y,XXYY/q_3/B,XXYYB/q_4/B}
%% FIN RECOPIE
}



\newcommand\DESSINALGO[3]{
% argument 1= ou
% argument 2= texte 
% argument 3= nom du noeud ex m
\draw node[draw,#1,minimum size=1cm] (#3) {#2};

% pour fitting box
\draw (#3.west) +(-0.2cm,0) node (#3c2) {};
\draw (#3.north) +(0,+0.2cm) node (#3c3) {};
\draw (#3.south) +(0,-0.2cm) node (#3c4) {};
}

\newcommand\DESSINALGOa[3]{
\DESSINALGO{#1}{#2}{#3}
\draw (#3.east) +(-0.1cm,0.2cm) node(#3a) {};
\draw[->] (#3a) -- ++(0.5cm,0) node (#3aa) {};
\draw (#3aa) node[anchor=west] (#3aaa) {\LANGUE{Accept}{Accepte}};
\draw (#3aaa.east) node(#3aaaa) {};
}

\newcommand\DESSINALGOe[4]{
\DESSINALGO{#1}{#2}{#3}
\draw (#3.east) +(-0.1cm,0.2cm) node(#3a) {};
\draw[->] (#3a) -- ++(0.5cm,0) node (#3aa) {};
\draw (#3aa) node[anchor=west] (#3aaa) {#4};
\draw (#3aaa.east) node(#3aaaa) {};
}


\newcommand\DESSINALGOar[3]{
\DESSINALGOa{#1}{#2}{#3}
\draw (#3.east) +(-0.1cm,-0.2cm) node(#3r) {};
\draw[->] (#3r) -- ++(0.5cm,0) node (#3rr) {};
\draw (#3rr) node[anchor=west] (#3rrr) {\LANGUE{Reject}{Refuse}};
\draw (#3rrr.east) node(#3rrrr) {};
}


\newcommand\WINDOWarg[7] {
%#1: position
%#2..7: contenu du tableau
\draw (#1) node[case]  {#2} \dcase;
\draw (#1) ++ (\lcase,0) node[case] {#3} \dcase;
\draw (#1) ++ (2*\lcase,0) node[case] {#4} \dcase;
\draw (#1) ++ (0,-\lcase) node[case] {#5} \dcase;
\draw (#1) ++ (\lcase,-\lcase) node[case] {#6} \dcase;
\draw (#1) ++ (2*\lcase,-\lcase) node[case] {#7} \dcase;
}

\newcommand\WINDOW[3]{
%#1: positioin
%#2: contenu
%#3: étiquette
\WINDOWarg{#1}
{\StrChar{#2}{1}}{\StrChar{#2}{2}}{\StrChar{#2}{3}}
{\StrChar{#2}{4}}{\StrChar{#2}{5}}{\StrChar{#2}{6}}
\draw (#1) +(-\lcase*1.2,-\lcase*0.5) node {#3};
}





%%%%%%%%%%%%%%%%%%%%%%%%%%%%%%%%%%%%%%%%%%%%%%%%%%%%%%%%%%%%%%%%%%%%%%%%%%%%
%                 LISTINGS
%%%%%%%%%%%%%%%%%%%%%%%%%%%%%%%%%%%%%%%%%%%%%%%%%%%%%%%%%%%%%%%%%%%%%%%%%%%%


% LISTINGS
\RequirePackage{listings}
 \lstset{language=Java,
 mathescape=true,
 showspaces=false,
 showstringspaces=false,
 frame=single,
morekeywords={recursive,then,div,function,integer,read,out,else,par,endpar,seq,endseq,and,not,read,write,repeat,until,undef,let,to,endfor,endif,endwhile,guess,in,parallel},
% basicstyle=\ttfamily\small,
% basewidth={1.1ex},
% keywordstyle=\bf\color{blue},
% %stringstyle=\bf\ttfamily\color{green},
% commentstyle=\bfseries\color{magenta},
 literate={:=}{{$\gets$ }}1 {<=}{{$\leq$}}2  {>=}{{$\geq$}}2
 {!=}{{$\neq$}}1 
 }
% \lstset{escapechar=\%}
\let\lst\lstinline
\def\beginlisting{\begin{lstlisting}}
\def\endlisting{\end{lstlisting}}  




\newenvironment{algo}[1]{
  \lstset{escapechar=\@}
  \def\input{$#1$}
  \def\out{out}
  \def\BEGIN{\textbf{read} \input }
 \def\BEGINN{\textbf{repeat} }
  \def\END{\textbf{until out!=undef}}
  \def\ENDD{\textbf{write \out}}
  \def\ENDT{\textbf{until $\ctl$ = $q_a$ or $\ctl$ = $q_r$}}
  \def\ENDR{\textbf{until $\ctl$ = $q_a$}}
  \def\ENDP{\textbf{until $\ctl$ = $q_a$}}
}{}
\newenvironment{algon}[1]{
  \lstset{numbers=left}
  \begin{algo}{#1}
}{\end{algo}}
\newcommand\FSM[3]{FSM(#1,\lst{#2},#3)}
\newcommand\FSMi[3]{FSM(#1,{#2},#3)}




%%%%%%%%%%%%%%%%%%%%%%%%%%%%%%%%%%%%%%%%%%%%%%%%%
%
%
%                      TRADUCTION EN COURS
%
%%%%%%%%%%%%%%%%%%%%%%%%%%%%%%%%%%%%%%%%%%%%%%%%%

\newcommand\ENTREPOT[1]{\NDLR{ENTREPOT ICI: #1}}
\newcommand\scite[2][]{\cite[#1]{#2}}
\newcommand\nospellcheck[1]{{#1}}
\newcommand\nd{nd}
\newcommand\NL{
 
}

\newenvironment{FRANCAIS}{
\small
\color{orange}
}
{\normalsize
\color{black}
}


\newenvironment{SEMIFRANCAIS}{
\small
\color{violet}
}
{\normalsize
\color{black}
}

\newcommand\QFRANCAIS[1]{
\begin{FRANCAIS}
#1
\end{FRANCAIS}
}


%%%%%%%%%%%%%%%%%%%%%%%%%%%%%%%%%%%%%%%%%%%%%%%%%
%
%
%                      Wikipedia
%
%%%%%%%%%%%%%%%%%%%%%%%%%%%%%%%%%%%%%%%%%%%%%%%%%

\providecommand\tightlist{}
\providecommand\Complex{\C}
\newenvironment{myalgie}{\begin{array}{lll}}{\end{array}}

\newcommand\nl{\ \\ }




%%%%%%%%%%%%%%%%%%%%%%%%%%%%%%%%%%%%%%%%%%%%%%%%%
%
%
%                      Pour citations précises
%
%%%%%%%%%%%%%%%%%%%%%%%%%%%%%%%%%%%%%%%%%%%%%%%%%



\newcommand\citeprecis[2]{{\cite[#1]{#2}}}
\newcommand\citeprecisinv[2]{\citeprecis{#2}{#1}}

\usepackage{csquotes}



%%%%%%%%%%%%%%%%%%%%%%%%%%%%%%%%%%%%%%%%%%%%%%%%%
%
%
%                      Travail sur couleurs
%
%   toutes les réponses pointent dans le même sens, et le schéma correspond exactement 
%. au profil perceptif d’un protanope** (déficit des cônes L / longueurs d’onde longues)
%  MAIS SEVERE
%
%%%%%%%%%%%%%%%%%%%%%%%%%%%%%%%%%%%%%%%%%%%%%%%%%


% ================================
% GROUPE 1 — Verts sombres
% ================================
\definecolor{deepgreen}{RGB}{0,160,40}
\definecolor{forestgreen}{RGB}{0,150,70}
\definecolor{mossgreen}{RGB}{0,140,90}
\definecolor{greenwave}{RGB}{0,150,90}
\definecolor{forestmist}{RGB}{0,160,90}
\definecolor{darkteal}{RGB}{0,130,70}

% ================================
% GROUPE 2 — Teal / verts bleutés
% ================================
\definecolor{seagreen}{RGB}{0,130,100}
\definecolor{freshteal}{RGB}{0,140,130}
\definecolor{greenteal}{RGB}{0,140,120}
\definecolor{blueteal}{RGB}{0,150,120}      % ancien "teal" (doublon corrigé)
\definecolor{lagoon}{RGB}{0,180,120}
\definecolor{greenmist}{RGB}{0,170,110}

% ================================
% GROUPE 3 — Turquoises / Aqua
% ================================
\definecolor{springgreen}{RGB}{0,190,100}
\definecolor{mintgreen}{RGB}{10,200,150}
\definecolor{oceanfoam}{RGB}{0,200,150}
\definecolor{aquagreen}{RGB}{30,210,170}
\definecolor{coastalaqua}{RGB}{40,210,180}
\definecolor{softaqua}{RGB}{40,210,190}

% ================================
% GROUPE 4 — Cyan / Aqua clairs
% ================================
\definecolor{emeraldteal}{RGB}{0,180,140}
\definecolor{deepcyan}{RGB}{0,180,160}
\definecolor{lightcyan}{RGB}{0,200,160}
\definecolor{aquamarine}{RGB}{0,200,160}
\definecolor{lightaqua}{RGB}{60,230,200}
\definecolor{mistblue}{RGB}{80,240,210}

% ================================
% GROUPE 5 — Bleus clairs
% ================================
\definecolor{paleaqua}{RGB}{80,200,190}
\definecolor{glacier}{RGB}{90,150,200}
\definecolor{skyblueA}{RGB}{60,150,220}      % ancien "skyblue"
\definecolor{lightsky}{RGB}{90,170,230}
\definecolor{iceblue}{RGB}{120,190,240}
\definecolor{frost}{RGB}{150,210,250}

% ================================
% GROUPE 6 — Bleus moyens
% ================================
\definecolor{steelblue}{RGB}{30,100,160}
\definecolor{cadetblue}{RGB}{70,120,180}
\definecolor{coolblue}{RGB}{110,170,220}
\definecolor{morningblue}{RGB}{140,200,240}
\definecolor{blueiceA}{RGB}{40,150,210}       % ancien "blueice" (doublon corrigé)
\definecolor{blueair}{RGB}{60,170,225}

% ================================
% GROUPE 7 — Azurs / Bleu saturé
% ================================
\definecolor{tealblue}{RGB}{0,110,130}
\definecolor{deepbluecyan}{RGB}{0,110,170}
\definecolor{azur}{RGB}{0,70,180}
\definecolor{brightazure}{RGB}{20,130,200}
\definecolor{lightazure}{RGB}{80,180,235}
\definecolor{skylight}{RGB}{100,200,245}

% ================================
% GROUPE 8 — Bleus sombres / violets
% ================================
\definecolor{nightblue}{RGB}{40,0,110}
\definecolor{slateblue}{RGB}{20,40,120}
\definecolor{violetA}{RGB}{60,50,160}         % ancien "violet" (distinct de violetB ci-dessous)
\definecolor{blueviolet}{RGB}{80,60,180}
\definecolor{purpleA}{RGB}{100,70,200}        % ancien "purple"
\definecolor{lightpurple}{RGB}{130,90,210}

% ================================
% GROUPE 9 — Violets saturés / indigos
% ================================
\definecolor{lilac}{RGB}{160,110,220}
\definecolor{indigoA}{RGB}{50,0,130}          % ancien "indigo"
\definecolor{deepindigo}{RGB}{70,10,150}
\definecolor{royalpurple}{RGB}{90,20,170}

% ================================
% GROUPE 10 — Fermeture cercle perceptif
% ================================
\definecolor{coolmint}{RGB}{0,130,70}
\definecolor{deepsea}{RGB}{0,130,70}

% ================================
% PALETTE DEUTÉRANOPE — noms uniques
% ================================
\definecolor{crimson}{RGB}{200,30,60}
\definecolor{brickred}{RGB}{180,40,40}
\definecolor{coral}{RGB}{220,90,70}
\definecolor{salmon}{RGB}{240,130,120}

\definecolor{magenta}{RGB}{200,40,150}
\definecolor{fuchsia}{RGB}{220,60,170}
\definecolor{violetB}{RGB}{140,70,190}        % ancien "violet"
\definecolor{indigoB}{RGB}{90,50,180}         % ancien "indigo"

\definecolor{royalblueA}{RGB}{50,80,200}      % ancien "royalblue"
\definecolor{azure}{RGB}{30,120,230}
\definecolor{skyblueB}{RGB}{80,160,230}       % second "skyblue"
\definecolor{tealB}{RGB}{0,150,170}           % second "teal"

\definecolor{marine}{RGB}{0,80,140}




%% Alternative

% === Générateur automatique fill/text pour protanopes ===
% Usage : \ProtaColorPair{mistblue}
% Suppose que mistblue est déjà défini par \definecolor{mistblue}{RGB}{...}

\newcommand{\ProtaColorPair}[1]{%
  % fill = couleur originale
  \colorlet{#1fill}{#1}%
  % text = version assombrie idéalement lisible
  \colorlet{#1text}{#1!80!black}%
}

%Utilisation
%\ProtaColorPair{mistblue}

%\textcolor{mistbluefill}{Couleur claire} \quad
%\textcolor{mistbluetext}{Couleur pour texte}

	\foreach \i/\name in {
	  1/deepgreen,
	  2/greenwave,
	  3/lagoon,
	  4/oceanfoam,
	  5/softaqua,
	  6/mistblue,
 	 7/lightsky,
	  8/blueiceA,   % <-- corrigé
	  9/deepcyan,
	  10/azur,
	  11/blueviolet,
	  12/royalpurple}
	  {\ProtaColorPair{\name}}
	  


\newcommand\ROUEPROTANOPEDOUZE{
	% Roue chromatique protanope — 12 couleurs
	\begin{tikzpicture}[scale=3]

	\foreach \i/\name/\nametext in {
	  1/deepgreen/deepgreen,
	  2/greenwave/greenwave,
	  3/lagoon/lagoontext,
	  4/oceanfoam/oceanfoam,
	  5/softaqua/softaquatext,
	  6/mistblue/mistbluetext,
 	 7/lightsky/lightskytext,
	  8/blueiceA/blueiceA,   % <-- corrigé
	  9/deepcyan/blueiceA,
	  10/azur/azur,
	  11/blueviolet/blueviolet,
	  12/royalpurple/royalpurple}
	{
	  % secteur coloré
 	 \filldraw[fill=\name, draw=black, line width=0.2pt]
	    ({90-30*\i}:1)
	    -- ({90-30*(\i-1)}:1)
	    -- (0,0) -- cycle;

 	 % label dans la bonne couleur
	  \node at ({90-30*(\i-0.5)}:1.25)
	    {\textcolor{\nametext}{\small \nametext}};
	}

	\end{tikzpicture}
}


\newcommand\ROUEPROTANOPEVINGT{
\begin{tikzpicture}[scale=3]

\foreach \i/\name/\nametext in {
  1/deepgreen/deepgreen,
  2/greenwave/greenwave,
  3/lagoon/lagoon,
  4/oceanfoam/oceanfoam,
  5/softaqua/softaquatext,
  6/mistblue/mistbluetext,
  7/lightsky/lightsky,
  8/blueiceA/blueiceA,      % corrigé
  9/deepcyan/deepcyan,
  10/azur/azur,
  11/brightazure/brightazure,
  12/skyblueA/skyblueA,     % corrigé
  13/coolblue/coolblue,
  14/glacier/glacier,
  15/marine/marine,
  16/violetA/violetA,      % corrigé
  17/blueviolet/blueviolet,
  18/purpleA/purpleA,      % corrigé
  19/royalpurple/royalpurple,
  20/indigoA/indigoA}      % corrigé
{
  % Secteur coloré
  \filldraw[fill=\name, draw=black, line width=0.2pt]
    ({90-18*\i}:1)
    -- ({90-18*(\i-1)}:1)
    -- (0,0) -- cycle;

  % Étiquette
  \node at ({90-18*(\i-0.5)}:1.25)
    {\textcolor{\nametext}{\small \nametext}};
}

\end{tikzpicture}
}

\newcommand\ROUEDEUTERANOPEDOUZE{
% Roue chromatique deutéranope — 12 couleurs
\begin{tikzpicture}[scale=3]

\foreach \i/\name in {
  1/crimson,
  2/brickred,
  3/coral,
  4/salmon,
  5/magenta,
  6/fuchsia,
  7/violetB,       % corrigé
  8/indigoB,       % corrigé
  9/royalblueA,    % corrigé
  10/azure,
  11/skyblueB,     % corrigé
  12/tealB}        % corrigé
{
  % Secteur coloré
  \filldraw[fill=\name, draw=black, line width=0.2pt]
    ({90-30*\i}:1)
    -- ({90-30*(\i-1)}:1)
    -- (0,0) -- cycle;

  % Label dans la couleur
  \node at ({90-30*(\i-0.5)}:1.25)
    {\textcolor{\name}{\small \name}};
}

\end{tikzpicture}
}



% === Macro intelligente pour foncer une couleur de façon lisible pour protanope ===
% Usage : \protadarken{mistblue}{40}  -> crée mistblue@dark
% Le 40 signifie : 40% de mélange avec du noir

\newcommand{\DarkenForText}[2]{%
  % #1 = nom de la couleur
  % #2 = pourcentage (40 conseillé si tu ne sais pas)
  %
  % On crée une nouvelle couleur #1@dark
  \colorlet{#1text}{black!#2!#1}%
}





%%%%

% Groupe 1
\DarkenForText{deepgreen}{20}
\DarkenForText{forestgreen}{20}
\DarkenForText{mossgreen}{20}
\DarkenForText{greenwave}{20}
\DarkenForText{forestmist}{20}
\DarkenForText{darkteal}{20}

% Groupe 2
\DarkenForText{seagreen}{20}
\DarkenForText{freshteal}{20}
\DarkenForText{greenteal}{20}
\DarkenForText{blueteal}{20}
\DarkenForText{lagoon}{20}
\DarkenForText{greenmist}{20}

% Groupe 3
\DarkenForText{springgreen}{20}
\DarkenForText{mintgreen}{20}
\DarkenForText{oceanfoam}{20}
\DarkenForText{aquagreen}{20}
\DarkenForText{coastalaqua}{20}
\DarkenForText{softaqua}{20}

% Groupe 4
\DarkenForText{emeraldteal}{20}
\DarkenForText{deepcyan}{20}
\DarkenForText{lightcyan}{20}
\DarkenForText{aquamarine}{20}
\DarkenForText{lightaqua}{20}
\DarkenForText{mistblue}{30}

% Groupe 5
\DarkenForText{paleaqua}{20}
\DarkenForText{glacier}{20}
\DarkenForText{skyblueA}{20}
\DarkenForText{lightsky}{20}
\DarkenForText{iceblue}{20}
\DarkenForText{frost}{20}

% Groupe 6
\DarkenForText{steelblue}{20}
\DarkenForText{cadetblue}{20}
\DarkenForText{coolblue}{20}
\DarkenForText{morningblue}{20}
\DarkenForText{blueiceA}{20}
\DarkenForText{blueair}{20}

% Groupe 7
\DarkenForText{tealblue}{20}
\DarkenForText{deepbluecyan}{20}
\DarkenForText{azur}{20}
\DarkenForText{brightazure}{20}
\DarkenForText{lightazure}{20}
\DarkenForText{skylight}{20}

% Groupe 8
\DarkenForText{nightblue}{20}
\DarkenForText{slateblue}{20}
\DarkenForText{violetA}{20}
\DarkenForText{blueviolet}{20}
\DarkenForText{purpleA}{20}
\DarkenForText{lightpurple}{20}

% Groupe 9
\DarkenForText{lilac}{20}
\DarkenForText{indigoA}{20}
\DarkenForText{deepindigo}{20}
\DarkenForText{royalpurple}{20}

% Groupe 10
\DarkenForText{coolmint}{20}
\DarkenForText{deepsea}{20}

% Palette deutéranope
\DarkenForText{crimson}{20}
\DarkenForText{brickred}{20}
\DarkenForText{coral}{20}
\DarkenForText{salmon}{20}

\DarkenForText{magenta}{20}
\DarkenForText{fuchsia}{20}
\DarkenForText{violetB}{20}
\DarkenForText{indigoB}{20}

\DarkenForText{royalblueA}{20}
\DarkenForText{azure}{20}
\DarkenForText{skyblueB}{20}
\DarkenForText{tealB}{20}

% Couleur supplémentaire
\DarkenForText{marine}{20}

%%%%%%%%%%%%%%%%%%%%%
%
%.   apx proof
%
%%%%%%%%%%%%%%%%%%%%%



% do
\newcommand\PREUVESENAPPENDIXCOMPATIBLE{
	\newenvironment{appendixproof}{{\bf \LANGUE{Proof}{Démonstration}}:}{\hfill $\Box$}
	\newtheorem{theoremrep}[theorem]{Theorem}
	\newtheorem{lemmarep}[theorem]{Lemma}
	\newtheorem{propositionrep}[theorem]{Proposition}
	\newtheorem{corollaryrep}[theorem]{Corollary}
	\newtheorem{summaryrep}[theorem]{Summary}	
}

\providecommand\apxparameter{}

%% proof inline
%\providecommand\apxparameter{appendix=inline}
%% proof in appendix
%\providecommand\apxparameter{}


\apxproofmode{
\usepackage[\apxparameter]{apxproof}
\theoremstyle{plain}

\newtheoremrep{theorem}{Theorem}
%\newtheoremrep{theoremconjecture}[theorem]{Theorem-To-Be-Proved=Conjecture}
%\newtheoremrep{lemmaconjecture}[theorem]{Conjecture-To-Be-Proved=Conjecture}
\newtheoremrep{claim}[theorem]{Theorem}
\newtheoremrep{proposition}[theorem]{Proposition}
\newtheoremrep{lemma}[theorem]{Lemma}
\newtheoremrep{corollary}[theorem]{Corollary}
%%bug mars 26: \newtheoremrep{summary}[theorem]{Summary}
% Du coup:
\newtheoremrep{resume}[theorem]{Summary}
%bug mars 26: \newtheoremrep{openproblem}[theorem]{Open Problem}
%\newtheoremrep{claim}[theorem]{Claim}
}


%\apxproofmodesubsection{
%%apx proof, pour que compteur correct
%\makeatletter
%\AtBeginDocument{%
%\let\axp@section\subsection
%  \let\axp@oldsection\subsection
%}
%\makeatother
%\renewcommand\appendixprelim{\section{Missing proofs}}
%}


\makeatletter
\apxproofmodesubsection{
\let\axp@section\subsection
%\let\axp@oldsection\subsection
\renewcommand\appendixprelim{\section{Proofs from main documents}}
}
\makeatother


\makeatletter
\newcommand\appendixapxproofsubsection{
\appendix
\let\apx@orig@appendix\appendix
\renewcommand{\appendix}{}
}
\makeatother


%%% Pour checker/vérifier que bon


\newcommand\docheck[1]{\textcolor{check}{#1}}

\usepackage{xcolor}
%\definecolor{chose}{RGB}{0,110,50}
%\definecolor{forestgreen}{RGB}{0,150,70}
%\definecolor{ctruc}{RGB}{164,164,164}
%\definecolor{oceanfoam}{RGB}{0,200,150}
\definecolor{check}{RGB}{160,160,0}




 via \providecommand,
%% ou après via \renewcommand) :
\providecommand\BIBFILES{@@reference-biblio}
\providecommand\FIGCOMMONS{/Users/bournez/lib/LaTeX/Perso/FIG-COMMONS}

\newcommand\affiliationofficielleolivier{LIX, CNRS, École polytechnique, Institut Polytechnique de Paris, Palaiseau, France}
\newcommand\affiliationofficielleolivierslides{LIX, CNRS, École polytechnique, \\ Institut Polytechnique de Paris, Palaiseau, France}


%%%%%%%%%%%%%%%%%%%%%%%%
%
%              begin. Switch modes compatibilités
%
%%%%%%%%%%%%%%%%%%%%%%%

%% Apxproof mode
\providecommand\apxproofmode[1]{}
	% sinon: \newcommand\apxproof[1]{#1}
	
\providecommand\apxproofmodesubsection[1]{}
	% sinon: \newcommand\apxproofmodesubsection[1]{#1}
% utiliser: \appendixapxproofsubsection au lieu de \appendix dans ce cas. 
		

%% Lipics mode
\providecommand\lipicsmode{}
	% sinon: \newcommand\lipicsmode[1]{#1}
		% en vrai sert qu'à définir \nolipics
		
%%%%%%%%%%%%%%%%%%%%%%%%
%
%              end. Switch modes compatibilités
%
%%%%%%%%%%%%%%%%%%%%%%%


%% begin gestion des switchs

% en fait \nopilcs est la seule commande utilisée

\newcommand\nolipics[1]{#1}
\lipicsmode{\renewcommand\nolipics[1]{}}

%% Décorations (fourier + cadres colorés mdframed autour de definition,
%% theorem, exercice, example, remark, ...) : actives par défaut hors slides.
%% Pour les désactiver dans un document : \newcommand\NODECORATION[1]{}
%% avant 
\providecommand\nolncs[1]{#1}	

%% Chemins par défaut (surchargables avant 
\providecommand\nolncs[1]{#1}	

%% Chemins par défaut (surchargables avant \input{macros} via \providecommand,
%% ou après via \renewcommand) :
\providecommand\BIBFILES{@@reference-biblio}
\providecommand\FIGCOMMONS{/Users/bournez/lib/LaTeX/Perso/FIG-COMMONS}

\newcommand\affiliationofficielleolivier{LIX, CNRS, École polytechnique, Institut Polytechnique de Paris, Palaiseau, France}
\newcommand\affiliationofficielleolivierslides{LIX, CNRS, École polytechnique, \\ Institut Polytechnique de Paris, Palaiseau, France}


%%%%%%%%%%%%%%%%%%%%%%%%
%
%              begin. Switch modes compatibilités
%
%%%%%%%%%%%%%%%%%%%%%%%

%% Apxproof mode
\providecommand\apxproofmode[1]{}
	% sinon: \newcommand\apxproof[1]{#1}
	
\providecommand\apxproofmodesubsection[1]{}
	% sinon: \newcommand\apxproofmodesubsection[1]{#1}
% utiliser: \appendixapxproofsubsection au lieu de \appendix dans ce cas. 
		

%% Lipics mode
\providecommand\lipicsmode{}
	% sinon: \newcommand\lipicsmode[1]{#1}
		% en vrai sert qu'à définir \nolipics
		
%%%%%%%%%%%%%%%%%%%%%%%%
%
%              end. Switch modes compatibilités
%
%%%%%%%%%%%%%%%%%%%%%%%


%% begin gestion des switchs

% en fait \nopilcs est la seule commande utilisée

\newcommand\nolipics[1]{#1}
\lipicsmode{\renewcommand\nolipics[1]{}}

%% Décorations (fourier + cadres colorés mdframed autour de definition,
%% theorem, exercice, example, remark, ...) : actives par défaut hors slides.
%% Pour les désactiver dans un document : \newcommand\NODECORATION[1]{}
%% avant \input{macros}. Retirées automatiquement en mode lipics.
\providecommand\NODECORATION[1]{\ifnotSLIDE{#1}}
\lipicsmode{\renewcommand\NODECORATION[1]{}}

\providecommand\PASSIAPXPROOF[1]{#1}
\apxproofmode{\renewcommand\PASSIAPXPROOF[1]{}}



%% end gestion des switchs

\def\MACROSPOINTEXOUVERT{}


\providecommand\ifnotTDEXAM[1]{#1}
\providecommand\ifnotSLIDE[1]{#1}


  
% pour preuves en appendix
\providecommand\PREUVESENAPPENDIX[1]{}



%%%% STYLES: APRES VOLS


\NODECORATION{
\usepackage{fourier}
}

%% Repli si fourier (et donc fourier-orns, qui définit \danger) n'est pas
%% chargé ; doit rester APRÈS le bloc ci-dessus, sinon fourier-orns échoue
%% avec « \danger already defined ».
\providecommand\danger{}

%%%% Box around environment
\usepackage[framemethod=tikz]{mdframed}

\newcommand\ENVCOLOR[1]{#1}
\newcommand\BEGINTEMPENVCOLOR[1]{\let\envcolorwas\ENVCOLOR\renewcommand\ENVCOLOR[1]{#1}}
\newcommand\ENDTEMPENVCOLOR{\let\ENVCOLOR\envcolorwas}


\NODECORATION{
\surroundwithmdframed[hidealllines=true,backgroundcolor=\ENVCOLOR{blue!10},roundcorner=8pt]{exercice}
\surroundwithmdframed[hidealllines=true,backgroundcolor=\ENVCOLOR{blue!10},roundcorner=8pt]{exercise}
\surroundwithmdframed[hidealllines=true,backgroundcolor=\ENVCOLOR{blue!10},roundcorner=8pt]{exercicee}
\surroundwithmdframed[hidealllines=true,backgroundcolor=\ENVCOLOR{blue!10},roundcorner=8pt]{exerciced}
\surroundwithmdframed[hidealllines=true,backgroundcolor=\ENVCOLOR{blue!10},roundcorner=8pt]{exercicec}
\surroundwithmdframed[hidealllines=true,backgroundcolor=\ENVCOLOR{blue!10},roundcorner=8pt]{exercicedc}
\surroundwithmdframed[hidealllines=true,backgroundcolor=\ENVCOLOR{orange!20},roundcorner=8pt]{example}
\surroundwithmdframed[hidealllines=true,backgroundcolor=\ENVCOLOR{orange!20},roundcorner=8pt]{exemple}
\surroundwithmdframed[hidealllines=true,backgroundcolor=\ENVCOLOR{orange!10},roundcorner=8pt]{remark}
\ifnotSLIDE{\surroundwithmdframed[backgroundcolor=\ENVCOLOR{green!15},roundcorner=8pt]{definition}}
\surroundwithmdframed[hidealllines=true,backgroundcolor=\ENVCOLOR{gray!35}]{proposition}
\ifnotSLIDE{\surroundwithmdframed[backgroundcolor=\ENVCOLOR{gray!35},roundcorner=8pt]{theorem}}
\surroundwithmdframed[hidealllines=true,backgroundcolor=\ENVCOLOR{gray!35}]{lemma}
\surroundwithmdframed[hidealllines=true,backgroundcolor=\ENVCOLOR{gray!35}]{corollary}


\surroundwithmdframed[backgroundcolor=red!35,roundcorner=8pt]{theoremaprouver}
\surroundwithmdframed[backgroundcolor=red!35,roundcorner=8pt]{conjectureaprouver}
}

%%% tcolorbox

\usepackage{tcolorbox}
\tcbuselibrary{skins}


  %%. Commande magique
%% \tcolorboxenvironment{⟨name⟩}{⟨options⟩}: attach options de tcolorbox à environnement
%% alternative: je cree des boites moi meme: voici des exemples


\newtcolorbox{maboiteconfiance}[2][]{colback=red!5!white, colframe=red!75!black,fonttitle=\bfseries, colbacktitle=red!85!black,enhanced,
attach boxed title to top center={yshift=-2mm},
  title={#2},#1}
  
\newtcolorbox{maboiteresume}[2][]{colback=red!5!white, colframe=red!75!black,fonttitle=\bfseries, colbacktitle=red!85!black,enhanced,
attach boxed title to top center={yshift=-2mm},
  title={#2},#1}
  \newtcolorbox{maboiteresumeperso}[2][]{colback=red!5!white, colframe=red!75!black,fonttitle=\bfseries, colbacktitle=red!85!black,enhanced,
attach boxed title to top center={yshift=-2mm},
  title={#2},#1}
  \newtcolorbox{maboitecommentaire}[2][]{colback=red!5!white, colframe=red!75!black,fonttitle=\bfseries, colbacktitle=red!85!black,enhanced,
attach boxed title to top center={yshift=-2mm},
  title={#2},#1}
\newtcolorbox{maboitetruc}[2][]{colback=red!5!white, colframe=red!75!black,fonttitle=\bfseries, colbacktitle=red!85!black,enhanced,
attach boxed title to top left={yshift=-2mm},
  title={#2},#1}

%% Boîte des mises en évidence CARE. Couleurs pensées pour la
%% protanopie : l'axe bleu-jaune reste discriminant (jaune = \IMPORTANT
%% d'Olivier, bleu = CARE), et la distinction ne repose pas que sur la
%% couleur (cadre + bandeau étiqueté, lisibles même en niveaux de gris).
\newtcolorbox{maboiteimportantcare}[2][]{colback=cyan!10!white, colframe=blue!60!black,fonttitle=\bfseries, colbacktitle=blue!70!black,enhanced,
attach boxed title to top center={yshift=-2mm},
  title={#2},#1}
  


%%%% FIN STYLES



%%%%%%%%%%%%%%%%%%%%%%%%%%%%%%%%%%%%%%%%%%%%%%%%%
%%%%%%%%%%%%%%%%%%%%%%%%%%%%%%%%%%%%%%%%%%%%%%%%%
%
%
%                      Version ENSEIGNEMENT 2020
%
%
%%%%%%%%%%%%%%%%%%%%%%%%%%%%%%%%%%%%%%%%%%%%%%%%%
%%%%%%%%%%%%%%%%%%%%%%%%%%%%%%%%%%%%%%%%%%%%%%%%%

%english, french
\ifdefined\CESTDUFRANCAIS
	\providecommand\LANGUE[2]{#2}
	\usepackage[french]{babel}
\else
	\providecommand\LANGUE[2]{#1}
\fi
\providecommand\LANGUEinv[2]{\LANGUE{#2}{#1}}

%%%%%%%%%%%%%%%%%%%%%%%%%%%%%%%%%%%%%%%%%%%%%%%%%
%
%
%                      Packages
%
%
%%%%%%%%%%%%%%%%%%%%%%%%%%%%%%%%%%%%%%%%%%%%%%%%%


%%% SURLIGNEUR
\providecommand\ifnotmodeDIFFUSIONFINALE[1]{
%% 2021: Comprend pas mais assure
\ifdefined\SAFEMODE
\else
#1
\fi
%#1
}
\providecommand\ifmodeDIFFUSIONFINALE[1]{
\ifdefined\SAFEMODE
#1
\else
\fi
}
\newcommand\SURLIGNE[1]{
\ifnotmodeDIFFUSIONFINALE{
\BEGINTEMPENVCOLOR{blue!30}
{                        \colorbox{blue!30}{

\noindent \begin{minipage}{\dimexpr\textwidth-2\fboxsep\relax}
#1
\end{minipage}
}
}
\ENDTEMPENVCOLOR}
\ifmodeDIFFUSIONFINALE{
\noindent \begin{minipage}{\textwidth}
#1
\end{minipage}
}
}


\newenvironment{confiance}{\begin{maboiteresume}[colback=brown!50]{Résumé ici}}{\end{maboiteresume}}
\newcommand\CONFIANCE[2][Confiance ici]{
\begin{maboiteconfiance}[colback=yellow!50]{#1}
\textbf{ATTENTION : #2}
\end{maboiteconfiance}
}

\newenvironment{resume}{\begin{maboiteresume}[colback=brown!50]{Résumé ici}}{\end{maboiteresume}}
\newcommand\RESUME[2][Résumé ici]{
\begin{maboiteresume}[colback=brown!50]{#1}
#2
\end{maboiteresume}
}

\newenvironment{resumeperso}{\begin{maboiteresumeperso}[colback=brown!50]{Résumé perso ici}}{\end{maboiteresumeperso}}
\newcommand\RESUMEPERSO[2][Résumé perso  ici]{
\begin{maboiteresumeperso}[colback=brown!50]{#1}
#2
\end{maboiteresumeperso}
%
%
%\todo[inline, color=brown]{résumé: #1}
%\colorbox{brown}{
%\noindent \begin{minipage}{\textwidth}
%{
%\bf 
%#2
%}
%\end{minipage}
%}
}

\newcommand\DANSLAMARGE[1]{
\ifnotmodeDIFFUSIONFINALE{
\marginpar{\textcolor{brown}{#1}}}
}

\newenvironment{commentaireperso}{\begin{maboitecommentaire}[colback=brown!30]{Commentaire perso} Commentaire perso:}{\end{maboitecommentaire}}
\newcommand\COMMENTAIREPERSO[2][Commentaire perso]{
\begin{maboitecommentaire}[colback=brown!30]{#1}
#2
\end{maboitecommentaire}
}


\newcommand\SURSURLIGNE[2][truc surligné ici]{
%\todo[inline, color=orange]{surligné: #1}
\ifnotmodeDIFFUSIONFINALE{
\BEGINTEMPENVCOLOR{blue!60}
	                 \colorbox{blue!60}{
\noindent \begin{minipage}{\dimexpr\textwidth-2\fboxsep\relax}
{
\bf 
#2
}
\end{minipage}
}
\ENDTEMPENVCOLOR
}
\ifmodeDIFFUSIONFINALE{#2}
}

\newcommand\IMPORTANT[1]{
\ifnotmodeDIFFUSIONFINALE{
\BEGINTEMPENVCOLOR{yellow!100}
		        \colorbox{yellow!100}{

\noindent \begin{minipage}{\dimexpr\textwidth-2\fboxsep\relax}
{ 
#1
}
\end{minipage}
}
\ENDTEMPENVCOLOR
}
\ifmodeDIFFUSIONFINALE{#1}
}

%% Pendant de \IMPORTANT, réservé aux mises en évidence CARE ;
%% \IMPORTANT reste réservé aux annotations d'Olivier. Comme \IMPORTANT,
%% redevient du texte simple en mode DIFFUSIONFINALE.
\newcommand\IMPORTANTCARE[2][Important (CARE)]{
\ifnotmodeDIFFUSIONFINALE{
\begin{maboiteimportantcare}{#1}
#2
\end{maboiteimportantcare}
}
\ifmodeDIFFUSIONFINALE{#2}
}


\newcommand\QUESTION[1]{
\ifnotmodeDIFFUSIONFINALE{
\BEGINTEMPENVCOLOR{oceanfoam!100}
		        \colorbox{oceanfoam!100}{

\noindent \begin{minipage}{\dimexpr\textwidth-2\fboxsep\relax}
{ 
#1
}
\end{minipage}
}
\ENDTEMPENVCOLOR
}
\ifmodeDIFFUSIONFINALE{#1}
}




\newcommand\SOUSLIGNE[1]{
\colorbox{lightgray}{

\noindent \begin{minipage}{\dimexpr\textwidth-2\fboxsep\relax}
{ \color{gray}
#1
}
\end{minipage}
}
}



%% Gestion des cross-references cross-referencing
\usepackage{xr-hyper}
\usepackage{hyperref}

%%
\usepackage{versions}
\usepackage{amsmath,amssymb,mathrsfs,amsfonts}
%\usepackage{mathabx}
\usepackage{listings}
\usepackage{url}
\usepackage{versions}
\usepackage{color}
\usepackage[colorinlistoftodos]{todonotes}
%\usepackage{todonotes}
\usepackage{proof}
\usepackage{xstring}

%% Index
\usepackage{makeidx}
\makeindex
\providecommand\DOINDEXPRINT{}
\providecommand\SIGNALEMOTNOUV[1]{}
\ifdefined\INDEXPRINT
	\ifnotSAFEMODE{
%	\usepackage{showidx}
%	\let\oldindex\index
%	\renewcommand{\index}[1]{\oldindex{#1}\marginpar{LA: \tiny#1}}
	\renewcommand\DOINDEXPRINT{
	      	\let\oldindex\index
		\renewcommand{\index}[1]{\oldindex{##1}\marginpar{IDXLA: \tiny##1}}
		\renewcommand\SIGNALEMOTNOUV[1]{\marginpar{MNV: \small\textbf{##1}}}
	}
	}
\fi

% === gestion des accents


%soit iso latin1
%\usepackage[cyr]{aeguill}
%soit utf8
\usepackage[utf8]{inputenc}
\usepackage[T1]{fontenc}


%%%  Pour citation \Cref (cref, ref)	

\usepackage{cleveref}


%%%%%%%%%%%%%%%%%%%%%%%%%%%%%%%%%%%%%%%%%%%%%%%%%
%
%
%                      Gestion des accents spéciaux.
%
%
%%%%%%%%%%%%%%%%%%%%%%%%%%%%%%%%%%%%%%%%%%%%%%%%%

\input{macros-accents}
%% Tolérant tant que macros-markdown.tex n'est pas dans le dépôt :
\InputIfFileExists{macros-markdown}{}{\typeout{*** macros-markdown.tex introuvable: ignore ***}}


%\usepackage{PDFdraftcopy}


%-- Pour état des chapitres


%\newcommand\ETAT{\draftstring{BROUILLON}}
%
%\newcommand\brouillon{\renewcommand\ETAT{\draftstring{BROUILLON}}}
%
%\newcommand\final{\renewcommand\ETAT{\draftstring{\rien}}}
%
%\newcommand\travail{\renewcommand\ETAT{\draftstring{CHANTIER}}}
%
%\newcommand\bazar{\renewcommand\ETAT{\draftstring{BAZAR TOTAL}}}
%
%\final



%%%%%%%%%%%%%%%%%%%%%%%%%%%%%%%%%%%%%%%%%%%%%%%%%
%
%
%                      TRADUCTION EN COURS
%
%%%%%%%%%%%%%%%%%%%%%%%%%%%%%%%%%%%%%%%%%%%%%%%%%


% === volé Amaury



\usepackage{stmaryrd}
\usepackage{ifthen}
\usepackage{mathtools}
\usepackage{dsfont}
\PASSIAPXPROOF{\usepackage{thmtools}}
\usepackage{longtable}
\usepackage{mathdots}
\usepackage{booktabs}


%%%%%%%%%%%%%%%%%%%%%%%%%%%%%%%%%%%%%%%%%%%%%%%%%
%
%
%                      TIKZ STUFFS
%
%
%%%%%%%%%%%%%%%%%%%%%%%%%%%%%%%%%%%%%%%%%%%%%%%%%

%%% volé pour dessin des GPAC: PAS SUR QUE BESOIN de TOUT

\usepackage{tikz}
\usetikzlibrary{calc}
\usepgflibrary{arrows}
\usetikzlibrary{fit}
\usetikzlibrary{patterns}
\usetikzlibrary{decorations.pathreplacing}
\usetikzlibrary{shapes.multipart}
\usetikzlibrary{shapes.geometric}
\usetikzlibrary{shapes.misc}
\usetikzlibrary{positioning}
\usetikzlibrary{graphs}
\usetikzlibrary{topaths}
\usetikzlibrary{arrows.meta}
\usetikzlibrary{decorations.pathreplacing}
\usetikzlibrary{decorations.markings}
\usetikzlibrary{decorations.pathmorphing}
%
 \usetikzlibrary{calc}
 \usetikzlibrary{automata}
\usetikzlibrary{chains}
\usetikzlibrary{scopes}
 \usetikzlibrary{positioning}
 \usetikzlibrary{through,backgrounds}
\usepackage{circuitikz}
% % pour accolades




%%%%%%%%%%%%%%%%%%%%%%%%%%%%%%%%%%%%%%%%%%%%%%%%%
%
%
%                      MACROS locales
%
%
%%%%%%%%%%%%%%%%%%%%%%%%%%%%%%%%%%%%%%%%%%%%%%%%%


%====================
%                      Moi
%=====================


\newcommand\myaddress{ 
{\sc LIX, CNRS, Ecole Polytechnique, Institut Polytechnique de Paris, }\\ 91128 Palaiseau Cedex, FRANCE \\
{\small Olivier.Bournez@lix.polytechnique.fr}}



%====================
%                      un peu de trucs sales
%=====================

%=  pour style lipics et ses trucs

\providecommand\ifnotPOLYCHAPTERMODE[1]{#1}
\providecommand\ifPOLYCHAPTERMODE[1]{}

\ifnotPOLYCHAPTERMODE{
\ifdefined\thechapter
\providecommand\spnewtheorem[4]{\newtheorem{#1}{#2}[chapter]}
\providecommand\spnewtheoremi[5]{\newtheorem{#1}[#5]{#2}}
\else
\providecommand\spnewtheorem[4]{\newtheorem{#1}{#2}}
\providecommand\spnewtheoremi[5]{\newtheorem{#1}[#5]{#2}}
\fi
}
\ifPOLYCHAPTERMODE{
\providecommand\spnewtheorem[4]{\newtheorem{#1}{#2}}
\providecommand\spnewtheoremi[5]{\newtheorem{#1}[#5]{#2}}
}
\providecommand\nolipics[1]{#1}




%====================
%                      un peu de trucs franchement sales
%=====================


% on l'inclus pas ici.
%\input{macros-moins-propres}

%%%%%%%%%%%%%%%%%%%%%%%%%%%%%%%%%%%%%%%%%%%%%%%%%
%
%
%                      Page de Garde Cours X
%
%
%%%%%%%%%%%%%%%%%%%%%%%%%%%%%%%%%%%%%%%%%%%%%%%%%


\includeversion{metdate}

\newcommand\XPAGEDEGARDE[2]{
\thispagestyle{empty}
\title{{#1} \\[2cm] %\\[3cm] 
{\large #2}
}
\author{
 Olivier Bournez% \inst{1}
\\[0.5cm] 
bournez@lix.polytechnique.fr%{\myaddress}
}
%\date{
%\vspace{2cm}
%\includegraphics[height=2.5cm]{/Users/bournez/lib/LaTeX/Perso/FIG-COMMONS/logo_X_new}
%}
\begin{metdate}
\date{
\vspace{2cm}
Version \LANGUE{of}{du} \today \\[1cm] 
\includegraphics[height=3.5cm]{/Users/bournez/lib/LaTeX/Perso/FIG-COMMONS/logo_X_new}
}
\end{metdate}
\maketitle
}


%%%%%%%%%%%%%%%%%%%%%%%%%%%%%%%%%%%%%%%%%%%%%%%%%
%
%
%                      MACROS UTILES
%
%
%%%%%%%%%%%%%%%%%%%%%%%%%%%%%%%%%%%%%%%%%%%%%%%%%

%%%                      Macros
 
%-------------
%--- vectors  ---
%-------------

 
\newcommand\vectorl[1]{{\mathbf#1}}
\newcommand\tu[1]{\vectorl{#1}}

\newcommand\vp{\vectorl{p}}
\newcommand\va{\vectorl{a}}
\newcommand\vb{\vectorl{b}}
\newcommand\vc{\vectorl{c}}
\newcommand\vf{\vectorl{f}}
\newcommand\vn{\vectorl{n}}
\newcommand\vx{\vectorl{x}}
\newcommand\vX{\vectorl{X}}
\newcommand\vy{\vectorl{y}}
\newcommand\vz{\vectorl{z}}
\newcommand\ve{\vectorl{e}}
\newcommand\vv{\vectorl{v}}
\newcommand\vr{\vectorl{r}}
\newcommand\vu{\vectorl{u}}
\newcommand\vA{\vectorl{A}}
\newcommand\vB{\vectorl{B}}
\newcommand\vC{\vectorl{C}}
\newcommand\vD{\vectorl{D}}

%-------------
%--- overline ---
%-------------


\newcommand\ox{\overline{x}}
\newcommand\oy{\overline{y}}
\newcommand\oc{\overline{c}}
\newcommand\oa{\overline{a}}
\newcommand\ow{\overline{w}}
\newcommand\od{\overline{d}}
\newcommand\ob{\overline{b}}
\newcommand\oP{\overline{P}}
\newcommand\oee{\overline{e}}
\newcommand\ot{\overline{t}}
\newcommand\ou{\overline{u}}
\newcommand\ov{\overline{v}}
\newcommand\oz{\overline{z}}
\newcommand\am{\overline{a}}
\newcommand\om{\overline{m}}
\newcommand\oj{\overline{j}}

%----------------
%--- ensembles ---
%----------------

\newcommand\Q{\mathbb{Q}}
\newcommand\R{\mathbb{R}}
\newcommand\Z{\mathbb{Z}}
\providecommand\C{\mathbb{C}}
\newcommand\N{\mathbb{N}}
\newcommand\K{\mathbb{K}}
\newcommand\F{\mathbb{F}}

\newcommand\dyadic{\mathbb{D}}
\newcommand\Zd{{\Z_2}}
\newcommand\Zp{{\Z_p}}
%
\newcommand\RP{\mathbb{R}^{>0}}
\newcommand\QP{\mathbb{Q}^{>0}}


%----------------
%--- lettres cal ---
%----------------

\newcommand\mK{\mathcal{K}}
\newcommand\mD{\mathcal{D}}
\newcommand\mA{\mathcal{A}}
\newcommand\mL{\mathcal{L}}
\newcommand\mX{\mathcal{X}}
\newcommand\mG{\mathcal{G}}
\newcommand\mE{\mathcal{E}}
\newcommand\mH{\mathcal{H}}
\newcommand\mR{\mathcal{R}}
\newcommand\mF{\mathcal{F}}
\newcommand\mJ{\mathcal{J}}
\newcommand\mO{\mathcal{O}}
\newcommand\mP{\mathcal{P}}
\newcommand\mN{\mathcal{N}}
\newcommand\mS{\mathcal{S}}
\newcommand\mM{\mathcal{M}}
\newcommand\mC{\mathcal{C}}
\newcommand\mV{\mathcal{V}}
\newcommand\mI{\mathcal{I}}
\newcommand\mT{\mathcal{T}}

\newcommand\fC{\mathcal{C}}

%----------------
%--- lettres frak ---
%----------------


\newcommand\kM{\mathfrak{M}}
\newcommand\kN{\mathfrak{N}}
\newcommand\kU{\mathfrak{U}}
\newcommand\kg{\mathfrak{g}}



%--------------------------
%--- Façon noter noms classes ---
%---------------------------

\newcommand{\cp}[1]{\operatorname{#1}}


%--------------------------
%--- Calculabilité  ---
%---------------------------

% -- classes
\LANGUE{
\newcommand\RE{\cp{CE}}
}
{\newcommand\RE{\cp{RE}}
}

\newcommand\Rec{\cp{D}}
\newcommand\PR{\cp{PR}}
\newcommand\Gregn{\mE_n}


%--------------------------
%--- Réductions  ---
%---------------------------

\newcommand\lem{\le_{m}}
\newcommand\lelog{\le_{L}}
\newcommand\equivlog{\equiv_{L}}



%--------------------------
%--- Complexité  ---
%---------------------------

% COMPLEXITE

% -- P
\newcommand{\Ptime}{\cp{P}}      % P
\newcommand\PTIME{\Ptime}
\newcommand\PTIMEs[1]{{\PTIME}_{#1}}
\newcommand\PTIMEsp[1]{{\PTIME}_{#1}^0}
\renewcommand\P{\PTIME}
\newcommand{\FPtime}{\cp{FPTIME}}

% -- NP
\newcommand\NP{\cp{NP}}
\newcommand{\NPtime}{\NP}
\newcommand\NPs[1]{\cp{NP}_{#1}}
\newcommand{\FNP}{\cp{FNP}}

\newcommand\NDP{\cp{NDP}}
\newcommand\coNDP{\cp{coNDP}}
\newcommand\NDPs[1]{\cp{NDP}_{#1}}

% -- coNP
\newcommand\coNP{\cp{coNP}}

% -- PSPACE
\newcommand\PSPACE{\cp{ PSPACE}}
\newcommand\NPSPACE{\cp{ NPSPACE}}
\newcommand{\Pspace}{\PSPACE}
\newcommand{\FPspace}{\cp{FPSPACE}}


% -- LOGSPACE
\newcommand\LSPACE{\cp{LOGSPACE}}
\newcommand\NLSPACE{\cp{NLOGSPACE}}
\newcommand\coNLSPACE{\cp{coNLOGSPACE}}
\newcommand\coNL{\cp{coNL}}

% -- Autres
\newcommand\EXPTIME{\cp{EXPTIME}}
\newcommand\NEXPTIME{\cp{NEXPTIME}}
\newcommand\EXPSPACE{\cp{EXPSPACE}}
\newcommand\PH{\cp{PH}}
\newcommand\NC{\cp{NC}}
\newcommand\AC{\cp{AC}}
\newcommand\PPAD{\cp{\mathbf{PPAD}}}
\newcommand\PLS{\cp{\mathbf{PLS}}}

% -- Reverse math
\newcommand{\WKL}{\ensuremath{\docheck{\mathsf{WKL}_0}}}
\newcommand{\ACA}{\ensuremath{\docheck{\mathsf{ACA}_0}}}
\newcommand{\ATR}{\ensuremath{\docheck{\mathsf{ATR}_0}}}
\newcommand{\RCAO}{\ensuremath{\docheck{\mathsf{RCA}_0}}}
\newcommand{\PiCA}{\ensuremath{\docheck{\Pi^1_1\!\text{-}\mathsf{CA}_0}}}


% -- circuits
\newcommand\CIRCUIT[2]{\cp{CIRCUIT(#1,#2)}}
\newcommand\UCIRCUIT[2]{\cp{UCIRCUIT(#1,#2)}}

% -- Classes proba
\newcommand\PP{\cp{PP}}
\newcommand\RPP{\cp{RP}}
\newcommand\ZPP{\cp{ZPP}}
\newcommand\coRP{\cp{coRP}}
\newcommand\BPP{{\cp{BPP}}}

% -- IP
\newcommand\IP{\cp{IP}}
\newcommand\PCP{\cp{PCP}}

% -- non-uniforme
\newcommand\nuP{\mathbb{P}}
\newcommand\nuPs[1]{\mathbb{P}_{#1}}
\newcommand\nuPsp[1]{\mathbb{P}_{#1}^0}
\newcommand\nuNP{\mathbb{N}\mathbb{P}}
%\newcommand\poly{/{poly}}
\newcommand\Ppoly{\PTIME_{\poly}}
\newcommand\NPpoly{\NP_{\poly}}

% -- conditions d'uniformité
\newcommand\Luniforme{\cp{L}}
\newcommand\BCuniforme{\cp{BC}}


% Classes de circuits
\newcommand\TC{\cp{TC}}
\newcommand\LFP{\cp{LFP}}
\newcommand\FAC{\cp{FAC}}
\newcommand\FACC{\cp{FACC}}
\newcommand\FTC{\cp{FTC}}
\newcommand\FNC{\cp{FNC}}

% Logiques temorelles
\newcommand\CTL{\cp{CTL}}
\newcommand\LTL{\cp{LTL}}

% -- mesure temps/espace/taille

\newcommand\DTIME{\cp{DTIME}}
\newcommand\TIME[1]{\cp{TIME(#1)}}
\newcommand\TIMEs[2]{\cp{TIME_{#2}(#1)}}
\newcommand\TIMEsp[2]{\cp{TIME^0_{#2}(#1)}}
\newcommand\NTIME[1]{\cp{NTIME(#1)}}
\newcommand\NTIMEs[2]{\cp{NTIME_{#2}(#1)}}
\newcommand\NTIMEsp[2]{\cp{NTIME_{#2}^0(#1)}}
\newcommand\DSPACE{\cp{DSPACE}}
\newcommand\SPACE[1]{\cp{SPACE(#1)}}
\newcommand\SPACEs[2]{\cp{SPACE_{#2}(#1)}}
\newcommand\NSPACE[1]{\cp{NSPACE(#1)}}
\newcommand\NSPACEs[2]{\cp{NSPACE_{#2}(#1)}}
\newcommand\coNSPACE[1]{\cp{coNSPACE(#1)}}
\newcommand\TIMEON[2]{\cp{TIME(#1,#2)}}
\newcommand\SPACEON[2]{\cp{SPACE(#1,#2)}}
\newcommand\SIZE[1]{\cp{ size(#1)}}
\newcommand\ATIME{\cp{ATIME}}
\newcommand\ASPACE{\cp{ASPACE}}


%% autre:
\newcommand\LH{\cp{LH}}
\newcommand\FO{\cp{FO}}
\newcommand\SO{\cp{SO}}
\newcommand\ESOhorn{\cp{ESO_{horn}}}
\newcommand\SOhorn{\cp{SO_{horn}}}
\newcommand\ACzero{\cp{AC_0}}
\newcommand{\LINSPACE}{\cp{LINSPACE}}
\newcommand{\mFLINSPACE}{\mF_{\cp{LINSPACE}}}
\newcommand{\LOGSPACE}{\cp{LOGSPACE}}
\newcommand\ALOGTIME{\cp{ALOGTIME}}
\newcommand\RF{\cp{RF}}


\newcommand\BOOL[1]{\cp{{BOOLEAN(#1)}}}
\newcommand\DET{\cp{ DET}}

% % -- dirty


\newcommand{\cPspace}{\cp{\sharp PSPACE}}
\newcommand{\cAPtime}{\cp{\sharp APTIME}}
%\newcommand{\FPspace}{\cp{FPSPACE}}
%\newcommand{\FPspaceN}{\cp{FPSPACE}^{+}}
%\newcommand{\FPspaceN}{\cp{FPSPACE}}



%-------------------------
%--- Problèmes de décision  ---
%--------------------------

\newcommand\decisionname[1]{\cp{#1}}


\newcommand\Luniv{\decisionname{L_{univ}}}
\newcommand\HALTINGPROBLEM{\decisionname{HALTING-PROBLEM}}
\newcommand\SAT{\decisionname{SAT}}
\newcommand\MAXtSAT{\decisionname{MAX3SAT}}
\newcommand\REACH{\decisionname{REACH}}
\newcommand\CIRCUITVALUE{\decisionname{CIRCUITVALUE}}
\newcommand\MONOTONECIRCUITVALUE{\decisionname{MONOTONECIRCUITVALUE}}
\newcommand\THEOREMS{\decisionname{THEOREMS}}
\newcommand\QCIRCUITSAT{\decisionname{QCIRCUITSAT}}
\newcommand\QSAT{\decisionname{QSAT}}
\newcommand\QBF{\decisionname{QBF}}
\newcommand\SATVALUE{\decisionname{SATVALUE}}
\newcommand\SUCCINTSAT{\decisionname{SUCCINTSAT}}
\newcommand\SUCCINTSATVALUE{\decisionname{SUCCINTSATVALUE}}
\newcommand\GEOGRAPHY{\decisionname{GEOGRAPHY}}
\newcommand\PARITY{\decisionname{PARITY}}
\newcommand\MAJ{\decisionname{MAJ}}
\newcommand\CVP{\CIRCUITVALUE}
\newcommand\BOOLEANPROGRAMVALUE{BOOLEAN~PROGRAM~VALUE}

\newcommand\CLIQUE{\cp{CLIQUE}}
\LANGUEinv{
\newcommand\RECOUVREMENTDESOMMETS{\cp{RECOUVREMENT~DE~SOMMETS}}
\newcommand\RS{\RECOUVREMENTDESOMMETS}
}{
\newcommand\RECOUVREMENTDESOMMETSeng{\cp{VERTEX~COVER}}
\newcommand\RSeng{\RECOUVREMENTDESOMMETSeng}
}

\LANGUEinv{\newcommand\COUPUREMAXIMALE{\cp{COUPURE~MAXIMALE}}}
{\newcommand\COUPUREMAXIMALEeng{\cp{MAXIMAL~CUT}}}
\newcommand\CIRCUITHAMILTONIEN{\CHAMILTONIEN}
\LANGUEinv{\newcommand\VOYAGEURDECOMMERCE{\VCOMMERCE}}
{\newcommand\VOYAGEURDECOMMERCEeng{\VCOMMERCEeng}}

\LANGUEinv{\newcommand\CIRCUITLEPLUSLONG{\cp{CIRCUIT~LE~PLUS~LONG}}}
{\newcommand\CIRCUITLEPLUSLONGeng{\cp{LONGEST~CIRCUIT}}}
\LANGUEinv{\newcommand\SOMMEDESOUSENSEMBLE{\SUBSETSUM}}{
\newcommand\SOMMEDESOUSENSEMBLEeng{\SUBSETSUMeng}}
\LANGUEinv{
\newcommand\SACADOS{\cp{SAC~A~DOS}}}
{\newcommand\SACADOSeng{\cp{KNAPSACK}}}

%français
\LANGUEinv{
\newcommand\CHAMILTONIEN{\cp{CIRCUIT~HAMILTONIEN}}}
{\newcommand\CHAMILTONIENeng{\cp{HAMILTONIAN~CIRCUIT}}}
\LANGUEinv{
\newcommand\VCOMMERCE{\cp{VOYAGEUR~DE~COMMERCE}}}
{\newcommand\VCOMMERCEeng{\cp{TRAVELING~SALESMAN}}}
\LANGUEinv{
\newcommand\kCOLORABILITE{\cp{k-COLORABILITE}}}
{\newcommand\kCOLORABILITEeng{\cp{k-COLORABILITY}}}
\LANGUEinv{
\newcommand\COLORIAGE{\cp{BON~COLORIAGE}}}
{\newcommand\COLORIAGEeng{\cp{GOOD~COLORING}}}
\LANGUEinv{
\newcommand\tCOLORABILITE{\cp{3-COLORABILITE}}}
{\newcommand\tCOLORABILITEeng{\cp{3-COLORABILITY}}}
\LANGUEinv{
\newcommand\SUBSETSUM{\cp{SOMME~DE~SOUS~ENSEMBLE}}}
{\newcommand\SUBSETSUMeng{\cp{SUBSET~SUM}}}
\newcommand\PARTITION{\cp{PARTITION}}
\LANGUEinv{
\newcommand\NOMBREPREMIER{\decisionname{NOMBRE~ PREMIER}}}
{\newcommand\NOMBREPREMIEReng{\decisionname{PRIME~NUMBER}}}
\LANGUEinv{
\newcommand\CODAGE{\decisionname{CODAGE}}}
{\newcommand\CODAGEeng{\decisionname{ENCODING}}}
\newcommand\kSAT{\cp{k-SAT}}
\newcommand\tSAT{\cp{3-SAT}}
\newcommand\NAESAT{\cp{NAESAT}}
\newcommand\NAEtSAT{\cp{NAE3SAT}}
\newcommand\NAEqSAT{\cp{NAE4SAT}}
\newcommand\STABLE{\cp{\LANGUE{INDEPENDANT~SET}{STABLE}}}

%\newcommand\CM{\cp{COUPURE MAXIMALE}}
\LANGUEinv{
\newcommand\POSTpb{\cp{Problème~de~ la~correspondance~de~Post}}}
{\newcommand\POSTpbeng{\cp{Post~correspondence~problem}}}

\LANGUEinv{
\newcommand\PARITE{\cp{PARITE}}}
{\newcommand\PARITEeng{\cp{PARITY}}}


%-------------------------
%--- Modèles parallèles  ---
%--------------------------


% PARALLELISME
\newcommand\RAM{\cp{RAM}}
%
\newcommand\PRAM{\cp{PRAM}}
\newcommand\CRCW{\cp{CRCW}}
\newcommand\CREW{\cp{CREW}}
\newcommand\EREW{\cp{EREW}}


%------------------------
%--- codages  ---
%------------------------


\newcommand\codage[1]{\langle #1 \rangle}
\newcommand\tuple[1]{\codage{#1}}
\newcommand\paire[2]{\codage{#1,#2}}
\newcommand\triplet[3]{\codage{#1,#2,#3}}
\newcommand\quadruplet[4]{\codage{#1,#2,#3,#4}}
\newcommand\cinquplet[5]{\codage{#1,#2,#3,#4,#5}}


%------------------------
%--- descriptive set theory  ---
%------------------------

\newcommand\baire{\mathcal{N}}
\newcommand\peano{\overline{\mathcal{N}}} %% A MOI
\newcommand\cantor{\mathcal{C}}
%
\newcommand\bSigma{\mathbf{\Sigma}}
\newcommand\bPi{\mathbf{\Pi}}
\newcommand\bDelta{\mathbf{\Delta}}
%


%%%%%% ABOUT NOMS QUI CHANGES

\newcommand\WF{\cp{WF}}
\newcommand\WO{\WF}
\newcommand\LO{\cp{LO}}
\newcommand\DLO{\cp{DLO}}
% 
\newcommand\rk{\cp{rank}}


%% Moins propre:

\newcommand\I{\mathbb{I}}
\renewcommand\H{\mathbb{H}}
%
\newcommand\leKB{<_{KB}}
%
\newcommand\finiomega{{<\omega}}
\newcommand\compl{^\frown}
\newcommand\prefix{ ~ prefix~  } 
\newcommand\X{{X}}
\newcommand\Y{\mathcal{Y}}
\newcommand\subsetstrict{\subset}

\renewcommand\complement[1]{#1^{c}}


%------------------------
%--- schémas de récursion ---
%------------------------

%%                      Source papier MFCS 2019


% COMPLEXITY

%% MFCS differe de Clote:
%% Version à taille restreinte

\newcommand\co{\cp{co}}

\newcommand\mFPTIME{\mF_{PTIME}}
\newcommand{\mFPSPACE}{\mF_{SPACE}}
\newcommand{\mFPH}{\mF_{PH}}

\newcommand\SigmaP{\Sigma^P}
\newcommand\DeltaP{\Delta^{P}}
\newcommand\PiP{\Pi^{P}}
\newcommand\coSigmaP{\co\SigmaP}
\newcommand\coPiP{\mF{\co\PiP}}

%% autres classes:


%% Schemas pour complexité implicite

\newcommand\COMP{\cp{ COMP}}
\newcommand\CRN{\cp{ CRN}}
\newcommand\BRN{\cp{ BRN}}
\newcommand\SBRN{\cp{ SBRN}}
\newcommand\BR{\cp{ BR}}
\newcommand\kBR{k-\cp{ BR}}
\newcommand \WBPR{\cp{ WBPR}}
\newcommand \BMIN{\cp{ BMIN}}
\newcommand \BSUM{\cp{ BSUM}}
\newcommand\SBSUM{\cp{ SBSUM}}
\newcommand \BPROD{\cp{ BPROD}}
\newcommand \SBPROD{\cp{ SBPROD}}
\newcommand \SCOMP{\cp{ SCOMP}}
\newcommand \SR{\cp{ SR}}
\newcommand \SRN{\cp{ SRN}}

%% devrait etre défini
\newcommand\WBRN{\cp{ WBRN}}



%-- godel
\newcommand\INTDIV{\cp{INTDIV}}
\newcommand\DIV{\cp{DIV}}
\newcommand\EVEN{\cp{EVEN}}
\newcommand\ODD{\cp{ODD}}
\newcommand\PRIME{\cp{PRIME}}
\newcommand\POWER{\cp{POWER}}
\newcommand\BIT{\cp{BIT}}
\newcommand\negHP{\overline{\cp{HP}}}
\newcommand\HP{\cp{HP}}
\newcommand\negLuniv{\overline{\Luniv}}
\newcommand\VALCOMP{\cp{VALCOMP}}
\newcommand\mWINDOW{\cp{LEGAL}}
\newcommand\mmWINDOW{\cp{WINDOW}}
\newcommand\LENGTH{\cp{LENGTH}}
\newcommand\DIGIT{\cp{DIGIT}}
\newcommand\MATCH{\cp{MATCH}}
\newcommand\MOVE{\cp{MOVE}}
\newcommand\START{\cp{START}}
\newcommand\CELL{\cp{CELL}}
\newcommand\HALT{\cp{HALT}}
\newcommand\SUBST{\cp{SUBST}}
\newcommand\PROOF{\cp{PROOF}}
\newcommand\PROVABLE{\cp{PROVABLE}}
\newcommand\CONSIST{\cp{CONSIST}}



%---------------------
%--- calculus (AP) ---
%---------------------

% \jacobian{f}{x}
\newcommand{\jacobian}[1]{J_{#1}}
%\newcommand{\grad}[1]{\overrightarrow{\operatorname{grad}}~{#1}}
\newcommand{\grad}[1]{\nabla{#1}}
\newcommand{\scalarprod}[2]{{#1}\cdot{#2}}
\newcommand{\bigO}[1]{\mathcal{O}\left(#1\right)}
\newcommand\smallO{o}
\newcommand{\softO}[1]{\tilde{\mathcal{O}}\left(#1\right)}
\newcommand{\taylor}[3]{{T_{#1}^{#2}#3}}
\newcommand{\transpose}[1]{{#1}^T}
\newcommand{\pastsup}[2]{{\sup}_{#1}#2}

\DeclareMathOperator\arctanh{arctanh}



%---------------
%--- norms (AP) ---
%---------------


\newcommand{\inorm}[2]{\left\lVert{#1}\right\rVert_{#2}}
\newcommand{\twonorm}[1]{\inorm{#1}{2}}
\newcommand{\infnorm}[1]{\inorm{#1}{}}
\newcommand{\normsup}[1]{\infnorm{#1}{}}
%
\newcommand\hausdorff{d_H}
%
\newcommand\norm[1]{\infnorm{#1}}


%----------------
%--- topology ---
%----------------

% \openball{pt}{radius}
 \newcommand{\openball}[2]{B_{#2}(#1)}
\newcommand{\closedball}[2]{\ovl{B}_{#2}(#1)}
\newcommand{\closure}[1]{\ovl{#1}}
% \newcommand{\interior}[1]{\mathring{#1}}
% \newcommand{\border}[1]{\partial{#1}}




%-----------------
%--- functions (AP) ---
%-----------------
 
\newcommand{\lambdafun}[2]{(#1)\mapsto{#2}}
\newcommand{\lambdafunex}[2]{#1\mapsto{#2}}
\newcommand{\argmin}{\operatornamewithlimits{argmin}}
\newcommand{\argmax}{\operatornamewithlimits{argmax}}


 
%-----------------------
%--- Machines de Turing  ---
%-----------------------

\newcommand\blanc{\vectorl{B}}



%-----------------
%--- Logique ---
%-----------------
 
\newcommand\sem[1]{{[\semspace[}#1{]\semspace]}}
\newcommand\sems[1]{{[\semspace[}#1{]\semspace]}}
\newcommand\semss[1]{\sem{#1}_{\sigma'}}
\newcommand\semspace{\kern -.25ex}
\newcommand\true{true}
\newcommand\false{false}
%extension élémentaire:
\newcommand\elem{\preceq}
\newcommand\godel[1]{\lceil #1 \rceil}

%-----------------
%--- Model theory ---
%-----------------

\newcommand\hull[1]{\langle #1 \rangle}

%-----------------
%--- Ordinaux ---
%-----------------

\newcommand\omegaunCK{\omega_1^{CK}}
\newcommand\omegaunck{\omegaunCK}

\DeclareMathOperator{\cof}{cof}

%-----------------
%--- Surréels (QG) ---
%-----------------

\newcommand{\On}{\textbf{On}}
\newcommand{\Nobf}{\textbf{No}}
\newcommand{\Nobb}{\ensuremath{\mathbbmss{No}}}
\newcommand{\Nolt}[1]{\ensuremath{\Nobf_{< #1}}}
\newcommand{\Noleq}[1]{\ensuremath{\Nobf_{\leq #1}}}
\newcommand{\Noltbb}[1]{\ensuremath{\Nobb_{< #1}}}
\providecommand\No{\Nobf}
\newcommand\Noalpha{\Nolt{\alpha}}
\newcommand\Nolambda{\Nolt{\lambda}}

\DeclareMathOperator{\length}{length}
\DeclareMathOperator{\supp}{supp} % operateur support


%----------------------
%---Maths  de base (AP)  ---
%----------------------

\DeclareMathOperator{\dom}{dom} % operateur support

\DeclareMathOperator{\size}{size}

\newcommand{\powerset}[1]{\mathcal{P}(#1)}
\newcommand{\card}[1]{{\##1}}
\newcommand\priv{ - }
\newcommand\prive{-}
\renewcommand{\restriction}{\mathord{\upharpoonright}}
\newcommand\restrict[1]{_{\restriction #1}}

% \providecommand\mod{}
%% UNDEFINE:
\let\mod\relax
\let\div\relax
\DeclareMathOperator{\mod}{mod} 
\DeclareMathOperator{\div}{div} 

% partieentier
\newcommand\entieres[1]{\lceil #1 \rceil}
\newcommand\entierei[1]{\lfloor #1 \rfloor}


%----------------------
%---Symbole danger ---
%----------------------


%\providecommand*{\TakeFourierOrnament}[1]{{%
%\fontencoding{U}\fontfamily{futs}\selectfont\char#1}}
%\providecommand*{\danger}{{\Huge \TakeFourierOrnament{66}}}



%%%%%%%%%%%%%%%%%%%%%%%%%%%%%%%%%%%%%%%%%%%%%%%%%
%
%
%                      ASMeries
%
%
%%%%%%%%%%%%%%%%%%%%%%%%%%%%%%%%%%%%%%%%%%%%%%%%%



% CIRCUITS
%-- particuliers
\newcommand\REPLACE{\cp{REPLACE}}
\newcommand\EQUAL{\cp{EQUAL}}
\newcommand\EVALUE{\cp{EVALUE}}
\newcommand\TVALUE{\cp{TVALUE}}
\newcommand\UPDATE{\cp{UPDATE}}
\newcommand\ASSIGN{\cp{ASSIGN}}
\newcommand\PAR{\cp{PAR}}
\newcommand\DO{\cp{DO}}


% ASM
\newcommand\emplacement[2]{(#1,#2)}
\newcommand\contenu[2]{\sem{(#1,#2)}}
\newcommand\affecte[3]{(#1,#2)\leftarrow#3}
\newcommand\ctl{ctl\_state}


%%%%%%%%%%%%%%%%%%%%%%%%%%%%%%%%%%%%%%%%%%%%%%%%%
%
%
%                      GPACqueries
%
%
%%%%%%%%%%%%%%%%%%%%%%%%%%%%%%%%%%%%%%%%%%%%%%%%%



%------------------------
%--- generable zoo (GPAC) ---
%------------------------

\newcommand{\myop}[1]{\operatorname{#1}}
\newcommand{\abs}{\myop{abs}}
\newcommand{\sg}{\myop{sg}}
\newcommand{\ip}[1]{\myop{ip}_{#1}}
\newcommand{\rnd}{\myop{rnd}}
\newcommand{\lil}{\myop{lil}}
\newcommand{\lxh}{\myop{lxh}}
\newcommand{\hxl}{\myop{hxl}}
\newcommand{\plil}{\myop{plil}}
\newcommand{\reach}{\myop{reach}}
%\newcommand{\mreach}{\myop{mreach}}
\newcommand{\watchdog}{\myop{watchdog}}
\newcommand{\monitor}{\myop{monitor}}
%\newcommand{\norm}{\myop{norm}}
\newcommand{\mx}{\myop{mx}}
\newcommand{\mn}{\myop{mn}}
\newcommand{\sample}{\myop{sample}}
\newcommand{\select}{\myop{select}}
\newcommand{\ssine}[1]{\sin_{#1}}
\newcommand{\clamp}{\myop{clamp}}

\newcommand{\fnsg}[3]{tanh((#1)*(#2)*(#3))}
\newcommand{\fnipone}[3]{(1+\fnsg{(#1)-1}{#2}{#3})/2.}
\newcommand{\fnipinf}[1]{#1-sin(2*pi*(#1))/2./pi}%
\newcommand{\fnlil}[5]{%
\fnipone{(#3)-(#1)+1-((#2)-(#1))/8.}{#4+log(1+4./(#2-#1)*(#5)*(#5))+log(1+4)}{8./(#2-#1)}%
*\fnipone{#2-#3+1-(#2-#1)/8.}{#4+log(1+4/(#2-#1)*(#5)*(#5))+log(1+4)}{8./(#2-#1)}%
*4./(#2-#1)*(#5)}
\newcommand{\fnlxh}[5]{\fnipone{(#3)-(#1+#2)/2.+1}{#4+log(1+(#5)*(#5))}{2./(#2-#1)}*(#5)}
\newcommand{\fnhxl}[5]{\fnipone{(#1+#2)/2.-(#3)+1}{#4+log(1+(#5)*(#5))}{2./(#2-#1)}*(#5)}
\newcommand{\fnplilf}[4]{sin(((#4)-((#1)+(#2))/2.+(#3)/4.)*2.*pi/(#3))}
\newcommand{\fnplil}[6]{(#6)*\fnlxh{\fnplilf{#1}{#2}{#3}{#1}}%
{\fnplilf{#1}{#2}{#3}{#1+((#2)-(#1))/4.}}%
{\fnplilf{#1}{#2}{#3}{#4}}%
{#5+2.+log(1+(#6)*(#6))}%
{1./4.+2./((#2)-(#1))}}
\newcommand{\fnssine}[2]{sin(#2)*(1-cos(#2)/cos(#1))}



%--------------------
%--- complexity GPAC ---
%--------------------


%\ifthenelse{\equal{#1}{}}{Empty.}{Nonempty.}
\newcommand{\myclass}[1]{\operatorname{#1}}
% \newcommand{\gcomp}[2]{\ensuremath{\myclass{GCOMP}(#1,#2)}}
% \newcommand{\ggenerable}[1]{\textcolor{red}{\ensuremath{#1-}\text{generable}}}
 \newcommand{\gval}[2][]{\ensuremath{\myclass{GVAL}_{#1}\ifthenelse{\equal{#2}{}}{}{[#2]}}}
 \newcommand{\gpval}[1][]{\ensuremath{\myclass{GPVAL}_{#1}}}
 \newcommand{\gc}[3][]{\ensuremath{\myclass{ATSC}(#2,#3)}}
 \newcommand{\gpc}[1][]{\ensuremath{\myclass{ATSP}}}
\newcommand{\cgc}[1][]{\ensuremath{\myclass{ATS}}}
\newcommand{\gexpc}[1][]{\ensuremath{\myclass{AEXP}}}
\newcommand{\gwc}[2]{\ensuremath{\myclass{AW}(#1,#2)}}
\newcommand{\gpwc}{\ensuremath{\myclass{AWP}}}
\newcommand{\cgwc}{\ensuremath{\myclass{AW}}}
\newcommand{\grc}[3]{\ensuremath{\myclass{AR}(#1,#2,#3)}}
\newcommand{\gprc}{\ensuremath{\myclass{ARP}}}
\newcommand{\gsc}[3]{\ensuremath{\myclass{AS}(#1,#2,#3)}}
\newcommand{\gpsc}{\ensuremath{\myclass{ASP}}}
\newcommand{\goc}[3]{\ensuremath{\myclass{AO}(#1,#2,#3)}}
\newcommand{\gpoc}{\ensuremath{\myclass{AOP}}}
\newcommand{\cgoc}{\ensuremath{\myclass{AO}}}
\newcommand{\guc}[4]{\ensuremath{\myclass{AX}(#1,#2,#3,#4)}}
\newcommand{\gpuc}{\ensuremath{\myclass{AXP}}}
\newcommand{\glc}[2][]{\ensuremath{\myclass{AL}(#2)}}
\newcommand{\gplc}[1][]{\ensuremath{\myclass{ALP}}}
\newcommand{\cglc}[1][]{\ensuremath{\myclass{AL}}}

%\newcommand{\intinterv}[2]{\{#1,\ldots,#2\}}
\newcommand{\intinterv}[2]{\llbracket#1,#2\rrbracket}

\newcommand{\Rgen}{\R_{G}}
\newcommand{\Rpoly}{\R_{P}}

%%%%%%%%%%%%%%%%%%%%%%%%%%%%%%%%%%%%%%%%%%%%%%%%%
%
%
%                      Mes codages obscures
%
%
%%%%%%%%%%%%%%%%%%%%%%%%%%%%%%%%%%%%%%%%%%%%%%%%%




%%% Codages

% \newcommand\gammaoc{\gamma_{[0,1]}}
% \newcommand\gammaonc{\gamma_{\N}}
\newcommand\gammaTMc{\gamma^{TM}_{[0,1]}}
\newcommand\gammaTMcatan{\gamma^{TM}_{(0,1)^2}}
\newcommand\gammaTMnc{\gamma^{TM}_{\N^3}}
\newcommand\gammaTMncd{\gamma^{TM}_{\N^2}}
% \newcommand\gammaTMe{\gamma^{TM}_{*}}
% \newcommand\gammas{\gamma_{sab}}
% \newcommand\gammaord{\gamma_{cantor}}


\newcommand\enccompact{\nu_{X}}
 \newcommand\encinfinite{\nu_\N}
 \newcommand\enc{\nu}
 
%%%%%%%%%%%%%%%%%%%%%%%%%%%%%%%%%%%%%%%%%%%%%%%%%
%
%
%                      DESSINS ET IMAGES
%
%
%%%%%%%%%%%%%%%%%%%%%%%%%%%%%%%%%%%%%%%%%%%%%%%%%


\newcommand\IMAGE[2]{ \begin{center} 
    \includegraphics[#1]{#2}
  \end{center}
  }


  
%%%% POUR FAIRE DESSIN.
\newenvironment{dessinpp}[1]{
\begin{tikzpicture}[baseline=(current bounding
      box.west),#1]
}
{\end{tikzpicture}}




\newenvironment{dessinp}[1]{
\begin{center}\begin{tikzpicture}[baseline=(current bounding
      box.north),#1]
}
{\end{tikzpicture}
\end{center}}


\newenvironment{dessin}{
\begin{center}\begin{tikzpicture}[baseline=(current bounding
      box.north)]
}
{\end{tikzpicture}
\end{center}}



\newenvironment{dessind}{
\begin{tikzpicture}[baseline=(current bounding
      box.north)]
}
{\end{tikzpicture}
}



%%%%%%%%%%%%%%%%%%%%%%%%%%%%%%%%%%%%%%%%%%%%%%%%%
%
%
%                      POUR DESSINS GPAC
%
%
%%%%%%%%%%%%%%%%%%%%%%%%%%%%%%%%%%%%%%%%%%%%%%%%%


%%%% MACROS DE BASE.


\newcommand\constant[2]
{
node (#1)  [shape=circle,draw] {#2};
\draw (#1.east) -- +(0.3,0) coordinate (#1-o) ; 
\draw (#1.west) -- +(-0.3,0) coordinate (#1-i) ; 
}

\newcommand\basicproduct[2]
{
node (#1) {#2};
\draw (#1) +(-0.4,0.4) -- +(0.4,0.4) -- +(0.8,0) -- +(0.4,-0.4) -- +(-0.4,-0.4) --
+(-0.4,0.4);
\draw (#1) ++(0.8,0) -- ++(+0.3,0) coordinate (#1-o) ;
\draw (#1) ++ (-0.4,0.2) -- +(-0.3,0) coordinate (#1-i1) ;
\draw (#1) ++(-0.4,-0.2) -- +(-0.3,0) coordinate (#1-i2) ;
}

\newcommand\product[1]{\basicproduct{#1}{$+\Pi$}}
\newcommand\negativeproduct[1]{\basicproduct{#1}{$-\Pi$}}


%%BEGIN INTERNAL
\newcommand\basicsummer[1]{
coordinate (#1) {};
\draw (#1) +(-0.5,0.6) -- +(0.5,0) -- +(-0.5,-0.6) -- +(-0.5,0.6);
\draw (#1) ++(0.5,0) -- +(+0.3,0) coordinate (#1-o) ;
}

\newcommand\basicintegrator[1]{
coordinate (#1) {};
\draw (#1) +(-0.5,0.6) -- +(0.5,0) -- +(-0.5,-0.6) -- +(-0.5,0.6);
\draw (#1) ++(0.5,0) -- +(+0.3,0) coordinate (#1-o) ;
\draw (#1) +(-0.5,-0.6) -| +(-0.8,0.6) -- +(-0.5,0.6) ;
 \draw (#1) ++(-0.65,0.6) -- +(0,+0.3) coordinate (#1-init) ;
}

\newcommand\splitfour[2]{
\draw (#2) ++ (#1,0.4) -- +(-0.3,0) coordinate (#2-i1) ;
\draw (#2) ++ (#1,0.15) -- +(-0.3,0) coordinate (#2-i2) ;
\draw (#2) ++ (#1,-0.15) -- +(-0.3,0) coordinate (#2-i3) ;
\draw (#2) ++ (#1,-0.4) -- +(-0.3,0) coordinate (#2-i4) ;
}

\newcommand\splitthree[2]{
\draw  (#2) ++(#1,0.3) -- +(-0.3,0) coordinate (#2-i1) ;
\draw  (#2) ++(#1,0) -- +(-0.3,0) coordinate (#2-i2) ;
\draw  (#2) ++(#1,-0.3) -- +(-0.3,0) coordinate (#2-i3) ;
}

\newcommand\splittwo[2]{
\draw  (#2) ++(#1,0.2) -- +(-0.3,0) coordinate (#2-i1) ;
\draw  (#2) ++(#1,-0.2) -- +(-0.3,0) coordinate (#2-i2) ;
}

\newcommand\splitone[2]{
\draw (#2) ++(#1,0) -- +(-0.3,0) coordinate (#2-i1) ;
}

%% END INTERNAL

\newcommand\summerfour[1]{
\basicsummer{#1}
\splitfour{-0.5}{#1}
}

\newcommand\summerthree[1]{
\basicsummer{#1}
\splitthree{-0.5}{#1}
}

\newcommand\summertwo[1]{
\basicsummer{#1}
\splittwo{-0.5}{#1}
}

\newcommand\summerone[1]{
\basicsummer{#1}
\splitone{-0.5}{#1}
}


\newcommand\integratorfour[1]{
\basicintegrator{#1}
\splitfour{-0.8}{#1}
<}


\newcommand\integratorthree[1]{
\basicintegrator{#1}
\splitthree{-0.8}{#1}
}


\newcommand\integratortwo[1]{
\basicintegrator{#1}
\splittwo{-0.8}{#1}
}


\newcommand\integratorone[1]{
\basicintegrator{#1}
\splitone{-0.8}{#1}
}

%%%%% FIN MACROS




%%%%%%%%%%%%%%%%%%%%%%%%%%%%%%%%%%%%%%%%%%%%%%%%%
%
%
%                      Trucs d'Amaury
%
%
%%%%%%%%%%%%%%%%%%%%%%%%%%%%%%%%%%%%%%%%%%%%%%%%%


%----------------
%--- ref (AP) ---
%----------------

% \newcommand{\defref}[1]{Definition~\ref{#1}}
\newcommand{\lemref}[1]{Lemma~\ref{#1}}
\newcommand{\thref}[1]{Theorem~\ref{#1}}
\newcommand{\amaurycorref}[1]{Corollary~\ref{#1}}
 \newcommand{\figref}[1]{Figure~\ref{#1}}
% \newcommand{\figrefmult}[1]{Figures~{#1}}
% \newcommand{\secref}[1]{Section~\ref{#1}}
% %\renewcommand{\algref}[1]{Algorithm~\ref{#1}}
 \newcommand{\propref}[1]{Proposition~\ref{#1}}
% \newcommand{\exref}[1]{Example~\ref{#1}}
\newcommand{\remref}[1]{Remark~\ref{#1}}


 %-------------------
%--- polynomials ---
%-------------------

\newcommand{\sigmap}[1]{{\Sigma{#1}}}
\newcommand{\degp}[1]{{\operatorname{deg}(#1)}}
\newcommand{\poly}{\operatorname{poly}}


 %----------------------
%--- misc functions ---
%----------------------

\newcommand{\sgn}[1]{\operatorname{sgn}(#1)}
\newcommand{\intp}{\operatorname{int}}
 \newcommand{\fracp}{\operatorname{frac}}
\newcommand{\fiter}[2]{#1^{[#2]}}
\newcommand{\frestrict}[2]{#1\restriction_{#2}}
\newcommand{\indicator}[1]{\mathds{1}_{#1}}
\newcommand{\idfun}{\operatorname{id}}
% \newcommand{\floor}[1]{\left\lfloor#1\right\rfloor}
 \newcommand{\ceil}[1]{\left\lceil#1\right\rceil}
\newcommand{\round}[1]{\left\lfloor#1\right\rceil}
\DeclareMathOperator\erf{erf}


%------------
%--- sets ---
%------------

\newcommand{\A}{\mathbb{A}}
%\newcommand{\D}{\mathbb{D}}
\newcommand{\Ns}{\mathbb{N}^*}
%\newcommand{\C}{\mathbb{C}}
%\newcommand{\K}{\mathbb{K}}
% \MAT{height}{[width]}{[field]}
\newcommand{\MAT}[3]{M_{#1\ifthenelse{\equal{#2}{}}{}{,#2}}\ifthenelse{\equal{#3}{}}{}{\left(#3\right)}}

%\newcommand{\Rcomp}{\R_{c}}
\newcommand{\Rp}{\R_{+}}
\newcommand{\Rs}{\R^{*}}
\newcommand{\Rps}{\R_{+}^{*}}

%\newcommand{\Analytic}{C^\omega}
%\newcommand{\PTIME}{\ensuremath{\operatorname{P}}}
%\newcommand{\EXPTIME}{\ensuremath{\operatorname{EXPTIME}}}
%\newcommand{\NP}{\ensuremath{\operatorname{NP}}}
%\newcommand{\NPTIME}{\NP}
\newcommand{\FP}{\ensuremath{\operatorname{FP}}}
%\newcommand{\PSPACE}{\ensuremath{\operatorname{PSPACE}}}
% \newcommand{\FPSPACE}{\ensuremath{\operatorname{FPSPACE}}}



%%%%%%%%%%%%%%%%%%%%%%%%%%%%%%%%%%%%%%%%%%%%%%%%%
%
%
%                      Trucs de Quentin
%
%
%%%%%%%%%%%%%%%%%%%%%%%%%%%%%%%%%%%%%%%%%%%%%%%%%


\newcommand{\fct}[5]{\ensuremath{#1 : \left\{\begin{array}{c c c}
#2 & \rightarrow & #3\\
#4 & \mapsto & #5
                                             \end{array}\right.}}

%Nouvelle fraction, plus jolie
\newcommand{\f}[2]{	
	\mathchoice%
		{\dfrac{#1}{#2}}
    	{\dfrac{#1}{#2}}
		{\frac{#1}{#2}}
		{\frac{#1}{#2}}
}


%%Ensembles et combinatoires

%Ensemble avec propriété à drotie \{x | blabla\}
\newcommand{\enstq}[2]{\left\{ \left.#1\ \vphantom{#2}\right|\ #2  \right\}}
%meme chose mais avec des crochets
\newcommand{\crotq}[2]{\left[ \left.#1\ \vphantom{#2}\right|\ #2  \right]}
%même chose mais avec des brackets
\newcommand{\bratq}[2]{\left\langle \left.#1\ \vphantom{#2}\right|\ #2  
  \right\rangle}

\newcommand{\intn}[2]{\ensuremath{
		\left\llbracket\, #1\ ;\ #2\,\right\rrbracket
              }}


%%%%%%%%%%%%%%%%%%%%%%%%%%%%%%%%%%%%%%%%%%%%%%%%%
%
%
%               ENVIRONNEMENTS THEOREMS & CO
%
%
%%%%%%%%%%%%%%%%%%%%%%%%%%%%%%%%%%%%%%%%%%%%%%%%%


%----------------
%--- pas dans lipics ---
%-----------------


\ifnotSLIDE{
% commente car bug Fevrier 2026: \ifnotTDEXAM{\spnewtheorem{solution}{Solution}{\bfseries}{\rmfamily}}
\spnewtheorem{openproblem}{Open Problem}{\bfseries}{\rmfamily}
\spnewtheorem{remarkolivier}{Commentaire d'Olivier: }{\bfseries}{\rmfamily}
\spnewtheorem{postulate}{Postulate}{\bfseries}{\rmfamily}
}

%----------------
%--- dans lipics ---
%-----------------
%
\newenvironment{exercice}[1]{\begin{exercise} \label{exo:#1-stt}}{\end{exercise}}
\newenvironment{exerciced}[1]{\begin{exercisee}\label{exo:#1-stt}}{\end{exercisee}}
\newenvironment{exercicec}[1]{\begin{exercise} \label{exo:#1-stt} (\LANGUE{solution on}{corrigé} page \pageref{exo:#1-cor})
}{\end{exercise}}
\newenvironment{exercicedc}[1]{\begin{exercisee} \label{exo:#1-stt} (\LANGUE{solution on}{corrigé} page \pageref{exo:#1-cor})
}{\end{exercisee}}

\newenvironment{slideproposition}{{\bf Proposition}}{}

\ifnotSLIDE{
\nolipics{
\nolncs{\spnewtheorem{theorem}{\LANGUE{Theorem}{Théorème}}{\bfseries}{\rmfamily}}
\nolncs{\spnewtheoremi{definition}{\LANGUE{Definition}{Définition}}{\bfseries}{\rmfamily}{theorem}}
\nolncs{\spnewtheoremi{lemma}{\LANGUE{Lemma}{Lemme}}{\bfseries}{\rmfamily}{theorem}}
\nolncs{\spnewtheoremi{corollary}{\LANGUE{Corollary}{Corollaire}}{\bfseries}{\rmfamily}{theorem}}
\nolncs{\spnewtheoremi{proposition}{Proposition}{\bfseries}{\rmfamily}{theorem}}
\nolncs{\spnewtheoremi{example}{\LANGUE{Example}{Exemple}}{\bfseries}{\rmfamily}{theorem}}
%\spnewtheorem{sketch}{{\bf \LANGUE{Sketch of proof}{Démonstration (principe)}}}{\bfseries}{\rmfamily}
\nolncs{\spnewtheoremi{remark}{\LANGUE{Remark}{Remarque}}{\bfseries}{\rmfamily}{theorem}}
\nolncs{\spnewtheorem{exercise}{\LANGUE{Exercise}{Exercice}}{\bfseries}{\rmfamily}}
\spnewtheorem{exercisee}{\LANGUE{*Exercise}{*Exercice}}{\bfseries}{\rmfamily}
\spnewtheorem{summary}{\LANGUE{Summary}{Résumé}}{\bfseries}{\rmfamily}
 %\spnewtheorem{exerciceenglish}{Exercise}{\bfseries}{\rmfamily}
 %
 \newenvironment{sketch}{{\bf \LANGUE{Sketch of proof}{Démonstration (principe)}}:}{\hfill $\Box$ }
\nolncs{\spnewtheorem{conjecture}{\LANGUE{Conjecture}{Conjecture}}{\bfseries}{\rmfamily}}
 %%
 \spnewtheorem{theoremaprouver}{\LANGUE{Theorem TO BE PROVED}{Théorème (A PROUVER)}}{\bfseries}{\rmfamily}
  \spnewtheorem{conjectureaprouver}{\LANGUE{Conjecture TO BE VERIFIED}{Conjecture (A VERIFIER)}}{\bfseries}{\rmfamily}
 }}
            
 \makeatletter
 \ifnotSLIDE{
  \ifnotTDEXAM{
\@ifundefined{proof}{
	\newenvironment{proof}{{\bf \LANGUE{Proof}{Démonstration}}:}{\hfill $\Box$ 
	}{}
	}}}
\makeatother

 
%%%%%%%%%%%%%%%%%%%%%%%%%%%%%%%%%%%%%%%%%%%%%%%%%
%
%
%                      Macros dont pas fier du tout
%
%
%%%%%%%%%%%%%%%%%%%%%%%%%%%%%%%%%%%%%%%%%%%%%%%%%



\newcommand\equations[1]{
\begin{array}{lll}
#1
\end{array}
}



\newcommand\systeme[1]{
\left\{
\begin{array}{lll} 
#1
\end{array}
\right.
}



%------------------
%--- shorthands (AP)---
%------------------
\newcommand{\mtt}[1]{\mathtt{#1}}
\newcommand{\ovl}[1]{\overline{#1}}
\newcommand{\panic}[1]{\textcolor{red}{#1}}
%\NewEnviron{epanic}{{\color{red}\BODY}}
\newcommand{\idest}{\emph{i.e.\thinspace}}


%------------------
%--- mes trucs  ---
%------------------


\newcommand\implique{\Rightarrow}
\newcommand\ssi{\Leftrightarrow}
\newcommand\ac[1]{\{#1\}}
\newcommand\zero{\vectorl{0}}
\newcommand\un{\vectorl{1}}

\newcommand\technique[1]{{\scriptsize #1}}


%% Preuve and Co

\newcommand\sensdirect{ Direction $\Rightarrow$. }
\newcommand\sensindirect{ Direction $\Leftarrow$. }




\newcommand\donnee{\mbox{<donnée>}}



% pour aller vite

\newcommand\gM{\mathfrak{M}}
\newcommand\gMM{(M,f_1,\cdots,f_u,r_1,\cdots,r_v)}
%
\newcommand\mygM{\mK}
\newcommand\mygMM{\left ( \K, \{op_i\}_{i\in I}, r_1 \ldots,
r_\ell, \zero,\un \right )}
\newcommand\myM{\K} 

\newcommand\gMMM{(M,f_1,\cdots,f_u,c_1,\cdots,c_k,r_1,\cdots,r_v)}
\newcommand\gMMMM{(f_1,\cdots,f_u,r_1,\cdots,r_v)}
\newcommand\gMME{(M,f_1,\cdots,f_u,f'_1,\cdots,f'_\ell,r_1,\cdots,r_v)}
\newcommand\gMMS{(f_1,\cdots,f_u,f'_1,\cdots,f'_\ell,r_1,\cdots,r_v)}

\newcommand\bool{\mathfrak{B}}
\newcommand\gbool{(\{0,1\},\zero,\un,=)}

\newcommand\osg{\overline{sg}}
\newcommand\moins{\mathrel{\dot{-}}}


%%% TRUC DE BARWISE


\newcommand\UM{\kU_{\kM}}
\newcommand\VV{\mathbb{V}}
% Hereditarily finite
\newcommand\IHF{\mathbb{I}HF}
\newcommand\IHP{\IHYP}
\newcommand\IHYP{\mathbb{I}HYP}

\providecommand\citep[1]{\cite{#1}}

%% Cofinali


% probas
% défini dans amsmath
\renewcommand\Pr{\cp{Pr}}
\newcommand\Var{\cp{Var}}


% -- mesures
%\newcommand\card{card}
\newcommand\taille{\LANGUE{size}{taille}}
\newcommand\entree{e}
\newcommand\profondeur{\LANGUE{depth}{profondeur}}


%%%%%%%%%%%%%%%%%%%%%%%%%%%%%%%%%%%%%%%%%%%%%%%%%%%%%%%%%%%%%%%%%%%%%%%%%%%%
%                 MACHINES DE TURING (graphique)
%%%%%%%%%%%%%%%%%%%%%%%%%%%%%%%%%%%%%%%%%%%%%%%%%%%%%%%%%%%%%%%%%%%%%%%%%%%%



% Configuration de machine de Turing
\newcommand\configTuring[3]{#1\textbf{#2}#3}


%%% DESSIN CROIX PORU AIDER  A DESSINER

\newcommand\croix{  +(-1,-1) -- +(1,1) +(-1,1) --
  +(1,-1) }

%% DESSIN ACCOLADE

\newcommand\ACCOLADE[3]{
%#1: coordonnées de départ
%#2: coordonnées arrivée
%#3: texte
\draw[decorate,decoration={brace}]  (#1) -- (#2)
 node[midway,above=0.1cm] {#3};
}

%%%


\def\lcase{3.5 ex}
\def\hcase{3.5ex}
\tikzstyle{case} = []
%[draw,minimum width = \lcase,minimum height = 3.6ex,baseline]
\newcommand\dcase{ +(-0.5*\lcase,-0.5*\hcase) rectangle +(0.5*\lcase,0.5*\hcase)}
%\newcommand\lettre[1]{\node [name=l,case,on chain=1] {#1};}


\newcommand\RUBANMACHINEDETURING[6]{
%#1: texte
%#2: position de la tête
%#3: nb de blancs à gauche
%#4: nb de cases totales
%#5: texte au dessus du ruban
%#6: position du coin
    \draw (#6) +(\lcase,4ex ) node {#5};
    \draw (#6) +(-0.6,0)  node  {\ldots};  
    \draw (#6) node[coordinate] (t) {};
    \foreach \x in {1,2,...,#3} {
         \draw (t) node [case]   {$\blanc$} \dcase;
         \draw (t) ++(\lcase,0) node [coordinate] (t) {};
       };
   \StrLen{#1}[\Resultat] % ATTENTION IL FAUT CETTE SYNTAXE POUR FAIRE
                         % EQUIVALENT D UN EDEF
  
\foreach \x in {1,...,\Resultat} 
             { 
        \draw (t) node [case]   {\StrChar{#1}{\x}} \dcase;
        \draw (t) ++(\lcase,0) node [coordinate] (t) {};
             }
\pgfmathtruncatemacro{\z}{#4 - \Resultat - #3}

  \foreach \x in {1,2,..., \z} {
        \draw (t) node [case]   {$\blanc$} \dcase;
        \draw (t) ++(\lcase,0) node [coordinate] (t) {};
}
%% Convention: graphiquement, la case numéro n se trouve à (n*\lcase,0)
%% en comptant les case à partir de $0$

   \draw  (t) +(0.1,0) node [name=r] {\ldots}; 
}

\newcommand\RUBANMACHINEDETURINGtexte[6]{
%#1: texte
%#2: position de la tête
%#3: nb de blancs à gauche
%#4: nb de cases totales (équivalent)
%#5: texte au dessus du ruban
%#6: position du coin
    \draw (#6) +(\lcase,4ex ) node {#5};
    \draw (#6) +(-0.6,0)  node  {\ldots};  
    \draw (#6) node[coordinate] (t) {};
   \foreach \x in {1,2,...,#3} {
        \draw (t) node [case]   {$\blanc$} \dcase;
        \draw (t) ++(\lcase,0) node [coordinate] (t) {};
      };
   \draw (#6) ++(-2*\lcase,0.5*\lcase)-- +(\lcase*#4,0);
   \draw (#6) ++(-2*\lcase,-0.5*\lcase) -- +(\lcase*#4,0);
   
    \draw (t) ++(-0.5*\lcase,0) node [case,anchor=west]  (te) {#1};
    \draw (te.east)  node [anchor=west,coordinate] (t) {};

  \draw (t) ++(0.2*\lcase,0) node [coordinate] (t) {};
   \foreach \x in {1,2,...,#3} {
        \draw (t) node [case]   {$\blanc$} \dcase;
        \draw (t) ++(\lcase,0) node [coordinate] (t) {};
      };

    \draw  (t) +(0.1,0) node [anchor=west] {\ldots}; 
}


\newcommand\RUBANMACHINEDETURINGTETE[6]{
%
% dessine la tete
%
% output: (tete)
%
\RUBANMACHINEDETURING{#1}{#2}{#3}{#4}{#5}{#6}
% TETE
     \draw (#6) + ( #2 * \lcase + #3  *\lcase,- 0.5*\lcase )
     node[coordinate] (k) {};
     \draw[->] (k) +(0,-0.5) -- (k);
     \draw (k) +(0,-0.5) node[coordinate] (tete) {};
}


\newcommand\RUBANMACHINEDETURINGTETEtexte[6]{
%
% dessine la tete
%
% output: (tete)
%
\RUBANMACHINEDETURINGtexte{#1}{#2}{#3}{#4}{#5}{#6}
% TETE
     \draw (#6) + ( #2 * \lcase + #3  *\lcase,- 0.5*\lcase )
     node[coordinate] (k) {};
     \draw[->] (k) +(0,-0.5) -- (k);
     \draw (k) +(0,-0.5) node[coordinate] (tete) {};
}


\newcommand\MACHINEDETURING[6]{
%#1: texte
%#2: position de la tête
%#3: état
%#4: nb de blancs à gauche
%#5: nb de cases totales
%#6: texte en haut du ruban
%\begin{tikzpicture}[node distance=-0.15mm]

\RUBANMACHINEDETURINGTETE{#1}{#2}{#4}{#5}{#6}{0,0}
%
%% ETAT:
     \draw (tete) node [name=etat,draw,circle,anchor=north]  {#3};
     \draw (etat) -- (tete);
%\end{tikzpicture}
}

\newcommand\MACHINEDETURINGtexte[6]{
%#1: texte
%#2: position de la tête
%#3: état
%#4: nb de blancs à gauche
%#5: nb de cases totales
%#6: texte en haut du ruban
%\begin{tikzpicture}[node distance=-0.15mm]

\RUBANMACHINEDETURINGTETEtexte{#1}{#2}{#4}{#5}{#6}{0,0}
%
%% ETAT:
     \draw (tete) node [name=etat,draw,circle,anchor=north]  {#3};
     \draw (etat) -- (tete);
%\end{tikzpicture}
}

\newcommand\MACHINEDETURINGTROISRUBANS[9]{
%#1: texte1
%#2: position de la tête 1
%#3: texte2
%#4: position de la tête 2
%#5: texte3
%#6: position de la tête 3
%#7: état
%#8: nb de blancs à gauche 1
%#9: nb de cases totales
%\begin{tikzpicture}[node distance=-0.15mm]

\RUBANMACHINEDETURINGTETE{#1}{#2}{#8}{#9}{\LANGUE{tape}{ruban} $1$}{0,0}
\draw (tete) node[coordinate] (tete1) {};
\draw (tete1) +(8,0) node[coordinate] (tete1b) {};
\draw (tete1) -| (tete1b) ;

\RUBANMACHINEDETURINGTETE{#3}{#4}{#8}{#9}{\LANGUE{tape}{ruban} $2$}{0,-2}
\draw (tete) node[coordinate] (tete2) {};
\draw (tete2) +(-8.5,0) node[coordinate] (tete2b) {};
\draw (tete2) -| (tete2b) ;

\RUBANMACHINEDETURINGTETE{#5}{#6}{#8}{#9}{\LANGUE{tape}{ruban} $3$}{0,-4}
\draw (tete) node[coordinate] (tete3) {};

%% ETAT:
     \draw (tete3) node [name=etat,draw,circle,anchor=north]  {#7};
     \draw (etat) -- (tete3);
    \draw  (etat) -|  (tete2b);
\draw  (etat) -|  (tete1b);
%\end{tikzpicture}
}

\newcommand\MACHINEDETURINGTROISRUBANSENGLISH[9]{
%#1: texte1
%#2: position de la tête 1
%#3: texte2
%#4: position de la tête 2
%#5: texte3
%#6: position de la tête 3
%#7: état
%#8: nb de blancs à gauche 1
%#9: nb de cases totales
%\begin{tikzpicture}[node distance=-0.15mm]

\RUBANMACHINEDETURINGTETE{#1}{#2}{#8}{#9}{tape $1$}{0,0}
\draw (tete) node[coordinate] (tete1) {};
\draw (tete1) +(8,0) node[coordinate] (tete1b) {};
\draw (tete1) -| (tete1b) ;

\RUBANMACHINEDETURINGTETE{#3}{#4}{#8}{#9}{tape $2$}{0,-2}
\draw (tete) node[coordinate] (tete2) {};
\draw (tete2) +(-8.5,0) node[coordinate] (tete2b) {};
\draw (tete2) -| (tete2b) ;

\RUBANMACHINEDETURINGTETE{#5}{#6}{#8}{#9}{tape $3$}{0,-4}
\draw (tete) node[coordinate] (tete3) {};

%% ETAT:
     \draw (tete3) node [name=etat,draw,circle,anchor=north]  {#7};
     \draw (etat) -- (tete3);
    \draw  (etat) -|  (tete2b);
\draw  (etat) -|  (tete1b);
%\end{tikzpicture}
}



\newcommand\MACHINEDETURINGTROISRUBANStexte[9]{
%#1: texte1
%#2: position de la tête 1
%#3: texte2
%#4: position de la tête 2
%#5: texte3
%#6: position de la tête 3
%#7: état
%#8: nb de blancs à gauche 1
%#9: nb de cases totales
% \begin{tikzpicture}[node distance=-0.15mm]

\RUBANMACHINEDETURINGTETEtexte{#1}{#2}{#8}{#9}{\LANGUE{tape}{ruban} $1$}{0,0}
\draw (tete) node[coordinate] (tete1) {};
\draw (tete1) +(8,0) node[coordinate] (tete1b) {};
\draw (tete1) -| (tete1b) ;

\RUBANMACHINEDETURINGTETEtexte{#3}{#4}{#8}{#9}{\LANGUE{tape}{ruban} $2$}{0,-2}
\draw (tete) node[coordinate] (tete2) {};
\draw (tete2) +(-8.5,0) node[coordinate] (tete2b) {};
\draw (tete2) -| (tete2b) ;

\RUBANMACHINEDETURINGTETEtexte{#5}{#6}{#8}{#9}{\LANGUE{tape}{ruban} $3$}{0,-4}
\draw (tete) node[coordinate] (tete3) {};

%% ETAT:
     \draw (tete3) node [name=etat,draw,circle,anchor=north]  {#7};
     \draw (etat) -- (tete3);
    \draw  (etat) -|  (tete2b);
\draw  (etat) -|  (tete1b);
%\end{tikzpicture}
}


\newcommand\MACHINEDETURINGDEUXRUBANStexte[9]{
%#1: texte1  
%#2: position de la tête 1 
%#3: texte2 inutilisée
%#4: position de la tête 2 inutilisée
%#5: texte3
%#6: position de la tête 3
%#7: état
%#8: nb de blancs à gauche 1
%#9: nb de cases totales
% \begin{tikzpicture}[node distance=-0.15mm]

\RUBANMACHINEDETURINGTETEtexte{#1}{#2}{#8}{#9}{\LANGUE{tape}{ruban} $1$}{0,0}
\draw (tete) node[coordinate] (tete1) {};
\draw (tete1) +(8,0) node[coordinate] (tete1b) {};
\draw (tete1) -| (tete1b) ;

% \RUBANMACHINEDETURINGTETEtexte{#3}{#4}{#8}{#9}{\LANGUE{tape}{ruban} $2$}{0,-1.5}
% \draw (tete) node[coordinate] (tete2) {};
% \draw (tete2) +(-5.5,0) node[coordinate] (tete2b) {};
% \draw (tete2) -| (tete2b) ;

\RUBANMACHINEDETURINGTETEtexte{#5}{#6}{#8}{#9}{\LANGUE{tape}{ruban} $2$}{0,-1.5}
\draw (tete) node[coordinate] (tete3) {};

%% ETAT:
     \draw (tete3) node [name=etat,draw,circle,anchor=north]  {#7};
     \draw (etat) -- (tete3);
%   \draw  (etat) -|  (tete2b);
\draw  (etat) -|  (tete1b);
%\end{tikzpicture}
}


\def\argdix{2}
\def\argonze{25}
\newcommand\MACHINEDETURINGQUATRERUBANStextep[9]{
%#1: texte1
%#2: position de la tête 1
%#3: texte2
%#4: position de la tête 2
%#5: texte3
%#6: position de la tête 3
%#7: texte4
%#8: position de la tête 3
%#9: état
%#10: nb de blancs à gauche 1
%#11: nb de cases totales
% \begin{tikzpicture}[node distance=-0.15mm]

%% BUG sur #10 devenu \argdix
%% BUG SUR #11 devenu \argonze

\RUBANMACHINEDETURINGTETEtexte{#1}{#2}{\argdix}{\argonze}{\LANGUE{tape}{ruban} $1$}{0,0}
% \draw (tete) node[coordinate] (tete1) {};
% \draw (tete1) +(8,0) node[coordinate] (tete1b) {};
% \draw (tete1) -| (tete1b) ;

\RUBANMACHINEDETURINGTETEtexte{#3}{#4}{\argdix}{\argonze}{\LANGUE{tape}{ruban} $2$}{0,-1.5}
% \draw (tete) node[coordinate] (tete2) {};
% \draw (tete2) +(-5.5,0) node[coordinate] (tete2b) {};
% \draw (tete2) -| (tete2b) ;

\RUBANMACHINEDETURINGTETEtexte{#5}{#6}{\argdix}{\argonze}{\LANGUE{tape}{ruban} $3$}{0,-3.2}
\draw (tete) node[coordinate] (tete3) {};

\RUBANMACHINEDETURINGTETEtexte{#7}{#8}{\argdix}{\argonze}{\LANGUE{tape}{ruban} $4$}{0,-4.9}
\draw (tete) node[coordinate] (tete4) {};

%% ETAT:
     \draw (tete4) node [name=etat,draw,circle,anchor=north]  {#9};
%      \draw (etat) -- (tete3);
%     \draw  (etat) -|  (tete2b);
% \draw  (etat) -|  (tete1b);
%\end{tikzpicture}
}


%%%%%%%%%%%%%%%%%%%%%%%%%%%%%%%%%%%%%%%%%%%%%%%%%%%%%%%%%%%%%%%%%%%%%%%%%%%%
%                 MACHINES DE TURING (automate)
%%%%%%%%%%%%%%%%%%%%%%%%%%%%%%%%%%%%%%%%%%%%%%%%%%%%%%%%%%%%%%%%%%%%%%%%%%%%



\newcommand\dec{1ex}
\newcommand\myACCOLADE[4]{
\ACCOLADE{#1*\lcase+#4*\lcase-0.5*\lcase,0.5*\lcase+\dec}{#1*\lcase+#4*\lcase-0.5*\lcase+#3*\lcase,0.5*\lcase+\dec}{#2}
}


\newcommand\M{\Sigma}


\newcommand\sommet[2]{\node[state] (#1) #2 {$#1$};}
\newcommand\sommetinitial[2]{\node[state,initial] (#1) #2 {$#1$};}
\newcommand\sommetfinal[2]{\node[state,accepting] (#1) #2 {$#1$};}
\newcommand\gtransition[5]{\draw (#1) edge node {#2/#4 $#5$} (#3);}
\newcommand\gtransitionba[5]{\draw (#1) [loop above] edge node {#2/#4
    $#5$} (#3);}
\newcommand\gtransitionbb[5]{\draw (#1) [loop below] edge node {#2/#4
    $#5$} (#3);}
\newcommand\gtransitionbleft[5]{\draw (#1) [loop left] edge node {#2/#4
    $#5$} (#3);}
\newcommand\gtransitionbright[5]{\draw (#1) [loop right] edge node {#2/#4
    $#5$} (#3);}
\newcommand\gtransitionbl[5]{\draw (#1) [bend left] edge node {#2/#4
    $#5$} (#3);}
\newcommand\gtransitionbr[5]{\draw (#1) [bend right] edge node {#2/#4
    $#5$} (#3);}
\newcommand\gtransitionbadouble[9]{\draw (#1) [loop above] edge node {
\begin{tabular}{l}
#2/#4 $#5$ \\
  #7/#8 $#9$ \\
\end{tabular}
} (#3);}

\newcommand\SCALE{}

\newcommand\dessinmtp[2]{
\begin{center}
\begin{tikzpicture}[->,>=stealth',shorten >=1pt,auto,node distance=2.3cm,semithick,#2] 
#1
\end{tikzpicture}
\end{center}
}

\newcommand\dessinmt[1]{\dessinmtp{#1}{scale=1}}


\newcommand\ANIMATIONMT[1]{
\begin{overprint}
\begin{dessinp}{scale=0.52}
\foreach \av/\q/\b  [count=\xi]
in  
{#1}
{ 
\only<\xi| handout 0>{ 
 \StrLen{\av}[\Resintm] % ATTENTION IL FAUT CETTE SYNTAXE POUR FAIRE
                         % EQUIVALENT D UN EDEF
 \MACHINEDETURING{\av \b}{\Resintm}{$\q$}{5}{20}{}
}}
 \end{dessinp}
\end{overprint}
}

\newcommand\ANIMATIONMTd[1]{
\begin{overprint}
\begin{dessinp}{scale=0.52}
\foreach \av/\q/\b  [count=\xi]
in  
{#1}
{ 
\only<\xi| handout 0>{ 
 \StrLen{\av}[\Resintm] % ATTENTION IL FAUT CETTE SYNTAXE POUR FAIRE
                         % EQUIVALENT D UN EDEF
 \hspace{-1cm}\MACHINEDETURING{\av \b}{\Resintm}{$\q$}{5}{33}{}
}}
 \end{dessinp}
\end{overprint}
}

\def\MAXDET{20}

\newcommand\DIAGRAMMEESPACETEMPS[1]{
\draw node[coordinate] (tt) {};
\foreach \av/\q/\b  in  
{#1}
{ 

\RUBANMACHINEDETURING{\av{$\q$}\b}{0}{4}{\MAXDET}{}{tt};
\draw (tt) +(0,-\hcase) node[coordinate] (tt) {};
}
%\foreach \x in {0,...,19} {\draw (tt) + (\x*\lcase,0) node {$\vdots$};}
}



\newcommand\progun{
/q_0/0011, % initial
X/q_1/011,X0/q_1/11,X/q_2/0Y1,/q_2/X0Y1,
X/q_0/0Y1,XX/q_1/Y1,XXY/q_1/1, XX/q_2/YY, X/q_2/XYY,
XX/q_0/YY,XXY/q_3/Y,XXYY/q_3/B,XXYYB/q_4/B
}

\newcommand\DIAGRAMMEESPACETEMPSun{
\DIAGRAMMEESPACETEMPS
%% RECOPIE progun
{/q_0/0011, % initial
X/q_1/011,X0/q_1/11,X/q_2/0Y1,/q_2/X0Y1,
X/q_0/0Y1,XX/q_1/Y1,XXY/q_1/1, XX/q_2/YY, X/q_2/XYY,
XX/q_0/YY,XXY/q_3/Y,XXYY/q_3/B,XXYYB/q_4/B}
%% FIN RECOPIE
}

\newcommand\progdeux{
/q_0/0010, X/q_1/010, X0/q_1/10, X/q_2/0Y0,/q_2/X0Y0,
X/q_0/0Y0,XX/q_1/Y0,XXY/q_1/0,XXY0/q_1/%B,
}


\newcommand\DIAGRAMMEESPACETEMPSdeux{
\DIAGRAMMEESPACETEMPS
{
%% RECOPIE progdeux
/q_0/0010, X/q_1/010, X0/q_1/10, X/q_2/0Y0,/q_2/X0Y0,
X/q_0/0Y0,XX/q_1/Y0,XXY/q_1/0,XXY0/q_1/%B,
%% FIN RECOPIE
}
}

\newcommand\DIAGRAMMEESPACETEMPSdeuxvariante{
\def\memolcase{\lcase}
\def\memoMAXDET{20}
\def\MAXDET{15}
\def\lcase{6 ex}
\DIAGRAMMEESPACETEMPS
{
%% RECOPIE progdeux
/q_00/010, X/q_10/10, X0/q_11/0, X/q_20/Y0,/q_2X/0Y0,
X/q_00/Y0,XX/q_1Y/0,XXY/q_10/,XXY0/q_1B/%B,
%% FIN RECOPIE
}
\def\lcase{\memolcase}
\def\MAXDET{\memoMAXDET}
}


\newcommand\ANIMATIONun{
\ANIMATIONMT
%% RECOPIE progun
{/q_0/0011, % initial
X/q_1/011,X0/q_1/11,X/q_2/0Y1,/q_2/X0Y1,
X/q_0/0Y1,XX/q_1/Y1,XXY/q_1/1, XX/q_2/YY, X/q_2/XYY,
XX/q_0/YY,XXY/q_3/Y,XXYY/q_3/B,XXYYB/q_4/B}
%% FIN RECOPIE
}



\newcommand\DESSINALGO[3]{
% argument 1= ou
% argument 2= texte 
% argument 3= nom du noeud ex m
\draw node[draw,#1,minimum size=1cm] (#3) {#2};

% pour fitting box
\draw (#3.west) +(-0.2cm,0) node (#3c2) {};
\draw (#3.north) +(0,+0.2cm) node (#3c3) {};
\draw (#3.south) +(0,-0.2cm) node (#3c4) {};
}

\newcommand\DESSINALGOa[3]{
\DESSINALGO{#1}{#2}{#3}
\draw (#3.east) +(-0.1cm,0.2cm) node(#3a) {};
\draw[->] (#3a) -- ++(0.5cm,0) node (#3aa) {};
\draw (#3aa) node[anchor=west] (#3aaa) {\LANGUE{Accept}{Accepte}};
\draw (#3aaa.east) node(#3aaaa) {};
}

\newcommand\DESSINALGOe[4]{
\DESSINALGO{#1}{#2}{#3}
\draw (#3.east) +(-0.1cm,0.2cm) node(#3a) {};
\draw[->] (#3a) -- ++(0.5cm,0) node (#3aa) {};
\draw (#3aa) node[anchor=west] (#3aaa) {#4};
\draw (#3aaa.east) node(#3aaaa) {};
}


\newcommand\DESSINALGOar[3]{
\DESSINALGOa{#1}{#2}{#3}
\draw (#3.east) +(-0.1cm,-0.2cm) node(#3r) {};
\draw[->] (#3r) -- ++(0.5cm,0) node (#3rr) {};
\draw (#3rr) node[anchor=west] (#3rrr) {\LANGUE{Reject}{Refuse}};
\draw (#3rrr.east) node(#3rrrr) {};
}


\newcommand\WINDOWarg[7] {
%#1: position
%#2..7: contenu du tableau
\draw (#1) node[case]  {#2} \dcase;
\draw (#1) ++ (\lcase,0) node[case] {#3} \dcase;
\draw (#1) ++ (2*\lcase,0) node[case] {#4} \dcase;
\draw (#1) ++ (0,-\lcase) node[case] {#5} \dcase;
\draw (#1) ++ (\lcase,-\lcase) node[case] {#6} \dcase;
\draw (#1) ++ (2*\lcase,-\lcase) node[case] {#7} \dcase;
}

\newcommand\WINDOW[3]{
%#1: positioin
%#2: contenu
%#3: étiquette
\WINDOWarg{#1}
{\StrChar{#2}{1}}{\StrChar{#2}{2}}{\StrChar{#2}{3}}
{\StrChar{#2}{4}}{\StrChar{#2}{5}}{\StrChar{#2}{6}}
\draw (#1) +(-\lcase*1.2,-\lcase*0.5) node {#3};
}





%%%%%%%%%%%%%%%%%%%%%%%%%%%%%%%%%%%%%%%%%%%%%%%%%%%%%%%%%%%%%%%%%%%%%%%%%%%%
%                 LISTINGS
%%%%%%%%%%%%%%%%%%%%%%%%%%%%%%%%%%%%%%%%%%%%%%%%%%%%%%%%%%%%%%%%%%%%%%%%%%%%


% LISTINGS
\RequirePackage{listings}
 \lstset{language=Java,
 mathescape=true,
 showspaces=false,
 showstringspaces=false,
 frame=single,
morekeywords={recursive,then,div,function,integer,read,out,else,par,endpar,seq,endseq,and,not,read,write,repeat,until,undef,let,to,endfor,endif,endwhile,guess,in,parallel},
% basicstyle=\ttfamily\small,
% basewidth={1.1ex},
% keywordstyle=\bf\color{blue},
% %stringstyle=\bf\ttfamily\color{green},
% commentstyle=\bfseries\color{magenta},
 literate={:=}{{$\gets$ }}1 {<=}{{$\leq$}}2  {>=}{{$\geq$}}2
 {!=}{{$\neq$}}1 
 }
% \lstset{escapechar=\%}
\let\lst\lstinline
\def\beginlisting{\begin{lstlisting}}
\def\endlisting{\end{lstlisting}}  




\newenvironment{algo}[1]{
  \lstset{escapechar=\@}
  \def\input{$#1$}
  \def\out{out}
  \def\BEGIN{\textbf{read} \input }
 \def\BEGINN{\textbf{repeat} }
  \def\END{\textbf{until out!=undef}}
  \def\ENDD{\textbf{write \out}}
  \def\ENDT{\textbf{until $\ctl$ = $q_a$ or $\ctl$ = $q_r$}}
  \def\ENDR{\textbf{until $\ctl$ = $q_a$}}
  \def\ENDP{\textbf{until $\ctl$ = $q_a$}}
}{}
\newenvironment{algon}[1]{
  \lstset{numbers=left}
  \begin{algo}{#1}
}{\end{algo}}
\newcommand\FSM[3]{FSM(#1,\lst{#2},#3)}
\newcommand\FSMi[3]{FSM(#1,{#2},#3)}




%%%%%%%%%%%%%%%%%%%%%%%%%%%%%%%%%%%%%%%%%%%%%%%%%
%
%
%                      TRADUCTION EN COURS
%
%%%%%%%%%%%%%%%%%%%%%%%%%%%%%%%%%%%%%%%%%%%%%%%%%

\newcommand\ENTREPOT[1]{\NDLR{ENTREPOT ICI: #1}}
\newcommand\scite[2][]{\cite[#1]{#2}}
\newcommand\nospellcheck[1]{{#1}}
\newcommand\nd{nd}
\newcommand\NL{
 
}

\newenvironment{FRANCAIS}{
\small
\color{orange}
}
{\normalsize
\color{black}
}


\newenvironment{SEMIFRANCAIS}{
\small
\color{violet}
}
{\normalsize
\color{black}
}

\newcommand\QFRANCAIS[1]{
\begin{FRANCAIS}
#1
\end{FRANCAIS}
}


%%%%%%%%%%%%%%%%%%%%%%%%%%%%%%%%%%%%%%%%%%%%%%%%%
%
%
%                      Wikipedia
%
%%%%%%%%%%%%%%%%%%%%%%%%%%%%%%%%%%%%%%%%%%%%%%%%%

\providecommand\tightlist{}
\providecommand\Complex{\C}
\newenvironment{myalgie}{\begin{array}{lll}}{\end{array}}

\newcommand\nl{\ \\ }




%%%%%%%%%%%%%%%%%%%%%%%%%%%%%%%%%%%%%%%%%%%%%%%%%
%
%
%                      Pour citations précises
%
%%%%%%%%%%%%%%%%%%%%%%%%%%%%%%%%%%%%%%%%%%%%%%%%%



\newcommand\citeprecis[2]{{\cite[#1]{#2}}}
\newcommand\citeprecisinv[2]{\citeprecis{#2}{#1}}

\usepackage{csquotes}



%%%%%%%%%%%%%%%%%%%%%%%%%%%%%%%%%%%%%%%%%%%%%%%%%
%
%
%                      Travail sur couleurs
%
%   toutes les réponses pointent dans le même sens, et le schéma correspond exactement 
%. au profil perceptif d’un protanope** (déficit des cônes L / longueurs d’onde longues)
%  MAIS SEVERE
%
%%%%%%%%%%%%%%%%%%%%%%%%%%%%%%%%%%%%%%%%%%%%%%%%%


% ================================
% GROUPE 1 — Verts sombres
% ================================
\definecolor{deepgreen}{RGB}{0,160,40}
\definecolor{forestgreen}{RGB}{0,150,70}
\definecolor{mossgreen}{RGB}{0,140,90}
\definecolor{greenwave}{RGB}{0,150,90}
\definecolor{forestmist}{RGB}{0,160,90}
\definecolor{darkteal}{RGB}{0,130,70}

% ================================
% GROUPE 2 — Teal / verts bleutés
% ================================
\definecolor{seagreen}{RGB}{0,130,100}
\definecolor{freshteal}{RGB}{0,140,130}
\definecolor{greenteal}{RGB}{0,140,120}
\definecolor{blueteal}{RGB}{0,150,120}      % ancien "teal" (doublon corrigé)
\definecolor{lagoon}{RGB}{0,180,120}
\definecolor{greenmist}{RGB}{0,170,110}

% ================================
% GROUPE 3 — Turquoises / Aqua
% ================================
\definecolor{springgreen}{RGB}{0,190,100}
\definecolor{mintgreen}{RGB}{10,200,150}
\definecolor{oceanfoam}{RGB}{0,200,150}
\definecolor{aquagreen}{RGB}{30,210,170}
\definecolor{coastalaqua}{RGB}{40,210,180}
\definecolor{softaqua}{RGB}{40,210,190}

% ================================
% GROUPE 4 — Cyan / Aqua clairs
% ================================
\definecolor{emeraldteal}{RGB}{0,180,140}
\definecolor{deepcyan}{RGB}{0,180,160}
\definecolor{lightcyan}{RGB}{0,200,160}
\definecolor{aquamarine}{RGB}{0,200,160}
\definecolor{lightaqua}{RGB}{60,230,200}
\definecolor{mistblue}{RGB}{80,240,210}

% ================================
% GROUPE 5 — Bleus clairs
% ================================
\definecolor{paleaqua}{RGB}{80,200,190}
\definecolor{glacier}{RGB}{90,150,200}
\definecolor{skyblueA}{RGB}{60,150,220}      % ancien "skyblue"
\definecolor{lightsky}{RGB}{90,170,230}
\definecolor{iceblue}{RGB}{120,190,240}
\definecolor{frost}{RGB}{150,210,250}

% ================================
% GROUPE 6 — Bleus moyens
% ================================
\definecolor{steelblue}{RGB}{30,100,160}
\definecolor{cadetblue}{RGB}{70,120,180}
\definecolor{coolblue}{RGB}{110,170,220}
\definecolor{morningblue}{RGB}{140,200,240}
\definecolor{blueiceA}{RGB}{40,150,210}       % ancien "blueice" (doublon corrigé)
\definecolor{blueair}{RGB}{60,170,225}

% ================================
% GROUPE 7 — Azurs / Bleu saturé
% ================================
\definecolor{tealblue}{RGB}{0,110,130}
\definecolor{deepbluecyan}{RGB}{0,110,170}
\definecolor{azur}{RGB}{0,70,180}
\definecolor{brightazure}{RGB}{20,130,200}
\definecolor{lightazure}{RGB}{80,180,235}
\definecolor{skylight}{RGB}{100,200,245}

% ================================
% GROUPE 8 — Bleus sombres / violets
% ================================
\definecolor{nightblue}{RGB}{40,0,110}
\definecolor{slateblue}{RGB}{20,40,120}
\definecolor{violetA}{RGB}{60,50,160}         % ancien "violet" (distinct de violetB ci-dessous)
\definecolor{blueviolet}{RGB}{80,60,180}
\definecolor{purpleA}{RGB}{100,70,200}        % ancien "purple"
\definecolor{lightpurple}{RGB}{130,90,210}

% ================================
% GROUPE 9 — Violets saturés / indigos
% ================================
\definecolor{lilac}{RGB}{160,110,220}
\definecolor{indigoA}{RGB}{50,0,130}          % ancien "indigo"
\definecolor{deepindigo}{RGB}{70,10,150}
\definecolor{royalpurple}{RGB}{90,20,170}

% ================================
% GROUPE 10 — Fermeture cercle perceptif
% ================================
\definecolor{coolmint}{RGB}{0,130,70}
\definecolor{deepsea}{RGB}{0,130,70}

% ================================
% PALETTE DEUTÉRANOPE — noms uniques
% ================================
\definecolor{crimson}{RGB}{200,30,60}
\definecolor{brickred}{RGB}{180,40,40}
\definecolor{coral}{RGB}{220,90,70}
\definecolor{salmon}{RGB}{240,130,120}

\definecolor{magenta}{RGB}{200,40,150}
\definecolor{fuchsia}{RGB}{220,60,170}
\definecolor{violetB}{RGB}{140,70,190}        % ancien "violet"
\definecolor{indigoB}{RGB}{90,50,180}         % ancien "indigo"

\definecolor{royalblueA}{RGB}{50,80,200}      % ancien "royalblue"
\definecolor{azure}{RGB}{30,120,230}
\definecolor{skyblueB}{RGB}{80,160,230}       % second "skyblue"
\definecolor{tealB}{RGB}{0,150,170}           % second "teal"

\definecolor{marine}{RGB}{0,80,140}




%% Alternative

% === Générateur automatique fill/text pour protanopes ===
% Usage : \ProtaColorPair{mistblue}
% Suppose que mistblue est déjà défini par \definecolor{mistblue}{RGB}{...}

\newcommand{\ProtaColorPair}[1]{%
  % fill = couleur originale
  \colorlet{#1fill}{#1}%
  % text = version assombrie idéalement lisible
  \colorlet{#1text}{#1!80!black}%
}

%Utilisation
%\ProtaColorPair{mistblue}

%\textcolor{mistbluefill}{Couleur claire} \quad
%\textcolor{mistbluetext}{Couleur pour texte}

	\foreach \i/\name in {
	  1/deepgreen,
	  2/greenwave,
	  3/lagoon,
	  4/oceanfoam,
	  5/softaqua,
	  6/mistblue,
 	 7/lightsky,
	  8/blueiceA,   % <-- corrigé
	  9/deepcyan,
	  10/azur,
	  11/blueviolet,
	  12/royalpurple}
	  {\ProtaColorPair{\name}}
	  


\newcommand\ROUEPROTANOPEDOUZE{
	% Roue chromatique protanope — 12 couleurs
	\begin{tikzpicture}[scale=3]

	\foreach \i/\name/\nametext in {
	  1/deepgreen/deepgreen,
	  2/greenwave/greenwave,
	  3/lagoon/lagoontext,
	  4/oceanfoam/oceanfoam,
	  5/softaqua/softaquatext,
	  6/mistblue/mistbluetext,
 	 7/lightsky/lightskytext,
	  8/blueiceA/blueiceA,   % <-- corrigé
	  9/deepcyan/blueiceA,
	  10/azur/azur,
	  11/blueviolet/blueviolet,
	  12/royalpurple/royalpurple}
	{
	  % secteur coloré
 	 \filldraw[fill=\name, draw=black, line width=0.2pt]
	    ({90-30*\i}:1)
	    -- ({90-30*(\i-1)}:1)
	    -- (0,0) -- cycle;

 	 % label dans la bonne couleur
	  \node at ({90-30*(\i-0.5)}:1.25)
	    {\textcolor{\nametext}{\small \nametext}};
	}

	\end{tikzpicture}
}


\newcommand\ROUEPROTANOPEVINGT{
\begin{tikzpicture}[scale=3]

\foreach \i/\name/\nametext in {
  1/deepgreen/deepgreen,
  2/greenwave/greenwave,
  3/lagoon/lagoon,
  4/oceanfoam/oceanfoam,
  5/softaqua/softaquatext,
  6/mistblue/mistbluetext,
  7/lightsky/lightsky,
  8/blueiceA/blueiceA,      % corrigé
  9/deepcyan/deepcyan,
  10/azur/azur,
  11/brightazure/brightazure,
  12/skyblueA/skyblueA,     % corrigé
  13/coolblue/coolblue,
  14/glacier/glacier,
  15/marine/marine,
  16/violetA/violetA,      % corrigé
  17/blueviolet/blueviolet,
  18/purpleA/purpleA,      % corrigé
  19/royalpurple/royalpurple,
  20/indigoA/indigoA}      % corrigé
{
  % Secteur coloré
  \filldraw[fill=\name, draw=black, line width=0.2pt]
    ({90-18*\i}:1)
    -- ({90-18*(\i-1)}:1)
    -- (0,0) -- cycle;

  % Étiquette
  \node at ({90-18*(\i-0.5)}:1.25)
    {\textcolor{\nametext}{\small \nametext}};
}

\end{tikzpicture}
}

\newcommand\ROUEDEUTERANOPEDOUZE{
% Roue chromatique deutéranope — 12 couleurs
\begin{tikzpicture}[scale=3]

\foreach \i/\name in {
  1/crimson,
  2/brickred,
  3/coral,
  4/salmon,
  5/magenta,
  6/fuchsia,
  7/violetB,       % corrigé
  8/indigoB,       % corrigé
  9/royalblueA,    % corrigé
  10/azure,
  11/skyblueB,     % corrigé
  12/tealB}        % corrigé
{
  % Secteur coloré
  \filldraw[fill=\name, draw=black, line width=0.2pt]
    ({90-30*\i}:1)
    -- ({90-30*(\i-1)}:1)
    -- (0,0) -- cycle;

  % Label dans la couleur
  \node at ({90-30*(\i-0.5)}:1.25)
    {\textcolor{\name}{\small \name}};
}

\end{tikzpicture}
}



% === Macro intelligente pour foncer une couleur de façon lisible pour protanope ===
% Usage : \protadarken{mistblue}{40}  -> crée mistblue@dark
% Le 40 signifie : 40% de mélange avec du noir

\newcommand{\DarkenForText}[2]{%
  % #1 = nom de la couleur
  % #2 = pourcentage (40 conseillé si tu ne sais pas)
  %
  % On crée une nouvelle couleur #1@dark
  \colorlet{#1text}{black!#2!#1}%
}





%%%%

% Groupe 1
\DarkenForText{deepgreen}{20}
\DarkenForText{forestgreen}{20}
\DarkenForText{mossgreen}{20}
\DarkenForText{greenwave}{20}
\DarkenForText{forestmist}{20}
\DarkenForText{darkteal}{20}

% Groupe 2
\DarkenForText{seagreen}{20}
\DarkenForText{freshteal}{20}
\DarkenForText{greenteal}{20}
\DarkenForText{blueteal}{20}
\DarkenForText{lagoon}{20}
\DarkenForText{greenmist}{20}

% Groupe 3
\DarkenForText{springgreen}{20}
\DarkenForText{mintgreen}{20}
\DarkenForText{oceanfoam}{20}
\DarkenForText{aquagreen}{20}
\DarkenForText{coastalaqua}{20}
\DarkenForText{softaqua}{20}

% Groupe 4
\DarkenForText{emeraldteal}{20}
\DarkenForText{deepcyan}{20}
\DarkenForText{lightcyan}{20}
\DarkenForText{aquamarine}{20}
\DarkenForText{lightaqua}{20}
\DarkenForText{mistblue}{30}

% Groupe 5
\DarkenForText{paleaqua}{20}
\DarkenForText{glacier}{20}
\DarkenForText{skyblueA}{20}
\DarkenForText{lightsky}{20}
\DarkenForText{iceblue}{20}
\DarkenForText{frost}{20}

% Groupe 6
\DarkenForText{steelblue}{20}
\DarkenForText{cadetblue}{20}
\DarkenForText{coolblue}{20}
\DarkenForText{morningblue}{20}
\DarkenForText{blueiceA}{20}
\DarkenForText{blueair}{20}

% Groupe 7
\DarkenForText{tealblue}{20}
\DarkenForText{deepbluecyan}{20}
\DarkenForText{azur}{20}
\DarkenForText{brightazure}{20}
\DarkenForText{lightazure}{20}
\DarkenForText{skylight}{20}

% Groupe 8
\DarkenForText{nightblue}{20}
\DarkenForText{slateblue}{20}
\DarkenForText{violetA}{20}
\DarkenForText{blueviolet}{20}
\DarkenForText{purpleA}{20}
\DarkenForText{lightpurple}{20}

% Groupe 9
\DarkenForText{lilac}{20}
\DarkenForText{indigoA}{20}
\DarkenForText{deepindigo}{20}
\DarkenForText{royalpurple}{20}

% Groupe 10
\DarkenForText{coolmint}{20}
\DarkenForText{deepsea}{20}

% Palette deutéranope
\DarkenForText{crimson}{20}
\DarkenForText{brickred}{20}
\DarkenForText{coral}{20}
\DarkenForText{salmon}{20}

\DarkenForText{magenta}{20}
\DarkenForText{fuchsia}{20}
\DarkenForText{violetB}{20}
\DarkenForText{indigoB}{20}

\DarkenForText{royalblueA}{20}
\DarkenForText{azure}{20}
\DarkenForText{skyblueB}{20}
\DarkenForText{tealB}{20}

% Couleur supplémentaire
\DarkenForText{marine}{20}

%%%%%%%%%%%%%%%%%%%%%
%
%.   apx proof
%
%%%%%%%%%%%%%%%%%%%%%



% do
\newcommand\PREUVESENAPPENDIXCOMPATIBLE{
	\newenvironment{appendixproof}{{\bf \LANGUE{Proof}{Démonstration}}:}{\hfill $\Box$}
	\newtheorem{theoremrep}[theorem]{Theorem}
	\newtheorem{lemmarep}[theorem]{Lemma}
	\newtheorem{propositionrep}[theorem]{Proposition}
	\newtheorem{corollaryrep}[theorem]{Corollary}
	\newtheorem{summaryrep}[theorem]{Summary}	
}

\providecommand\apxparameter{}

%% proof inline
%\providecommand\apxparameter{appendix=inline}
%% proof in appendix
%\providecommand\apxparameter{}


\apxproofmode{
\usepackage[\apxparameter]{apxproof}
\theoremstyle{plain}

\newtheoremrep{theorem}{Theorem}
%\newtheoremrep{theoremconjecture}[theorem]{Theorem-To-Be-Proved=Conjecture}
%\newtheoremrep{lemmaconjecture}[theorem]{Conjecture-To-Be-Proved=Conjecture}
\newtheoremrep{claim}[theorem]{Theorem}
\newtheoremrep{proposition}[theorem]{Proposition}
\newtheoremrep{lemma}[theorem]{Lemma}
\newtheoremrep{corollary}[theorem]{Corollary}
%%bug mars 26: \newtheoremrep{summary}[theorem]{Summary}
% Du coup:
\newtheoremrep{resume}[theorem]{Summary}
%bug mars 26: \newtheoremrep{openproblem}[theorem]{Open Problem}
%\newtheoremrep{claim}[theorem]{Claim}
}


%\apxproofmodesubsection{
%%apx proof, pour que compteur correct
%\makeatletter
%\AtBeginDocument{%
%\let\axp@section\subsection
%  \let\axp@oldsection\subsection
%}
%\makeatother
%\renewcommand\appendixprelim{\section{Missing proofs}}
%}


\makeatletter
\apxproofmodesubsection{
\let\axp@section\subsection
%\let\axp@oldsection\subsection
\renewcommand\appendixprelim{\section{Proofs from main documents}}
}
\makeatother


\makeatletter
\newcommand\appendixapxproofsubsection{
\appendix
\let\apx@orig@appendix\appendix
\renewcommand{\appendix}{}
}
\makeatother


%%% Pour checker/vérifier que bon


\newcommand\docheck[1]{\textcolor{check}{#1}}

\usepackage{xcolor}
%\definecolor{chose}{RGB}{0,110,50}
%\definecolor{forestgreen}{RGB}{0,150,70}
%\definecolor{ctruc}{RGB}{164,164,164}
%\definecolor{oceanfoam}{RGB}{0,200,150}
\definecolor{check}{RGB}{160,160,0}




 via \providecommand,
%% ou après via \renewcommand) :
\providecommand\BIBFILES{@@reference-biblio}
\providecommand\FIGCOMMONS{/Users/bournez/lib/LaTeX/Perso/FIG-COMMONS}

\newcommand\affiliationofficielleolivier{LIX, CNRS, École polytechnique, Institut Polytechnique de Paris, Palaiseau, France}
\newcommand\affiliationofficielleolivierslides{LIX, CNRS, École polytechnique, \\ Institut Polytechnique de Paris, Palaiseau, France}


%%%%%%%%%%%%%%%%%%%%%%%%
%
%              begin. Switch modes compatibilités
%
%%%%%%%%%%%%%%%%%%%%%%%

%% Apxproof mode
\providecommand\apxproofmode[1]{}
	% sinon: \newcommand\apxproof[1]{#1}
	
\providecommand\apxproofmodesubsection[1]{}
	% sinon: \newcommand\apxproofmodesubsection[1]{#1}
% utiliser: \appendixapxproofsubsection au lieu de \appendix dans ce cas. 
		

%% Lipics mode
\providecommand\lipicsmode{}
	% sinon: \newcommand\lipicsmode[1]{#1}
		% en vrai sert qu'à définir \nolipics
		
%%%%%%%%%%%%%%%%%%%%%%%%
%
%              end. Switch modes compatibilités
%
%%%%%%%%%%%%%%%%%%%%%%%


%% begin gestion des switchs

% en fait \nopilcs est la seule commande utilisée

\newcommand\nolipics[1]{#1}
\lipicsmode{\renewcommand\nolipics[1]{}}

%% Décorations (fourier + cadres colorés mdframed autour de definition,
%% theorem, exercice, example, remark, ...) : actives par défaut hors slides.
%% Pour les désactiver dans un document : \newcommand\NODECORATION[1]{}
%% avant 
\providecommand\nolncs[1]{#1}	

%% Chemins par défaut (surchargables avant \input{macros} via \providecommand,
%% ou après via \renewcommand) :
\providecommand\BIBFILES{@@reference-biblio}
\providecommand\FIGCOMMONS{/Users/bournez/lib/LaTeX/Perso/FIG-COMMONS}

\newcommand\affiliationofficielleolivier{LIX, CNRS, École polytechnique, Institut Polytechnique de Paris, Palaiseau, France}
\newcommand\affiliationofficielleolivierslides{LIX, CNRS, École polytechnique, \\ Institut Polytechnique de Paris, Palaiseau, France}


%%%%%%%%%%%%%%%%%%%%%%%%
%
%              begin. Switch modes compatibilités
%
%%%%%%%%%%%%%%%%%%%%%%%

%% Apxproof mode
\providecommand\apxproofmode[1]{}
	% sinon: \newcommand\apxproof[1]{#1}
	
\providecommand\apxproofmodesubsection[1]{}
	% sinon: \newcommand\apxproofmodesubsection[1]{#1}
% utiliser: \appendixapxproofsubsection au lieu de \appendix dans ce cas. 
		

%% Lipics mode
\providecommand\lipicsmode{}
	% sinon: \newcommand\lipicsmode[1]{#1}
		% en vrai sert qu'à définir \nolipics
		
%%%%%%%%%%%%%%%%%%%%%%%%
%
%              end. Switch modes compatibilités
%
%%%%%%%%%%%%%%%%%%%%%%%


%% begin gestion des switchs

% en fait \nopilcs est la seule commande utilisée

\newcommand\nolipics[1]{#1}
\lipicsmode{\renewcommand\nolipics[1]{}}

%% Décorations (fourier + cadres colorés mdframed autour de definition,
%% theorem, exercice, example, remark, ...) : actives par défaut hors slides.
%% Pour les désactiver dans un document : \newcommand\NODECORATION[1]{}
%% avant \input{macros}. Retirées automatiquement en mode lipics.
\providecommand\NODECORATION[1]{\ifnotSLIDE{#1}}
\lipicsmode{\renewcommand\NODECORATION[1]{}}

\providecommand\PASSIAPXPROOF[1]{#1}
\apxproofmode{\renewcommand\PASSIAPXPROOF[1]{}}



%% end gestion des switchs

\def\MACROSPOINTEXOUVERT{}


\providecommand\ifnotTDEXAM[1]{#1}
\providecommand\ifnotSLIDE[1]{#1}


  
% pour preuves en appendix
\providecommand\PREUVESENAPPENDIX[1]{}



%%%% STYLES: APRES VOLS


\NODECORATION{
\usepackage{fourier}
}

%% Repli si fourier (et donc fourier-orns, qui définit \danger) n'est pas
%% chargé ; doit rester APRÈS le bloc ci-dessus, sinon fourier-orns échoue
%% avec « \danger already defined ».
\providecommand\danger{}

%%%% Box around environment
\usepackage[framemethod=tikz]{mdframed}

\newcommand\ENVCOLOR[1]{#1}
\newcommand\BEGINTEMPENVCOLOR[1]{\let\envcolorwas\ENVCOLOR\renewcommand\ENVCOLOR[1]{#1}}
\newcommand\ENDTEMPENVCOLOR{\let\ENVCOLOR\envcolorwas}


\NODECORATION{
\surroundwithmdframed[hidealllines=true,backgroundcolor=\ENVCOLOR{blue!10},roundcorner=8pt]{exercice}
\surroundwithmdframed[hidealllines=true,backgroundcolor=\ENVCOLOR{blue!10},roundcorner=8pt]{exercise}
\surroundwithmdframed[hidealllines=true,backgroundcolor=\ENVCOLOR{blue!10},roundcorner=8pt]{exercicee}
\surroundwithmdframed[hidealllines=true,backgroundcolor=\ENVCOLOR{blue!10},roundcorner=8pt]{exerciced}
\surroundwithmdframed[hidealllines=true,backgroundcolor=\ENVCOLOR{blue!10},roundcorner=8pt]{exercicec}
\surroundwithmdframed[hidealllines=true,backgroundcolor=\ENVCOLOR{blue!10},roundcorner=8pt]{exercicedc}
\surroundwithmdframed[hidealllines=true,backgroundcolor=\ENVCOLOR{orange!20},roundcorner=8pt]{example}
\surroundwithmdframed[hidealllines=true,backgroundcolor=\ENVCOLOR{orange!20},roundcorner=8pt]{exemple}
\surroundwithmdframed[hidealllines=true,backgroundcolor=\ENVCOLOR{orange!10},roundcorner=8pt]{remark}
\ifnotSLIDE{\surroundwithmdframed[backgroundcolor=\ENVCOLOR{green!15},roundcorner=8pt]{definition}}
\surroundwithmdframed[hidealllines=true,backgroundcolor=\ENVCOLOR{gray!35}]{proposition}
\ifnotSLIDE{\surroundwithmdframed[backgroundcolor=\ENVCOLOR{gray!35},roundcorner=8pt]{theorem}}
\surroundwithmdframed[hidealllines=true,backgroundcolor=\ENVCOLOR{gray!35}]{lemma}
\surroundwithmdframed[hidealllines=true,backgroundcolor=\ENVCOLOR{gray!35}]{corollary}


\surroundwithmdframed[backgroundcolor=red!35,roundcorner=8pt]{theoremaprouver}
\surroundwithmdframed[backgroundcolor=red!35,roundcorner=8pt]{conjectureaprouver}
}

%%% tcolorbox

\usepackage{tcolorbox}
\tcbuselibrary{skins}


  %%. Commande magique
%% \tcolorboxenvironment{⟨name⟩}{⟨options⟩}: attach options de tcolorbox à environnement
%% alternative: je cree des boites moi meme: voici des exemples


\newtcolorbox{maboiteconfiance}[2][]{colback=red!5!white, colframe=red!75!black,fonttitle=\bfseries, colbacktitle=red!85!black,enhanced,
attach boxed title to top center={yshift=-2mm},
  title={#2},#1}
  
\newtcolorbox{maboiteresume}[2][]{colback=red!5!white, colframe=red!75!black,fonttitle=\bfseries, colbacktitle=red!85!black,enhanced,
attach boxed title to top center={yshift=-2mm},
  title={#2},#1}
  \newtcolorbox{maboiteresumeperso}[2][]{colback=red!5!white, colframe=red!75!black,fonttitle=\bfseries, colbacktitle=red!85!black,enhanced,
attach boxed title to top center={yshift=-2mm},
  title={#2},#1}
  \newtcolorbox{maboitecommentaire}[2][]{colback=red!5!white, colframe=red!75!black,fonttitle=\bfseries, colbacktitle=red!85!black,enhanced,
attach boxed title to top center={yshift=-2mm},
  title={#2},#1}
\newtcolorbox{maboitetruc}[2][]{colback=red!5!white, colframe=red!75!black,fonttitle=\bfseries, colbacktitle=red!85!black,enhanced,
attach boxed title to top left={yshift=-2mm},
  title={#2},#1}

%% Boîte des mises en évidence CARE. Couleurs pensées pour la
%% protanopie : l'axe bleu-jaune reste discriminant (jaune = \IMPORTANT
%% d'Olivier, bleu = CARE), et la distinction ne repose pas que sur la
%% couleur (cadre + bandeau étiqueté, lisibles même en niveaux de gris).
\newtcolorbox{maboiteimportantcare}[2][]{colback=cyan!10!white, colframe=blue!60!black,fonttitle=\bfseries, colbacktitle=blue!70!black,enhanced,
attach boxed title to top center={yshift=-2mm},
  title={#2},#1}
  


%%%% FIN STYLES



%%%%%%%%%%%%%%%%%%%%%%%%%%%%%%%%%%%%%%%%%%%%%%%%%
%%%%%%%%%%%%%%%%%%%%%%%%%%%%%%%%%%%%%%%%%%%%%%%%%
%
%
%                      Version ENSEIGNEMENT 2020
%
%
%%%%%%%%%%%%%%%%%%%%%%%%%%%%%%%%%%%%%%%%%%%%%%%%%
%%%%%%%%%%%%%%%%%%%%%%%%%%%%%%%%%%%%%%%%%%%%%%%%%

%english, french
\ifdefined\CESTDUFRANCAIS
	\providecommand\LANGUE[2]{#2}
	\usepackage[french]{babel}
\else
	\providecommand\LANGUE[2]{#1}
\fi
\providecommand\LANGUEinv[2]{\LANGUE{#2}{#1}}

%%%%%%%%%%%%%%%%%%%%%%%%%%%%%%%%%%%%%%%%%%%%%%%%%
%
%
%                      Packages
%
%
%%%%%%%%%%%%%%%%%%%%%%%%%%%%%%%%%%%%%%%%%%%%%%%%%


%%% SURLIGNEUR
\providecommand\ifnotmodeDIFFUSIONFINALE[1]{
%% 2021: Comprend pas mais assure
\ifdefined\SAFEMODE
\else
#1
\fi
%#1
}
\providecommand\ifmodeDIFFUSIONFINALE[1]{
\ifdefined\SAFEMODE
#1
\else
\fi
}
\newcommand\SURLIGNE[1]{
\ifnotmodeDIFFUSIONFINALE{
\BEGINTEMPENVCOLOR{blue!30}
{                        \colorbox{blue!30}{

\noindent \begin{minipage}{\dimexpr\textwidth-2\fboxsep\relax}
#1
\end{minipage}
}
}
\ENDTEMPENVCOLOR}
\ifmodeDIFFUSIONFINALE{
\noindent \begin{minipage}{\textwidth}
#1
\end{minipage}
}
}


\newenvironment{confiance}{\begin{maboiteresume}[colback=brown!50]{Résumé ici}}{\end{maboiteresume}}
\newcommand\CONFIANCE[2][Confiance ici]{
\begin{maboiteconfiance}[colback=yellow!50]{#1}
\textbf{ATTENTION : #2}
\end{maboiteconfiance}
}

\newenvironment{resume}{\begin{maboiteresume}[colback=brown!50]{Résumé ici}}{\end{maboiteresume}}
\newcommand\RESUME[2][Résumé ici]{
\begin{maboiteresume}[colback=brown!50]{#1}
#2
\end{maboiteresume}
}

\newenvironment{resumeperso}{\begin{maboiteresumeperso}[colback=brown!50]{Résumé perso ici}}{\end{maboiteresumeperso}}
\newcommand\RESUMEPERSO[2][Résumé perso  ici]{
\begin{maboiteresumeperso}[colback=brown!50]{#1}
#2
\end{maboiteresumeperso}
%
%
%\todo[inline, color=brown]{résumé: #1}
%\colorbox{brown}{
%\noindent \begin{minipage}{\textwidth}
%{
%\bf 
%#2
%}
%\end{minipage}
%}
}

\newcommand\DANSLAMARGE[1]{
\ifnotmodeDIFFUSIONFINALE{
\marginpar{\textcolor{brown}{#1}}}
}

\newenvironment{commentaireperso}{\begin{maboitecommentaire}[colback=brown!30]{Commentaire perso} Commentaire perso:}{\end{maboitecommentaire}}
\newcommand\COMMENTAIREPERSO[2][Commentaire perso]{
\begin{maboitecommentaire}[colback=brown!30]{#1}
#2
\end{maboitecommentaire}
}


\newcommand\SURSURLIGNE[2][truc surligné ici]{
%\todo[inline, color=orange]{surligné: #1}
\ifnotmodeDIFFUSIONFINALE{
\BEGINTEMPENVCOLOR{blue!60}
	                 \colorbox{blue!60}{
\noindent \begin{minipage}{\dimexpr\textwidth-2\fboxsep\relax}
{
\bf 
#2
}
\end{minipage}
}
\ENDTEMPENVCOLOR
}
\ifmodeDIFFUSIONFINALE{#2}
}

\newcommand\IMPORTANT[1]{
\ifnotmodeDIFFUSIONFINALE{
\BEGINTEMPENVCOLOR{yellow!100}
		        \colorbox{yellow!100}{

\noindent \begin{minipage}{\dimexpr\textwidth-2\fboxsep\relax}
{ 
#1
}
\end{minipage}
}
\ENDTEMPENVCOLOR
}
\ifmodeDIFFUSIONFINALE{#1}
}

%% Pendant de \IMPORTANT, réservé aux mises en évidence CARE ;
%% \IMPORTANT reste réservé aux annotations d'Olivier. Comme \IMPORTANT,
%% redevient du texte simple en mode DIFFUSIONFINALE.
\newcommand\IMPORTANTCARE[2][Important (CARE)]{
\ifnotmodeDIFFUSIONFINALE{
\begin{maboiteimportantcare}{#1}
#2
\end{maboiteimportantcare}
}
\ifmodeDIFFUSIONFINALE{#2}
}


\newcommand\QUESTION[1]{
\ifnotmodeDIFFUSIONFINALE{
\BEGINTEMPENVCOLOR{oceanfoam!100}
		        \colorbox{oceanfoam!100}{

\noindent \begin{minipage}{\dimexpr\textwidth-2\fboxsep\relax}
{ 
#1
}
\end{minipage}
}
\ENDTEMPENVCOLOR
}
\ifmodeDIFFUSIONFINALE{#1}
}




\newcommand\SOUSLIGNE[1]{
\colorbox{lightgray}{

\noindent \begin{minipage}{\dimexpr\textwidth-2\fboxsep\relax}
{ \color{gray}
#1
}
\end{minipage}
}
}



%% Gestion des cross-references cross-referencing
\usepackage{xr-hyper}
\usepackage{hyperref}

%%
\usepackage{versions}
\usepackage{amsmath,amssymb,mathrsfs,amsfonts}
%\usepackage{mathabx}
\usepackage{listings}
\usepackage{url}
\usepackage{versions}
\usepackage{color}
\usepackage[colorinlistoftodos]{todonotes}
%\usepackage{todonotes}
\usepackage{proof}
\usepackage{xstring}

%% Index
\usepackage{makeidx}
\makeindex
\providecommand\DOINDEXPRINT{}
\providecommand\SIGNALEMOTNOUV[1]{}
\ifdefined\INDEXPRINT
	\ifnotSAFEMODE{
%	\usepackage{showidx}
%	\let\oldindex\index
%	\renewcommand{\index}[1]{\oldindex{#1}\marginpar{LA: \tiny#1}}
	\renewcommand\DOINDEXPRINT{
	      	\let\oldindex\index
		\renewcommand{\index}[1]{\oldindex{##1}\marginpar{IDXLA: \tiny##1}}
		\renewcommand\SIGNALEMOTNOUV[1]{\marginpar{MNV: \small\textbf{##1}}}
	}
	}
\fi

% === gestion des accents


%soit iso latin1
%\usepackage[cyr]{aeguill}
%soit utf8
\usepackage[utf8]{inputenc}
\usepackage[T1]{fontenc}


%%%  Pour citation \Cref (cref, ref)	

\usepackage{cleveref}


%%%%%%%%%%%%%%%%%%%%%%%%%%%%%%%%%%%%%%%%%%%%%%%%%
%
%
%                      Gestion des accents spéciaux.
%
%
%%%%%%%%%%%%%%%%%%%%%%%%%%%%%%%%%%%%%%%%%%%%%%%%%

\input{macros-accents}
%% Tolérant tant que macros-markdown.tex n'est pas dans le dépôt :
\InputIfFileExists{macros-markdown}{}{\typeout{*** macros-markdown.tex introuvable: ignore ***}}


%\usepackage{PDFdraftcopy}


%-- Pour état des chapitres


%\newcommand\ETAT{\draftstring{BROUILLON}}
%
%\newcommand\brouillon{\renewcommand\ETAT{\draftstring{BROUILLON}}}
%
%\newcommand\final{\renewcommand\ETAT{\draftstring{\rien}}}
%
%\newcommand\travail{\renewcommand\ETAT{\draftstring{CHANTIER}}}
%
%\newcommand\bazar{\renewcommand\ETAT{\draftstring{BAZAR TOTAL}}}
%
%\final



%%%%%%%%%%%%%%%%%%%%%%%%%%%%%%%%%%%%%%%%%%%%%%%%%
%
%
%                      TRADUCTION EN COURS
%
%%%%%%%%%%%%%%%%%%%%%%%%%%%%%%%%%%%%%%%%%%%%%%%%%


% === volé Amaury



\usepackage{stmaryrd}
\usepackage{ifthen}
\usepackage{mathtools}
\usepackage{dsfont}
\PASSIAPXPROOF{\usepackage{thmtools}}
\usepackage{longtable}
\usepackage{mathdots}
\usepackage{booktabs}


%%%%%%%%%%%%%%%%%%%%%%%%%%%%%%%%%%%%%%%%%%%%%%%%%
%
%
%                      TIKZ STUFFS
%
%
%%%%%%%%%%%%%%%%%%%%%%%%%%%%%%%%%%%%%%%%%%%%%%%%%

%%% volé pour dessin des GPAC: PAS SUR QUE BESOIN de TOUT

\usepackage{tikz}
\usetikzlibrary{calc}
\usepgflibrary{arrows}
\usetikzlibrary{fit}
\usetikzlibrary{patterns}
\usetikzlibrary{decorations.pathreplacing}
\usetikzlibrary{shapes.multipart}
\usetikzlibrary{shapes.geometric}
\usetikzlibrary{shapes.misc}
\usetikzlibrary{positioning}
\usetikzlibrary{graphs}
\usetikzlibrary{topaths}
\usetikzlibrary{arrows.meta}
\usetikzlibrary{decorations.pathreplacing}
\usetikzlibrary{decorations.markings}
\usetikzlibrary{decorations.pathmorphing}
%
 \usetikzlibrary{calc}
 \usetikzlibrary{automata}
\usetikzlibrary{chains}
\usetikzlibrary{scopes}
 \usetikzlibrary{positioning}
 \usetikzlibrary{through,backgrounds}
\usepackage{circuitikz}
% % pour accolades




%%%%%%%%%%%%%%%%%%%%%%%%%%%%%%%%%%%%%%%%%%%%%%%%%
%
%
%                      MACROS locales
%
%
%%%%%%%%%%%%%%%%%%%%%%%%%%%%%%%%%%%%%%%%%%%%%%%%%


%====================
%                      Moi
%=====================


\newcommand\myaddress{ 
{\sc LIX, CNRS, Ecole Polytechnique, Institut Polytechnique de Paris, }\\ 91128 Palaiseau Cedex, FRANCE \\
{\small Olivier.Bournez@lix.polytechnique.fr}}



%====================
%                      un peu de trucs sales
%=====================

%=  pour style lipics et ses trucs

\providecommand\ifnotPOLYCHAPTERMODE[1]{#1}
\providecommand\ifPOLYCHAPTERMODE[1]{}

\ifnotPOLYCHAPTERMODE{
\ifdefined\thechapter
\providecommand\spnewtheorem[4]{\newtheorem{#1}{#2}[chapter]}
\providecommand\spnewtheoremi[5]{\newtheorem{#1}[#5]{#2}}
\else
\providecommand\spnewtheorem[4]{\newtheorem{#1}{#2}}
\providecommand\spnewtheoremi[5]{\newtheorem{#1}[#5]{#2}}
\fi
}
\ifPOLYCHAPTERMODE{
\providecommand\spnewtheorem[4]{\newtheorem{#1}{#2}}
\providecommand\spnewtheoremi[5]{\newtheorem{#1}[#5]{#2}}
}
\providecommand\nolipics[1]{#1}




%====================
%                      un peu de trucs franchement sales
%=====================


% on l'inclus pas ici.
%\input{macros-moins-propres}

%%%%%%%%%%%%%%%%%%%%%%%%%%%%%%%%%%%%%%%%%%%%%%%%%
%
%
%                      Page de Garde Cours X
%
%
%%%%%%%%%%%%%%%%%%%%%%%%%%%%%%%%%%%%%%%%%%%%%%%%%


\includeversion{metdate}

\newcommand\XPAGEDEGARDE[2]{
\thispagestyle{empty}
\title{{#1} \\[2cm] %\\[3cm] 
{\large #2}
}
\author{
 Olivier Bournez% \inst{1}
\\[0.5cm] 
bournez@lix.polytechnique.fr%{\myaddress}
}
%\date{
%\vspace{2cm}
%\includegraphics[height=2.5cm]{/Users/bournez/lib/LaTeX/Perso/FIG-COMMONS/logo_X_new}
%}
\begin{metdate}
\date{
\vspace{2cm}
Version \LANGUE{of}{du} \today \\[1cm] 
\includegraphics[height=3.5cm]{/Users/bournez/lib/LaTeX/Perso/FIG-COMMONS/logo_X_new}
}
\end{metdate}
\maketitle
}


%%%%%%%%%%%%%%%%%%%%%%%%%%%%%%%%%%%%%%%%%%%%%%%%%
%
%
%                      MACROS UTILES
%
%
%%%%%%%%%%%%%%%%%%%%%%%%%%%%%%%%%%%%%%%%%%%%%%%%%

%%%                      Macros
 
%-------------
%--- vectors  ---
%-------------

 
\newcommand\vectorl[1]{{\mathbf#1}}
\newcommand\tu[1]{\vectorl{#1}}

\newcommand\vp{\vectorl{p}}
\newcommand\va{\vectorl{a}}
\newcommand\vb{\vectorl{b}}
\newcommand\vc{\vectorl{c}}
\newcommand\vf{\vectorl{f}}
\newcommand\vn{\vectorl{n}}
\newcommand\vx{\vectorl{x}}
\newcommand\vX{\vectorl{X}}
\newcommand\vy{\vectorl{y}}
\newcommand\vz{\vectorl{z}}
\newcommand\ve{\vectorl{e}}
\newcommand\vv{\vectorl{v}}
\newcommand\vr{\vectorl{r}}
\newcommand\vu{\vectorl{u}}
\newcommand\vA{\vectorl{A}}
\newcommand\vB{\vectorl{B}}
\newcommand\vC{\vectorl{C}}
\newcommand\vD{\vectorl{D}}

%-------------
%--- overline ---
%-------------


\newcommand\ox{\overline{x}}
\newcommand\oy{\overline{y}}
\newcommand\oc{\overline{c}}
\newcommand\oa{\overline{a}}
\newcommand\ow{\overline{w}}
\newcommand\od{\overline{d}}
\newcommand\ob{\overline{b}}
\newcommand\oP{\overline{P}}
\newcommand\oee{\overline{e}}
\newcommand\ot{\overline{t}}
\newcommand\ou{\overline{u}}
\newcommand\ov{\overline{v}}
\newcommand\oz{\overline{z}}
\newcommand\am{\overline{a}}
\newcommand\om{\overline{m}}
\newcommand\oj{\overline{j}}

%----------------
%--- ensembles ---
%----------------

\newcommand\Q{\mathbb{Q}}
\newcommand\R{\mathbb{R}}
\newcommand\Z{\mathbb{Z}}
\providecommand\C{\mathbb{C}}
\newcommand\N{\mathbb{N}}
\newcommand\K{\mathbb{K}}
\newcommand\F{\mathbb{F}}

\newcommand\dyadic{\mathbb{D}}
\newcommand\Zd{{\Z_2}}
\newcommand\Zp{{\Z_p}}
%
\newcommand\RP{\mathbb{R}^{>0}}
\newcommand\QP{\mathbb{Q}^{>0}}


%----------------
%--- lettres cal ---
%----------------

\newcommand\mK{\mathcal{K}}
\newcommand\mD{\mathcal{D}}
\newcommand\mA{\mathcal{A}}
\newcommand\mL{\mathcal{L}}
\newcommand\mX{\mathcal{X}}
\newcommand\mG{\mathcal{G}}
\newcommand\mE{\mathcal{E}}
\newcommand\mH{\mathcal{H}}
\newcommand\mR{\mathcal{R}}
\newcommand\mF{\mathcal{F}}
\newcommand\mJ{\mathcal{J}}
\newcommand\mO{\mathcal{O}}
\newcommand\mP{\mathcal{P}}
\newcommand\mN{\mathcal{N}}
\newcommand\mS{\mathcal{S}}
\newcommand\mM{\mathcal{M}}
\newcommand\mC{\mathcal{C}}
\newcommand\mV{\mathcal{V}}
\newcommand\mI{\mathcal{I}}
\newcommand\mT{\mathcal{T}}

\newcommand\fC{\mathcal{C}}

%----------------
%--- lettres frak ---
%----------------


\newcommand\kM{\mathfrak{M}}
\newcommand\kN{\mathfrak{N}}
\newcommand\kU{\mathfrak{U}}
\newcommand\kg{\mathfrak{g}}



%--------------------------
%--- Façon noter noms classes ---
%---------------------------

\newcommand{\cp}[1]{\operatorname{#1}}


%--------------------------
%--- Calculabilité  ---
%---------------------------

% -- classes
\LANGUE{
\newcommand\RE{\cp{CE}}
}
{\newcommand\RE{\cp{RE}}
}

\newcommand\Rec{\cp{D}}
\newcommand\PR{\cp{PR}}
\newcommand\Gregn{\mE_n}


%--------------------------
%--- Réductions  ---
%---------------------------

\newcommand\lem{\le_{m}}
\newcommand\lelog{\le_{L}}
\newcommand\equivlog{\equiv_{L}}



%--------------------------
%--- Complexité  ---
%---------------------------

% COMPLEXITE

% -- P
\newcommand{\Ptime}{\cp{P}}      % P
\newcommand\PTIME{\Ptime}
\newcommand\PTIMEs[1]{{\PTIME}_{#1}}
\newcommand\PTIMEsp[1]{{\PTIME}_{#1}^0}
\renewcommand\P{\PTIME}
\newcommand{\FPtime}{\cp{FPTIME}}

% -- NP
\newcommand\NP{\cp{NP}}
\newcommand{\NPtime}{\NP}
\newcommand\NPs[1]{\cp{NP}_{#1}}
\newcommand{\FNP}{\cp{FNP}}

\newcommand\NDP{\cp{NDP}}
\newcommand\coNDP{\cp{coNDP}}
\newcommand\NDPs[1]{\cp{NDP}_{#1}}

% -- coNP
\newcommand\coNP{\cp{coNP}}

% -- PSPACE
\newcommand\PSPACE{\cp{ PSPACE}}
\newcommand\NPSPACE{\cp{ NPSPACE}}
\newcommand{\Pspace}{\PSPACE}
\newcommand{\FPspace}{\cp{FPSPACE}}


% -- LOGSPACE
\newcommand\LSPACE{\cp{LOGSPACE}}
\newcommand\NLSPACE{\cp{NLOGSPACE}}
\newcommand\coNLSPACE{\cp{coNLOGSPACE}}
\newcommand\coNL{\cp{coNL}}

% -- Autres
\newcommand\EXPTIME{\cp{EXPTIME}}
\newcommand\NEXPTIME{\cp{NEXPTIME}}
\newcommand\EXPSPACE{\cp{EXPSPACE}}
\newcommand\PH{\cp{PH}}
\newcommand\NC{\cp{NC}}
\newcommand\AC{\cp{AC}}
\newcommand\PPAD{\cp{\mathbf{PPAD}}}
\newcommand\PLS{\cp{\mathbf{PLS}}}

% -- Reverse math
\newcommand{\WKL}{\ensuremath{\docheck{\mathsf{WKL}_0}}}
\newcommand{\ACA}{\ensuremath{\docheck{\mathsf{ACA}_0}}}
\newcommand{\ATR}{\ensuremath{\docheck{\mathsf{ATR}_0}}}
\newcommand{\RCAO}{\ensuremath{\docheck{\mathsf{RCA}_0}}}
\newcommand{\PiCA}{\ensuremath{\docheck{\Pi^1_1\!\text{-}\mathsf{CA}_0}}}


% -- circuits
\newcommand\CIRCUIT[2]{\cp{CIRCUIT(#1,#2)}}
\newcommand\UCIRCUIT[2]{\cp{UCIRCUIT(#1,#2)}}

% -- Classes proba
\newcommand\PP{\cp{PP}}
\newcommand\RPP{\cp{RP}}
\newcommand\ZPP{\cp{ZPP}}
\newcommand\coRP{\cp{coRP}}
\newcommand\BPP{{\cp{BPP}}}

% -- IP
\newcommand\IP{\cp{IP}}
\newcommand\PCP{\cp{PCP}}

% -- non-uniforme
\newcommand\nuP{\mathbb{P}}
\newcommand\nuPs[1]{\mathbb{P}_{#1}}
\newcommand\nuPsp[1]{\mathbb{P}_{#1}^0}
\newcommand\nuNP{\mathbb{N}\mathbb{P}}
%\newcommand\poly{/{poly}}
\newcommand\Ppoly{\PTIME_{\poly}}
\newcommand\NPpoly{\NP_{\poly}}

% -- conditions d'uniformité
\newcommand\Luniforme{\cp{L}}
\newcommand\BCuniforme{\cp{BC}}


% Classes de circuits
\newcommand\TC{\cp{TC}}
\newcommand\LFP{\cp{LFP}}
\newcommand\FAC{\cp{FAC}}
\newcommand\FACC{\cp{FACC}}
\newcommand\FTC{\cp{FTC}}
\newcommand\FNC{\cp{FNC}}

% Logiques temorelles
\newcommand\CTL{\cp{CTL}}
\newcommand\LTL{\cp{LTL}}

% -- mesure temps/espace/taille

\newcommand\DTIME{\cp{DTIME}}
\newcommand\TIME[1]{\cp{TIME(#1)}}
\newcommand\TIMEs[2]{\cp{TIME_{#2}(#1)}}
\newcommand\TIMEsp[2]{\cp{TIME^0_{#2}(#1)}}
\newcommand\NTIME[1]{\cp{NTIME(#1)}}
\newcommand\NTIMEs[2]{\cp{NTIME_{#2}(#1)}}
\newcommand\NTIMEsp[2]{\cp{NTIME_{#2}^0(#1)}}
\newcommand\DSPACE{\cp{DSPACE}}
\newcommand\SPACE[1]{\cp{SPACE(#1)}}
\newcommand\SPACEs[2]{\cp{SPACE_{#2}(#1)}}
\newcommand\NSPACE[1]{\cp{NSPACE(#1)}}
\newcommand\NSPACEs[2]{\cp{NSPACE_{#2}(#1)}}
\newcommand\coNSPACE[1]{\cp{coNSPACE(#1)}}
\newcommand\TIMEON[2]{\cp{TIME(#1,#2)}}
\newcommand\SPACEON[2]{\cp{SPACE(#1,#2)}}
\newcommand\SIZE[1]{\cp{ size(#1)}}
\newcommand\ATIME{\cp{ATIME}}
\newcommand\ASPACE{\cp{ASPACE}}


%% autre:
\newcommand\LH{\cp{LH}}
\newcommand\FO{\cp{FO}}
\newcommand\SO{\cp{SO}}
\newcommand\ESOhorn{\cp{ESO_{horn}}}
\newcommand\SOhorn{\cp{SO_{horn}}}
\newcommand\ACzero{\cp{AC_0}}
\newcommand{\LINSPACE}{\cp{LINSPACE}}
\newcommand{\mFLINSPACE}{\mF_{\cp{LINSPACE}}}
\newcommand{\LOGSPACE}{\cp{LOGSPACE}}
\newcommand\ALOGTIME{\cp{ALOGTIME}}
\newcommand\RF{\cp{RF}}


\newcommand\BOOL[1]{\cp{{BOOLEAN(#1)}}}
\newcommand\DET{\cp{ DET}}

% % -- dirty


\newcommand{\cPspace}{\cp{\sharp PSPACE}}
\newcommand{\cAPtime}{\cp{\sharp APTIME}}
%\newcommand{\FPspace}{\cp{FPSPACE}}
%\newcommand{\FPspaceN}{\cp{FPSPACE}^{+}}
%\newcommand{\FPspaceN}{\cp{FPSPACE}}



%-------------------------
%--- Problèmes de décision  ---
%--------------------------

\newcommand\decisionname[1]{\cp{#1}}


\newcommand\Luniv{\decisionname{L_{univ}}}
\newcommand\HALTINGPROBLEM{\decisionname{HALTING-PROBLEM}}
\newcommand\SAT{\decisionname{SAT}}
\newcommand\MAXtSAT{\decisionname{MAX3SAT}}
\newcommand\REACH{\decisionname{REACH}}
\newcommand\CIRCUITVALUE{\decisionname{CIRCUITVALUE}}
\newcommand\MONOTONECIRCUITVALUE{\decisionname{MONOTONECIRCUITVALUE}}
\newcommand\THEOREMS{\decisionname{THEOREMS}}
\newcommand\QCIRCUITSAT{\decisionname{QCIRCUITSAT}}
\newcommand\QSAT{\decisionname{QSAT}}
\newcommand\QBF{\decisionname{QBF}}
\newcommand\SATVALUE{\decisionname{SATVALUE}}
\newcommand\SUCCINTSAT{\decisionname{SUCCINTSAT}}
\newcommand\SUCCINTSATVALUE{\decisionname{SUCCINTSATVALUE}}
\newcommand\GEOGRAPHY{\decisionname{GEOGRAPHY}}
\newcommand\PARITY{\decisionname{PARITY}}
\newcommand\MAJ{\decisionname{MAJ}}
\newcommand\CVP{\CIRCUITVALUE}
\newcommand\BOOLEANPROGRAMVALUE{BOOLEAN~PROGRAM~VALUE}

\newcommand\CLIQUE{\cp{CLIQUE}}
\LANGUEinv{
\newcommand\RECOUVREMENTDESOMMETS{\cp{RECOUVREMENT~DE~SOMMETS}}
\newcommand\RS{\RECOUVREMENTDESOMMETS}
}{
\newcommand\RECOUVREMENTDESOMMETSeng{\cp{VERTEX~COVER}}
\newcommand\RSeng{\RECOUVREMENTDESOMMETSeng}
}

\LANGUEinv{\newcommand\COUPUREMAXIMALE{\cp{COUPURE~MAXIMALE}}}
{\newcommand\COUPUREMAXIMALEeng{\cp{MAXIMAL~CUT}}}
\newcommand\CIRCUITHAMILTONIEN{\CHAMILTONIEN}
\LANGUEinv{\newcommand\VOYAGEURDECOMMERCE{\VCOMMERCE}}
{\newcommand\VOYAGEURDECOMMERCEeng{\VCOMMERCEeng}}

\LANGUEinv{\newcommand\CIRCUITLEPLUSLONG{\cp{CIRCUIT~LE~PLUS~LONG}}}
{\newcommand\CIRCUITLEPLUSLONGeng{\cp{LONGEST~CIRCUIT}}}
\LANGUEinv{\newcommand\SOMMEDESOUSENSEMBLE{\SUBSETSUM}}{
\newcommand\SOMMEDESOUSENSEMBLEeng{\SUBSETSUMeng}}
\LANGUEinv{
\newcommand\SACADOS{\cp{SAC~A~DOS}}}
{\newcommand\SACADOSeng{\cp{KNAPSACK}}}

%français
\LANGUEinv{
\newcommand\CHAMILTONIEN{\cp{CIRCUIT~HAMILTONIEN}}}
{\newcommand\CHAMILTONIENeng{\cp{HAMILTONIAN~CIRCUIT}}}
\LANGUEinv{
\newcommand\VCOMMERCE{\cp{VOYAGEUR~DE~COMMERCE}}}
{\newcommand\VCOMMERCEeng{\cp{TRAVELING~SALESMAN}}}
\LANGUEinv{
\newcommand\kCOLORABILITE{\cp{k-COLORABILITE}}}
{\newcommand\kCOLORABILITEeng{\cp{k-COLORABILITY}}}
\LANGUEinv{
\newcommand\COLORIAGE{\cp{BON~COLORIAGE}}}
{\newcommand\COLORIAGEeng{\cp{GOOD~COLORING}}}
\LANGUEinv{
\newcommand\tCOLORABILITE{\cp{3-COLORABILITE}}}
{\newcommand\tCOLORABILITEeng{\cp{3-COLORABILITY}}}
\LANGUEinv{
\newcommand\SUBSETSUM{\cp{SOMME~DE~SOUS~ENSEMBLE}}}
{\newcommand\SUBSETSUMeng{\cp{SUBSET~SUM}}}
\newcommand\PARTITION{\cp{PARTITION}}
\LANGUEinv{
\newcommand\NOMBREPREMIER{\decisionname{NOMBRE~ PREMIER}}}
{\newcommand\NOMBREPREMIEReng{\decisionname{PRIME~NUMBER}}}
\LANGUEinv{
\newcommand\CODAGE{\decisionname{CODAGE}}}
{\newcommand\CODAGEeng{\decisionname{ENCODING}}}
\newcommand\kSAT{\cp{k-SAT}}
\newcommand\tSAT{\cp{3-SAT}}
\newcommand\NAESAT{\cp{NAESAT}}
\newcommand\NAEtSAT{\cp{NAE3SAT}}
\newcommand\NAEqSAT{\cp{NAE4SAT}}
\newcommand\STABLE{\cp{\LANGUE{INDEPENDANT~SET}{STABLE}}}

%\newcommand\CM{\cp{COUPURE MAXIMALE}}
\LANGUEinv{
\newcommand\POSTpb{\cp{Problème~de~ la~correspondance~de~Post}}}
{\newcommand\POSTpbeng{\cp{Post~correspondence~problem}}}

\LANGUEinv{
\newcommand\PARITE{\cp{PARITE}}}
{\newcommand\PARITEeng{\cp{PARITY}}}


%-------------------------
%--- Modèles parallèles  ---
%--------------------------


% PARALLELISME
\newcommand\RAM{\cp{RAM}}
%
\newcommand\PRAM{\cp{PRAM}}
\newcommand\CRCW{\cp{CRCW}}
\newcommand\CREW{\cp{CREW}}
\newcommand\EREW{\cp{EREW}}


%------------------------
%--- codages  ---
%------------------------


\newcommand\codage[1]{\langle #1 \rangle}
\newcommand\tuple[1]{\codage{#1}}
\newcommand\paire[2]{\codage{#1,#2}}
\newcommand\triplet[3]{\codage{#1,#2,#3}}
\newcommand\quadruplet[4]{\codage{#1,#2,#3,#4}}
\newcommand\cinquplet[5]{\codage{#1,#2,#3,#4,#5}}


%------------------------
%--- descriptive set theory  ---
%------------------------

\newcommand\baire{\mathcal{N}}
\newcommand\peano{\overline{\mathcal{N}}} %% A MOI
\newcommand\cantor{\mathcal{C}}
%
\newcommand\bSigma{\mathbf{\Sigma}}
\newcommand\bPi{\mathbf{\Pi}}
\newcommand\bDelta{\mathbf{\Delta}}
%


%%%%%% ABOUT NOMS QUI CHANGES

\newcommand\WF{\cp{WF}}
\newcommand\WO{\WF}
\newcommand\LO{\cp{LO}}
\newcommand\DLO{\cp{DLO}}
% 
\newcommand\rk{\cp{rank}}


%% Moins propre:

\newcommand\I{\mathbb{I}}
\renewcommand\H{\mathbb{H}}
%
\newcommand\leKB{<_{KB}}
%
\newcommand\finiomega{{<\omega}}
\newcommand\compl{^\frown}
\newcommand\prefix{ ~ prefix~  } 
\newcommand\X{{X}}
\newcommand\Y{\mathcal{Y}}
\newcommand\subsetstrict{\subset}

\renewcommand\complement[1]{#1^{c}}


%------------------------
%--- schémas de récursion ---
%------------------------

%%                      Source papier MFCS 2019


% COMPLEXITY

%% MFCS differe de Clote:
%% Version à taille restreinte

\newcommand\co{\cp{co}}

\newcommand\mFPTIME{\mF_{PTIME}}
\newcommand{\mFPSPACE}{\mF_{SPACE}}
\newcommand{\mFPH}{\mF_{PH}}

\newcommand\SigmaP{\Sigma^P}
\newcommand\DeltaP{\Delta^{P}}
\newcommand\PiP{\Pi^{P}}
\newcommand\coSigmaP{\co\SigmaP}
\newcommand\coPiP{\mF{\co\PiP}}

%% autres classes:


%% Schemas pour complexité implicite

\newcommand\COMP{\cp{ COMP}}
\newcommand\CRN{\cp{ CRN}}
\newcommand\BRN{\cp{ BRN}}
\newcommand\SBRN{\cp{ SBRN}}
\newcommand\BR{\cp{ BR}}
\newcommand\kBR{k-\cp{ BR}}
\newcommand \WBPR{\cp{ WBPR}}
\newcommand \BMIN{\cp{ BMIN}}
\newcommand \BSUM{\cp{ BSUM}}
\newcommand\SBSUM{\cp{ SBSUM}}
\newcommand \BPROD{\cp{ BPROD}}
\newcommand \SBPROD{\cp{ SBPROD}}
\newcommand \SCOMP{\cp{ SCOMP}}
\newcommand \SR{\cp{ SR}}
\newcommand \SRN{\cp{ SRN}}

%% devrait etre défini
\newcommand\WBRN{\cp{ WBRN}}



%-- godel
\newcommand\INTDIV{\cp{INTDIV}}
\newcommand\DIV{\cp{DIV}}
\newcommand\EVEN{\cp{EVEN}}
\newcommand\ODD{\cp{ODD}}
\newcommand\PRIME{\cp{PRIME}}
\newcommand\POWER{\cp{POWER}}
\newcommand\BIT{\cp{BIT}}
\newcommand\negHP{\overline{\cp{HP}}}
\newcommand\HP{\cp{HP}}
\newcommand\negLuniv{\overline{\Luniv}}
\newcommand\VALCOMP{\cp{VALCOMP}}
\newcommand\mWINDOW{\cp{LEGAL}}
\newcommand\mmWINDOW{\cp{WINDOW}}
\newcommand\LENGTH{\cp{LENGTH}}
\newcommand\DIGIT{\cp{DIGIT}}
\newcommand\MATCH{\cp{MATCH}}
\newcommand\MOVE{\cp{MOVE}}
\newcommand\START{\cp{START}}
\newcommand\CELL{\cp{CELL}}
\newcommand\HALT{\cp{HALT}}
\newcommand\SUBST{\cp{SUBST}}
\newcommand\PROOF{\cp{PROOF}}
\newcommand\PROVABLE{\cp{PROVABLE}}
\newcommand\CONSIST{\cp{CONSIST}}



%---------------------
%--- calculus (AP) ---
%---------------------

% \jacobian{f}{x}
\newcommand{\jacobian}[1]{J_{#1}}
%\newcommand{\grad}[1]{\overrightarrow{\operatorname{grad}}~{#1}}
\newcommand{\grad}[1]{\nabla{#1}}
\newcommand{\scalarprod}[2]{{#1}\cdot{#2}}
\newcommand{\bigO}[1]{\mathcal{O}\left(#1\right)}
\newcommand\smallO{o}
\newcommand{\softO}[1]{\tilde{\mathcal{O}}\left(#1\right)}
\newcommand{\taylor}[3]{{T_{#1}^{#2}#3}}
\newcommand{\transpose}[1]{{#1}^T}
\newcommand{\pastsup}[2]{{\sup}_{#1}#2}

\DeclareMathOperator\arctanh{arctanh}



%---------------
%--- norms (AP) ---
%---------------


\newcommand{\inorm}[2]{\left\lVert{#1}\right\rVert_{#2}}
\newcommand{\twonorm}[1]{\inorm{#1}{2}}
\newcommand{\infnorm}[1]{\inorm{#1}{}}
\newcommand{\normsup}[1]{\infnorm{#1}{}}
%
\newcommand\hausdorff{d_H}
%
\newcommand\norm[1]{\infnorm{#1}}


%----------------
%--- topology ---
%----------------

% \openball{pt}{radius}
 \newcommand{\openball}[2]{B_{#2}(#1)}
\newcommand{\closedball}[2]{\ovl{B}_{#2}(#1)}
\newcommand{\closure}[1]{\ovl{#1}}
% \newcommand{\interior}[1]{\mathring{#1}}
% \newcommand{\border}[1]{\partial{#1}}




%-----------------
%--- functions (AP) ---
%-----------------
 
\newcommand{\lambdafun}[2]{(#1)\mapsto{#2}}
\newcommand{\lambdafunex}[2]{#1\mapsto{#2}}
\newcommand{\argmin}{\operatornamewithlimits{argmin}}
\newcommand{\argmax}{\operatornamewithlimits{argmax}}


 
%-----------------------
%--- Machines de Turing  ---
%-----------------------

\newcommand\blanc{\vectorl{B}}



%-----------------
%--- Logique ---
%-----------------
 
\newcommand\sem[1]{{[\semspace[}#1{]\semspace]}}
\newcommand\sems[1]{{[\semspace[}#1{]\semspace]}}
\newcommand\semss[1]{\sem{#1}_{\sigma'}}
\newcommand\semspace{\kern -.25ex}
\newcommand\true{true}
\newcommand\false{false}
%extension élémentaire:
\newcommand\elem{\preceq}
\newcommand\godel[1]{\lceil #1 \rceil}

%-----------------
%--- Model theory ---
%-----------------

\newcommand\hull[1]{\langle #1 \rangle}

%-----------------
%--- Ordinaux ---
%-----------------

\newcommand\omegaunCK{\omega_1^{CK}}
\newcommand\omegaunck{\omegaunCK}

\DeclareMathOperator{\cof}{cof}

%-----------------
%--- Surréels (QG) ---
%-----------------

\newcommand{\On}{\textbf{On}}
\newcommand{\Nobf}{\textbf{No}}
\newcommand{\Nobb}{\ensuremath{\mathbbmss{No}}}
\newcommand{\Nolt}[1]{\ensuremath{\Nobf_{< #1}}}
\newcommand{\Noleq}[1]{\ensuremath{\Nobf_{\leq #1}}}
\newcommand{\Noltbb}[1]{\ensuremath{\Nobb_{< #1}}}
\providecommand\No{\Nobf}
\newcommand\Noalpha{\Nolt{\alpha}}
\newcommand\Nolambda{\Nolt{\lambda}}

\DeclareMathOperator{\length}{length}
\DeclareMathOperator{\supp}{supp} % operateur support


%----------------------
%---Maths  de base (AP)  ---
%----------------------

\DeclareMathOperator{\dom}{dom} % operateur support

\DeclareMathOperator{\size}{size}

\newcommand{\powerset}[1]{\mathcal{P}(#1)}
\newcommand{\card}[1]{{\##1}}
\newcommand\priv{ - }
\newcommand\prive{-}
\renewcommand{\restriction}{\mathord{\upharpoonright}}
\newcommand\restrict[1]{_{\restriction #1}}

% \providecommand\mod{}
%% UNDEFINE:
\let\mod\relax
\let\div\relax
\DeclareMathOperator{\mod}{mod} 
\DeclareMathOperator{\div}{div} 

% partieentier
\newcommand\entieres[1]{\lceil #1 \rceil}
\newcommand\entierei[1]{\lfloor #1 \rfloor}


%----------------------
%---Symbole danger ---
%----------------------


%\providecommand*{\TakeFourierOrnament}[1]{{%
%\fontencoding{U}\fontfamily{futs}\selectfont\char#1}}
%\providecommand*{\danger}{{\Huge \TakeFourierOrnament{66}}}



%%%%%%%%%%%%%%%%%%%%%%%%%%%%%%%%%%%%%%%%%%%%%%%%%
%
%
%                      ASMeries
%
%
%%%%%%%%%%%%%%%%%%%%%%%%%%%%%%%%%%%%%%%%%%%%%%%%%



% CIRCUITS
%-- particuliers
\newcommand\REPLACE{\cp{REPLACE}}
\newcommand\EQUAL{\cp{EQUAL}}
\newcommand\EVALUE{\cp{EVALUE}}
\newcommand\TVALUE{\cp{TVALUE}}
\newcommand\UPDATE{\cp{UPDATE}}
\newcommand\ASSIGN{\cp{ASSIGN}}
\newcommand\PAR{\cp{PAR}}
\newcommand\DO{\cp{DO}}


% ASM
\newcommand\emplacement[2]{(#1,#2)}
\newcommand\contenu[2]{\sem{(#1,#2)}}
\newcommand\affecte[3]{(#1,#2)\leftarrow#3}
\newcommand\ctl{ctl\_state}


%%%%%%%%%%%%%%%%%%%%%%%%%%%%%%%%%%%%%%%%%%%%%%%%%
%
%
%                      GPACqueries
%
%
%%%%%%%%%%%%%%%%%%%%%%%%%%%%%%%%%%%%%%%%%%%%%%%%%



%------------------------
%--- generable zoo (GPAC) ---
%------------------------

\newcommand{\myop}[1]{\operatorname{#1}}
\newcommand{\abs}{\myop{abs}}
\newcommand{\sg}{\myop{sg}}
\newcommand{\ip}[1]{\myop{ip}_{#1}}
\newcommand{\rnd}{\myop{rnd}}
\newcommand{\lil}{\myop{lil}}
\newcommand{\lxh}{\myop{lxh}}
\newcommand{\hxl}{\myop{hxl}}
\newcommand{\plil}{\myop{plil}}
\newcommand{\reach}{\myop{reach}}
%\newcommand{\mreach}{\myop{mreach}}
\newcommand{\watchdog}{\myop{watchdog}}
\newcommand{\monitor}{\myop{monitor}}
%\newcommand{\norm}{\myop{norm}}
\newcommand{\mx}{\myop{mx}}
\newcommand{\mn}{\myop{mn}}
\newcommand{\sample}{\myop{sample}}
\newcommand{\select}{\myop{select}}
\newcommand{\ssine}[1]{\sin_{#1}}
\newcommand{\clamp}{\myop{clamp}}

\newcommand{\fnsg}[3]{tanh((#1)*(#2)*(#3))}
\newcommand{\fnipone}[3]{(1+\fnsg{(#1)-1}{#2}{#3})/2.}
\newcommand{\fnipinf}[1]{#1-sin(2*pi*(#1))/2./pi}%
\newcommand{\fnlil}[5]{%
\fnipone{(#3)-(#1)+1-((#2)-(#1))/8.}{#4+log(1+4./(#2-#1)*(#5)*(#5))+log(1+4)}{8./(#2-#1)}%
*\fnipone{#2-#3+1-(#2-#1)/8.}{#4+log(1+4/(#2-#1)*(#5)*(#5))+log(1+4)}{8./(#2-#1)}%
*4./(#2-#1)*(#5)}
\newcommand{\fnlxh}[5]{\fnipone{(#3)-(#1+#2)/2.+1}{#4+log(1+(#5)*(#5))}{2./(#2-#1)}*(#5)}
\newcommand{\fnhxl}[5]{\fnipone{(#1+#2)/2.-(#3)+1}{#4+log(1+(#5)*(#5))}{2./(#2-#1)}*(#5)}
\newcommand{\fnplilf}[4]{sin(((#4)-((#1)+(#2))/2.+(#3)/4.)*2.*pi/(#3))}
\newcommand{\fnplil}[6]{(#6)*\fnlxh{\fnplilf{#1}{#2}{#3}{#1}}%
{\fnplilf{#1}{#2}{#3}{#1+((#2)-(#1))/4.}}%
{\fnplilf{#1}{#2}{#3}{#4}}%
{#5+2.+log(1+(#6)*(#6))}%
{1./4.+2./((#2)-(#1))}}
\newcommand{\fnssine}[2]{sin(#2)*(1-cos(#2)/cos(#1))}



%--------------------
%--- complexity GPAC ---
%--------------------


%\ifthenelse{\equal{#1}{}}{Empty.}{Nonempty.}
\newcommand{\myclass}[1]{\operatorname{#1}}
% \newcommand{\gcomp}[2]{\ensuremath{\myclass{GCOMP}(#1,#2)}}
% \newcommand{\ggenerable}[1]{\textcolor{red}{\ensuremath{#1-}\text{generable}}}
 \newcommand{\gval}[2][]{\ensuremath{\myclass{GVAL}_{#1}\ifthenelse{\equal{#2}{}}{}{[#2]}}}
 \newcommand{\gpval}[1][]{\ensuremath{\myclass{GPVAL}_{#1}}}
 \newcommand{\gc}[3][]{\ensuremath{\myclass{ATSC}(#2,#3)}}
 \newcommand{\gpc}[1][]{\ensuremath{\myclass{ATSP}}}
\newcommand{\cgc}[1][]{\ensuremath{\myclass{ATS}}}
\newcommand{\gexpc}[1][]{\ensuremath{\myclass{AEXP}}}
\newcommand{\gwc}[2]{\ensuremath{\myclass{AW}(#1,#2)}}
\newcommand{\gpwc}{\ensuremath{\myclass{AWP}}}
\newcommand{\cgwc}{\ensuremath{\myclass{AW}}}
\newcommand{\grc}[3]{\ensuremath{\myclass{AR}(#1,#2,#3)}}
\newcommand{\gprc}{\ensuremath{\myclass{ARP}}}
\newcommand{\gsc}[3]{\ensuremath{\myclass{AS}(#1,#2,#3)}}
\newcommand{\gpsc}{\ensuremath{\myclass{ASP}}}
\newcommand{\goc}[3]{\ensuremath{\myclass{AO}(#1,#2,#3)}}
\newcommand{\gpoc}{\ensuremath{\myclass{AOP}}}
\newcommand{\cgoc}{\ensuremath{\myclass{AO}}}
\newcommand{\guc}[4]{\ensuremath{\myclass{AX}(#1,#2,#3,#4)}}
\newcommand{\gpuc}{\ensuremath{\myclass{AXP}}}
\newcommand{\glc}[2][]{\ensuremath{\myclass{AL}(#2)}}
\newcommand{\gplc}[1][]{\ensuremath{\myclass{ALP}}}
\newcommand{\cglc}[1][]{\ensuremath{\myclass{AL}}}

%\newcommand{\intinterv}[2]{\{#1,\ldots,#2\}}
\newcommand{\intinterv}[2]{\llbracket#1,#2\rrbracket}

\newcommand{\Rgen}{\R_{G}}
\newcommand{\Rpoly}{\R_{P}}

%%%%%%%%%%%%%%%%%%%%%%%%%%%%%%%%%%%%%%%%%%%%%%%%%
%
%
%                      Mes codages obscures
%
%
%%%%%%%%%%%%%%%%%%%%%%%%%%%%%%%%%%%%%%%%%%%%%%%%%




%%% Codages

% \newcommand\gammaoc{\gamma_{[0,1]}}
% \newcommand\gammaonc{\gamma_{\N}}
\newcommand\gammaTMc{\gamma^{TM}_{[0,1]}}
\newcommand\gammaTMcatan{\gamma^{TM}_{(0,1)^2}}
\newcommand\gammaTMnc{\gamma^{TM}_{\N^3}}
\newcommand\gammaTMncd{\gamma^{TM}_{\N^2}}
% \newcommand\gammaTMe{\gamma^{TM}_{*}}
% \newcommand\gammas{\gamma_{sab}}
% \newcommand\gammaord{\gamma_{cantor}}


\newcommand\enccompact{\nu_{X}}
 \newcommand\encinfinite{\nu_\N}
 \newcommand\enc{\nu}
 
%%%%%%%%%%%%%%%%%%%%%%%%%%%%%%%%%%%%%%%%%%%%%%%%%
%
%
%                      DESSINS ET IMAGES
%
%
%%%%%%%%%%%%%%%%%%%%%%%%%%%%%%%%%%%%%%%%%%%%%%%%%


\newcommand\IMAGE[2]{ \begin{center} 
    \includegraphics[#1]{#2}
  \end{center}
  }


  
%%%% POUR FAIRE DESSIN.
\newenvironment{dessinpp}[1]{
\begin{tikzpicture}[baseline=(current bounding
      box.west),#1]
}
{\end{tikzpicture}}




\newenvironment{dessinp}[1]{
\begin{center}\begin{tikzpicture}[baseline=(current bounding
      box.north),#1]
}
{\end{tikzpicture}
\end{center}}


\newenvironment{dessin}{
\begin{center}\begin{tikzpicture}[baseline=(current bounding
      box.north)]
}
{\end{tikzpicture}
\end{center}}



\newenvironment{dessind}{
\begin{tikzpicture}[baseline=(current bounding
      box.north)]
}
{\end{tikzpicture}
}



%%%%%%%%%%%%%%%%%%%%%%%%%%%%%%%%%%%%%%%%%%%%%%%%%
%
%
%                      POUR DESSINS GPAC
%
%
%%%%%%%%%%%%%%%%%%%%%%%%%%%%%%%%%%%%%%%%%%%%%%%%%


%%%% MACROS DE BASE.


\newcommand\constant[2]
{
node (#1)  [shape=circle,draw] {#2};
\draw (#1.east) -- +(0.3,0) coordinate (#1-o) ; 
\draw (#1.west) -- +(-0.3,0) coordinate (#1-i) ; 
}

\newcommand\basicproduct[2]
{
node (#1) {#2};
\draw (#1) +(-0.4,0.4) -- +(0.4,0.4) -- +(0.8,0) -- +(0.4,-0.4) -- +(-0.4,-0.4) --
+(-0.4,0.4);
\draw (#1) ++(0.8,0) -- ++(+0.3,0) coordinate (#1-o) ;
\draw (#1) ++ (-0.4,0.2) -- +(-0.3,0) coordinate (#1-i1) ;
\draw (#1) ++(-0.4,-0.2) -- +(-0.3,0) coordinate (#1-i2) ;
}

\newcommand\product[1]{\basicproduct{#1}{$+\Pi$}}
\newcommand\negativeproduct[1]{\basicproduct{#1}{$-\Pi$}}


%%BEGIN INTERNAL
\newcommand\basicsummer[1]{
coordinate (#1) {};
\draw (#1) +(-0.5,0.6) -- +(0.5,0) -- +(-0.5,-0.6) -- +(-0.5,0.6);
\draw (#1) ++(0.5,0) -- +(+0.3,0) coordinate (#1-o) ;
}

\newcommand\basicintegrator[1]{
coordinate (#1) {};
\draw (#1) +(-0.5,0.6) -- +(0.5,0) -- +(-0.5,-0.6) -- +(-0.5,0.6);
\draw (#1) ++(0.5,0) -- +(+0.3,0) coordinate (#1-o) ;
\draw (#1) +(-0.5,-0.6) -| +(-0.8,0.6) -- +(-0.5,0.6) ;
 \draw (#1) ++(-0.65,0.6) -- +(0,+0.3) coordinate (#1-init) ;
}

\newcommand\splitfour[2]{
\draw (#2) ++ (#1,0.4) -- +(-0.3,0) coordinate (#2-i1) ;
\draw (#2) ++ (#1,0.15) -- +(-0.3,0) coordinate (#2-i2) ;
\draw (#2) ++ (#1,-0.15) -- +(-0.3,0) coordinate (#2-i3) ;
\draw (#2) ++ (#1,-0.4) -- +(-0.3,0) coordinate (#2-i4) ;
}

\newcommand\splitthree[2]{
\draw  (#2) ++(#1,0.3) -- +(-0.3,0) coordinate (#2-i1) ;
\draw  (#2) ++(#1,0) -- +(-0.3,0) coordinate (#2-i2) ;
\draw  (#2) ++(#1,-0.3) -- +(-0.3,0) coordinate (#2-i3) ;
}

\newcommand\splittwo[2]{
\draw  (#2) ++(#1,0.2) -- +(-0.3,0) coordinate (#2-i1) ;
\draw  (#2) ++(#1,-0.2) -- +(-0.3,0) coordinate (#2-i2) ;
}

\newcommand\splitone[2]{
\draw (#2) ++(#1,0) -- +(-0.3,0) coordinate (#2-i1) ;
}

%% END INTERNAL

\newcommand\summerfour[1]{
\basicsummer{#1}
\splitfour{-0.5}{#1}
}

\newcommand\summerthree[1]{
\basicsummer{#1}
\splitthree{-0.5}{#1}
}

\newcommand\summertwo[1]{
\basicsummer{#1}
\splittwo{-0.5}{#1}
}

\newcommand\summerone[1]{
\basicsummer{#1}
\splitone{-0.5}{#1}
}


\newcommand\integratorfour[1]{
\basicintegrator{#1}
\splitfour{-0.8}{#1}
<}


\newcommand\integratorthree[1]{
\basicintegrator{#1}
\splitthree{-0.8}{#1}
}


\newcommand\integratortwo[1]{
\basicintegrator{#1}
\splittwo{-0.8}{#1}
}


\newcommand\integratorone[1]{
\basicintegrator{#1}
\splitone{-0.8}{#1}
}

%%%%% FIN MACROS




%%%%%%%%%%%%%%%%%%%%%%%%%%%%%%%%%%%%%%%%%%%%%%%%%
%
%
%                      Trucs d'Amaury
%
%
%%%%%%%%%%%%%%%%%%%%%%%%%%%%%%%%%%%%%%%%%%%%%%%%%


%----------------
%--- ref (AP) ---
%----------------

% \newcommand{\defref}[1]{Definition~\ref{#1}}
\newcommand{\lemref}[1]{Lemma~\ref{#1}}
\newcommand{\thref}[1]{Theorem~\ref{#1}}
\newcommand{\amaurycorref}[1]{Corollary~\ref{#1}}
 \newcommand{\figref}[1]{Figure~\ref{#1}}
% \newcommand{\figrefmult}[1]{Figures~{#1}}
% \newcommand{\secref}[1]{Section~\ref{#1}}
% %\renewcommand{\algref}[1]{Algorithm~\ref{#1}}
 \newcommand{\propref}[1]{Proposition~\ref{#1}}
% \newcommand{\exref}[1]{Example~\ref{#1}}
\newcommand{\remref}[1]{Remark~\ref{#1}}


 %-------------------
%--- polynomials ---
%-------------------

\newcommand{\sigmap}[1]{{\Sigma{#1}}}
\newcommand{\degp}[1]{{\operatorname{deg}(#1)}}
\newcommand{\poly}{\operatorname{poly}}


 %----------------------
%--- misc functions ---
%----------------------

\newcommand{\sgn}[1]{\operatorname{sgn}(#1)}
\newcommand{\intp}{\operatorname{int}}
 \newcommand{\fracp}{\operatorname{frac}}
\newcommand{\fiter}[2]{#1^{[#2]}}
\newcommand{\frestrict}[2]{#1\restriction_{#2}}
\newcommand{\indicator}[1]{\mathds{1}_{#1}}
\newcommand{\idfun}{\operatorname{id}}
% \newcommand{\floor}[1]{\left\lfloor#1\right\rfloor}
 \newcommand{\ceil}[1]{\left\lceil#1\right\rceil}
\newcommand{\round}[1]{\left\lfloor#1\right\rceil}
\DeclareMathOperator\erf{erf}


%------------
%--- sets ---
%------------

\newcommand{\A}{\mathbb{A}}
%\newcommand{\D}{\mathbb{D}}
\newcommand{\Ns}{\mathbb{N}^*}
%\newcommand{\C}{\mathbb{C}}
%\newcommand{\K}{\mathbb{K}}
% \MAT{height}{[width]}{[field]}
\newcommand{\MAT}[3]{M_{#1\ifthenelse{\equal{#2}{}}{}{,#2}}\ifthenelse{\equal{#3}{}}{}{\left(#3\right)}}

%\newcommand{\Rcomp}{\R_{c}}
\newcommand{\Rp}{\R_{+}}
\newcommand{\Rs}{\R^{*}}
\newcommand{\Rps}{\R_{+}^{*}}

%\newcommand{\Analytic}{C^\omega}
%\newcommand{\PTIME}{\ensuremath{\operatorname{P}}}
%\newcommand{\EXPTIME}{\ensuremath{\operatorname{EXPTIME}}}
%\newcommand{\NP}{\ensuremath{\operatorname{NP}}}
%\newcommand{\NPTIME}{\NP}
\newcommand{\FP}{\ensuremath{\operatorname{FP}}}
%\newcommand{\PSPACE}{\ensuremath{\operatorname{PSPACE}}}
% \newcommand{\FPSPACE}{\ensuremath{\operatorname{FPSPACE}}}



%%%%%%%%%%%%%%%%%%%%%%%%%%%%%%%%%%%%%%%%%%%%%%%%%
%
%
%                      Trucs de Quentin
%
%
%%%%%%%%%%%%%%%%%%%%%%%%%%%%%%%%%%%%%%%%%%%%%%%%%


\newcommand{\fct}[5]{\ensuremath{#1 : \left\{\begin{array}{c c c}
#2 & \rightarrow & #3\\
#4 & \mapsto & #5
                                             \end{array}\right.}}

%Nouvelle fraction, plus jolie
\newcommand{\f}[2]{	
	\mathchoice%
		{\dfrac{#1}{#2}}
    	{\dfrac{#1}{#2}}
		{\frac{#1}{#2}}
		{\frac{#1}{#2}}
}


%%Ensembles et combinatoires

%Ensemble avec propriété à drotie \{x | blabla\}
\newcommand{\enstq}[2]{\left\{ \left.#1\ \vphantom{#2}\right|\ #2  \right\}}
%meme chose mais avec des crochets
\newcommand{\crotq}[2]{\left[ \left.#1\ \vphantom{#2}\right|\ #2  \right]}
%même chose mais avec des brackets
\newcommand{\bratq}[2]{\left\langle \left.#1\ \vphantom{#2}\right|\ #2  
  \right\rangle}

\newcommand{\intn}[2]{\ensuremath{
		\left\llbracket\, #1\ ;\ #2\,\right\rrbracket
              }}


%%%%%%%%%%%%%%%%%%%%%%%%%%%%%%%%%%%%%%%%%%%%%%%%%
%
%
%               ENVIRONNEMENTS THEOREMS & CO
%
%
%%%%%%%%%%%%%%%%%%%%%%%%%%%%%%%%%%%%%%%%%%%%%%%%%


%----------------
%--- pas dans lipics ---
%-----------------


\ifnotSLIDE{
% commente car bug Fevrier 2026: \ifnotTDEXAM{\spnewtheorem{solution}{Solution}{\bfseries}{\rmfamily}}
\spnewtheorem{openproblem}{Open Problem}{\bfseries}{\rmfamily}
\spnewtheorem{remarkolivier}{Commentaire d'Olivier: }{\bfseries}{\rmfamily}
\spnewtheorem{postulate}{Postulate}{\bfseries}{\rmfamily}
}

%----------------
%--- dans lipics ---
%-----------------
%
\newenvironment{exercice}[1]{\begin{exercise} \label{exo:#1-stt}}{\end{exercise}}
\newenvironment{exerciced}[1]{\begin{exercisee}\label{exo:#1-stt}}{\end{exercisee}}
\newenvironment{exercicec}[1]{\begin{exercise} \label{exo:#1-stt} (\LANGUE{solution on}{corrigé} page \pageref{exo:#1-cor})
}{\end{exercise}}
\newenvironment{exercicedc}[1]{\begin{exercisee} \label{exo:#1-stt} (\LANGUE{solution on}{corrigé} page \pageref{exo:#1-cor})
}{\end{exercisee}}

\newenvironment{slideproposition}{{\bf Proposition}}{}

\ifnotSLIDE{
\nolipics{
\nolncs{\spnewtheorem{theorem}{\LANGUE{Theorem}{Théorème}}{\bfseries}{\rmfamily}}
\nolncs{\spnewtheoremi{definition}{\LANGUE{Definition}{Définition}}{\bfseries}{\rmfamily}{theorem}}
\nolncs{\spnewtheoremi{lemma}{\LANGUE{Lemma}{Lemme}}{\bfseries}{\rmfamily}{theorem}}
\nolncs{\spnewtheoremi{corollary}{\LANGUE{Corollary}{Corollaire}}{\bfseries}{\rmfamily}{theorem}}
\nolncs{\spnewtheoremi{proposition}{Proposition}{\bfseries}{\rmfamily}{theorem}}
\nolncs{\spnewtheoremi{example}{\LANGUE{Example}{Exemple}}{\bfseries}{\rmfamily}{theorem}}
%\spnewtheorem{sketch}{{\bf \LANGUE{Sketch of proof}{Démonstration (principe)}}}{\bfseries}{\rmfamily}
\nolncs{\spnewtheoremi{remark}{\LANGUE{Remark}{Remarque}}{\bfseries}{\rmfamily}{theorem}}
\nolncs{\spnewtheorem{exercise}{\LANGUE{Exercise}{Exercice}}{\bfseries}{\rmfamily}}
\spnewtheorem{exercisee}{\LANGUE{*Exercise}{*Exercice}}{\bfseries}{\rmfamily}
\spnewtheorem{summary}{\LANGUE{Summary}{Résumé}}{\bfseries}{\rmfamily}
 %\spnewtheorem{exerciceenglish}{Exercise}{\bfseries}{\rmfamily}
 %
 \newenvironment{sketch}{{\bf \LANGUE{Sketch of proof}{Démonstration (principe)}}:}{\hfill $\Box$ }
\nolncs{\spnewtheorem{conjecture}{\LANGUE{Conjecture}{Conjecture}}{\bfseries}{\rmfamily}}
 %%
 \spnewtheorem{theoremaprouver}{\LANGUE{Theorem TO BE PROVED}{Théorème (A PROUVER)}}{\bfseries}{\rmfamily}
  \spnewtheorem{conjectureaprouver}{\LANGUE{Conjecture TO BE VERIFIED}{Conjecture (A VERIFIER)}}{\bfseries}{\rmfamily}
 }}
            
 \makeatletter
 \ifnotSLIDE{
  \ifnotTDEXAM{
\@ifundefined{proof}{
	\newenvironment{proof}{{\bf \LANGUE{Proof}{Démonstration}}:}{\hfill $\Box$ 
	}{}
	}}}
\makeatother

 
%%%%%%%%%%%%%%%%%%%%%%%%%%%%%%%%%%%%%%%%%%%%%%%%%
%
%
%                      Macros dont pas fier du tout
%
%
%%%%%%%%%%%%%%%%%%%%%%%%%%%%%%%%%%%%%%%%%%%%%%%%%



\newcommand\equations[1]{
\begin{array}{lll}
#1
\end{array}
}



\newcommand\systeme[1]{
\left\{
\begin{array}{lll} 
#1
\end{array}
\right.
}



%------------------
%--- shorthands (AP)---
%------------------
\newcommand{\mtt}[1]{\mathtt{#1}}
\newcommand{\ovl}[1]{\overline{#1}}
\newcommand{\panic}[1]{\textcolor{red}{#1}}
%\NewEnviron{epanic}{{\color{red}\BODY}}
\newcommand{\idest}{\emph{i.e.\thinspace}}


%------------------
%--- mes trucs  ---
%------------------


\newcommand\implique{\Rightarrow}
\newcommand\ssi{\Leftrightarrow}
\newcommand\ac[1]{\{#1\}}
\newcommand\zero{\vectorl{0}}
\newcommand\un{\vectorl{1}}

\newcommand\technique[1]{{\scriptsize #1}}


%% Preuve and Co

\newcommand\sensdirect{ Direction $\Rightarrow$. }
\newcommand\sensindirect{ Direction $\Leftarrow$. }




\newcommand\donnee{\mbox{<donnée>}}



% pour aller vite

\newcommand\gM{\mathfrak{M}}
\newcommand\gMM{(M,f_1,\cdots,f_u,r_1,\cdots,r_v)}
%
\newcommand\mygM{\mK}
\newcommand\mygMM{\left ( \K, \{op_i\}_{i\in I}, r_1 \ldots,
r_\ell, \zero,\un \right )}
\newcommand\myM{\K} 

\newcommand\gMMM{(M,f_1,\cdots,f_u,c_1,\cdots,c_k,r_1,\cdots,r_v)}
\newcommand\gMMMM{(f_1,\cdots,f_u,r_1,\cdots,r_v)}
\newcommand\gMME{(M,f_1,\cdots,f_u,f'_1,\cdots,f'_\ell,r_1,\cdots,r_v)}
\newcommand\gMMS{(f_1,\cdots,f_u,f'_1,\cdots,f'_\ell,r_1,\cdots,r_v)}

\newcommand\bool{\mathfrak{B}}
\newcommand\gbool{(\{0,1\},\zero,\un,=)}

\newcommand\osg{\overline{sg}}
\newcommand\moins{\mathrel{\dot{-}}}


%%% TRUC DE BARWISE


\newcommand\UM{\kU_{\kM}}
\newcommand\VV{\mathbb{V}}
% Hereditarily finite
\newcommand\IHF{\mathbb{I}HF}
\newcommand\IHP{\IHYP}
\newcommand\IHYP{\mathbb{I}HYP}

\providecommand\citep[1]{\cite{#1}}

%% Cofinali


% probas
% défini dans amsmath
\renewcommand\Pr{\cp{Pr}}
\newcommand\Var{\cp{Var}}


% -- mesures
%\newcommand\card{card}
\newcommand\taille{\LANGUE{size}{taille}}
\newcommand\entree{e}
\newcommand\profondeur{\LANGUE{depth}{profondeur}}


%%%%%%%%%%%%%%%%%%%%%%%%%%%%%%%%%%%%%%%%%%%%%%%%%%%%%%%%%%%%%%%%%%%%%%%%%%%%
%                 MACHINES DE TURING (graphique)
%%%%%%%%%%%%%%%%%%%%%%%%%%%%%%%%%%%%%%%%%%%%%%%%%%%%%%%%%%%%%%%%%%%%%%%%%%%%



% Configuration de machine de Turing
\newcommand\configTuring[3]{#1\textbf{#2}#3}


%%% DESSIN CROIX PORU AIDER  A DESSINER

\newcommand\croix{  +(-1,-1) -- +(1,1) +(-1,1) --
  +(1,-1) }

%% DESSIN ACCOLADE

\newcommand\ACCOLADE[3]{
%#1: coordonnées de départ
%#2: coordonnées arrivée
%#3: texte
\draw[decorate,decoration={brace}]  (#1) -- (#2)
 node[midway,above=0.1cm] {#3};
}

%%%


\def\lcase{3.5 ex}
\def\hcase{3.5ex}
\tikzstyle{case} = []
%[draw,minimum width = \lcase,minimum height = 3.6ex,baseline]
\newcommand\dcase{ +(-0.5*\lcase,-0.5*\hcase) rectangle +(0.5*\lcase,0.5*\hcase)}
%\newcommand\lettre[1]{\node [name=l,case,on chain=1] {#1};}


\newcommand\RUBANMACHINEDETURING[6]{
%#1: texte
%#2: position de la tête
%#3: nb de blancs à gauche
%#4: nb de cases totales
%#5: texte au dessus du ruban
%#6: position du coin
    \draw (#6) +(\lcase,4ex ) node {#5};
    \draw (#6) +(-0.6,0)  node  {\ldots};  
    \draw (#6) node[coordinate] (t) {};
    \foreach \x in {1,2,...,#3} {
         \draw (t) node [case]   {$\blanc$} \dcase;
         \draw (t) ++(\lcase,0) node [coordinate] (t) {};
       };
   \StrLen{#1}[\Resultat] % ATTENTION IL FAUT CETTE SYNTAXE POUR FAIRE
                         % EQUIVALENT D UN EDEF
  
\foreach \x in {1,...,\Resultat} 
             { 
        \draw (t) node [case]   {\StrChar{#1}{\x}} \dcase;
        \draw (t) ++(\lcase,0) node [coordinate] (t) {};
             }
\pgfmathtruncatemacro{\z}{#4 - \Resultat - #3}

  \foreach \x in {1,2,..., \z} {
        \draw (t) node [case]   {$\blanc$} \dcase;
        \draw (t) ++(\lcase,0) node [coordinate] (t) {};
}
%% Convention: graphiquement, la case numéro n se trouve à (n*\lcase,0)
%% en comptant les case à partir de $0$

   \draw  (t) +(0.1,0) node [name=r] {\ldots}; 
}

\newcommand\RUBANMACHINEDETURINGtexte[6]{
%#1: texte
%#2: position de la tête
%#3: nb de blancs à gauche
%#4: nb de cases totales (équivalent)
%#5: texte au dessus du ruban
%#6: position du coin
    \draw (#6) +(\lcase,4ex ) node {#5};
    \draw (#6) +(-0.6,0)  node  {\ldots};  
    \draw (#6) node[coordinate] (t) {};
   \foreach \x in {1,2,...,#3} {
        \draw (t) node [case]   {$\blanc$} \dcase;
        \draw (t) ++(\lcase,0) node [coordinate] (t) {};
      };
   \draw (#6) ++(-2*\lcase,0.5*\lcase)-- +(\lcase*#4,0);
   \draw (#6) ++(-2*\lcase,-0.5*\lcase) -- +(\lcase*#4,0);
   
    \draw (t) ++(-0.5*\lcase,0) node [case,anchor=west]  (te) {#1};
    \draw (te.east)  node [anchor=west,coordinate] (t) {};

  \draw (t) ++(0.2*\lcase,0) node [coordinate] (t) {};
   \foreach \x in {1,2,...,#3} {
        \draw (t) node [case]   {$\blanc$} \dcase;
        \draw (t) ++(\lcase,0) node [coordinate] (t) {};
      };

    \draw  (t) +(0.1,0) node [anchor=west] {\ldots}; 
}


\newcommand\RUBANMACHINEDETURINGTETE[6]{
%
% dessine la tete
%
% output: (tete)
%
\RUBANMACHINEDETURING{#1}{#2}{#3}{#4}{#5}{#6}
% TETE
     \draw (#6) + ( #2 * \lcase + #3  *\lcase,- 0.5*\lcase )
     node[coordinate] (k) {};
     \draw[->] (k) +(0,-0.5) -- (k);
     \draw (k) +(0,-0.5) node[coordinate] (tete) {};
}


\newcommand\RUBANMACHINEDETURINGTETEtexte[6]{
%
% dessine la tete
%
% output: (tete)
%
\RUBANMACHINEDETURINGtexte{#1}{#2}{#3}{#4}{#5}{#6}
% TETE
     \draw (#6) + ( #2 * \lcase + #3  *\lcase,- 0.5*\lcase )
     node[coordinate] (k) {};
     \draw[->] (k) +(0,-0.5) -- (k);
     \draw (k) +(0,-0.5) node[coordinate] (tete) {};
}


\newcommand\MACHINEDETURING[6]{
%#1: texte
%#2: position de la tête
%#3: état
%#4: nb de blancs à gauche
%#5: nb de cases totales
%#6: texte en haut du ruban
%\begin{tikzpicture}[node distance=-0.15mm]

\RUBANMACHINEDETURINGTETE{#1}{#2}{#4}{#5}{#6}{0,0}
%
%% ETAT:
     \draw (tete) node [name=etat,draw,circle,anchor=north]  {#3};
     \draw (etat) -- (tete);
%\end{tikzpicture}
}

\newcommand\MACHINEDETURINGtexte[6]{
%#1: texte
%#2: position de la tête
%#3: état
%#4: nb de blancs à gauche
%#5: nb de cases totales
%#6: texte en haut du ruban
%\begin{tikzpicture}[node distance=-0.15mm]

\RUBANMACHINEDETURINGTETEtexte{#1}{#2}{#4}{#5}{#6}{0,0}
%
%% ETAT:
     \draw (tete) node [name=etat,draw,circle,anchor=north]  {#3};
     \draw (etat) -- (tete);
%\end{tikzpicture}
}

\newcommand\MACHINEDETURINGTROISRUBANS[9]{
%#1: texte1
%#2: position de la tête 1
%#3: texte2
%#4: position de la tête 2
%#5: texte3
%#6: position de la tête 3
%#7: état
%#8: nb de blancs à gauche 1
%#9: nb de cases totales
%\begin{tikzpicture}[node distance=-0.15mm]

\RUBANMACHINEDETURINGTETE{#1}{#2}{#8}{#9}{\LANGUE{tape}{ruban} $1$}{0,0}
\draw (tete) node[coordinate] (tete1) {};
\draw (tete1) +(8,0) node[coordinate] (tete1b) {};
\draw (tete1) -| (tete1b) ;

\RUBANMACHINEDETURINGTETE{#3}{#4}{#8}{#9}{\LANGUE{tape}{ruban} $2$}{0,-2}
\draw (tete) node[coordinate] (tete2) {};
\draw (tete2) +(-8.5,0) node[coordinate] (tete2b) {};
\draw (tete2) -| (tete2b) ;

\RUBANMACHINEDETURINGTETE{#5}{#6}{#8}{#9}{\LANGUE{tape}{ruban} $3$}{0,-4}
\draw (tete) node[coordinate] (tete3) {};

%% ETAT:
     \draw (tete3) node [name=etat,draw,circle,anchor=north]  {#7};
     \draw (etat) -- (tete3);
    \draw  (etat) -|  (tete2b);
\draw  (etat) -|  (tete1b);
%\end{tikzpicture}
}

\newcommand\MACHINEDETURINGTROISRUBANSENGLISH[9]{
%#1: texte1
%#2: position de la tête 1
%#3: texte2
%#4: position de la tête 2
%#5: texte3
%#6: position de la tête 3
%#7: état
%#8: nb de blancs à gauche 1
%#9: nb de cases totales
%\begin{tikzpicture}[node distance=-0.15mm]

\RUBANMACHINEDETURINGTETE{#1}{#2}{#8}{#9}{tape $1$}{0,0}
\draw (tete) node[coordinate] (tete1) {};
\draw (tete1) +(8,0) node[coordinate] (tete1b) {};
\draw (tete1) -| (tete1b) ;

\RUBANMACHINEDETURINGTETE{#3}{#4}{#8}{#9}{tape $2$}{0,-2}
\draw (tete) node[coordinate] (tete2) {};
\draw (tete2) +(-8.5,0) node[coordinate] (tete2b) {};
\draw (tete2) -| (tete2b) ;

\RUBANMACHINEDETURINGTETE{#5}{#6}{#8}{#9}{tape $3$}{0,-4}
\draw (tete) node[coordinate] (tete3) {};

%% ETAT:
     \draw (tete3) node [name=etat,draw,circle,anchor=north]  {#7};
     \draw (etat) -- (tete3);
    \draw  (etat) -|  (tete2b);
\draw  (etat) -|  (tete1b);
%\end{tikzpicture}
}



\newcommand\MACHINEDETURINGTROISRUBANStexte[9]{
%#1: texte1
%#2: position de la tête 1
%#3: texte2
%#4: position de la tête 2
%#5: texte3
%#6: position de la tête 3
%#7: état
%#8: nb de blancs à gauche 1
%#9: nb de cases totales
% \begin{tikzpicture}[node distance=-0.15mm]

\RUBANMACHINEDETURINGTETEtexte{#1}{#2}{#8}{#9}{\LANGUE{tape}{ruban} $1$}{0,0}
\draw (tete) node[coordinate] (tete1) {};
\draw (tete1) +(8,0) node[coordinate] (tete1b) {};
\draw (tete1) -| (tete1b) ;

\RUBANMACHINEDETURINGTETEtexte{#3}{#4}{#8}{#9}{\LANGUE{tape}{ruban} $2$}{0,-2}
\draw (tete) node[coordinate] (tete2) {};
\draw (tete2) +(-8.5,0) node[coordinate] (tete2b) {};
\draw (tete2) -| (tete2b) ;

\RUBANMACHINEDETURINGTETEtexte{#5}{#6}{#8}{#9}{\LANGUE{tape}{ruban} $3$}{0,-4}
\draw (tete) node[coordinate] (tete3) {};

%% ETAT:
     \draw (tete3) node [name=etat,draw,circle,anchor=north]  {#7};
     \draw (etat) -- (tete3);
    \draw  (etat) -|  (tete2b);
\draw  (etat) -|  (tete1b);
%\end{tikzpicture}
}


\newcommand\MACHINEDETURINGDEUXRUBANStexte[9]{
%#1: texte1  
%#2: position de la tête 1 
%#3: texte2 inutilisée
%#4: position de la tête 2 inutilisée
%#5: texte3
%#6: position de la tête 3
%#7: état
%#8: nb de blancs à gauche 1
%#9: nb de cases totales
% \begin{tikzpicture}[node distance=-0.15mm]

\RUBANMACHINEDETURINGTETEtexte{#1}{#2}{#8}{#9}{\LANGUE{tape}{ruban} $1$}{0,0}
\draw (tete) node[coordinate] (tete1) {};
\draw (tete1) +(8,0) node[coordinate] (tete1b) {};
\draw (tete1) -| (tete1b) ;

% \RUBANMACHINEDETURINGTETEtexte{#3}{#4}{#8}{#9}{\LANGUE{tape}{ruban} $2$}{0,-1.5}
% \draw (tete) node[coordinate] (tete2) {};
% \draw (tete2) +(-5.5,0) node[coordinate] (tete2b) {};
% \draw (tete2) -| (tete2b) ;

\RUBANMACHINEDETURINGTETEtexte{#5}{#6}{#8}{#9}{\LANGUE{tape}{ruban} $2$}{0,-1.5}
\draw (tete) node[coordinate] (tete3) {};

%% ETAT:
     \draw (tete3) node [name=etat,draw,circle,anchor=north]  {#7};
     \draw (etat) -- (tete3);
%   \draw  (etat) -|  (tete2b);
\draw  (etat) -|  (tete1b);
%\end{tikzpicture}
}


\def\argdix{2}
\def\argonze{25}
\newcommand\MACHINEDETURINGQUATRERUBANStextep[9]{
%#1: texte1
%#2: position de la tête 1
%#3: texte2
%#4: position de la tête 2
%#5: texte3
%#6: position de la tête 3
%#7: texte4
%#8: position de la tête 3
%#9: état
%#10: nb de blancs à gauche 1
%#11: nb de cases totales
% \begin{tikzpicture}[node distance=-0.15mm]

%% BUG sur #10 devenu \argdix
%% BUG SUR #11 devenu \argonze

\RUBANMACHINEDETURINGTETEtexte{#1}{#2}{\argdix}{\argonze}{\LANGUE{tape}{ruban} $1$}{0,0}
% \draw (tete) node[coordinate] (tete1) {};
% \draw (tete1) +(8,0) node[coordinate] (tete1b) {};
% \draw (tete1) -| (tete1b) ;

\RUBANMACHINEDETURINGTETEtexte{#3}{#4}{\argdix}{\argonze}{\LANGUE{tape}{ruban} $2$}{0,-1.5}
% \draw (tete) node[coordinate] (tete2) {};
% \draw (tete2) +(-5.5,0) node[coordinate] (tete2b) {};
% \draw (tete2) -| (tete2b) ;

\RUBANMACHINEDETURINGTETEtexte{#5}{#6}{\argdix}{\argonze}{\LANGUE{tape}{ruban} $3$}{0,-3.2}
\draw (tete) node[coordinate] (tete3) {};

\RUBANMACHINEDETURINGTETEtexte{#7}{#8}{\argdix}{\argonze}{\LANGUE{tape}{ruban} $4$}{0,-4.9}
\draw (tete) node[coordinate] (tete4) {};

%% ETAT:
     \draw (tete4) node [name=etat,draw,circle,anchor=north]  {#9};
%      \draw (etat) -- (tete3);
%     \draw  (etat) -|  (tete2b);
% \draw  (etat) -|  (tete1b);
%\end{tikzpicture}
}


%%%%%%%%%%%%%%%%%%%%%%%%%%%%%%%%%%%%%%%%%%%%%%%%%%%%%%%%%%%%%%%%%%%%%%%%%%%%
%                 MACHINES DE TURING (automate)
%%%%%%%%%%%%%%%%%%%%%%%%%%%%%%%%%%%%%%%%%%%%%%%%%%%%%%%%%%%%%%%%%%%%%%%%%%%%



\newcommand\dec{1ex}
\newcommand\myACCOLADE[4]{
\ACCOLADE{#1*\lcase+#4*\lcase-0.5*\lcase,0.5*\lcase+\dec}{#1*\lcase+#4*\lcase-0.5*\lcase+#3*\lcase,0.5*\lcase+\dec}{#2}
}


\newcommand\M{\Sigma}


\newcommand\sommet[2]{\node[state] (#1) #2 {$#1$};}
\newcommand\sommetinitial[2]{\node[state,initial] (#1) #2 {$#1$};}
\newcommand\sommetfinal[2]{\node[state,accepting] (#1) #2 {$#1$};}
\newcommand\gtransition[5]{\draw (#1) edge node {#2/#4 $#5$} (#3);}
\newcommand\gtransitionba[5]{\draw (#1) [loop above] edge node {#2/#4
    $#5$} (#3);}
\newcommand\gtransitionbb[5]{\draw (#1) [loop below] edge node {#2/#4
    $#5$} (#3);}
\newcommand\gtransitionbleft[5]{\draw (#1) [loop left] edge node {#2/#4
    $#5$} (#3);}
\newcommand\gtransitionbright[5]{\draw (#1) [loop right] edge node {#2/#4
    $#5$} (#3);}
\newcommand\gtransitionbl[5]{\draw (#1) [bend left] edge node {#2/#4
    $#5$} (#3);}
\newcommand\gtransitionbr[5]{\draw (#1) [bend right] edge node {#2/#4
    $#5$} (#3);}
\newcommand\gtransitionbadouble[9]{\draw (#1) [loop above] edge node {
\begin{tabular}{l}
#2/#4 $#5$ \\
  #7/#8 $#9$ \\
\end{tabular}
} (#3);}

\newcommand\SCALE{}

\newcommand\dessinmtp[2]{
\begin{center}
\begin{tikzpicture}[->,>=stealth',shorten >=1pt,auto,node distance=2.3cm,semithick,#2] 
#1
\end{tikzpicture}
\end{center}
}

\newcommand\dessinmt[1]{\dessinmtp{#1}{scale=1}}


\newcommand\ANIMATIONMT[1]{
\begin{overprint}
\begin{dessinp}{scale=0.52}
\foreach \av/\q/\b  [count=\xi]
in  
{#1}
{ 
\only<\xi| handout 0>{ 
 \StrLen{\av}[\Resintm] % ATTENTION IL FAUT CETTE SYNTAXE POUR FAIRE
                         % EQUIVALENT D UN EDEF
 \MACHINEDETURING{\av \b}{\Resintm}{$\q$}{5}{20}{}
}}
 \end{dessinp}
\end{overprint}
}

\newcommand\ANIMATIONMTd[1]{
\begin{overprint}
\begin{dessinp}{scale=0.52}
\foreach \av/\q/\b  [count=\xi]
in  
{#1}
{ 
\only<\xi| handout 0>{ 
 \StrLen{\av}[\Resintm] % ATTENTION IL FAUT CETTE SYNTAXE POUR FAIRE
                         % EQUIVALENT D UN EDEF
 \hspace{-1cm}\MACHINEDETURING{\av \b}{\Resintm}{$\q$}{5}{33}{}
}}
 \end{dessinp}
\end{overprint}
}

\def\MAXDET{20}

\newcommand\DIAGRAMMEESPACETEMPS[1]{
\draw node[coordinate] (tt) {};
\foreach \av/\q/\b  in  
{#1}
{ 

\RUBANMACHINEDETURING{\av{$\q$}\b}{0}{4}{\MAXDET}{}{tt};
\draw (tt) +(0,-\hcase) node[coordinate] (tt) {};
}
%\foreach \x in {0,...,19} {\draw (tt) + (\x*\lcase,0) node {$\vdots$};}
}



\newcommand\progun{
/q_0/0011, % initial
X/q_1/011,X0/q_1/11,X/q_2/0Y1,/q_2/X0Y1,
X/q_0/0Y1,XX/q_1/Y1,XXY/q_1/1, XX/q_2/YY, X/q_2/XYY,
XX/q_0/YY,XXY/q_3/Y,XXYY/q_3/B,XXYYB/q_4/B
}

\newcommand\DIAGRAMMEESPACETEMPSun{
\DIAGRAMMEESPACETEMPS
%% RECOPIE progun
{/q_0/0011, % initial
X/q_1/011,X0/q_1/11,X/q_2/0Y1,/q_2/X0Y1,
X/q_0/0Y1,XX/q_1/Y1,XXY/q_1/1, XX/q_2/YY, X/q_2/XYY,
XX/q_0/YY,XXY/q_3/Y,XXYY/q_3/B,XXYYB/q_4/B}
%% FIN RECOPIE
}

\newcommand\progdeux{
/q_0/0010, X/q_1/010, X0/q_1/10, X/q_2/0Y0,/q_2/X0Y0,
X/q_0/0Y0,XX/q_1/Y0,XXY/q_1/0,XXY0/q_1/%B,
}


\newcommand\DIAGRAMMEESPACETEMPSdeux{
\DIAGRAMMEESPACETEMPS
{
%% RECOPIE progdeux
/q_0/0010, X/q_1/010, X0/q_1/10, X/q_2/0Y0,/q_2/X0Y0,
X/q_0/0Y0,XX/q_1/Y0,XXY/q_1/0,XXY0/q_1/%B,
%% FIN RECOPIE
}
}

\newcommand\DIAGRAMMEESPACETEMPSdeuxvariante{
\def\memolcase{\lcase}
\def\memoMAXDET{20}
\def\MAXDET{15}
\def\lcase{6 ex}
\DIAGRAMMEESPACETEMPS
{
%% RECOPIE progdeux
/q_00/010, X/q_10/10, X0/q_11/0, X/q_20/Y0,/q_2X/0Y0,
X/q_00/Y0,XX/q_1Y/0,XXY/q_10/,XXY0/q_1B/%B,
%% FIN RECOPIE
}
\def\lcase{\memolcase}
\def\MAXDET{\memoMAXDET}
}


\newcommand\ANIMATIONun{
\ANIMATIONMT
%% RECOPIE progun
{/q_0/0011, % initial
X/q_1/011,X0/q_1/11,X/q_2/0Y1,/q_2/X0Y1,
X/q_0/0Y1,XX/q_1/Y1,XXY/q_1/1, XX/q_2/YY, X/q_2/XYY,
XX/q_0/YY,XXY/q_3/Y,XXYY/q_3/B,XXYYB/q_4/B}
%% FIN RECOPIE
}



\newcommand\DESSINALGO[3]{
% argument 1= ou
% argument 2= texte 
% argument 3= nom du noeud ex m
\draw node[draw,#1,minimum size=1cm] (#3) {#2};

% pour fitting box
\draw (#3.west) +(-0.2cm,0) node (#3c2) {};
\draw (#3.north) +(0,+0.2cm) node (#3c3) {};
\draw (#3.south) +(0,-0.2cm) node (#3c4) {};
}

\newcommand\DESSINALGOa[3]{
\DESSINALGO{#1}{#2}{#3}
\draw (#3.east) +(-0.1cm,0.2cm) node(#3a) {};
\draw[->] (#3a) -- ++(0.5cm,0) node (#3aa) {};
\draw (#3aa) node[anchor=west] (#3aaa) {\LANGUE{Accept}{Accepte}};
\draw (#3aaa.east) node(#3aaaa) {};
}

\newcommand\DESSINALGOe[4]{
\DESSINALGO{#1}{#2}{#3}
\draw (#3.east) +(-0.1cm,0.2cm) node(#3a) {};
\draw[->] (#3a) -- ++(0.5cm,0) node (#3aa) {};
\draw (#3aa) node[anchor=west] (#3aaa) {#4};
\draw (#3aaa.east) node(#3aaaa) {};
}


\newcommand\DESSINALGOar[3]{
\DESSINALGOa{#1}{#2}{#3}
\draw (#3.east) +(-0.1cm,-0.2cm) node(#3r) {};
\draw[->] (#3r) -- ++(0.5cm,0) node (#3rr) {};
\draw (#3rr) node[anchor=west] (#3rrr) {\LANGUE{Reject}{Refuse}};
\draw (#3rrr.east) node(#3rrrr) {};
}


\newcommand\WINDOWarg[7] {
%#1: position
%#2..7: contenu du tableau
\draw (#1) node[case]  {#2} \dcase;
\draw (#1) ++ (\lcase,0) node[case] {#3} \dcase;
\draw (#1) ++ (2*\lcase,0) node[case] {#4} \dcase;
\draw (#1) ++ (0,-\lcase) node[case] {#5} \dcase;
\draw (#1) ++ (\lcase,-\lcase) node[case] {#6} \dcase;
\draw (#1) ++ (2*\lcase,-\lcase) node[case] {#7} \dcase;
}

\newcommand\WINDOW[3]{
%#1: positioin
%#2: contenu
%#3: étiquette
\WINDOWarg{#1}
{\StrChar{#2}{1}}{\StrChar{#2}{2}}{\StrChar{#2}{3}}
{\StrChar{#2}{4}}{\StrChar{#2}{5}}{\StrChar{#2}{6}}
\draw (#1) +(-\lcase*1.2,-\lcase*0.5) node {#3};
}





%%%%%%%%%%%%%%%%%%%%%%%%%%%%%%%%%%%%%%%%%%%%%%%%%%%%%%%%%%%%%%%%%%%%%%%%%%%%
%                 LISTINGS
%%%%%%%%%%%%%%%%%%%%%%%%%%%%%%%%%%%%%%%%%%%%%%%%%%%%%%%%%%%%%%%%%%%%%%%%%%%%


% LISTINGS
\RequirePackage{listings}
 \lstset{language=Java,
 mathescape=true,
 showspaces=false,
 showstringspaces=false,
 frame=single,
morekeywords={recursive,then,div,function,integer,read,out,else,par,endpar,seq,endseq,and,not,read,write,repeat,until,undef,let,to,endfor,endif,endwhile,guess,in,parallel},
% basicstyle=\ttfamily\small,
% basewidth={1.1ex},
% keywordstyle=\bf\color{blue},
% %stringstyle=\bf\ttfamily\color{green},
% commentstyle=\bfseries\color{magenta},
 literate={:=}{{$\gets$ }}1 {<=}{{$\leq$}}2  {>=}{{$\geq$}}2
 {!=}{{$\neq$}}1 
 }
% \lstset{escapechar=\%}
\let\lst\lstinline
\def\beginlisting{\begin{lstlisting}}
\def\endlisting{\end{lstlisting}}  




\newenvironment{algo}[1]{
  \lstset{escapechar=\@}
  \def\input{$#1$}
  \def\out{out}
  \def\BEGIN{\textbf{read} \input }
 \def\BEGINN{\textbf{repeat} }
  \def\END{\textbf{until out!=undef}}
  \def\ENDD{\textbf{write \out}}
  \def\ENDT{\textbf{until $\ctl$ = $q_a$ or $\ctl$ = $q_r$}}
  \def\ENDR{\textbf{until $\ctl$ = $q_a$}}
  \def\ENDP{\textbf{until $\ctl$ = $q_a$}}
}{}
\newenvironment{algon}[1]{
  \lstset{numbers=left}
  \begin{algo}{#1}
}{\end{algo}}
\newcommand\FSM[3]{FSM(#1,\lst{#2},#3)}
\newcommand\FSMi[3]{FSM(#1,{#2},#3)}




%%%%%%%%%%%%%%%%%%%%%%%%%%%%%%%%%%%%%%%%%%%%%%%%%
%
%
%                      TRADUCTION EN COURS
%
%%%%%%%%%%%%%%%%%%%%%%%%%%%%%%%%%%%%%%%%%%%%%%%%%

\newcommand\ENTREPOT[1]{\NDLR{ENTREPOT ICI: #1}}
\newcommand\scite[2][]{\cite[#1]{#2}}
\newcommand\nospellcheck[1]{{#1}}
\newcommand\nd{nd}
\newcommand\NL{
 
}

\newenvironment{FRANCAIS}{
\small
\color{orange}
}
{\normalsize
\color{black}
}


\newenvironment{SEMIFRANCAIS}{
\small
\color{violet}
}
{\normalsize
\color{black}
}

\newcommand\QFRANCAIS[1]{
\begin{FRANCAIS}
#1
\end{FRANCAIS}
}


%%%%%%%%%%%%%%%%%%%%%%%%%%%%%%%%%%%%%%%%%%%%%%%%%
%
%
%                      Wikipedia
%
%%%%%%%%%%%%%%%%%%%%%%%%%%%%%%%%%%%%%%%%%%%%%%%%%

\providecommand\tightlist{}
\providecommand\Complex{\C}
\newenvironment{myalgie}{\begin{array}{lll}}{\end{array}}

\newcommand\nl{\ \\ }




%%%%%%%%%%%%%%%%%%%%%%%%%%%%%%%%%%%%%%%%%%%%%%%%%
%
%
%                      Pour citations précises
%
%%%%%%%%%%%%%%%%%%%%%%%%%%%%%%%%%%%%%%%%%%%%%%%%%



\newcommand\citeprecis[2]{{\cite[#1]{#2}}}
\newcommand\citeprecisinv[2]{\citeprecis{#2}{#1}}

\usepackage{csquotes}



%%%%%%%%%%%%%%%%%%%%%%%%%%%%%%%%%%%%%%%%%%%%%%%%%
%
%
%                      Travail sur couleurs
%
%   toutes les réponses pointent dans le même sens, et le schéma correspond exactement 
%. au profil perceptif d’un protanope** (déficit des cônes L / longueurs d’onde longues)
%  MAIS SEVERE
%
%%%%%%%%%%%%%%%%%%%%%%%%%%%%%%%%%%%%%%%%%%%%%%%%%


% ================================
% GROUPE 1 — Verts sombres
% ================================
\definecolor{deepgreen}{RGB}{0,160,40}
\definecolor{forestgreen}{RGB}{0,150,70}
\definecolor{mossgreen}{RGB}{0,140,90}
\definecolor{greenwave}{RGB}{0,150,90}
\definecolor{forestmist}{RGB}{0,160,90}
\definecolor{darkteal}{RGB}{0,130,70}

% ================================
% GROUPE 2 — Teal / verts bleutés
% ================================
\definecolor{seagreen}{RGB}{0,130,100}
\definecolor{freshteal}{RGB}{0,140,130}
\definecolor{greenteal}{RGB}{0,140,120}
\definecolor{blueteal}{RGB}{0,150,120}      % ancien "teal" (doublon corrigé)
\definecolor{lagoon}{RGB}{0,180,120}
\definecolor{greenmist}{RGB}{0,170,110}

% ================================
% GROUPE 3 — Turquoises / Aqua
% ================================
\definecolor{springgreen}{RGB}{0,190,100}
\definecolor{mintgreen}{RGB}{10,200,150}
\definecolor{oceanfoam}{RGB}{0,200,150}
\definecolor{aquagreen}{RGB}{30,210,170}
\definecolor{coastalaqua}{RGB}{40,210,180}
\definecolor{softaqua}{RGB}{40,210,190}

% ================================
% GROUPE 4 — Cyan / Aqua clairs
% ================================
\definecolor{emeraldteal}{RGB}{0,180,140}
\definecolor{deepcyan}{RGB}{0,180,160}
\definecolor{lightcyan}{RGB}{0,200,160}
\definecolor{aquamarine}{RGB}{0,200,160}
\definecolor{lightaqua}{RGB}{60,230,200}
\definecolor{mistblue}{RGB}{80,240,210}

% ================================
% GROUPE 5 — Bleus clairs
% ================================
\definecolor{paleaqua}{RGB}{80,200,190}
\definecolor{glacier}{RGB}{90,150,200}
\definecolor{skyblueA}{RGB}{60,150,220}      % ancien "skyblue"
\definecolor{lightsky}{RGB}{90,170,230}
\definecolor{iceblue}{RGB}{120,190,240}
\definecolor{frost}{RGB}{150,210,250}

% ================================
% GROUPE 6 — Bleus moyens
% ================================
\definecolor{steelblue}{RGB}{30,100,160}
\definecolor{cadetblue}{RGB}{70,120,180}
\definecolor{coolblue}{RGB}{110,170,220}
\definecolor{morningblue}{RGB}{140,200,240}
\definecolor{blueiceA}{RGB}{40,150,210}       % ancien "blueice" (doublon corrigé)
\definecolor{blueair}{RGB}{60,170,225}

% ================================
% GROUPE 7 — Azurs / Bleu saturé
% ================================
\definecolor{tealblue}{RGB}{0,110,130}
\definecolor{deepbluecyan}{RGB}{0,110,170}
\definecolor{azur}{RGB}{0,70,180}
\definecolor{brightazure}{RGB}{20,130,200}
\definecolor{lightazure}{RGB}{80,180,235}
\definecolor{skylight}{RGB}{100,200,245}

% ================================
% GROUPE 8 — Bleus sombres / violets
% ================================
\definecolor{nightblue}{RGB}{40,0,110}
\definecolor{slateblue}{RGB}{20,40,120}
\definecolor{violetA}{RGB}{60,50,160}         % ancien "violet" (distinct de violetB ci-dessous)
\definecolor{blueviolet}{RGB}{80,60,180}
\definecolor{purpleA}{RGB}{100,70,200}        % ancien "purple"
\definecolor{lightpurple}{RGB}{130,90,210}

% ================================
% GROUPE 9 — Violets saturés / indigos
% ================================
\definecolor{lilac}{RGB}{160,110,220}
\definecolor{indigoA}{RGB}{50,0,130}          % ancien "indigo"
\definecolor{deepindigo}{RGB}{70,10,150}
\definecolor{royalpurple}{RGB}{90,20,170}

% ================================
% GROUPE 10 — Fermeture cercle perceptif
% ================================
\definecolor{coolmint}{RGB}{0,130,70}
\definecolor{deepsea}{RGB}{0,130,70}

% ================================
% PALETTE DEUTÉRANOPE — noms uniques
% ================================
\definecolor{crimson}{RGB}{200,30,60}
\definecolor{brickred}{RGB}{180,40,40}
\definecolor{coral}{RGB}{220,90,70}
\definecolor{salmon}{RGB}{240,130,120}

\definecolor{magenta}{RGB}{200,40,150}
\definecolor{fuchsia}{RGB}{220,60,170}
\definecolor{violetB}{RGB}{140,70,190}        % ancien "violet"
\definecolor{indigoB}{RGB}{90,50,180}         % ancien "indigo"

\definecolor{royalblueA}{RGB}{50,80,200}      % ancien "royalblue"
\definecolor{azure}{RGB}{30,120,230}
\definecolor{skyblueB}{RGB}{80,160,230}       % second "skyblue"
\definecolor{tealB}{RGB}{0,150,170}           % second "teal"

\definecolor{marine}{RGB}{0,80,140}




%% Alternative

% === Générateur automatique fill/text pour protanopes ===
% Usage : \ProtaColorPair{mistblue}
% Suppose que mistblue est déjà défini par \definecolor{mistblue}{RGB}{...}

\newcommand{\ProtaColorPair}[1]{%
  % fill = couleur originale
  \colorlet{#1fill}{#1}%
  % text = version assombrie idéalement lisible
  \colorlet{#1text}{#1!80!black}%
}

%Utilisation
%\ProtaColorPair{mistblue}

%\textcolor{mistbluefill}{Couleur claire} \quad
%\textcolor{mistbluetext}{Couleur pour texte}

	\foreach \i/\name in {
	  1/deepgreen,
	  2/greenwave,
	  3/lagoon,
	  4/oceanfoam,
	  5/softaqua,
	  6/mistblue,
 	 7/lightsky,
	  8/blueiceA,   % <-- corrigé
	  9/deepcyan,
	  10/azur,
	  11/blueviolet,
	  12/royalpurple}
	  {\ProtaColorPair{\name}}
	  


\newcommand\ROUEPROTANOPEDOUZE{
	% Roue chromatique protanope — 12 couleurs
	\begin{tikzpicture}[scale=3]

	\foreach \i/\name/\nametext in {
	  1/deepgreen/deepgreen,
	  2/greenwave/greenwave,
	  3/lagoon/lagoontext,
	  4/oceanfoam/oceanfoam,
	  5/softaqua/softaquatext,
	  6/mistblue/mistbluetext,
 	 7/lightsky/lightskytext,
	  8/blueiceA/blueiceA,   % <-- corrigé
	  9/deepcyan/blueiceA,
	  10/azur/azur,
	  11/blueviolet/blueviolet,
	  12/royalpurple/royalpurple}
	{
	  % secteur coloré
 	 \filldraw[fill=\name, draw=black, line width=0.2pt]
	    ({90-30*\i}:1)
	    -- ({90-30*(\i-1)}:1)
	    -- (0,0) -- cycle;

 	 % label dans la bonne couleur
	  \node at ({90-30*(\i-0.5)}:1.25)
	    {\textcolor{\nametext}{\small \nametext}};
	}

	\end{tikzpicture}
}


\newcommand\ROUEPROTANOPEVINGT{
\begin{tikzpicture}[scale=3]

\foreach \i/\name/\nametext in {
  1/deepgreen/deepgreen,
  2/greenwave/greenwave,
  3/lagoon/lagoon,
  4/oceanfoam/oceanfoam,
  5/softaqua/softaquatext,
  6/mistblue/mistbluetext,
  7/lightsky/lightsky,
  8/blueiceA/blueiceA,      % corrigé
  9/deepcyan/deepcyan,
  10/azur/azur,
  11/brightazure/brightazure,
  12/skyblueA/skyblueA,     % corrigé
  13/coolblue/coolblue,
  14/glacier/glacier,
  15/marine/marine,
  16/violetA/violetA,      % corrigé
  17/blueviolet/blueviolet,
  18/purpleA/purpleA,      % corrigé
  19/royalpurple/royalpurple,
  20/indigoA/indigoA}      % corrigé
{
  % Secteur coloré
  \filldraw[fill=\name, draw=black, line width=0.2pt]
    ({90-18*\i}:1)
    -- ({90-18*(\i-1)}:1)
    -- (0,0) -- cycle;

  % Étiquette
  \node at ({90-18*(\i-0.5)}:1.25)
    {\textcolor{\nametext}{\small \nametext}};
}

\end{tikzpicture}
}

\newcommand\ROUEDEUTERANOPEDOUZE{
% Roue chromatique deutéranope — 12 couleurs
\begin{tikzpicture}[scale=3]

\foreach \i/\name in {
  1/crimson,
  2/brickred,
  3/coral,
  4/salmon,
  5/magenta,
  6/fuchsia,
  7/violetB,       % corrigé
  8/indigoB,       % corrigé
  9/royalblueA,    % corrigé
  10/azure,
  11/skyblueB,     % corrigé
  12/tealB}        % corrigé
{
  % Secteur coloré
  \filldraw[fill=\name, draw=black, line width=0.2pt]
    ({90-30*\i}:1)
    -- ({90-30*(\i-1)}:1)
    -- (0,0) -- cycle;

  % Label dans la couleur
  \node at ({90-30*(\i-0.5)}:1.25)
    {\textcolor{\name}{\small \name}};
}

\end{tikzpicture}
}



% === Macro intelligente pour foncer une couleur de façon lisible pour protanope ===
% Usage : \protadarken{mistblue}{40}  -> crée mistblue@dark
% Le 40 signifie : 40% de mélange avec du noir

\newcommand{\DarkenForText}[2]{%
  % #1 = nom de la couleur
  % #2 = pourcentage (40 conseillé si tu ne sais pas)
  %
  % On crée une nouvelle couleur #1@dark
  \colorlet{#1text}{black!#2!#1}%
}





%%%%

% Groupe 1
\DarkenForText{deepgreen}{20}
\DarkenForText{forestgreen}{20}
\DarkenForText{mossgreen}{20}
\DarkenForText{greenwave}{20}
\DarkenForText{forestmist}{20}
\DarkenForText{darkteal}{20}

% Groupe 2
\DarkenForText{seagreen}{20}
\DarkenForText{freshteal}{20}
\DarkenForText{greenteal}{20}
\DarkenForText{blueteal}{20}
\DarkenForText{lagoon}{20}
\DarkenForText{greenmist}{20}

% Groupe 3
\DarkenForText{springgreen}{20}
\DarkenForText{mintgreen}{20}
\DarkenForText{oceanfoam}{20}
\DarkenForText{aquagreen}{20}
\DarkenForText{coastalaqua}{20}
\DarkenForText{softaqua}{20}

% Groupe 4
\DarkenForText{emeraldteal}{20}
\DarkenForText{deepcyan}{20}
\DarkenForText{lightcyan}{20}
\DarkenForText{aquamarine}{20}
\DarkenForText{lightaqua}{20}
\DarkenForText{mistblue}{30}

% Groupe 5
\DarkenForText{paleaqua}{20}
\DarkenForText{glacier}{20}
\DarkenForText{skyblueA}{20}
\DarkenForText{lightsky}{20}
\DarkenForText{iceblue}{20}
\DarkenForText{frost}{20}

% Groupe 6
\DarkenForText{steelblue}{20}
\DarkenForText{cadetblue}{20}
\DarkenForText{coolblue}{20}
\DarkenForText{morningblue}{20}
\DarkenForText{blueiceA}{20}
\DarkenForText{blueair}{20}

% Groupe 7
\DarkenForText{tealblue}{20}
\DarkenForText{deepbluecyan}{20}
\DarkenForText{azur}{20}
\DarkenForText{brightazure}{20}
\DarkenForText{lightazure}{20}
\DarkenForText{skylight}{20}

% Groupe 8
\DarkenForText{nightblue}{20}
\DarkenForText{slateblue}{20}
\DarkenForText{violetA}{20}
\DarkenForText{blueviolet}{20}
\DarkenForText{purpleA}{20}
\DarkenForText{lightpurple}{20}

% Groupe 9
\DarkenForText{lilac}{20}
\DarkenForText{indigoA}{20}
\DarkenForText{deepindigo}{20}
\DarkenForText{royalpurple}{20}

% Groupe 10
\DarkenForText{coolmint}{20}
\DarkenForText{deepsea}{20}

% Palette deutéranope
\DarkenForText{crimson}{20}
\DarkenForText{brickred}{20}
\DarkenForText{coral}{20}
\DarkenForText{salmon}{20}

\DarkenForText{magenta}{20}
\DarkenForText{fuchsia}{20}
\DarkenForText{violetB}{20}
\DarkenForText{indigoB}{20}

\DarkenForText{royalblueA}{20}
\DarkenForText{azure}{20}
\DarkenForText{skyblueB}{20}
\DarkenForText{tealB}{20}

% Couleur supplémentaire
\DarkenForText{marine}{20}

%%%%%%%%%%%%%%%%%%%%%
%
%.   apx proof
%
%%%%%%%%%%%%%%%%%%%%%



% do
\newcommand\PREUVESENAPPENDIXCOMPATIBLE{
	\newenvironment{appendixproof}{{\bf \LANGUE{Proof}{Démonstration}}:}{\hfill $\Box$}
	\newtheorem{theoremrep}[theorem]{Theorem}
	\newtheorem{lemmarep}[theorem]{Lemma}
	\newtheorem{propositionrep}[theorem]{Proposition}
	\newtheorem{corollaryrep}[theorem]{Corollary}
	\newtheorem{summaryrep}[theorem]{Summary}	
}

\providecommand\apxparameter{}

%% proof inline
%\providecommand\apxparameter{appendix=inline}
%% proof in appendix
%\providecommand\apxparameter{}


\apxproofmode{
\usepackage[\apxparameter]{apxproof}
\theoremstyle{plain}

\newtheoremrep{theorem}{Theorem}
%\newtheoremrep{theoremconjecture}[theorem]{Theorem-To-Be-Proved=Conjecture}
%\newtheoremrep{lemmaconjecture}[theorem]{Conjecture-To-Be-Proved=Conjecture}
\newtheoremrep{claim}[theorem]{Theorem}
\newtheoremrep{proposition}[theorem]{Proposition}
\newtheoremrep{lemma}[theorem]{Lemma}
\newtheoremrep{corollary}[theorem]{Corollary}
%%bug mars 26: \newtheoremrep{summary}[theorem]{Summary}
% Du coup:
\newtheoremrep{resume}[theorem]{Summary}
%bug mars 26: \newtheoremrep{openproblem}[theorem]{Open Problem}
%\newtheoremrep{claim}[theorem]{Claim}
}


%\apxproofmodesubsection{
%%apx proof, pour que compteur correct
%\makeatletter
%\AtBeginDocument{%
%\let\axp@section\subsection
%  \let\axp@oldsection\subsection
%}
%\makeatother
%\renewcommand\appendixprelim{\section{Missing proofs}}
%}


\makeatletter
\apxproofmodesubsection{
\let\axp@section\subsection
%\let\axp@oldsection\subsection
\renewcommand\appendixprelim{\section{Proofs from main documents}}
}
\makeatother


\makeatletter
\newcommand\appendixapxproofsubsection{
\appendix
\let\apx@orig@appendix\appendix
\renewcommand{\appendix}{}
}
\makeatother


%%% Pour checker/vérifier que bon


\newcommand\docheck[1]{\textcolor{check}{#1}}

\usepackage{xcolor}
%\definecolor{chose}{RGB}{0,110,50}
%\definecolor{forestgreen}{RGB}{0,150,70}
%\definecolor{ctruc}{RGB}{164,164,164}
%\definecolor{oceanfoam}{RGB}{0,200,150}
\definecolor{check}{RGB}{160,160,0}




. Retirées automatiquement en mode lipics.
\providecommand\NODECORATION[1]{\ifnotSLIDE{#1}}
\lipicsmode{\renewcommand\NODECORATION[1]{}}

\providecommand\PASSIAPXPROOF[1]{#1}
\apxproofmode{\renewcommand\PASSIAPXPROOF[1]{}}



%% end gestion des switchs

\def\MACROSPOINTEXOUVERT{}


\providecommand\ifnotTDEXAM[1]{#1}
\providecommand\ifnotSLIDE[1]{#1}


  
% pour preuves en appendix
\providecommand\PREUVESENAPPENDIX[1]{}



%%%% STYLES: APRES VOLS


\NODECORATION{
\usepackage{fourier}
}

%% Repli si fourier (et donc fourier-orns, qui définit \danger) n'est pas
%% chargé ; doit rester APRÈS le bloc ci-dessus, sinon fourier-orns échoue
%% avec « \danger already defined ».
\providecommand\danger{}

%%%% Box around environment
\usepackage[framemethod=tikz]{mdframed}

\newcommand\ENVCOLOR[1]{#1}
\newcommand\BEGINTEMPENVCOLOR[1]{\let\envcolorwas\ENVCOLOR\renewcommand\ENVCOLOR[1]{#1}}
\newcommand\ENDTEMPENVCOLOR{\let\ENVCOLOR\envcolorwas}


\NODECORATION{
\surroundwithmdframed[hidealllines=true,backgroundcolor=\ENVCOLOR{blue!10},roundcorner=8pt]{exercice}
\surroundwithmdframed[hidealllines=true,backgroundcolor=\ENVCOLOR{blue!10},roundcorner=8pt]{exercise}
\surroundwithmdframed[hidealllines=true,backgroundcolor=\ENVCOLOR{blue!10},roundcorner=8pt]{exercicee}
\surroundwithmdframed[hidealllines=true,backgroundcolor=\ENVCOLOR{blue!10},roundcorner=8pt]{exerciced}
\surroundwithmdframed[hidealllines=true,backgroundcolor=\ENVCOLOR{blue!10},roundcorner=8pt]{exercicec}
\surroundwithmdframed[hidealllines=true,backgroundcolor=\ENVCOLOR{blue!10},roundcorner=8pt]{exercicedc}
\surroundwithmdframed[hidealllines=true,backgroundcolor=\ENVCOLOR{orange!20},roundcorner=8pt]{example}
\surroundwithmdframed[hidealllines=true,backgroundcolor=\ENVCOLOR{orange!20},roundcorner=8pt]{exemple}
\surroundwithmdframed[hidealllines=true,backgroundcolor=\ENVCOLOR{orange!10},roundcorner=8pt]{remark}
\ifnotSLIDE{\surroundwithmdframed[backgroundcolor=\ENVCOLOR{green!15},roundcorner=8pt]{definition}}
\surroundwithmdframed[hidealllines=true,backgroundcolor=\ENVCOLOR{gray!35}]{proposition}
\ifnotSLIDE{\surroundwithmdframed[backgroundcolor=\ENVCOLOR{gray!35},roundcorner=8pt]{theorem}}
\surroundwithmdframed[hidealllines=true,backgroundcolor=\ENVCOLOR{gray!35}]{lemma}
\surroundwithmdframed[hidealllines=true,backgroundcolor=\ENVCOLOR{gray!35}]{corollary}


\surroundwithmdframed[backgroundcolor=red!35,roundcorner=8pt]{theoremaprouver}
\surroundwithmdframed[backgroundcolor=red!35,roundcorner=8pt]{conjectureaprouver}
}

%%% tcolorbox

\usepackage{tcolorbox}
\tcbuselibrary{skins}


  %%. Commande magique
%% \tcolorboxenvironment{⟨name⟩}{⟨options⟩}: attach options de tcolorbox à environnement
%% alternative: je cree des boites moi meme: voici des exemples


\newtcolorbox{maboiteconfiance}[2][]{colback=red!5!white, colframe=red!75!black,fonttitle=\bfseries, colbacktitle=red!85!black,enhanced,
attach boxed title to top center={yshift=-2mm},
  title={#2},#1}
  
\newtcolorbox{maboiteresume}[2][]{colback=red!5!white, colframe=red!75!black,fonttitle=\bfseries, colbacktitle=red!85!black,enhanced,
attach boxed title to top center={yshift=-2mm},
  title={#2},#1}
  \newtcolorbox{maboiteresumeperso}[2][]{colback=red!5!white, colframe=red!75!black,fonttitle=\bfseries, colbacktitle=red!85!black,enhanced,
attach boxed title to top center={yshift=-2mm},
  title={#2},#1}
  \newtcolorbox{maboitecommentaire}[2][]{colback=red!5!white, colframe=red!75!black,fonttitle=\bfseries, colbacktitle=red!85!black,enhanced,
attach boxed title to top center={yshift=-2mm},
  title={#2},#1}
\newtcolorbox{maboitetruc}[2][]{colback=red!5!white, colframe=red!75!black,fonttitle=\bfseries, colbacktitle=red!85!black,enhanced,
attach boxed title to top left={yshift=-2mm},
  title={#2},#1}

%% Boîte des mises en évidence CARE. Couleurs pensées pour la
%% protanopie : l'axe bleu-jaune reste discriminant (jaune = \IMPORTANT
%% d'Olivier, bleu = CARE), et la distinction ne repose pas que sur la
%% couleur (cadre + bandeau étiqueté, lisibles même en niveaux de gris).
\newtcolorbox{maboiteimportantcare}[2][]{colback=cyan!10!white, colframe=blue!60!black,fonttitle=\bfseries, colbacktitle=blue!70!black,enhanced,
attach boxed title to top center={yshift=-2mm},
  title={#2},#1}
  


%%%% FIN STYLES



%%%%%%%%%%%%%%%%%%%%%%%%%%%%%%%%%%%%%%%%%%%%%%%%%
%%%%%%%%%%%%%%%%%%%%%%%%%%%%%%%%%%%%%%%%%%%%%%%%%
%
%
%                      Version ENSEIGNEMENT 2020
%
%
%%%%%%%%%%%%%%%%%%%%%%%%%%%%%%%%%%%%%%%%%%%%%%%%%
%%%%%%%%%%%%%%%%%%%%%%%%%%%%%%%%%%%%%%%%%%%%%%%%%

%english, french
\ifdefined\CESTDUFRANCAIS
	\providecommand\LANGUE[2]{#2}
	\usepackage[french]{babel}
\else
	\providecommand\LANGUE[2]{#1}
\fi
\providecommand\LANGUEinv[2]{\LANGUE{#2}{#1}}

%%%%%%%%%%%%%%%%%%%%%%%%%%%%%%%%%%%%%%%%%%%%%%%%%
%
%
%                      Packages
%
%
%%%%%%%%%%%%%%%%%%%%%%%%%%%%%%%%%%%%%%%%%%%%%%%%%


%%% SURLIGNEUR
\providecommand\ifnotmodeDIFFUSIONFINALE[1]{
%% 2021: Comprend pas mais assure
\ifdefined\SAFEMODE
\else
#1
\fi
%#1
}
\providecommand\ifmodeDIFFUSIONFINALE[1]{
\ifdefined\SAFEMODE
#1
\else
\fi
}
\newcommand\SURLIGNE[1]{
\ifnotmodeDIFFUSIONFINALE{
\BEGINTEMPENVCOLOR{blue!30}
{                        \colorbox{blue!30}{

\noindent \begin{minipage}{\dimexpr\textwidth-2\fboxsep\relax}
#1
\end{minipage}
}
}
\ENDTEMPENVCOLOR}
\ifmodeDIFFUSIONFINALE{
\noindent \begin{minipage}{\textwidth}
#1
\end{minipage}
}
}


\newenvironment{confiance}{\begin{maboiteresume}[colback=brown!50]{Résumé ici}}{\end{maboiteresume}}
\newcommand\CONFIANCE[2][Confiance ici]{
\begin{maboiteconfiance}[colback=yellow!50]{#1}
\textbf{ATTENTION : #2}
\end{maboiteconfiance}
}

\newenvironment{resume}{\begin{maboiteresume}[colback=brown!50]{Résumé ici}}{\end{maboiteresume}}
\newcommand\RESUME[2][Résumé ici]{
\begin{maboiteresume}[colback=brown!50]{#1}
#2
\end{maboiteresume}
}

\newenvironment{resumeperso}{\begin{maboiteresumeperso}[colback=brown!50]{Résumé perso ici}}{\end{maboiteresumeperso}}
\newcommand\RESUMEPERSO[2][Résumé perso  ici]{
\begin{maboiteresumeperso}[colback=brown!50]{#1}
#2
\end{maboiteresumeperso}
%
%
%\todo[inline, color=brown]{résumé: #1}
%\colorbox{brown}{
%\noindent \begin{minipage}{\textwidth}
%{
%\bf 
%#2
%}
%\end{minipage}
%}
}

\newcommand\DANSLAMARGE[1]{
\ifnotmodeDIFFUSIONFINALE{
\marginpar{\textcolor{brown}{#1}}}
}

\newenvironment{commentaireperso}{\begin{maboitecommentaire}[colback=brown!30]{Commentaire perso} Commentaire perso:}{\end{maboitecommentaire}}
\newcommand\COMMENTAIREPERSO[2][Commentaire perso]{
\begin{maboitecommentaire}[colback=brown!30]{#1}
#2
\end{maboitecommentaire}
}


\newcommand\SURSURLIGNE[2][truc surligné ici]{
%\todo[inline, color=orange]{surligné: #1}
\ifnotmodeDIFFUSIONFINALE{
\BEGINTEMPENVCOLOR{blue!60}
	                 \colorbox{blue!60}{
\noindent \begin{minipage}{\dimexpr\textwidth-2\fboxsep\relax}
{
\bf 
#2
}
\end{minipage}
}
\ENDTEMPENVCOLOR
}
\ifmodeDIFFUSIONFINALE{#2}
}

\newcommand\IMPORTANT[1]{
\ifnotmodeDIFFUSIONFINALE{
\BEGINTEMPENVCOLOR{yellow!100}
		        \colorbox{yellow!100}{

\noindent \begin{minipage}{\dimexpr\textwidth-2\fboxsep\relax}
{ 
#1
}
\end{minipage}
}
\ENDTEMPENVCOLOR
}
\ifmodeDIFFUSIONFINALE{#1}
}

%% Pendant de \IMPORTANT, réservé aux mises en évidence CARE ;
%% \IMPORTANT reste réservé aux annotations d'Olivier. Comme \IMPORTANT,
%% redevient du texte simple en mode DIFFUSIONFINALE.
\newcommand\IMPORTANTCARE[2][Important (CARE)]{
\ifnotmodeDIFFUSIONFINALE{
\begin{maboiteimportantcare}{#1}
#2
\end{maboiteimportantcare}
}
\ifmodeDIFFUSIONFINALE{#2}
}


\newcommand\QUESTION[1]{
\ifnotmodeDIFFUSIONFINALE{
\BEGINTEMPENVCOLOR{oceanfoam!100}
		        \colorbox{oceanfoam!100}{

\noindent \begin{minipage}{\dimexpr\textwidth-2\fboxsep\relax}
{ 
#1
}
\end{minipage}
}
\ENDTEMPENVCOLOR
}
\ifmodeDIFFUSIONFINALE{#1}
}




\newcommand\SOUSLIGNE[1]{
\colorbox{lightgray}{

\noindent \begin{minipage}{\dimexpr\textwidth-2\fboxsep\relax}
{ \color{gray}
#1
}
\end{minipage}
}
}



%% Gestion des cross-references cross-referencing
\usepackage{xr-hyper}
\usepackage{hyperref}

%%
\usepackage{versions}
\usepackage{amsmath,amssymb,mathrsfs,amsfonts}
%\usepackage{mathabx}
\usepackage{listings}
\usepackage{url}
\usepackage{versions}
\usepackage{color}
\usepackage[colorinlistoftodos]{todonotes}
%\usepackage{todonotes}
\usepackage{proof}
\usepackage{xstring}

%% Index
\usepackage{makeidx}
\makeindex
\providecommand\DOINDEXPRINT{}
\providecommand\SIGNALEMOTNOUV[1]{}
\ifdefined\INDEXPRINT
	\ifnotSAFEMODE{
%	\usepackage{showidx}
%	\let\oldindex\index
%	\renewcommand{\index}[1]{\oldindex{#1}\marginpar{LA: \tiny#1}}
	\renewcommand\DOINDEXPRINT{
	      	\let\oldindex\index
		\renewcommand{\index}[1]{\oldindex{##1}\marginpar{IDXLA: \tiny##1}}
		\renewcommand\SIGNALEMOTNOUV[1]{\marginpar{MNV: \small\textbf{##1}}}
	}
	}
\fi

% === gestion des accents


%soit iso latin1
%\usepackage[cyr]{aeguill}
%soit utf8
\usepackage[utf8]{inputenc}
\usepackage[T1]{fontenc}


%%%  Pour citation \Cref (cref, ref)	

\usepackage{cleveref}


%%%%%%%%%%%%%%%%%%%%%%%%%%%%%%%%%%%%%%%%%%%%%%%%%
%
%
%                      Gestion des accents spéciaux.
%
%
%%%%%%%%%%%%%%%%%%%%%%%%%%%%%%%%%%%%%%%%%%%%%%%%%


%% MOI A LA MAIN:
\usepackage{bbm}

\DeclareUnicodeCharacter{2212}{-}
\DeclareUnicodeCharacter{2297}{\ensuremath{\otimes}}       
\DeclareUnicodeCharacter{2A33}{\ensuremath{\smashtimes}}
\DeclareUnicodeCharacter{1D56F}{\ensuremath{\mathfrak{D}}}  
\DeclareUnicodeCharacter{1D40E}{\ensuremath{\mathbf{O}}}  
\DeclareUnicodeCharacter{1D427}{\ensuremath{\mathbf{n}}}  
\DeclareUnicodeCharacter{3008}{\ensuremath{\langle}} % 〈 
\DeclareUnicodeCharacter{3009}{\ensuremath{\rangle}} %  〉

% (possible de voir: http://tug.ctan.org/info/symbols/comprehensive/symbols-letter.pdf)

%% SOURCE:  https://github.com/agda/agda/blob/master/doc/user-manual/unicode-symbols-sphinx.tex.txt

%%%%%%%%%%%%%%%%%%%%%%%%%%%%%%%%%%%%%%%%%%%%%%%%%%%%%%%%%%%%%%%%%%%%%%%%%%%%%%
% Please keep the below list ordered by code points!

\DeclareUnicodeCharacter{00AC}{\ensuremath{\neg}}                     % Symbol ¬
\DeclareUnicodeCharacter{00B1}{\ensuremath{\pm}}                      % Symbol ±
\DeclareUnicodeCharacter{00B2}{\ensuremath{^2}}                       % Symbol ²
\DeclareUnicodeCharacter{00B3}{\ensuremath{^3}}                       % Symbol ³
\DeclareUnicodeCharacter{00B7}{\ensuremath{\cdot}}                    % Symbol ·
\DeclareUnicodeCharacter{00B9}{\ensuremath{^1}}                       % Symbol ¹
\DeclareUnicodeCharacter{00D7}{\ensuremath{\times}}                   % Symbol ×
\DeclareUnicodeCharacter{02B0}{\ensuremath{^h}}                       % Symbol ʰ
\DeclareUnicodeCharacter{02B3}{\ensuremath{^r}}                       % Symbol ʳ
\DeclareUnicodeCharacter{02E2}{\ensuremath{^s}}                       % Symbol ˢ
\DeclareUnicodeCharacter{0302}{\ensuremath{\hat{\phantom{x}}}}        % Symbol ̂
\DeclareUnicodeCharacter{0332}{\ensuremath{\underline{\phantom{x}}}}  % Symbol ̲
\DeclareUnicodeCharacter{0308}{\ensuremath{\ddot{\phantom{x}}}}       % Symbol ̈
\DeclareUnicodeCharacter{0393}{\ensuremath{\Gamma}}                   % Symbol Γ
\DeclareUnicodeCharacter{0394}{\ensuremath{\Delta}}                   % Symbol Δ
\DeclareUnicodeCharacter{0398}{\ensuremath{\Theta}}                   % Symbol Θ
\DeclareUnicodeCharacter{039B}{\ensuremath{\Lambda}}                  % Symbol Λ
\DeclareUnicodeCharacter{039E}{\ensuremath{\Xi}}                      % Symbol Ξ
\DeclareUnicodeCharacter{03A0}{\ensuremath{\Pi}}                      % Symbol Π
\DeclareUnicodeCharacter{03A3}{\ensuremath{\Sigma}}                   % Symbol Σ


% There is not `\Tau` command in LaTeX, so we use `\mathrm{T}`. See
% http://tex.stackexchange.com/questions/1368/why-can-i-only-use-some-capital-greek-letters-inside-my-equations.
\DeclareUnicodeCharacter{03A4}{\ensuremath{\mathrm{T}}}  % Symbol Τ

\DeclareUnicodeCharacter{03A5}{\ensuremath{\Upsilon}}   % Symbol Υ
\DeclareUnicodeCharacter{03A6}{\ensuremath{\Phi}}       % Symbol Φ
\DeclareUnicodeCharacter{03A8}{\ensuremath{\Psi}}       % Symbol Ψ
\DeclareUnicodeCharacter{03A9}{\ensuremath{\Omega}}     % Symbol Ω
\DeclareUnicodeCharacter{03B1}{\ensuremath{\alpha}}     % Symbol α
\DeclareUnicodeCharacter{03B2}{\ensuremath{\beta}}      % Symbol β
\DeclareUnicodeCharacter{03B3}{\ensuremath{\gamma}}     % Symbol γ
\DeclareUnicodeCharacter{03B4}{\ensuremath{\delta}}     % Symbol δ
\DeclareUnicodeCharacter{03B5}{\ensuremath{\epsilon}}   % Symbol ε
\DeclareUnicodeCharacter{03B6}{\ensuremath{\zeta}}      % Symbol ζ
\DeclareUnicodeCharacter{03B7}{\ensuremath{\eta}}       % Symbol η
\DeclareUnicodeCharacter{03B8}{\ensuremath{\theta}}     % Symbol θ
\DeclareUnicodeCharacter{03B9}{\ensuremath{\iota}}      % Symbol ι
\DeclareUnicodeCharacter{03BA}{\ensuremath{\kappa}}     % Symbol κ
\DeclareUnicodeCharacter{03BB}{\ensuremath{\lambda}}    % Symbol λ
\DeclareUnicodeCharacter{03BC}{\ensuremath{\mu}}        % Symbol μ
\DeclareUnicodeCharacter{03BD}{\ensuremath{\nu}}        % Symbol ν
\DeclareUnicodeCharacter{03BE}{\ensuremath{\xi}}        % Symbol ξ
\DeclareUnicodeCharacter{03C0}{\ensuremath{\pi}}        % Symbol π
\DeclareUnicodeCharacter{03C1}{\ensuremath{\rho}}       % Symbol ρ
\DeclareUnicodeCharacter{03C2}{\ensuremath{\varsigma}}  % Symbol ς
\DeclareUnicodeCharacter{03C3}{\ensuremath{\sigma}}     % Symbol σ
\DeclareUnicodeCharacter{03C4}{\ensuremath{\tau}}       % Symbol τ
\DeclareUnicodeCharacter{03C5}{\ensuremath{\upsilon}}   % Symbol υ
\DeclareUnicodeCharacter{03C6}{\ensuremath{\varphi}}    % Symbol φ
\DeclareUnicodeCharacter{03C7}{\ensuremath{\chi}}       % Symbol χ
\DeclareUnicodeCharacter{03C8}{\ensuremath{\psi}}       % Symbol ψ
\DeclareUnicodeCharacter{03C9}{\ensuremath{\omega}}     % Symbol ω
\DeclareUnicodeCharacter{03D1}{\ensuremath{\vartheta}}  % Symbol ϑ

\DeclareUnicodeCharacter{0432}{\ensuremath{\mathrm{PDF\;TODO}}}  % Symbol в
\DeclareUnicodeCharacter{0435}{\ensuremath{\mathrm{PDF\;TODO}}}  % Symbol е
\DeclareUnicodeCharacter{0438}{\ensuremath{\mathrm{PDF\;TODO}}}  % Symbol и
\DeclareUnicodeCharacter{043C}{\ensuremath{\mathrm{PDF\;TODO}}}  % Symbol м
\DeclareUnicodeCharacter{0440}{\ensuremath{\mathrm{PDF\;TODO}}}  % Symbol р
\DeclareUnicodeCharacter{0442}{\ensuremath{\mathrm{PDF\;TODO}}}  % Symbol т
\DeclareUnicodeCharacter{041F}{\ensuremath{\mathrm{PDF\;TODO}}}  % Symbol П
\DeclareUnicodeCharacter{0BA8}{\ensuremath{\mathrm{PDF\;TODO}}}  % Symbol ந
\DeclareUnicodeCharacter{0BBF}{\ensuremath{\mathrm{PDF\;TODO}}}  % Symbol  ி
\DeclareUnicodeCharacter{1100}{\ensuremath{\mathrm{PDF\;TODO}}}  % Symbol ᄀ
\DeclareUnicodeCharacter{11F9}{\ensuremath{\mathrm{PDF\;TODO}}}  % Symbol ᇹ

\DeclareUnicodeCharacter{1D48}{\ensuremath{^d}}  % Symbol ᵈ
\DeclareUnicodeCharacter{1D50}{\ensuremath{^m}}  % Symbol ᵐ
\DeclareUnicodeCharacter{1D56}{\ensuremath{^p}}  % Symbol ᵖ
\DeclareUnicodeCharacter{1D62}{\ensuremath{_i}}  % Symbol ᵢ
\DeclareUnicodeCharacter{1D63}{\ensuremath{_r}}  % Symbol ᵣ
\DeclareUnicodeCharacter{1D64}{\ensuremath{_u}}  % Symbol ᵤ
\DeclareUnicodeCharacter{1D9C}{\ensuremath{^c}}  % Symbol ᶜ

% The command `\text` requires `\usepackage{amsmath}` and the command
% `\textasciiacute` requires `\usepackage{textcomp}`.

\DeclareUnicodeCharacter{2032}{\ensuremath{\text{\textasciiacute}}}  % Symbol ′

\DeclareUnicodeCharacter{2070}{\ensuremath{^0}}  % Symbol ⁰
\DeclareUnicodeCharacter{2074}{\ensuremath{^4}}  % Symbol ⁴
\DeclareUnicodeCharacter{2075}{\ensuremath{^5}}  % Symbol ⁵
\DeclareUnicodeCharacter{2076}{\ensuremath{^6}}  % Symbol ⁶
\DeclareUnicodeCharacter{2077}{\ensuremath{^7}}  % Symbol ⁷
\DeclareUnicodeCharacter{2078}{\ensuremath{^8}}  % Symbol ⁸
\DeclareUnicodeCharacter{2079}{\ensuremath{^9}}  % Symbol ⁹
\DeclareUnicodeCharacter{207A}{\ensuremath{^+}}  % Symbol ⁺
\DeclareUnicodeCharacter{207B}{\ensuremath{^-}}  % Symbol ⁻
\DeclareUnicodeCharacter{207F}{\ensuremath{^n}}  % Symbol ⁿ
\DeclareUnicodeCharacter{2080}{\ensuremath{_0}}  % Symbol ₀
\DeclareUnicodeCharacter{2081}{\ensuremath{_1}}  % Symbol ₁
\DeclareUnicodeCharacter{2082}{\ensuremath{_2}}  % Symbol ₂
\DeclareUnicodeCharacter{2083}{\ensuremath{_3}}  % Symbol ₃
\DeclareUnicodeCharacter{2084}{\ensuremath{_4}}  % Symbol ₄
\DeclareUnicodeCharacter{2085}{\ensuremath{_5}}  % Symbol ₅
\DeclareUnicodeCharacter{2086}{\ensuremath{_6}}  % Symbol ₆
\DeclareUnicodeCharacter{2087}{\ensuremath{_7}}  % Symbol ₇
\DeclareUnicodeCharacter{2088}{\ensuremath{_8}}  % Symbol ₈
\DeclareUnicodeCharacter{2089}{\ensuremath{_9}}  % Symbol ₉
\DeclareUnicodeCharacter{208A}{\ensuremath{_+}}  % Symbol ₊
\DeclareUnicodeCharacter{208B}{\ensuremath{_-}}  % Symbol ₋
\DeclareUnicodeCharacter{2091}{\ensuremath{_e}}  % Symbol ₑ
\DeclareUnicodeCharacter{2095}{\ensuremath{_h}}  % Symbol ₕ
\DeclareUnicodeCharacter{2096}{\ensuremath{_k}}  % Symbol ₖ
\DeclareUnicodeCharacter{2097}{\ensuremath{_l}}  % Symbol ₗ
\DeclareUnicodeCharacter{2098}{\ensuremath{_m}}  % Symbol ₘ
\DeclareUnicodeCharacter{2099}{\ensuremath{_n}}  % Symbol ₙ
\DeclareUnicodeCharacter{209A}{\ensuremath{_p}}  % Symbol ₚ
\DeclareUnicodeCharacter{209B}{\ensuremath{_s}}  % Symbol ₛ
\DeclareUnicodeCharacter{209C}{\ensuremath{_t}}  % Symbol ₜ

\DeclareUnicodeCharacter{2113}{\ensuremath{\mathpzc{l}}} % Symbol ℓ

% The command `\mathbbm` requires `\usepackage{bbm}`.
\DeclareUnicodeCharacter{2115}{\ensuremath{\mathbbm{N}}}  % Symbol ℕ
\DeclareUnicodeCharacter{211A}{\ensuremath{\mathbbm{Q}}}  % Symbol ℚ
\DeclareUnicodeCharacter{211D}{\ensuremath{\mathbbm{R}}}  % Symbol ℝ
\DeclareUnicodeCharacter{2124}{\ensuremath{\mathbbm{Z}}}  % Symbol ℤ

\DeclareUnicodeCharacter{2135}{\ensuremath{\aleph}}           % Symbol ℵ
\DeclareUnicodeCharacter{2190}{\ensuremath{\leftarrow}}       % Symbol ←
\DeclareUnicodeCharacter{2192}{\ensuremath{\rightarrow}}      % Symbol →
\DeclareUnicodeCharacter{2194}{\ensuremath{\leftrightarrow}}  % Symbol ↔

% The command `\twoheadrightarrow` requires `\usepackage{amssymb}`.
\DeclareUnicodeCharacter{21A0}{\ensuremath{\twoheadrightarrow}}  % Symbol ↠

\DeclareUnicodeCharacter{21A6}{\ensuremath{\mapsto}}  % Symbol ↦
\DeclareUnicodeCharacter{21A6}{\ensuremath{\mapsto}}  % Symbol ↦

% The command `\restriction` requires `\usepackage{amssymb}`.
\DeclareUnicodeCharacter{21BE}{\ensuremath{\restriction}}  % Symbol ↾

\DeclareUnicodeCharacter{21D0}{\ensuremath{\Leftarrow}}       % Symbol ⇐
\DeclareUnicodeCharacter{21D2}{\ensuremath{\Rightarrow}}      % Symbol ⇒
\DeclareUnicodeCharacter{21D4}{\ensuremath{\Leftrightarrow}}  % Symbol ⇔

\DeclareUnicodeCharacter{2200}{\ensuremath{\forall}}      % Symbol ∀
\DeclareUnicodeCharacter{2203}{\ensuremath{\exists}}      % Symbol ∃
\DeclareUnicodeCharacter{2204}{\ensuremath{\not\exists}}  % Symbol ∃
\DeclareUnicodeCharacter{2205}{\ensuremath{\emptyset}}    % Symbol ∅
\DeclareUnicodeCharacter{2208}{\ensuremath{\in}}          % Symbol ∈
\DeclareUnicodeCharacter{2209}{\ensuremath{\not\in}}      % Symbol ∉
\DeclareUnicodeCharacter{220B}{\ensuremath{\ni}}          % Symbol ∋
\DeclareUnicodeCharacter{220C}{\ensuremath{\not\ni}}      % Symbol ∌

% The command `\blacksquare` requires `\usepackage{amssymb}`.
\DeclareUnicodeCharacter{220E}{{\tiny \ensuremath{\blacksquare}}} % Symbol ∎

\DeclareUnicodeCharacter{220F}{{\tiny \ensuremath{\prod}}}  % Symbol ∏
\DeclareUnicodeCharacter{2211}{\ensuremath{\sum}}           % Symbol ∑
\DeclareUnicodeCharacter{2218}{\ensuremath{\circ}}          % Symbol ∘
\DeclareUnicodeCharacter{2219}{\ensuremath{\bullet}}        % Symbol ∙
\DeclareUnicodeCharacter{221E}{\ensuremath{\infty}}         % Symbol ∞
\DeclareUnicodeCharacter{2223}{\ensuremath{\mid}}           % Symbol ∣
\DeclareUnicodeCharacter{2225}{\ensuremath{\parallel}}      % Symbol ∥
\DeclareUnicodeCharacter{2227}{\ensuremath{\wedge}}         % Symbol ∧
\DeclareUnicodeCharacter{2228}{\ensuremath{\vee}}           % Symbol ∨
\DeclareUnicodeCharacter{2229}{\ensuremath{\cap}}           % Symbol ∩
\DeclareUnicodeCharacter{222A}{\ensuremath{\cup}}           % Symbol ∪
\DeclareUnicodeCharacter{222B}{\ensuremath{\int}}           % Symbol ∫
\DeclareUnicodeCharacter{2234}{\ensuremath{\therefore}}     % Symbol ∴

% The command `\dblcolon` requires `\usepackage{mathtools}`.
\DeclareUnicodeCharacter{2237}{\ensuremath{\dblcolon}} % Symbol ∷

\newcommand{\dotminus}{\mathop{\mbox{$-^{\hspace{-.5em}\cdot}\,$}}}
\DeclareUnicodeCharacter{2238}{\ensuremath{\dotminus}}  % Symbol ∸

\DeclareUnicodeCharacter{223C}{\ensuremath{\sim}}       % Symbol ∼
\DeclareUnicodeCharacter{2243}{\ensuremath{\simeq}}     % Symbol ≃
\DeclareUnicodeCharacter{2245}{\ensuremath{\cong}}      % Symbol ≅
\DeclareUnicodeCharacter{2248}{\ensuremath{\approx}}    % Symbol ≈
\DeclareUnicodeCharacter{224D}{\ensuremath{\asymp}}       % Symbol ≍ : OLIVIER
\DeclareUnicodeCharacter{2250}{\ensuremath{\doteq}}     % Symbol ≐

% The command `\coloneqq` requires `\usepackage{mathtools}`.
\DeclareUnicodeCharacter{2254}{\ensuremath{\coloneqq}}  % Symbol ≔

\newcommand{\eqq}{\stackrel{\text{\tiny ?}}{=}}
\DeclareUnicodeCharacter{225F}{\ensuremath{\eqq}}  % Symbol ≟

\DeclareUnicodeCharacter{2260}{\ensuremath{\neq}}         % Symbol ≠
\DeclareUnicodeCharacter{2261}{\ensuremath{\equiv}}       % Symbol ≡
\DeclareUnicodeCharacter{2262}{\ensuremath{\not\equiv}}   % Symbol ≢
\DeclareUnicodeCharacter{2264}{\ensuremath{\le}}          % Symbol ≤
\DeclareUnicodeCharacter{2265}{\ensuremath{\ge}}          % Symbol ≥
\DeclareUnicodeCharacter{226E}{\ensuremath{\nless}}       % Symbol ≮
\DeclareUnicodeCharacter{226F}{\ensuremath{\ngtr}}        % Symbol ≯
\DeclareUnicodeCharacter{2270}{\ensuremath{\nleq}}        % Symbol ≰
\DeclareUnicodeCharacter{2271}{\ensuremath{\ngeq}}        % Symbol ≱
\DeclareUnicodeCharacter{227A}{\ensuremath{\prec}}        % Symbol ≺
\DeclareUnicodeCharacter{227C}{\ensuremath{\preceq}}      % Symbol ≼
\DeclareUnicodeCharacter{2282}{\ensuremath{\subset}}      % Symbol ⊂
\DeclareUnicodeCharacter{2284}{\ensuremath{\not\subset}} % Symbol ⊄
\DeclareUnicodeCharacter{2283}{\ensuremath{\supset}}      % Symbol ⊃
\DeclareUnicodeCharacter{2286}{\ensuremath{\subseteq}}    % Symbol ⊆
\DeclareUnicodeCharacter{228E}{\ensuremath{\uplus}}       % Symbol ⊎
\DeclareUnicodeCharacter{2291}{\ensuremath{\sqsubseteq}}  % Symbol ⊑
\DeclareUnicodeCharacter{2293}{\ensuremath{\sqcap}}       % Symbol ⊓
\DeclareUnicodeCharacter{2294}{\ensuremath{\sqcup}}       % Symbol ⊔
\DeclareUnicodeCharacter{2295}{\ensuremath{\oplus}}       % Symbol ⊕
\DeclareUnicodeCharacter{22A2}{\ensuremath{\vdash}}       % Symbol ⊢
\DeclareUnicodeCharacter{22A4}{\ensuremath{\top}}         % Symbol ⊤
\DeclareUnicodeCharacter{22A5}{\ensuremath{\bot}}         % Symbol ⊥
\DeclareUnicodeCharacter{22C0}{\ensuremath{\bigwedge}}    % Symbol ⋀
\DeclareUnicodeCharacter{22C2}{\ensuremath{\bigcap}}      % Symbol ⋂
\DeclareUnicodeCharacter{22C3}{\ensuremath{\bigcup}}      % Symbol ⋃
\DeclareUnicodeCharacter{22EE}{\ensuremath{\vdots}}       % Symbol ⋮
\DeclareUnicodeCharacter{22EF}{\ensuremath{\cdots}}       % Symbol ⋯

% The command `\iddots` requires `\usepackage{mathdots}`.
\DeclareUnicodeCharacter{22F0}{\ensuremath{\iddots}} % Symbol ⋰

\DeclareUnicodeCharacter{230A}{\ensuremath{\lfloor}}  % Symbol ⌊
\DeclareUnicodeCharacter{230B}{\ensuremath{\rfloor}}  % Symbol ⌋
\DeclareUnicodeCharacter{2308}{\ensuremath{\lceil}}   % Symbol ⌈
\DeclareUnicodeCharacter{2309}{\ensuremath{\rceil}}   % Symbol ⌉
\DeclareUnicodeCharacter{2329}{\ensuremath{\langle}}  % Symbol 〈
\DeclareUnicodeCharacter{232A}{\ensuremath{\rangle}}  % Symbol 〉

% The command `\Return` requires `\usepackage{keystroke}`.
\DeclareUnicodeCharacter{23CE}{\ensuremath{\Return}}  % Symbol ⏎

% The command `\text` requires `\usepackage{amsmath}` and the command
% `\ding` requires `\usepackage{pifont}`.
\DeclareUnicodeCharacter{2460}{\ensuremath{\text{\ding{172}}}} % Symbol ①
\DeclareUnicodeCharacter{2461}{\ensuremath{\text{\ding{173}}}} % Symbol ②
\DeclareUnicodeCharacter{2462}{\ensuremath{\text{\ding{174}}}} % Symbol ③
\DeclareUnicodeCharacter{2463}{\ensuremath{\text{\ding{175}}}} % Symbol ④
\DeclareUnicodeCharacter{2464}{\ensuremath{\text{\ding{176}}}} % Symbol ⑤
\DeclareUnicodeCharacter{2465}{\ensuremath{\text{\ding{177}}}} % Symbol ⑥
\DeclareUnicodeCharacter{2466}{\ensuremath{\text{\ding{178}}}} % Symbol ⑦
\DeclareUnicodeCharacter{2467}{\ensuremath{\text{\ding{179}}}} % Symbol ⑧
\DeclareUnicodeCharacter{2468}{\ensuremath{\text{\ding{180}}}} % Symbol ⑨

\DeclareUnicodeCharacter{25C5}{\ensuremath{\triangleleft}}  % Symbol ◅

\DeclareUnicodeCharacter{2660}{\ensuremath{\spadesuit}} % Symbol ♠
\DeclareUnicodeCharacter{266D}{\ensuremath{\flat}}      % Symbol ♭
\DeclareUnicodeCharacter{266F}{\ensuremath{\sharp}}     % Symbol ♯

\DeclareUnicodeCharacter{2713}{\ensuremath{\checkmark}}  % Symbol ✓

% The commands `\llbracket` and `\rrbracket` require
% `\usepackage{stmaryrd}`.
\DeclareUnicodeCharacter{27E6}{\ensuremath{\llbracket}}  % Symbol ⟦
\DeclareUnicodeCharacter{27E7}{\ensuremath{\rrbracket}}  % Symbol ⟧

\DeclareUnicodeCharacter{27E8}{\ensuremath{\langle}}          % Symbol ⟨
\DeclareUnicodeCharacter{27E9}{\ensuremath{\rangle}}          % Symbol 〉
\DeclareUnicodeCharacter{27F6}{\ensuremath{\longrightarrow}}  % Symbol % ⟶

\DeclareUnicodeCharacter{2983}{\ensuremath{\mathrm{PDF\;TODO}}}  % Symbol ⦃
\DeclareUnicodeCharacter{2984}{\ensuremath{\mathrm{PDF\;TODO}}}  % Symbol ⦄
\DeclareUnicodeCharacter{2985}{\ensuremath{\mathrm{PDF\;TODO}}}  % Symbol ⦅
\DeclareUnicodeCharacter{2986}{\ensuremath{\mathrm{PDF\;TODO}}}  % Symbol ⦆

% The commands `\llparenthesis` and `\rrparenthesis` require
% `\usepackage{stmaryrd}`.
\DeclareUnicodeCharacter{2987}{\ensuremath{\llparenthesis}}  % Symbol ⦇
\DeclareUnicodeCharacter{2988}{\ensuremath{\rrparenthesis}}  % Symbol ⦈

\DeclareUnicodeCharacter{29F5}{\ensuremath{\setminus}}   % Symbol ⧵

\DeclareUnicodeCharacter{2A06}{\ensuremath{\bigcup}}  % Symbol ⋃
\DeclareUnicodeCharacter{2C7C}{\ensuremath{_j}}       % Symbol ⱼ

\DeclareUnicodeCharacter{D7B0}{\ensuremath{\mathrm{PDF\;TODO}}} % Symbol ힰ

% The command `\mathbbm{1}` requires `\usepackage{bbm}`.
\DeclareUnicodeCharacter{1D7D9}{\ensuremath{\mathbbm{1}}}  % Symbol 𝟙


%%%% MOI


\DeclareUnicodeCharacter{22AC}{\ensuremath{\neg\vdash}}
\DeclareUnicodeCharacter{279C}{\ensuremath{\to}}
\DeclareUnicodeCharacter{1D442}{ }
\DeclareUnicodeCharacter{2000}{ }
\DeclareUnicodeCharacter{1D440}{M}
\DeclareUnicodeCharacter{1D442}{O}
\DeclareUnicodeCharacter{20E3}{2}
\DeclareUnicodeCharacter{1D45A}{m}
\DeclareUnicodeCharacter{1D45B}{n}

\usepackage{pifont}
\DeclareUnicodeCharacter{2714}{\ding{52}}
\DeclareUnicodeCharacter{2705}{\ding{52}}
\DeclareUnicodeCharacter{2717}{\ding{55}}
\DeclareUnicodeCharacter{274C}{\ding{53}}


%% Tolérant tant que macros-markdown.tex n'est pas dans le dépôt :
\InputIfFileExists{macros-markdown}{}{\typeout{*** macros-markdown.tex introuvable: ignore ***}}


%\usepackage{PDFdraftcopy}


%-- Pour état des chapitres


%\newcommand\ETAT{\draftstring{BROUILLON}}
%
%\newcommand\brouillon{\renewcommand\ETAT{\draftstring{BROUILLON}}}
%
%\newcommand\final{\renewcommand\ETAT{\draftstring{\rien}}}
%
%\newcommand\travail{\renewcommand\ETAT{\draftstring{CHANTIER}}}
%
%\newcommand\bazar{\renewcommand\ETAT{\draftstring{BAZAR TOTAL}}}
%
%\final



%%%%%%%%%%%%%%%%%%%%%%%%%%%%%%%%%%%%%%%%%%%%%%%%%
%
%
%                      TRADUCTION EN COURS
%
%%%%%%%%%%%%%%%%%%%%%%%%%%%%%%%%%%%%%%%%%%%%%%%%%


% === volé Amaury



\usepackage{stmaryrd}
\usepackage{ifthen}
\usepackage{mathtools}
\usepackage{dsfont}
\PASSIAPXPROOF{\usepackage{thmtools}}
\usepackage{longtable}
\usepackage{mathdots}
\usepackage{booktabs}


%%%%%%%%%%%%%%%%%%%%%%%%%%%%%%%%%%%%%%%%%%%%%%%%%
%
%
%                      TIKZ STUFFS
%
%
%%%%%%%%%%%%%%%%%%%%%%%%%%%%%%%%%%%%%%%%%%%%%%%%%

%%% volé pour dessin des GPAC: PAS SUR QUE BESOIN de TOUT

\usepackage{tikz}
\usetikzlibrary{calc}
\usepgflibrary{arrows}
\usetikzlibrary{fit}
\usetikzlibrary{patterns}
\usetikzlibrary{decorations.pathreplacing}
\usetikzlibrary{shapes.multipart}
\usetikzlibrary{shapes.geometric}
\usetikzlibrary{shapes.misc}
\usetikzlibrary{positioning}
\usetikzlibrary{graphs}
\usetikzlibrary{topaths}
\usetikzlibrary{arrows.meta}
\usetikzlibrary{decorations.pathreplacing}
\usetikzlibrary{decorations.markings}
\usetikzlibrary{decorations.pathmorphing}
%
 \usetikzlibrary{calc}
 \usetikzlibrary{automata}
\usetikzlibrary{chains}
\usetikzlibrary{scopes}
 \usetikzlibrary{positioning}
 \usetikzlibrary{through,backgrounds}
\usepackage{circuitikz}
% % pour accolades




%%%%%%%%%%%%%%%%%%%%%%%%%%%%%%%%%%%%%%%%%%%%%%%%%
%
%
%                      MACROS locales
%
%
%%%%%%%%%%%%%%%%%%%%%%%%%%%%%%%%%%%%%%%%%%%%%%%%%


%====================
%                      Moi
%=====================


\newcommand\myaddress{ 
{\sc LIX, CNRS, Ecole Polytechnique, Institut Polytechnique de Paris, }\\ 91128 Palaiseau Cedex, FRANCE \\
{\small Olivier.Bournez@lix.polytechnique.fr}}



%====================
%                      un peu de trucs sales
%=====================

%=  pour style lipics et ses trucs

\providecommand\ifnotPOLYCHAPTERMODE[1]{#1}
\providecommand\ifPOLYCHAPTERMODE[1]{}

\ifnotPOLYCHAPTERMODE{
\ifdefined\thechapter
\providecommand\spnewtheorem[4]{\newtheorem{#1}{#2}[chapter]}
\providecommand\spnewtheoremi[5]{\newtheorem{#1}[#5]{#2}}
\else
\providecommand\spnewtheorem[4]{\newtheorem{#1}{#2}}
\providecommand\spnewtheoremi[5]{\newtheorem{#1}[#5]{#2}}
\fi
}
\ifPOLYCHAPTERMODE{
\providecommand\spnewtheorem[4]{\newtheorem{#1}{#2}}
\providecommand\spnewtheoremi[5]{\newtheorem{#1}[#5]{#2}}
}
\providecommand\nolipics[1]{#1}




%====================
%                      un peu de trucs franchement sales
%=====================


% on l'inclus pas ici.
%
%% defaults

\includeversion{NOTLFMI2020}
\includeversion{NOTCSE304}
\excludeversion{NOTDIFFUSIONFINALE}

%% Booleans (dont teste si existe ou pas, pour construire des macros)
% \SAFEMODE
% -> \DIFFUSIONFINALE (via entête-cours)
% -> \PAGEDEGARDE (via entête-cours)
% \POLYMODE
% \POLYCHAPTERMODE

%%%%%%%%%%%%%%%%%%%%%%%%%%%%%%%%%%%%%%%%%%%%%%%%%
%
%
%                      Hack pour DIFFUSION FINALE/ LFMI / Fin provisoire
%
%
%%%%%%%%%%%%%%%%%%%%%%%%%%%%%%%%%%%%%%%%%%%%%%%%%

%\def\DIFFUSIONFINALE{}

\providecommand\ifnotmodeDIFFUSIONFINALE[1]{
\ifdefined\DIFFUSIONFINALE
\else
#1
\fi
}

\providecommand\ifmodeDIFFUSIONFINALE[1]{
\ifdefined\DIFFUSIONFINALE
%% 2021: Comprend pas mais assure
\ifdefined\SAFEMODE
\else
#1
\fi
%#1
\fi
}



\ifnotmodeDIFFUSIONFINALE{\includeversion{NOTDIFFUSIONFINALE}}
\ifnotmodeDIFFUSIONFINALE{\newcommand\FINPROVISOIRELFMI[1]{}}

\providecommand\FINPROVISOIRELFMI{
  \section{Bibliographic notes}

  The current text is highly based (sometimes copied) from:
  
  \bibliographystyle{alpha}
  \bibliography{/Users/bournez/bibliographie/Bib-Files/bournez-avantbib2DOI,/Users/bournez/bibliographie/Bib-Files/perso}
  \end{document}
  }

%\def\POLYCHAPTERMODE{}

\providecommand\ifnotPOLYCHAPTERMODE[1]{
\ifdefined\POLYCHAPTERMODE
\else
#1
\fi
}

\providecommand\ifmodePOLYCHAPTERMODE[1]{
\ifdefined\POLYCHAPTERMODE
#1
\fi
}


%%%%%%%%%%%%%%%%%%%%%%%%%%%%%%%%%%%%%%%%%%%%%%%%%
%
%
%                      POUR INDICATIONS: ETAT
%
%
%%%%%%%%%%%%%%%%%%%%%%%%%%%%%%%%%%%%%%%%%%%%%%%%%


\newcommand\ABSTRACTINTERNAL[1]{
  \montruc[color=black!30]{
  \hrule
  \hrule
  \danger   ABSTRACT:\\[0.5cm] #1

  \vspace{0.5cm}
  \hrule
  \hrule
  }
  }

\newcommand\ABSTRACT[1]{
  \newpage
  \montruc[color=black!20]{
  \hrule
  \hrule
  \danger   ABSTRACT:\\[0.5cm] #1

  \vspace{0.5cm}
  \hrule
  \hrule
  }
  }

  
\newcommand\PARTIEFINALISEE{
  \newpage
  \montruc[color=black!15]{
  \hrule
  \hrule
  \danger   FINALISE A PARTIR ICI
  \hrule
  \hrule
  }
  }
\newcommand\PARTIESEMIFINALISEE{
  \newpage
  \montruc[color=black!15]{
  \hrule
  \hrule
  \danger   1/2 FINALISE A PARTIR ICI
  \hrule
  \hrule
  }
  }


\newcommand\PARTIENONFINALISEE{
  \newpage
  \montruc[color=black!30]{
  \hrule
  \hrule
  \danger   NON FINALISE A PARTIR ICI
  \hrule
  \hrule
  }
  }


%%%%%%%%%%%%%%%%%%%%%%%%%%%%%%%%%%%%%%%%%%%%%%%%%
%
%
%                      Affichage titre du cours: générique
%
%
%%%%%%%%%%%%%%%%%%%%%%%%%%%%%%%%%%%%%%%%%%%%%%%%%


  \newcommand\POLYWARNING[1]{
    
    \chapter*{Preliminaries}

    
\begin{center} 
  These are notes for the course: \\[1cm]
  \textbf{#1}
\end{center}

\bigskip
\fbox{\fbox{
\begin{minipage}{10cm}
\begin{center}
These notes change from days to days.  \\
Please do  not distribute. \\
\end{center}
\end{minipage}
}}

\bigskip
\begin{center}
  \textbf{
    Version of \today.
    }
\end{center}

 % This document is not stable.

\vfill
\begin{center}
Any comment (even about typos, orthography) is welcome: \\
send an email to bournez@lix.polytechnique.fr
\end{center}
}


\newcommand\MYCOURSEON[2]{ 
	%% External self-references
	\ifnotPOLYMODE{
		\ifdefined\MAINAUXFILE
			\externaldocument{\MAINAUXFILE}
		\fi
		}
	
	\ifdefined\PASLAPREMIEREFOIS
		\newpage
	\else
		\makeindex

     		\begin{document}
	\fi
	\gdef\PASLAPREMIEREFOIS{}

	\ifdefined\XPOLY
		\XPAGEDEGARDE{\TITRECOURS\ifPOLYCHAPTERMODE{\\[0.5cm] \Large \bf \LANGUE{Chapter:}{Chapitre:}  #2}}{\SUBTITRECOURS}
		\ifPOLYCHAPTERMODE{
		\renewcommand\thesection{\arabic{section}}
		\chapter*{#2}
		}
	\else
	
      \title{ \ifnotmodeDIFFUSIONFINALE{Course  #1 -}
      \ifPOLYCHAPTERMODE{\LANGUE{Chapter:}{Chapitre:} }
      #2}
      \author{
        Olivier Bournez% \inst{1}
        \\[0.5cm] %{\myaddress}
      }
      \date{Version of \today}
      % \institute{\myaddress}
      




      \ifnotPOLYCHAPTERMODE{
      	\maketitle
		}
      \ifPOLYCHAPTERMODE{\ifdefined\DOPAGEDEGARDE
      			\maketitle
		\fi
		\renewcommand\thesection{\arabic{section}}
%		\renewcommand\thesubsection{\arabic{section}.\ara}
      		\chapter*{#2}}

      % \begin{abstract}
      % \end{abstract}
      \fi
      


      \DOINDEXPRINT
      
      %% 2021: (meme si je comprends pas)
      \ifmodeDIFFUSIONFINALE{
		\excludeversion{recherche}
	}


    }



%%%%%%%%%%%%%%%%%%%%%%%%%%%%%%%%%%%%%%%%%%%%%%%%%
%
%
%                      Affichage titre du cours: variantes
%
%
%%%%%%%%%%%%%%%%%%%%%%%%%%%%%%%%%%%%%%%%%%%%%%%%%

    
\newcommand\MYCOURSE[1]{
  \MYCOURSEON{\csname numerocours#1\endcsname}{\csname nomcours#1\endcsname}
  
    \ifnotmodeDIFFUSIONFINALE{}%\listoftodos}

  }
  
  \newcommand\MYCOURSEnew[1]{
  \MYCOURSEON{???}{#1}
  
  \ifnotmodeDIFFUSIONFINALE{}%\listoftodos}

  }

%   ---------------------------
%   Pour compatilité chapter
%  ----------------------------

\newcommand\PASPREMIEREPAGE[1]{}
\providecommand\mychapter[1]{% \PASPREMIEREPAGE{\newpage} 
  \renewcommand\PASPREMIEREPAGE[1]{##1}  \hrule\hrule \section*{#1}
  \hrule\hrule \  \\[1cm]}
\providecommand\chapter[1]{\PASPREMIEREPAGE{\newpage} 
  \renewcommand\PASPREMIEREPAGE[1]{##1}  \hrule\hrule \section*{#1}
  \hrule\hrule \  \\[1cm]}

%   ---------------------------
%   if not sub
%  ----------------------------


\providecommand\IFNOTSUB[1]{#1}
\IFNOTSUB{
  
%\documentclass{article}
\usepackage{versions}
\includeversion{pasweb}
%\input{macros}

}


 
%%%%%%%%%%%%%%%%%%%%%%%%%%%%%%%%%%%%%%%%%%%%%%%%%
%
%
%               Déclaration de cours
%
%  (l'argument 1 = numéro sert pas)
%
%%%%%%%%%%%%%%%%%%%%%%%%%%%%%%%%%%%%%%%%%%%%%%%%%


\newcommand\COURS[3]{
  \expandafter\newcommand\csname numerocours#2\endcsname{#1}
  \expandafter\newcommand\csname nomcours#2\endcsname{#3}
}



\newcommand\NANCY[3]{
  \expandafter\newcommand\csname numerocours#2\endcsname{#1}
  \expandafter\newcommand\csname nomcours#2\endcsname{#3}
}


\newcommand\INFquatrecentdouze[3]{
  \expandafter\newcommand\csname numerocours#2\endcsname{#1}
  \expandafter\newcommand\csname nomcours#2\endcsname{#3}
}

\newcommand\INFcinqsixun[3]{
  \expandafter\newcommand\csname numerocours#2\endcsname{#1}
  \expandafter\newcommand\csname nomcours#2\endcsname{#3}
}



\newcommand\MPRI[3]{
  \expandafter\newcommand\csname numerocours#2\endcsname{#1}
  \expandafter\newcommand\csname nomcours#2\endcsname{#3}
}

\newcommand\COURSEINCLUDE[2][]{
  \renewcommand\IFNOTSUB[1]{}
  \renewcommand\MYCOURSE[1]{\chapter{\ifnotmodeDIFFUSIONFINALE{#1} \csname nomcours##1\endcsname}
    \ifnotmodeDIFFUSIONFINALE{\mychapter{Fichier: #2}}
  }
   \renewcommand\MYCOURSEnew[1]{\chapter{\ifnotmodeDIFFUSIONFINALE{#1} ##1 }
    \ifnotmodeDIFFUSIONFINALE{\mychapter{Fichier: #2}}
  }
%  \renewcommand\FINSUBMPRI{}
  \input{#2}
}

\newcommand\STRESSCOURSEINCLUDE[1]{
  \COURSEINCLUDE[!STRESS!]{#1}
  }

  \newcommand\subCOURSEINCLUDE[1] {
    \NDLR{BEGIN:  J'inclus/ Je répète ICI #1 }
    \input{INCLUDE/#1}
    \NDLR{END:  J'inclus/ Je répète ICI #1 }
}

%%%%%%%%%%%%%%%%%%%%%%%%%%%%%%%%%%%%%%%%%%%%%%%%%
%
%
%                      POUR INDICATIONS: MESSAGES
%
%
%%%%%%%%%%%%%%%%%%%%%%%%%%%%%%%%%%%%%%%%%%%%%%%%%


% ancienne version
% \newcommand\montruc[1]{\begin{center} {\fbox{\fbox{{#1}
%         }}} \end{center}}




\newcommand\montruc[2][]{
  \ifnotmodeDIFFUSIONFINALE{
  \begin{maboitetruc}[#1]{\danger} #2
  \end{maboitetruc}
%%  \todo[inline,color=green!20,caption={2do}, ssss#1]{\begin{minipage}{\textwidth-4pt}
%\todo[inline,color=green!5%,caption={2do}
%]{\begin{minipage}{\textwidth-4pt}
%    #2\end{minipage}}
}
}

\newcommand\montrucdanger[3][]{
  \ifnotmodeDIFFUSIONFINALE{
                    \montruc[#1]{\danger #2: #3}
                    }
}

\newcommand\montrucdangermargin[3][]{
  \ifnotmodeDIFFUSIONFINALE{
  \montruc[#1]{\danger #2: #3}
\marginpar{\danger #2: #3}
}
 }



%%%%%%%%%%%%%%%%%%%%%%%%%%%%%%%%%%%%%%%%%%%%%%%%%
%
%
%                      GESTION DES TRUCS MARQUES VERSIONS
%
%
%%%%%%%%%%%%%%%%%%%%%%%%%%%%%%%%%%%%%%%%%%%%%%%%%

\makeatletter
\renewcommand\beginmarkversion{\montrucdanger[colback=yellow!10]{=== begin NOT}{}\@Vs@sffbox{\@currenvir$>$}}
 \renewcommand\endmarkversion{\@Vs@sffbox{$<$\@currenvir}\montrucdanger[colback=yellow!10]{end=== NOT}{}}
\makeatother 


%%%%%%%%%%%%%%%%%%%%%%%%%%%%%%%%%%%%%%%%%%%%%%%%%
%
%
%                      VARIANTES DE MESSAGES
%
%
%%%%%%%%%%%%%%%%%%%%%%%%%%%%%%%%%%%%%%%%%%%%%%%%%

 \newcommand\NDLR[1]{\montrucdanger[colback=blue!20]{NDLR}{#1}}
 \newcommand\NLDR[1]{\NDLR{#1}}
 \newcommand\VOIRAUSSI[1]{\montrucdanger[colback=yellow!20]{VOIR AUSSI}{#1}}
 \newcommand\TODO[1]{\montrucdanger[colback=green!20]{TODO}{#1}}
 \newcommand\ATTENTION[1]{\montruc[colback=brown!20]{\danger ATTENTION#1}}
 \newcommand\COMPRENDPAS[1]{\montrucdanger[colback=yellow!20]{JE
     COMPRENDS PAS CA}{#1}}
  \newcommand\POSSIBLE[1]{\montrucdanger[colback=yellow!10]{
      Serait possible}{#1}}
 \newcommand\ENLEVE[1]{\montrucdanger[colback=gray!20]{
      ENLEVE}{#1}}
 
 
  \newcommand\DEPLACE[1]{
    \vspace{0.5cm}
    \hrule
    \hrule
      \vspace{0.5cm}
    \montrucdangermargin[colback=black!40]{DEPLACE/MISSING CONCEPT
    }{#1}
      \vspace{0.5cm}
    \hrule
    \hrule
      \vspace{0.5cm}
  }

\newcommand\BESOINDE[1]{\montrucdangermargin[colback=pink!20]{ATTTENTION
     BESOIN DE}{#1}}
\newcommand\MISSINGCONCEPT[1]{\montrucdangermargin[colback=pink!20]{MISSING CONCEPT
     }{#1}}
\newcommand\SAUTE[1]{\montrucdangermargin[colback=pink!20]{SAUTE
  }{#1}}
\newcommand\SHOULDBEHERE[1]{\BEGINTEMPENVCOLOR{blue!30}\montrucdanger[colback=pink!10]{
SHOULD
    BE HERE
  }{    
    \vspace{1cm}
    
    \centerline{$\vdots$}
    
   \centerline{ #1}
    
      \centerline{$\vdots$}
          
    \vspace{1cm}
    }\ENDTEMPENVCOLOR}
\newcommand\COULDBEHERE[1]{\montrucdangermargin[colback=pink!20]{COULD
    BE HERE
  }{#1}}



%%%%%%%%%%%%%%%%%%%%%%%%%%%%%%%%%%%%%%%%%%%%%%%%%
%
%
%                      POUR INDICATIONS: SOURCES
%
%
%%%%%%%%%%%%%%%%%%%%%%%%%%%%%%%%%%%%%%%%%%%%%%%%%


\newcommand\SOURCEURL[1]{
  \ifnotmodeDIFFUSIONFINALE{
    \marginpar{ \textcolor{magenta}{Source: \url{#1}} }
    }
}

\newcommand\SOURCETXT[1]{
  \ifnotmodeDIFFUSIONFINALE{
    \marginpar{ \textcolor{magenta}{Source: {#1}} }
    }
}

\newcommand\SOURCE[1]{
  \nocite{#1}
  \ifnotmodeDIFFUSIONFINALE{
    \marginpar{ \textcolor{magenta}{Source: \cite{#1}} }
    }
}


%--------------------------------
%--- Raccourci Sources particulières  ---
%--------------------------------

\newcommand\SOURCEJECH{\SOURCE{jech2003set}}
\newcommand\SOURCEKECHRIS{\SOURCE{kechris2012classical}}
\newcommand\SOURCENOTESSETTHEORY{\SOURCE{moschovakis2006notes}}
\newcommand\SOURCEKOZENCT{\SOURCE{LivreKozenTC}}
\newcommand\SOURCEMARKER{\SOURCE{marker2002descriptive}}
\newcommand\SOURCEMOSCHO{\SOURCE{moschovakis2009descriptive}}



%%%%%%%%%%%%%%%%%%%%%%%%%%%%%%%%%%%%%%%%%%%%%%%%%
%
%
%                      POUR INDICATIONS: DELIRES/COMMENTAIRES
%
%
%%%%%%%%%%%%%%%%%%%%%%%%%%%%%%%%%%%%%%%%%%%%%%%%%


%%%%%% ABOUT PLAN

%\newcommand\PLANBUT[1]{\montruc{\begin{minipage}{0.8\textwidth} #1\end{minipage}}}


%%%%%% ABOUT RECHERCHE


\ifnotmodeDIFFUSIONFINALE{
\newenvironment{recherche}
{}{}
%{\color{blue} BEGIN-RECHERCHE:
%%\ifmodeDIFFUSIONFINALE{diffusion finale}
%%\ifnotmodeDIFFUSIONFINALE{not diffusion finale}
%%	\ifdefined\SAFEMODE
%%		SAFEMODE
%%	\else
%%		NOTSAFEMODE
%%	\fi
%  \todo[inline, color=blue]{recherche ici} }{
%  END-RECHERCHE \color{black} }
\tcolorboxenvironment{recherche}{colback=blue!5!white, colframe=red!75!black,fonttitle=\bfseries, colbacktitle=blue!85!black,enhanced,
attach boxed title to top center={yshift=-2mm},
  title={FAIRE RECHERCHE ICI}}
}


%%%%%% ABOUT DELIRE


%\ifnotmodeDIFFUSIONFINALE{
%\newenvironment{delire}{\color{violet} BEGIN-DELIRE: }{
%  END-DELIRE \color{black} }
%}


%%%%%% ABOUT PERSO


%\ifnotmodeDIFFUSIONFINALE{
%\newenvironment{perso}{\color{brown}
%  BEGIN-PERSO: }{
%  END-PERSO\color{black} }
%}



%%%%%% ABOUT COMMENTAIRE

\ifnotmodeDIFFUSIONFINALE{
\newenvironment{commentaire}{\begin{maboitecommentaire}[colback=brown!30]{Commentaire} Commentaire :}{\end{maboitecommentaire}}
\newcommand\COMMENTAIRE[2][Commentaire ]{
\begin{maboitecommentaire}[colback=brown!30]{#1}
Commentaire : #2
\end{maboitecommentaire}
}
}

%%%%%% ABOUT HORSUJET


%\ifnotmodeDIFFUSIONFINALE{
%\newenvironment{probablementhorsujet}{\color{gray}
%  BEGIN-PROBABLEMENT HORS SUJET: }{
%  END-PROBABLEMENT HORS SUJET \color{black} }
%}

%%%%%% ABOUT BONUS


\ifnotmodeDIFFUSIONFINALE{
\newenvironment{pasbesoinpourlinstant}{NOT NEEDEED FOR
  THIS SECTION: }{ }
    \tcolorboxenvironment{pasbesoinpourlinstant}{colback=green!5!white, colframe=gray!75!black,fonttitle=\bfseries, colbacktitle=green!85!black,enhanced,
attach boxed title to top left={yshift=-2mm},
  title={PAS BESOIN POUR L'INSTANT}}
}

%%%%%% ABOUT BONUS


\ifnotmodeDIFFUSIONFINALE{
\newenvironment{bonus}{BONUS: }{
  }
  \tcolorboxenvironment{bonus}{colback=gray!5!white, colframe=gray!75!black,fonttitle=\bfseries, colbacktitle=gray!85!black,enhanced,
attach boxed title to top left={yshift=-2mm},
  title={BONUS}}
}

%%%%%% ABOUT BONUS


\ifnotmodeDIFFUSIONFINALE{
\newcommand\ATTENTIONREPETE[1]{
  \ATTENTION{Repete (attention modif): #1}
  \begin{remarkolivier}
Repete (attention modif): #1
\end{remarkolivier}
}
}


%%%%%%%%%%%%%%%%%%%%%%%%%%%%%%%%%%%%%%%%%%%%%%%%%
%
%
%                      MACROS Papiers Paulin
%
%
%%%%%%%%%%%%%%%%%%%%%%%%%%%%%%%%%%%%%%%%%%%%%%%%%


\newcommand{\x}{\times}
\newcommand{\s}{\sigma}
\newcommand\sep{.}
\newcommand\bzero{\zero}
\newcommand\bun{\un}
\renewcommand{\a}{\alpha}
\def\hd{{\sf hd}}
\def\tl{{\sf tl}}
\def\cons{{\sf cons}}
\def\Cons{{\sf Cons}}
\def\Pr{{\sf Pr}}
\def\Op{{\sf Op}}
\def\Rel{{\sf Rel}}
\def\Equal{{\sf Equal}}
\def\Eval{{\sf Eval}}
\def\Gate{{\sf Gate}}
\def\Selection{{\sf Selection}}
\def\Result{{\sf Result}}
\def\Circuit{{\sf Circuit}}
\def\next{{\sf next}}
\def\comp{{\sf comp}}
\def\finalcomp{{\sf finalcomp}}
\def\Pick{{\sf Pick}}
\def\Hd{{\sf Hd}}
\def\Tl{{\sf Tl}}
\def\RevHd{{\sf RevHd}}
\def\Node{{\sf Node}}
\def\Length{{\sf Length}}
\def\Tape{{\sf Tape}}
\def\Head{{\sf Head}}
\def\Time{{\sf Time}}
\def\unpad{{\sf unpad}}
\def\concat{{\sf concat}}
\def\rg{{\sf right}}
\def\lf{{\sf left}}
\def\st{{\sf state}}
\def\tope{{\sf top}}


%%%%%%%%%%%%%%%%%%%%%%%%%%%%%%%%%%%%%%%%%%%%%%%%%
%
%
%                      MACROS POLY INF561
%
%
%%%%%%%%%%%%%%%%%%%%%%%%%%%%%%%%%%%%%%%%%%%%%%%%%

\newcommand\buglst[1]{#1}

%\newcommand\mydefref[1]{\expandafter\newcommand\csname #1\endcsname{\ref{#1}}}
%\mydefref{refthrisc}
%\mydefref{refthcinqhuit}
%\mydefref{refthcinqsix}
%\mydefref{refthfondamcaract}
%\mydefref{refdefcircuit}
%%% sinon, fait à la main
%\newcommand\refchapalgo{\ref{chapalgo}}
%\newcommand\refchapmodeles{\ref{chapmodeles}}
%\newcommand\refchappreliminaires{\ref{chappreliminaires}}
%\newcommand\refchapcircuits{\ref{chapcircuits}}
%\newcommand \refpropfondamen{\ref{propfondamen}}
%\newcommand\refthhierspace{ref{thhierspace}}
%\newcommand\refchapcomplexitetemps{\ref{chapcomplexitetemps}}
%\newcommand\refdefverifun{\ref{defverifun}}
%\newcommand\refdefverif{\ref{defverif}}
%\newcommand\refthmachinevs{ref{thmachinevs}}
%\newcommand\reflemnotationconfig{\ref{lemnotationconfig}}
%\newcommand\reflemfondam{\ref{lemfondam}}
%\newcommand\refthtroisat{\ref{thtroisat}}
%\newcommand\refthfondam{\ref{thfondam}}
%\newcommand\refthsansram{\ref{thsansram}}
%\newcommand\refdefsram{\ref{defsram}}
%!TEX encoding = UTF-8 Unicode

%% TRUCS DU POLY INF561

\ifnotSLIDE{
\LANGUE{
\spnewtheoremi{exemple}{\NDLR{ECRIT EXEMPLE AU LIEU DE
  EXAMPLE}}{\bfseries}{\rmfamily}{theorem} 
\spnewtheoremi{remarque}{\NDLR{ECRIT REMARQUE AU LIEU DE REMARK}}{\bfseries}{\rmfamily}{theorem} 
}
{
\spnewtheoremi{exemple}{Exemple}{\bfseries}{\rmfamily}{theorem} 
\spnewtheoremi{remarque}{Remarque}{\bfseries}{\rmfamily}{theorem} 
}
}
            
            \newcommand\PERSOSOURCEPOSSIBLE[1]{\PERSOSOURCE{#1}\PERSOCOMMENTAIRE{#1}}
\newcommand\PERSOPOSSIBLE[1]{\POSSIBLE{#1}}
% \newcommand\PERSOPOSSIBLEbug[1]{\NDLR{\textbf{-- possible --
%       enleve car pb compilation}}}
\newcommand\PERSOETAIT[1]{\NDLR{ETAIT: #1}}
%   \begin{minipage}{10cm}
% \NDLR{\textbf{\tiny  -- était --      #1 -- était
%     --
%   \end{minipage}
% }}}

\providecommand\nl{}

\newcommand\PERSOSOURCE[1] { \SOURCETXT{#1}}
\newcommand\PERSOCOMMENTAIRE[1]{\NDLR{#1}}
\newcommand\PERSOCOMMENTAIREd[2]{\NDLR{#1 \textbf{#2}}}


\newenvironment{dessinpb}
{%\shorthandoff{:}
%\selectlanguage{english}
\begin{center}\begin{tikzpicture}[baseline=(current bounding
      box.north)]
%babelbug with tikz
}
{\end{tikzpicture}
%\selectlanguage{frenchb}
\end{center}}



\newcommand\PERSOMANQUE[1]{\NDLR{\textbf{-- manque --
      #1 --manque--}}}

\newcommand\PERSOENLEVE[1]{\ENLEVE{#1}}






%%%%%%%%%%%%%%%%%%%%%%%%%%%%%%%%%%%%%%%%%%%%%%%%%
%
%
%                      MACROS POLY INF412
%
%
%%%%%%%%%%%%%%%%%%%%%%%%%%%%%%%%%%%%%%%%%%%%%%%%%

%\mydefref{refEEXTtruc}
%\mydefref{refEEXTchappreliminaires}
%\mydefref{refEEXTsecarbrebinaire}
%\mydefref{refEEXTsectermes}
%\mydefref{refEEXTchappredicats}
%\mydefref{refEXTchappredicats}
%\mydefref{refEXTchappreliminaires}
%\mydefref{refEEXTdefTuring}
%\mydefref{refEEXTpropnondet}
%\mydefref{refEEXTchapcompletude}
%\mydefref{refEEXTdefconfigalternative}
%\mydefref{refEEXTchapmachinesdeTuring}
%\mydefref{refEEXTpbdedecision}
%\mydefref{refEEXTchapcalc}
%\mydefref{refEEXTincompletude}
%\mydefref{refEEXTthhierspace}
%\mydefref{refEEXTchapcomplexitetemps}

\tikzstyle{nd} = [shape=circle,draw]


\newcommand\mmod{mod}

  %%% VALUATIONS
\newcommand\mmV{v}
\newcommand\mmVb{\overline{\mmV}}
\newcommand\mmVe{\mmVb}

%%% Corrections

\newcommand\exoref[1]{\ref{exo:#1-stt}}

\newenvironment{correction}[1]
{\par\medskip\noindent \label{exo:#1-cor} \hbox{\bf Exercice \ref{exo:#1-stt} 
  \   (page \pageref{exo:#1-stt}). }}
{\par\medskip}

% PROBLEMES DE DECISION

\newcommand\DECISIONint[3]{
%Le problème de décision #1
\begin{itemize}
\item[\textbf{Donnée}:] #2
\item[\textbf{Réponse}:] #3
\end{itemize}
}

\newcommand\INDECIDABLE[5]{
\DECISION{#1}{#2}{#3}

\begin{#4} \label{#5}
Le problème $#1$ % ATTENTION ENLEVE $#1$ en 2020. REMIS 2023
est indécidable.
\end{#4}
}


\newcommand\INDECIDABLEenglish[5]{
\DECISIONenglish{#1}{#2}{#3}

\begin{#4} \label{#5}
The problem $#1$
is undecidable.
\end{#4}
}

\newcommand\INDECIDABLEq[5]{
\DECISION{#1}{#2}{#3}

\begin{#4} \label{#5}
Le problème $#1$
est ?
\end{#4}
}

\newcommand\INDECIDABLEqenglish[5]{
\DECISION{#1}{#2}{#3}

\begin{#4} \label{#5}
The problem $#1$
is ?
\end{#4}
}



\newcommand\DECISIONintenglish[3]{
%Le problème de décision #1
\begin{itemize}
\item[\textbf{Input}:] #2
\item[\textbf{Answer}:] #3
\end{itemize}
}

\newcommand\DECISIONi[4]{
\item $#1$
%Le problème de décision #1
\DECISIONint{#1}{#2}{#3}
}

\newcommand\DECISIONienglish[4]{
\item $#1$
%Le problème de décision #1
\DECISIONintenglish{#1}{#2}{#3}
}


\newcommand\DECISION[3]{
\begin{definition}[$#1$]\ 

%Le problème de décision #1
\DECISIONint{#1}{#2}{#3}
\end{definition}
}


\newcommand\DECISIONenglish[3]{
\begin{definition}[$#1$]\ 
%
%Le problème de décision #1
\DECISIONintenglish{#1}{#2}{#3}
\end{definition}
}



\newcommand\NPC[4]{ 
\DECISION{#1}{#2}{#3}

\begin{#4}
Le problème $#1$
est $\NP$-complet. \LANGUE{\myindex{NP-completeness}}{\myindex{NP-completude@$\NP$-complétude}}
\end{#4}
}


\newcommand\NPCenglish[4]{ 
\DECISIONenglish{#1}{#2}{#3}
\begin{#4}
The problem $#1$
is $\NP$-complete. \LANGUE{\myindex{NP-completeness}}{\myindex{NP-completude@$\NP$-complétude}}
\end{#4}
}


\newcommand\NPCp[5]{
\DECISION{#1}{#2}{#3}

\begin{#4}[#5]
Le problème $#1$
est $\NP$-complet.
\end{#4}
}


\newcommand\NPCpenglish[5]{
\DECISIONenglish{#1}{#2}{#3}

\begin{#4}[#5]
The problem $#1$
is $\NP$-complete.
\end{#4}
}




\newcommand\BRUNOCOMMENTAIRE[1]{\NDLR{Bruno commentaire: #1}}

\ifnotSLIDE{
%\spnewtheorem{exercicec}{Exercice}{\bfseries}{\rmfamily}
%\spnewtheorem{exerciced}{Exercice}{\bfseries}{\rmfamily}
%\spnewtheorem{exercicedc}{Exercice}{\bfseries}{\rmfamily}
%\spnewtheorem{exercicecenglish}{Exercice}{\bfseries}{\rmfamily}
%\spnewtheorem{exercicedenglish}{Exercice}{\bfseries}{\rmfamily}
%\spnewtheorem{exercicedcenglish}{Exercice}{\bfseries}{\rmfamily}
\spnewtheorem{notation}{Notation}{\bfseries}{\rmfamily}
}

%%% INDEX
 \newcommand\motnouv[1]{\emph{#1}\index{#1}}%\SIGNALEMOTNOUV{#1}}%\index{#1}}
\newcommand\INTmotnouv[1]{\emph{#1}}
\newcommand\myindex[1]{\index{#1}}
%\marginpar{-#1=}

\newcommand\synonyme[2]{\kindex{#1|see{#2}}}
\newcommand\indexM[1]{\index{$#1$}}
\newcommand\idx[1]{#1\index{#1}}
\newcommand\idxm[1]{#1\index{$#1$}}
\newcommand\idxM[1]{$#1$\index{$#1$}}
%\newcommand\idxmBUG[1]{#1}
\newcommand\idxmm[2]{#1\index{#2@$#1$}}
\newcommand\idxi[2]{#1\index{#2}}
%\newcommand\motnouv[1]{\INTmotnouv{#1}\index{#1}}
 \newcommand\motnouvi[2]{\INTmotnouv{#1}\index{#2}} % indexage manuel
 \newcommand\motnouva[2]{\INTmotnouv{#1}\index{#2@#1}} % indexage: 2ième
%                                 % argumet = version sans accent
\newcommand\motnouvm[1]{\INTmotnouv{#1}} %indexage manuel fait ailleurs


\newenvironment{PERSOVERBATIM}{}{}
 \newcommand{\joNP}[3]{{ Problème \textbf{ { \motnouv{##1}}:}
            \begin{itemize}
            \item  \textbf{ Instance:} {##2} 
            \item   \textbf{ Question:} {##3}
            \end{itemize}
          }}

%%%%%%%%%%%%%%%%%%%%%%%%%%%%%%%%%%%%%%%%%%%%%%%%%
%
%
%                      POUR COMPATIBILITE AVEC SLIDES & AUTRE
%
%
%%%%%%%%%%%%%%%%%%%%%%%%%%%%%%%%%%%%%%%%%%%%%%%%%


\newcommand\defin[1]{\textbf{#1}}


%%%   COMPATIBILITE AVEC SLIDES

\newcommand\soustitre[1]{\textbf{#1}}


%% COMPATIBILITE AVEC MA THESE

\providecommand\motnouv[1]{\emph{#1}}
\providecommand\motnouvi[2]{\emph{#1}}

\newcommand\papiernorm[1]{\myop{norm}}

%%% OLIVIER
\newcommand\olivier[1]{ \ATTENTION{Olivier: #1}}


%%%%%%%%%%%%%%%%%%%%%%%%%%%%%%%%%%%%%%%%%%%%%%%%%
%
%
%                      Page de Garde Cours X
%
%
%%%%%%%%%%%%%%%%%%%%%%%%%%%%%%%%%%%%%%%%%%%%%%%%%


\includeversion{metdate}

\newcommand\XPAGEDEGARDE[2]{
\thispagestyle{empty}
\title{{#1} \\[2cm] %\\[3cm] 
{\large #2}
}
\author{
 Olivier Bournez% \inst{1}
\\[0.5cm] 
bournez@lix.polytechnique.fr%{\myaddress}
}
%\date{
%\vspace{2cm}
%\includegraphics[height=2.5cm]{/Users/bournez/lib/LaTeX/Perso/FIG-COMMONS/logo_X_new}
%}
\begin{metdate}
\date{
\vspace{2cm}
Version \LANGUE{of}{du} \today \\[1cm] 
\includegraphics[height=3.5cm]{/Users/bournez/lib/LaTeX/Perso/FIG-COMMONS/logo_X_new}
}
\end{metdate}
\maketitle
}


%%%%%%%%%%%%%%%%%%%%%%%%%%%%%%%%%%%%%%%%%%%%%%%%%
%
%
%                      MACROS UTILES
%
%
%%%%%%%%%%%%%%%%%%%%%%%%%%%%%%%%%%%%%%%%%%%%%%%%%

%%%                      Macros
 
%-------------
%--- vectors  ---
%-------------

 
\newcommand\vectorl[1]{{\mathbf#1}}
\newcommand\tu[1]{\vectorl{#1}}

\newcommand\vp{\vectorl{p}}
\newcommand\va{\vectorl{a}}
\newcommand\vb{\vectorl{b}}
\newcommand\vc{\vectorl{c}}
\newcommand\vf{\vectorl{f}}
\newcommand\vn{\vectorl{n}}
\newcommand\vx{\vectorl{x}}
\newcommand\vX{\vectorl{X}}
\newcommand\vy{\vectorl{y}}
\newcommand\vz{\vectorl{z}}
\newcommand\ve{\vectorl{e}}
\newcommand\vv{\vectorl{v}}
\newcommand\vr{\vectorl{r}}
\newcommand\vu{\vectorl{u}}
\newcommand\vA{\vectorl{A}}
\newcommand\vB{\vectorl{B}}
\newcommand\vC{\vectorl{C}}
\newcommand\vD{\vectorl{D}}

%-------------
%--- overline ---
%-------------


\newcommand\ox{\overline{x}}
\newcommand\oy{\overline{y}}
\newcommand\oc{\overline{c}}
\newcommand\oa{\overline{a}}
\newcommand\ow{\overline{w}}
\newcommand\od{\overline{d}}
\newcommand\ob{\overline{b}}
\newcommand\oP{\overline{P}}
\newcommand\oee{\overline{e}}
\newcommand\ot{\overline{t}}
\newcommand\ou{\overline{u}}
\newcommand\ov{\overline{v}}
\newcommand\oz{\overline{z}}
\newcommand\am{\overline{a}}
\newcommand\om{\overline{m}}
\newcommand\oj{\overline{j}}

%----------------
%--- ensembles ---
%----------------

\newcommand\Q{\mathbb{Q}}
\newcommand\R{\mathbb{R}}
\newcommand\Z{\mathbb{Z}}
\providecommand\C{\mathbb{C}}
\newcommand\N{\mathbb{N}}
\newcommand\K{\mathbb{K}}
\newcommand\F{\mathbb{F}}

\newcommand\dyadic{\mathbb{D}}
\newcommand\Zd{{\Z_2}}
\newcommand\Zp{{\Z_p}}
%
\newcommand\RP{\mathbb{R}^{>0}}
\newcommand\QP{\mathbb{Q}^{>0}}


%----------------
%--- lettres cal ---
%----------------

\newcommand\mK{\mathcal{K}}
\newcommand\mD{\mathcal{D}}
\newcommand\mA{\mathcal{A}}
\newcommand\mL{\mathcal{L}}
\newcommand\mX{\mathcal{X}}
\newcommand\mG{\mathcal{G}}
\newcommand\mE{\mathcal{E}}
\newcommand\mH{\mathcal{H}}
\newcommand\mR{\mathcal{R}}
\newcommand\mF{\mathcal{F}}
\newcommand\mJ{\mathcal{J}}
\newcommand\mO{\mathcal{O}}
\newcommand\mP{\mathcal{P}}
\newcommand\mN{\mathcal{N}}
\newcommand\mS{\mathcal{S}}
\newcommand\mM{\mathcal{M}}
\newcommand\mC{\mathcal{C}}
\newcommand\mV{\mathcal{V}}
\newcommand\mI{\mathcal{I}}
\newcommand\mT{\mathcal{T}}

\newcommand\fC{\mathcal{C}}

%----------------
%--- lettres frak ---
%----------------


\newcommand\kM{\mathfrak{M}}
\newcommand\kN{\mathfrak{N}}
\newcommand\kU{\mathfrak{U}}
\newcommand\kg{\mathfrak{g}}



%--------------------------
%--- Façon noter noms classes ---
%---------------------------

\newcommand{\cp}[1]{\operatorname{#1}}


%--------------------------
%--- Calculabilité  ---
%---------------------------

% -- classes
\LANGUE{
\newcommand\RE{\cp{CE}}
}
{\newcommand\RE{\cp{RE}}
}

\newcommand\Rec{\cp{D}}
\newcommand\PR{\cp{PR}}
\newcommand\Gregn{\mE_n}


%--------------------------
%--- Réductions  ---
%---------------------------

\newcommand\lem{\le_{m}}
\newcommand\lelog{\le_{L}}
\newcommand\equivlog{\equiv_{L}}



%--------------------------
%--- Complexité  ---
%---------------------------

% COMPLEXITE

% -- P
\newcommand{\Ptime}{\cp{P}}      % P
\newcommand\PTIME{\Ptime}
\newcommand\PTIMEs[1]{{\PTIME}_{#1}}
\newcommand\PTIMEsp[1]{{\PTIME}_{#1}^0}
\renewcommand\P{\PTIME}
\newcommand{\FPtime}{\cp{FPTIME}}

% -- NP
\newcommand\NP{\cp{NP}}
\newcommand{\NPtime}{\NP}
\newcommand\NPs[1]{\cp{NP}_{#1}}
\newcommand{\FNP}{\cp{FNP}}

\newcommand\NDP{\cp{NDP}}
\newcommand\coNDP{\cp{coNDP}}
\newcommand\NDPs[1]{\cp{NDP}_{#1}}

% -- coNP
\newcommand\coNP{\cp{coNP}}

% -- PSPACE
\newcommand\PSPACE{\cp{ PSPACE}}
\newcommand\NPSPACE{\cp{ NPSPACE}}
\newcommand{\Pspace}{\PSPACE}
\newcommand{\FPspace}{\cp{FPSPACE}}


% -- LOGSPACE
\newcommand\LSPACE{\cp{LOGSPACE}}
\newcommand\NLSPACE{\cp{NLOGSPACE}}
\newcommand\coNLSPACE{\cp{coNLOGSPACE}}
\newcommand\coNL{\cp{coNL}}

% -- Autres
\newcommand\EXPTIME{\cp{EXPTIME}}
\newcommand\NEXPTIME{\cp{NEXPTIME}}
\newcommand\EXPSPACE{\cp{EXPSPACE}}
\newcommand\PH{\cp{PH}}
\newcommand\NC{\cp{NC}}
\newcommand\AC{\cp{AC}}
\newcommand\PPAD{\cp{\mathbf{PPAD}}}
\newcommand\PLS{\cp{\mathbf{PLS}}}

% -- Reverse math
\newcommand{\WKL}{\ensuremath{\docheck{\mathsf{WKL}_0}}}
\newcommand{\ACA}{\ensuremath{\docheck{\mathsf{ACA}_0}}}
\newcommand{\ATR}{\ensuremath{\docheck{\mathsf{ATR}_0}}}
\newcommand{\RCAO}{\ensuremath{\docheck{\mathsf{RCA}_0}}}
\newcommand{\PiCA}{\ensuremath{\docheck{\Pi^1_1\!\text{-}\mathsf{CA}_0}}}


% -- circuits
\newcommand\CIRCUIT[2]{\cp{CIRCUIT(#1,#2)}}
\newcommand\UCIRCUIT[2]{\cp{UCIRCUIT(#1,#2)}}

% -- Classes proba
\newcommand\PP{\cp{PP}}
\newcommand\RPP{\cp{RP}}
\newcommand\ZPP{\cp{ZPP}}
\newcommand\coRP{\cp{coRP}}
\newcommand\BPP{{\cp{BPP}}}

% -- IP
\newcommand\IP{\cp{IP}}
\newcommand\PCP{\cp{PCP}}

% -- non-uniforme
\newcommand\nuP{\mathbb{P}}
\newcommand\nuPs[1]{\mathbb{P}_{#1}}
\newcommand\nuPsp[1]{\mathbb{P}_{#1}^0}
\newcommand\nuNP{\mathbb{N}\mathbb{P}}
%\newcommand\poly{/{poly}}
\newcommand\Ppoly{\PTIME_{\poly}}
\newcommand\NPpoly{\NP_{\poly}}

% -- conditions d'uniformité
\newcommand\Luniforme{\cp{L}}
\newcommand\BCuniforme{\cp{BC}}


% Classes de circuits
\newcommand\TC{\cp{TC}}
\newcommand\LFP{\cp{LFP}}
\newcommand\FAC{\cp{FAC}}
\newcommand\FACC{\cp{FACC}}
\newcommand\FTC{\cp{FTC}}
\newcommand\FNC{\cp{FNC}}

% Logiques temorelles
\newcommand\CTL{\cp{CTL}}
\newcommand\LTL{\cp{LTL}}

% -- mesure temps/espace/taille

\newcommand\DTIME{\cp{DTIME}}
\newcommand\TIME[1]{\cp{TIME(#1)}}
\newcommand\TIMEs[2]{\cp{TIME_{#2}(#1)}}
\newcommand\TIMEsp[2]{\cp{TIME^0_{#2}(#1)}}
\newcommand\NTIME[1]{\cp{NTIME(#1)}}
\newcommand\NTIMEs[2]{\cp{NTIME_{#2}(#1)}}
\newcommand\NTIMEsp[2]{\cp{NTIME_{#2}^0(#1)}}
\newcommand\DSPACE{\cp{DSPACE}}
\newcommand\SPACE[1]{\cp{SPACE(#1)}}
\newcommand\SPACEs[2]{\cp{SPACE_{#2}(#1)}}
\newcommand\NSPACE[1]{\cp{NSPACE(#1)}}
\newcommand\NSPACEs[2]{\cp{NSPACE_{#2}(#1)}}
\newcommand\coNSPACE[1]{\cp{coNSPACE(#1)}}
\newcommand\TIMEON[2]{\cp{TIME(#1,#2)}}
\newcommand\SPACEON[2]{\cp{SPACE(#1,#2)}}
\newcommand\SIZE[1]{\cp{ size(#1)}}
\newcommand\ATIME{\cp{ATIME}}
\newcommand\ASPACE{\cp{ASPACE}}


%% autre:
\newcommand\LH{\cp{LH}}
\newcommand\FO{\cp{FO}}
\newcommand\SO{\cp{SO}}
\newcommand\ESOhorn{\cp{ESO_{horn}}}
\newcommand\SOhorn{\cp{SO_{horn}}}
\newcommand\ACzero{\cp{AC_0}}
\newcommand{\LINSPACE}{\cp{LINSPACE}}
\newcommand{\mFLINSPACE}{\mF_{\cp{LINSPACE}}}
\newcommand{\LOGSPACE}{\cp{LOGSPACE}}
\newcommand\ALOGTIME{\cp{ALOGTIME}}
\newcommand\RF{\cp{RF}}


\newcommand\BOOL[1]{\cp{{BOOLEAN(#1)}}}
\newcommand\DET{\cp{ DET}}

% % -- dirty


\newcommand{\cPspace}{\cp{\sharp PSPACE}}
\newcommand{\cAPtime}{\cp{\sharp APTIME}}
%\newcommand{\FPspace}{\cp{FPSPACE}}
%\newcommand{\FPspaceN}{\cp{FPSPACE}^{+}}
%\newcommand{\FPspaceN}{\cp{FPSPACE}}



%-------------------------
%--- Problèmes de décision  ---
%--------------------------

\newcommand\decisionname[1]{\cp{#1}}


\newcommand\Luniv{\decisionname{L_{univ}}}
\newcommand\HALTINGPROBLEM{\decisionname{HALTING-PROBLEM}}
\newcommand\SAT{\decisionname{SAT}}
\newcommand\MAXtSAT{\decisionname{MAX3SAT}}
\newcommand\REACH{\decisionname{REACH}}
\newcommand\CIRCUITVALUE{\decisionname{CIRCUITVALUE}}
\newcommand\MONOTONECIRCUITVALUE{\decisionname{MONOTONECIRCUITVALUE}}
\newcommand\THEOREMS{\decisionname{THEOREMS}}
\newcommand\QCIRCUITSAT{\decisionname{QCIRCUITSAT}}
\newcommand\QSAT{\decisionname{QSAT}}
\newcommand\QBF{\decisionname{QBF}}
\newcommand\SATVALUE{\decisionname{SATVALUE}}
\newcommand\SUCCINTSAT{\decisionname{SUCCINTSAT}}
\newcommand\SUCCINTSATVALUE{\decisionname{SUCCINTSATVALUE}}
\newcommand\GEOGRAPHY{\decisionname{GEOGRAPHY}}
\newcommand\PARITY{\decisionname{PARITY}}
\newcommand\MAJ{\decisionname{MAJ}}
\newcommand\CVP{\CIRCUITVALUE}
\newcommand\BOOLEANPROGRAMVALUE{BOOLEAN~PROGRAM~VALUE}

\newcommand\CLIQUE{\cp{CLIQUE}}
\LANGUEinv{
\newcommand\RECOUVREMENTDESOMMETS{\cp{RECOUVREMENT~DE~SOMMETS}}
\newcommand\RS{\RECOUVREMENTDESOMMETS}
}{
\newcommand\RECOUVREMENTDESOMMETSeng{\cp{VERTEX~COVER}}
\newcommand\RSeng{\RECOUVREMENTDESOMMETSeng}
}

\LANGUEinv{\newcommand\COUPUREMAXIMALE{\cp{COUPURE~MAXIMALE}}}
{\newcommand\COUPUREMAXIMALEeng{\cp{MAXIMAL~CUT}}}
\newcommand\CIRCUITHAMILTONIEN{\CHAMILTONIEN}
\LANGUEinv{\newcommand\VOYAGEURDECOMMERCE{\VCOMMERCE}}
{\newcommand\VOYAGEURDECOMMERCEeng{\VCOMMERCEeng}}

\LANGUEinv{\newcommand\CIRCUITLEPLUSLONG{\cp{CIRCUIT~LE~PLUS~LONG}}}
{\newcommand\CIRCUITLEPLUSLONGeng{\cp{LONGEST~CIRCUIT}}}
\LANGUEinv{\newcommand\SOMMEDESOUSENSEMBLE{\SUBSETSUM}}{
\newcommand\SOMMEDESOUSENSEMBLEeng{\SUBSETSUMeng}}
\LANGUEinv{
\newcommand\SACADOS{\cp{SAC~A~DOS}}}
{\newcommand\SACADOSeng{\cp{KNAPSACK}}}

%français
\LANGUEinv{
\newcommand\CHAMILTONIEN{\cp{CIRCUIT~HAMILTONIEN}}}
{\newcommand\CHAMILTONIENeng{\cp{HAMILTONIAN~CIRCUIT}}}
\LANGUEinv{
\newcommand\VCOMMERCE{\cp{VOYAGEUR~DE~COMMERCE}}}
{\newcommand\VCOMMERCEeng{\cp{TRAVELING~SALESMAN}}}
\LANGUEinv{
\newcommand\kCOLORABILITE{\cp{k-COLORABILITE}}}
{\newcommand\kCOLORABILITEeng{\cp{k-COLORABILITY}}}
\LANGUEinv{
\newcommand\COLORIAGE{\cp{BON~COLORIAGE}}}
{\newcommand\COLORIAGEeng{\cp{GOOD~COLORING}}}
\LANGUEinv{
\newcommand\tCOLORABILITE{\cp{3-COLORABILITE}}}
{\newcommand\tCOLORABILITEeng{\cp{3-COLORABILITY}}}
\LANGUEinv{
\newcommand\SUBSETSUM{\cp{SOMME~DE~SOUS~ENSEMBLE}}}
{\newcommand\SUBSETSUMeng{\cp{SUBSET~SUM}}}
\newcommand\PARTITION{\cp{PARTITION}}
\LANGUEinv{
\newcommand\NOMBREPREMIER{\decisionname{NOMBRE~ PREMIER}}}
{\newcommand\NOMBREPREMIEReng{\decisionname{PRIME~NUMBER}}}
\LANGUEinv{
\newcommand\CODAGE{\decisionname{CODAGE}}}
{\newcommand\CODAGEeng{\decisionname{ENCODING}}}
\newcommand\kSAT{\cp{k-SAT}}
\newcommand\tSAT{\cp{3-SAT}}
\newcommand\NAESAT{\cp{NAESAT}}
\newcommand\NAEtSAT{\cp{NAE3SAT}}
\newcommand\NAEqSAT{\cp{NAE4SAT}}
\newcommand\STABLE{\cp{\LANGUE{INDEPENDANT~SET}{STABLE}}}

%\newcommand\CM{\cp{COUPURE MAXIMALE}}
\LANGUEinv{
\newcommand\POSTpb{\cp{Problème~de~ la~correspondance~de~Post}}}
{\newcommand\POSTpbeng{\cp{Post~correspondence~problem}}}

\LANGUEinv{
\newcommand\PARITE{\cp{PARITE}}}
{\newcommand\PARITEeng{\cp{PARITY}}}


%-------------------------
%--- Modèles parallèles  ---
%--------------------------


% PARALLELISME
\newcommand\RAM{\cp{RAM}}
%
\newcommand\PRAM{\cp{PRAM}}
\newcommand\CRCW{\cp{CRCW}}
\newcommand\CREW{\cp{CREW}}
\newcommand\EREW{\cp{EREW}}


%------------------------
%--- codages  ---
%------------------------


\newcommand\codage[1]{\langle #1 \rangle}
\newcommand\tuple[1]{\codage{#1}}
\newcommand\paire[2]{\codage{#1,#2}}
\newcommand\triplet[3]{\codage{#1,#2,#3}}
\newcommand\quadruplet[4]{\codage{#1,#2,#3,#4}}
\newcommand\cinquplet[5]{\codage{#1,#2,#3,#4,#5}}


%------------------------
%--- descriptive set theory  ---
%------------------------

\newcommand\baire{\mathcal{N}}
\newcommand\peano{\overline{\mathcal{N}}} %% A MOI
\newcommand\cantor{\mathcal{C}}
%
\newcommand\bSigma{\mathbf{\Sigma}}
\newcommand\bPi{\mathbf{\Pi}}
\newcommand\bDelta{\mathbf{\Delta}}
%


%%%%%% ABOUT NOMS QUI CHANGES

\newcommand\WF{\cp{WF}}
\newcommand\WO{\WF}
\newcommand\LO{\cp{LO}}
\newcommand\DLO{\cp{DLO}}
% 
\newcommand\rk{\cp{rank}}


%% Moins propre:

\newcommand\I{\mathbb{I}}
\renewcommand\H{\mathbb{H}}
%
\newcommand\leKB{<_{KB}}
%
\newcommand\finiomega{{<\omega}}
\newcommand\compl{^\frown}
\newcommand\prefix{ ~ prefix~  } 
\newcommand\X{{X}}
\newcommand\Y{\mathcal{Y}}
\newcommand\subsetstrict{\subset}

\renewcommand\complement[1]{#1^{c}}


%------------------------
%--- schémas de récursion ---
%------------------------

%%                      Source papier MFCS 2019


% COMPLEXITY

%% MFCS differe de Clote:
%% Version à taille restreinte

\newcommand\co{\cp{co}}

\newcommand\mFPTIME{\mF_{PTIME}}
\newcommand{\mFPSPACE}{\mF_{SPACE}}
\newcommand{\mFPH}{\mF_{PH}}

\newcommand\SigmaP{\Sigma^P}
\newcommand\DeltaP{\Delta^{P}}
\newcommand\PiP{\Pi^{P}}
\newcommand\coSigmaP{\co\SigmaP}
\newcommand\coPiP{\mF{\co\PiP}}

%% autres classes:


%% Schemas pour complexité implicite

\newcommand\COMP{\cp{ COMP}}
\newcommand\CRN{\cp{ CRN}}
\newcommand\BRN{\cp{ BRN}}
\newcommand\SBRN{\cp{ SBRN}}
\newcommand\BR{\cp{ BR}}
\newcommand\kBR{k-\cp{ BR}}
\newcommand \WBPR{\cp{ WBPR}}
\newcommand \BMIN{\cp{ BMIN}}
\newcommand \BSUM{\cp{ BSUM}}
\newcommand\SBSUM{\cp{ SBSUM}}
\newcommand \BPROD{\cp{ BPROD}}
\newcommand \SBPROD{\cp{ SBPROD}}
\newcommand \SCOMP{\cp{ SCOMP}}
\newcommand \SR{\cp{ SR}}
\newcommand \SRN{\cp{ SRN}}

%% devrait etre défini
\newcommand\WBRN{\cp{ WBRN}}



%-- godel
\newcommand\INTDIV{\cp{INTDIV}}
\newcommand\DIV{\cp{DIV}}
\newcommand\EVEN{\cp{EVEN}}
\newcommand\ODD{\cp{ODD}}
\newcommand\PRIME{\cp{PRIME}}
\newcommand\POWER{\cp{POWER}}
\newcommand\BIT{\cp{BIT}}
\newcommand\negHP{\overline{\cp{HP}}}
\newcommand\HP{\cp{HP}}
\newcommand\negLuniv{\overline{\Luniv}}
\newcommand\VALCOMP{\cp{VALCOMP}}
\newcommand\mWINDOW{\cp{LEGAL}}
\newcommand\mmWINDOW{\cp{WINDOW}}
\newcommand\LENGTH{\cp{LENGTH}}
\newcommand\DIGIT{\cp{DIGIT}}
\newcommand\MATCH{\cp{MATCH}}
\newcommand\MOVE{\cp{MOVE}}
\newcommand\START{\cp{START}}
\newcommand\CELL{\cp{CELL}}
\newcommand\HALT{\cp{HALT}}
\newcommand\SUBST{\cp{SUBST}}
\newcommand\PROOF{\cp{PROOF}}
\newcommand\PROVABLE{\cp{PROVABLE}}
\newcommand\CONSIST{\cp{CONSIST}}



%---------------------
%--- calculus (AP) ---
%---------------------

% \jacobian{f}{x}
\newcommand{\jacobian}[1]{J_{#1}}
%\newcommand{\grad}[1]{\overrightarrow{\operatorname{grad}}~{#1}}
\newcommand{\grad}[1]{\nabla{#1}}
\newcommand{\scalarprod}[2]{{#1}\cdot{#2}}
\newcommand{\bigO}[1]{\mathcal{O}\left(#1\right)}
\newcommand\smallO{o}
\newcommand{\softO}[1]{\tilde{\mathcal{O}}\left(#1\right)}
\newcommand{\taylor}[3]{{T_{#1}^{#2}#3}}
\newcommand{\transpose}[1]{{#1}^T}
\newcommand{\pastsup}[2]{{\sup}_{#1}#2}

\DeclareMathOperator\arctanh{arctanh}



%---------------
%--- norms (AP) ---
%---------------


\newcommand{\inorm}[2]{\left\lVert{#1}\right\rVert_{#2}}
\newcommand{\twonorm}[1]{\inorm{#1}{2}}
\newcommand{\infnorm}[1]{\inorm{#1}{}}
\newcommand{\normsup}[1]{\infnorm{#1}{}}
%
\newcommand\hausdorff{d_H}
%
\newcommand\norm[1]{\infnorm{#1}}


%----------------
%--- topology ---
%----------------

% \openball{pt}{radius}
 \newcommand{\openball}[2]{B_{#2}(#1)}
\newcommand{\closedball}[2]{\ovl{B}_{#2}(#1)}
\newcommand{\closure}[1]{\ovl{#1}}
% \newcommand{\interior}[1]{\mathring{#1}}
% \newcommand{\border}[1]{\partial{#1}}




%-----------------
%--- functions (AP) ---
%-----------------
 
\newcommand{\lambdafun}[2]{(#1)\mapsto{#2}}
\newcommand{\lambdafunex}[2]{#1\mapsto{#2}}
\newcommand{\argmin}{\operatornamewithlimits{argmin}}
\newcommand{\argmax}{\operatornamewithlimits{argmax}}


 
%-----------------------
%--- Machines de Turing  ---
%-----------------------

\newcommand\blanc{\vectorl{B}}



%-----------------
%--- Logique ---
%-----------------
 
\newcommand\sem[1]{{[\semspace[}#1{]\semspace]}}
\newcommand\sems[1]{{[\semspace[}#1{]\semspace]}}
\newcommand\semss[1]{\sem{#1}_{\sigma'}}
\newcommand\semspace{\kern -.25ex}
\newcommand\true{true}
\newcommand\false{false}
%extension élémentaire:
\newcommand\elem{\preceq}
\newcommand\godel[1]{\lceil #1 \rceil}

%-----------------
%--- Model theory ---
%-----------------

\newcommand\hull[1]{\langle #1 \rangle}

%-----------------
%--- Ordinaux ---
%-----------------

\newcommand\omegaunCK{\omega_1^{CK}}
\newcommand\omegaunck{\omegaunCK}

\DeclareMathOperator{\cof}{cof}

%-----------------
%--- Surréels (QG) ---
%-----------------

\newcommand{\On}{\textbf{On}}
\newcommand{\Nobf}{\textbf{No}}
\newcommand{\Nobb}{\ensuremath{\mathbbmss{No}}}
\newcommand{\Nolt}[1]{\ensuremath{\Nobf_{< #1}}}
\newcommand{\Noleq}[1]{\ensuremath{\Nobf_{\leq #1}}}
\newcommand{\Noltbb}[1]{\ensuremath{\Nobb_{< #1}}}
\providecommand\No{\Nobf}
\newcommand\Noalpha{\Nolt{\alpha}}
\newcommand\Nolambda{\Nolt{\lambda}}

\DeclareMathOperator{\length}{length}
\DeclareMathOperator{\supp}{supp} % operateur support


%----------------------
%---Maths  de base (AP)  ---
%----------------------

\DeclareMathOperator{\dom}{dom} % operateur support

\DeclareMathOperator{\size}{size}

\newcommand{\powerset}[1]{\mathcal{P}(#1)}
\newcommand{\card}[1]{{\##1}}
\newcommand\priv{ - }
\newcommand\prive{-}
\renewcommand{\restriction}{\mathord{\upharpoonright}}
\newcommand\restrict[1]{_{\restriction #1}}

% \providecommand\mod{}
%% UNDEFINE:
\let\mod\relax
\let\div\relax
\DeclareMathOperator{\mod}{mod} 
\DeclareMathOperator{\div}{div} 

% partieentier
\newcommand\entieres[1]{\lceil #1 \rceil}
\newcommand\entierei[1]{\lfloor #1 \rfloor}


%----------------------
%---Symbole danger ---
%----------------------


%\providecommand*{\TakeFourierOrnament}[1]{{%
%\fontencoding{U}\fontfamily{futs}\selectfont\char#1}}
%\providecommand*{\danger}{{\Huge \TakeFourierOrnament{66}}}



%%%%%%%%%%%%%%%%%%%%%%%%%%%%%%%%%%%%%%%%%%%%%%%%%
%
%
%                      ASMeries
%
%
%%%%%%%%%%%%%%%%%%%%%%%%%%%%%%%%%%%%%%%%%%%%%%%%%



% CIRCUITS
%-- particuliers
\newcommand\REPLACE{\cp{REPLACE}}
\newcommand\EQUAL{\cp{EQUAL}}
\newcommand\EVALUE{\cp{EVALUE}}
\newcommand\TVALUE{\cp{TVALUE}}
\newcommand\UPDATE{\cp{UPDATE}}
\newcommand\ASSIGN{\cp{ASSIGN}}
\newcommand\PAR{\cp{PAR}}
\newcommand\DO{\cp{DO}}


% ASM
\newcommand\emplacement[2]{(#1,#2)}
\newcommand\contenu[2]{\sem{(#1,#2)}}
\newcommand\affecte[3]{(#1,#2)\leftarrow#3}
\newcommand\ctl{ctl\_state}


%%%%%%%%%%%%%%%%%%%%%%%%%%%%%%%%%%%%%%%%%%%%%%%%%
%
%
%                      GPACqueries
%
%
%%%%%%%%%%%%%%%%%%%%%%%%%%%%%%%%%%%%%%%%%%%%%%%%%



%------------------------
%--- generable zoo (GPAC) ---
%------------------------

\newcommand{\myop}[1]{\operatorname{#1}}
\newcommand{\abs}{\myop{abs}}
\newcommand{\sg}{\myop{sg}}
\newcommand{\ip}[1]{\myop{ip}_{#1}}
\newcommand{\rnd}{\myop{rnd}}
\newcommand{\lil}{\myop{lil}}
\newcommand{\lxh}{\myop{lxh}}
\newcommand{\hxl}{\myop{hxl}}
\newcommand{\plil}{\myop{plil}}
\newcommand{\reach}{\myop{reach}}
%\newcommand{\mreach}{\myop{mreach}}
\newcommand{\watchdog}{\myop{watchdog}}
\newcommand{\monitor}{\myop{monitor}}
%\newcommand{\norm}{\myop{norm}}
\newcommand{\mx}{\myop{mx}}
\newcommand{\mn}{\myop{mn}}
\newcommand{\sample}{\myop{sample}}
\newcommand{\select}{\myop{select}}
\newcommand{\ssine}[1]{\sin_{#1}}
\newcommand{\clamp}{\myop{clamp}}

\newcommand{\fnsg}[3]{tanh((#1)*(#2)*(#3))}
\newcommand{\fnipone}[3]{(1+\fnsg{(#1)-1}{#2}{#3})/2.}
\newcommand{\fnipinf}[1]{#1-sin(2*pi*(#1))/2./pi}%
\newcommand{\fnlil}[5]{%
\fnipone{(#3)-(#1)+1-((#2)-(#1))/8.}{#4+log(1+4./(#2-#1)*(#5)*(#5))+log(1+4)}{8./(#2-#1)}%
*\fnipone{#2-#3+1-(#2-#1)/8.}{#4+log(1+4/(#2-#1)*(#5)*(#5))+log(1+4)}{8./(#2-#1)}%
*4./(#2-#1)*(#5)}
\newcommand{\fnlxh}[5]{\fnipone{(#3)-(#1+#2)/2.+1}{#4+log(1+(#5)*(#5))}{2./(#2-#1)}*(#5)}
\newcommand{\fnhxl}[5]{\fnipone{(#1+#2)/2.-(#3)+1}{#4+log(1+(#5)*(#5))}{2./(#2-#1)}*(#5)}
\newcommand{\fnplilf}[4]{sin(((#4)-((#1)+(#2))/2.+(#3)/4.)*2.*pi/(#3))}
\newcommand{\fnplil}[6]{(#6)*\fnlxh{\fnplilf{#1}{#2}{#3}{#1}}%
{\fnplilf{#1}{#2}{#3}{#1+((#2)-(#1))/4.}}%
{\fnplilf{#1}{#2}{#3}{#4}}%
{#5+2.+log(1+(#6)*(#6))}%
{1./4.+2./((#2)-(#1))}}
\newcommand{\fnssine}[2]{sin(#2)*(1-cos(#2)/cos(#1))}



%--------------------
%--- complexity GPAC ---
%--------------------


%\ifthenelse{\equal{#1}{}}{Empty.}{Nonempty.}
\newcommand{\myclass}[1]{\operatorname{#1}}
% \newcommand{\gcomp}[2]{\ensuremath{\myclass{GCOMP}(#1,#2)}}
% \newcommand{\ggenerable}[1]{\textcolor{red}{\ensuremath{#1-}\text{generable}}}
 \newcommand{\gval}[2][]{\ensuremath{\myclass{GVAL}_{#1}\ifthenelse{\equal{#2}{}}{}{[#2]}}}
 \newcommand{\gpval}[1][]{\ensuremath{\myclass{GPVAL}_{#1}}}
 \newcommand{\gc}[3][]{\ensuremath{\myclass{ATSC}(#2,#3)}}
 \newcommand{\gpc}[1][]{\ensuremath{\myclass{ATSP}}}
\newcommand{\cgc}[1][]{\ensuremath{\myclass{ATS}}}
\newcommand{\gexpc}[1][]{\ensuremath{\myclass{AEXP}}}
\newcommand{\gwc}[2]{\ensuremath{\myclass{AW}(#1,#2)}}
\newcommand{\gpwc}{\ensuremath{\myclass{AWP}}}
\newcommand{\cgwc}{\ensuremath{\myclass{AW}}}
\newcommand{\grc}[3]{\ensuremath{\myclass{AR}(#1,#2,#3)}}
\newcommand{\gprc}{\ensuremath{\myclass{ARP}}}
\newcommand{\gsc}[3]{\ensuremath{\myclass{AS}(#1,#2,#3)}}
\newcommand{\gpsc}{\ensuremath{\myclass{ASP}}}
\newcommand{\goc}[3]{\ensuremath{\myclass{AO}(#1,#2,#3)}}
\newcommand{\gpoc}{\ensuremath{\myclass{AOP}}}
\newcommand{\cgoc}{\ensuremath{\myclass{AO}}}
\newcommand{\guc}[4]{\ensuremath{\myclass{AX}(#1,#2,#3,#4)}}
\newcommand{\gpuc}{\ensuremath{\myclass{AXP}}}
\newcommand{\glc}[2][]{\ensuremath{\myclass{AL}(#2)}}
\newcommand{\gplc}[1][]{\ensuremath{\myclass{ALP}}}
\newcommand{\cglc}[1][]{\ensuremath{\myclass{AL}}}

%\newcommand{\intinterv}[2]{\{#1,\ldots,#2\}}
\newcommand{\intinterv}[2]{\llbracket#1,#2\rrbracket}

\newcommand{\Rgen}{\R_{G}}
\newcommand{\Rpoly}{\R_{P}}

%%%%%%%%%%%%%%%%%%%%%%%%%%%%%%%%%%%%%%%%%%%%%%%%%
%
%
%                      Mes codages obscures
%
%
%%%%%%%%%%%%%%%%%%%%%%%%%%%%%%%%%%%%%%%%%%%%%%%%%




%%% Codages

% \newcommand\gammaoc{\gamma_{[0,1]}}
% \newcommand\gammaonc{\gamma_{\N}}
\newcommand\gammaTMc{\gamma^{TM}_{[0,1]}}
\newcommand\gammaTMcatan{\gamma^{TM}_{(0,1)^2}}
\newcommand\gammaTMnc{\gamma^{TM}_{\N^3}}
\newcommand\gammaTMncd{\gamma^{TM}_{\N^2}}
% \newcommand\gammaTMe{\gamma^{TM}_{*}}
% \newcommand\gammas{\gamma_{sab}}
% \newcommand\gammaord{\gamma_{cantor}}


\newcommand\enccompact{\nu_{X}}
 \newcommand\encinfinite{\nu_\N}
 \newcommand\enc{\nu}
 
%%%%%%%%%%%%%%%%%%%%%%%%%%%%%%%%%%%%%%%%%%%%%%%%%
%
%
%                      DESSINS ET IMAGES
%
%
%%%%%%%%%%%%%%%%%%%%%%%%%%%%%%%%%%%%%%%%%%%%%%%%%


\newcommand\IMAGE[2]{ \begin{center} 
    \includegraphics[#1]{#2}
  \end{center}
  }


  
%%%% POUR FAIRE DESSIN.
\newenvironment{dessinpp}[1]{
\begin{tikzpicture}[baseline=(current bounding
      box.west),#1]
}
{\end{tikzpicture}}




\newenvironment{dessinp}[1]{
\begin{center}\begin{tikzpicture}[baseline=(current bounding
      box.north),#1]
}
{\end{tikzpicture}
\end{center}}


\newenvironment{dessin}{
\begin{center}\begin{tikzpicture}[baseline=(current bounding
      box.north)]
}
{\end{tikzpicture}
\end{center}}



\newenvironment{dessind}{
\begin{tikzpicture}[baseline=(current bounding
      box.north)]
}
{\end{tikzpicture}
}



%%%%%%%%%%%%%%%%%%%%%%%%%%%%%%%%%%%%%%%%%%%%%%%%%
%
%
%                      POUR DESSINS GPAC
%
%
%%%%%%%%%%%%%%%%%%%%%%%%%%%%%%%%%%%%%%%%%%%%%%%%%


%%%% MACROS DE BASE.


\newcommand\constant[2]
{
node (#1)  [shape=circle,draw] {#2};
\draw (#1.east) -- +(0.3,0) coordinate (#1-o) ; 
\draw (#1.west) -- +(-0.3,0) coordinate (#1-i) ; 
}

\newcommand\basicproduct[2]
{
node (#1) {#2};
\draw (#1) +(-0.4,0.4) -- +(0.4,0.4) -- +(0.8,0) -- +(0.4,-0.4) -- +(-0.4,-0.4) --
+(-0.4,0.4);
\draw (#1) ++(0.8,0) -- ++(+0.3,0) coordinate (#1-o) ;
\draw (#1) ++ (-0.4,0.2) -- +(-0.3,0) coordinate (#1-i1) ;
\draw (#1) ++(-0.4,-0.2) -- +(-0.3,0) coordinate (#1-i2) ;
}

\newcommand\product[1]{\basicproduct{#1}{$+\Pi$}}
\newcommand\negativeproduct[1]{\basicproduct{#1}{$-\Pi$}}


%%BEGIN INTERNAL
\newcommand\basicsummer[1]{
coordinate (#1) {};
\draw (#1) +(-0.5,0.6) -- +(0.5,0) -- +(-0.5,-0.6) -- +(-0.5,0.6);
\draw (#1) ++(0.5,0) -- +(+0.3,0) coordinate (#1-o) ;
}

\newcommand\basicintegrator[1]{
coordinate (#1) {};
\draw (#1) +(-0.5,0.6) -- +(0.5,0) -- +(-0.5,-0.6) -- +(-0.5,0.6);
\draw (#1) ++(0.5,0) -- +(+0.3,0) coordinate (#1-o) ;
\draw (#1) +(-0.5,-0.6) -| +(-0.8,0.6) -- +(-0.5,0.6) ;
 \draw (#1) ++(-0.65,0.6) -- +(0,+0.3) coordinate (#1-init) ;
}

\newcommand\splitfour[2]{
\draw (#2) ++ (#1,0.4) -- +(-0.3,0) coordinate (#2-i1) ;
\draw (#2) ++ (#1,0.15) -- +(-0.3,0) coordinate (#2-i2) ;
\draw (#2) ++ (#1,-0.15) -- +(-0.3,0) coordinate (#2-i3) ;
\draw (#2) ++ (#1,-0.4) -- +(-0.3,0) coordinate (#2-i4) ;
}

\newcommand\splitthree[2]{
\draw  (#2) ++(#1,0.3) -- +(-0.3,0) coordinate (#2-i1) ;
\draw  (#2) ++(#1,0) -- +(-0.3,0) coordinate (#2-i2) ;
\draw  (#2) ++(#1,-0.3) -- +(-0.3,0) coordinate (#2-i3) ;
}

\newcommand\splittwo[2]{
\draw  (#2) ++(#1,0.2) -- +(-0.3,0) coordinate (#2-i1) ;
\draw  (#2) ++(#1,-0.2) -- +(-0.3,0) coordinate (#2-i2) ;
}

\newcommand\splitone[2]{
\draw (#2) ++(#1,0) -- +(-0.3,0) coordinate (#2-i1) ;
}

%% END INTERNAL

\newcommand\summerfour[1]{
\basicsummer{#1}
\splitfour{-0.5}{#1}
}

\newcommand\summerthree[1]{
\basicsummer{#1}
\splitthree{-0.5}{#1}
}

\newcommand\summertwo[1]{
\basicsummer{#1}
\splittwo{-0.5}{#1}
}

\newcommand\summerone[1]{
\basicsummer{#1}
\splitone{-0.5}{#1}
}


\newcommand\integratorfour[1]{
\basicintegrator{#1}
\splitfour{-0.8}{#1}
<}


\newcommand\integratorthree[1]{
\basicintegrator{#1}
\splitthree{-0.8}{#1}
}


\newcommand\integratortwo[1]{
\basicintegrator{#1}
\splittwo{-0.8}{#1}
}


\newcommand\integratorone[1]{
\basicintegrator{#1}
\splitone{-0.8}{#1}
}

%%%%% FIN MACROS




%%%%%%%%%%%%%%%%%%%%%%%%%%%%%%%%%%%%%%%%%%%%%%%%%
%
%
%                      Trucs d'Amaury
%
%
%%%%%%%%%%%%%%%%%%%%%%%%%%%%%%%%%%%%%%%%%%%%%%%%%


%----------------
%--- ref (AP) ---
%----------------

% \newcommand{\defref}[1]{Definition~\ref{#1}}
\newcommand{\lemref}[1]{Lemma~\ref{#1}}
\newcommand{\thref}[1]{Theorem~\ref{#1}}
\newcommand{\amaurycorref}[1]{Corollary~\ref{#1}}
 \newcommand{\figref}[1]{Figure~\ref{#1}}
% \newcommand{\figrefmult}[1]{Figures~{#1}}
% \newcommand{\secref}[1]{Section~\ref{#1}}
% %\renewcommand{\algref}[1]{Algorithm~\ref{#1}}
 \newcommand{\propref}[1]{Proposition~\ref{#1}}
% \newcommand{\exref}[1]{Example~\ref{#1}}
\newcommand{\remref}[1]{Remark~\ref{#1}}


 %-------------------
%--- polynomials ---
%-------------------

\newcommand{\sigmap}[1]{{\Sigma{#1}}}
\newcommand{\degp}[1]{{\operatorname{deg}(#1)}}
\newcommand{\poly}{\operatorname{poly}}


 %----------------------
%--- misc functions ---
%----------------------

\newcommand{\sgn}[1]{\operatorname{sgn}(#1)}
\newcommand{\intp}{\operatorname{int}}
 \newcommand{\fracp}{\operatorname{frac}}
\newcommand{\fiter}[2]{#1^{[#2]}}
\newcommand{\frestrict}[2]{#1\restriction_{#2}}
\newcommand{\indicator}[1]{\mathds{1}_{#1}}
\newcommand{\idfun}{\operatorname{id}}
% \newcommand{\floor}[1]{\left\lfloor#1\right\rfloor}
 \newcommand{\ceil}[1]{\left\lceil#1\right\rceil}
\newcommand{\round}[1]{\left\lfloor#1\right\rceil}
\DeclareMathOperator\erf{erf}


%------------
%--- sets ---
%------------

\newcommand{\A}{\mathbb{A}}
%\newcommand{\D}{\mathbb{D}}
\newcommand{\Ns}{\mathbb{N}^*}
%\newcommand{\C}{\mathbb{C}}
%\newcommand{\K}{\mathbb{K}}
% \MAT{height}{[width]}{[field]}
\newcommand{\MAT}[3]{M_{#1\ifthenelse{\equal{#2}{}}{}{,#2}}\ifthenelse{\equal{#3}{}}{}{\left(#3\right)}}

%\newcommand{\Rcomp}{\R_{c}}
\newcommand{\Rp}{\R_{+}}
\newcommand{\Rs}{\R^{*}}
\newcommand{\Rps}{\R_{+}^{*}}

%\newcommand{\Analytic}{C^\omega}
%\newcommand{\PTIME}{\ensuremath{\operatorname{P}}}
%\newcommand{\EXPTIME}{\ensuremath{\operatorname{EXPTIME}}}
%\newcommand{\NP}{\ensuremath{\operatorname{NP}}}
%\newcommand{\NPTIME}{\NP}
\newcommand{\FP}{\ensuremath{\operatorname{FP}}}
%\newcommand{\PSPACE}{\ensuremath{\operatorname{PSPACE}}}
% \newcommand{\FPSPACE}{\ensuremath{\operatorname{FPSPACE}}}



%%%%%%%%%%%%%%%%%%%%%%%%%%%%%%%%%%%%%%%%%%%%%%%%%
%
%
%                      Trucs de Quentin
%
%
%%%%%%%%%%%%%%%%%%%%%%%%%%%%%%%%%%%%%%%%%%%%%%%%%


\newcommand{\fct}[5]{\ensuremath{#1 : \left\{\begin{array}{c c c}
#2 & \rightarrow & #3\\
#4 & \mapsto & #5
                                             \end{array}\right.}}

%Nouvelle fraction, plus jolie
\newcommand{\f}[2]{	
	\mathchoice%
		{\dfrac{#1}{#2}}
    	{\dfrac{#1}{#2}}
		{\frac{#1}{#2}}
		{\frac{#1}{#2}}
}


%%Ensembles et combinatoires

%Ensemble avec propriété à drotie \{x | blabla\}
\newcommand{\enstq}[2]{\left\{ \left.#1\ \vphantom{#2}\right|\ #2  \right\}}
%meme chose mais avec des crochets
\newcommand{\crotq}[2]{\left[ \left.#1\ \vphantom{#2}\right|\ #2  \right]}
%même chose mais avec des brackets
\newcommand{\bratq}[2]{\left\langle \left.#1\ \vphantom{#2}\right|\ #2  
  \right\rangle}

\newcommand{\intn}[2]{\ensuremath{
		\left\llbracket\, #1\ ;\ #2\,\right\rrbracket
              }}


%%%%%%%%%%%%%%%%%%%%%%%%%%%%%%%%%%%%%%%%%%%%%%%%%
%
%
%               ENVIRONNEMENTS THEOREMS & CO
%
%
%%%%%%%%%%%%%%%%%%%%%%%%%%%%%%%%%%%%%%%%%%%%%%%%%


%----------------
%--- pas dans lipics ---
%-----------------


\ifnotSLIDE{
% commente car bug Fevrier 2026: \ifnotTDEXAM{\spnewtheorem{solution}{Solution}{\bfseries}{\rmfamily}}
\spnewtheorem{openproblem}{Open Problem}{\bfseries}{\rmfamily}
\spnewtheorem{remarkolivier}{Commentaire d'Olivier: }{\bfseries}{\rmfamily}
\spnewtheorem{postulate}{Postulate}{\bfseries}{\rmfamily}
}

%----------------
%--- dans lipics ---
%-----------------
%
\newenvironment{exercice}[1]{\begin{exercise} \label{exo:#1-stt}}{\end{exercise}}
\newenvironment{exerciced}[1]{\begin{exercisee}\label{exo:#1-stt}}{\end{exercisee}}
\newenvironment{exercicec}[1]{\begin{exercise} \label{exo:#1-stt} (\LANGUE{solution on}{corrigé} page \pageref{exo:#1-cor})
}{\end{exercise}}
\newenvironment{exercicedc}[1]{\begin{exercisee} \label{exo:#1-stt} (\LANGUE{solution on}{corrigé} page \pageref{exo:#1-cor})
}{\end{exercisee}}

\newenvironment{slideproposition}{{\bf Proposition}}{}

\ifnotSLIDE{
\nolipics{
\nolncs{\spnewtheorem{theorem}{\LANGUE{Theorem}{Théorème}}{\bfseries}{\rmfamily}}
\nolncs{\spnewtheoremi{definition}{\LANGUE{Definition}{Définition}}{\bfseries}{\rmfamily}{theorem}}
\nolncs{\spnewtheoremi{lemma}{\LANGUE{Lemma}{Lemme}}{\bfseries}{\rmfamily}{theorem}}
\nolncs{\spnewtheoremi{corollary}{\LANGUE{Corollary}{Corollaire}}{\bfseries}{\rmfamily}{theorem}}
\nolncs{\spnewtheoremi{proposition}{Proposition}{\bfseries}{\rmfamily}{theorem}}
\nolncs{\spnewtheoremi{example}{\LANGUE{Example}{Exemple}}{\bfseries}{\rmfamily}{theorem}}
%\spnewtheorem{sketch}{{\bf \LANGUE{Sketch of proof}{Démonstration (principe)}}}{\bfseries}{\rmfamily}
\nolncs{\spnewtheoremi{remark}{\LANGUE{Remark}{Remarque}}{\bfseries}{\rmfamily}{theorem}}
\nolncs{\spnewtheorem{exercise}{\LANGUE{Exercise}{Exercice}}{\bfseries}{\rmfamily}}
\spnewtheorem{exercisee}{\LANGUE{*Exercise}{*Exercice}}{\bfseries}{\rmfamily}
\spnewtheorem{summary}{\LANGUE{Summary}{Résumé}}{\bfseries}{\rmfamily}
 %\spnewtheorem{exerciceenglish}{Exercise}{\bfseries}{\rmfamily}
 %
 \newenvironment{sketch}{{\bf \LANGUE{Sketch of proof}{Démonstration (principe)}}:}{\hfill $\Box$ }
\nolncs{\spnewtheorem{conjecture}{\LANGUE{Conjecture}{Conjecture}}{\bfseries}{\rmfamily}}
 %%
 \spnewtheorem{theoremaprouver}{\LANGUE{Theorem TO BE PROVED}{Théorème (A PROUVER)}}{\bfseries}{\rmfamily}
  \spnewtheorem{conjectureaprouver}{\LANGUE{Conjecture TO BE VERIFIED}{Conjecture (A VERIFIER)}}{\bfseries}{\rmfamily}
 }}
            
 \makeatletter
 \ifnotSLIDE{
  \ifnotTDEXAM{
\@ifundefined{proof}{
	\newenvironment{proof}{{\bf \LANGUE{Proof}{Démonstration}}:}{\hfill $\Box$ 
	}{}
	}}}
\makeatother

 
%%%%%%%%%%%%%%%%%%%%%%%%%%%%%%%%%%%%%%%%%%%%%%%%%
%
%
%                      Macros dont pas fier du tout
%
%
%%%%%%%%%%%%%%%%%%%%%%%%%%%%%%%%%%%%%%%%%%%%%%%%%



\newcommand\equations[1]{
\begin{array}{lll}
#1
\end{array}
}



\newcommand\systeme[1]{
\left\{
\begin{array}{lll} 
#1
\end{array}
\right.
}



%------------------
%--- shorthands (AP)---
%------------------
\newcommand{\mtt}[1]{\mathtt{#1}}
\newcommand{\ovl}[1]{\overline{#1}}
\newcommand{\panic}[1]{\textcolor{red}{#1}}
%\NewEnviron{epanic}{{\color{red}\BODY}}
\newcommand{\idest}{\emph{i.e.\thinspace}}


%------------------
%--- mes trucs  ---
%------------------


\newcommand\implique{\Rightarrow}
\newcommand\ssi{\Leftrightarrow}
\newcommand\ac[1]{\{#1\}}
\newcommand\zero{\vectorl{0}}
\newcommand\un{\vectorl{1}}

\newcommand\technique[1]{{\scriptsize #1}}


%% Preuve and Co

\newcommand\sensdirect{ Direction $\Rightarrow$. }
\newcommand\sensindirect{ Direction $\Leftarrow$. }




\newcommand\donnee{\mbox{<donnée>}}



% pour aller vite

\newcommand\gM{\mathfrak{M}}
\newcommand\gMM{(M,f_1,\cdots,f_u,r_1,\cdots,r_v)}
%
\newcommand\mygM{\mK}
\newcommand\mygMM{\left ( \K, \{op_i\}_{i\in I}, r_1 \ldots,
r_\ell, \zero,\un \right )}
\newcommand\myM{\K} 

\newcommand\gMMM{(M,f_1,\cdots,f_u,c_1,\cdots,c_k,r_1,\cdots,r_v)}
\newcommand\gMMMM{(f_1,\cdots,f_u,r_1,\cdots,r_v)}
\newcommand\gMME{(M,f_1,\cdots,f_u,f'_1,\cdots,f'_\ell,r_1,\cdots,r_v)}
\newcommand\gMMS{(f_1,\cdots,f_u,f'_1,\cdots,f'_\ell,r_1,\cdots,r_v)}

\newcommand\bool{\mathfrak{B}}
\newcommand\gbool{(\{0,1\},\zero,\un,=)}

\newcommand\osg{\overline{sg}}
\newcommand\moins{\mathrel{\dot{-}}}


%%% TRUC DE BARWISE


\newcommand\UM{\kU_{\kM}}
\newcommand\VV{\mathbb{V}}
% Hereditarily finite
\newcommand\IHF{\mathbb{I}HF}
\newcommand\IHP{\IHYP}
\newcommand\IHYP{\mathbb{I}HYP}

\providecommand\citep[1]{\cite{#1}}

%% Cofinali


% probas
% défini dans amsmath
\renewcommand\Pr{\cp{Pr}}
\newcommand\Var{\cp{Var}}


% -- mesures
%\newcommand\card{card}
\newcommand\taille{\LANGUE{size}{taille}}
\newcommand\entree{e}
\newcommand\profondeur{\LANGUE{depth}{profondeur}}


%%%%%%%%%%%%%%%%%%%%%%%%%%%%%%%%%%%%%%%%%%%%%%%%%%%%%%%%%%%%%%%%%%%%%%%%%%%%
%                 MACHINES DE TURING (graphique)
%%%%%%%%%%%%%%%%%%%%%%%%%%%%%%%%%%%%%%%%%%%%%%%%%%%%%%%%%%%%%%%%%%%%%%%%%%%%



% Configuration de machine de Turing
\newcommand\configTuring[3]{#1\textbf{#2}#3}


%%% DESSIN CROIX PORU AIDER  A DESSINER

\newcommand\croix{  +(-1,-1) -- +(1,1) +(-1,1) --
  +(1,-1) }

%% DESSIN ACCOLADE

\newcommand\ACCOLADE[3]{
%#1: coordonnées de départ
%#2: coordonnées arrivée
%#3: texte
\draw[decorate,decoration={brace}]  (#1) -- (#2)
 node[midway,above=0.1cm] {#3};
}

%%%


\def\lcase{3.5 ex}
\def\hcase{3.5ex}
\tikzstyle{case} = []
%[draw,minimum width = \lcase,minimum height = 3.6ex,baseline]
\newcommand\dcase{ +(-0.5*\lcase,-0.5*\hcase) rectangle +(0.5*\lcase,0.5*\hcase)}
%\newcommand\lettre[1]{\node [name=l,case,on chain=1] {#1};}


\newcommand\RUBANMACHINEDETURING[6]{
%#1: texte
%#2: position de la tête
%#3: nb de blancs à gauche
%#4: nb de cases totales
%#5: texte au dessus du ruban
%#6: position du coin
    \draw (#6) +(\lcase,4ex ) node {#5};
    \draw (#6) +(-0.6,0)  node  {\ldots};  
    \draw (#6) node[coordinate] (t) {};
    \foreach \x in {1,2,...,#3} {
         \draw (t) node [case]   {$\blanc$} \dcase;
         \draw (t) ++(\lcase,0) node [coordinate] (t) {};
       };
   \StrLen{#1}[\Resultat] % ATTENTION IL FAUT CETTE SYNTAXE POUR FAIRE
                         % EQUIVALENT D UN EDEF
  
\foreach \x in {1,...,\Resultat} 
             { 
        \draw (t) node [case]   {\StrChar{#1}{\x}} \dcase;
        \draw (t) ++(\lcase,0) node [coordinate] (t) {};
             }
\pgfmathtruncatemacro{\z}{#4 - \Resultat - #3}

  \foreach \x in {1,2,..., \z} {
        \draw (t) node [case]   {$\blanc$} \dcase;
        \draw (t) ++(\lcase,0) node [coordinate] (t) {};
}
%% Convention: graphiquement, la case numéro n se trouve à (n*\lcase,0)
%% en comptant les case à partir de $0$

   \draw  (t) +(0.1,0) node [name=r] {\ldots}; 
}

\newcommand\RUBANMACHINEDETURINGtexte[6]{
%#1: texte
%#2: position de la tête
%#3: nb de blancs à gauche
%#4: nb de cases totales (équivalent)
%#5: texte au dessus du ruban
%#6: position du coin
    \draw (#6) +(\lcase,4ex ) node {#5};
    \draw (#6) +(-0.6,0)  node  {\ldots};  
    \draw (#6) node[coordinate] (t) {};
   \foreach \x in {1,2,...,#3} {
        \draw (t) node [case]   {$\blanc$} \dcase;
        \draw (t) ++(\lcase,0) node [coordinate] (t) {};
      };
   \draw (#6) ++(-2*\lcase,0.5*\lcase)-- +(\lcase*#4,0);
   \draw (#6) ++(-2*\lcase,-0.5*\lcase) -- +(\lcase*#4,0);
   
    \draw (t) ++(-0.5*\lcase,0) node [case,anchor=west]  (te) {#1};
    \draw (te.east)  node [anchor=west,coordinate] (t) {};

  \draw (t) ++(0.2*\lcase,0) node [coordinate] (t) {};
   \foreach \x in {1,2,...,#3} {
        \draw (t) node [case]   {$\blanc$} \dcase;
        \draw (t) ++(\lcase,0) node [coordinate] (t) {};
      };

    \draw  (t) +(0.1,0) node [anchor=west] {\ldots}; 
}


\newcommand\RUBANMACHINEDETURINGTETE[6]{
%
% dessine la tete
%
% output: (tete)
%
\RUBANMACHINEDETURING{#1}{#2}{#3}{#4}{#5}{#6}
% TETE
     \draw (#6) + ( #2 * \lcase + #3  *\lcase,- 0.5*\lcase )
     node[coordinate] (k) {};
     \draw[->] (k) +(0,-0.5) -- (k);
     \draw (k) +(0,-0.5) node[coordinate] (tete) {};
}


\newcommand\RUBANMACHINEDETURINGTETEtexte[6]{
%
% dessine la tete
%
% output: (tete)
%
\RUBANMACHINEDETURINGtexte{#1}{#2}{#3}{#4}{#5}{#6}
% TETE
     \draw (#6) + ( #2 * \lcase + #3  *\lcase,- 0.5*\lcase )
     node[coordinate] (k) {};
     \draw[->] (k) +(0,-0.5) -- (k);
     \draw (k) +(0,-0.5) node[coordinate] (tete) {};
}


\newcommand\MACHINEDETURING[6]{
%#1: texte
%#2: position de la tête
%#3: état
%#4: nb de blancs à gauche
%#5: nb de cases totales
%#6: texte en haut du ruban
%\begin{tikzpicture}[node distance=-0.15mm]

\RUBANMACHINEDETURINGTETE{#1}{#2}{#4}{#5}{#6}{0,0}
%
%% ETAT:
     \draw (tete) node [name=etat,draw,circle,anchor=north]  {#3};
     \draw (etat) -- (tete);
%\end{tikzpicture}
}

\newcommand\MACHINEDETURINGtexte[6]{
%#1: texte
%#2: position de la tête
%#3: état
%#4: nb de blancs à gauche
%#5: nb de cases totales
%#6: texte en haut du ruban
%\begin{tikzpicture}[node distance=-0.15mm]

\RUBANMACHINEDETURINGTETEtexte{#1}{#2}{#4}{#5}{#6}{0,0}
%
%% ETAT:
     \draw (tete) node [name=etat,draw,circle,anchor=north]  {#3};
     \draw (etat) -- (tete);
%\end{tikzpicture}
}

\newcommand\MACHINEDETURINGTROISRUBANS[9]{
%#1: texte1
%#2: position de la tête 1
%#3: texte2
%#4: position de la tête 2
%#5: texte3
%#6: position de la tête 3
%#7: état
%#8: nb de blancs à gauche 1
%#9: nb de cases totales
%\begin{tikzpicture}[node distance=-0.15mm]

\RUBANMACHINEDETURINGTETE{#1}{#2}{#8}{#9}{\LANGUE{tape}{ruban} $1$}{0,0}
\draw (tete) node[coordinate] (tete1) {};
\draw (tete1) +(8,0) node[coordinate] (tete1b) {};
\draw (tete1) -| (tete1b) ;

\RUBANMACHINEDETURINGTETE{#3}{#4}{#8}{#9}{\LANGUE{tape}{ruban} $2$}{0,-2}
\draw (tete) node[coordinate] (tete2) {};
\draw (tete2) +(-8.5,0) node[coordinate] (tete2b) {};
\draw (tete2) -| (tete2b) ;

\RUBANMACHINEDETURINGTETE{#5}{#6}{#8}{#9}{\LANGUE{tape}{ruban} $3$}{0,-4}
\draw (tete) node[coordinate] (tete3) {};

%% ETAT:
     \draw (tete3) node [name=etat,draw,circle,anchor=north]  {#7};
     \draw (etat) -- (tete3);
    \draw  (etat) -|  (tete2b);
\draw  (etat) -|  (tete1b);
%\end{tikzpicture}
}

\newcommand\MACHINEDETURINGTROISRUBANSENGLISH[9]{
%#1: texte1
%#2: position de la tête 1
%#3: texte2
%#4: position de la tête 2
%#5: texte3
%#6: position de la tête 3
%#7: état
%#8: nb de blancs à gauche 1
%#9: nb de cases totales
%\begin{tikzpicture}[node distance=-0.15mm]

\RUBANMACHINEDETURINGTETE{#1}{#2}{#8}{#9}{tape $1$}{0,0}
\draw (tete) node[coordinate] (tete1) {};
\draw (tete1) +(8,0) node[coordinate] (tete1b) {};
\draw (tete1) -| (tete1b) ;

\RUBANMACHINEDETURINGTETE{#3}{#4}{#8}{#9}{tape $2$}{0,-2}
\draw (tete) node[coordinate] (tete2) {};
\draw (tete2) +(-8.5,0) node[coordinate] (tete2b) {};
\draw (tete2) -| (tete2b) ;

\RUBANMACHINEDETURINGTETE{#5}{#6}{#8}{#9}{tape $3$}{0,-4}
\draw (tete) node[coordinate] (tete3) {};

%% ETAT:
     \draw (tete3) node [name=etat,draw,circle,anchor=north]  {#7};
     \draw (etat) -- (tete3);
    \draw  (etat) -|  (tete2b);
\draw  (etat) -|  (tete1b);
%\end{tikzpicture}
}



\newcommand\MACHINEDETURINGTROISRUBANStexte[9]{
%#1: texte1
%#2: position de la tête 1
%#3: texte2
%#4: position de la tête 2
%#5: texte3
%#6: position de la tête 3
%#7: état
%#8: nb de blancs à gauche 1
%#9: nb de cases totales
% \begin{tikzpicture}[node distance=-0.15mm]

\RUBANMACHINEDETURINGTETEtexte{#1}{#2}{#8}{#9}{\LANGUE{tape}{ruban} $1$}{0,0}
\draw (tete) node[coordinate] (tete1) {};
\draw (tete1) +(8,0) node[coordinate] (tete1b) {};
\draw (tete1) -| (tete1b) ;

\RUBANMACHINEDETURINGTETEtexte{#3}{#4}{#8}{#9}{\LANGUE{tape}{ruban} $2$}{0,-2}
\draw (tete) node[coordinate] (tete2) {};
\draw (tete2) +(-8.5,0) node[coordinate] (tete2b) {};
\draw (tete2) -| (tete2b) ;

\RUBANMACHINEDETURINGTETEtexte{#5}{#6}{#8}{#9}{\LANGUE{tape}{ruban} $3$}{0,-4}
\draw (tete) node[coordinate] (tete3) {};

%% ETAT:
     \draw (tete3) node [name=etat,draw,circle,anchor=north]  {#7};
     \draw (etat) -- (tete3);
    \draw  (etat) -|  (tete2b);
\draw  (etat) -|  (tete1b);
%\end{tikzpicture}
}


\newcommand\MACHINEDETURINGDEUXRUBANStexte[9]{
%#1: texte1  
%#2: position de la tête 1 
%#3: texte2 inutilisée
%#4: position de la tête 2 inutilisée
%#5: texte3
%#6: position de la tête 3
%#7: état
%#8: nb de blancs à gauche 1
%#9: nb de cases totales
% \begin{tikzpicture}[node distance=-0.15mm]

\RUBANMACHINEDETURINGTETEtexte{#1}{#2}{#8}{#9}{\LANGUE{tape}{ruban} $1$}{0,0}
\draw (tete) node[coordinate] (tete1) {};
\draw (tete1) +(8,0) node[coordinate] (tete1b) {};
\draw (tete1) -| (tete1b) ;

% \RUBANMACHINEDETURINGTETEtexte{#3}{#4}{#8}{#9}{\LANGUE{tape}{ruban} $2$}{0,-1.5}
% \draw (tete) node[coordinate] (tete2) {};
% \draw (tete2) +(-5.5,0) node[coordinate] (tete2b) {};
% \draw (tete2) -| (tete2b) ;

\RUBANMACHINEDETURINGTETEtexte{#5}{#6}{#8}{#9}{\LANGUE{tape}{ruban} $2$}{0,-1.5}
\draw (tete) node[coordinate] (tete3) {};

%% ETAT:
     \draw (tete3) node [name=etat,draw,circle,anchor=north]  {#7};
     \draw (etat) -- (tete3);
%   \draw  (etat) -|  (tete2b);
\draw  (etat) -|  (tete1b);
%\end{tikzpicture}
}


\def\argdix{2}
\def\argonze{25}
\newcommand\MACHINEDETURINGQUATRERUBANStextep[9]{
%#1: texte1
%#2: position de la tête 1
%#3: texte2
%#4: position de la tête 2
%#5: texte3
%#6: position de la tête 3
%#7: texte4
%#8: position de la tête 3
%#9: état
%#10: nb de blancs à gauche 1
%#11: nb de cases totales
% \begin{tikzpicture}[node distance=-0.15mm]

%% BUG sur #10 devenu \argdix
%% BUG SUR #11 devenu \argonze

\RUBANMACHINEDETURINGTETEtexte{#1}{#2}{\argdix}{\argonze}{\LANGUE{tape}{ruban} $1$}{0,0}
% \draw (tete) node[coordinate] (tete1) {};
% \draw (tete1) +(8,0) node[coordinate] (tete1b) {};
% \draw (tete1) -| (tete1b) ;

\RUBANMACHINEDETURINGTETEtexte{#3}{#4}{\argdix}{\argonze}{\LANGUE{tape}{ruban} $2$}{0,-1.5}
% \draw (tete) node[coordinate] (tete2) {};
% \draw (tete2) +(-5.5,0) node[coordinate] (tete2b) {};
% \draw (tete2) -| (tete2b) ;

\RUBANMACHINEDETURINGTETEtexte{#5}{#6}{\argdix}{\argonze}{\LANGUE{tape}{ruban} $3$}{0,-3.2}
\draw (tete) node[coordinate] (tete3) {};

\RUBANMACHINEDETURINGTETEtexte{#7}{#8}{\argdix}{\argonze}{\LANGUE{tape}{ruban} $4$}{0,-4.9}
\draw (tete) node[coordinate] (tete4) {};

%% ETAT:
     \draw (tete4) node [name=etat,draw,circle,anchor=north]  {#9};
%      \draw (etat) -- (tete3);
%     \draw  (etat) -|  (tete2b);
% \draw  (etat) -|  (tete1b);
%\end{tikzpicture}
}


%%%%%%%%%%%%%%%%%%%%%%%%%%%%%%%%%%%%%%%%%%%%%%%%%%%%%%%%%%%%%%%%%%%%%%%%%%%%
%                 MACHINES DE TURING (automate)
%%%%%%%%%%%%%%%%%%%%%%%%%%%%%%%%%%%%%%%%%%%%%%%%%%%%%%%%%%%%%%%%%%%%%%%%%%%%



\newcommand\dec{1ex}
\newcommand\myACCOLADE[4]{
\ACCOLADE{#1*\lcase+#4*\lcase-0.5*\lcase,0.5*\lcase+\dec}{#1*\lcase+#4*\lcase-0.5*\lcase+#3*\lcase,0.5*\lcase+\dec}{#2}
}


\newcommand\M{\Sigma}


\newcommand\sommet[2]{\node[state] (#1) #2 {$#1$};}
\newcommand\sommetinitial[2]{\node[state,initial] (#1) #2 {$#1$};}
\newcommand\sommetfinal[2]{\node[state,accepting] (#1) #2 {$#1$};}
\newcommand\gtransition[5]{\draw (#1) edge node {#2/#4 $#5$} (#3);}
\newcommand\gtransitionba[5]{\draw (#1) [loop above] edge node {#2/#4
    $#5$} (#3);}
\newcommand\gtransitionbb[5]{\draw (#1) [loop below] edge node {#2/#4
    $#5$} (#3);}
\newcommand\gtransitionbleft[5]{\draw (#1) [loop left] edge node {#2/#4
    $#5$} (#3);}
\newcommand\gtransitionbright[5]{\draw (#1) [loop right] edge node {#2/#4
    $#5$} (#3);}
\newcommand\gtransitionbl[5]{\draw (#1) [bend left] edge node {#2/#4
    $#5$} (#3);}
\newcommand\gtransitionbr[5]{\draw (#1) [bend right] edge node {#2/#4
    $#5$} (#3);}
\newcommand\gtransitionbadouble[9]{\draw (#1) [loop above] edge node {
\begin{tabular}{l}
#2/#4 $#5$ \\
  #7/#8 $#9$ \\
\end{tabular}
} (#3);}

\newcommand\SCALE{}

\newcommand\dessinmtp[2]{
\begin{center}
\begin{tikzpicture}[->,>=stealth',shorten >=1pt,auto,node distance=2.3cm,semithick,#2] 
#1
\end{tikzpicture}
\end{center}
}

\newcommand\dessinmt[1]{\dessinmtp{#1}{scale=1}}


\newcommand\ANIMATIONMT[1]{
\begin{overprint}
\begin{dessinp}{scale=0.52}
\foreach \av/\q/\b  [count=\xi]
in  
{#1}
{ 
\only<\xi| handout 0>{ 
 \StrLen{\av}[\Resintm] % ATTENTION IL FAUT CETTE SYNTAXE POUR FAIRE
                         % EQUIVALENT D UN EDEF
 \MACHINEDETURING{\av \b}{\Resintm}{$\q$}{5}{20}{}
}}
 \end{dessinp}
\end{overprint}
}

\newcommand\ANIMATIONMTd[1]{
\begin{overprint}
\begin{dessinp}{scale=0.52}
\foreach \av/\q/\b  [count=\xi]
in  
{#1}
{ 
\only<\xi| handout 0>{ 
 \StrLen{\av}[\Resintm] % ATTENTION IL FAUT CETTE SYNTAXE POUR FAIRE
                         % EQUIVALENT D UN EDEF
 \hspace{-1cm}\MACHINEDETURING{\av \b}{\Resintm}{$\q$}{5}{33}{}
}}
 \end{dessinp}
\end{overprint}
}

\def\MAXDET{20}

\newcommand\DIAGRAMMEESPACETEMPS[1]{
\draw node[coordinate] (tt) {};
\foreach \av/\q/\b  in  
{#1}
{ 

\RUBANMACHINEDETURING{\av{$\q$}\b}{0}{4}{\MAXDET}{}{tt};
\draw (tt) +(0,-\hcase) node[coordinate] (tt) {};
}
%\foreach \x in {0,...,19} {\draw (tt) + (\x*\lcase,0) node {$\vdots$};}
}



\newcommand\progun{
/q_0/0011, % initial
X/q_1/011,X0/q_1/11,X/q_2/0Y1,/q_2/X0Y1,
X/q_0/0Y1,XX/q_1/Y1,XXY/q_1/1, XX/q_2/YY, X/q_2/XYY,
XX/q_0/YY,XXY/q_3/Y,XXYY/q_3/B,XXYYB/q_4/B
}

\newcommand\DIAGRAMMEESPACETEMPSun{
\DIAGRAMMEESPACETEMPS
%% RECOPIE progun
{/q_0/0011, % initial
X/q_1/011,X0/q_1/11,X/q_2/0Y1,/q_2/X0Y1,
X/q_0/0Y1,XX/q_1/Y1,XXY/q_1/1, XX/q_2/YY, X/q_2/XYY,
XX/q_0/YY,XXY/q_3/Y,XXYY/q_3/B,XXYYB/q_4/B}
%% FIN RECOPIE
}

\newcommand\progdeux{
/q_0/0010, X/q_1/010, X0/q_1/10, X/q_2/0Y0,/q_2/X0Y0,
X/q_0/0Y0,XX/q_1/Y0,XXY/q_1/0,XXY0/q_1/%B,
}


\newcommand\DIAGRAMMEESPACETEMPSdeux{
\DIAGRAMMEESPACETEMPS
{
%% RECOPIE progdeux
/q_0/0010, X/q_1/010, X0/q_1/10, X/q_2/0Y0,/q_2/X0Y0,
X/q_0/0Y0,XX/q_1/Y0,XXY/q_1/0,XXY0/q_1/%B,
%% FIN RECOPIE
}
}

\newcommand\DIAGRAMMEESPACETEMPSdeuxvariante{
\def\memolcase{\lcase}
\def\memoMAXDET{20}
\def\MAXDET{15}
\def\lcase{6 ex}
\DIAGRAMMEESPACETEMPS
{
%% RECOPIE progdeux
/q_00/010, X/q_10/10, X0/q_11/0, X/q_20/Y0,/q_2X/0Y0,
X/q_00/Y0,XX/q_1Y/0,XXY/q_10/,XXY0/q_1B/%B,
%% FIN RECOPIE
}
\def\lcase{\memolcase}
\def\MAXDET{\memoMAXDET}
}


\newcommand\ANIMATIONun{
\ANIMATIONMT
%% RECOPIE progun
{/q_0/0011, % initial
X/q_1/011,X0/q_1/11,X/q_2/0Y1,/q_2/X0Y1,
X/q_0/0Y1,XX/q_1/Y1,XXY/q_1/1, XX/q_2/YY, X/q_2/XYY,
XX/q_0/YY,XXY/q_3/Y,XXYY/q_3/B,XXYYB/q_4/B}
%% FIN RECOPIE
}



\newcommand\DESSINALGO[3]{
% argument 1= ou
% argument 2= texte 
% argument 3= nom du noeud ex m
\draw node[draw,#1,minimum size=1cm] (#3) {#2};

% pour fitting box
\draw (#3.west) +(-0.2cm,0) node (#3c2) {};
\draw (#3.north) +(0,+0.2cm) node (#3c3) {};
\draw (#3.south) +(0,-0.2cm) node (#3c4) {};
}

\newcommand\DESSINALGOa[3]{
\DESSINALGO{#1}{#2}{#3}
\draw (#3.east) +(-0.1cm,0.2cm) node(#3a) {};
\draw[->] (#3a) -- ++(0.5cm,0) node (#3aa) {};
\draw (#3aa) node[anchor=west] (#3aaa) {\LANGUE{Accept}{Accepte}};
\draw (#3aaa.east) node(#3aaaa) {};
}

\newcommand\DESSINALGOe[4]{
\DESSINALGO{#1}{#2}{#3}
\draw (#3.east) +(-0.1cm,0.2cm) node(#3a) {};
\draw[->] (#3a) -- ++(0.5cm,0) node (#3aa) {};
\draw (#3aa) node[anchor=west] (#3aaa) {#4};
\draw (#3aaa.east) node(#3aaaa) {};
}


\newcommand\DESSINALGOar[3]{
\DESSINALGOa{#1}{#2}{#3}
\draw (#3.east) +(-0.1cm,-0.2cm) node(#3r) {};
\draw[->] (#3r) -- ++(0.5cm,0) node (#3rr) {};
\draw (#3rr) node[anchor=west] (#3rrr) {\LANGUE{Reject}{Refuse}};
\draw (#3rrr.east) node(#3rrrr) {};
}


\newcommand\WINDOWarg[7] {
%#1: position
%#2..7: contenu du tableau
\draw (#1) node[case]  {#2} \dcase;
\draw (#1) ++ (\lcase,0) node[case] {#3} \dcase;
\draw (#1) ++ (2*\lcase,0) node[case] {#4} \dcase;
\draw (#1) ++ (0,-\lcase) node[case] {#5} \dcase;
\draw (#1) ++ (\lcase,-\lcase) node[case] {#6} \dcase;
\draw (#1) ++ (2*\lcase,-\lcase) node[case] {#7} \dcase;
}

\newcommand\WINDOW[3]{
%#1: positioin
%#2: contenu
%#3: étiquette
\WINDOWarg{#1}
{\StrChar{#2}{1}}{\StrChar{#2}{2}}{\StrChar{#2}{3}}
{\StrChar{#2}{4}}{\StrChar{#2}{5}}{\StrChar{#2}{6}}
\draw (#1) +(-\lcase*1.2,-\lcase*0.5) node {#3};
}





%%%%%%%%%%%%%%%%%%%%%%%%%%%%%%%%%%%%%%%%%%%%%%%%%%%%%%%%%%%%%%%%%%%%%%%%%%%%
%                 LISTINGS
%%%%%%%%%%%%%%%%%%%%%%%%%%%%%%%%%%%%%%%%%%%%%%%%%%%%%%%%%%%%%%%%%%%%%%%%%%%%


% LISTINGS
\RequirePackage{listings}
 \lstset{language=Java,
 mathescape=true,
 showspaces=false,
 showstringspaces=false,
 frame=single,
morekeywords={recursive,then,div,function,integer,read,out,else,par,endpar,seq,endseq,and,not,read,write,repeat,until,undef,let,to,endfor,endif,endwhile,guess,in,parallel},
% basicstyle=\ttfamily\small,
% basewidth={1.1ex},
% keywordstyle=\bf\color{blue},
% %stringstyle=\bf\ttfamily\color{green},
% commentstyle=\bfseries\color{magenta},
 literate={:=}{{$\gets$ }}1 {<=}{{$\leq$}}2  {>=}{{$\geq$}}2
 {!=}{{$\neq$}}1 
 }
% \lstset{escapechar=\%}
\let\lst\lstinline
\def\beginlisting{\begin{lstlisting}}
\def\endlisting{\end{lstlisting}}  




\newenvironment{algo}[1]{
  \lstset{escapechar=\@}
  \def\input{$#1$}
  \def\out{out}
  \def\BEGIN{\textbf{read} \input }
 \def\BEGINN{\textbf{repeat} }
  \def\END{\textbf{until out!=undef}}
  \def\ENDD{\textbf{write \out}}
  \def\ENDT{\textbf{until $\ctl$ = $q_a$ or $\ctl$ = $q_r$}}
  \def\ENDR{\textbf{until $\ctl$ = $q_a$}}
  \def\ENDP{\textbf{until $\ctl$ = $q_a$}}
}{}
\newenvironment{algon}[1]{
  \lstset{numbers=left}
  \begin{algo}{#1}
}{\end{algo}}
\newcommand\FSM[3]{FSM(#1,\lst{#2},#3)}
\newcommand\FSMi[3]{FSM(#1,{#2},#3)}




%%%%%%%%%%%%%%%%%%%%%%%%%%%%%%%%%%%%%%%%%%%%%%%%%
%
%
%                      TRADUCTION EN COURS
%
%%%%%%%%%%%%%%%%%%%%%%%%%%%%%%%%%%%%%%%%%%%%%%%%%

\newcommand\ENTREPOT[1]{\NDLR{ENTREPOT ICI: #1}}
\newcommand\scite[2][]{\cite[#1]{#2}}
\newcommand\nospellcheck[1]{{#1}}
\newcommand\nd{nd}
\newcommand\NL{
 
}

\newenvironment{FRANCAIS}{
\small
\color{orange}
}
{\normalsize
\color{black}
}


\newenvironment{SEMIFRANCAIS}{
\small
\color{violet}
}
{\normalsize
\color{black}
}

\newcommand\QFRANCAIS[1]{
\begin{FRANCAIS}
#1
\end{FRANCAIS}
}


%%%%%%%%%%%%%%%%%%%%%%%%%%%%%%%%%%%%%%%%%%%%%%%%%
%
%
%                      Wikipedia
%
%%%%%%%%%%%%%%%%%%%%%%%%%%%%%%%%%%%%%%%%%%%%%%%%%

\providecommand\tightlist{}
\providecommand\Complex{\C}
\newenvironment{myalgie}{\begin{array}{lll}}{\end{array}}

\newcommand\nl{\ \\ }




%%%%%%%%%%%%%%%%%%%%%%%%%%%%%%%%%%%%%%%%%%%%%%%%%
%
%
%                      Pour citations précises
%
%%%%%%%%%%%%%%%%%%%%%%%%%%%%%%%%%%%%%%%%%%%%%%%%%



\newcommand\citeprecis[2]{{\cite[#1]{#2}}}
\newcommand\citeprecisinv[2]{\citeprecis{#2}{#1}}

\usepackage{csquotes}



%%%%%%%%%%%%%%%%%%%%%%%%%%%%%%%%%%%%%%%%%%%%%%%%%
%
%
%                      Travail sur couleurs
%
%   toutes les réponses pointent dans le même sens, et le schéma correspond exactement 
%. au profil perceptif d’un protanope** (déficit des cônes L / longueurs d’onde longues)
%  MAIS SEVERE
%
%%%%%%%%%%%%%%%%%%%%%%%%%%%%%%%%%%%%%%%%%%%%%%%%%


% ================================
% GROUPE 1 — Verts sombres
% ================================
\definecolor{deepgreen}{RGB}{0,160,40}
\definecolor{forestgreen}{RGB}{0,150,70}
\definecolor{mossgreen}{RGB}{0,140,90}
\definecolor{greenwave}{RGB}{0,150,90}
\definecolor{forestmist}{RGB}{0,160,90}
\definecolor{darkteal}{RGB}{0,130,70}

% ================================
% GROUPE 2 — Teal / verts bleutés
% ================================
\definecolor{seagreen}{RGB}{0,130,100}
\definecolor{freshteal}{RGB}{0,140,130}
\definecolor{greenteal}{RGB}{0,140,120}
\definecolor{blueteal}{RGB}{0,150,120}      % ancien "teal" (doublon corrigé)
\definecolor{lagoon}{RGB}{0,180,120}
\definecolor{greenmist}{RGB}{0,170,110}

% ================================
% GROUPE 3 — Turquoises / Aqua
% ================================
\definecolor{springgreen}{RGB}{0,190,100}
\definecolor{mintgreen}{RGB}{10,200,150}
\definecolor{oceanfoam}{RGB}{0,200,150}
\definecolor{aquagreen}{RGB}{30,210,170}
\definecolor{coastalaqua}{RGB}{40,210,180}
\definecolor{softaqua}{RGB}{40,210,190}

% ================================
% GROUPE 4 — Cyan / Aqua clairs
% ================================
\definecolor{emeraldteal}{RGB}{0,180,140}
\definecolor{deepcyan}{RGB}{0,180,160}
\definecolor{lightcyan}{RGB}{0,200,160}
\definecolor{aquamarine}{RGB}{0,200,160}
\definecolor{lightaqua}{RGB}{60,230,200}
\definecolor{mistblue}{RGB}{80,240,210}

% ================================
% GROUPE 5 — Bleus clairs
% ================================
\definecolor{paleaqua}{RGB}{80,200,190}
\definecolor{glacier}{RGB}{90,150,200}
\definecolor{skyblueA}{RGB}{60,150,220}      % ancien "skyblue"
\definecolor{lightsky}{RGB}{90,170,230}
\definecolor{iceblue}{RGB}{120,190,240}
\definecolor{frost}{RGB}{150,210,250}

% ================================
% GROUPE 6 — Bleus moyens
% ================================
\definecolor{steelblue}{RGB}{30,100,160}
\definecolor{cadetblue}{RGB}{70,120,180}
\definecolor{coolblue}{RGB}{110,170,220}
\definecolor{morningblue}{RGB}{140,200,240}
\definecolor{blueiceA}{RGB}{40,150,210}       % ancien "blueice" (doublon corrigé)
\definecolor{blueair}{RGB}{60,170,225}

% ================================
% GROUPE 7 — Azurs / Bleu saturé
% ================================
\definecolor{tealblue}{RGB}{0,110,130}
\definecolor{deepbluecyan}{RGB}{0,110,170}
\definecolor{azur}{RGB}{0,70,180}
\definecolor{brightazure}{RGB}{20,130,200}
\definecolor{lightazure}{RGB}{80,180,235}
\definecolor{skylight}{RGB}{100,200,245}

% ================================
% GROUPE 8 — Bleus sombres / violets
% ================================
\definecolor{nightblue}{RGB}{40,0,110}
\definecolor{slateblue}{RGB}{20,40,120}
\definecolor{violetA}{RGB}{60,50,160}         % ancien "violet" (distinct de violetB ci-dessous)
\definecolor{blueviolet}{RGB}{80,60,180}
\definecolor{purpleA}{RGB}{100,70,200}        % ancien "purple"
\definecolor{lightpurple}{RGB}{130,90,210}

% ================================
% GROUPE 9 — Violets saturés / indigos
% ================================
\definecolor{lilac}{RGB}{160,110,220}
\definecolor{indigoA}{RGB}{50,0,130}          % ancien "indigo"
\definecolor{deepindigo}{RGB}{70,10,150}
\definecolor{royalpurple}{RGB}{90,20,170}

% ================================
% GROUPE 10 — Fermeture cercle perceptif
% ================================
\definecolor{coolmint}{RGB}{0,130,70}
\definecolor{deepsea}{RGB}{0,130,70}

% ================================
% PALETTE DEUTÉRANOPE — noms uniques
% ================================
\definecolor{crimson}{RGB}{200,30,60}
\definecolor{brickred}{RGB}{180,40,40}
\definecolor{coral}{RGB}{220,90,70}
\definecolor{salmon}{RGB}{240,130,120}

\definecolor{magenta}{RGB}{200,40,150}
\definecolor{fuchsia}{RGB}{220,60,170}
\definecolor{violetB}{RGB}{140,70,190}        % ancien "violet"
\definecolor{indigoB}{RGB}{90,50,180}         % ancien "indigo"

\definecolor{royalblueA}{RGB}{50,80,200}      % ancien "royalblue"
\definecolor{azure}{RGB}{30,120,230}
\definecolor{skyblueB}{RGB}{80,160,230}       % second "skyblue"
\definecolor{tealB}{RGB}{0,150,170}           % second "teal"

\definecolor{marine}{RGB}{0,80,140}




%% Alternative

% === Générateur automatique fill/text pour protanopes ===
% Usage : \ProtaColorPair{mistblue}
% Suppose que mistblue est déjà défini par \definecolor{mistblue}{RGB}{...}

\newcommand{\ProtaColorPair}[1]{%
  % fill = couleur originale
  \colorlet{#1fill}{#1}%
  % text = version assombrie idéalement lisible
  \colorlet{#1text}{#1!80!black}%
}

%Utilisation
%\ProtaColorPair{mistblue}

%\textcolor{mistbluefill}{Couleur claire} \quad
%\textcolor{mistbluetext}{Couleur pour texte}

	\foreach \i/\name in {
	  1/deepgreen,
	  2/greenwave,
	  3/lagoon,
	  4/oceanfoam,
	  5/softaqua,
	  6/mistblue,
 	 7/lightsky,
	  8/blueiceA,   % <-- corrigé
	  9/deepcyan,
	  10/azur,
	  11/blueviolet,
	  12/royalpurple}
	  {\ProtaColorPair{\name}}
	  


\newcommand\ROUEPROTANOPEDOUZE{
	% Roue chromatique protanope — 12 couleurs
	\begin{tikzpicture}[scale=3]

	\foreach \i/\name/\nametext in {
	  1/deepgreen/deepgreen,
	  2/greenwave/greenwave,
	  3/lagoon/lagoontext,
	  4/oceanfoam/oceanfoam,
	  5/softaqua/softaquatext,
	  6/mistblue/mistbluetext,
 	 7/lightsky/lightskytext,
	  8/blueiceA/blueiceA,   % <-- corrigé
	  9/deepcyan/blueiceA,
	  10/azur/azur,
	  11/blueviolet/blueviolet,
	  12/royalpurple/royalpurple}
	{
	  % secteur coloré
 	 \filldraw[fill=\name, draw=black, line width=0.2pt]
	    ({90-30*\i}:1)
	    -- ({90-30*(\i-1)}:1)
	    -- (0,0) -- cycle;

 	 % label dans la bonne couleur
	  \node at ({90-30*(\i-0.5)}:1.25)
	    {\textcolor{\nametext}{\small \nametext}};
	}

	\end{tikzpicture}
}


\newcommand\ROUEPROTANOPEVINGT{
\begin{tikzpicture}[scale=3]

\foreach \i/\name/\nametext in {
  1/deepgreen/deepgreen,
  2/greenwave/greenwave,
  3/lagoon/lagoon,
  4/oceanfoam/oceanfoam,
  5/softaqua/softaquatext,
  6/mistblue/mistbluetext,
  7/lightsky/lightsky,
  8/blueiceA/blueiceA,      % corrigé
  9/deepcyan/deepcyan,
  10/azur/azur,
  11/brightazure/brightazure,
  12/skyblueA/skyblueA,     % corrigé
  13/coolblue/coolblue,
  14/glacier/glacier,
  15/marine/marine,
  16/violetA/violetA,      % corrigé
  17/blueviolet/blueviolet,
  18/purpleA/purpleA,      % corrigé
  19/royalpurple/royalpurple,
  20/indigoA/indigoA}      % corrigé
{
  % Secteur coloré
  \filldraw[fill=\name, draw=black, line width=0.2pt]
    ({90-18*\i}:1)
    -- ({90-18*(\i-1)}:1)
    -- (0,0) -- cycle;

  % Étiquette
  \node at ({90-18*(\i-0.5)}:1.25)
    {\textcolor{\nametext}{\small \nametext}};
}

\end{tikzpicture}
}

\newcommand\ROUEDEUTERANOPEDOUZE{
% Roue chromatique deutéranope — 12 couleurs
\begin{tikzpicture}[scale=3]

\foreach \i/\name in {
  1/crimson,
  2/brickred,
  3/coral,
  4/salmon,
  5/magenta,
  6/fuchsia,
  7/violetB,       % corrigé
  8/indigoB,       % corrigé
  9/royalblueA,    % corrigé
  10/azure,
  11/skyblueB,     % corrigé
  12/tealB}        % corrigé
{
  % Secteur coloré
  \filldraw[fill=\name, draw=black, line width=0.2pt]
    ({90-30*\i}:1)
    -- ({90-30*(\i-1)}:1)
    -- (0,0) -- cycle;

  % Label dans la couleur
  \node at ({90-30*(\i-0.5)}:1.25)
    {\textcolor{\name}{\small \name}};
}

\end{tikzpicture}
}



% === Macro intelligente pour foncer une couleur de façon lisible pour protanope ===
% Usage : \protadarken{mistblue}{40}  -> crée mistblue@dark
% Le 40 signifie : 40% de mélange avec du noir

\newcommand{\DarkenForText}[2]{%
  % #1 = nom de la couleur
  % #2 = pourcentage (40 conseillé si tu ne sais pas)
  %
  % On crée une nouvelle couleur #1@dark
  \colorlet{#1text}{black!#2!#1}%
}





%%%%

% Groupe 1
\DarkenForText{deepgreen}{20}
\DarkenForText{forestgreen}{20}
\DarkenForText{mossgreen}{20}
\DarkenForText{greenwave}{20}
\DarkenForText{forestmist}{20}
\DarkenForText{darkteal}{20}

% Groupe 2
\DarkenForText{seagreen}{20}
\DarkenForText{freshteal}{20}
\DarkenForText{greenteal}{20}
\DarkenForText{blueteal}{20}
\DarkenForText{lagoon}{20}
\DarkenForText{greenmist}{20}

% Groupe 3
\DarkenForText{springgreen}{20}
\DarkenForText{mintgreen}{20}
\DarkenForText{oceanfoam}{20}
\DarkenForText{aquagreen}{20}
\DarkenForText{coastalaqua}{20}
\DarkenForText{softaqua}{20}

% Groupe 4
\DarkenForText{emeraldteal}{20}
\DarkenForText{deepcyan}{20}
\DarkenForText{lightcyan}{20}
\DarkenForText{aquamarine}{20}
\DarkenForText{lightaqua}{20}
\DarkenForText{mistblue}{30}

% Groupe 5
\DarkenForText{paleaqua}{20}
\DarkenForText{glacier}{20}
\DarkenForText{skyblueA}{20}
\DarkenForText{lightsky}{20}
\DarkenForText{iceblue}{20}
\DarkenForText{frost}{20}

% Groupe 6
\DarkenForText{steelblue}{20}
\DarkenForText{cadetblue}{20}
\DarkenForText{coolblue}{20}
\DarkenForText{morningblue}{20}
\DarkenForText{blueiceA}{20}
\DarkenForText{blueair}{20}

% Groupe 7
\DarkenForText{tealblue}{20}
\DarkenForText{deepbluecyan}{20}
\DarkenForText{azur}{20}
\DarkenForText{brightazure}{20}
\DarkenForText{lightazure}{20}
\DarkenForText{skylight}{20}

% Groupe 8
\DarkenForText{nightblue}{20}
\DarkenForText{slateblue}{20}
\DarkenForText{violetA}{20}
\DarkenForText{blueviolet}{20}
\DarkenForText{purpleA}{20}
\DarkenForText{lightpurple}{20}

% Groupe 9
\DarkenForText{lilac}{20}
\DarkenForText{indigoA}{20}
\DarkenForText{deepindigo}{20}
\DarkenForText{royalpurple}{20}

% Groupe 10
\DarkenForText{coolmint}{20}
\DarkenForText{deepsea}{20}

% Palette deutéranope
\DarkenForText{crimson}{20}
\DarkenForText{brickred}{20}
\DarkenForText{coral}{20}
\DarkenForText{salmon}{20}

\DarkenForText{magenta}{20}
\DarkenForText{fuchsia}{20}
\DarkenForText{violetB}{20}
\DarkenForText{indigoB}{20}

\DarkenForText{royalblueA}{20}
\DarkenForText{azure}{20}
\DarkenForText{skyblueB}{20}
\DarkenForText{tealB}{20}

% Couleur supplémentaire
\DarkenForText{marine}{20}

%%%%%%%%%%%%%%%%%%%%%
%
%.   apx proof
%
%%%%%%%%%%%%%%%%%%%%%



% do
\newcommand\PREUVESENAPPENDIXCOMPATIBLE{
	\newenvironment{appendixproof}{{\bf \LANGUE{Proof}{Démonstration}}:}{\hfill $\Box$}
	\newtheorem{theoremrep}[theorem]{Theorem}
	\newtheorem{lemmarep}[theorem]{Lemma}
	\newtheorem{propositionrep}[theorem]{Proposition}
	\newtheorem{corollaryrep}[theorem]{Corollary}
	\newtheorem{summaryrep}[theorem]{Summary}	
}

\providecommand\apxparameter{}

%% proof inline
%\providecommand\apxparameter{appendix=inline}
%% proof in appendix
%\providecommand\apxparameter{}


\apxproofmode{
\usepackage[\apxparameter]{apxproof}
\theoremstyle{plain}

\newtheoremrep{theorem}{Theorem}
%\newtheoremrep{theoremconjecture}[theorem]{Theorem-To-Be-Proved=Conjecture}
%\newtheoremrep{lemmaconjecture}[theorem]{Conjecture-To-Be-Proved=Conjecture}
\newtheoremrep{claim}[theorem]{Theorem}
\newtheoremrep{proposition}[theorem]{Proposition}
\newtheoremrep{lemma}[theorem]{Lemma}
\newtheoremrep{corollary}[theorem]{Corollary}
%%bug mars 26: \newtheoremrep{summary}[theorem]{Summary}
% Du coup:
\newtheoremrep{resume}[theorem]{Summary}
%bug mars 26: \newtheoremrep{openproblem}[theorem]{Open Problem}
%\newtheoremrep{claim}[theorem]{Claim}
}


%\apxproofmodesubsection{
%%apx proof, pour que compteur correct
%\makeatletter
%\AtBeginDocument{%
%\let\axp@section\subsection
%  \let\axp@oldsection\subsection
%}
%\makeatother
%\renewcommand\appendixprelim{\section{Missing proofs}}
%}


\makeatletter
\apxproofmodesubsection{
\let\axp@section\subsection
%\let\axp@oldsection\subsection
\renewcommand\appendixprelim{\section{Proofs from main documents}}
}
\makeatother


\makeatletter
\newcommand\appendixapxproofsubsection{
\appendix
\let\apx@orig@appendix\appendix
\renewcommand{\appendix}{}
}
\makeatother


%%% Pour checker/vérifier que bon


\newcommand\docheck[1]{\textcolor{check}{#1}}

\usepackage{xcolor}
%\definecolor{chose}{RGB}{0,110,50}
%\definecolor{forestgreen}{RGB}{0,150,70}
%\definecolor{ctruc}{RGB}{164,164,164}
%\definecolor{oceanfoam}{RGB}{0,200,150}
\definecolor{check}{RGB}{160,160,0}




. Retirées automatiquement en mode lipics.
\providecommand\NODECORATION[1]{\ifnotSLIDE{#1}}
\lipicsmode{\renewcommand\NODECORATION[1]{}}

\providecommand\PASSIAPXPROOF[1]{#1}
\apxproofmode{\renewcommand\PASSIAPXPROOF[1]{}}



%% end gestion des switchs

\def\MACROSPOINTEXOUVERT{}


\providecommand\ifnotTDEXAM[1]{#1}
\providecommand\ifnotSLIDE[1]{#1}


  
% pour preuves en appendix
\providecommand\PREUVESENAPPENDIX[1]{}



%%%% STYLES: APRES VOLS


\NODECORATION{
\usepackage{fourier}
}

%% Repli si fourier (et donc fourier-orns, qui définit \danger) n'est pas
%% chargé ; doit rester APRÈS le bloc ci-dessus, sinon fourier-orns échoue
%% avec « \danger already defined ».
\providecommand\danger{}

%%%% Box around environment
\usepackage[framemethod=tikz]{mdframed}

\newcommand\ENVCOLOR[1]{#1}
\newcommand\BEGINTEMPENVCOLOR[1]{\let\envcolorwas\ENVCOLOR\renewcommand\ENVCOLOR[1]{#1}}
\newcommand\ENDTEMPENVCOLOR{\let\ENVCOLOR\envcolorwas}


\NODECORATION{
\surroundwithmdframed[hidealllines=true,backgroundcolor=\ENVCOLOR{blue!10},roundcorner=8pt]{exercice}
\surroundwithmdframed[hidealllines=true,backgroundcolor=\ENVCOLOR{blue!10},roundcorner=8pt]{exercise}
\surroundwithmdframed[hidealllines=true,backgroundcolor=\ENVCOLOR{blue!10},roundcorner=8pt]{exercicee}
\surroundwithmdframed[hidealllines=true,backgroundcolor=\ENVCOLOR{blue!10},roundcorner=8pt]{exerciced}
\surroundwithmdframed[hidealllines=true,backgroundcolor=\ENVCOLOR{blue!10},roundcorner=8pt]{exercicec}
\surroundwithmdframed[hidealllines=true,backgroundcolor=\ENVCOLOR{blue!10},roundcorner=8pt]{exercicedc}
\surroundwithmdframed[hidealllines=true,backgroundcolor=\ENVCOLOR{orange!20},roundcorner=8pt]{example}
\surroundwithmdframed[hidealllines=true,backgroundcolor=\ENVCOLOR{orange!20},roundcorner=8pt]{exemple}
\surroundwithmdframed[hidealllines=true,backgroundcolor=\ENVCOLOR{orange!10},roundcorner=8pt]{remark}
\ifnotSLIDE{\surroundwithmdframed[backgroundcolor=\ENVCOLOR{green!15},roundcorner=8pt]{definition}}
\surroundwithmdframed[hidealllines=true,backgroundcolor=\ENVCOLOR{gray!35}]{proposition}
\ifnotSLIDE{\surroundwithmdframed[backgroundcolor=\ENVCOLOR{gray!35},roundcorner=8pt]{theorem}}
\surroundwithmdframed[hidealllines=true,backgroundcolor=\ENVCOLOR{gray!35}]{lemma}
\surroundwithmdframed[hidealllines=true,backgroundcolor=\ENVCOLOR{gray!35}]{corollary}


\surroundwithmdframed[backgroundcolor=red!35,roundcorner=8pt]{theoremaprouver}
\surroundwithmdframed[backgroundcolor=red!35,roundcorner=8pt]{conjectureaprouver}
}

%%% tcolorbox

\usepackage{tcolorbox}
\tcbuselibrary{skins}


  %%. Commande magique
%% \tcolorboxenvironment{⟨name⟩}{⟨options⟩}: attach options de tcolorbox à environnement
%% alternative: je cree des boites moi meme: voici des exemples


\newtcolorbox{maboiteconfiance}[2][]{colback=red!5!white, colframe=red!75!black,fonttitle=\bfseries, colbacktitle=red!85!black,enhanced,
attach boxed title to top center={yshift=-2mm},
  title={#2},#1}
  
\newtcolorbox{maboiteresume}[2][]{colback=red!5!white, colframe=red!75!black,fonttitle=\bfseries, colbacktitle=red!85!black,enhanced,
attach boxed title to top center={yshift=-2mm},
  title={#2},#1}
  \newtcolorbox{maboiteresumeperso}[2][]{colback=red!5!white, colframe=red!75!black,fonttitle=\bfseries, colbacktitle=red!85!black,enhanced,
attach boxed title to top center={yshift=-2mm},
  title={#2},#1}
  \newtcolorbox{maboitecommentaire}[2][]{colback=red!5!white, colframe=red!75!black,fonttitle=\bfseries, colbacktitle=red!85!black,enhanced,
attach boxed title to top center={yshift=-2mm},
  title={#2},#1}
\newtcolorbox{maboitetruc}[2][]{colback=red!5!white, colframe=red!75!black,fonttitle=\bfseries, colbacktitle=red!85!black,enhanced,
attach boxed title to top left={yshift=-2mm},
  title={#2},#1}

%% Boîte des mises en évidence CARE. Couleurs pensées pour la
%% protanopie : l'axe bleu-jaune reste discriminant (jaune = \IMPORTANT
%% d'Olivier, bleu = CARE), et la distinction ne repose pas que sur la
%% couleur (cadre + bandeau étiqueté, lisibles même en niveaux de gris).
\newtcolorbox{maboiteimportantcare}[2][]{colback=cyan!10!white, colframe=blue!60!black,fonttitle=\bfseries, colbacktitle=blue!70!black,enhanced,
attach boxed title to top center={yshift=-2mm},
  title={#2},#1}
  


%%%% FIN STYLES



%%%%%%%%%%%%%%%%%%%%%%%%%%%%%%%%%%%%%%%%%%%%%%%%%
%%%%%%%%%%%%%%%%%%%%%%%%%%%%%%%%%%%%%%%%%%%%%%%%%
%
%
%                      Version ENSEIGNEMENT 2020
%
%
%%%%%%%%%%%%%%%%%%%%%%%%%%%%%%%%%%%%%%%%%%%%%%%%%
%%%%%%%%%%%%%%%%%%%%%%%%%%%%%%%%%%%%%%%%%%%%%%%%%

%english, french
\ifdefined\CESTDUFRANCAIS
	\providecommand\LANGUE[2]{#2}
	\usepackage[french]{babel}
\else
	\providecommand\LANGUE[2]{#1}
\fi
\providecommand\LANGUEinv[2]{\LANGUE{#2}{#1}}

%%%%%%%%%%%%%%%%%%%%%%%%%%%%%%%%%%%%%%%%%%%%%%%%%
%
%
%                      Packages
%
%
%%%%%%%%%%%%%%%%%%%%%%%%%%%%%%%%%%%%%%%%%%%%%%%%%


%%% SURLIGNEUR
\providecommand\ifnotmodeDIFFUSIONFINALE[1]{
%% 2021: Comprend pas mais assure
\ifdefined\SAFEMODE
\else
#1
\fi
%#1
}
\providecommand\ifmodeDIFFUSIONFINALE[1]{
\ifdefined\SAFEMODE
#1
\else
\fi
}
\newcommand\SURLIGNE[1]{
\ifnotmodeDIFFUSIONFINALE{
\BEGINTEMPENVCOLOR{blue!30}
{                        \colorbox{blue!30}{

\noindent \begin{minipage}{\dimexpr\textwidth-2\fboxsep\relax}
#1
\end{minipage}
}
}
\ENDTEMPENVCOLOR}
\ifmodeDIFFUSIONFINALE{
\noindent \begin{minipage}{\textwidth}
#1
\end{minipage}
}
}


\newenvironment{confiance}{\begin{maboiteresume}[colback=brown!50]{Résumé ici}}{\end{maboiteresume}}
\newcommand\CONFIANCE[2][Confiance ici]{
\begin{maboiteconfiance}[colback=yellow!50]{#1}
\textbf{ATTENTION : #2}
\end{maboiteconfiance}
}

\newenvironment{resume}{\begin{maboiteresume}[colback=brown!50]{Résumé ici}}{\end{maboiteresume}}
\newcommand\RESUME[2][Résumé ici]{
\begin{maboiteresume}[colback=brown!50]{#1}
#2
\end{maboiteresume}
}

\newenvironment{resumeperso}{\begin{maboiteresumeperso}[colback=brown!50]{Résumé perso ici}}{\end{maboiteresumeperso}}
\newcommand\RESUMEPERSO[2][Résumé perso  ici]{
\begin{maboiteresumeperso}[colback=brown!50]{#1}
#2
\end{maboiteresumeperso}
%
%
%\todo[inline, color=brown]{résumé: #1}
%\colorbox{brown}{
%\noindent \begin{minipage}{\textwidth}
%{
%\bf 
%#2
%}
%\end{minipage}
%}
}

\newcommand\DANSLAMARGE[1]{
\ifnotmodeDIFFUSIONFINALE{
\marginpar{\textcolor{brown}{#1}}}
}

\newenvironment{commentaireperso}{\begin{maboitecommentaire}[colback=brown!30]{Commentaire perso} Commentaire perso:}{\end{maboitecommentaire}}
\newcommand\COMMENTAIREPERSO[2][Commentaire perso]{
\begin{maboitecommentaire}[colback=brown!30]{#1}
#2
\end{maboitecommentaire}
}


\newcommand\SURSURLIGNE[2][truc surligné ici]{
%\todo[inline, color=orange]{surligné: #1}
\ifnotmodeDIFFUSIONFINALE{
\BEGINTEMPENVCOLOR{blue!60}
	                 \colorbox{blue!60}{
\noindent \begin{minipage}{\dimexpr\textwidth-2\fboxsep\relax}
{
\bf 
#2
}
\end{minipage}
}
\ENDTEMPENVCOLOR
}
\ifmodeDIFFUSIONFINALE{#2}
}

\newcommand\IMPORTANT[1]{
\ifnotmodeDIFFUSIONFINALE{
\BEGINTEMPENVCOLOR{yellow!100}
		        \colorbox{yellow!100}{

\noindent \begin{minipage}{\dimexpr\textwidth-2\fboxsep\relax}
{ 
#1
}
\end{minipage}
}
\ENDTEMPENVCOLOR
}
\ifmodeDIFFUSIONFINALE{#1}
}

%% Pendant de \IMPORTANT, réservé aux mises en évidence CARE ;
%% \IMPORTANT reste réservé aux annotations d'Olivier. Comme \IMPORTANT,
%% redevient du texte simple en mode DIFFUSIONFINALE.
\newcommand\IMPORTANTCARE[2][Important (CARE)]{
\ifnotmodeDIFFUSIONFINALE{
\begin{maboiteimportantcare}{#1}
#2
\end{maboiteimportantcare}
}
\ifmodeDIFFUSIONFINALE{#2}
}


\newcommand\QUESTION[1]{
\ifnotmodeDIFFUSIONFINALE{
\BEGINTEMPENVCOLOR{oceanfoam!100}
		        \colorbox{oceanfoam!100}{

\noindent \begin{minipage}{\dimexpr\textwidth-2\fboxsep\relax}
{ 
#1
}
\end{minipage}
}
\ENDTEMPENVCOLOR
}
\ifmodeDIFFUSIONFINALE{#1}
}




\newcommand\SOUSLIGNE[1]{
\colorbox{lightgray}{

\noindent \begin{minipage}{\dimexpr\textwidth-2\fboxsep\relax}
{ \color{gray}
#1
}
\end{minipage}
}
}



%% Gestion des cross-references cross-referencing
\usepackage{xr-hyper}
\usepackage{hyperref}

%%
\usepackage{versions}
\usepackage{amsmath,amssymb,mathrsfs,amsfonts}
%\usepackage{mathabx}
\usepackage{listings}
\usepackage{url}
\usepackage{versions}
\usepackage{color}
\usepackage[colorinlistoftodos]{todonotes}
%\usepackage{todonotes}
\usepackage{proof}
\usepackage{xstring}

%% Index
\usepackage{makeidx}
\makeindex
\providecommand\DOINDEXPRINT{}
\providecommand\SIGNALEMOTNOUV[1]{}
\ifdefined\INDEXPRINT
	\ifnotSAFEMODE{
%	\usepackage{showidx}
%	\let\oldindex\index
%	\renewcommand{\index}[1]{\oldindex{#1}\marginpar{LA: \tiny#1}}
	\renewcommand\DOINDEXPRINT{
	      	\let\oldindex\index
		\renewcommand{\index}[1]{\oldindex{##1}\marginpar{IDXLA: \tiny##1}}
		\renewcommand\SIGNALEMOTNOUV[1]{\marginpar{MNV: \small\textbf{##1}}}
	}
	}
\fi

% === gestion des accents


%soit iso latin1
%\usepackage[cyr]{aeguill}
%soit utf8
\usepackage[utf8]{inputenc}
\usepackage[T1]{fontenc}


%%%  Pour citation \Cref (cref, ref)	

\usepackage{cleveref}


%%%%%%%%%%%%%%%%%%%%%%%%%%%%%%%%%%%%%%%%%%%%%%%%%
%
%
%                      Gestion des accents spéciaux.
%
%
%%%%%%%%%%%%%%%%%%%%%%%%%%%%%%%%%%%%%%%%%%%%%%%%%


%% MOI A LA MAIN:
\usepackage{bbm}

\DeclareUnicodeCharacter{2212}{-}
\DeclareUnicodeCharacter{2297}{\ensuremath{\otimes}}       
\DeclareUnicodeCharacter{2A33}{\ensuremath{\smashtimes}}
\DeclareUnicodeCharacter{1D56F}{\ensuremath{\mathfrak{D}}}  
\DeclareUnicodeCharacter{1D40E}{\ensuremath{\mathbf{O}}}  
\DeclareUnicodeCharacter{1D427}{\ensuremath{\mathbf{n}}}  
\DeclareUnicodeCharacter{3008}{\ensuremath{\langle}} % 〈 
\DeclareUnicodeCharacter{3009}{\ensuremath{\rangle}} %  〉

% (possible de voir: http://tug.ctan.org/info/symbols/comprehensive/symbols-letter.pdf)

%% SOURCE:  https://github.com/agda/agda/blob/master/doc/user-manual/unicode-symbols-sphinx.tex.txt

%%%%%%%%%%%%%%%%%%%%%%%%%%%%%%%%%%%%%%%%%%%%%%%%%%%%%%%%%%%%%%%%%%%%%%%%%%%%%%
% Please keep the below list ordered by code points!

\DeclareUnicodeCharacter{00AC}{\ensuremath{\neg}}                     % Symbol ¬
\DeclareUnicodeCharacter{00B1}{\ensuremath{\pm}}                      % Symbol ±
\DeclareUnicodeCharacter{00B2}{\ensuremath{^2}}                       % Symbol ²
\DeclareUnicodeCharacter{00B3}{\ensuremath{^3}}                       % Symbol ³
\DeclareUnicodeCharacter{00B7}{\ensuremath{\cdot}}                    % Symbol ·
\DeclareUnicodeCharacter{00B9}{\ensuremath{^1}}                       % Symbol ¹
\DeclareUnicodeCharacter{00D7}{\ensuremath{\times}}                   % Symbol ×
\DeclareUnicodeCharacter{02B0}{\ensuremath{^h}}                       % Symbol ʰ
\DeclareUnicodeCharacter{02B3}{\ensuremath{^r}}                       % Symbol ʳ
\DeclareUnicodeCharacter{02E2}{\ensuremath{^s}}                       % Symbol ˢ
\DeclareUnicodeCharacter{0302}{\ensuremath{\hat{\phantom{x}}}}        % Symbol ̂
\DeclareUnicodeCharacter{0332}{\ensuremath{\underline{\phantom{x}}}}  % Symbol ̲
\DeclareUnicodeCharacter{0308}{\ensuremath{\ddot{\phantom{x}}}}       % Symbol ̈
\DeclareUnicodeCharacter{0393}{\ensuremath{\Gamma}}                   % Symbol Γ
\DeclareUnicodeCharacter{0394}{\ensuremath{\Delta}}                   % Symbol Δ
\DeclareUnicodeCharacter{0398}{\ensuremath{\Theta}}                   % Symbol Θ
\DeclareUnicodeCharacter{039B}{\ensuremath{\Lambda}}                  % Symbol Λ
\DeclareUnicodeCharacter{039E}{\ensuremath{\Xi}}                      % Symbol Ξ
\DeclareUnicodeCharacter{03A0}{\ensuremath{\Pi}}                      % Symbol Π
\DeclareUnicodeCharacter{03A3}{\ensuremath{\Sigma}}                   % Symbol Σ


% There is not `\Tau` command in LaTeX, so we use `\mathrm{T}`. See
% http://tex.stackexchange.com/questions/1368/why-can-i-only-use-some-capital-greek-letters-inside-my-equations.
\DeclareUnicodeCharacter{03A4}{\ensuremath{\mathrm{T}}}  % Symbol Τ

\DeclareUnicodeCharacter{03A5}{\ensuremath{\Upsilon}}   % Symbol Υ
\DeclareUnicodeCharacter{03A6}{\ensuremath{\Phi}}       % Symbol Φ
\DeclareUnicodeCharacter{03A8}{\ensuremath{\Psi}}       % Symbol Ψ
\DeclareUnicodeCharacter{03A9}{\ensuremath{\Omega}}     % Symbol Ω
\DeclareUnicodeCharacter{03B1}{\ensuremath{\alpha}}     % Symbol α
\DeclareUnicodeCharacter{03B2}{\ensuremath{\beta}}      % Symbol β
\DeclareUnicodeCharacter{03B3}{\ensuremath{\gamma}}     % Symbol γ
\DeclareUnicodeCharacter{03B4}{\ensuremath{\delta}}     % Symbol δ
\DeclareUnicodeCharacter{03B5}{\ensuremath{\epsilon}}   % Symbol ε
\DeclareUnicodeCharacter{03B6}{\ensuremath{\zeta}}      % Symbol ζ
\DeclareUnicodeCharacter{03B7}{\ensuremath{\eta}}       % Symbol η
\DeclareUnicodeCharacter{03B8}{\ensuremath{\theta}}     % Symbol θ
\DeclareUnicodeCharacter{03B9}{\ensuremath{\iota}}      % Symbol ι
\DeclareUnicodeCharacter{03BA}{\ensuremath{\kappa}}     % Symbol κ
\DeclareUnicodeCharacter{03BB}{\ensuremath{\lambda}}    % Symbol λ
\DeclareUnicodeCharacter{03BC}{\ensuremath{\mu}}        % Symbol μ
\DeclareUnicodeCharacter{03BD}{\ensuremath{\nu}}        % Symbol ν
\DeclareUnicodeCharacter{03BE}{\ensuremath{\xi}}        % Symbol ξ
\DeclareUnicodeCharacter{03C0}{\ensuremath{\pi}}        % Symbol π
\DeclareUnicodeCharacter{03C1}{\ensuremath{\rho}}       % Symbol ρ
\DeclareUnicodeCharacter{03C2}{\ensuremath{\varsigma}}  % Symbol ς
\DeclareUnicodeCharacter{03C3}{\ensuremath{\sigma}}     % Symbol σ
\DeclareUnicodeCharacter{03C4}{\ensuremath{\tau}}       % Symbol τ
\DeclareUnicodeCharacter{03C5}{\ensuremath{\upsilon}}   % Symbol υ
\DeclareUnicodeCharacter{03C6}{\ensuremath{\varphi}}    % Symbol φ
\DeclareUnicodeCharacter{03C7}{\ensuremath{\chi}}       % Symbol χ
\DeclareUnicodeCharacter{03C8}{\ensuremath{\psi}}       % Symbol ψ
\DeclareUnicodeCharacter{03C9}{\ensuremath{\omega}}     % Symbol ω
\DeclareUnicodeCharacter{03D1}{\ensuremath{\vartheta}}  % Symbol ϑ

\DeclareUnicodeCharacter{0432}{\ensuremath{\mathrm{PDF\;TODO}}}  % Symbol в
\DeclareUnicodeCharacter{0435}{\ensuremath{\mathrm{PDF\;TODO}}}  % Symbol е
\DeclareUnicodeCharacter{0438}{\ensuremath{\mathrm{PDF\;TODO}}}  % Symbol и
\DeclareUnicodeCharacter{043C}{\ensuremath{\mathrm{PDF\;TODO}}}  % Symbol м
\DeclareUnicodeCharacter{0440}{\ensuremath{\mathrm{PDF\;TODO}}}  % Symbol р
\DeclareUnicodeCharacter{0442}{\ensuremath{\mathrm{PDF\;TODO}}}  % Symbol т
\DeclareUnicodeCharacter{041F}{\ensuremath{\mathrm{PDF\;TODO}}}  % Symbol П
\DeclareUnicodeCharacter{0BA8}{\ensuremath{\mathrm{PDF\;TODO}}}  % Symbol ந
\DeclareUnicodeCharacter{0BBF}{\ensuremath{\mathrm{PDF\;TODO}}}  % Symbol  ி
\DeclareUnicodeCharacter{1100}{\ensuremath{\mathrm{PDF\;TODO}}}  % Symbol ᄀ
\DeclareUnicodeCharacter{11F9}{\ensuremath{\mathrm{PDF\;TODO}}}  % Symbol ᇹ

\DeclareUnicodeCharacter{1D48}{\ensuremath{^d}}  % Symbol ᵈ
\DeclareUnicodeCharacter{1D50}{\ensuremath{^m}}  % Symbol ᵐ
\DeclareUnicodeCharacter{1D56}{\ensuremath{^p}}  % Symbol ᵖ
\DeclareUnicodeCharacter{1D62}{\ensuremath{_i}}  % Symbol ᵢ
\DeclareUnicodeCharacter{1D63}{\ensuremath{_r}}  % Symbol ᵣ
\DeclareUnicodeCharacter{1D64}{\ensuremath{_u}}  % Symbol ᵤ
\DeclareUnicodeCharacter{1D9C}{\ensuremath{^c}}  % Symbol ᶜ

% The command `\text` requires `\usepackage{amsmath}` and the command
% `\textasciiacute` requires `\usepackage{textcomp}`.

\DeclareUnicodeCharacter{2032}{\ensuremath{\text{\textasciiacute}}}  % Symbol ′

\DeclareUnicodeCharacter{2070}{\ensuremath{^0}}  % Symbol ⁰
\DeclareUnicodeCharacter{2074}{\ensuremath{^4}}  % Symbol ⁴
\DeclareUnicodeCharacter{2075}{\ensuremath{^5}}  % Symbol ⁵
\DeclareUnicodeCharacter{2076}{\ensuremath{^6}}  % Symbol ⁶
\DeclareUnicodeCharacter{2077}{\ensuremath{^7}}  % Symbol ⁷
\DeclareUnicodeCharacter{2078}{\ensuremath{^8}}  % Symbol ⁸
\DeclareUnicodeCharacter{2079}{\ensuremath{^9}}  % Symbol ⁹
\DeclareUnicodeCharacter{207A}{\ensuremath{^+}}  % Symbol ⁺
\DeclareUnicodeCharacter{207B}{\ensuremath{^-}}  % Symbol ⁻
\DeclareUnicodeCharacter{207F}{\ensuremath{^n}}  % Symbol ⁿ
\DeclareUnicodeCharacter{2080}{\ensuremath{_0}}  % Symbol ₀
\DeclareUnicodeCharacter{2081}{\ensuremath{_1}}  % Symbol ₁
\DeclareUnicodeCharacter{2082}{\ensuremath{_2}}  % Symbol ₂
\DeclareUnicodeCharacter{2083}{\ensuremath{_3}}  % Symbol ₃
\DeclareUnicodeCharacter{2084}{\ensuremath{_4}}  % Symbol ₄
\DeclareUnicodeCharacter{2085}{\ensuremath{_5}}  % Symbol ₅
\DeclareUnicodeCharacter{2086}{\ensuremath{_6}}  % Symbol ₆
\DeclareUnicodeCharacter{2087}{\ensuremath{_7}}  % Symbol ₇
\DeclareUnicodeCharacter{2088}{\ensuremath{_8}}  % Symbol ₈
\DeclareUnicodeCharacter{2089}{\ensuremath{_9}}  % Symbol ₉
\DeclareUnicodeCharacter{208A}{\ensuremath{_+}}  % Symbol ₊
\DeclareUnicodeCharacter{208B}{\ensuremath{_-}}  % Symbol ₋
\DeclareUnicodeCharacter{2091}{\ensuremath{_e}}  % Symbol ₑ
\DeclareUnicodeCharacter{2095}{\ensuremath{_h}}  % Symbol ₕ
\DeclareUnicodeCharacter{2096}{\ensuremath{_k}}  % Symbol ₖ
\DeclareUnicodeCharacter{2097}{\ensuremath{_l}}  % Symbol ₗ
\DeclareUnicodeCharacter{2098}{\ensuremath{_m}}  % Symbol ₘ
\DeclareUnicodeCharacter{2099}{\ensuremath{_n}}  % Symbol ₙ
\DeclareUnicodeCharacter{209A}{\ensuremath{_p}}  % Symbol ₚ
\DeclareUnicodeCharacter{209B}{\ensuremath{_s}}  % Symbol ₛ
\DeclareUnicodeCharacter{209C}{\ensuremath{_t}}  % Symbol ₜ

\DeclareUnicodeCharacter{2113}{\ensuremath{\mathpzc{l}}} % Symbol ℓ

% The command `\mathbbm` requires `\usepackage{bbm}`.
\DeclareUnicodeCharacter{2115}{\ensuremath{\mathbbm{N}}}  % Symbol ℕ
\DeclareUnicodeCharacter{211A}{\ensuremath{\mathbbm{Q}}}  % Symbol ℚ
\DeclareUnicodeCharacter{211D}{\ensuremath{\mathbbm{R}}}  % Symbol ℝ
\DeclareUnicodeCharacter{2124}{\ensuremath{\mathbbm{Z}}}  % Symbol ℤ

\DeclareUnicodeCharacter{2135}{\ensuremath{\aleph}}           % Symbol ℵ
\DeclareUnicodeCharacter{2190}{\ensuremath{\leftarrow}}       % Symbol ←
\DeclareUnicodeCharacter{2192}{\ensuremath{\rightarrow}}      % Symbol →
\DeclareUnicodeCharacter{2194}{\ensuremath{\leftrightarrow}}  % Symbol ↔

% The command `\twoheadrightarrow` requires `\usepackage{amssymb}`.
\DeclareUnicodeCharacter{21A0}{\ensuremath{\twoheadrightarrow}}  % Symbol ↠

\DeclareUnicodeCharacter{21A6}{\ensuremath{\mapsto}}  % Symbol ↦
\DeclareUnicodeCharacter{21A6}{\ensuremath{\mapsto}}  % Symbol ↦

% The command `\restriction` requires `\usepackage{amssymb}`.
\DeclareUnicodeCharacter{21BE}{\ensuremath{\restriction}}  % Symbol ↾

\DeclareUnicodeCharacter{21D0}{\ensuremath{\Leftarrow}}       % Symbol ⇐
\DeclareUnicodeCharacter{21D2}{\ensuremath{\Rightarrow}}      % Symbol ⇒
\DeclareUnicodeCharacter{21D4}{\ensuremath{\Leftrightarrow}}  % Symbol ⇔

\DeclareUnicodeCharacter{2200}{\ensuremath{\forall}}      % Symbol ∀
\DeclareUnicodeCharacter{2203}{\ensuremath{\exists}}      % Symbol ∃
\DeclareUnicodeCharacter{2204}{\ensuremath{\not\exists}}  % Symbol ∃
\DeclareUnicodeCharacter{2205}{\ensuremath{\emptyset}}    % Symbol ∅
\DeclareUnicodeCharacter{2208}{\ensuremath{\in}}          % Symbol ∈
\DeclareUnicodeCharacter{2209}{\ensuremath{\not\in}}      % Symbol ∉
\DeclareUnicodeCharacter{220B}{\ensuremath{\ni}}          % Symbol ∋
\DeclareUnicodeCharacter{220C}{\ensuremath{\not\ni}}      % Symbol ∌

% The command `\blacksquare` requires `\usepackage{amssymb}`.
\DeclareUnicodeCharacter{220E}{{\tiny \ensuremath{\blacksquare}}} % Symbol ∎

\DeclareUnicodeCharacter{220F}{{\tiny \ensuremath{\prod}}}  % Symbol ∏
\DeclareUnicodeCharacter{2211}{\ensuremath{\sum}}           % Symbol ∑
\DeclareUnicodeCharacter{2218}{\ensuremath{\circ}}          % Symbol ∘
\DeclareUnicodeCharacter{2219}{\ensuremath{\bullet}}        % Symbol ∙
\DeclareUnicodeCharacter{221E}{\ensuremath{\infty}}         % Symbol ∞
\DeclareUnicodeCharacter{2223}{\ensuremath{\mid}}           % Symbol ∣
\DeclareUnicodeCharacter{2225}{\ensuremath{\parallel}}      % Symbol ∥
\DeclareUnicodeCharacter{2227}{\ensuremath{\wedge}}         % Symbol ∧
\DeclareUnicodeCharacter{2228}{\ensuremath{\vee}}           % Symbol ∨
\DeclareUnicodeCharacter{2229}{\ensuremath{\cap}}           % Symbol ∩
\DeclareUnicodeCharacter{222A}{\ensuremath{\cup}}           % Symbol ∪
\DeclareUnicodeCharacter{222B}{\ensuremath{\int}}           % Symbol ∫
\DeclareUnicodeCharacter{2234}{\ensuremath{\therefore}}     % Symbol ∴

% The command `\dblcolon` requires `\usepackage{mathtools}`.
\DeclareUnicodeCharacter{2237}{\ensuremath{\dblcolon}} % Symbol ∷

\newcommand{\dotminus}{\mathop{\mbox{$-^{\hspace{-.5em}\cdot}\,$}}}
\DeclareUnicodeCharacter{2238}{\ensuremath{\dotminus}}  % Symbol ∸

\DeclareUnicodeCharacter{223C}{\ensuremath{\sim}}       % Symbol ∼
\DeclareUnicodeCharacter{2243}{\ensuremath{\simeq}}     % Symbol ≃
\DeclareUnicodeCharacter{2245}{\ensuremath{\cong}}      % Symbol ≅
\DeclareUnicodeCharacter{2248}{\ensuremath{\approx}}    % Symbol ≈
\DeclareUnicodeCharacter{224D}{\ensuremath{\asymp}}       % Symbol ≍ : OLIVIER
\DeclareUnicodeCharacter{2250}{\ensuremath{\doteq}}     % Symbol ≐

% The command `\coloneqq` requires `\usepackage{mathtools}`.
\DeclareUnicodeCharacter{2254}{\ensuremath{\coloneqq}}  % Symbol ≔

\newcommand{\eqq}{\stackrel{\text{\tiny ?}}{=}}
\DeclareUnicodeCharacter{225F}{\ensuremath{\eqq}}  % Symbol ≟

\DeclareUnicodeCharacter{2260}{\ensuremath{\neq}}         % Symbol ≠
\DeclareUnicodeCharacter{2261}{\ensuremath{\equiv}}       % Symbol ≡
\DeclareUnicodeCharacter{2262}{\ensuremath{\not\equiv}}   % Symbol ≢
\DeclareUnicodeCharacter{2264}{\ensuremath{\le}}          % Symbol ≤
\DeclareUnicodeCharacter{2265}{\ensuremath{\ge}}          % Symbol ≥
\DeclareUnicodeCharacter{226E}{\ensuremath{\nless}}       % Symbol ≮
\DeclareUnicodeCharacter{226F}{\ensuremath{\ngtr}}        % Symbol ≯
\DeclareUnicodeCharacter{2270}{\ensuremath{\nleq}}        % Symbol ≰
\DeclareUnicodeCharacter{2271}{\ensuremath{\ngeq}}        % Symbol ≱
\DeclareUnicodeCharacter{227A}{\ensuremath{\prec}}        % Symbol ≺
\DeclareUnicodeCharacter{227C}{\ensuremath{\preceq}}      % Symbol ≼
\DeclareUnicodeCharacter{2282}{\ensuremath{\subset}}      % Symbol ⊂
\DeclareUnicodeCharacter{2284}{\ensuremath{\not\subset}} % Symbol ⊄
\DeclareUnicodeCharacter{2283}{\ensuremath{\supset}}      % Symbol ⊃
\DeclareUnicodeCharacter{2286}{\ensuremath{\subseteq}}    % Symbol ⊆
\DeclareUnicodeCharacter{228E}{\ensuremath{\uplus}}       % Symbol ⊎
\DeclareUnicodeCharacter{2291}{\ensuremath{\sqsubseteq}}  % Symbol ⊑
\DeclareUnicodeCharacter{2293}{\ensuremath{\sqcap}}       % Symbol ⊓
\DeclareUnicodeCharacter{2294}{\ensuremath{\sqcup}}       % Symbol ⊔
\DeclareUnicodeCharacter{2295}{\ensuremath{\oplus}}       % Symbol ⊕
\DeclareUnicodeCharacter{22A2}{\ensuremath{\vdash}}       % Symbol ⊢
\DeclareUnicodeCharacter{22A4}{\ensuremath{\top}}         % Symbol ⊤
\DeclareUnicodeCharacter{22A5}{\ensuremath{\bot}}         % Symbol ⊥
\DeclareUnicodeCharacter{22C0}{\ensuremath{\bigwedge}}    % Symbol ⋀
\DeclareUnicodeCharacter{22C2}{\ensuremath{\bigcap}}      % Symbol ⋂
\DeclareUnicodeCharacter{22C3}{\ensuremath{\bigcup}}      % Symbol ⋃
\DeclareUnicodeCharacter{22EE}{\ensuremath{\vdots}}       % Symbol ⋮
\DeclareUnicodeCharacter{22EF}{\ensuremath{\cdots}}       % Symbol ⋯

% The command `\iddots` requires `\usepackage{mathdots}`.
\DeclareUnicodeCharacter{22F0}{\ensuremath{\iddots}} % Symbol ⋰

\DeclareUnicodeCharacter{230A}{\ensuremath{\lfloor}}  % Symbol ⌊
\DeclareUnicodeCharacter{230B}{\ensuremath{\rfloor}}  % Symbol ⌋
\DeclareUnicodeCharacter{2308}{\ensuremath{\lceil}}   % Symbol ⌈
\DeclareUnicodeCharacter{2309}{\ensuremath{\rceil}}   % Symbol ⌉
\DeclareUnicodeCharacter{2329}{\ensuremath{\langle}}  % Symbol 〈
\DeclareUnicodeCharacter{232A}{\ensuremath{\rangle}}  % Symbol 〉

% The command `\Return` requires `\usepackage{keystroke}`.
\DeclareUnicodeCharacter{23CE}{\ensuremath{\Return}}  % Symbol ⏎

% The command `\text` requires `\usepackage{amsmath}` and the command
% `\ding` requires `\usepackage{pifont}`.
\DeclareUnicodeCharacter{2460}{\ensuremath{\text{\ding{172}}}} % Symbol ①
\DeclareUnicodeCharacter{2461}{\ensuremath{\text{\ding{173}}}} % Symbol ②
\DeclareUnicodeCharacter{2462}{\ensuremath{\text{\ding{174}}}} % Symbol ③
\DeclareUnicodeCharacter{2463}{\ensuremath{\text{\ding{175}}}} % Symbol ④
\DeclareUnicodeCharacter{2464}{\ensuremath{\text{\ding{176}}}} % Symbol ⑤
\DeclareUnicodeCharacter{2465}{\ensuremath{\text{\ding{177}}}} % Symbol ⑥
\DeclareUnicodeCharacter{2466}{\ensuremath{\text{\ding{178}}}} % Symbol ⑦
\DeclareUnicodeCharacter{2467}{\ensuremath{\text{\ding{179}}}} % Symbol ⑧
\DeclareUnicodeCharacter{2468}{\ensuremath{\text{\ding{180}}}} % Symbol ⑨

\DeclareUnicodeCharacter{25C5}{\ensuremath{\triangleleft}}  % Symbol ◅

\DeclareUnicodeCharacter{2660}{\ensuremath{\spadesuit}} % Symbol ♠
\DeclareUnicodeCharacter{266D}{\ensuremath{\flat}}      % Symbol ♭
\DeclareUnicodeCharacter{266F}{\ensuremath{\sharp}}     % Symbol ♯

\DeclareUnicodeCharacter{2713}{\ensuremath{\checkmark}}  % Symbol ✓

% The commands `\llbracket` and `\rrbracket` require
% `\usepackage{stmaryrd}`.
\DeclareUnicodeCharacter{27E6}{\ensuremath{\llbracket}}  % Symbol ⟦
\DeclareUnicodeCharacter{27E7}{\ensuremath{\rrbracket}}  % Symbol ⟧

\DeclareUnicodeCharacter{27E8}{\ensuremath{\langle}}          % Symbol ⟨
\DeclareUnicodeCharacter{27E9}{\ensuremath{\rangle}}          % Symbol 〉
\DeclareUnicodeCharacter{27F6}{\ensuremath{\longrightarrow}}  % Symbol % ⟶

\DeclareUnicodeCharacter{2983}{\ensuremath{\mathrm{PDF\;TODO}}}  % Symbol ⦃
\DeclareUnicodeCharacter{2984}{\ensuremath{\mathrm{PDF\;TODO}}}  % Symbol ⦄
\DeclareUnicodeCharacter{2985}{\ensuremath{\mathrm{PDF\;TODO}}}  % Symbol ⦅
\DeclareUnicodeCharacter{2986}{\ensuremath{\mathrm{PDF\;TODO}}}  % Symbol ⦆

% The commands `\llparenthesis` and `\rrparenthesis` require
% `\usepackage{stmaryrd}`.
\DeclareUnicodeCharacter{2987}{\ensuremath{\llparenthesis}}  % Symbol ⦇
\DeclareUnicodeCharacter{2988}{\ensuremath{\rrparenthesis}}  % Symbol ⦈

\DeclareUnicodeCharacter{29F5}{\ensuremath{\setminus}}   % Symbol ⧵

\DeclareUnicodeCharacter{2A06}{\ensuremath{\bigcup}}  % Symbol ⋃
\DeclareUnicodeCharacter{2C7C}{\ensuremath{_j}}       % Symbol ⱼ

\DeclareUnicodeCharacter{D7B0}{\ensuremath{\mathrm{PDF\;TODO}}} % Symbol ힰ

% The command `\mathbbm{1}` requires `\usepackage{bbm}`.
\DeclareUnicodeCharacter{1D7D9}{\ensuremath{\mathbbm{1}}}  % Symbol 𝟙


%%%% MOI


\DeclareUnicodeCharacter{22AC}{\ensuremath{\neg\vdash}}
\DeclareUnicodeCharacter{279C}{\ensuremath{\to}}
\DeclareUnicodeCharacter{1D442}{ }
\DeclareUnicodeCharacter{2000}{ }
\DeclareUnicodeCharacter{1D440}{M}
\DeclareUnicodeCharacter{1D442}{O}
\DeclareUnicodeCharacter{20E3}{2}
\DeclareUnicodeCharacter{1D45A}{m}
\DeclareUnicodeCharacter{1D45B}{n}

\usepackage{pifont}
\DeclareUnicodeCharacter{2714}{\ding{52}}
\DeclareUnicodeCharacter{2705}{\ding{52}}
\DeclareUnicodeCharacter{2717}{\ding{55}}
\DeclareUnicodeCharacter{274C}{\ding{53}}


%% Tolérant tant que macros-markdown.tex n'est pas dans le dépôt :
\InputIfFileExists{macros-markdown}{}{\typeout{*** macros-markdown.tex introuvable: ignore ***}}


%\usepackage{PDFdraftcopy}


%-- Pour état des chapitres


%\newcommand\ETAT{\draftstring{BROUILLON}}
%
%\newcommand\brouillon{\renewcommand\ETAT{\draftstring{BROUILLON}}}
%
%\newcommand\final{\renewcommand\ETAT{\draftstring{\rien}}}
%
%\newcommand\travail{\renewcommand\ETAT{\draftstring{CHANTIER}}}
%
%\newcommand\bazar{\renewcommand\ETAT{\draftstring{BAZAR TOTAL}}}
%
%\final



%%%%%%%%%%%%%%%%%%%%%%%%%%%%%%%%%%%%%%%%%%%%%%%%%
%
%
%                      TRADUCTION EN COURS
%
%%%%%%%%%%%%%%%%%%%%%%%%%%%%%%%%%%%%%%%%%%%%%%%%%


% === volé Amaury



\usepackage{stmaryrd}
\usepackage{ifthen}
\usepackage{mathtools}
\usepackage{dsfont}
\PASSIAPXPROOF{\usepackage{thmtools}}
\usepackage{longtable}
\usepackage{mathdots}
\usepackage{booktabs}


%%%%%%%%%%%%%%%%%%%%%%%%%%%%%%%%%%%%%%%%%%%%%%%%%
%
%
%                      TIKZ STUFFS
%
%
%%%%%%%%%%%%%%%%%%%%%%%%%%%%%%%%%%%%%%%%%%%%%%%%%

%%% volé pour dessin des GPAC: PAS SUR QUE BESOIN de TOUT

\usepackage{tikz}
\usetikzlibrary{calc}
\usepgflibrary{arrows}
\usetikzlibrary{fit}
\usetikzlibrary{patterns}
\usetikzlibrary{decorations.pathreplacing}
\usetikzlibrary{shapes.multipart}
\usetikzlibrary{shapes.geometric}
\usetikzlibrary{shapes.misc}
\usetikzlibrary{positioning}
\usetikzlibrary{graphs}
\usetikzlibrary{topaths}
\usetikzlibrary{arrows.meta}
\usetikzlibrary{decorations.pathreplacing}
\usetikzlibrary{decorations.markings}
\usetikzlibrary{decorations.pathmorphing}
%
 \usetikzlibrary{calc}
 \usetikzlibrary{automata}
\usetikzlibrary{chains}
\usetikzlibrary{scopes}
 \usetikzlibrary{positioning}
 \usetikzlibrary{through,backgrounds}
\usepackage{circuitikz}
% % pour accolades




%%%%%%%%%%%%%%%%%%%%%%%%%%%%%%%%%%%%%%%%%%%%%%%%%
%
%
%                      MACROS locales
%
%
%%%%%%%%%%%%%%%%%%%%%%%%%%%%%%%%%%%%%%%%%%%%%%%%%


%====================
%                      Moi
%=====================


\newcommand\myaddress{ 
{\sc LIX, CNRS, Ecole Polytechnique, Institut Polytechnique de Paris, }\\ 91128 Palaiseau Cedex, FRANCE \\
{\small Olivier.Bournez@lix.polytechnique.fr}}



%====================
%                      un peu de trucs sales
%=====================

%=  pour style lipics et ses trucs

\providecommand\ifnotPOLYCHAPTERMODE[1]{#1}
\providecommand\ifPOLYCHAPTERMODE[1]{}

\ifnotPOLYCHAPTERMODE{
\ifdefined\thechapter
\providecommand\spnewtheorem[4]{\newtheorem{#1}{#2}[chapter]}
\providecommand\spnewtheoremi[5]{\newtheorem{#1}[#5]{#2}}
\else
\providecommand\spnewtheorem[4]{\newtheorem{#1}{#2}}
\providecommand\spnewtheoremi[5]{\newtheorem{#1}[#5]{#2}}
\fi
}
\ifPOLYCHAPTERMODE{
\providecommand\spnewtheorem[4]{\newtheorem{#1}{#2}}
\providecommand\spnewtheoremi[5]{\newtheorem{#1}[#5]{#2}}
}
\providecommand\nolipics[1]{#1}




%====================
%                      un peu de trucs franchement sales
%=====================


% on l'inclus pas ici.
%
%% defaults

\includeversion{NOTLFMI2020}
\includeversion{NOTCSE304}
\excludeversion{NOTDIFFUSIONFINALE}

%% Booleans (dont teste si existe ou pas, pour construire des macros)
% \SAFEMODE
% -> \DIFFUSIONFINALE (via entête-cours)
% -> \PAGEDEGARDE (via entête-cours)
% \POLYMODE
% \POLYCHAPTERMODE

%%%%%%%%%%%%%%%%%%%%%%%%%%%%%%%%%%%%%%%%%%%%%%%%%
%
%
%                      Hack pour DIFFUSION FINALE/ LFMI / Fin provisoire
%
%
%%%%%%%%%%%%%%%%%%%%%%%%%%%%%%%%%%%%%%%%%%%%%%%%%

%\def\DIFFUSIONFINALE{}

\providecommand\ifnotmodeDIFFUSIONFINALE[1]{
\ifdefined\DIFFUSIONFINALE
\else
#1
\fi
}

\providecommand\ifmodeDIFFUSIONFINALE[1]{
\ifdefined\DIFFUSIONFINALE
%% 2021: Comprend pas mais assure
\ifdefined\SAFEMODE
\else
#1
\fi
%#1
\fi
}



\ifnotmodeDIFFUSIONFINALE{\includeversion{NOTDIFFUSIONFINALE}}
\ifnotmodeDIFFUSIONFINALE{\newcommand\FINPROVISOIRELFMI[1]{}}

\providecommand\FINPROVISOIRELFMI{
  \section{Bibliographic notes}

  The current text is highly based (sometimes copied) from:
  
  \bibliographystyle{alpha}
  \bibliography{/Users/bournez/bibliographie/Bib-Files/bournez-avantbib2DOI,/Users/bournez/bibliographie/Bib-Files/perso}
  \end{document}
  }

%\def\POLYCHAPTERMODE{}

\providecommand\ifnotPOLYCHAPTERMODE[1]{
\ifdefined\POLYCHAPTERMODE
\else
#1
\fi
}

\providecommand\ifmodePOLYCHAPTERMODE[1]{
\ifdefined\POLYCHAPTERMODE
#1
\fi
}


%%%%%%%%%%%%%%%%%%%%%%%%%%%%%%%%%%%%%%%%%%%%%%%%%
%
%
%                      POUR INDICATIONS: ETAT
%
%
%%%%%%%%%%%%%%%%%%%%%%%%%%%%%%%%%%%%%%%%%%%%%%%%%


\newcommand\ABSTRACTINTERNAL[1]{
  \montruc[color=black!30]{
  \hrule
  \hrule
  \danger   ABSTRACT:\\[0.5cm] #1

  \vspace{0.5cm}
  \hrule
  \hrule
  }
  }

\newcommand\ABSTRACT[1]{
  \newpage
  \montruc[color=black!20]{
  \hrule
  \hrule
  \danger   ABSTRACT:\\[0.5cm] #1

  \vspace{0.5cm}
  \hrule
  \hrule
  }
  }

  
\newcommand\PARTIEFINALISEE{
  \newpage
  \montruc[color=black!15]{
  \hrule
  \hrule
  \danger   FINALISE A PARTIR ICI
  \hrule
  \hrule
  }
  }
\newcommand\PARTIESEMIFINALISEE{
  \newpage
  \montruc[color=black!15]{
  \hrule
  \hrule
  \danger   1/2 FINALISE A PARTIR ICI
  \hrule
  \hrule
  }
  }


\newcommand\PARTIENONFINALISEE{
  \newpage
  \montruc[color=black!30]{
  \hrule
  \hrule
  \danger   NON FINALISE A PARTIR ICI
  \hrule
  \hrule
  }
  }


%%%%%%%%%%%%%%%%%%%%%%%%%%%%%%%%%%%%%%%%%%%%%%%%%
%
%
%                      Affichage titre du cours: générique
%
%
%%%%%%%%%%%%%%%%%%%%%%%%%%%%%%%%%%%%%%%%%%%%%%%%%


  \newcommand\POLYWARNING[1]{
    
    \chapter*{Preliminaries}

    
\begin{center} 
  These are notes for the course: \\[1cm]
  \textbf{#1}
\end{center}

\bigskip
\fbox{\fbox{
\begin{minipage}{10cm}
\begin{center}
These notes change from days to days.  \\
Please do  not distribute. \\
\end{center}
\end{minipage}
}}

\bigskip
\begin{center}
  \textbf{
    Version of \today.
    }
\end{center}

 % This document is not stable.

\vfill
\begin{center}
Any comment (even about typos, orthography) is welcome: \\
send an email to bournez@lix.polytechnique.fr
\end{center}
}


\newcommand\MYCOURSEON[2]{ 
	%% External self-references
	\ifnotPOLYMODE{
		\ifdefined\MAINAUXFILE
			\externaldocument{\MAINAUXFILE}
		\fi
		}
	
	\ifdefined\PASLAPREMIEREFOIS
		\newpage
	\else
		\makeindex

     		\begin{document}
	\fi
	\gdef\PASLAPREMIEREFOIS{}

	\ifdefined\XPOLY
		\XPAGEDEGARDE{\TITRECOURS\ifPOLYCHAPTERMODE{\\[0.5cm] \Large \bf \LANGUE{Chapter:}{Chapitre:}  #2}}{\SUBTITRECOURS}
		\ifPOLYCHAPTERMODE{
		\renewcommand\thesection{\arabic{section}}
		\chapter*{#2}
		}
	\else
	
      \title{ \ifnotmodeDIFFUSIONFINALE{Course  #1 -}
      \ifPOLYCHAPTERMODE{\LANGUE{Chapter:}{Chapitre:} }
      #2}
      \author{
        Olivier Bournez% \inst{1}
        \\[0.5cm] %{\myaddress}
      }
      \date{Version of \today}
      % \institute{\myaddress}
      




      \ifnotPOLYCHAPTERMODE{
      	\maketitle
		}
      \ifPOLYCHAPTERMODE{\ifdefined\DOPAGEDEGARDE
      			\maketitle
		\fi
		\renewcommand\thesection{\arabic{section}}
%		\renewcommand\thesubsection{\arabic{section}.\ara}
      		\chapter*{#2}}

      % \begin{abstract}
      % \end{abstract}
      \fi
      


      \DOINDEXPRINT
      
      %% 2021: (meme si je comprends pas)
      \ifmodeDIFFUSIONFINALE{
		\excludeversion{recherche}
	}


    }



%%%%%%%%%%%%%%%%%%%%%%%%%%%%%%%%%%%%%%%%%%%%%%%%%
%
%
%                      Affichage titre du cours: variantes
%
%
%%%%%%%%%%%%%%%%%%%%%%%%%%%%%%%%%%%%%%%%%%%%%%%%%

    
\newcommand\MYCOURSE[1]{
  \MYCOURSEON{\csname numerocours#1\endcsname}{\csname nomcours#1\endcsname}
  
    \ifnotmodeDIFFUSIONFINALE{}%\listoftodos}

  }
  
  \newcommand\MYCOURSEnew[1]{
  \MYCOURSEON{???}{#1}
  
  \ifnotmodeDIFFUSIONFINALE{}%\listoftodos}

  }

%   ---------------------------
%   Pour compatilité chapter
%  ----------------------------

\newcommand\PASPREMIEREPAGE[1]{}
\providecommand\mychapter[1]{% \PASPREMIEREPAGE{\newpage} 
  \renewcommand\PASPREMIEREPAGE[1]{##1}  \hrule\hrule \section*{#1}
  \hrule\hrule \  \\[1cm]}
\providecommand\chapter[1]{\PASPREMIEREPAGE{\newpage} 
  \renewcommand\PASPREMIEREPAGE[1]{##1}  \hrule\hrule \section*{#1}
  \hrule\hrule \  \\[1cm]}

%   ---------------------------
%   if not sub
%  ----------------------------


\providecommand\IFNOTSUB[1]{#1}
\IFNOTSUB{
  
%\documentclass{article}
\usepackage{versions}
\includeversion{pasweb}
%
\providecommand\nolncs[1]{#1}	

%% Chemins par défaut (surchargables avant \input{macros} via \providecommand,
%% ou après via \renewcommand) :
\providecommand\BIBFILES{@@reference-biblio}
\providecommand\FIGCOMMONS{/Users/bournez/lib/LaTeX/Perso/FIG-COMMONS}

\newcommand\affiliationofficielleolivier{LIX, CNRS, École polytechnique, Institut Polytechnique de Paris, Palaiseau, France}
\newcommand\affiliationofficielleolivierslides{LIX, CNRS, École polytechnique, \\ Institut Polytechnique de Paris, Palaiseau, France}


%%%%%%%%%%%%%%%%%%%%%%%%
%
%              begin. Switch modes compatibilités
%
%%%%%%%%%%%%%%%%%%%%%%%

%% Apxproof mode
\providecommand\apxproofmode[1]{}
	% sinon: \newcommand\apxproof[1]{#1}
	
\providecommand\apxproofmodesubsection[1]{}
	% sinon: \newcommand\apxproofmodesubsection[1]{#1}
% utiliser: \appendixapxproofsubsection au lieu de \appendix dans ce cas. 
		

%% Lipics mode
\providecommand\lipicsmode{}
	% sinon: \newcommand\lipicsmode[1]{#1}
		% en vrai sert qu'à définir \nolipics
		
%%%%%%%%%%%%%%%%%%%%%%%%
%
%              end. Switch modes compatibilités
%
%%%%%%%%%%%%%%%%%%%%%%%


%% begin gestion des switchs

% en fait \nopilcs est la seule commande utilisée

\newcommand\nolipics[1]{#1}
\lipicsmode{\renewcommand\nolipics[1]{}}

%% Décorations (fourier + cadres colorés mdframed autour de definition,
%% theorem, exercice, example, remark, ...) : actives par défaut hors slides.
%% Pour les désactiver dans un document : \newcommand\NODECORATION[1]{}
%% avant \input{macros}. Retirées automatiquement en mode lipics.
\providecommand\NODECORATION[1]{\ifnotSLIDE{#1}}
\lipicsmode{\renewcommand\NODECORATION[1]{}}

\providecommand\PASSIAPXPROOF[1]{#1}
\apxproofmode{\renewcommand\PASSIAPXPROOF[1]{}}



%% end gestion des switchs

\def\MACROSPOINTEXOUVERT{}


\providecommand\ifnotTDEXAM[1]{#1}
\providecommand\ifnotSLIDE[1]{#1}


  
% pour preuves en appendix
\providecommand\PREUVESENAPPENDIX[1]{}



%%%% STYLES: APRES VOLS


\NODECORATION{
\usepackage{fourier}
}

%% Repli si fourier (et donc fourier-orns, qui définit \danger) n'est pas
%% chargé ; doit rester APRÈS le bloc ci-dessus, sinon fourier-orns échoue
%% avec « \danger already defined ».
\providecommand\danger{}

%%%% Box around environment
\usepackage[framemethod=tikz]{mdframed}

\newcommand\ENVCOLOR[1]{#1}
\newcommand\BEGINTEMPENVCOLOR[1]{\let\envcolorwas\ENVCOLOR\renewcommand\ENVCOLOR[1]{#1}}
\newcommand\ENDTEMPENVCOLOR{\let\ENVCOLOR\envcolorwas}


\NODECORATION{
\surroundwithmdframed[hidealllines=true,backgroundcolor=\ENVCOLOR{blue!10},roundcorner=8pt]{exercice}
\surroundwithmdframed[hidealllines=true,backgroundcolor=\ENVCOLOR{blue!10},roundcorner=8pt]{exercise}
\surroundwithmdframed[hidealllines=true,backgroundcolor=\ENVCOLOR{blue!10},roundcorner=8pt]{exercicee}
\surroundwithmdframed[hidealllines=true,backgroundcolor=\ENVCOLOR{blue!10},roundcorner=8pt]{exerciced}
\surroundwithmdframed[hidealllines=true,backgroundcolor=\ENVCOLOR{blue!10},roundcorner=8pt]{exercicec}
\surroundwithmdframed[hidealllines=true,backgroundcolor=\ENVCOLOR{blue!10},roundcorner=8pt]{exercicedc}
\surroundwithmdframed[hidealllines=true,backgroundcolor=\ENVCOLOR{orange!20},roundcorner=8pt]{example}
\surroundwithmdframed[hidealllines=true,backgroundcolor=\ENVCOLOR{orange!20},roundcorner=8pt]{exemple}
\surroundwithmdframed[hidealllines=true,backgroundcolor=\ENVCOLOR{orange!10},roundcorner=8pt]{remark}
\ifnotSLIDE{\surroundwithmdframed[backgroundcolor=\ENVCOLOR{green!15},roundcorner=8pt]{definition}}
\surroundwithmdframed[hidealllines=true,backgroundcolor=\ENVCOLOR{gray!35}]{proposition}
\ifnotSLIDE{\surroundwithmdframed[backgroundcolor=\ENVCOLOR{gray!35},roundcorner=8pt]{theorem}}
\surroundwithmdframed[hidealllines=true,backgroundcolor=\ENVCOLOR{gray!35}]{lemma}
\surroundwithmdframed[hidealllines=true,backgroundcolor=\ENVCOLOR{gray!35}]{corollary}


\surroundwithmdframed[backgroundcolor=red!35,roundcorner=8pt]{theoremaprouver}
\surroundwithmdframed[backgroundcolor=red!35,roundcorner=8pt]{conjectureaprouver}
}

%%% tcolorbox

\usepackage{tcolorbox}
\tcbuselibrary{skins}


  %%. Commande magique
%% \tcolorboxenvironment{⟨name⟩}{⟨options⟩}: attach options de tcolorbox à environnement
%% alternative: je cree des boites moi meme: voici des exemples


\newtcolorbox{maboiteconfiance}[2][]{colback=red!5!white, colframe=red!75!black,fonttitle=\bfseries, colbacktitle=red!85!black,enhanced,
attach boxed title to top center={yshift=-2mm},
  title={#2},#1}
  
\newtcolorbox{maboiteresume}[2][]{colback=red!5!white, colframe=red!75!black,fonttitle=\bfseries, colbacktitle=red!85!black,enhanced,
attach boxed title to top center={yshift=-2mm},
  title={#2},#1}
  \newtcolorbox{maboiteresumeperso}[2][]{colback=red!5!white, colframe=red!75!black,fonttitle=\bfseries, colbacktitle=red!85!black,enhanced,
attach boxed title to top center={yshift=-2mm},
  title={#2},#1}
  \newtcolorbox{maboitecommentaire}[2][]{colback=red!5!white, colframe=red!75!black,fonttitle=\bfseries, colbacktitle=red!85!black,enhanced,
attach boxed title to top center={yshift=-2mm},
  title={#2},#1}
\newtcolorbox{maboitetruc}[2][]{colback=red!5!white, colframe=red!75!black,fonttitle=\bfseries, colbacktitle=red!85!black,enhanced,
attach boxed title to top left={yshift=-2mm},
  title={#2},#1}

%% Boîte des mises en évidence CARE. Couleurs pensées pour la
%% protanopie : l'axe bleu-jaune reste discriminant (jaune = \IMPORTANT
%% d'Olivier, bleu = CARE), et la distinction ne repose pas que sur la
%% couleur (cadre + bandeau étiqueté, lisibles même en niveaux de gris).
\newtcolorbox{maboiteimportantcare}[2][]{colback=cyan!10!white, colframe=blue!60!black,fonttitle=\bfseries, colbacktitle=blue!70!black,enhanced,
attach boxed title to top center={yshift=-2mm},
  title={#2},#1}
  


%%%% FIN STYLES



%%%%%%%%%%%%%%%%%%%%%%%%%%%%%%%%%%%%%%%%%%%%%%%%%
%%%%%%%%%%%%%%%%%%%%%%%%%%%%%%%%%%%%%%%%%%%%%%%%%
%
%
%                      Version ENSEIGNEMENT 2020
%
%
%%%%%%%%%%%%%%%%%%%%%%%%%%%%%%%%%%%%%%%%%%%%%%%%%
%%%%%%%%%%%%%%%%%%%%%%%%%%%%%%%%%%%%%%%%%%%%%%%%%

%english, french
\ifdefined\CESTDUFRANCAIS
	\providecommand\LANGUE[2]{#2}
	\usepackage[french]{babel}
\else
	\providecommand\LANGUE[2]{#1}
\fi
\providecommand\LANGUEinv[2]{\LANGUE{#2}{#1}}

%%%%%%%%%%%%%%%%%%%%%%%%%%%%%%%%%%%%%%%%%%%%%%%%%
%
%
%                      Packages
%
%
%%%%%%%%%%%%%%%%%%%%%%%%%%%%%%%%%%%%%%%%%%%%%%%%%


%%% SURLIGNEUR
\providecommand\ifnotmodeDIFFUSIONFINALE[1]{
%% 2021: Comprend pas mais assure
\ifdefined\SAFEMODE
\else
#1
\fi
%#1
}
\providecommand\ifmodeDIFFUSIONFINALE[1]{
\ifdefined\SAFEMODE
#1
\else
\fi
}
\newcommand\SURLIGNE[1]{
\ifnotmodeDIFFUSIONFINALE{
\BEGINTEMPENVCOLOR{blue!30}
{                        \colorbox{blue!30}{

\noindent \begin{minipage}{\dimexpr\textwidth-2\fboxsep\relax}
#1
\end{minipage}
}
}
\ENDTEMPENVCOLOR}
\ifmodeDIFFUSIONFINALE{
\noindent \begin{minipage}{\textwidth}
#1
\end{minipage}
}
}


\newenvironment{confiance}{\begin{maboiteresume}[colback=brown!50]{Résumé ici}}{\end{maboiteresume}}
\newcommand\CONFIANCE[2][Confiance ici]{
\begin{maboiteconfiance}[colback=yellow!50]{#1}
\textbf{ATTENTION : #2}
\end{maboiteconfiance}
}

\newenvironment{resume}{\begin{maboiteresume}[colback=brown!50]{Résumé ici}}{\end{maboiteresume}}
\newcommand\RESUME[2][Résumé ici]{
\begin{maboiteresume}[colback=brown!50]{#1}
#2
\end{maboiteresume}
}

\newenvironment{resumeperso}{\begin{maboiteresumeperso}[colback=brown!50]{Résumé perso ici}}{\end{maboiteresumeperso}}
\newcommand\RESUMEPERSO[2][Résumé perso  ici]{
\begin{maboiteresumeperso}[colback=brown!50]{#1}
#2
\end{maboiteresumeperso}
%
%
%\todo[inline, color=brown]{résumé: #1}
%\colorbox{brown}{
%\noindent \begin{minipage}{\textwidth}
%{
%\bf 
%#2
%}
%\end{minipage}
%}
}

\newcommand\DANSLAMARGE[1]{
\ifnotmodeDIFFUSIONFINALE{
\marginpar{\textcolor{brown}{#1}}}
}

\newenvironment{commentaireperso}{\begin{maboitecommentaire}[colback=brown!30]{Commentaire perso} Commentaire perso:}{\end{maboitecommentaire}}
\newcommand\COMMENTAIREPERSO[2][Commentaire perso]{
\begin{maboitecommentaire}[colback=brown!30]{#1}
#2
\end{maboitecommentaire}
}


\newcommand\SURSURLIGNE[2][truc surligné ici]{
%\todo[inline, color=orange]{surligné: #1}
\ifnotmodeDIFFUSIONFINALE{
\BEGINTEMPENVCOLOR{blue!60}
	                 \colorbox{blue!60}{
\noindent \begin{minipage}{\dimexpr\textwidth-2\fboxsep\relax}
{
\bf 
#2
}
\end{minipage}
}
\ENDTEMPENVCOLOR
}
\ifmodeDIFFUSIONFINALE{#2}
}

\newcommand\IMPORTANT[1]{
\ifnotmodeDIFFUSIONFINALE{
\BEGINTEMPENVCOLOR{yellow!100}
		        \colorbox{yellow!100}{

\noindent \begin{minipage}{\dimexpr\textwidth-2\fboxsep\relax}
{ 
#1
}
\end{minipage}
}
\ENDTEMPENVCOLOR
}
\ifmodeDIFFUSIONFINALE{#1}
}

%% Pendant de \IMPORTANT, réservé aux mises en évidence CARE ;
%% \IMPORTANT reste réservé aux annotations d'Olivier. Comme \IMPORTANT,
%% redevient du texte simple en mode DIFFUSIONFINALE.
\newcommand\IMPORTANTCARE[2][Important (CARE)]{
\ifnotmodeDIFFUSIONFINALE{
\begin{maboiteimportantcare}{#1}
#2
\end{maboiteimportantcare}
}
\ifmodeDIFFUSIONFINALE{#2}
}


\newcommand\QUESTION[1]{
\ifnotmodeDIFFUSIONFINALE{
\BEGINTEMPENVCOLOR{oceanfoam!100}
		        \colorbox{oceanfoam!100}{

\noindent \begin{minipage}{\dimexpr\textwidth-2\fboxsep\relax}
{ 
#1
}
\end{minipage}
}
\ENDTEMPENVCOLOR
}
\ifmodeDIFFUSIONFINALE{#1}
}




\newcommand\SOUSLIGNE[1]{
\colorbox{lightgray}{

\noindent \begin{minipage}{\dimexpr\textwidth-2\fboxsep\relax}
{ \color{gray}
#1
}
\end{minipage}
}
}



%% Gestion des cross-references cross-referencing
\usepackage{xr-hyper}
\usepackage{hyperref}

%%
\usepackage{versions}
\usepackage{amsmath,amssymb,mathrsfs,amsfonts}
%\usepackage{mathabx}
\usepackage{listings}
\usepackage{url}
\usepackage{versions}
\usepackage{color}
\usepackage[colorinlistoftodos]{todonotes}
%\usepackage{todonotes}
\usepackage{proof}
\usepackage{xstring}

%% Index
\usepackage{makeidx}
\makeindex
\providecommand\DOINDEXPRINT{}
\providecommand\SIGNALEMOTNOUV[1]{}
\ifdefined\INDEXPRINT
	\ifnotSAFEMODE{
%	\usepackage{showidx}
%	\let\oldindex\index
%	\renewcommand{\index}[1]{\oldindex{#1}\marginpar{LA: \tiny#1}}
	\renewcommand\DOINDEXPRINT{
	      	\let\oldindex\index
		\renewcommand{\index}[1]{\oldindex{##1}\marginpar{IDXLA: \tiny##1}}
		\renewcommand\SIGNALEMOTNOUV[1]{\marginpar{MNV: \small\textbf{##1}}}
	}
	}
\fi

% === gestion des accents


%soit iso latin1
%\usepackage[cyr]{aeguill}
%soit utf8
\usepackage[utf8]{inputenc}
\usepackage[T1]{fontenc}


%%%  Pour citation \Cref (cref, ref)	

\usepackage{cleveref}


%%%%%%%%%%%%%%%%%%%%%%%%%%%%%%%%%%%%%%%%%%%%%%%%%
%
%
%                      Gestion des accents spéciaux.
%
%
%%%%%%%%%%%%%%%%%%%%%%%%%%%%%%%%%%%%%%%%%%%%%%%%%

\input{macros-accents}
%% Tolérant tant que macros-markdown.tex n'est pas dans le dépôt :
\InputIfFileExists{macros-markdown}{}{\typeout{*** macros-markdown.tex introuvable: ignore ***}}


%\usepackage{PDFdraftcopy}


%-- Pour état des chapitres


%\newcommand\ETAT{\draftstring{BROUILLON}}
%
%\newcommand\brouillon{\renewcommand\ETAT{\draftstring{BROUILLON}}}
%
%\newcommand\final{\renewcommand\ETAT{\draftstring{\rien}}}
%
%\newcommand\travail{\renewcommand\ETAT{\draftstring{CHANTIER}}}
%
%\newcommand\bazar{\renewcommand\ETAT{\draftstring{BAZAR TOTAL}}}
%
%\final



%%%%%%%%%%%%%%%%%%%%%%%%%%%%%%%%%%%%%%%%%%%%%%%%%
%
%
%                      TRADUCTION EN COURS
%
%%%%%%%%%%%%%%%%%%%%%%%%%%%%%%%%%%%%%%%%%%%%%%%%%


% === volé Amaury



\usepackage{stmaryrd}
\usepackage{ifthen}
\usepackage{mathtools}
\usepackage{dsfont}
\PASSIAPXPROOF{\usepackage{thmtools}}
\usepackage{longtable}
\usepackage{mathdots}
\usepackage{booktabs}


%%%%%%%%%%%%%%%%%%%%%%%%%%%%%%%%%%%%%%%%%%%%%%%%%
%
%
%                      TIKZ STUFFS
%
%
%%%%%%%%%%%%%%%%%%%%%%%%%%%%%%%%%%%%%%%%%%%%%%%%%

%%% volé pour dessin des GPAC: PAS SUR QUE BESOIN de TOUT

\usepackage{tikz}
\usetikzlibrary{calc}
\usepgflibrary{arrows}
\usetikzlibrary{fit}
\usetikzlibrary{patterns}
\usetikzlibrary{decorations.pathreplacing}
\usetikzlibrary{shapes.multipart}
\usetikzlibrary{shapes.geometric}
\usetikzlibrary{shapes.misc}
\usetikzlibrary{positioning}
\usetikzlibrary{graphs}
\usetikzlibrary{topaths}
\usetikzlibrary{arrows.meta}
\usetikzlibrary{decorations.pathreplacing}
\usetikzlibrary{decorations.markings}
\usetikzlibrary{decorations.pathmorphing}
%
 \usetikzlibrary{calc}
 \usetikzlibrary{automata}
\usetikzlibrary{chains}
\usetikzlibrary{scopes}
 \usetikzlibrary{positioning}
 \usetikzlibrary{through,backgrounds}
\usepackage{circuitikz}
% % pour accolades




%%%%%%%%%%%%%%%%%%%%%%%%%%%%%%%%%%%%%%%%%%%%%%%%%
%
%
%                      MACROS locales
%
%
%%%%%%%%%%%%%%%%%%%%%%%%%%%%%%%%%%%%%%%%%%%%%%%%%


%====================
%                      Moi
%=====================


\newcommand\myaddress{ 
{\sc LIX, CNRS, Ecole Polytechnique, Institut Polytechnique de Paris, }\\ 91128 Palaiseau Cedex, FRANCE \\
{\small Olivier.Bournez@lix.polytechnique.fr}}



%====================
%                      un peu de trucs sales
%=====================

%=  pour style lipics et ses trucs

\providecommand\ifnotPOLYCHAPTERMODE[1]{#1}
\providecommand\ifPOLYCHAPTERMODE[1]{}

\ifnotPOLYCHAPTERMODE{
\ifdefined\thechapter
\providecommand\spnewtheorem[4]{\newtheorem{#1}{#2}[chapter]}
\providecommand\spnewtheoremi[5]{\newtheorem{#1}[#5]{#2}}
\else
\providecommand\spnewtheorem[4]{\newtheorem{#1}{#2}}
\providecommand\spnewtheoremi[5]{\newtheorem{#1}[#5]{#2}}
\fi
}
\ifPOLYCHAPTERMODE{
\providecommand\spnewtheorem[4]{\newtheorem{#1}{#2}}
\providecommand\spnewtheoremi[5]{\newtheorem{#1}[#5]{#2}}
}
\providecommand\nolipics[1]{#1}




%====================
%                      un peu de trucs franchement sales
%=====================


% on l'inclus pas ici.
%\input{macros-moins-propres}

%%%%%%%%%%%%%%%%%%%%%%%%%%%%%%%%%%%%%%%%%%%%%%%%%
%
%
%                      Page de Garde Cours X
%
%
%%%%%%%%%%%%%%%%%%%%%%%%%%%%%%%%%%%%%%%%%%%%%%%%%


\includeversion{metdate}

\newcommand\XPAGEDEGARDE[2]{
\thispagestyle{empty}
\title{{#1} \\[2cm] %\\[3cm] 
{\large #2}
}
\author{
 Olivier Bournez% \inst{1}
\\[0.5cm] 
bournez@lix.polytechnique.fr%{\myaddress}
}
%\date{
%\vspace{2cm}
%\includegraphics[height=2.5cm]{/Users/bournez/lib/LaTeX/Perso/FIG-COMMONS/logo_X_new}
%}
\begin{metdate}
\date{
\vspace{2cm}
Version \LANGUE{of}{du} \today \\[1cm] 
\includegraphics[height=3.5cm]{/Users/bournez/lib/LaTeX/Perso/FIG-COMMONS/logo_X_new}
}
\end{metdate}
\maketitle
}


%%%%%%%%%%%%%%%%%%%%%%%%%%%%%%%%%%%%%%%%%%%%%%%%%
%
%
%                      MACROS UTILES
%
%
%%%%%%%%%%%%%%%%%%%%%%%%%%%%%%%%%%%%%%%%%%%%%%%%%

%%%                      Macros
 
%-------------
%--- vectors  ---
%-------------

 
\newcommand\vectorl[1]{{\mathbf#1}}
\newcommand\tu[1]{\vectorl{#1}}

\newcommand\vp{\vectorl{p}}
\newcommand\va{\vectorl{a}}
\newcommand\vb{\vectorl{b}}
\newcommand\vc{\vectorl{c}}
\newcommand\vf{\vectorl{f}}
\newcommand\vn{\vectorl{n}}
\newcommand\vx{\vectorl{x}}
\newcommand\vX{\vectorl{X}}
\newcommand\vy{\vectorl{y}}
\newcommand\vz{\vectorl{z}}
\newcommand\ve{\vectorl{e}}
\newcommand\vv{\vectorl{v}}
\newcommand\vr{\vectorl{r}}
\newcommand\vu{\vectorl{u}}
\newcommand\vA{\vectorl{A}}
\newcommand\vB{\vectorl{B}}
\newcommand\vC{\vectorl{C}}
\newcommand\vD{\vectorl{D}}

%-------------
%--- overline ---
%-------------


\newcommand\ox{\overline{x}}
\newcommand\oy{\overline{y}}
\newcommand\oc{\overline{c}}
\newcommand\oa{\overline{a}}
\newcommand\ow{\overline{w}}
\newcommand\od{\overline{d}}
\newcommand\ob{\overline{b}}
\newcommand\oP{\overline{P}}
\newcommand\oee{\overline{e}}
\newcommand\ot{\overline{t}}
\newcommand\ou{\overline{u}}
\newcommand\ov{\overline{v}}
\newcommand\oz{\overline{z}}
\newcommand\am{\overline{a}}
\newcommand\om{\overline{m}}
\newcommand\oj{\overline{j}}

%----------------
%--- ensembles ---
%----------------

\newcommand\Q{\mathbb{Q}}
\newcommand\R{\mathbb{R}}
\newcommand\Z{\mathbb{Z}}
\providecommand\C{\mathbb{C}}
\newcommand\N{\mathbb{N}}
\newcommand\K{\mathbb{K}}
\newcommand\F{\mathbb{F}}

\newcommand\dyadic{\mathbb{D}}
\newcommand\Zd{{\Z_2}}
\newcommand\Zp{{\Z_p}}
%
\newcommand\RP{\mathbb{R}^{>0}}
\newcommand\QP{\mathbb{Q}^{>0}}


%----------------
%--- lettres cal ---
%----------------

\newcommand\mK{\mathcal{K}}
\newcommand\mD{\mathcal{D}}
\newcommand\mA{\mathcal{A}}
\newcommand\mL{\mathcal{L}}
\newcommand\mX{\mathcal{X}}
\newcommand\mG{\mathcal{G}}
\newcommand\mE{\mathcal{E}}
\newcommand\mH{\mathcal{H}}
\newcommand\mR{\mathcal{R}}
\newcommand\mF{\mathcal{F}}
\newcommand\mJ{\mathcal{J}}
\newcommand\mO{\mathcal{O}}
\newcommand\mP{\mathcal{P}}
\newcommand\mN{\mathcal{N}}
\newcommand\mS{\mathcal{S}}
\newcommand\mM{\mathcal{M}}
\newcommand\mC{\mathcal{C}}
\newcommand\mV{\mathcal{V}}
\newcommand\mI{\mathcal{I}}
\newcommand\mT{\mathcal{T}}

\newcommand\fC{\mathcal{C}}

%----------------
%--- lettres frak ---
%----------------


\newcommand\kM{\mathfrak{M}}
\newcommand\kN{\mathfrak{N}}
\newcommand\kU{\mathfrak{U}}
\newcommand\kg{\mathfrak{g}}



%--------------------------
%--- Façon noter noms classes ---
%---------------------------

\newcommand{\cp}[1]{\operatorname{#1}}


%--------------------------
%--- Calculabilité  ---
%---------------------------

% -- classes
\LANGUE{
\newcommand\RE{\cp{CE}}
}
{\newcommand\RE{\cp{RE}}
}

\newcommand\Rec{\cp{D}}
\newcommand\PR{\cp{PR}}
\newcommand\Gregn{\mE_n}


%--------------------------
%--- Réductions  ---
%---------------------------

\newcommand\lem{\le_{m}}
\newcommand\lelog{\le_{L}}
\newcommand\equivlog{\equiv_{L}}



%--------------------------
%--- Complexité  ---
%---------------------------

% COMPLEXITE

% -- P
\newcommand{\Ptime}{\cp{P}}      % P
\newcommand\PTIME{\Ptime}
\newcommand\PTIMEs[1]{{\PTIME}_{#1}}
\newcommand\PTIMEsp[1]{{\PTIME}_{#1}^0}
\renewcommand\P{\PTIME}
\newcommand{\FPtime}{\cp{FPTIME}}

% -- NP
\newcommand\NP{\cp{NP}}
\newcommand{\NPtime}{\NP}
\newcommand\NPs[1]{\cp{NP}_{#1}}
\newcommand{\FNP}{\cp{FNP}}

\newcommand\NDP{\cp{NDP}}
\newcommand\coNDP{\cp{coNDP}}
\newcommand\NDPs[1]{\cp{NDP}_{#1}}

% -- coNP
\newcommand\coNP{\cp{coNP}}

% -- PSPACE
\newcommand\PSPACE{\cp{ PSPACE}}
\newcommand\NPSPACE{\cp{ NPSPACE}}
\newcommand{\Pspace}{\PSPACE}
\newcommand{\FPspace}{\cp{FPSPACE}}


% -- LOGSPACE
\newcommand\LSPACE{\cp{LOGSPACE}}
\newcommand\NLSPACE{\cp{NLOGSPACE}}
\newcommand\coNLSPACE{\cp{coNLOGSPACE}}
\newcommand\coNL{\cp{coNL}}

% -- Autres
\newcommand\EXPTIME{\cp{EXPTIME}}
\newcommand\NEXPTIME{\cp{NEXPTIME}}
\newcommand\EXPSPACE{\cp{EXPSPACE}}
\newcommand\PH{\cp{PH}}
\newcommand\NC{\cp{NC}}
\newcommand\AC{\cp{AC}}
\newcommand\PPAD{\cp{\mathbf{PPAD}}}
\newcommand\PLS{\cp{\mathbf{PLS}}}

% -- Reverse math
\newcommand{\WKL}{\ensuremath{\docheck{\mathsf{WKL}_0}}}
\newcommand{\ACA}{\ensuremath{\docheck{\mathsf{ACA}_0}}}
\newcommand{\ATR}{\ensuremath{\docheck{\mathsf{ATR}_0}}}
\newcommand{\RCAO}{\ensuremath{\docheck{\mathsf{RCA}_0}}}
\newcommand{\PiCA}{\ensuremath{\docheck{\Pi^1_1\!\text{-}\mathsf{CA}_0}}}


% -- circuits
\newcommand\CIRCUIT[2]{\cp{CIRCUIT(#1,#2)}}
\newcommand\UCIRCUIT[2]{\cp{UCIRCUIT(#1,#2)}}

% -- Classes proba
\newcommand\PP{\cp{PP}}
\newcommand\RPP{\cp{RP}}
\newcommand\ZPP{\cp{ZPP}}
\newcommand\coRP{\cp{coRP}}
\newcommand\BPP{{\cp{BPP}}}

% -- IP
\newcommand\IP{\cp{IP}}
\newcommand\PCP{\cp{PCP}}

% -- non-uniforme
\newcommand\nuP{\mathbb{P}}
\newcommand\nuPs[1]{\mathbb{P}_{#1}}
\newcommand\nuPsp[1]{\mathbb{P}_{#1}^0}
\newcommand\nuNP{\mathbb{N}\mathbb{P}}
%\newcommand\poly{/{poly}}
\newcommand\Ppoly{\PTIME_{\poly}}
\newcommand\NPpoly{\NP_{\poly}}

% -- conditions d'uniformité
\newcommand\Luniforme{\cp{L}}
\newcommand\BCuniforme{\cp{BC}}


% Classes de circuits
\newcommand\TC{\cp{TC}}
\newcommand\LFP{\cp{LFP}}
\newcommand\FAC{\cp{FAC}}
\newcommand\FACC{\cp{FACC}}
\newcommand\FTC{\cp{FTC}}
\newcommand\FNC{\cp{FNC}}

% Logiques temorelles
\newcommand\CTL{\cp{CTL}}
\newcommand\LTL{\cp{LTL}}

% -- mesure temps/espace/taille

\newcommand\DTIME{\cp{DTIME}}
\newcommand\TIME[1]{\cp{TIME(#1)}}
\newcommand\TIMEs[2]{\cp{TIME_{#2}(#1)}}
\newcommand\TIMEsp[2]{\cp{TIME^0_{#2}(#1)}}
\newcommand\NTIME[1]{\cp{NTIME(#1)}}
\newcommand\NTIMEs[2]{\cp{NTIME_{#2}(#1)}}
\newcommand\NTIMEsp[2]{\cp{NTIME_{#2}^0(#1)}}
\newcommand\DSPACE{\cp{DSPACE}}
\newcommand\SPACE[1]{\cp{SPACE(#1)}}
\newcommand\SPACEs[2]{\cp{SPACE_{#2}(#1)}}
\newcommand\NSPACE[1]{\cp{NSPACE(#1)}}
\newcommand\NSPACEs[2]{\cp{NSPACE_{#2}(#1)}}
\newcommand\coNSPACE[1]{\cp{coNSPACE(#1)}}
\newcommand\TIMEON[2]{\cp{TIME(#1,#2)}}
\newcommand\SPACEON[2]{\cp{SPACE(#1,#2)}}
\newcommand\SIZE[1]{\cp{ size(#1)}}
\newcommand\ATIME{\cp{ATIME}}
\newcommand\ASPACE{\cp{ASPACE}}


%% autre:
\newcommand\LH{\cp{LH}}
\newcommand\FO{\cp{FO}}
\newcommand\SO{\cp{SO}}
\newcommand\ESOhorn{\cp{ESO_{horn}}}
\newcommand\SOhorn{\cp{SO_{horn}}}
\newcommand\ACzero{\cp{AC_0}}
\newcommand{\LINSPACE}{\cp{LINSPACE}}
\newcommand{\mFLINSPACE}{\mF_{\cp{LINSPACE}}}
\newcommand{\LOGSPACE}{\cp{LOGSPACE}}
\newcommand\ALOGTIME{\cp{ALOGTIME}}
\newcommand\RF{\cp{RF}}


\newcommand\BOOL[1]{\cp{{BOOLEAN(#1)}}}
\newcommand\DET{\cp{ DET}}

% % -- dirty


\newcommand{\cPspace}{\cp{\sharp PSPACE}}
\newcommand{\cAPtime}{\cp{\sharp APTIME}}
%\newcommand{\FPspace}{\cp{FPSPACE}}
%\newcommand{\FPspaceN}{\cp{FPSPACE}^{+}}
%\newcommand{\FPspaceN}{\cp{FPSPACE}}



%-------------------------
%--- Problèmes de décision  ---
%--------------------------

\newcommand\decisionname[1]{\cp{#1}}


\newcommand\Luniv{\decisionname{L_{univ}}}
\newcommand\HALTINGPROBLEM{\decisionname{HALTING-PROBLEM}}
\newcommand\SAT{\decisionname{SAT}}
\newcommand\MAXtSAT{\decisionname{MAX3SAT}}
\newcommand\REACH{\decisionname{REACH}}
\newcommand\CIRCUITVALUE{\decisionname{CIRCUITVALUE}}
\newcommand\MONOTONECIRCUITVALUE{\decisionname{MONOTONECIRCUITVALUE}}
\newcommand\THEOREMS{\decisionname{THEOREMS}}
\newcommand\QCIRCUITSAT{\decisionname{QCIRCUITSAT}}
\newcommand\QSAT{\decisionname{QSAT}}
\newcommand\QBF{\decisionname{QBF}}
\newcommand\SATVALUE{\decisionname{SATVALUE}}
\newcommand\SUCCINTSAT{\decisionname{SUCCINTSAT}}
\newcommand\SUCCINTSATVALUE{\decisionname{SUCCINTSATVALUE}}
\newcommand\GEOGRAPHY{\decisionname{GEOGRAPHY}}
\newcommand\PARITY{\decisionname{PARITY}}
\newcommand\MAJ{\decisionname{MAJ}}
\newcommand\CVP{\CIRCUITVALUE}
\newcommand\BOOLEANPROGRAMVALUE{BOOLEAN~PROGRAM~VALUE}

\newcommand\CLIQUE{\cp{CLIQUE}}
\LANGUEinv{
\newcommand\RECOUVREMENTDESOMMETS{\cp{RECOUVREMENT~DE~SOMMETS}}
\newcommand\RS{\RECOUVREMENTDESOMMETS}
}{
\newcommand\RECOUVREMENTDESOMMETSeng{\cp{VERTEX~COVER}}
\newcommand\RSeng{\RECOUVREMENTDESOMMETSeng}
}

\LANGUEinv{\newcommand\COUPUREMAXIMALE{\cp{COUPURE~MAXIMALE}}}
{\newcommand\COUPUREMAXIMALEeng{\cp{MAXIMAL~CUT}}}
\newcommand\CIRCUITHAMILTONIEN{\CHAMILTONIEN}
\LANGUEinv{\newcommand\VOYAGEURDECOMMERCE{\VCOMMERCE}}
{\newcommand\VOYAGEURDECOMMERCEeng{\VCOMMERCEeng}}

\LANGUEinv{\newcommand\CIRCUITLEPLUSLONG{\cp{CIRCUIT~LE~PLUS~LONG}}}
{\newcommand\CIRCUITLEPLUSLONGeng{\cp{LONGEST~CIRCUIT}}}
\LANGUEinv{\newcommand\SOMMEDESOUSENSEMBLE{\SUBSETSUM}}{
\newcommand\SOMMEDESOUSENSEMBLEeng{\SUBSETSUMeng}}
\LANGUEinv{
\newcommand\SACADOS{\cp{SAC~A~DOS}}}
{\newcommand\SACADOSeng{\cp{KNAPSACK}}}

%français
\LANGUEinv{
\newcommand\CHAMILTONIEN{\cp{CIRCUIT~HAMILTONIEN}}}
{\newcommand\CHAMILTONIENeng{\cp{HAMILTONIAN~CIRCUIT}}}
\LANGUEinv{
\newcommand\VCOMMERCE{\cp{VOYAGEUR~DE~COMMERCE}}}
{\newcommand\VCOMMERCEeng{\cp{TRAVELING~SALESMAN}}}
\LANGUEinv{
\newcommand\kCOLORABILITE{\cp{k-COLORABILITE}}}
{\newcommand\kCOLORABILITEeng{\cp{k-COLORABILITY}}}
\LANGUEinv{
\newcommand\COLORIAGE{\cp{BON~COLORIAGE}}}
{\newcommand\COLORIAGEeng{\cp{GOOD~COLORING}}}
\LANGUEinv{
\newcommand\tCOLORABILITE{\cp{3-COLORABILITE}}}
{\newcommand\tCOLORABILITEeng{\cp{3-COLORABILITY}}}
\LANGUEinv{
\newcommand\SUBSETSUM{\cp{SOMME~DE~SOUS~ENSEMBLE}}}
{\newcommand\SUBSETSUMeng{\cp{SUBSET~SUM}}}
\newcommand\PARTITION{\cp{PARTITION}}
\LANGUEinv{
\newcommand\NOMBREPREMIER{\decisionname{NOMBRE~ PREMIER}}}
{\newcommand\NOMBREPREMIEReng{\decisionname{PRIME~NUMBER}}}
\LANGUEinv{
\newcommand\CODAGE{\decisionname{CODAGE}}}
{\newcommand\CODAGEeng{\decisionname{ENCODING}}}
\newcommand\kSAT{\cp{k-SAT}}
\newcommand\tSAT{\cp{3-SAT}}
\newcommand\NAESAT{\cp{NAESAT}}
\newcommand\NAEtSAT{\cp{NAE3SAT}}
\newcommand\NAEqSAT{\cp{NAE4SAT}}
\newcommand\STABLE{\cp{\LANGUE{INDEPENDANT~SET}{STABLE}}}

%\newcommand\CM{\cp{COUPURE MAXIMALE}}
\LANGUEinv{
\newcommand\POSTpb{\cp{Problème~de~ la~correspondance~de~Post}}}
{\newcommand\POSTpbeng{\cp{Post~correspondence~problem}}}

\LANGUEinv{
\newcommand\PARITE{\cp{PARITE}}}
{\newcommand\PARITEeng{\cp{PARITY}}}


%-------------------------
%--- Modèles parallèles  ---
%--------------------------


% PARALLELISME
\newcommand\RAM{\cp{RAM}}
%
\newcommand\PRAM{\cp{PRAM}}
\newcommand\CRCW{\cp{CRCW}}
\newcommand\CREW{\cp{CREW}}
\newcommand\EREW{\cp{EREW}}


%------------------------
%--- codages  ---
%------------------------


\newcommand\codage[1]{\langle #1 \rangle}
\newcommand\tuple[1]{\codage{#1}}
\newcommand\paire[2]{\codage{#1,#2}}
\newcommand\triplet[3]{\codage{#1,#2,#3}}
\newcommand\quadruplet[4]{\codage{#1,#2,#3,#4}}
\newcommand\cinquplet[5]{\codage{#1,#2,#3,#4,#5}}


%------------------------
%--- descriptive set theory  ---
%------------------------

\newcommand\baire{\mathcal{N}}
\newcommand\peano{\overline{\mathcal{N}}} %% A MOI
\newcommand\cantor{\mathcal{C}}
%
\newcommand\bSigma{\mathbf{\Sigma}}
\newcommand\bPi{\mathbf{\Pi}}
\newcommand\bDelta{\mathbf{\Delta}}
%


%%%%%% ABOUT NOMS QUI CHANGES

\newcommand\WF{\cp{WF}}
\newcommand\WO{\WF}
\newcommand\LO{\cp{LO}}
\newcommand\DLO{\cp{DLO}}
% 
\newcommand\rk{\cp{rank}}


%% Moins propre:

\newcommand\I{\mathbb{I}}
\renewcommand\H{\mathbb{H}}
%
\newcommand\leKB{<_{KB}}
%
\newcommand\finiomega{{<\omega}}
\newcommand\compl{^\frown}
\newcommand\prefix{ ~ prefix~  } 
\newcommand\X{{X}}
\newcommand\Y{\mathcal{Y}}
\newcommand\subsetstrict{\subset}

\renewcommand\complement[1]{#1^{c}}


%------------------------
%--- schémas de récursion ---
%------------------------

%%                      Source papier MFCS 2019


% COMPLEXITY

%% MFCS differe de Clote:
%% Version à taille restreinte

\newcommand\co{\cp{co}}

\newcommand\mFPTIME{\mF_{PTIME}}
\newcommand{\mFPSPACE}{\mF_{SPACE}}
\newcommand{\mFPH}{\mF_{PH}}

\newcommand\SigmaP{\Sigma^P}
\newcommand\DeltaP{\Delta^{P}}
\newcommand\PiP{\Pi^{P}}
\newcommand\coSigmaP{\co\SigmaP}
\newcommand\coPiP{\mF{\co\PiP}}

%% autres classes:


%% Schemas pour complexité implicite

\newcommand\COMP{\cp{ COMP}}
\newcommand\CRN{\cp{ CRN}}
\newcommand\BRN{\cp{ BRN}}
\newcommand\SBRN{\cp{ SBRN}}
\newcommand\BR{\cp{ BR}}
\newcommand\kBR{k-\cp{ BR}}
\newcommand \WBPR{\cp{ WBPR}}
\newcommand \BMIN{\cp{ BMIN}}
\newcommand \BSUM{\cp{ BSUM}}
\newcommand\SBSUM{\cp{ SBSUM}}
\newcommand \BPROD{\cp{ BPROD}}
\newcommand \SBPROD{\cp{ SBPROD}}
\newcommand \SCOMP{\cp{ SCOMP}}
\newcommand \SR{\cp{ SR}}
\newcommand \SRN{\cp{ SRN}}

%% devrait etre défini
\newcommand\WBRN{\cp{ WBRN}}



%-- godel
\newcommand\INTDIV{\cp{INTDIV}}
\newcommand\DIV{\cp{DIV}}
\newcommand\EVEN{\cp{EVEN}}
\newcommand\ODD{\cp{ODD}}
\newcommand\PRIME{\cp{PRIME}}
\newcommand\POWER{\cp{POWER}}
\newcommand\BIT{\cp{BIT}}
\newcommand\negHP{\overline{\cp{HP}}}
\newcommand\HP{\cp{HP}}
\newcommand\negLuniv{\overline{\Luniv}}
\newcommand\VALCOMP{\cp{VALCOMP}}
\newcommand\mWINDOW{\cp{LEGAL}}
\newcommand\mmWINDOW{\cp{WINDOW}}
\newcommand\LENGTH{\cp{LENGTH}}
\newcommand\DIGIT{\cp{DIGIT}}
\newcommand\MATCH{\cp{MATCH}}
\newcommand\MOVE{\cp{MOVE}}
\newcommand\START{\cp{START}}
\newcommand\CELL{\cp{CELL}}
\newcommand\HALT{\cp{HALT}}
\newcommand\SUBST{\cp{SUBST}}
\newcommand\PROOF{\cp{PROOF}}
\newcommand\PROVABLE{\cp{PROVABLE}}
\newcommand\CONSIST{\cp{CONSIST}}



%---------------------
%--- calculus (AP) ---
%---------------------

% \jacobian{f}{x}
\newcommand{\jacobian}[1]{J_{#1}}
%\newcommand{\grad}[1]{\overrightarrow{\operatorname{grad}}~{#1}}
\newcommand{\grad}[1]{\nabla{#1}}
\newcommand{\scalarprod}[2]{{#1}\cdot{#2}}
\newcommand{\bigO}[1]{\mathcal{O}\left(#1\right)}
\newcommand\smallO{o}
\newcommand{\softO}[1]{\tilde{\mathcal{O}}\left(#1\right)}
\newcommand{\taylor}[3]{{T_{#1}^{#2}#3}}
\newcommand{\transpose}[1]{{#1}^T}
\newcommand{\pastsup}[2]{{\sup}_{#1}#2}

\DeclareMathOperator\arctanh{arctanh}



%---------------
%--- norms (AP) ---
%---------------


\newcommand{\inorm}[2]{\left\lVert{#1}\right\rVert_{#2}}
\newcommand{\twonorm}[1]{\inorm{#1}{2}}
\newcommand{\infnorm}[1]{\inorm{#1}{}}
\newcommand{\normsup}[1]{\infnorm{#1}{}}
%
\newcommand\hausdorff{d_H}
%
\newcommand\norm[1]{\infnorm{#1}}


%----------------
%--- topology ---
%----------------

% \openball{pt}{radius}
 \newcommand{\openball}[2]{B_{#2}(#1)}
\newcommand{\closedball}[2]{\ovl{B}_{#2}(#1)}
\newcommand{\closure}[1]{\ovl{#1}}
% \newcommand{\interior}[1]{\mathring{#1}}
% \newcommand{\border}[1]{\partial{#1}}




%-----------------
%--- functions (AP) ---
%-----------------
 
\newcommand{\lambdafun}[2]{(#1)\mapsto{#2}}
\newcommand{\lambdafunex}[2]{#1\mapsto{#2}}
\newcommand{\argmin}{\operatornamewithlimits{argmin}}
\newcommand{\argmax}{\operatornamewithlimits{argmax}}


 
%-----------------------
%--- Machines de Turing  ---
%-----------------------

\newcommand\blanc{\vectorl{B}}



%-----------------
%--- Logique ---
%-----------------
 
\newcommand\sem[1]{{[\semspace[}#1{]\semspace]}}
\newcommand\sems[1]{{[\semspace[}#1{]\semspace]}}
\newcommand\semss[1]{\sem{#1}_{\sigma'}}
\newcommand\semspace{\kern -.25ex}
\newcommand\true{true}
\newcommand\false{false}
%extension élémentaire:
\newcommand\elem{\preceq}
\newcommand\godel[1]{\lceil #1 \rceil}

%-----------------
%--- Model theory ---
%-----------------

\newcommand\hull[1]{\langle #1 \rangle}

%-----------------
%--- Ordinaux ---
%-----------------

\newcommand\omegaunCK{\omega_1^{CK}}
\newcommand\omegaunck{\omegaunCK}

\DeclareMathOperator{\cof}{cof}

%-----------------
%--- Surréels (QG) ---
%-----------------

\newcommand{\On}{\textbf{On}}
\newcommand{\Nobf}{\textbf{No}}
\newcommand{\Nobb}{\ensuremath{\mathbbmss{No}}}
\newcommand{\Nolt}[1]{\ensuremath{\Nobf_{< #1}}}
\newcommand{\Noleq}[1]{\ensuremath{\Nobf_{\leq #1}}}
\newcommand{\Noltbb}[1]{\ensuremath{\Nobb_{< #1}}}
\providecommand\No{\Nobf}
\newcommand\Noalpha{\Nolt{\alpha}}
\newcommand\Nolambda{\Nolt{\lambda}}

\DeclareMathOperator{\length}{length}
\DeclareMathOperator{\supp}{supp} % operateur support


%----------------------
%---Maths  de base (AP)  ---
%----------------------

\DeclareMathOperator{\dom}{dom} % operateur support

\DeclareMathOperator{\size}{size}

\newcommand{\powerset}[1]{\mathcal{P}(#1)}
\newcommand{\card}[1]{{\##1}}
\newcommand\priv{ - }
\newcommand\prive{-}
\renewcommand{\restriction}{\mathord{\upharpoonright}}
\newcommand\restrict[1]{_{\restriction #1}}

% \providecommand\mod{}
%% UNDEFINE:
\let\mod\relax
\let\div\relax
\DeclareMathOperator{\mod}{mod} 
\DeclareMathOperator{\div}{div} 

% partieentier
\newcommand\entieres[1]{\lceil #1 \rceil}
\newcommand\entierei[1]{\lfloor #1 \rfloor}


%----------------------
%---Symbole danger ---
%----------------------


%\providecommand*{\TakeFourierOrnament}[1]{{%
%\fontencoding{U}\fontfamily{futs}\selectfont\char#1}}
%\providecommand*{\danger}{{\Huge \TakeFourierOrnament{66}}}



%%%%%%%%%%%%%%%%%%%%%%%%%%%%%%%%%%%%%%%%%%%%%%%%%
%
%
%                      ASMeries
%
%
%%%%%%%%%%%%%%%%%%%%%%%%%%%%%%%%%%%%%%%%%%%%%%%%%



% CIRCUITS
%-- particuliers
\newcommand\REPLACE{\cp{REPLACE}}
\newcommand\EQUAL{\cp{EQUAL}}
\newcommand\EVALUE{\cp{EVALUE}}
\newcommand\TVALUE{\cp{TVALUE}}
\newcommand\UPDATE{\cp{UPDATE}}
\newcommand\ASSIGN{\cp{ASSIGN}}
\newcommand\PAR{\cp{PAR}}
\newcommand\DO{\cp{DO}}


% ASM
\newcommand\emplacement[2]{(#1,#2)}
\newcommand\contenu[2]{\sem{(#1,#2)}}
\newcommand\affecte[3]{(#1,#2)\leftarrow#3}
\newcommand\ctl{ctl\_state}


%%%%%%%%%%%%%%%%%%%%%%%%%%%%%%%%%%%%%%%%%%%%%%%%%
%
%
%                      GPACqueries
%
%
%%%%%%%%%%%%%%%%%%%%%%%%%%%%%%%%%%%%%%%%%%%%%%%%%



%------------------------
%--- generable zoo (GPAC) ---
%------------------------

\newcommand{\myop}[1]{\operatorname{#1}}
\newcommand{\abs}{\myop{abs}}
\newcommand{\sg}{\myop{sg}}
\newcommand{\ip}[1]{\myop{ip}_{#1}}
\newcommand{\rnd}{\myop{rnd}}
\newcommand{\lil}{\myop{lil}}
\newcommand{\lxh}{\myop{lxh}}
\newcommand{\hxl}{\myop{hxl}}
\newcommand{\plil}{\myop{plil}}
\newcommand{\reach}{\myop{reach}}
%\newcommand{\mreach}{\myop{mreach}}
\newcommand{\watchdog}{\myop{watchdog}}
\newcommand{\monitor}{\myop{monitor}}
%\newcommand{\norm}{\myop{norm}}
\newcommand{\mx}{\myop{mx}}
\newcommand{\mn}{\myop{mn}}
\newcommand{\sample}{\myop{sample}}
\newcommand{\select}{\myop{select}}
\newcommand{\ssine}[1]{\sin_{#1}}
\newcommand{\clamp}{\myop{clamp}}

\newcommand{\fnsg}[3]{tanh((#1)*(#2)*(#3))}
\newcommand{\fnipone}[3]{(1+\fnsg{(#1)-1}{#2}{#3})/2.}
\newcommand{\fnipinf}[1]{#1-sin(2*pi*(#1))/2./pi}%
\newcommand{\fnlil}[5]{%
\fnipone{(#3)-(#1)+1-((#2)-(#1))/8.}{#4+log(1+4./(#2-#1)*(#5)*(#5))+log(1+4)}{8./(#2-#1)}%
*\fnipone{#2-#3+1-(#2-#1)/8.}{#4+log(1+4/(#2-#1)*(#5)*(#5))+log(1+4)}{8./(#2-#1)}%
*4./(#2-#1)*(#5)}
\newcommand{\fnlxh}[5]{\fnipone{(#3)-(#1+#2)/2.+1}{#4+log(1+(#5)*(#5))}{2./(#2-#1)}*(#5)}
\newcommand{\fnhxl}[5]{\fnipone{(#1+#2)/2.-(#3)+1}{#4+log(1+(#5)*(#5))}{2./(#2-#1)}*(#5)}
\newcommand{\fnplilf}[4]{sin(((#4)-((#1)+(#2))/2.+(#3)/4.)*2.*pi/(#3))}
\newcommand{\fnplil}[6]{(#6)*\fnlxh{\fnplilf{#1}{#2}{#3}{#1}}%
{\fnplilf{#1}{#2}{#3}{#1+((#2)-(#1))/4.}}%
{\fnplilf{#1}{#2}{#3}{#4}}%
{#5+2.+log(1+(#6)*(#6))}%
{1./4.+2./((#2)-(#1))}}
\newcommand{\fnssine}[2]{sin(#2)*(1-cos(#2)/cos(#1))}



%--------------------
%--- complexity GPAC ---
%--------------------


%\ifthenelse{\equal{#1}{}}{Empty.}{Nonempty.}
\newcommand{\myclass}[1]{\operatorname{#1}}
% \newcommand{\gcomp}[2]{\ensuremath{\myclass{GCOMP}(#1,#2)}}
% \newcommand{\ggenerable}[1]{\textcolor{red}{\ensuremath{#1-}\text{generable}}}
 \newcommand{\gval}[2][]{\ensuremath{\myclass{GVAL}_{#1}\ifthenelse{\equal{#2}{}}{}{[#2]}}}
 \newcommand{\gpval}[1][]{\ensuremath{\myclass{GPVAL}_{#1}}}
 \newcommand{\gc}[3][]{\ensuremath{\myclass{ATSC}(#2,#3)}}
 \newcommand{\gpc}[1][]{\ensuremath{\myclass{ATSP}}}
\newcommand{\cgc}[1][]{\ensuremath{\myclass{ATS}}}
\newcommand{\gexpc}[1][]{\ensuremath{\myclass{AEXP}}}
\newcommand{\gwc}[2]{\ensuremath{\myclass{AW}(#1,#2)}}
\newcommand{\gpwc}{\ensuremath{\myclass{AWP}}}
\newcommand{\cgwc}{\ensuremath{\myclass{AW}}}
\newcommand{\grc}[3]{\ensuremath{\myclass{AR}(#1,#2,#3)}}
\newcommand{\gprc}{\ensuremath{\myclass{ARP}}}
\newcommand{\gsc}[3]{\ensuremath{\myclass{AS}(#1,#2,#3)}}
\newcommand{\gpsc}{\ensuremath{\myclass{ASP}}}
\newcommand{\goc}[3]{\ensuremath{\myclass{AO}(#1,#2,#3)}}
\newcommand{\gpoc}{\ensuremath{\myclass{AOP}}}
\newcommand{\cgoc}{\ensuremath{\myclass{AO}}}
\newcommand{\guc}[4]{\ensuremath{\myclass{AX}(#1,#2,#3,#4)}}
\newcommand{\gpuc}{\ensuremath{\myclass{AXP}}}
\newcommand{\glc}[2][]{\ensuremath{\myclass{AL}(#2)}}
\newcommand{\gplc}[1][]{\ensuremath{\myclass{ALP}}}
\newcommand{\cglc}[1][]{\ensuremath{\myclass{AL}}}

%\newcommand{\intinterv}[2]{\{#1,\ldots,#2\}}
\newcommand{\intinterv}[2]{\llbracket#1,#2\rrbracket}

\newcommand{\Rgen}{\R_{G}}
\newcommand{\Rpoly}{\R_{P}}

%%%%%%%%%%%%%%%%%%%%%%%%%%%%%%%%%%%%%%%%%%%%%%%%%
%
%
%                      Mes codages obscures
%
%
%%%%%%%%%%%%%%%%%%%%%%%%%%%%%%%%%%%%%%%%%%%%%%%%%




%%% Codages

% \newcommand\gammaoc{\gamma_{[0,1]}}
% \newcommand\gammaonc{\gamma_{\N}}
\newcommand\gammaTMc{\gamma^{TM}_{[0,1]}}
\newcommand\gammaTMcatan{\gamma^{TM}_{(0,1)^2}}
\newcommand\gammaTMnc{\gamma^{TM}_{\N^3}}
\newcommand\gammaTMncd{\gamma^{TM}_{\N^2}}
% \newcommand\gammaTMe{\gamma^{TM}_{*}}
% \newcommand\gammas{\gamma_{sab}}
% \newcommand\gammaord{\gamma_{cantor}}


\newcommand\enccompact{\nu_{X}}
 \newcommand\encinfinite{\nu_\N}
 \newcommand\enc{\nu}
 
%%%%%%%%%%%%%%%%%%%%%%%%%%%%%%%%%%%%%%%%%%%%%%%%%
%
%
%                      DESSINS ET IMAGES
%
%
%%%%%%%%%%%%%%%%%%%%%%%%%%%%%%%%%%%%%%%%%%%%%%%%%


\newcommand\IMAGE[2]{ \begin{center} 
    \includegraphics[#1]{#2}
  \end{center}
  }


  
%%%% POUR FAIRE DESSIN.
\newenvironment{dessinpp}[1]{
\begin{tikzpicture}[baseline=(current bounding
      box.west),#1]
}
{\end{tikzpicture}}




\newenvironment{dessinp}[1]{
\begin{center}\begin{tikzpicture}[baseline=(current bounding
      box.north),#1]
}
{\end{tikzpicture}
\end{center}}


\newenvironment{dessin}{
\begin{center}\begin{tikzpicture}[baseline=(current bounding
      box.north)]
}
{\end{tikzpicture}
\end{center}}



\newenvironment{dessind}{
\begin{tikzpicture}[baseline=(current bounding
      box.north)]
}
{\end{tikzpicture}
}



%%%%%%%%%%%%%%%%%%%%%%%%%%%%%%%%%%%%%%%%%%%%%%%%%
%
%
%                      POUR DESSINS GPAC
%
%
%%%%%%%%%%%%%%%%%%%%%%%%%%%%%%%%%%%%%%%%%%%%%%%%%


%%%% MACROS DE BASE.


\newcommand\constant[2]
{
node (#1)  [shape=circle,draw] {#2};
\draw (#1.east) -- +(0.3,0) coordinate (#1-o) ; 
\draw (#1.west) -- +(-0.3,0) coordinate (#1-i) ; 
}

\newcommand\basicproduct[2]
{
node (#1) {#2};
\draw (#1) +(-0.4,0.4) -- +(0.4,0.4) -- +(0.8,0) -- +(0.4,-0.4) -- +(-0.4,-0.4) --
+(-0.4,0.4);
\draw (#1) ++(0.8,0) -- ++(+0.3,0) coordinate (#1-o) ;
\draw (#1) ++ (-0.4,0.2) -- +(-0.3,0) coordinate (#1-i1) ;
\draw (#1) ++(-0.4,-0.2) -- +(-0.3,0) coordinate (#1-i2) ;
}

\newcommand\product[1]{\basicproduct{#1}{$+\Pi$}}
\newcommand\negativeproduct[1]{\basicproduct{#1}{$-\Pi$}}


%%BEGIN INTERNAL
\newcommand\basicsummer[1]{
coordinate (#1) {};
\draw (#1) +(-0.5,0.6) -- +(0.5,0) -- +(-0.5,-0.6) -- +(-0.5,0.6);
\draw (#1) ++(0.5,0) -- +(+0.3,0) coordinate (#1-o) ;
}

\newcommand\basicintegrator[1]{
coordinate (#1) {};
\draw (#1) +(-0.5,0.6) -- +(0.5,0) -- +(-0.5,-0.6) -- +(-0.5,0.6);
\draw (#1) ++(0.5,0) -- +(+0.3,0) coordinate (#1-o) ;
\draw (#1) +(-0.5,-0.6) -| +(-0.8,0.6) -- +(-0.5,0.6) ;
 \draw (#1) ++(-0.65,0.6) -- +(0,+0.3) coordinate (#1-init) ;
}

\newcommand\splitfour[2]{
\draw (#2) ++ (#1,0.4) -- +(-0.3,0) coordinate (#2-i1) ;
\draw (#2) ++ (#1,0.15) -- +(-0.3,0) coordinate (#2-i2) ;
\draw (#2) ++ (#1,-0.15) -- +(-0.3,0) coordinate (#2-i3) ;
\draw (#2) ++ (#1,-0.4) -- +(-0.3,0) coordinate (#2-i4) ;
}

\newcommand\splitthree[2]{
\draw  (#2) ++(#1,0.3) -- +(-0.3,0) coordinate (#2-i1) ;
\draw  (#2) ++(#1,0) -- +(-0.3,0) coordinate (#2-i2) ;
\draw  (#2) ++(#1,-0.3) -- +(-0.3,0) coordinate (#2-i3) ;
}

\newcommand\splittwo[2]{
\draw  (#2) ++(#1,0.2) -- +(-0.3,0) coordinate (#2-i1) ;
\draw  (#2) ++(#1,-0.2) -- +(-0.3,0) coordinate (#2-i2) ;
}

\newcommand\splitone[2]{
\draw (#2) ++(#1,0) -- +(-0.3,0) coordinate (#2-i1) ;
}

%% END INTERNAL

\newcommand\summerfour[1]{
\basicsummer{#1}
\splitfour{-0.5}{#1}
}

\newcommand\summerthree[1]{
\basicsummer{#1}
\splitthree{-0.5}{#1}
}

\newcommand\summertwo[1]{
\basicsummer{#1}
\splittwo{-0.5}{#1}
}

\newcommand\summerone[1]{
\basicsummer{#1}
\splitone{-0.5}{#1}
}


\newcommand\integratorfour[1]{
\basicintegrator{#1}
\splitfour{-0.8}{#1}
<}


\newcommand\integratorthree[1]{
\basicintegrator{#1}
\splitthree{-0.8}{#1}
}


\newcommand\integratortwo[1]{
\basicintegrator{#1}
\splittwo{-0.8}{#1}
}


\newcommand\integratorone[1]{
\basicintegrator{#1}
\splitone{-0.8}{#1}
}

%%%%% FIN MACROS




%%%%%%%%%%%%%%%%%%%%%%%%%%%%%%%%%%%%%%%%%%%%%%%%%
%
%
%                      Trucs d'Amaury
%
%
%%%%%%%%%%%%%%%%%%%%%%%%%%%%%%%%%%%%%%%%%%%%%%%%%


%----------------
%--- ref (AP) ---
%----------------

% \newcommand{\defref}[1]{Definition~\ref{#1}}
\newcommand{\lemref}[1]{Lemma~\ref{#1}}
\newcommand{\thref}[1]{Theorem~\ref{#1}}
\newcommand{\amaurycorref}[1]{Corollary~\ref{#1}}
 \newcommand{\figref}[1]{Figure~\ref{#1}}
% \newcommand{\figrefmult}[1]{Figures~{#1}}
% \newcommand{\secref}[1]{Section~\ref{#1}}
% %\renewcommand{\algref}[1]{Algorithm~\ref{#1}}
 \newcommand{\propref}[1]{Proposition~\ref{#1}}
% \newcommand{\exref}[1]{Example~\ref{#1}}
\newcommand{\remref}[1]{Remark~\ref{#1}}


 %-------------------
%--- polynomials ---
%-------------------

\newcommand{\sigmap}[1]{{\Sigma{#1}}}
\newcommand{\degp}[1]{{\operatorname{deg}(#1)}}
\newcommand{\poly}{\operatorname{poly}}


 %----------------------
%--- misc functions ---
%----------------------

\newcommand{\sgn}[1]{\operatorname{sgn}(#1)}
\newcommand{\intp}{\operatorname{int}}
 \newcommand{\fracp}{\operatorname{frac}}
\newcommand{\fiter}[2]{#1^{[#2]}}
\newcommand{\frestrict}[2]{#1\restriction_{#2}}
\newcommand{\indicator}[1]{\mathds{1}_{#1}}
\newcommand{\idfun}{\operatorname{id}}
% \newcommand{\floor}[1]{\left\lfloor#1\right\rfloor}
 \newcommand{\ceil}[1]{\left\lceil#1\right\rceil}
\newcommand{\round}[1]{\left\lfloor#1\right\rceil}
\DeclareMathOperator\erf{erf}


%------------
%--- sets ---
%------------

\newcommand{\A}{\mathbb{A}}
%\newcommand{\D}{\mathbb{D}}
\newcommand{\Ns}{\mathbb{N}^*}
%\newcommand{\C}{\mathbb{C}}
%\newcommand{\K}{\mathbb{K}}
% \MAT{height}{[width]}{[field]}
\newcommand{\MAT}[3]{M_{#1\ifthenelse{\equal{#2}{}}{}{,#2}}\ifthenelse{\equal{#3}{}}{}{\left(#3\right)}}

%\newcommand{\Rcomp}{\R_{c}}
\newcommand{\Rp}{\R_{+}}
\newcommand{\Rs}{\R^{*}}
\newcommand{\Rps}{\R_{+}^{*}}

%\newcommand{\Analytic}{C^\omega}
%\newcommand{\PTIME}{\ensuremath{\operatorname{P}}}
%\newcommand{\EXPTIME}{\ensuremath{\operatorname{EXPTIME}}}
%\newcommand{\NP}{\ensuremath{\operatorname{NP}}}
%\newcommand{\NPTIME}{\NP}
\newcommand{\FP}{\ensuremath{\operatorname{FP}}}
%\newcommand{\PSPACE}{\ensuremath{\operatorname{PSPACE}}}
% \newcommand{\FPSPACE}{\ensuremath{\operatorname{FPSPACE}}}



%%%%%%%%%%%%%%%%%%%%%%%%%%%%%%%%%%%%%%%%%%%%%%%%%
%
%
%                      Trucs de Quentin
%
%
%%%%%%%%%%%%%%%%%%%%%%%%%%%%%%%%%%%%%%%%%%%%%%%%%


\newcommand{\fct}[5]{\ensuremath{#1 : \left\{\begin{array}{c c c}
#2 & \rightarrow & #3\\
#4 & \mapsto & #5
                                             \end{array}\right.}}

%Nouvelle fraction, plus jolie
\newcommand{\f}[2]{	
	\mathchoice%
		{\dfrac{#1}{#2}}
    	{\dfrac{#1}{#2}}
		{\frac{#1}{#2}}
		{\frac{#1}{#2}}
}


%%Ensembles et combinatoires

%Ensemble avec propriété à drotie \{x | blabla\}
\newcommand{\enstq}[2]{\left\{ \left.#1\ \vphantom{#2}\right|\ #2  \right\}}
%meme chose mais avec des crochets
\newcommand{\crotq}[2]{\left[ \left.#1\ \vphantom{#2}\right|\ #2  \right]}
%même chose mais avec des brackets
\newcommand{\bratq}[2]{\left\langle \left.#1\ \vphantom{#2}\right|\ #2  
  \right\rangle}

\newcommand{\intn}[2]{\ensuremath{
		\left\llbracket\, #1\ ;\ #2\,\right\rrbracket
              }}


%%%%%%%%%%%%%%%%%%%%%%%%%%%%%%%%%%%%%%%%%%%%%%%%%
%
%
%               ENVIRONNEMENTS THEOREMS & CO
%
%
%%%%%%%%%%%%%%%%%%%%%%%%%%%%%%%%%%%%%%%%%%%%%%%%%


%----------------
%--- pas dans lipics ---
%-----------------


\ifnotSLIDE{
% commente car bug Fevrier 2026: \ifnotTDEXAM{\spnewtheorem{solution}{Solution}{\bfseries}{\rmfamily}}
\spnewtheorem{openproblem}{Open Problem}{\bfseries}{\rmfamily}
\spnewtheorem{remarkolivier}{Commentaire d'Olivier: }{\bfseries}{\rmfamily}
\spnewtheorem{postulate}{Postulate}{\bfseries}{\rmfamily}
}

%----------------
%--- dans lipics ---
%-----------------
%
\newenvironment{exercice}[1]{\begin{exercise} \label{exo:#1-stt}}{\end{exercise}}
\newenvironment{exerciced}[1]{\begin{exercisee}\label{exo:#1-stt}}{\end{exercisee}}
\newenvironment{exercicec}[1]{\begin{exercise} \label{exo:#1-stt} (\LANGUE{solution on}{corrigé} page \pageref{exo:#1-cor})
}{\end{exercise}}
\newenvironment{exercicedc}[1]{\begin{exercisee} \label{exo:#1-stt} (\LANGUE{solution on}{corrigé} page \pageref{exo:#1-cor})
}{\end{exercisee}}

\newenvironment{slideproposition}{{\bf Proposition}}{}

\ifnotSLIDE{
\nolipics{
\nolncs{\spnewtheorem{theorem}{\LANGUE{Theorem}{Théorème}}{\bfseries}{\rmfamily}}
\nolncs{\spnewtheoremi{definition}{\LANGUE{Definition}{Définition}}{\bfseries}{\rmfamily}{theorem}}
\nolncs{\spnewtheoremi{lemma}{\LANGUE{Lemma}{Lemme}}{\bfseries}{\rmfamily}{theorem}}
\nolncs{\spnewtheoremi{corollary}{\LANGUE{Corollary}{Corollaire}}{\bfseries}{\rmfamily}{theorem}}
\nolncs{\spnewtheoremi{proposition}{Proposition}{\bfseries}{\rmfamily}{theorem}}
\nolncs{\spnewtheoremi{example}{\LANGUE{Example}{Exemple}}{\bfseries}{\rmfamily}{theorem}}
%\spnewtheorem{sketch}{{\bf \LANGUE{Sketch of proof}{Démonstration (principe)}}}{\bfseries}{\rmfamily}
\nolncs{\spnewtheoremi{remark}{\LANGUE{Remark}{Remarque}}{\bfseries}{\rmfamily}{theorem}}
\nolncs{\spnewtheorem{exercise}{\LANGUE{Exercise}{Exercice}}{\bfseries}{\rmfamily}}
\spnewtheorem{exercisee}{\LANGUE{*Exercise}{*Exercice}}{\bfseries}{\rmfamily}
\spnewtheorem{summary}{\LANGUE{Summary}{Résumé}}{\bfseries}{\rmfamily}
 %\spnewtheorem{exerciceenglish}{Exercise}{\bfseries}{\rmfamily}
 %
 \newenvironment{sketch}{{\bf \LANGUE{Sketch of proof}{Démonstration (principe)}}:}{\hfill $\Box$ }
\nolncs{\spnewtheorem{conjecture}{\LANGUE{Conjecture}{Conjecture}}{\bfseries}{\rmfamily}}
 %%
 \spnewtheorem{theoremaprouver}{\LANGUE{Theorem TO BE PROVED}{Théorème (A PROUVER)}}{\bfseries}{\rmfamily}
  \spnewtheorem{conjectureaprouver}{\LANGUE{Conjecture TO BE VERIFIED}{Conjecture (A VERIFIER)}}{\bfseries}{\rmfamily}
 }}
            
 \makeatletter
 \ifnotSLIDE{
  \ifnotTDEXAM{
\@ifundefined{proof}{
	\newenvironment{proof}{{\bf \LANGUE{Proof}{Démonstration}}:}{\hfill $\Box$ 
	}{}
	}}}
\makeatother

 
%%%%%%%%%%%%%%%%%%%%%%%%%%%%%%%%%%%%%%%%%%%%%%%%%
%
%
%                      Macros dont pas fier du tout
%
%
%%%%%%%%%%%%%%%%%%%%%%%%%%%%%%%%%%%%%%%%%%%%%%%%%



\newcommand\equations[1]{
\begin{array}{lll}
#1
\end{array}
}



\newcommand\systeme[1]{
\left\{
\begin{array}{lll} 
#1
\end{array}
\right.
}



%------------------
%--- shorthands (AP)---
%------------------
\newcommand{\mtt}[1]{\mathtt{#1}}
\newcommand{\ovl}[1]{\overline{#1}}
\newcommand{\panic}[1]{\textcolor{red}{#1}}
%\NewEnviron{epanic}{{\color{red}\BODY}}
\newcommand{\idest}{\emph{i.e.\thinspace}}


%------------------
%--- mes trucs  ---
%------------------


\newcommand\implique{\Rightarrow}
\newcommand\ssi{\Leftrightarrow}
\newcommand\ac[1]{\{#1\}}
\newcommand\zero{\vectorl{0}}
\newcommand\un{\vectorl{1}}

\newcommand\technique[1]{{\scriptsize #1}}


%% Preuve and Co

\newcommand\sensdirect{ Direction $\Rightarrow$. }
\newcommand\sensindirect{ Direction $\Leftarrow$. }




\newcommand\donnee{\mbox{<donnée>}}



% pour aller vite

\newcommand\gM{\mathfrak{M}}
\newcommand\gMM{(M,f_1,\cdots,f_u,r_1,\cdots,r_v)}
%
\newcommand\mygM{\mK}
\newcommand\mygMM{\left ( \K, \{op_i\}_{i\in I}, r_1 \ldots,
r_\ell, \zero,\un \right )}
\newcommand\myM{\K} 

\newcommand\gMMM{(M,f_1,\cdots,f_u,c_1,\cdots,c_k,r_1,\cdots,r_v)}
\newcommand\gMMMM{(f_1,\cdots,f_u,r_1,\cdots,r_v)}
\newcommand\gMME{(M,f_1,\cdots,f_u,f'_1,\cdots,f'_\ell,r_1,\cdots,r_v)}
\newcommand\gMMS{(f_1,\cdots,f_u,f'_1,\cdots,f'_\ell,r_1,\cdots,r_v)}

\newcommand\bool{\mathfrak{B}}
\newcommand\gbool{(\{0,1\},\zero,\un,=)}

\newcommand\osg{\overline{sg}}
\newcommand\moins{\mathrel{\dot{-}}}


%%% TRUC DE BARWISE


\newcommand\UM{\kU_{\kM}}
\newcommand\VV{\mathbb{V}}
% Hereditarily finite
\newcommand\IHF{\mathbb{I}HF}
\newcommand\IHP{\IHYP}
\newcommand\IHYP{\mathbb{I}HYP}

\providecommand\citep[1]{\cite{#1}}

%% Cofinali


% probas
% défini dans amsmath
\renewcommand\Pr{\cp{Pr}}
\newcommand\Var{\cp{Var}}


% -- mesures
%\newcommand\card{card}
\newcommand\taille{\LANGUE{size}{taille}}
\newcommand\entree{e}
\newcommand\profondeur{\LANGUE{depth}{profondeur}}


%%%%%%%%%%%%%%%%%%%%%%%%%%%%%%%%%%%%%%%%%%%%%%%%%%%%%%%%%%%%%%%%%%%%%%%%%%%%
%                 MACHINES DE TURING (graphique)
%%%%%%%%%%%%%%%%%%%%%%%%%%%%%%%%%%%%%%%%%%%%%%%%%%%%%%%%%%%%%%%%%%%%%%%%%%%%



% Configuration de machine de Turing
\newcommand\configTuring[3]{#1\textbf{#2}#3}


%%% DESSIN CROIX PORU AIDER  A DESSINER

\newcommand\croix{  +(-1,-1) -- +(1,1) +(-1,1) --
  +(1,-1) }

%% DESSIN ACCOLADE

\newcommand\ACCOLADE[3]{
%#1: coordonnées de départ
%#2: coordonnées arrivée
%#3: texte
\draw[decorate,decoration={brace}]  (#1) -- (#2)
 node[midway,above=0.1cm] {#3};
}

%%%


\def\lcase{3.5 ex}
\def\hcase{3.5ex}
\tikzstyle{case} = []
%[draw,minimum width = \lcase,minimum height = 3.6ex,baseline]
\newcommand\dcase{ +(-0.5*\lcase,-0.5*\hcase) rectangle +(0.5*\lcase,0.5*\hcase)}
%\newcommand\lettre[1]{\node [name=l,case,on chain=1] {#1};}


\newcommand\RUBANMACHINEDETURING[6]{
%#1: texte
%#2: position de la tête
%#3: nb de blancs à gauche
%#4: nb de cases totales
%#5: texte au dessus du ruban
%#6: position du coin
    \draw (#6) +(\lcase,4ex ) node {#5};
    \draw (#6) +(-0.6,0)  node  {\ldots};  
    \draw (#6) node[coordinate] (t) {};
    \foreach \x in {1,2,...,#3} {
         \draw (t) node [case]   {$\blanc$} \dcase;
         \draw (t) ++(\lcase,0) node [coordinate] (t) {};
       };
   \StrLen{#1}[\Resultat] % ATTENTION IL FAUT CETTE SYNTAXE POUR FAIRE
                         % EQUIVALENT D UN EDEF
  
\foreach \x in {1,...,\Resultat} 
             { 
        \draw (t) node [case]   {\StrChar{#1}{\x}} \dcase;
        \draw (t) ++(\lcase,0) node [coordinate] (t) {};
             }
\pgfmathtruncatemacro{\z}{#4 - \Resultat - #3}

  \foreach \x in {1,2,..., \z} {
        \draw (t) node [case]   {$\blanc$} \dcase;
        \draw (t) ++(\lcase,0) node [coordinate] (t) {};
}
%% Convention: graphiquement, la case numéro n se trouve à (n*\lcase,0)
%% en comptant les case à partir de $0$

   \draw  (t) +(0.1,0) node [name=r] {\ldots}; 
}

\newcommand\RUBANMACHINEDETURINGtexte[6]{
%#1: texte
%#2: position de la tête
%#3: nb de blancs à gauche
%#4: nb de cases totales (équivalent)
%#5: texte au dessus du ruban
%#6: position du coin
    \draw (#6) +(\lcase,4ex ) node {#5};
    \draw (#6) +(-0.6,0)  node  {\ldots};  
    \draw (#6) node[coordinate] (t) {};
   \foreach \x in {1,2,...,#3} {
        \draw (t) node [case]   {$\blanc$} \dcase;
        \draw (t) ++(\lcase,0) node [coordinate] (t) {};
      };
   \draw (#6) ++(-2*\lcase,0.5*\lcase)-- +(\lcase*#4,0);
   \draw (#6) ++(-2*\lcase,-0.5*\lcase) -- +(\lcase*#4,0);
   
    \draw (t) ++(-0.5*\lcase,0) node [case,anchor=west]  (te) {#1};
    \draw (te.east)  node [anchor=west,coordinate] (t) {};

  \draw (t) ++(0.2*\lcase,0) node [coordinate] (t) {};
   \foreach \x in {1,2,...,#3} {
        \draw (t) node [case]   {$\blanc$} \dcase;
        \draw (t) ++(\lcase,0) node [coordinate] (t) {};
      };

    \draw  (t) +(0.1,0) node [anchor=west] {\ldots}; 
}


\newcommand\RUBANMACHINEDETURINGTETE[6]{
%
% dessine la tete
%
% output: (tete)
%
\RUBANMACHINEDETURING{#1}{#2}{#3}{#4}{#5}{#6}
% TETE
     \draw (#6) + ( #2 * \lcase + #3  *\lcase,- 0.5*\lcase )
     node[coordinate] (k) {};
     \draw[->] (k) +(0,-0.5) -- (k);
     \draw (k) +(0,-0.5) node[coordinate] (tete) {};
}


\newcommand\RUBANMACHINEDETURINGTETEtexte[6]{
%
% dessine la tete
%
% output: (tete)
%
\RUBANMACHINEDETURINGtexte{#1}{#2}{#3}{#4}{#5}{#6}
% TETE
     \draw (#6) + ( #2 * \lcase + #3  *\lcase,- 0.5*\lcase )
     node[coordinate] (k) {};
     \draw[->] (k) +(0,-0.5) -- (k);
     \draw (k) +(0,-0.5) node[coordinate] (tete) {};
}


\newcommand\MACHINEDETURING[6]{
%#1: texte
%#2: position de la tête
%#3: état
%#4: nb de blancs à gauche
%#5: nb de cases totales
%#6: texte en haut du ruban
%\begin{tikzpicture}[node distance=-0.15mm]

\RUBANMACHINEDETURINGTETE{#1}{#2}{#4}{#5}{#6}{0,0}
%
%% ETAT:
     \draw (tete) node [name=etat,draw,circle,anchor=north]  {#3};
     \draw (etat) -- (tete);
%\end{tikzpicture}
}

\newcommand\MACHINEDETURINGtexte[6]{
%#1: texte
%#2: position de la tête
%#3: état
%#4: nb de blancs à gauche
%#5: nb de cases totales
%#6: texte en haut du ruban
%\begin{tikzpicture}[node distance=-0.15mm]

\RUBANMACHINEDETURINGTETEtexte{#1}{#2}{#4}{#5}{#6}{0,0}
%
%% ETAT:
     \draw (tete) node [name=etat,draw,circle,anchor=north]  {#3};
     \draw (etat) -- (tete);
%\end{tikzpicture}
}

\newcommand\MACHINEDETURINGTROISRUBANS[9]{
%#1: texte1
%#2: position de la tête 1
%#3: texte2
%#4: position de la tête 2
%#5: texte3
%#6: position de la tête 3
%#7: état
%#8: nb de blancs à gauche 1
%#9: nb de cases totales
%\begin{tikzpicture}[node distance=-0.15mm]

\RUBANMACHINEDETURINGTETE{#1}{#2}{#8}{#9}{\LANGUE{tape}{ruban} $1$}{0,0}
\draw (tete) node[coordinate] (tete1) {};
\draw (tete1) +(8,0) node[coordinate] (tete1b) {};
\draw (tete1) -| (tete1b) ;

\RUBANMACHINEDETURINGTETE{#3}{#4}{#8}{#9}{\LANGUE{tape}{ruban} $2$}{0,-2}
\draw (tete) node[coordinate] (tete2) {};
\draw (tete2) +(-8.5,0) node[coordinate] (tete2b) {};
\draw (tete2) -| (tete2b) ;

\RUBANMACHINEDETURINGTETE{#5}{#6}{#8}{#9}{\LANGUE{tape}{ruban} $3$}{0,-4}
\draw (tete) node[coordinate] (tete3) {};

%% ETAT:
     \draw (tete3) node [name=etat,draw,circle,anchor=north]  {#7};
     \draw (etat) -- (tete3);
    \draw  (etat) -|  (tete2b);
\draw  (etat) -|  (tete1b);
%\end{tikzpicture}
}

\newcommand\MACHINEDETURINGTROISRUBANSENGLISH[9]{
%#1: texte1
%#2: position de la tête 1
%#3: texte2
%#4: position de la tête 2
%#5: texte3
%#6: position de la tête 3
%#7: état
%#8: nb de blancs à gauche 1
%#9: nb de cases totales
%\begin{tikzpicture}[node distance=-0.15mm]

\RUBANMACHINEDETURINGTETE{#1}{#2}{#8}{#9}{tape $1$}{0,0}
\draw (tete) node[coordinate] (tete1) {};
\draw (tete1) +(8,0) node[coordinate] (tete1b) {};
\draw (tete1) -| (tete1b) ;

\RUBANMACHINEDETURINGTETE{#3}{#4}{#8}{#9}{tape $2$}{0,-2}
\draw (tete) node[coordinate] (tete2) {};
\draw (tete2) +(-8.5,0) node[coordinate] (tete2b) {};
\draw (tete2) -| (tete2b) ;

\RUBANMACHINEDETURINGTETE{#5}{#6}{#8}{#9}{tape $3$}{0,-4}
\draw (tete) node[coordinate] (tete3) {};

%% ETAT:
     \draw (tete3) node [name=etat,draw,circle,anchor=north]  {#7};
     \draw (etat) -- (tete3);
    \draw  (etat) -|  (tete2b);
\draw  (etat) -|  (tete1b);
%\end{tikzpicture}
}



\newcommand\MACHINEDETURINGTROISRUBANStexte[9]{
%#1: texte1
%#2: position de la tête 1
%#3: texte2
%#4: position de la tête 2
%#5: texte3
%#6: position de la tête 3
%#7: état
%#8: nb de blancs à gauche 1
%#9: nb de cases totales
% \begin{tikzpicture}[node distance=-0.15mm]

\RUBANMACHINEDETURINGTETEtexte{#1}{#2}{#8}{#9}{\LANGUE{tape}{ruban} $1$}{0,0}
\draw (tete) node[coordinate] (tete1) {};
\draw (tete1) +(8,0) node[coordinate] (tete1b) {};
\draw (tete1) -| (tete1b) ;

\RUBANMACHINEDETURINGTETEtexte{#3}{#4}{#8}{#9}{\LANGUE{tape}{ruban} $2$}{0,-2}
\draw (tete) node[coordinate] (tete2) {};
\draw (tete2) +(-8.5,0) node[coordinate] (tete2b) {};
\draw (tete2) -| (tete2b) ;

\RUBANMACHINEDETURINGTETEtexte{#5}{#6}{#8}{#9}{\LANGUE{tape}{ruban} $3$}{0,-4}
\draw (tete) node[coordinate] (tete3) {};

%% ETAT:
     \draw (tete3) node [name=etat,draw,circle,anchor=north]  {#7};
     \draw (etat) -- (tete3);
    \draw  (etat) -|  (tete2b);
\draw  (etat) -|  (tete1b);
%\end{tikzpicture}
}


\newcommand\MACHINEDETURINGDEUXRUBANStexte[9]{
%#1: texte1  
%#2: position de la tête 1 
%#3: texte2 inutilisée
%#4: position de la tête 2 inutilisée
%#5: texte3
%#6: position de la tête 3
%#7: état
%#8: nb de blancs à gauche 1
%#9: nb de cases totales
% \begin{tikzpicture}[node distance=-0.15mm]

\RUBANMACHINEDETURINGTETEtexte{#1}{#2}{#8}{#9}{\LANGUE{tape}{ruban} $1$}{0,0}
\draw (tete) node[coordinate] (tete1) {};
\draw (tete1) +(8,0) node[coordinate] (tete1b) {};
\draw (tete1) -| (tete1b) ;

% \RUBANMACHINEDETURINGTETEtexte{#3}{#4}{#8}{#9}{\LANGUE{tape}{ruban} $2$}{0,-1.5}
% \draw (tete) node[coordinate] (tete2) {};
% \draw (tete2) +(-5.5,0) node[coordinate] (tete2b) {};
% \draw (tete2) -| (tete2b) ;

\RUBANMACHINEDETURINGTETEtexte{#5}{#6}{#8}{#9}{\LANGUE{tape}{ruban} $2$}{0,-1.5}
\draw (tete) node[coordinate] (tete3) {};

%% ETAT:
     \draw (tete3) node [name=etat,draw,circle,anchor=north]  {#7};
     \draw (etat) -- (tete3);
%   \draw  (etat) -|  (tete2b);
\draw  (etat) -|  (tete1b);
%\end{tikzpicture}
}


\def\argdix{2}
\def\argonze{25}
\newcommand\MACHINEDETURINGQUATRERUBANStextep[9]{
%#1: texte1
%#2: position de la tête 1
%#3: texte2
%#4: position de la tête 2
%#5: texte3
%#6: position de la tête 3
%#7: texte4
%#8: position de la tête 3
%#9: état
%#10: nb de blancs à gauche 1
%#11: nb de cases totales
% \begin{tikzpicture}[node distance=-0.15mm]

%% BUG sur #10 devenu \argdix
%% BUG SUR #11 devenu \argonze

\RUBANMACHINEDETURINGTETEtexte{#1}{#2}{\argdix}{\argonze}{\LANGUE{tape}{ruban} $1$}{0,0}
% \draw (tete) node[coordinate] (tete1) {};
% \draw (tete1) +(8,0) node[coordinate] (tete1b) {};
% \draw (tete1) -| (tete1b) ;

\RUBANMACHINEDETURINGTETEtexte{#3}{#4}{\argdix}{\argonze}{\LANGUE{tape}{ruban} $2$}{0,-1.5}
% \draw (tete) node[coordinate] (tete2) {};
% \draw (tete2) +(-5.5,0) node[coordinate] (tete2b) {};
% \draw (tete2) -| (tete2b) ;

\RUBANMACHINEDETURINGTETEtexte{#5}{#6}{\argdix}{\argonze}{\LANGUE{tape}{ruban} $3$}{0,-3.2}
\draw (tete) node[coordinate] (tete3) {};

\RUBANMACHINEDETURINGTETEtexte{#7}{#8}{\argdix}{\argonze}{\LANGUE{tape}{ruban} $4$}{0,-4.9}
\draw (tete) node[coordinate] (tete4) {};

%% ETAT:
     \draw (tete4) node [name=etat,draw,circle,anchor=north]  {#9};
%      \draw (etat) -- (tete3);
%     \draw  (etat) -|  (tete2b);
% \draw  (etat) -|  (tete1b);
%\end{tikzpicture}
}


%%%%%%%%%%%%%%%%%%%%%%%%%%%%%%%%%%%%%%%%%%%%%%%%%%%%%%%%%%%%%%%%%%%%%%%%%%%%
%                 MACHINES DE TURING (automate)
%%%%%%%%%%%%%%%%%%%%%%%%%%%%%%%%%%%%%%%%%%%%%%%%%%%%%%%%%%%%%%%%%%%%%%%%%%%%



\newcommand\dec{1ex}
\newcommand\myACCOLADE[4]{
\ACCOLADE{#1*\lcase+#4*\lcase-0.5*\lcase,0.5*\lcase+\dec}{#1*\lcase+#4*\lcase-0.5*\lcase+#3*\lcase,0.5*\lcase+\dec}{#2}
}


\newcommand\M{\Sigma}


\newcommand\sommet[2]{\node[state] (#1) #2 {$#1$};}
\newcommand\sommetinitial[2]{\node[state,initial] (#1) #2 {$#1$};}
\newcommand\sommetfinal[2]{\node[state,accepting] (#1) #2 {$#1$};}
\newcommand\gtransition[5]{\draw (#1) edge node {#2/#4 $#5$} (#3);}
\newcommand\gtransitionba[5]{\draw (#1) [loop above] edge node {#2/#4
    $#5$} (#3);}
\newcommand\gtransitionbb[5]{\draw (#1) [loop below] edge node {#2/#4
    $#5$} (#3);}
\newcommand\gtransitionbleft[5]{\draw (#1) [loop left] edge node {#2/#4
    $#5$} (#3);}
\newcommand\gtransitionbright[5]{\draw (#1) [loop right] edge node {#2/#4
    $#5$} (#3);}
\newcommand\gtransitionbl[5]{\draw (#1) [bend left] edge node {#2/#4
    $#5$} (#3);}
\newcommand\gtransitionbr[5]{\draw (#1) [bend right] edge node {#2/#4
    $#5$} (#3);}
\newcommand\gtransitionbadouble[9]{\draw (#1) [loop above] edge node {
\begin{tabular}{l}
#2/#4 $#5$ \\
  #7/#8 $#9$ \\
\end{tabular}
} (#3);}

\newcommand\SCALE{}

\newcommand\dessinmtp[2]{
\begin{center}
\begin{tikzpicture}[->,>=stealth',shorten >=1pt,auto,node distance=2.3cm,semithick,#2] 
#1
\end{tikzpicture}
\end{center}
}

\newcommand\dessinmt[1]{\dessinmtp{#1}{scale=1}}


\newcommand\ANIMATIONMT[1]{
\begin{overprint}
\begin{dessinp}{scale=0.52}
\foreach \av/\q/\b  [count=\xi]
in  
{#1}
{ 
\only<\xi| handout 0>{ 
 \StrLen{\av}[\Resintm] % ATTENTION IL FAUT CETTE SYNTAXE POUR FAIRE
                         % EQUIVALENT D UN EDEF
 \MACHINEDETURING{\av \b}{\Resintm}{$\q$}{5}{20}{}
}}
 \end{dessinp}
\end{overprint}
}

\newcommand\ANIMATIONMTd[1]{
\begin{overprint}
\begin{dessinp}{scale=0.52}
\foreach \av/\q/\b  [count=\xi]
in  
{#1}
{ 
\only<\xi| handout 0>{ 
 \StrLen{\av}[\Resintm] % ATTENTION IL FAUT CETTE SYNTAXE POUR FAIRE
                         % EQUIVALENT D UN EDEF
 \hspace{-1cm}\MACHINEDETURING{\av \b}{\Resintm}{$\q$}{5}{33}{}
}}
 \end{dessinp}
\end{overprint}
}

\def\MAXDET{20}

\newcommand\DIAGRAMMEESPACETEMPS[1]{
\draw node[coordinate] (tt) {};
\foreach \av/\q/\b  in  
{#1}
{ 

\RUBANMACHINEDETURING{\av{$\q$}\b}{0}{4}{\MAXDET}{}{tt};
\draw (tt) +(0,-\hcase) node[coordinate] (tt) {};
}
%\foreach \x in {0,...,19} {\draw (tt) + (\x*\lcase,0) node {$\vdots$};}
}



\newcommand\progun{
/q_0/0011, % initial
X/q_1/011,X0/q_1/11,X/q_2/0Y1,/q_2/X0Y1,
X/q_0/0Y1,XX/q_1/Y1,XXY/q_1/1, XX/q_2/YY, X/q_2/XYY,
XX/q_0/YY,XXY/q_3/Y,XXYY/q_3/B,XXYYB/q_4/B
}

\newcommand\DIAGRAMMEESPACETEMPSun{
\DIAGRAMMEESPACETEMPS
%% RECOPIE progun
{/q_0/0011, % initial
X/q_1/011,X0/q_1/11,X/q_2/0Y1,/q_2/X0Y1,
X/q_0/0Y1,XX/q_1/Y1,XXY/q_1/1, XX/q_2/YY, X/q_2/XYY,
XX/q_0/YY,XXY/q_3/Y,XXYY/q_3/B,XXYYB/q_4/B}
%% FIN RECOPIE
}

\newcommand\progdeux{
/q_0/0010, X/q_1/010, X0/q_1/10, X/q_2/0Y0,/q_2/X0Y0,
X/q_0/0Y0,XX/q_1/Y0,XXY/q_1/0,XXY0/q_1/%B,
}


\newcommand\DIAGRAMMEESPACETEMPSdeux{
\DIAGRAMMEESPACETEMPS
{
%% RECOPIE progdeux
/q_0/0010, X/q_1/010, X0/q_1/10, X/q_2/0Y0,/q_2/X0Y0,
X/q_0/0Y0,XX/q_1/Y0,XXY/q_1/0,XXY0/q_1/%B,
%% FIN RECOPIE
}
}

\newcommand\DIAGRAMMEESPACETEMPSdeuxvariante{
\def\memolcase{\lcase}
\def\memoMAXDET{20}
\def\MAXDET{15}
\def\lcase{6 ex}
\DIAGRAMMEESPACETEMPS
{
%% RECOPIE progdeux
/q_00/010, X/q_10/10, X0/q_11/0, X/q_20/Y0,/q_2X/0Y0,
X/q_00/Y0,XX/q_1Y/0,XXY/q_10/,XXY0/q_1B/%B,
%% FIN RECOPIE
}
\def\lcase{\memolcase}
\def\MAXDET{\memoMAXDET}
}


\newcommand\ANIMATIONun{
\ANIMATIONMT
%% RECOPIE progun
{/q_0/0011, % initial
X/q_1/011,X0/q_1/11,X/q_2/0Y1,/q_2/X0Y1,
X/q_0/0Y1,XX/q_1/Y1,XXY/q_1/1, XX/q_2/YY, X/q_2/XYY,
XX/q_0/YY,XXY/q_3/Y,XXYY/q_3/B,XXYYB/q_4/B}
%% FIN RECOPIE
}



\newcommand\DESSINALGO[3]{
% argument 1= ou
% argument 2= texte 
% argument 3= nom du noeud ex m
\draw node[draw,#1,minimum size=1cm] (#3) {#2};

% pour fitting box
\draw (#3.west) +(-0.2cm,0) node (#3c2) {};
\draw (#3.north) +(0,+0.2cm) node (#3c3) {};
\draw (#3.south) +(0,-0.2cm) node (#3c4) {};
}

\newcommand\DESSINALGOa[3]{
\DESSINALGO{#1}{#2}{#3}
\draw (#3.east) +(-0.1cm,0.2cm) node(#3a) {};
\draw[->] (#3a) -- ++(0.5cm,0) node (#3aa) {};
\draw (#3aa) node[anchor=west] (#3aaa) {\LANGUE{Accept}{Accepte}};
\draw (#3aaa.east) node(#3aaaa) {};
}

\newcommand\DESSINALGOe[4]{
\DESSINALGO{#1}{#2}{#3}
\draw (#3.east) +(-0.1cm,0.2cm) node(#3a) {};
\draw[->] (#3a) -- ++(0.5cm,0) node (#3aa) {};
\draw (#3aa) node[anchor=west] (#3aaa) {#4};
\draw (#3aaa.east) node(#3aaaa) {};
}


\newcommand\DESSINALGOar[3]{
\DESSINALGOa{#1}{#2}{#3}
\draw (#3.east) +(-0.1cm,-0.2cm) node(#3r) {};
\draw[->] (#3r) -- ++(0.5cm,0) node (#3rr) {};
\draw (#3rr) node[anchor=west] (#3rrr) {\LANGUE{Reject}{Refuse}};
\draw (#3rrr.east) node(#3rrrr) {};
}


\newcommand\WINDOWarg[7] {
%#1: position
%#2..7: contenu du tableau
\draw (#1) node[case]  {#2} \dcase;
\draw (#1) ++ (\lcase,0) node[case] {#3} \dcase;
\draw (#1) ++ (2*\lcase,0) node[case] {#4} \dcase;
\draw (#1) ++ (0,-\lcase) node[case] {#5} \dcase;
\draw (#1) ++ (\lcase,-\lcase) node[case] {#6} \dcase;
\draw (#1) ++ (2*\lcase,-\lcase) node[case] {#7} \dcase;
}

\newcommand\WINDOW[3]{
%#1: positioin
%#2: contenu
%#3: étiquette
\WINDOWarg{#1}
{\StrChar{#2}{1}}{\StrChar{#2}{2}}{\StrChar{#2}{3}}
{\StrChar{#2}{4}}{\StrChar{#2}{5}}{\StrChar{#2}{6}}
\draw (#1) +(-\lcase*1.2,-\lcase*0.5) node {#3};
}





%%%%%%%%%%%%%%%%%%%%%%%%%%%%%%%%%%%%%%%%%%%%%%%%%%%%%%%%%%%%%%%%%%%%%%%%%%%%
%                 LISTINGS
%%%%%%%%%%%%%%%%%%%%%%%%%%%%%%%%%%%%%%%%%%%%%%%%%%%%%%%%%%%%%%%%%%%%%%%%%%%%


% LISTINGS
\RequirePackage{listings}
 \lstset{language=Java,
 mathescape=true,
 showspaces=false,
 showstringspaces=false,
 frame=single,
morekeywords={recursive,then,div,function,integer,read,out,else,par,endpar,seq,endseq,and,not,read,write,repeat,until,undef,let,to,endfor,endif,endwhile,guess,in,parallel},
% basicstyle=\ttfamily\small,
% basewidth={1.1ex},
% keywordstyle=\bf\color{blue},
% %stringstyle=\bf\ttfamily\color{green},
% commentstyle=\bfseries\color{magenta},
 literate={:=}{{$\gets$ }}1 {<=}{{$\leq$}}2  {>=}{{$\geq$}}2
 {!=}{{$\neq$}}1 
 }
% \lstset{escapechar=\%}
\let\lst\lstinline
\def\beginlisting{\begin{lstlisting}}
\def\endlisting{\end{lstlisting}}  




\newenvironment{algo}[1]{
  \lstset{escapechar=\@}
  \def\input{$#1$}
  \def\out{out}
  \def\BEGIN{\textbf{read} \input }
 \def\BEGINN{\textbf{repeat} }
  \def\END{\textbf{until out!=undef}}
  \def\ENDD{\textbf{write \out}}
  \def\ENDT{\textbf{until $\ctl$ = $q_a$ or $\ctl$ = $q_r$}}
  \def\ENDR{\textbf{until $\ctl$ = $q_a$}}
  \def\ENDP{\textbf{until $\ctl$ = $q_a$}}
}{}
\newenvironment{algon}[1]{
  \lstset{numbers=left}
  \begin{algo}{#1}
}{\end{algo}}
\newcommand\FSM[3]{FSM(#1,\lst{#2},#3)}
\newcommand\FSMi[3]{FSM(#1,{#2},#3)}




%%%%%%%%%%%%%%%%%%%%%%%%%%%%%%%%%%%%%%%%%%%%%%%%%
%
%
%                      TRADUCTION EN COURS
%
%%%%%%%%%%%%%%%%%%%%%%%%%%%%%%%%%%%%%%%%%%%%%%%%%

\newcommand\ENTREPOT[1]{\NDLR{ENTREPOT ICI: #1}}
\newcommand\scite[2][]{\cite[#1]{#2}}
\newcommand\nospellcheck[1]{{#1}}
\newcommand\nd{nd}
\newcommand\NL{
 
}

\newenvironment{FRANCAIS}{
\small
\color{orange}
}
{\normalsize
\color{black}
}


\newenvironment{SEMIFRANCAIS}{
\small
\color{violet}
}
{\normalsize
\color{black}
}

\newcommand\QFRANCAIS[1]{
\begin{FRANCAIS}
#1
\end{FRANCAIS}
}


%%%%%%%%%%%%%%%%%%%%%%%%%%%%%%%%%%%%%%%%%%%%%%%%%
%
%
%                      Wikipedia
%
%%%%%%%%%%%%%%%%%%%%%%%%%%%%%%%%%%%%%%%%%%%%%%%%%

\providecommand\tightlist{}
\providecommand\Complex{\C}
\newenvironment{myalgie}{\begin{array}{lll}}{\end{array}}

\newcommand\nl{\ \\ }




%%%%%%%%%%%%%%%%%%%%%%%%%%%%%%%%%%%%%%%%%%%%%%%%%
%
%
%                      Pour citations précises
%
%%%%%%%%%%%%%%%%%%%%%%%%%%%%%%%%%%%%%%%%%%%%%%%%%



\newcommand\citeprecis[2]{{\cite[#1]{#2}}}
\newcommand\citeprecisinv[2]{\citeprecis{#2}{#1}}

\usepackage{csquotes}



%%%%%%%%%%%%%%%%%%%%%%%%%%%%%%%%%%%%%%%%%%%%%%%%%
%
%
%                      Travail sur couleurs
%
%   toutes les réponses pointent dans le même sens, et le schéma correspond exactement 
%. au profil perceptif d’un protanope** (déficit des cônes L / longueurs d’onde longues)
%  MAIS SEVERE
%
%%%%%%%%%%%%%%%%%%%%%%%%%%%%%%%%%%%%%%%%%%%%%%%%%


% ================================
% GROUPE 1 — Verts sombres
% ================================
\definecolor{deepgreen}{RGB}{0,160,40}
\definecolor{forestgreen}{RGB}{0,150,70}
\definecolor{mossgreen}{RGB}{0,140,90}
\definecolor{greenwave}{RGB}{0,150,90}
\definecolor{forestmist}{RGB}{0,160,90}
\definecolor{darkteal}{RGB}{0,130,70}

% ================================
% GROUPE 2 — Teal / verts bleutés
% ================================
\definecolor{seagreen}{RGB}{0,130,100}
\definecolor{freshteal}{RGB}{0,140,130}
\definecolor{greenteal}{RGB}{0,140,120}
\definecolor{blueteal}{RGB}{0,150,120}      % ancien "teal" (doublon corrigé)
\definecolor{lagoon}{RGB}{0,180,120}
\definecolor{greenmist}{RGB}{0,170,110}

% ================================
% GROUPE 3 — Turquoises / Aqua
% ================================
\definecolor{springgreen}{RGB}{0,190,100}
\definecolor{mintgreen}{RGB}{10,200,150}
\definecolor{oceanfoam}{RGB}{0,200,150}
\definecolor{aquagreen}{RGB}{30,210,170}
\definecolor{coastalaqua}{RGB}{40,210,180}
\definecolor{softaqua}{RGB}{40,210,190}

% ================================
% GROUPE 4 — Cyan / Aqua clairs
% ================================
\definecolor{emeraldteal}{RGB}{0,180,140}
\definecolor{deepcyan}{RGB}{0,180,160}
\definecolor{lightcyan}{RGB}{0,200,160}
\definecolor{aquamarine}{RGB}{0,200,160}
\definecolor{lightaqua}{RGB}{60,230,200}
\definecolor{mistblue}{RGB}{80,240,210}

% ================================
% GROUPE 5 — Bleus clairs
% ================================
\definecolor{paleaqua}{RGB}{80,200,190}
\definecolor{glacier}{RGB}{90,150,200}
\definecolor{skyblueA}{RGB}{60,150,220}      % ancien "skyblue"
\definecolor{lightsky}{RGB}{90,170,230}
\definecolor{iceblue}{RGB}{120,190,240}
\definecolor{frost}{RGB}{150,210,250}

% ================================
% GROUPE 6 — Bleus moyens
% ================================
\definecolor{steelblue}{RGB}{30,100,160}
\definecolor{cadetblue}{RGB}{70,120,180}
\definecolor{coolblue}{RGB}{110,170,220}
\definecolor{morningblue}{RGB}{140,200,240}
\definecolor{blueiceA}{RGB}{40,150,210}       % ancien "blueice" (doublon corrigé)
\definecolor{blueair}{RGB}{60,170,225}

% ================================
% GROUPE 7 — Azurs / Bleu saturé
% ================================
\definecolor{tealblue}{RGB}{0,110,130}
\definecolor{deepbluecyan}{RGB}{0,110,170}
\definecolor{azur}{RGB}{0,70,180}
\definecolor{brightazure}{RGB}{20,130,200}
\definecolor{lightazure}{RGB}{80,180,235}
\definecolor{skylight}{RGB}{100,200,245}

% ================================
% GROUPE 8 — Bleus sombres / violets
% ================================
\definecolor{nightblue}{RGB}{40,0,110}
\definecolor{slateblue}{RGB}{20,40,120}
\definecolor{violetA}{RGB}{60,50,160}         % ancien "violet" (distinct de violetB ci-dessous)
\definecolor{blueviolet}{RGB}{80,60,180}
\definecolor{purpleA}{RGB}{100,70,200}        % ancien "purple"
\definecolor{lightpurple}{RGB}{130,90,210}

% ================================
% GROUPE 9 — Violets saturés / indigos
% ================================
\definecolor{lilac}{RGB}{160,110,220}
\definecolor{indigoA}{RGB}{50,0,130}          % ancien "indigo"
\definecolor{deepindigo}{RGB}{70,10,150}
\definecolor{royalpurple}{RGB}{90,20,170}

% ================================
% GROUPE 10 — Fermeture cercle perceptif
% ================================
\definecolor{coolmint}{RGB}{0,130,70}
\definecolor{deepsea}{RGB}{0,130,70}

% ================================
% PALETTE DEUTÉRANOPE — noms uniques
% ================================
\definecolor{crimson}{RGB}{200,30,60}
\definecolor{brickred}{RGB}{180,40,40}
\definecolor{coral}{RGB}{220,90,70}
\definecolor{salmon}{RGB}{240,130,120}

\definecolor{magenta}{RGB}{200,40,150}
\definecolor{fuchsia}{RGB}{220,60,170}
\definecolor{violetB}{RGB}{140,70,190}        % ancien "violet"
\definecolor{indigoB}{RGB}{90,50,180}         % ancien "indigo"

\definecolor{royalblueA}{RGB}{50,80,200}      % ancien "royalblue"
\definecolor{azure}{RGB}{30,120,230}
\definecolor{skyblueB}{RGB}{80,160,230}       % second "skyblue"
\definecolor{tealB}{RGB}{0,150,170}           % second "teal"

\definecolor{marine}{RGB}{0,80,140}




%% Alternative

% === Générateur automatique fill/text pour protanopes ===
% Usage : \ProtaColorPair{mistblue}
% Suppose que mistblue est déjà défini par \definecolor{mistblue}{RGB}{...}

\newcommand{\ProtaColorPair}[1]{%
  % fill = couleur originale
  \colorlet{#1fill}{#1}%
  % text = version assombrie idéalement lisible
  \colorlet{#1text}{#1!80!black}%
}

%Utilisation
%\ProtaColorPair{mistblue}

%\textcolor{mistbluefill}{Couleur claire} \quad
%\textcolor{mistbluetext}{Couleur pour texte}

	\foreach \i/\name in {
	  1/deepgreen,
	  2/greenwave,
	  3/lagoon,
	  4/oceanfoam,
	  5/softaqua,
	  6/mistblue,
 	 7/lightsky,
	  8/blueiceA,   % <-- corrigé
	  9/deepcyan,
	  10/azur,
	  11/blueviolet,
	  12/royalpurple}
	  {\ProtaColorPair{\name}}
	  


\newcommand\ROUEPROTANOPEDOUZE{
	% Roue chromatique protanope — 12 couleurs
	\begin{tikzpicture}[scale=3]

	\foreach \i/\name/\nametext in {
	  1/deepgreen/deepgreen,
	  2/greenwave/greenwave,
	  3/lagoon/lagoontext,
	  4/oceanfoam/oceanfoam,
	  5/softaqua/softaquatext,
	  6/mistblue/mistbluetext,
 	 7/lightsky/lightskytext,
	  8/blueiceA/blueiceA,   % <-- corrigé
	  9/deepcyan/blueiceA,
	  10/azur/azur,
	  11/blueviolet/blueviolet,
	  12/royalpurple/royalpurple}
	{
	  % secteur coloré
 	 \filldraw[fill=\name, draw=black, line width=0.2pt]
	    ({90-30*\i}:1)
	    -- ({90-30*(\i-1)}:1)
	    -- (0,0) -- cycle;

 	 % label dans la bonne couleur
	  \node at ({90-30*(\i-0.5)}:1.25)
	    {\textcolor{\nametext}{\small \nametext}};
	}

	\end{tikzpicture}
}


\newcommand\ROUEPROTANOPEVINGT{
\begin{tikzpicture}[scale=3]

\foreach \i/\name/\nametext in {
  1/deepgreen/deepgreen,
  2/greenwave/greenwave,
  3/lagoon/lagoon,
  4/oceanfoam/oceanfoam,
  5/softaqua/softaquatext,
  6/mistblue/mistbluetext,
  7/lightsky/lightsky,
  8/blueiceA/blueiceA,      % corrigé
  9/deepcyan/deepcyan,
  10/azur/azur,
  11/brightazure/brightazure,
  12/skyblueA/skyblueA,     % corrigé
  13/coolblue/coolblue,
  14/glacier/glacier,
  15/marine/marine,
  16/violetA/violetA,      % corrigé
  17/blueviolet/blueviolet,
  18/purpleA/purpleA,      % corrigé
  19/royalpurple/royalpurple,
  20/indigoA/indigoA}      % corrigé
{
  % Secteur coloré
  \filldraw[fill=\name, draw=black, line width=0.2pt]
    ({90-18*\i}:1)
    -- ({90-18*(\i-1)}:1)
    -- (0,0) -- cycle;

  % Étiquette
  \node at ({90-18*(\i-0.5)}:1.25)
    {\textcolor{\nametext}{\small \nametext}};
}

\end{tikzpicture}
}

\newcommand\ROUEDEUTERANOPEDOUZE{
% Roue chromatique deutéranope — 12 couleurs
\begin{tikzpicture}[scale=3]

\foreach \i/\name in {
  1/crimson,
  2/brickred,
  3/coral,
  4/salmon,
  5/magenta,
  6/fuchsia,
  7/violetB,       % corrigé
  8/indigoB,       % corrigé
  9/royalblueA,    % corrigé
  10/azure,
  11/skyblueB,     % corrigé
  12/tealB}        % corrigé
{
  % Secteur coloré
  \filldraw[fill=\name, draw=black, line width=0.2pt]
    ({90-30*\i}:1)
    -- ({90-30*(\i-1)}:1)
    -- (0,0) -- cycle;

  % Label dans la couleur
  \node at ({90-30*(\i-0.5)}:1.25)
    {\textcolor{\name}{\small \name}};
}

\end{tikzpicture}
}



% === Macro intelligente pour foncer une couleur de façon lisible pour protanope ===
% Usage : \protadarken{mistblue}{40}  -> crée mistblue@dark
% Le 40 signifie : 40% de mélange avec du noir

\newcommand{\DarkenForText}[2]{%
  % #1 = nom de la couleur
  % #2 = pourcentage (40 conseillé si tu ne sais pas)
  %
  % On crée une nouvelle couleur #1@dark
  \colorlet{#1text}{black!#2!#1}%
}





%%%%

% Groupe 1
\DarkenForText{deepgreen}{20}
\DarkenForText{forestgreen}{20}
\DarkenForText{mossgreen}{20}
\DarkenForText{greenwave}{20}
\DarkenForText{forestmist}{20}
\DarkenForText{darkteal}{20}

% Groupe 2
\DarkenForText{seagreen}{20}
\DarkenForText{freshteal}{20}
\DarkenForText{greenteal}{20}
\DarkenForText{blueteal}{20}
\DarkenForText{lagoon}{20}
\DarkenForText{greenmist}{20}

% Groupe 3
\DarkenForText{springgreen}{20}
\DarkenForText{mintgreen}{20}
\DarkenForText{oceanfoam}{20}
\DarkenForText{aquagreen}{20}
\DarkenForText{coastalaqua}{20}
\DarkenForText{softaqua}{20}

% Groupe 4
\DarkenForText{emeraldteal}{20}
\DarkenForText{deepcyan}{20}
\DarkenForText{lightcyan}{20}
\DarkenForText{aquamarine}{20}
\DarkenForText{lightaqua}{20}
\DarkenForText{mistblue}{30}

% Groupe 5
\DarkenForText{paleaqua}{20}
\DarkenForText{glacier}{20}
\DarkenForText{skyblueA}{20}
\DarkenForText{lightsky}{20}
\DarkenForText{iceblue}{20}
\DarkenForText{frost}{20}

% Groupe 6
\DarkenForText{steelblue}{20}
\DarkenForText{cadetblue}{20}
\DarkenForText{coolblue}{20}
\DarkenForText{morningblue}{20}
\DarkenForText{blueiceA}{20}
\DarkenForText{blueair}{20}

% Groupe 7
\DarkenForText{tealblue}{20}
\DarkenForText{deepbluecyan}{20}
\DarkenForText{azur}{20}
\DarkenForText{brightazure}{20}
\DarkenForText{lightazure}{20}
\DarkenForText{skylight}{20}

% Groupe 8
\DarkenForText{nightblue}{20}
\DarkenForText{slateblue}{20}
\DarkenForText{violetA}{20}
\DarkenForText{blueviolet}{20}
\DarkenForText{purpleA}{20}
\DarkenForText{lightpurple}{20}

% Groupe 9
\DarkenForText{lilac}{20}
\DarkenForText{indigoA}{20}
\DarkenForText{deepindigo}{20}
\DarkenForText{royalpurple}{20}

% Groupe 10
\DarkenForText{coolmint}{20}
\DarkenForText{deepsea}{20}

% Palette deutéranope
\DarkenForText{crimson}{20}
\DarkenForText{brickred}{20}
\DarkenForText{coral}{20}
\DarkenForText{salmon}{20}

\DarkenForText{magenta}{20}
\DarkenForText{fuchsia}{20}
\DarkenForText{violetB}{20}
\DarkenForText{indigoB}{20}

\DarkenForText{royalblueA}{20}
\DarkenForText{azure}{20}
\DarkenForText{skyblueB}{20}
\DarkenForText{tealB}{20}

% Couleur supplémentaire
\DarkenForText{marine}{20}

%%%%%%%%%%%%%%%%%%%%%
%
%.   apx proof
%
%%%%%%%%%%%%%%%%%%%%%



% do
\newcommand\PREUVESENAPPENDIXCOMPATIBLE{
	\newenvironment{appendixproof}{{\bf \LANGUE{Proof}{Démonstration}}:}{\hfill $\Box$}
	\newtheorem{theoremrep}[theorem]{Theorem}
	\newtheorem{lemmarep}[theorem]{Lemma}
	\newtheorem{propositionrep}[theorem]{Proposition}
	\newtheorem{corollaryrep}[theorem]{Corollary}
	\newtheorem{summaryrep}[theorem]{Summary}	
}

\providecommand\apxparameter{}

%% proof inline
%\providecommand\apxparameter{appendix=inline}
%% proof in appendix
%\providecommand\apxparameter{}


\apxproofmode{
\usepackage[\apxparameter]{apxproof}
\theoremstyle{plain}

\newtheoremrep{theorem}{Theorem}
%\newtheoremrep{theoremconjecture}[theorem]{Theorem-To-Be-Proved=Conjecture}
%\newtheoremrep{lemmaconjecture}[theorem]{Conjecture-To-Be-Proved=Conjecture}
\newtheoremrep{claim}[theorem]{Theorem}
\newtheoremrep{proposition}[theorem]{Proposition}
\newtheoremrep{lemma}[theorem]{Lemma}
\newtheoremrep{corollary}[theorem]{Corollary}
%%bug mars 26: \newtheoremrep{summary}[theorem]{Summary}
% Du coup:
\newtheoremrep{resume}[theorem]{Summary}
%bug mars 26: \newtheoremrep{openproblem}[theorem]{Open Problem}
%\newtheoremrep{claim}[theorem]{Claim}
}


%\apxproofmodesubsection{
%%apx proof, pour que compteur correct
%\makeatletter
%\AtBeginDocument{%
%\let\axp@section\subsection
%  \let\axp@oldsection\subsection
%}
%\makeatother
%\renewcommand\appendixprelim{\section{Missing proofs}}
%}


\makeatletter
\apxproofmodesubsection{
\let\axp@section\subsection
%\let\axp@oldsection\subsection
\renewcommand\appendixprelim{\section{Proofs from main documents}}
}
\makeatother


\makeatletter
\newcommand\appendixapxproofsubsection{
\appendix
\let\apx@orig@appendix\appendix
\renewcommand{\appendix}{}
}
\makeatother


%%% Pour checker/vérifier que bon


\newcommand\docheck[1]{\textcolor{check}{#1}}

\usepackage{xcolor}
%\definecolor{chose}{RGB}{0,110,50}
%\definecolor{forestgreen}{RGB}{0,150,70}
%\definecolor{ctruc}{RGB}{164,164,164}
%\definecolor{oceanfoam}{RGB}{0,200,150}
\definecolor{check}{RGB}{160,160,0}






}


 
%%%%%%%%%%%%%%%%%%%%%%%%%%%%%%%%%%%%%%%%%%%%%%%%%
%
%
%               Déclaration de cours
%
%  (l'argument 1 = numéro sert pas)
%
%%%%%%%%%%%%%%%%%%%%%%%%%%%%%%%%%%%%%%%%%%%%%%%%%


\newcommand\COURS[3]{
  \expandafter\newcommand\csname numerocours#2\endcsname{#1}
  \expandafter\newcommand\csname nomcours#2\endcsname{#3}
}



\newcommand\NANCY[3]{
  \expandafter\newcommand\csname numerocours#2\endcsname{#1}
  \expandafter\newcommand\csname nomcours#2\endcsname{#3}
}


\newcommand\INFquatrecentdouze[3]{
  \expandafter\newcommand\csname numerocours#2\endcsname{#1}
  \expandafter\newcommand\csname nomcours#2\endcsname{#3}
}

\newcommand\INFcinqsixun[3]{
  \expandafter\newcommand\csname numerocours#2\endcsname{#1}
  \expandafter\newcommand\csname nomcours#2\endcsname{#3}
}



\newcommand\MPRI[3]{
  \expandafter\newcommand\csname numerocours#2\endcsname{#1}
  \expandafter\newcommand\csname nomcours#2\endcsname{#3}
}

\newcommand\COURSEINCLUDE[2][]{
  \renewcommand\IFNOTSUB[1]{}
  \renewcommand\MYCOURSE[1]{\chapter{\ifnotmodeDIFFUSIONFINALE{#1} \csname nomcours##1\endcsname}
    \ifnotmodeDIFFUSIONFINALE{\mychapter{Fichier: #2}}
  }
   \renewcommand\MYCOURSEnew[1]{\chapter{\ifnotmodeDIFFUSIONFINALE{#1} ##1 }
    \ifnotmodeDIFFUSIONFINALE{\mychapter{Fichier: #2}}
  }
%  \renewcommand\FINSUBMPRI{}
  \input{#2}
}

\newcommand\STRESSCOURSEINCLUDE[1]{
  \COURSEINCLUDE[!STRESS!]{#1}
  }

  \newcommand\subCOURSEINCLUDE[1] {
    \NDLR{BEGIN:  J'inclus/ Je répète ICI #1 }
    \input{INCLUDE/#1}
    \NDLR{END:  J'inclus/ Je répète ICI #1 }
}

%%%%%%%%%%%%%%%%%%%%%%%%%%%%%%%%%%%%%%%%%%%%%%%%%
%
%
%                      POUR INDICATIONS: MESSAGES
%
%
%%%%%%%%%%%%%%%%%%%%%%%%%%%%%%%%%%%%%%%%%%%%%%%%%


% ancienne version
% \newcommand\montruc[1]{\begin{center} {\fbox{\fbox{{#1}
%         }}} \end{center}}




\newcommand\montruc[2][]{
  \ifnotmodeDIFFUSIONFINALE{
  \begin{maboitetruc}[#1]{\danger} #2
  \end{maboitetruc}
%%  \todo[inline,color=green!20,caption={2do}, ssss#1]{\begin{minipage}{\textwidth-4pt}
%\todo[inline,color=green!5%,caption={2do}
%]{\begin{minipage}{\textwidth-4pt}
%    #2\end{minipage}}
}
}

\newcommand\montrucdanger[3][]{
  \ifnotmodeDIFFUSIONFINALE{
                    \montruc[#1]{\danger #2: #3}
                    }
}

\newcommand\montrucdangermargin[3][]{
  \ifnotmodeDIFFUSIONFINALE{
  \montruc[#1]{\danger #2: #3}
\marginpar{\danger #2: #3}
}
 }



%%%%%%%%%%%%%%%%%%%%%%%%%%%%%%%%%%%%%%%%%%%%%%%%%
%
%
%                      GESTION DES TRUCS MARQUES VERSIONS
%
%
%%%%%%%%%%%%%%%%%%%%%%%%%%%%%%%%%%%%%%%%%%%%%%%%%

\makeatletter
\renewcommand\beginmarkversion{\montrucdanger[colback=yellow!10]{=== begin NOT}{}\@Vs@sffbox{\@currenvir$>$}}
 \renewcommand\endmarkversion{\@Vs@sffbox{$<$\@currenvir}\montrucdanger[colback=yellow!10]{end=== NOT}{}}
\makeatother 


%%%%%%%%%%%%%%%%%%%%%%%%%%%%%%%%%%%%%%%%%%%%%%%%%
%
%
%                      VARIANTES DE MESSAGES
%
%
%%%%%%%%%%%%%%%%%%%%%%%%%%%%%%%%%%%%%%%%%%%%%%%%%

 \newcommand\NDLR[1]{\montrucdanger[colback=blue!20]{NDLR}{#1}}
 \newcommand\NLDR[1]{\NDLR{#1}}
 \newcommand\VOIRAUSSI[1]{\montrucdanger[colback=yellow!20]{VOIR AUSSI}{#1}}
 \newcommand\TODO[1]{\montrucdanger[colback=green!20]{TODO}{#1}}
 \newcommand\ATTENTION[1]{\montruc[colback=brown!20]{\danger ATTENTION#1}}
 \newcommand\COMPRENDPAS[1]{\montrucdanger[colback=yellow!20]{JE
     COMPRENDS PAS CA}{#1}}
  \newcommand\POSSIBLE[1]{\montrucdanger[colback=yellow!10]{
      Serait possible}{#1}}
 \newcommand\ENLEVE[1]{\montrucdanger[colback=gray!20]{
      ENLEVE}{#1}}
 
 
  \newcommand\DEPLACE[1]{
    \vspace{0.5cm}
    \hrule
    \hrule
      \vspace{0.5cm}
    \montrucdangermargin[colback=black!40]{DEPLACE/MISSING CONCEPT
    }{#1}
      \vspace{0.5cm}
    \hrule
    \hrule
      \vspace{0.5cm}
  }

\newcommand\BESOINDE[1]{\montrucdangermargin[colback=pink!20]{ATTTENTION
     BESOIN DE}{#1}}
\newcommand\MISSINGCONCEPT[1]{\montrucdangermargin[colback=pink!20]{MISSING CONCEPT
     }{#1}}
\newcommand\SAUTE[1]{\montrucdangermargin[colback=pink!20]{SAUTE
  }{#1}}
\newcommand\SHOULDBEHERE[1]{\BEGINTEMPENVCOLOR{blue!30}\montrucdanger[colback=pink!10]{
SHOULD
    BE HERE
  }{    
    \vspace{1cm}
    
    \centerline{$\vdots$}
    
   \centerline{ #1}
    
      \centerline{$\vdots$}
          
    \vspace{1cm}
    }\ENDTEMPENVCOLOR}
\newcommand\COULDBEHERE[1]{\montrucdangermargin[colback=pink!20]{COULD
    BE HERE
  }{#1}}



%%%%%%%%%%%%%%%%%%%%%%%%%%%%%%%%%%%%%%%%%%%%%%%%%
%
%
%                      POUR INDICATIONS: SOURCES
%
%
%%%%%%%%%%%%%%%%%%%%%%%%%%%%%%%%%%%%%%%%%%%%%%%%%


\newcommand\SOURCEURL[1]{
  \ifnotmodeDIFFUSIONFINALE{
    \marginpar{ \textcolor{magenta}{Source: \url{#1}} }
    }
}

\newcommand\SOURCETXT[1]{
  \ifnotmodeDIFFUSIONFINALE{
    \marginpar{ \textcolor{magenta}{Source: {#1}} }
    }
}

\newcommand\SOURCE[1]{
  \nocite{#1}
  \ifnotmodeDIFFUSIONFINALE{
    \marginpar{ \textcolor{magenta}{Source: \cite{#1}} }
    }
}


%--------------------------------
%--- Raccourci Sources particulières  ---
%--------------------------------

\newcommand\SOURCEJECH{\SOURCE{jech2003set}}
\newcommand\SOURCEKECHRIS{\SOURCE{kechris2012classical}}
\newcommand\SOURCENOTESSETTHEORY{\SOURCE{moschovakis2006notes}}
\newcommand\SOURCEKOZENCT{\SOURCE{LivreKozenTC}}
\newcommand\SOURCEMARKER{\SOURCE{marker2002descriptive}}
\newcommand\SOURCEMOSCHO{\SOURCE{moschovakis2009descriptive}}



%%%%%%%%%%%%%%%%%%%%%%%%%%%%%%%%%%%%%%%%%%%%%%%%%
%
%
%                      POUR INDICATIONS: DELIRES/COMMENTAIRES
%
%
%%%%%%%%%%%%%%%%%%%%%%%%%%%%%%%%%%%%%%%%%%%%%%%%%


%%%%%% ABOUT PLAN

%\newcommand\PLANBUT[1]{\montruc{\begin{minipage}{0.8\textwidth} #1\end{minipage}}}


%%%%%% ABOUT RECHERCHE


\ifnotmodeDIFFUSIONFINALE{
\newenvironment{recherche}
{}{}
%{\color{blue} BEGIN-RECHERCHE:
%%\ifmodeDIFFUSIONFINALE{diffusion finale}
%%\ifnotmodeDIFFUSIONFINALE{not diffusion finale}
%%	\ifdefined\SAFEMODE
%%		SAFEMODE
%%	\else
%%		NOTSAFEMODE
%%	\fi
%  \todo[inline, color=blue]{recherche ici} }{
%  END-RECHERCHE \color{black} }
\tcolorboxenvironment{recherche}{colback=blue!5!white, colframe=red!75!black,fonttitle=\bfseries, colbacktitle=blue!85!black,enhanced,
attach boxed title to top center={yshift=-2mm},
  title={FAIRE RECHERCHE ICI}}
}


%%%%%% ABOUT DELIRE


%\ifnotmodeDIFFUSIONFINALE{
%\newenvironment{delire}{\color{violet} BEGIN-DELIRE: }{
%  END-DELIRE \color{black} }
%}


%%%%%% ABOUT PERSO


%\ifnotmodeDIFFUSIONFINALE{
%\newenvironment{perso}{\color{brown}
%  BEGIN-PERSO: }{
%  END-PERSO\color{black} }
%}



%%%%%% ABOUT COMMENTAIRE

\ifnotmodeDIFFUSIONFINALE{
\newenvironment{commentaire}{\begin{maboitecommentaire}[colback=brown!30]{Commentaire} Commentaire :}{\end{maboitecommentaire}}
\newcommand\COMMENTAIRE[2][Commentaire ]{
\begin{maboitecommentaire}[colback=brown!30]{#1}
Commentaire : #2
\end{maboitecommentaire}
}
}

%%%%%% ABOUT HORSUJET


%\ifnotmodeDIFFUSIONFINALE{
%\newenvironment{probablementhorsujet}{\color{gray}
%  BEGIN-PROBABLEMENT HORS SUJET: }{
%  END-PROBABLEMENT HORS SUJET \color{black} }
%}

%%%%%% ABOUT BONUS


\ifnotmodeDIFFUSIONFINALE{
\newenvironment{pasbesoinpourlinstant}{NOT NEEDEED FOR
  THIS SECTION: }{ }
    \tcolorboxenvironment{pasbesoinpourlinstant}{colback=green!5!white, colframe=gray!75!black,fonttitle=\bfseries, colbacktitle=green!85!black,enhanced,
attach boxed title to top left={yshift=-2mm},
  title={PAS BESOIN POUR L'INSTANT}}
}

%%%%%% ABOUT BONUS


\ifnotmodeDIFFUSIONFINALE{
\newenvironment{bonus}{BONUS: }{
  }
  \tcolorboxenvironment{bonus}{colback=gray!5!white, colframe=gray!75!black,fonttitle=\bfseries, colbacktitle=gray!85!black,enhanced,
attach boxed title to top left={yshift=-2mm},
  title={BONUS}}
}

%%%%%% ABOUT BONUS


\ifnotmodeDIFFUSIONFINALE{
\newcommand\ATTENTIONREPETE[1]{
  \ATTENTION{Repete (attention modif): #1}
  \begin{remarkolivier}
Repete (attention modif): #1
\end{remarkolivier}
}
}


%%%%%%%%%%%%%%%%%%%%%%%%%%%%%%%%%%%%%%%%%%%%%%%%%
%
%
%                      MACROS Papiers Paulin
%
%
%%%%%%%%%%%%%%%%%%%%%%%%%%%%%%%%%%%%%%%%%%%%%%%%%


\newcommand{\x}{\times}
\newcommand{\s}{\sigma}
\newcommand\sep{.}
\newcommand\bzero{\zero}
\newcommand\bun{\un}
\renewcommand{\a}{\alpha}
\def\hd{{\sf hd}}
\def\tl{{\sf tl}}
\def\cons{{\sf cons}}
\def\Cons{{\sf Cons}}
\def\Pr{{\sf Pr}}
\def\Op{{\sf Op}}
\def\Rel{{\sf Rel}}
\def\Equal{{\sf Equal}}
\def\Eval{{\sf Eval}}
\def\Gate{{\sf Gate}}
\def\Selection{{\sf Selection}}
\def\Result{{\sf Result}}
\def\Circuit{{\sf Circuit}}
\def\next{{\sf next}}
\def\comp{{\sf comp}}
\def\finalcomp{{\sf finalcomp}}
\def\Pick{{\sf Pick}}
\def\Hd{{\sf Hd}}
\def\Tl{{\sf Tl}}
\def\RevHd{{\sf RevHd}}
\def\Node{{\sf Node}}
\def\Length{{\sf Length}}
\def\Tape{{\sf Tape}}
\def\Head{{\sf Head}}
\def\Time{{\sf Time}}
\def\unpad{{\sf unpad}}
\def\concat{{\sf concat}}
\def\rg{{\sf right}}
\def\lf{{\sf left}}
\def\st{{\sf state}}
\def\tope{{\sf top}}


%%%%%%%%%%%%%%%%%%%%%%%%%%%%%%%%%%%%%%%%%%%%%%%%%
%
%
%                      MACROS POLY INF561
%
%
%%%%%%%%%%%%%%%%%%%%%%%%%%%%%%%%%%%%%%%%%%%%%%%%%

\newcommand\buglst[1]{#1}

%\newcommand\mydefref[1]{\expandafter\newcommand\csname #1\endcsname{\ref{#1}}}
%\mydefref{refthrisc}
%\mydefref{refthcinqhuit}
%\mydefref{refthcinqsix}
%\mydefref{refthfondamcaract}
%\mydefref{refdefcircuit}
%%% sinon, fait à la main
%\newcommand\refchapalgo{\ref{chapalgo}}
%\newcommand\refchapmodeles{\ref{chapmodeles}}
%\newcommand\refchappreliminaires{\ref{chappreliminaires}}
%\newcommand\refchapcircuits{\ref{chapcircuits}}
%\newcommand \refpropfondamen{\ref{propfondamen}}
%\newcommand\refthhierspace{ref{thhierspace}}
%\newcommand\refchapcomplexitetemps{\ref{chapcomplexitetemps}}
%\newcommand\refdefverifun{\ref{defverifun}}
%\newcommand\refdefverif{\ref{defverif}}
%\newcommand\refthmachinevs{ref{thmachinevs}}
%\newcommand\reflemnotationconfig{\ref{lemnotationconfig}}
%\newcommand\reflemfondam{\ref{lemfondam}}
%\newcommand\refthtroisat{\ref{thtroisat}}
%\newcommand\refthfondam{\ref{thfondam}}
%\newcommand\refthsansram{\ref{thsansram}}
%\newcommand\refdefsram{\ref{defsram}}
%!TEX encoding = UTF-8 Unicode

%% TRUCS DU POLY INF561

\ifnotSLIDE{
\LANGUE{
\spnewtheoremi{exemple}{\NDLR{ECRIT EXEMPLE AU LIEU DE
  EXAMPLE}}{\bfseries}{\rmfamily}{theorem} 
\spnewtheoremi{remarque}{\NDLR{ECRIT REMARQUE AU LIEU DE REMARK}}{\bfseries}{\rmfamily}{theorem} 
}
{
\spnewtheoremi{exemple}{Exemple}{\bfseries}{\rmfamily}{theorem} 
\spnewtheoremi{remarque}{Remarque}{\bfseries}{\rmfamily}{theorem} 
}
}
            
            \newcommand\PERSOSOURCEPOSSIBLE[1]{\PERSOSOURCE{#1}\PERSOCOMMENTAIRE{#1}}
\newcommand\PERSOPOSSIBLE[1]{\POSSIBLE{#1}}
% \newcommand\PERSOPOSSIBLEbug[1]{\NDLR{\textbf{-- possible --
%       enleve car pb compilation}}}
\newcommand\PERSOETAIT[1]{\NDLR{ETAIT: #1}}
%   \begin{minipage}{10cm}
% \NDLR{\textbf{\tiny  -- était --      #1 -- était
%     --
%   \end{minipage}
% }}}

\providecommand\nl{}

\newcommand\PERSOSOURCE[1] { \SOURCETXT{#1}}
\newcommand\PERSOCOMMENTAIRE[1]{\NDLR{#1}}
\newcommand\PERSOCOMMENTAIREd[2]{\NDLR{#1 \textbf{#2}}}


\newenvironment{dessinpb}
{%\shorthandoff{:}
%\selectlanguage{english}
\begin{center}\begin{tikzpicture}[baseline=(current bounding
      box.north)]
%babelbug with tikz
}
{\end{tikzpicture}
%\selectlanguage{frenchb}
\end{center}}



\newcommand\PERSOMANQUE[1]{\NDLR{\textbf{-- manque --
      #1 --manque--}}}

\newcommand\PERSOENLEVE[1]{\ENLEVE{#1}}






%%%%%%%%%%%%%%%%%%%%%%%%%%%%%%%%%%%%%%%%%%%%%%%%%
%
%
%                      MACROS POLY INF412
%
%
%%%%%%%%%%%%%%%%%%%%%%%%%%%%%%%%%%%%%%%%%%%%%%%%%

%\mydefref{refEEXTtruc}
%\mydefref{refEEXTchappreliminaires}
%\mydefref{refEEXTsecarbrebinaire}
%\mydefref{refEEXTsectermes}
%\mydefref{refEEXTchappredicats}
%\mydefref{refEXTchappredicats}
%\mydefref{refEXTchappreliminaires}
%\mydefref{refEEXTdefTuring}
%\mydefref{refEEXTpropnondet}
%\mydefref{refEEXTchapcompletude}
%\mydefref{refEEXTdefconfigalternative}
%\mydefref{refEEXTchapmachinesdeTuring}
%\mydefref{refEEXTpbdedecision}
%\mydefref{refEEXTchapcalc}
%\mydefref{refEEXTincompletude}
%\mydefref{refEEXTthhierspace}
%\mydefref{refEEXTchapcomplexitetemps}

\tikzstyle{nd} = [shape=circle,draw]


\newcommand\mmod{mod}

  %%% VALUATIONS
\newcommand\mmV{v}
\newcommand\mmVb{\overline{\mmV}}
\newcommand\mmVe{\mmVb}

%%% Corrections

\newcommand\exoref[1]{\ref{exo:#1-stt}}

\newenvironment{correction}[1]
{\par\medskip\noindent \label{exo:#1-cor} \hbox{\bf Exercice \ref{exo:#1-stt} 
  \   (page \pageref{exo:#1-stt}). }}
{\par\medskip}

% PROBLEMES DE DECISION

\newcommand\DECISIONint[3]{
%Le problème de décision #1
\begin{itemize}
\item[\textbf{Donnée}:] #2
\item[\textbf{Réponse}:] #3
\end{itemize}
}

\newcommand\INDECIDABLE[5]{
\DECISION{#1}{#2}{#3}

\begin{#4} \label{#5}
Le problème $#1$ % ATTENTION ENLEVE $#1$ en 2020. REMIS 2023
est indécidable.
\end{#4}
}


\newcommand\INDECIDABLEenglish[5]{
\DECISIONenglish{#1}{#2}{#3}

\begin{#4} \label{#5}
The problem $#1$
is undecidable.
\end{#4}
}

\newcommand\INDECIDABLEq[5]{
\DECISION{#1}{#2}{#3}

\begin{#4} \label{#5}
Le problème $#1$
est ?
\end{#4}
}

\newcommand\INDECIDABLEqenglish[5]{
\DECISION{#1}{#2}{#3}

\begin{#4} \label{#5}
The problem $#1$
is ?
\end{#4}
}



\newcommand\DECISIONintenglish[3]{
%Le problème de décision #1
\begin{itemize}
\item[\textbf{Input}:] #2
\item[\textbf{Answer}:] #3
\end{itemize}
}

\newcommand\DECISIONi[4]{
\item $#1$
%Le problème de décision #1
\DECISIONint{#1}{#2}{#3}
}

\newcommand\DECISIONienglish[4]{
\item $#1$
%Le problème de décision #1
\DECISIONintenglish{#1}{#2}{#3}
}


\newcommand\DECISION[3]{
\begin{definition}[$#1$]\ 

%Le problème de décision #1
\DECISIONint{#1}{#2}{#3}
\end{definition}
}


\newcommand\DECISIONenglish[3]{
\begin{definition}[$#1$]\ 
%
%Le problème de décision #1
\DECISIONintenglish{#1}{#2}{#3}
\end{definition}
}



\newcommand\NPC[4]{ 
\DECISION{#1}{#2}{#3}

\begin{#4}
Le problème $#1$
est $\NP$-complet. \LANGUE{\myindex{NP-completeness}}{\myindex{NP-completude@$\NP$-complétude}}
\end{#4}
}


\newcommand\NPCenglish[4]{ 
\DECISIONenglish{#1}{#2}{#3}
\begin{#4}
The problem $#1$
is $\NP$-complete. \LANGUE{\myindex{NP-completeness}}{\myindex{NP-completude@$\NP$-complétude}}
\end{#4}
}


\newcommand\NPCp[5]{
\DECISION{#1}{#2}{#3}

\begin{#4}[#5]
Le problème $#1$
est $\NP$-complet.
\end{#4}
}


\newcommand\NPCpenglish[5]{
\DECISIONenglish{#1}{#2}{#3}

\begin{#4}[#5]
The problem $#1$
is $\NP$-complete.
\end{#4}
}




\newcommand\BRUNOCOMMENTAIRE[1]{\NDLR{Bruno commentaire: #1}}

\ifnotSLIDE{
%\spnewtheorem{exercicec}{Exercice}{\bfseries}{\rmfamily}
%\spnewtheorem{exerciced}{Exercice}{\bfseries}{\rmfamily}
%\spnewtheorem{exercicedc}{Exercice}{\bfseries}{\rmfamily}
%\spnewtheorem{exercicecenglish}{Exercice}{\bfseries}{\rmfamily}
%\spnewtheorem{exercicedenglish}{Exercice}{\bfseries}{\rmfamily}
%\spnewtheorem{exercicedcenglish}{Exercice}{\bfseries}{\rmfamily}
\spnewtheorem{notation}{Notation}{\bfseries}{\rmfamily}
}

%%% INDEX
 \newcommand\motnouv[1]{\emph{#1}\index{#1}}%\SIGNALEMOTNOUV{#1}}%\index{#1}}
\newcommand\INTmotnouv[1]{\emph{#1}}
\newcommand\myindex[1]{\index{#1}}
%\marginpar{-#1=}

\newcommand\synonyme[2]{\kindex{#1|see{#2}}}
\newcommand\indexM[1]{\index{$#1$}}
\newcommand\idx[1]{#1\index{#1}}
\newcommand\idxm[1]{#1\index{$#1$}}
\newcommand\idxM[1]{$#1$\index{$#1$}}
%\newcommand\idxmBUG[1]{#1}
\newcommand\idxmm[2]{#1\index{#2@$#1$}}
\newcommand\idxi[2]{#1\index{#2}}
%\newcommand\motnouv[1]{\INTmotnouv{#1}\index{#1}}
 \newcommand\motnouvi[2]{\INTmotnouv{#1}\index{#2}} % indexage manuel
 \newcommand\motnouva[2]{\INTmotnouv{#1}\index{#2@#1}} % indexage: 2ième
%                                 % argumet = version sans accent
\newcommand\motnouvm[1]{\INTmotnouv{#1}} %indexage manuel fait ailleurs


\newenvironment{PERSOVERBATIM}{}{}
 \newcommand{\joNP}[3]{{ Problème \textbf{ { \motnouv{##1}}:}
            \begin{itemize}
            \item  \textbf{ Instance:} {##2} 
            \item   \textbf{ Question:} {##3}
            \end{itemize}
          }}

%%%%%%%%%%%%%%%%%%%%%%%%%%%%%%%%%%%%%%%%%%%%%%%%%
%
%
%                      POUR COMPATIBILITE AVEC SLIDES & AUTRE
%
%
%%%%%%%%%%%%%%%%%%%%%%%%%%%%%%%%%%%%%%%%%%%%%%%%%


\newcommand\defin[1]{\textbf{#1}}


%%%   COMPATIBILITE AVEC SLIDES

\newcommand\soustitre[1]{\textbf{#1}}


%% COMPATIBILITE AVEC MA THESE

\providecommand\motnouv[1]{\emph{#1}}
\providecommand\motnouvi[2]{\emph{#1}}

\newcommand\papiernorm[1]{\myop{norm}}

%%% OLIVIER
\newcommand\olivier[1]{ \ATTENTION{Olivier: #1}}


%%%%%%%%%%%%%%%%%%%%%%%%%%%%%%%%%%%%%%%%%%%%%%%%%
%
%
%                      Page de Garde Cours X
%
%
%%%%%%%%%%%%%%%%%%%%%%%%%%%%%%%%%%%%%%%%%%%%%%%%%


\includeversion{metdate}

\newcommand\XPAGEDEGARDE[2]{
\thispagestyle{empty}
\title{{#1} \\[2cm] %\\[3cm] 
{\large #2}
}
\author{
 Olivier Bournez% \inst{1}
\\[0.5cm] 
bournez@lix.polytechnique.fr%{\myaddress}
}
%\date{
%\vspace{2cm}
%\includegraphics[height=2.5cm]{/Users/bournez/lib/LaTeX/Perso/FIG-COMMONS/logo_X_new}
%}
\begin{metdate}
\date{
\vspace{2cm}
Version \LANGUE{of}{du} \today \\[1cm] 
\includegraphics[height=3.5cm]{/Users/bournez/lib/LaTeX/Perso/FIG-COMMONS/logo_X_new}
}
\end{metdate}
\maketitle
}


%%%%%%%%%%%%%%%%%%%%%%%%%%%%%%%%%%%%%%%%%%%%%%%%%
%
%
%                      MACROS UTILES
%
%
%%%%%%%%%%%%%%%%%%%%%%%%%%%%%%%%%%%%%%%%%%%%%%%%%

%%%                      Macros
 
%-------------
%--- vectors  ---
%-------------

 
\newcommand\vectorl[1]{{\mathbf#1}}
\newcommand\tu[1]{\vectorl{#1}}

\newcommand\vp{\vectorl{p}}
\newcommand\va{\vectorl{a}}
\newcommand\vb{\vectorl{b}}
\newcommand\vc{\vectorl{c}}
\newcommand\vf{\vectorl{f}}
\newcommand\vn{\vectorl{n}}
\newcommand\vx{\vectorl{x}}
\newcommand\vX{\vectorl{X}}
\newcommand\vy{\vectorl{y}}
\newcommand\vz{\vectorl{z}}
\newcommand\ve{\vectorl{e}}
\newcommand\vv{\vectorl{v}}
\newcommand\vr{\vectorl{r}}
\newcommand\vu{\vectorl{u}}
\newcommand\vA{\vectorl{A}}
\newcommand\vB{\vectorl{B}}
\newcommand\vC{\vectorl{C}}
\newcommand\vD{\vectorl{D}}

%-------------
%--- overline ---
%-------------


\newcommand\ox{\overline{x}}
\newcommand\oy{\overline{y}}
\newcommand\oc{\overline{c}}
\newcommand\oa{\overline{a}}
\newcommand\ow{\overline{w}}
\newcommand\od{\overline{d}}
\newcommand\ob{\overline{b}}
\newcommand\oP{\overline{P}}
\newcommand\oee{\overline{e}}
\newcommand\ot{\overline{t}}
\newcommand\ou{\overline{u}}
\newcommand\ov{\overline{v}}
\newcommand\oz{\overline{z}}
\newcommand\am{\overline{a}}
\newcommand\om{\overline{m}}
\newcommand\oj{\overline{j}}

%----------------
%--- ensembles ---
%----------------

\newcommand\Q{\mathbb{Q}}
\newcommand\R{\mathbb{R}}
\newcommand\Z{\mathbb{Z}}
\providecommand\C{\mathbb{C}}
\newcommand\N{\mathbb{N}}
\newcommand\K{\mathbb{K}}
\newcommand\F{\mathbb{F}}

\newcommand\dyadic{\mathbb{D}}
\newcommand\Zd{{\Z_2}}
\newcommand\Zp{{\Z_p}}
%
\newcommand\RP{\mathbb{R}^{>0}}
\newcommand\QP{\mathbb{Q}^{>0}}


%----------------
%--- lettres cal ---
%----------------

\newcommand\mK{\mathcal{K}}
\newcommand\mD{\mathcal{D}}
\newcommand\mA{\mathcal{A}}
\newcommand\mL{\mathcal{L}}
\newcommand\mX{\mathcal{X}}
\newcommand\mG{\mathcal{G}}
\newcommand\mE{\mathcal{E}}
\newcommand\mH{\mathcal{H}}
\newcommand\mR{\mathcal{R}}
\newcommand\mF{\mathcal{F}}
\newcommand\mJ{\mathcal{J}}
\newcommand\mO{\mathcal{O}}
\newcommand\mP{\mathcal{P}}
\newcommand\mN{\mathcal{N}}
\newcommand\mS{\mathcal{S}}
\newcommand\mM{\mathcal{M}}
\newcommand\mC{\mathcal{C}}
\newcommand\mV{\mathcal{V}}
\newcommand\mI{\mathcal{I}}
\newcommand\mT{\mathcal{T}}

\newcommand\fC{\mathcal{C}}

%----------------
%--- lettres frak ---
%----------------


\newcommand\kM{\mathfrak{M}}
\newcommand\kN{\mathfrak{N}}
\newcommand\kU{\mathfrak{U}}
\newcommand\kg{\mathfrak{g}}



%--------------------------
%--- Façon noter noms classes ---
%---------------------------

\newcommand{\cp}[1]{\operatorname{#1}}


%--------------------------
%--- Calculabilité  ---
%---------------------------

% -- classes
\LANGUE{
\newcommand\RE{\cp{CE}}
}
{\newcommand\RE{\cp{RE}}
}

\newcommand\Rec{\cp{D}}
\newcommand\PR{\cp{PR}}
\newcommand\Gregn{\mE_n}


%--------------------------
%--- Réductions  ---
%---------------------------

\newcommand\lem{\le_{m}}
\newcommand\lelog{\le_{L}}
\newcommand\equivlog{\equiv_{L}}



%--------------------------
%--- Complexité  ---
%---------------------------

% COMPLEXITE

% -- P
\newcommand{\Ptime}{\cp{P}}      % P
\newcommand\PTIME{\Ptime}
\newcommand\PTIMEs[1]{{\PTIME}_{#1}}
\newcommand\PTIMEsp[1]{{\PTIME}_{#1}^0}
\renewcommand\P{\PTIME}
\newcommand{\FPtime}{\cp{FPTIME}}

% -- NP
\newcommand\NP{\cp{NP}}
\newcommand{\NPtime}{\NP}
\newcommand\NPs[1]{\cp{NP}_{#1}}
\newcommand{\FNP}{\cp{FNP}}

\newcommand\NDP{\cp{NDP}}
\newcommand\coNDP{\cp{coNDP}}
\newcommand\NDPs[1]{\cp{NDP}_{#1}}

% -- coNP
\newcommand\coNP{\cp{coNP}}

% -- PSPACE
\newcommand\PSPACE{\cp{ PSPACE}}
\newcommand\NPSPACE{\cp{ NPSPACE}}
\newcommand{\Pspace}{\PSPACE}
\newcommand{\FPspace}{\cp{FPSPACE}}


% -- LOGSPACE
\newcommand\LSPACE{\cp{LOGSPACE}}
\newcommand\NLSPACE{\cp{NLOGSPACE}}
\newcommand\coNLSPACE{\cp{coNLOGSPACE}}
\newcommand\coNL{\cp{coNL}}

% -- Autres
\newcommand\EXPTIME{\cp{EXPTIME}}
\newcommand\NEXPTIME{\cp{NEXPTIME}}
\newcommand\EXPSPACE{\cp{EXPSPACE}}
\newcommand\PH{\cp{PH}}
\newcommand\NC{\cp{NC}}
\newcommand\AC{\cp{AC}}
\newcommand\PPAD{\cp{\mathbf{PPAD}}}
\newcommand\PLS{\cp{\mathbf{PLS}}}

% -- Reverse math
\newcommand{\WKL}{\ensuremath{\docheck{\mathsf{WKL}_0}}}
\newcommand{\ACA}{\ensuremath{\docheck{\mathsf{ACA}_0}}}
\newcommand{\ATR}{\ensuremath{\docheck{\mathsf{ATR}_0}}}
\newcommand{\RCAO}{\ensuremath{\docheck{\mathsf{RCA}_0}}}
\newcommand{\PiCA}{\ensuremath{\docheck{\Pi^1_1\!\text{-}\mathsf{CA}_0}}}


% -- circuits
\newcommand\CIRCUIT[2]{\cp{CIRCUIT(#1,#2)}}
\newcommand\UCIRCUIT[2]{\cp{UCIRCUIT(#1,#2)}}

% -- Classes proba
\newcommand\PP{\cp{PP}}
\newcommand\RPP{\cp{RP}}
\newcommand\ZPP{\cp{ZPP}}
\newcommand\coRP{\cp{coRP}}
\newcommand\BPP{{\cp{BPP}}}

% -- IP
\newcommand\IP{\cp{IP}}
\newcommand\PCP{\cp{PCP}}

% -- non-uniforme
\newcommand\nuP{\mathbb{P}}
\newcommand\nuPs[1]{\mathbb{P}_{#1}}
\newcommand\nuPsp[1]{\mathbb{P}_{#1}^0}
\newcommand\nuNP{\mathbb{N}\mathbb{P}}
%\newcommand\poly{/{poly}}
\newcommand\Ppoly{\PTIME_{\poly}}
\newcommand\NPpoly{\NP_{\poly}}

% -- conditions d'uniformité
\newcommand\Luniforme{\cp{L}}
\newcommand\BCuniforme{\cp{BC}}


% Classes de circuits
\newcommand\TC{\cp{TC}}
\newcommand\LFP{\cp{LFP}}
\newcommand\FAC{\cp{FAC}}
\newcommand\FACC{\cp{FACC}}
\newcommand\FTC{\cp{FTC}}
\newcommand\FNC{\cp{FNC}}

% Logiques temorelles
\newcommand\CTL{\cp{CTL}}
\newcommand\LTL{\cp{LTL}}

% -- mesure temps/espace/taille

\newcommand\DTIME{\cp{DTIME}}
\newcommand\TIME[1]{\cp{TIME(#1)}}
\newcommand\TIMEs[2]{\cp{TIME_{#2}(#1)}}
\newcommand\TIMEsp[2]{\cp{TIME^0_{#2}(#1)}}
\newcommand\NTIME[1]{\cp{NTIME(#1)}}
\newcommand\NTIMEs[2]{\cp{NTIME_{#2}(#1)}}
\newcommand\NTIMEsp[2]{\cp{NTIME_{#2}^0(#1)}}
\newcommand\DSPACE{\cp{DSPACE}}
\newcommand\SPACE[1]{\cp{SPACE(#1)}}
\newcommand\SPACEs[2]{\cp{SPACE_{#2}(#1)}}
\newcommand\NSPACE[1]{\cp{NSPACE(#1)}}
\newcommand\NSPACEs[2]{\cp{NSPACE_{#2}(#1)}}
\newcommand\coNSPACE[1]{\cp{coNSPACE(#1)}}
\newcommand\TIMEON[2]{\cp{TIME(#1,#2)}}
\newcommand\SPACEON[2]{\cp{SPACE(#1,#2)}}
\newcommand\SIZE[1]{\cp{ size(#1)}}
\newcommand\ATIME{\cp{ATIME}}
\newcommand\ASPACE{\cp{ASPACE}}


%% autre:
\newcommand\LH{\cp{LH}}
\newcommand\FO{\cp{FO}}
\newcommand\SO{\cp{SO}}
\newcommand\ESOhorn{\cp{ESO_{horn}}}
\newcommand\SOhorn{\cp{SO_{horn}}}
\newcommand\ACzero{\cp{AC_0}}
\newcommand{\LINSPACE}{\cp{LINSPACE}}
\newcommand{\mFLINSPACE}{\mF_{\cp{LINSPACE}}}
\newcommand{\LOGSPACE}{\cp{LOGSPACE}}
\newcommand\ALOGTIME{\cp{ALOGTIME}}
\newcommand\RF{\cp{RF}}


\newcommand\BOOL[1]{\cp{{BOOLEAN(#1)}}}
\newcommand\DET{\cp{ DET}}

% % -- dirty


\newcommand{\cPspace}{\cp{\sharp PSPACE}}
\newcommand{\cAPtime}{\cp{\sharp APTIME}}
%\newcommand{\FPspace}{\cp{FPSPACE}}
%\newcommand{\FPspaceN}{\cp{FPSPACE}^{+}}
%\newcommand{\FPspaceN}{\cp{FPSPACE}}



%-------------------------
%--- Problèmes de décision  ---
%--------------------------

\newcommand\decisionname[1]{\cp{#1}}


\newcommand\Luniv{\decisionname{L_{univ}}}
\newcommand\HALTINGPROBLEM{\decisionname{HALTING-PROBLEM}}
\newcommand\SAT{\decisionname{SAT}}
\newcommand\MAXtSAT{\decisionname{MAX3SAT}}
\newcommand\REACH{\decisionname{REACH}}
\newcommand\CIRCUITVALUE{\decisionname{CIRCUITVALUE}}
\newcommand\MONOTONECIRCUITVALUE{\decisionname{MONOTONECIRCUITVALUE}}
\newcommand\THEOREMS{\decisionname{THEOREMS}}
\newcommand\QCIRCUITSAT{\decisionname{QCIRCUITSAT}}
\newcommand\QSAT{\decisionname{QSAT}}
\newcommand\QBF{\decisionname{QBF}}
\newcommand\SATVALUE{\decisionname{SATVALUE}}
\newcommand\SUCCINTSAT{\decisionname{SUCCINTSAT}}
\newcommand\SUCCINTSATVALUE{\decisionname{SUCCINTSATVALUE}}
\newcommand\GEOGRAPHY{\decisionname{GEOGRAPHY}}
\newcommand\PARITY{\decisionname{PARITY}}
\newcommand\MAJ{\decisionname{MAJ}}
\newcommand\CVP{\CIRCUITVALUE}
\newcommand\BOOLEANPROGRAMVALUE{BOOLEAN~PROGRAM~VALUE}

\newcommand\CLIQUE{\cp{CLIQUE}}
\LANGUEinv{
\newcommand\RECOUVREMENTDESOMMETS{\cp{RECOUVREMENT~DE~SOMMETS}}
\newcommand\RS{\RECOUVREMENTDESOMMETS}
}{
\newcommand\RECOUVREMENTDESOMMETSeng{\cp{VERTEX~COVER}}
\newcommand\RSeng{\RECOUVREMENTDESOMMETSeng}
}

\LANGUEinv{\newcommand\COUPUREMAXIMALE{\cp{COUPURE~MAXIMALE}}}
{\newcommand\COUPUREMAXIMALEeng{\cp{MAXIMAL~CUT}}}
\newcommand\CIRCUITHAMILTONIEN{\CHAMILTONIEN}
\LANGUEinv{\newcommand\VOYAGEURDECOMMERCE{\VCOMMERCE}}
{\newcommand\VOYAGEURDECOMMERCEeng{\VCOMMERCEeng}}

\LANGUEinv{\newcommand\CIRCUITLEPLUSLONG{\cp{CIRCUIT~LE~PLUS~LONG}}}
{\newcommand\CIRCUITLEPLUSLONGeng{\cp{LONGEST~CIRCUIT}}}
\LANGUEinv{\newcommand\SOMMEDESOUSENSEMBLE{\SUBSETSUM}}{
\newcommand\SOMMEDESOUSENSEMBLEeng{\SUBSETSUMeng}}
\LANGUEinv{
\newcommand\SACADOS{\cp{SAC~A~DOS}}}
{\newcommand\SACADOSeng{\cp{KNAPSACK}}}

%français
\LANGUEinv{
\newcommand\CHAMILTONIEN{\cp{CIRCUIT~HAMILTONIEN}}}
{\newcommand\CHAMILTONIENeng{\cp{HAMILTONIAN~CIRCUIT}}}
\LANGUEinv{
\newcommand\VCOMMERCE{\cp{VOYAGEUR~DE~COMMERCE}}}
{\newcommand\VCOMMERCEeng{\cp{TRAVELING~SALESMAN}}}
\LANGUEinv{
\newcommand\kCOLORABILITE{\cp{k-COLORABILITE}}}
{\newcommand\kCOLORABILITEeng{\cp{k-COLORABILITY}}}
\LANGUEinv{
\newcommand\COLORIAGE{\cp{BON~COLORIAGE}}}
{\newcommand\COLORIAGEeng{\cp{GOOD~COLORING}}}
\LANGUEinv{
\newcommand\tCOLORABILITE{\cp{3-COLORABILITE}}}
{\newcommand\tCOLORABILITEeng{\cp{3-COLORABILITY}}}
\LANGUEinv{
\newcommand\SUBSETSUM{\cp{SOMME~DE~SOUS~ENSEMBLE}}}
{\newcommand\SUBSETSUMeng{\cp{SUBSET~SUM}}}
\newcommand\PARTITION{\cp{PARTITION}}
\LANGUEinv{
\newcommand\NOMBREPREMIER{\decisionname{NOMBRE~ PREMIER}}}
{\newcommand\NOMBREPREMIEReng{\decisionname{PRIME~NUMBER}}}
\LANGUEinv{
\newcommand\CODAGE{\decisionname{CODAGE}}}
{\newcommand\CODAGEeng{\decisionname{ENCODING}}}
\newcommand\kSAT{\cp{k-SAT}}
\newcommand\tSAT{\cp{3-SAT}}
\newcommand\NAESAT{\cp{NAESAT}}
\newcommand\NAEtSAT{\cp{NAE3SAT}}
\newcommand\NAEqSAT{\cp{NAE4SAT}}
\newcommand\STABLE{\cp{\LANGUE{INDEPENDANT~SET}{STABLE}}}

%\newcommand\CM{\cp{COUPURE MAXIMALE}}
\LANGUEinv{
\newcommand\POSTpb{\cp{Problème~de~ la~correspondance~de~Post}}}
{\newcommand\POSTpbeng{\cp{Post~correspondence~problem}}}

\LANGUEinv{
\newcommand\PARITE{\cp{PARITE}}}
{\newcommand\PARITEeng{\cp{PARITY}}}


%-------------------------
%--- Modèles parallèles  ---
%--------------------------


% PARALLELISME
\newcommand\RAM{\cp{RAM}}
%
\newcommand\PRAM{\cp{PRAM}}
\newcommand\CRCW{\cp{CRCW}}
\newcommand\CREW{\cp{CREW}}
\newcommand\EREW{\cp{EREW}}


%------------------------
%--- codages  ---
%------------------------


\newcommand\codage[1]{\langle #1 \rangle}
\newcommand\tuple[1]{\codage{#1}}
\newcommand\paire[2]{\codage{#1,#2}}
\newcommand\triplet[3]{\codage{#1,#2,#3}}
\newcommand\quadruplet[4]{\codage{#1,#2,#3,#4}}
\newcommand\cinquplet[5]{\codage{#1,#2,#3,#4,#5}}


%------------------------
%--- descriptive set theory  ---
%------------------------

\newcommand\baire{\mathcal{N}}
\newcommand\peano{\overline{\mathcal{N}}} %% A MOI
\newcommand\cantor{\mathcal{C}}
%
\newcommand\bSigma{\mathbf{\Sigma}}
\newcommand\bPi{\mathbf{\Pi}}
\newcommand\bDelta{\mathbf{\Delta}}
%


%%%%%% ABOUT NOMS QUI CHANGES

\newcommand\WF{\cp{WF}}
\newcommand\WO{\WF}
\newcommand\LO{\cp{LO}}
\newcommand\DLO{\cp{DLO}}
% 
\newcommand\rk{\cp{rank}}


%% Moins propre:

\newcommand\I{\mathbb{I}}
\renewcommand\H{\mathbb{H}}
%
\newcommand\leKB{<_{KB}}
%
\newcommand\finiomega{{<\omega}}
\newcommand\compl{^\frown}
\newcommand\prefix{ ~ prefix~  } 
\newcommand\X{{X}}
\newcommand\Y{\mathcal{Y}}
\newcommand\subsetstrict{\subset}

\renewcommand\complement[1]{#1^{c}}


%------------------------
%--- schémas de récursion ---
%------------------------

%%                      Source papier MFCS 2019


% COMPLEXITY

%% MFCS differe de Clote:
%% Version à taille restreinte

\newcommand\co{\cp{co}}

\newcommand\mFPTIME{\mF_{PTIME}}
\newcommand{\mFPSPACE}{\mF_{SPACE}}
\newcommand{\mFPH}{\mF_{PH}}

\newcommand\SigmaP{\Sigma^P}
\newcommand\DeltaP{\Delta^{P}}
\newcommand\PiP{\Pi^{P}}
\newcommand\coSigmaP{\co\SigmaP}
\newcommand\coPiP{\mF{\co\PiP}}

%% autres classes:


%% Schemas pour complexité implicite

\newcommand\COMP{\cp{ COMP}}
\newcommand\CRN{\cp{ CRN}}
\newcommand\BRN{\cp{ BRN}}
\newcommand\SBRN{\cp{ SBRN}}
\newcommand\BR{\cp{ BR}}
\newcommand\kBR{k-\cp{ BR}}
\newcommand \WBPR{\cp{ WBPR}}
\newcommand \BMIN{\cp{ BMIN}}
\newcommand \BSUM{\cp{ BSUM}}
\newcommand\SBSUM{\cp{ SBSUM}}
\newcommand \BPROD{\cp{ BPROD}}
\newcommand \SBPROD{\cp{ SBPROD}}
\newcommand \SCOMP{\cp{ SCOMP}}
\newcommand \SR{\cp{ SR}}
\newcommand \SRN{\cp{ SRN}}

%% devrait etre défini
\newcommand\WBRN{\cp{ WBRN}}



%-- godel
\newcommand\INTDIV{\cp{INTDIV}}
\newcommand\DIV{\cp{DIV}}
\newcommand\EVEN{\cp{EVEN}}
\newcommand\ODD{\cp{ODD}}
\newcommand\PRIME{\cp{PRIME}}
\newcommand\POWER{\cp{POWER}}
\newcommand\BIT{\cp{BIT}}
\newcommand\negHP{\overline{\cp{HP}}}
\newcommand\HP{\cp{HP}}
\newcommand\negLuniv{\overline{\Luniv}}
\newcommand\VALCOMP{\cp{VALCOMP}}
\newcommand\mWINDOW{\cp{LEGAL}}
\newcommand\mmWINDOW{\cp{WINDOW}}
\newcommand\LENGTH{\cp{LENGTH}}
\newcommand\DIGIT{\cp{DIGIT}}
\newcommand\MATCH{\cp{MATCH}}
\newcommand\MOVE{\cp{MOVE}}
\newcommand\START{\cp{START}}
\newcommand\CELL{\cp{CELL}}
\newcommand\HALT{\cp{HALT}}
\newcommand\SUBST{\cp{SUBST}}
\newcommand\PROOF{\cp{PROOF}}
\newcommand\PROVABLE{\cp{PROVABLE}}
\newcommand\CONSIST{\cp{CONSIST}}



%---------------------
%--- calculus (AP) ---
%---------------------

% \jacobian{f}{x}
\newcommand{\jacobian}[1]{J_{#1}}
%\newcommand{\grad}[1]{\overrightarrow{\operatorname{grad}}~{#1}}
\newcommand{\grad}[1]{\nabla{#1}}
\newcommand{\scalarprod}[2]{{#1}\cdot{#2}}
\newcommand{\bigO}[1]{\mathcal{O}\left(#1\right)}
\newcommand\smallO{o}
\newcommand{\softO}[1]{\tilde{\mathcal{O}}\left(#1\right)}
\newcommand{\taylor}[3]{{T_{#1}^{#2}#3}}
\newcommand{\transpose}[1]{{#1}^T}
\newcommand{\pastsup}[2]{{\sup}_{#1}#2}

\DeclareMathOperator\arctanh{arctanh}



%---------------
%--- norms (AP) ---
%---------------


\newcommand{\inorm}[2]{\left\lVert{#1}\right\rVert_{#2}}
\newcommand{\twonorm}[1]{\inorm{#1}{2}}
\newcommand{\infnorm}[1]{\inorm{#1}{}}
\newcommand{\normsup}[1]{\infnorm{#1}{}}
%
\newcommand\hausdorff{d_H}
%
\newcommand\norm[1]{\infnorm{#1}}


%----------------
%--- topology ---
%----------------

% \openball{pt}{radius}
 \newcommand{\openball}[2]{B_{#2}(#1)}
\newcommand{\closedball}[2]{\ovl{B}_{#2}(#1)}
\newcommand{\closure}[1]{\ovl{#1}}
% \newcommand{\interior}[1]{\mathring{#1}}
% \newcommand{\border}[1]{\partial{#1}}




%-----------------
%--- functions (AP) ---
%-----------------
 
\newcommand{\lambdafun}[2]{(#1)\mapsto{#2}}
\newcommand{\lambdafunex}[2]{#1\mapsto{#2}}
\newcommand{\argmin}{\operatornamewithlimits{argmin}}
\newcommand{\argmax}{\operatornamewithlimits{argmax}}


 
%-----------------------
%--- Machines de Turing  ---
%-----------------------

\newcommand\blanc{\vectorl{B}}



%-----------------
%--- Logique ---
%-----------------
 
\newcommand\sem[1]{{[\semspace[}#1{]\semspace]}}
\newcommand\sems[1]{{[\semspace[}#1{]\semspace]}}
\newcommand\semss[1]{\sem{#1}_{\sigma'}}
\newcommand\semspace{\kern -.25ex}
\newcommand\true{true}
\newcommand\false{false}
%extension élémentaire:
\newcommand\elem{\preceq}
\newcommand\godel[1]{\lceil #1 \rceil}

%-----------------
%--- Model theory ---
%-----------------

\newcommand\hull[1]{\langle #1 \rangle}

%-----------------
%--- Ordinaux ---
%-----------------

\newcommand\omegaunCK{\omega_1^{CK}}
\newcommand\omegaunck{\omegaunCK}

\DeclareMathOperator{\cof}{cof}

%-----------------
%--- Surréels (QG) ---
%-----------------

\newcommand{\On}{\textbf{On}}
\newcommand{\Nobf}{\textbf{No}}
\newcommand{\Nobb}{\ensuremath{\mathbbmss{No}}}
\newcommand{\Nolt}[1]{\ensuremath{\Nobf_{< #1}}}
\newcommand{\Noleq}[1]{\ensuremath{\Nobf_{\leq #1}}}
\newcommand{\Noltbb}[1]{\ensuremath{\Nobb_{< #1}}}
\providecommand\No{\Nobf}
\newcommand\Noalpha{\Nolt{\alpha}}
\newcommand\Nolambda{\Nolt{\lambda}}

\DeclareMathOperator{\length}{length}
\DeclareMathOperator{\supp}{supp} % operateur support


%----------------------
%---Maths  de base (AP)  ---
%----------------------

\DeclareMathOperator{\dom}{dom} % operateur support

\DeclareMathOperator{\size}{size}

\newcommand{\powerset}[1]{\mathcal{P}(#1)}
\newcommand{\card}[1]{{\##1}}
\newcommand\priv{ - }
\newcommand\prive{-}
\renewcommand{\restriction}{\mathord{\upharpoonright}}
\newcommand\restrict[1]{_{\restriction #1}}

% \providecommand\mod{}
%% UNDEFINE:
\let\mod\relax
\let\div\relax
\DeclareMathOperator{\mod}{mod} 
\DeclareMathOperator{\div}{div} 

% partieentier
\newcommand\entieres[1]{\lceil #1 \rceil}
\newcommand\entierei[1]{\lfloor #1 \rfloor}


%----------------------
%---Symbole danger ---
%----------------------


%\providecommand*{\TakeFourierOrnament}[1]{{%
%\fontencoding{U}\fontfamily{futs}\selectfont\char#1}}
%\providecommand*{\danger}{{\Huge \TakeFourierOrnament{66}}}



%%%%%%%%%%%%%%%%%%%%%%%%%%%%%%%%%%%%%%%%%%%%%%%%%
%
%
%                      ASMeries
%
%
%%%%%%%%%%%%%%%%%%%%%%%%%%%%%%%%%%%%%%%%%%%%%%%%%



% CIRCUITS
%-- particuliers
\newcommand\REPLACE{\cp{REPLACE}}
\newcommand\EQUAL{\cp{EQUAL}}
\newcommand\EVALUE{\cp{EVALUE}}
\newcommand\TVALUE{\cp{TVALUE}}
\newcommand\UPDATE{\cp{UPDATE}}
\newcommand\ASSIGN{\cp{ASSIGN}}
\newcommand\PAR{\cp{PAR}}
\newcommand\DO{\cp{DO}}


% ASM
\newcommand\emplacement[2]{(#1,#2)}
\newcommand\contenu[2]{\sem{(#1,#2)}}
\newcommand\affecte[3]{(#1,#2)\leftarrow#3}
\newcommand\ctl{ctl\_state}


%%%%%%%%%%%%%%%%%%%%%%%%%%%%%%%%%%%%%%%%%%%%%%%%%
%
%
%                      GPACqueries
%
%
%%%%%%%%%%%%%%%%%%%%%%%%%%%%%%%%%%%%%%%%%%%%%%%%%



%------------------------
%--- generable zoo (GPAC) ---
%------------------------

\newcommand{\myop}[1]{\operatorname{#1}}
\newcommand{\abs}{\myop{abs}}
\newcommand{\sg}{\myop{sg}}
\newcommand{\ip}[1]{\myop{ip}_{#1}}
\newcommand{\rnd}{\myop{rnd}}
\newcommand{\lil}{\myop{lil}}
\newcommand{\lxh}{\myop{lxh}}
\newcommand{\hxl}{\myop{hxl}}
\newcommand{\plil}{\myop{plil}}
\newcommand{\reach}{\myop{reach}}
%\newcommand{\mreach}{\myop{mreach}}
\newcommand{\watchdog}{\myop{watchdog}}
\newcommand{\monitor}{\myop{monitor}}
%\newcommand{\norm}{\myop{norm}}
\newcommand{\mx}{\myop{mx}}
\newcommand{\mn}{\myop{mn}}
\newcommand{\sample}{\myop{sample}}
\newcommand{\select}{\myop{select}}
\newcommand{\ssine}[1]{\sin_{#1}}
\newcommand{\clamp}{\myop{clamp}}

\newcommand{\fnsg}[3]{tanh((#1)*(#2)*(#3))}
\newcommand{\fnipone}[3]{(1+\fnsg{(#1)-1}{#2}{#3})/2.}
\newcommand{\fnipinf}[1]{#1-sin(2*pi*(#1))/2./pi}%
\newcommand{\fnlil}[5]{%
\fnipone{(#3)-(#1)+1-((#2)-(#1))/8.}{#4+log(1+4./(#2-#1)*(#5)*(#5))+log(1+4)}{8./(#2-#1)}%
*\fnipone{#2-#3+1-(#2-#1)/8.}{#4+log(1+4/(#2-#1)*(#5)*(#5))+log(1+4)}{8./(#2-#1)}%
*4./(#2-#1)*(#5)}
\newcommand{\fnlxh}[5]{\fnipone{(#3)-(#1+#2)/2.+1}{#4+log(1+(#5)*(#5))}{2./(#2-#1)}*(#5)}
\newcommand{\fnhxl}[5]{\fnipone{(#1+#2)/2.-(#3)+1}{#4+log(1+(#5)*(#5))}{2./(#2-#1)}*(#5)}
\newcommand{\fnplilf}[4]{sin(((#4)-((#1)+(#2))/2.+(#3)/4.)*2.*pi/(#3))}
\newcommand{\fnplil}[6]{(#6)*\fnlxh{\fnplilf{#1}{#2}{#3}{#1}}%
{\fnplilf{#1}{#2}{#3}{#1+((#2)-(#1))/4.}}%
{\fnplilf{#1}{#2}{#3}{#4}}%
{#5+2.+log(1+(#6)*(#6))}%
{1./4.+2./((#2)-(#1))}}
\newcommand{\fnssine}[2]{sin(#2)*(1-cos(#2)/cos(#1))}



%--------------------
%--- complexity GPAC ---
%--------------------


%\ifthenelse{\equal{#1}{}}{Empty.}{Nonempty.}
\newcommand{\myclass}[1]{\operatorname{#1}}
% \newcommand{\gcomp}[2]{\ensuremath{\myclass{GCOMP}(#1,#2)}}
% \newcommand{\ggenerable}[1]{\textcolor{red}{\ensuremath{#1-}\text{generable}}}
 \newcommand{\gval}[2][]{\ensuremath{\myclass{GVAL}_{#1}\ifthenelse{\equal{#2}{}}{}{[#2]}}}
 \newcommand{\gpval}[1][]{\ensuremath{\myclass{GPVAL}_{#1}}}
 \newcommand{\gc}[3][]{\ensuremath{\myclass{ATSC}(#2,#3)}}
 \newcommand{\gpc}[1][]{\ensuremath{\myclass{ATSP}}}
\newcommand{\cgc}[1][]{\ensuremath{\myclass{ATS}}}
\newcommand{\gexpc}[1][]{\ensuremath{\myclass{AEXP}}}
\newcommand{\gwc}[2]{\ensuremath{\myclass{AW}(#1,#2)}}
\newcommand{\gpwc}{\ensuremath{\myclass{AWP}}}
\newcommand{\cgwc}{\ensuremath{\myclass{AW}}}
\newcommand{\grc}[3]{\ensuremath{\myclass{AR}(#1,#2,#3)}}
\newcommand{\gprc}{\ensuremath{\myclass{ARP}}}
\newcommand{\gsc}[3]{\ensuremath{\myclass{AS}(#1,#2,#3)}}
\newcommand{\gpsc}{\ensuremath{\myclass{ASP}}}
\newcommand{\goc}[3]{\ensuremath{\myclass{AO}(#1,#2,#3)}}
\newcommand{\gpoc}{\ensuremath{\myclass{AOP}}}
\newcommand{\cgoc}{\ensuremath{\myclass{AO}}}
\newcommand{\guc}[4]{\ensuremath{\myclass{AX}(#1,#2,#3,#4)}}
\newcommand{\gpuc}{\ensuremath{\myclass{AXP}}}
\newcommand{\glc}[2][]{\ensuremath{\myclass{AL}(#2)}}
\newcommand{\gplc}[1][]{\ensuremath{\myclass{ALP}}}
\newcommand{\cglc}[1][]{\ensuremath{\myclass{AL}}}

%\newcommand{\intinterv}[2]{\{#1,\ldots,#2\}}
\newcommand{\intinterv}[2]{\llbracket#1,#2\rrbracket}

\newcommand{\Rgen}{\R_{G}}
\newcommand{\Rpoly}{\R_{P}}

%%%%%%%%%%%%%%%%%%%%%%%%%%%%%%%%%%%%%%%%%%%%%%%%%
%
%
%                      Mes codages obscures
%
%
%%%%%%%%%%%%%%%%%%%%%%%%%%%%%%%%%%%%%%%%%%%%%%%%%




%%% Codages

% \newcommand\gammaoc{\gamma_{[0,1]}}
% \newcommand\gammaonc{\gamma_{\N}}
\newcommand\gammaTMc{\gamma^{TM}_{[0,1]}}
\newcommand\gammaTMcatan{\gamma^{TM}_{(0,1)^2}}
\newcommand\gammaTMnc{\gamma^{TM}_{\N^3}}
\newcommand\gammaTMncd{\gamma^{TM}_{\N^2}}
% \newcommand\gammaTMe{\gamma^{TM}_{*}}
% \newcommand\gammas{\gamma_{sab}}
% \newcommand\gammaord{\gamma_{cantor}}


\newcommand\enccompact{\nu_{X}}
 \newcommand\encinfinite{\nu_\N}
 \newcommand\enc{\nu}
 
%%%%%%%%%%%%%%%%%%%%%%%%%%%%%%%%%%%%%%%%%%%%%%%%%
%
%
%                      DESSINS ET IMAGES
%
%
%%%%%%%%%%%%%%%%%%%%%%%%%%%%%%%%%%%%%%%%%%%%%%%%%


\newcommand\IMAGE[2]{ \begin{center} 
    \includegraphics[#1]{#2}
  \end{center}
  }


  
%%%% POUR FAIRE DESSIN.
\newenvironment{dessinpp}[1]{
\begin{tikzpicture}[baseline=(current bounding
      box.west),#1]
}
{\end{tikzpicture}}




\newenvironment{dessinp}[1]{
\begin{center}\begin{tikzpicture}[baseline=(current bounding
      box.north),#1]
}
{\end{tikzpicture}
\end{center}}


\newenvironment{dessin}{
\begin{center}\begin{tikzpicture}[baseline=(current bounding
      box.north)]
}
{\end{tikzpicture}
\end{center}}



\newenvironment{dessind}{
\begin{tikzpicture}[baseline=(current bounding
      box.north)]
}
{\end{tikzpicture}
}



%%%%%%%%%%%%%%%%%%%%%%%%%%%%%%%%%%%%%%%%%%%%%%%%%
%
%
%                      POUR DESSINS GPAC
%
%
%%%%%%%%%%%%%%%%%%%%%%%%%%%%%%%%%%%%%%%%%%%%%%%%%


%%%% MACROS DE BASE.


\newcommand\constant[2]
{
node (#1)  [shape=circle,draw] {#2};
\draw (#1.east) -- +(0.3,0) coordinate (#1-o) ; 
\draw (#1.west) -- +(-0.3,0) coordinate (#1-i) ; 
}

\newcommand\basicproduct[2]
{
node (#1) {#2};
\draw (#1) +(-0.4,0.4) -- +(0.4,0.4) -- +(0.8,0) -- +(0.4,-0.4) -- +(-0.4,-0.4) --
+(-0.4,0.4);
\draw (#1) ++(0.8,0) -- ++(+0.3,0) coordinate (#1-o) ;
\draw (#1) ++ (-0.4,0.2) -- +(-0.3,0) coordinate (#1-i1) ;
\draw (#1) ++(-0.4,-0.2) -- +(-0.3,0) coordinate (#1-i2) ;
}

\newcommand\product[1]{\basicproduct{#1}{$+\Pi$}}
\newcommand\negativeproduct[1]{\basicproduct{#1}{$-\Pi$}}


%%BEGIN INTERNAL
\newcommand\basicsummer[1]{
coordinate (#1) {};
\draw (#1) +(-0.5,0.6) -- +(0.5,0) -- +(-0.5,-0.6) -- +(-0.5,0.6);
\draw (#1) ++(0.5,0) -- +(+0.3,0) coordinate (#1-o) ;
}

\newcommand\basicintegrator[1]{
coordinate (#1) {};
\draw (#1) +(-0.5,0.6) -- +(0.5,0) -- +(-0.5,-0.6) -- +(-0.5,0.6);
\draw (#1) ++(0.5,0) -- +(+0.3,0) coordinate (#1-o) ;
\draw (#1) +(-0.5,-0.6) -| +(-0.8,0.6) -- +(-0.5,0.6) ;
 \draw (#1) ++(-0.65,0.6) -- +(0,+0.3) coordinate (#1-init) ;
}

\newcommand\splitfour[2]{
\draw (#2) ++ (#1,0.4) -- +(-0.3,0) coordinate (#2-i1) ;
\draw (#2) ++ (#1,0.15) -- +(-0.3,0) coordinate (#2-i2) ;
\draw (#2) ++ (#1,-0.15) -- +(-0.3,0) coordinate (#2-i3) ;
\draw (#2) ++ (#1,-0.4) -- +(-0.3,0) coordinate (#2-i4) ;
}

\newcommand\splitthree[2]{
\draw  (#2) ++(#1,0.3) -- +(-0.3,0) coordinate (#2-i1) ;
\draw  (#2) ++(#1,0) -- +(-0.3,0) coordinate (#2-i2) ;
\draw  (#2) ++(#1,-0.3) -- +(-0.3,0) coordinate (#2-i3) ;
}

\newcommand\splittwo[2]{
\draw  (#2) ++(#1,0.2) -- +(-0.3,0) coordinate (#2-i1) ;
\draw  (#2) ++(#1,-0.2) -- +(-0.3,0) coordinate (#2-i2) ;
}

\newcommand\splitone[2]{
\draw (#2) ++(#1,0) -- +(-0.3,0) coordinate (#2-i1) ;
}

%% END INTERNAL

\newcommand\summerfour[1]{
\basicsummer{#1}
\splitfour{-0.5}{#1}
}

\newcommand\summerthree[1]{
\basicsummer{#1}
\splitthree{-0.5}{#1}
}

\newcommand\summertwo[1]{
\basicsummer{#1}
\splittwo{-0.5}{#1}
}

\newcommand\summerone[1]{
\basicsummer{#1}
\splitone{-0.5}{#1}
}


\newcommand\integratorfour[1]{
\basicintegrator{#1}
\splitfour{-0.8}{#1}
<}


\newcommand\integratorthree[1]{
\basicintegrator{#1}
\splitthree{-0.8}{#1}
}


\newcommand\integratortwo[1]{
\basicintegrator{#1}
\splittwo{-0.8}{#1}
}


\newcommand\integratorone[1]{
\basicintegrator{#1}
\splitone{-0.8}{#1}
}

%%%%% FIN MACROS




%%%%%%%%%%%%%%%%%%%%%%%%%%%%%%%%%%%%%%%%%%%%%%%%%
%
%
%                      Trucs d'Amaury
%
%
%%%%%%%%%%%%%%%%%%%%%%%%%%%%%%%%%%%%%%%%%%%%%%%%%


%----------------
%--- ref (AP) ---
%----------------

% \newcommand{\defref}[1]{Definition~\ref{#1}}
\newcommand{\lemref}[1]{Lemma~\ref{#1}}
\newcommand{\thref}[1]{Theorem~\ref{#1}}
\newcommand{\amaurycorref}[1]{Corollary~\ref{#1}}
 \newcommand{\figref}[1]{Figure~\ref{#1}}
% \newcommand{\figrefmult}[1]{Figures~{#1}}
% \newcommand{\secref}[1]{Section~\ref{#1}}
% %\renewcommand{\algref}[1]{Algorithm~\ref{#1}}
 \newcommand{\propref}[1]{Proposition~\ref{#1}}
% \newcommand{\exref}[1]{Example~\ref{#1}}
\newcommand{\remref}[1]{Remark~\ref{#1}}


 %-------------------
%--- polynomials ---
%-------------------

\newcommand{\sigmap}[1]{{\Sigma{#1}}}
\newcommand{\degp}[1]{{\operatorname{deg}(#1)}}
\newcommand{\poly}{\operatorname{poly}}


 %----------------------
%--- misc functions ---
%----------------------

\newcommand{\sgn}[1]{\operatorname{sgn}(#1)}
\newcommand{\intp}{\operatorname{int}}
 \newcommand{\fracp}{\operatorname{frac}}
\newcommand{\fiter}[2]{#1^{[#2]}}
\newcommand{\frestrict}[2]{#1\restriction_{#2}}
\newcommand{\indicator}[1]{\mathds{1}_{#1}}
\newcommand{\idfun}{\operatorname{id}}
% \newcommand{\floor}[1]{\left\lfloor#1\right\rfloor}
 \newcommand{\ceil}[1]{\left\lceil#1\right\rceil}
\newcommand{\round}[1]{\left\lfloor#1\right\rceil}
\DeclareMathOperator\erf{erf}


%------------
%--- sets ---
%------------

\newcommand{\A}{\mathbb{A}}
%\newcommand{\D}{\mathbb{D}}
\newcommand{\Ns}{\mathbb{N}^*}
%\newcommand{\C}{\mathbb{C}}
%\newcommand{\K}{\mathbb{K}}
% \MAT{height}{[width]}{[field]}
\newcommand{\MAT}[3]{M_{#1\ifthenelse{\equal{#2}{}}{}{,#2}}\ifthenelse{\equal{#3}{}}{}{\left(#3\right)}}

%\newcommand{\Rcomp}{\R_{c}}
\newcommand{\Rp}{\R_{+}}
\newcommand{\Rs}{\R^{*}}
\newcommand{\Rps}{\R_{+}^{*}}

%\newcommand{\Analytic}{C^\omega}
%\newcommand{\PTIME}{\ensuremath{\operatorname{P}}}
%\newcommand{\EXPTIME}{\ensuremath{\operatorname{EXPTIME}}}
%\newcommand{\NP}{\ensuremath{\operatorname{NP}}}
%\newcommand{\NPTIME}{\NP}
\newcommand{\FP}{\ensuremath{\operatorname{FP}}}
%\newcommand{\PSPACE}{\ensuremath{\operatorname{PSPACE}}}
% \newcommand{\FPSPACE}{\ensuremath{\operatorname{FPSPACE}}}



%%%%%%%%%%%%%%%%%%%%%%%%%%%%%%%%%%%%%%%%%%%%%%%%%
%
%
%                      Trucs de Quentin
%
%
%%%%%%%%%%%%%%%%%%%%%%%%%%%%%%%%%%%%%%%%%%%%%%%%%


\newcommand{\fct}[5]{\ensuremath{#1 : \left\{\begin{array}{c c c}
#2 & \rightarrow & #3\\
#4 & \mapsto & #5
                                             \end{array}\right.}}

%Nouvelle fraction, plus jolie
\newcommand{\f}[2]{	
	\mathchoice%
		{\dfrac{#1}{#2}}
    	{\dfrac{#1}{#2}}
		{\frac{#1}{#2}}
		{\frac{#1}{#2}}
}


%%Ensembles et combinatoires

%Ensemble avec propriété à drotie \{x | blabla\}
\newcommand{\enstq}[2]{\left\{ \left.#1\ \vphantom{#2}\right|\ #2  \right\}}
%meme chose mais avec des crochets
\newcommand{\crotq}[2]{\left[ \left.#1\ \vphantom{#2}\right|\ #2  \right]}
%même chose mais avec des brackets
\newcommand{\bratq}[2]{\left\langle \left.#1\ \vphantom{#2}\right|\ #2  
  \right\rangle}

\newcommand{\intn}[2]{\ensuremath{
		\left\llbracket\, #1\ ;\ #2\,\right\rrbracket
              }}


%%%%%%%%%%%%%%%%%%%%%%%%%%%%%%%%%%%%%%%%%%%%%%%%%
%
%
%               ENVIRONNEMENTS THEOREMS & CO
%
%
%%%%%%%%%%%%%%%%%%%%%%%%%%%%%%%%%%%%%%%%%%%%%%%%%


%----------------
%--- pas dans lipics ---
%-----------------


\ifnotSLIDE{
% commente car bug Fevrier 2026: \ifnotTDEXAM{\spnewtheorem{solution}{Solution}{\bfseries}{\rmfamily}}
\spnewtheorem{openproblem}{Open Problem}{\bfseries}{\rmfamily}
\spnewtheorem{remarkolivier}{Commentaire d'Olivier: }{\bfseries}{\rmfamily}
\spnewtheorem{postulate}{Postulate}{\bfseries}{\rmfamily}
}

%----------------
%--- dans lipics ---
%-----------------
%
\newenvironment{exercice}[1]{\begin{exercise} \label{exo:#1-stt}}{\end{exercise}}
\newenvironment{exerciced}[1]{\begin{exercisee}\label{exo:#1-stt}}{\end{exercisee}}
\newenvironment{exercicec}[1]{\begin{exercise} \label{exo:#1-stt} (\LANGUE{solution on}{corrigé} page \pageref{exo:#1-cor})
}{\end{exercise}}
\newenvironment{exercicedc}[1]{\begin{exercisee} \label{exo:#1-stt} (\LANGUE{solution on}{corrigé} page \pageref{exo:#1-cor})
}{\end{exercisee}}

\newenvironment{slideproposition}{{\bf Proposition}}{}

\ifnotSLIDE{
\nolipics{
\nolncs{\spnewtheorem{theorem}{\LANGUE{Theorem}{Théorème}}{\bfseries}{\rmfamily}}
\nolncs{\spnewtheoremi{definition}{\LANGUE{Definition}{Définition}}{\bfseries}{\rmfamily}{theorem}}
\nolncs{\spnewtheoremi{lemma}{\LANGUE{Lemma}{Lemme}}{\bfseries}{\rmfamily}{theorem}}
\nolncs{\spnewtheoremi{corollary}{\LANGUE{Corollary}{Corollaire}}{\bfseries}{\rmfamily}{theorem}}
\nolncs{\spnewtheoremi{proposition}{Proposition}{\bfseries}{\rmfamily}{theorem}}
\nolncs{\spnewtheoremi{example}{\LANGUE{Example}{Exemple}}{\bfseries}{\rmfamily}{theorem}}
%\spnewtheorem{sketch}{{\bf \LANGUE{Sketch of proof}{Démonstration (principe)}}}{\bfseries}{\rmfamily}
\nolncs{\spnewtheoremi{remark}{\LANGUE{Remark}{Remarque}}{\bfseries}{\rmfamily}{theorem}}
\nolncs{\spnewtheorem{exercise}{\LANGUE{Exercise}{Exercice}}{\bfseries}{\rmfamily}}
\spnewtheorem{exercisee}{\LANGUE{*Exercise}{*Exercice}}{\bfseries}{\rmfamily}
\spnewtheorem{summary}{\LANGUE{Summary}{Résumé}}{\bfseries}{\rmfamily}
 %\spnewtheorem{exerciceenglish}{Exercise}{\bfseries}{\rmfamily}
 %
 \newenvironment{sketch}{{\bf \LANGUE{Sketch of proof}{Démonstration (principe)}}:}{\hfill $\Box$ }
\nolncs{\spnewtheorem{conjecture}{\LANGUE{Conjecture}{Conjecture}}{\bfseries}{\rmfamily}}
 %%
 \spnewtheorem{theoremaprouver}{\LANGUE{Theorem TO BE PROVED}{Théorème (A PROUVER)}}{\bfseries}{\rmfamily}
  \spnewtheorem{conjectureaprouver}{\LANGUE{Conjecture TO BE VERIFIED}{Conjecture (A VERIFIER)}}{\bfseries}{\rmfamily}
 }}
            
 \makeatletter
 \ifnotSLIDE{
  \ifnotTDEXAM{
\@ifundefined{proof}{
	\newenvironment{proof}{{\bf \LANGUE{Proof}{Démonstration}}:}{\hfill $\Box$ 
	}{}
	}}}
\makeatother

 
%%%%%%%%%%%%%%%%%%%%%%%%%%%%%%%%%%%%%%%%%%%%%%%%%
%
%
%                      Macros dont pas fier du tout
%
%
%%%%%%%%%%%%%%%%%%%%%%%%%%%%%%%%%%%%%%%%%%%%%%%%%



\newcommand\equations[1]{
\begin{array}{lll}
#1
\end{array}
}



\newcommand\systeme[1]{
\left\{
\begin{array}{lll} 
#1
\end{array}
\right.
}



%------------------
%--- shorthands (AP)---
%------------------
\newcommand{\mtt}[1]{\mathtt{#1}}
\newcommand{\ovl}[1]{\overline{#1}}
\newcommand{\panic}[1]{\textcolor{red}{#1}}
%\NewEnviron{epanic}{{\color{red}\BODY}}
\newcommand{\idest}{\emph{i.e.\thinspace}}


%------------------
%--- mes trucs  ---
%------------------


\newcommand\implique{\Rightarrow}
\newcommand\ssi{\Leftrightarrow}
\newcommand\ac[1]{\{#1\}}
\newcommand\zero{\vectorl{0}}
\newcommand\un{\vectorl{1}}

\newcommand\technique[1]{{\scriptsize #1}}


%% Preuve and Co

\newcommand\sensdirect{ Direction $\Rightarrow$. }
\newcommand\sensindirect{ Direction $\Leftarrow$. }




\newcommand\donnee{\mbox{<donnée>}}



% pour aller vite

\newcommand\gM{\mathfrak{M}}
\newcommand\gMM{(M,f_1,\cdots,f_u,r_1,\cdots,r_v)}
%
\newcommand\mygM{\mK}
\newcommand\mygMM{\left ( \K, \{op_i\}_{i\in I}, r_1 \ldots,
r_\ell, \zero,\un \right )}
\newcommand\myM{\K} 

\newcommand\gMMM{(M,f_1,\cdots,f_u,c_1,\cdots,c_k,r_1,\cdots,r_v)}
\newcommand\gMMMM{(f_1,\cdots,f_u,r_1,\cdots,r_v)}
\newcommand\gMME{(M,f_1,\cdots,f_u,f'_1,\cdots,f'_\ell,r_1,\cdots,r_v)}
\newcommand\gMMS{(f_1,\cdots,f_u,f'_1,\cdots,f'_\ell,r_1,\cdots,r_v)}

\newcommand\bool{\mathfrak{B}}
\newcommand\gbool{(\{0,1\},\zero,\un,=)}

\newcommand\osg{\overline{sg}}
\newcommand\moins{\mathrel{\dot{-}}}


%%% TRUC DE BARWISE


\newcommand\UM{\kU_{\kM}}
\newcommand\VV{\mathbb{V}}
% Hereditarily finite
\newcommand\IHF{\mathbb{I}HF}
\newcommand\IHP{\IHYP}
\newcommand\IHYP{\mathbb{I}HYP}

\providecommand\citep[1]{\cite{#1}}

%% Cofinali


% probas
% défini dans amsmath
\renewcommand\Pr{\cp{Pr}}
\newcommand\Var{\cp{Var}}


% -- mesures
%\newcommand\card{card}
\newcommand\taille{\LANGUE{size}{taille}}
\newcommand\entree{e}
\newcommand\profondeur{\LANGUE{depth}{profondeur}}


%%%%%%%%%%%%%%%%%%%%%%%%%%%%%%%%%%%%%%%%%%%%%%%%%%%%%%%%%%%%%%%%%%%%%%%%%%%%
%                 MACHINES DE TURING (graphique)
%%%%%%%%%%%%%%%%%%%%%%%%%%%%%%%%%%%%%%%%%%%%%%%%%%%%%%%%%%%%%%%%%%%%%%%%%%%%



% Configuration de machine de Turing
\newcommand\configTuring[3]{#1\textbf{#2}#3}


%%% DESSIN CROIX PORU AIDER  A DESSINER

\newcommand\croix{  +(-1,-1) -- +(1,1) +(-1,1) --
  +(1,-1) }

%% DESSIN ACCOLADE

\newcommand\ACCOLADE[3]{
%#1: coordonnées de départ
%#2: coordonnées arrivée
%#3: texte
\draw[decorate,decoration={brace}]  (#1) -- (#2)
 node[midway,above=0.1cm] {#3};
}

%%%


\def\lcase{3.5 ex}
\def\hcase{3.5ex}
\tikzstyle{case} = []
%[draw,minimum width = \lcase,minimum height = 3.6ex,baseline]
\newcommand\dcase{ +(-0.5*\lcase,-0.5*\hcase) rectangle +(0.5*\lcase,0.5*\hcase)}
%\newcommand\lettre[1]{\node [name=l,case,on chain=1] {#1};}


\newcommand\RUBANMACHINEDETURING[6]{
%#1: texte
%#2: position de la tête
%#3: nb de blancs à gauche
%#4: nb de cases totales
%#5: texte au dessus du ruban
%#6: position du coin
    \draw (#6) +(\lcase,4ex ) node {#5};
    \draw (#6) +(-0.6,0)  node  {\ldots};  
    \draw (#6) node[coordinate] (t) {};
    \foreach \x in {1,2,...,#3} {
         \draw (t) node [case]   {$\blanc$} \dcase;
         \draw (t) ++(\lcase,0) node [coordinate] (t) {};
       };
   \StrLen{#1}[\Resultat] % ATTENTION IL FAUT CETTE SYNTAXE POUR FAIRE
                         % EQUIVALENT D UN EDEF
  
\foreach \x in {1,...,\Resultat} 
             { 
        \draw (t) node [case]   {\StrChar{#1}{\x}} \dcase;
        \draw (t) ++(\lcase,0) node [coordinate] (t) {};
             }
\pgfmathtruncatemacro{\z}{#4 - \Resultat - #3}

  \foreach \x in {1,2,..., \z} {
        \draw (t) node [case]   {$\blanc$} \dcase;
        \draw (t) ++(\lcase,0) node [coordinate] (t) {};
}
%% Convention: graphiquement, la case numéro n se trouve à (n*\lcase,0)
%% en comptant les case à partir de $0$

   \draw  (t) +(0.1,0) node [name=r] {\ldots}; 
}

\newcommand\RUBANMACHINEDETURINGtexte[6]{
%#1: texte
%#2: position de la tête
%#3: nb de blancs à gauche
%#4: nb de cases totales (équivalent)
%#5: texte au dessus du ruban
%#6: position du coin
    \draw (#6) +(\lcase,4ex ) node {#5};
    \draw (#6) +(-0.6,0)  node  {\ldots};  
    \draw (#6) node[coordinate] (t) {};
   \foreach \x in {1,2,...,#3} {
        \draw (t) node [case]   {$\blanc$} \dcase;
        \draw (t) ++(\lcase,0) node [coordinate] (t) {};
      };
   \draw (#6) ++(-2*\lcase,0.5*\lcase)-- +(\lcase*#4,0);
   \draw (#6) ++(-2*\lcase,-0.5*\lcase) -- +(\lcase*#4,0);
   
    \draw (t) ++(-0.5*\lcase,0) node [case,anchor=west]  (te) {#1};
    \draw (te.east)  node [anchor=west,coordinate] (t) {};

  \draw (t) ++(0.2*\lcase,0) node [coordinate] (t) {};
   \foreach \x in {1,2,...,#3} {
        \draw (t) node [case]   {$\blanc$} \dcase;
        \draw (t) ++(\lcase,0) node [coordinate] (t) {};
      };

    \draw  (t) +(0.1,0) node [anchor=west] {\ldots}; 
}


\newcommand\RUBANMACHINEDETURINGTETE[6]{
%
% dessine la tete
%
% output: (tete)
%
\RUBANMACHINEDETURING{#1}{#2}{#3}{#4}{#5}{#6}
% TETE
     \draw (#6) + ( #2 * \lcase + #3  *\lcase,- 0.5*\lcase )
     node[coordinate] (k) {};
     \draw[->] (k) +(0,-0.5) -- (k);
     \draw (k) +(0,-0.5) node[coordinate] (tete) {};
}


\newcommand\RUBANMACHINEDETURINGTETEtexte[6]{
%
% dessine la tete
%
% output: (tete)
%
\RUBANMACHINEDETURINGtexte{#1}{#2}{#3}{#4}{#5}{#6}
% TETE
     \draw (#6) + ( #2 * \lcase + #3  *\lcase,- 0.5*\lcase )
     node[coordinate] (k) {};
     \draw[->] (k) +(0,-0.5) -- (k);
     \draw (k) +(0,-0.5) node[coordinate] (tete) {};
}


\newcommand\MACHINEDETURING[6]{
%#1: texte
%#2: position de la tête
%#3: état
%#4: nb de blancs à gauche
%#5: nb de cases totales
%#6: texte en haut du ruban
%\begin{tikzpicture}[node distance=-0.15mm]

\RUBANMACHINEDETURINGTETE{#1}{#2}{#4}{#5}{#6}{0,0}
%
%% ETAT:
     \draw (tete) node [name=etat,draw,circle,anchor=north]  {#3};
     \draw (etat) -- (tete);
%\end{tikzpicture}
}

\newcommand\MACHINEDETURINGtexte[6]{
%#1: texte
%#2: position de la tête
%#3: état
%#4: nb de blancs à gauche
%#5: nb de cases totales
%#6: texte en haut du ruban
%\begin{tikzpicture}[node distance=-0.15mm]

\RUBANMACHINEDETURINGTETEtexte{#1}{#2}{#4}{#5}{#6}{0,0}
%
%% ETAT:
     \draw (tete) node [name=etat,draw,circle,anchor=north]  {#3};
     \draw (etat) -- (tete);
%\end{tikzpicture}
}

\newcommand\MACHINEDETURINGTROISRUBANS[9]{
%#1: texte1
%#2: position de la tête 1
%#3: texte2
%#4: position de la tête 2
%#5: texte3
%#6: position de la tête 3
%#7: état
%#8: nb de blancs à gauche 1
%#9: nb de cases totales
%\begin{tikzpicture}[node distance=-0.15mm]

\RUBANMACHINEDETURINGTETE{#1}{#2}{#8}{#9}{\LANGUE{tape}{ruban} $1$}{0,0}
\draw (tete) node[coordinate] (tete1) {};
\draw (tete1) +(8,0) node[coordinate] (tete1b) {};
\draw (tete1) -| (tete1b) ;

\RUBANMACHINEDETURINGTETE{#3}{#4}{#8}{#9}{\LANGUE{tape}{ruban} $2$}{0,-2}
\draw (tete) node[coordinate] (tete2) {};
\draw (tete2) +(-8.5,0) node[coordinate] (tete2b) {};
\draw (tete2) -| (tete2b) ;

\RUBANMACHINEDETURINGTETE{#5}{#6}{#8}{#9}{\LANGUE{tape}{ruban} $3$}{0,-4}
\draw (tete) node[coordinate] (tete3) {};

%% ETAT:
     \draw (tete3) node [name=etat,draw,circle,anchor=north]  {#7};
     \draw (etat) -- (tete3);
    \draw  (etat) -|  (tete2b);
\draw  (etat) -|  (tete1b);
%\end{tikzpicture}
}

\newcommand\MACHINEDETURINGTROISRUBANSENGLISH[9]{
%#1: texte1
%#2: position de la tête 1
%#3: texte2
%#4: position de la tête 2
%#5: texte3
%#6: position de la tête 3
%#7: état
%#8: nb de blancs à gauche 1
%#9: nb de cases totales
%\begin{tikzpicture}[node distance=-0.15mm]

\RUBANMACHINEDETURINGTETE{#1}{#2}{#8}{#9}{tape $1$}{0,0}
\draw (tete) node[coordinate] (tete1) {};
\draw (tete1) +(8,0) node[coordinate] (tete1b) {};
\draw (tete1) -| (tete1b) ;

\RUBANMACHINEDETURINGTETE{#3}{#4}{#8}{#9}{tape $2$}{0,-2}
\draw (tete) node[coordinate] (tete2) {};
\draw (tete2) +(-8.5,0) node[coordinate] (tete2b) {};
\draw (tete2) -| (tete2b) ;

\RUBANMACHINEDETURINGTETE{#5}{#6}{#8}{#9}{tape $3$}{0,-4}
\draw (tete) node[coordinate] (tete3) {};

%% ETAT:
     \draw (tete3) node [name=etat,draw,circle,anchor=north]  {#7};
     \draw (etat) -- (tete3);
    \draw  (etat) -|  (tete2b);
\draw  (etat) -|  (tete1b);
%\end{tikzpicture}
}



\newcommand\MACHINEDETURINGTROISRUBANStexte[9]{
%#1: texte1
%#2: position de la tête 1
%#3: texte2
%#4: position de la tête 2
%#5: texte3
%#6: position de la tête 3
%#7: état
%#8: nb de blancs à gauche 1
%#9: nb de cases totales
% \begin{tikzpicture}[node distance=-0.15mm]

\RUBANMACHINEDETURINGTETEtexte{#1}{#2}{#8}{#9}{\LANGUE{tape}{ruban} $1$}{0,0}
\draw (tete) node[coordinate] (tete1) {};
\draw (tete1) +(8,0) node[coordinate] (tete1b) {};
\draw (tete1) -| (tete1b) ;

\RUBANMACHINEDETURINGTETEtexte{#3}{#4}{#8}{#9}{\LANGUE{tape}{ruban} $2$}{0,-2}
\draw (tete) node[coordinate] (tete2) {};
\draw (tete2) +(-8.5,0) node[coordinate] (tete2b) {};
\draw (tete2) -| (tete2b) ;

\RUBANMACHINEDETURINGTETEtexte{#5}{#6}{#8}{#9}{\LANGUE{tape}{ruban} $3$}{0,-4}
\draw (tete) node[coordinate] (tete3) {};

%% ETAT:
     \draw (tete3) node [name=etat,draw,circle,anchor=north]  {#7};
     \draw (etat) -- (tete3);
    \draw  (etat) -|  (tete2b);
\draw  (etat) -|  (tete1b);
%\end{tikzpicture}
}


\newcommand\MACHINEDETURINGDEUXRUBANStexte[9]{
%#1: texte1  
%#2: position de la tête 1 
%#3: texte2 inutilisée
%#4: position de la tête 2 inutilisée
%#5: texte3
%#6: position de la tête 3
%#7: état
%#8: nb de blancs à gauche 1
%#9: nb de cases totales
% \begin{tikzpicture}[node distance=-0.15mm]

\RUBANMACHINEDETURINGTETEtexte{#1}{#2}{#8}{#9}{\LANGUE{tape}{ruban} $1$}{0,0}
\draw (tete) node[coordinate] (tete1) {};
\draw (tete1) +(8,0) node[coordinate] (tete1b) {};
\draw (tete1) -| (tete1b) ;

% \RUBANMACHINEDETURINGTETEtexte{#3}{#4}{#8}{#9}{\LANGUE{tape}{ruban} $2$}{0,-1.5}
% \draw (tete) node[coordinate] (tete2) {};
% \draw (tete2) +(-5.5,0) node[coordinate] (tete2b) {};
% \draw (tete2) -| (tete2b) ;

\RUBANMACHINEDETURINGTETEtexte{#5}{#6}{#8}{#9}{\LANGUE{tape}{ruban} $2$}{0,-1.5}
\draw (tete) node[coordinate] (tete3) {};

%% ETAT:
     \draw (tete3) node [name=etat,draw,circle,anchor=north]  {#7};
     \draw (etat) -- (tete3);
%   \draw  (etat) -|  (tete2b);
\draw  (etat) -|  (tete1b);
%\end{tikzpicture}
}


\def\argdix{2}
\def\argonze{25}
\newcommand\MACHINEDETURINGQUATRERUBANStextep[9]{
%#1: texte1
%#2: position de la tête 1
%#3: texte2
%#4: position de la tête 2
%#5: texte3
%#6: position de la tête 3
%#7: texte4
%#8: position de la tête 3
%#9: état
%#10: nb de blancs à gauche 1
%#11: nb de cases totales
% \begin{tikzpicture}[node distance=-0.15mm]

%% BUG sur #10 devenu \argdix
%% BUG SUR #11 devenu \argonze

\RUBANMACHINEDETURINGTETEtexte{#1}{#2}{\argdix}{\argonze}{\LANGUE{tape}{ruban} $1$}{0,0}
% \draw (tete) node[coordinate] (tete1) {};
% \draw (tete1) +(8,0) node[coordinate] (tete1b) {};
% \draw (tete1) -| (tete1b) ;

\RUBANMACHINEDETURINGTETEtexte{#3}{#4}{\argdix}{\argonze}{\LANGUE{tape}{ruban} $2$}{0,-1.5}
% \draw (tete) node[coordinate] (tete2) {};
% \draw (tete2) +(-5.5,0) node[coordinate] (tete2b) {};
% \draw (tete2) -| (tete2b) ;

\RUBANMACHINEDETURINGTETEtexte{#5}{#6}{\argdix}{\argonze}{\LANGUE{tape}{ruban} $3$}{0,-3.2}
\draw (tete) node[coordinate] (tete3) {};

\RUBANMACHINEDETURINGTETEtexte{#7}{#8}{\argdix}{\argonze}{\LANGUE{tape}{ruban} $4$}{0,-4.9}
\draw (tete) node[coordinate] (tete4) {};

%% ETAT:
     \draw (tete4) node [name=etat,draw,circle,anchor=north]  {#9};
%      \draw (etat) -- (tete3);
%     \draw  (etat) -|  (tete2b);
% \draw  (etat) -|  (tete1b);
%\end{tikzpicture}
}


%%%%%%%%%%%%%%%%%%%%%%%%%%%%%%%%%%%%%%%%%%%%%%%%%%%%%%%%%%%%%%%%%%%%%%%%%%%%
%                 MACHINES DE TURING (automate)
%%%%%%%%%%%%%%%%%%%%%%%%%%%%%%%%%%%%%%%%%%%%%%%%%%%%%%%%%%%%%%%%%%%%%%%%%%%%



\newcommand\dec{1ex}
\newcommand\myACCOLADE[4]{
\ACCOLADE{#1*\lcase+#4*\lcase-0.5*\lcase,0.5*\lcase+\dec}{#1*\lcase+#4*\lcase-0.5*\lcase+#3*\lcase,0.5*\lcase+\dec}{#2}
}


\newcommand\M{\Sigma}


\newcommand\sommet[2]{\node[state] (#1) #2 {$#1$};}
\newcommand\sommetinitial[2]{\node[state,initial] (#1) #2 {$#1$};}
\newcommand\sommetfinal[2]{\node[state,accepting] (#1) #2 {$#1$};}
\newcommand\gtransition[5]{\draw (#1) edge node {#2/#4 $#5$} (#3);}
\newcommand\gtransitionba[5]{\draw (#1) [loop above] edge node {#2/#4
    $#5$} (#3);}
\newcommand\gtransitionbb[5]{\draw (#1) [loop below] edge node {#2/#4
    $#5$} (#3);}
\newcommand\gtransitionbleft[5]{\draw (#1) [loop left] edge node {#2/#4
    $#5$} (#3);}
\newcommand\gtransitionbright[5]{\draw (#1) [loop right] edge node {#2/#4
    $#5$} (#3);}
\newcommand\gtransitionbl[5]{\draw (#1) [bend left] edge node {#2/#4
    $#5$} (#3);}
\newcommand\gtransitionbr[5]{\draw (#1) [bend right] edge node {#2/#4
    $#5$} (#3);}
\newcommand\gtransitionbadouble[9]{\draw (#1) [loop above] edge node {
\begin{tabular}{l}
#2/#4 $#5$ \\
  #7/#8 $#9$ \\
\end{tabular}
} (#3);}

\newcommand\SCALE{}

\newcommand\dessinmtp[2]{
\begin{center}
\begin{tikzpicture}[->,>=stealth',shorten >=1pt,auto,node distance=2.3cm,semithick,#2] 
#1
\end{tikzpicture}
\end{center}
}

\newcommand\dessinmt[1]{\dessinmtp{#1}{scale=1}}


\newcommand\ANIMATIONMT[1]{
\begin{overprint}
\begin{dessinp}{scale=0.52}
\foreach \av/\q/\b  [count=\xi]
in  
{#1}
{ 
\only<\xi| handout 0>{ 
 \StrLen{\av}[\Resintm] % ATTENTION IL FAUT CETTE SYNTAXE POUR FAIRE
                         % EQUIVALENT D UN EDEF
 \MACHINEDETURING{\av \b}{\Resintm}{$\q$}{5}{20}{}
}}
 \end{dessinp}
\end{overprint}
}

\newcommand\ANIMATIONMTd[1]{
\begin{overprint}
\begin{dessinp}{scale=0.52}
\foreach \av/\q/\b  [count=\xi]
in  
{#1}
{ 
\only<\xi| handout 0>{ 
 \StrLen{\av}[\Resintm] % ATTENTION IL FAUT CETTE SYNTAXE POUR FAIRE
                         % EQUIVALENT D UN EDEF
 \hspace{-1cm}\MACHINEDETURING{\av \b}{\Resintm}{$\q$}{5}{33}{}
}}
 \end{dessinp}
\end{overprint}
}

\def\MAXDET{20}

\newcommand\DIAGRAMMEESPACETEMPS[1]{
\draw node[coordinate] (tt) {};
\foreach \av/\q/\b  in  
{#1}
{ 

\RUBANMACHINEDETURING{\av{$\q$}\b}{0}{4}{\MAXDET}{}{tt};
\draw (tt) +(0,-\hcase) node[coordinate] (tt) {};
}
%\foreach \x in {0,...,19} {\draw (tt) + (\x*\lcase,0) node {$\vdots$};}
}



\newcommand\progun{
/q_0/0011, % initial
X/q_1/011,X0/q_1/11,X/q_2/0Y1,/q_2/X0Y1,
X/q_0/0Y1,XX/q_1/Y1,XXY/q_1/1, XX/q_2/YY, X/q_2/XYY,
XX/q_0/YY,XXY/q_3/Y,XXYY/q_3/B,XXYYB/q_4/B
}

\newcommand\DIAGRAMMEESPACETEMPSun{
\DIAGRAMMEESPACETEMPS
%% RECOPIE progun
{/q_0/0011, % initial
X/q_1/011,X0/q_1/11,X/q_2/0Y1,/q_2/X0Y1,
X/q_0/0Y1,XX/q_1/Y1,XXY/q_1/1, XX/q_2/YY, X/q_2/XYY,
XX/q_0/YY,XXY/q_3/Y,XXYY/q_3/B,XXYYB/q_4/B}
%% FIN RECOPIE
}

\newcommand\progdeux{
/q_0/0010, X/q_1/010, X0/q_1/10, X/q_2/0Y0,/q_2/X0Y0,
X/q_0/0Y0,XX/q_1/Y0,XXY/q_1/0,XXY0/q_1/%B,
}


\newcommand\DIAGRAMMEESPACETEMPSdeux{
\DIAGRAMMEESPACETEMPS
{
%% RECOPIE progdeux
/q_0/0010, X/q_1/010, X0/q_1/10, X/q_2/0Y0,/q_2/X0Y0,
X/q_0/0Y0,XX/q_1/Y0,XXY/q_1/0,XXY0/q_1/%B,
%% FIN RECOPIE
}
}

\newcommand\DIAGRAMMEESPACETEMPSdeuxvariante{
\def\memolcase{\lcase}
\def\memoMAXDET{20}
\def\MAXDET{15}
\def\lcase{6 ex}
\DIAGRAMMEESPACETEMPS
{
%% RECOPIE progdeux
/q_00/010, X/q_10/10, X0/q_11/0, X/q_20/Y0,/q_2X/0Y0,
X/q_00/Y0,XX/q_1Y/0,XXY/q_10/,XXY0/q_1B/%B,
%% FIN RECOPIE
}
\def\lcase{\memolcase}
\def\MAXDET{\memoMAXDET}
}


\newcommand\ANIMATIONun{
\ANIMATIONMT
%% RECOPIE progun
{/q_0/0011, % initial
X/q_1/011,X0/q_1/11,X/q_2/0Y1,/q_2/X0Y1,
X/q_0/0Y1,XX/q_1/Y1,XXY/q_1/1, XX/q_2/YY, X/q_2/XYY,
XX/q_0/YY,XXY/q_3/Y,XXYY/q_3/B,XXYYB/q_4/B}
%% FIN RECOPIE
}



\newcommand\DESSINALGO[3]{
% argument 1= ou
% argument 2= texte 
% argument 3= nom du noeud ex m
\draw node[draw,#1,minimum size=1cm] (#3) {#2};

% pour fitting box
\draw (#3.west) +(-0.2cm,0) node (#3c2) {};
\draw (#3.north) +(0,+0.2cm) node (#3c3) {};
\draw (#3.south) +(0,-0.2cm) node (#3c4) {};
}

\newcommand\DESSINALGOa[3]{
\DESSINALGO{#1}{#2}{#3}
\draw (#3.east) +(-0.1cm,0.2cm) node(#3a) {};
\draw[->] (#3a) -- ++(0.5cm,0) node (#3aa) {};
\draw (#3aa) node[anchor=west] (#3aaa) {\LANGUE{Accept}{Accepte}};
\draw (#3aaa.east) node(#3aaaa) {};
}

\newcommand\DESSINALGOe[4]{
\DESSINALGO{#1}{#2}{#3}
\draw (#3.east) +(-0.1cm,0.2cm) node(#3a) {};
\draw[->] (#3a) -- ++(0.5cm,0) node (#3aa) {};
\draw (#3aa) node[anchor=west] (#3aaa) {#4};
\draw (#3aaa.east) node(#3aaaa) {};
}


\newcommand\DESSINALGOar[3]{
\DESSINALGOa{#1}{#2}{#3}
\draw (#3.east) +(-0.1cm,-0.2cm) node(#3r) {};
\draw[->] (#3r) -- ++(0.5cm,0) node (#3rr) {};
\draw (#3rr) node[anchor=west] (#3rrr) {\LANGUE{Reject}{Refuse}};
\draw (#3rrr.east) node(#3rrrr) {};
}


\newcommand\WINDOWarg[7] {
%#1: position
%#2..7: contenu du tableau
\draw (#1) node[case]  {#2} \dcase;
\draw (#1) ++ (\lcase,0) node[case] {#3} \dcase;
\draw (#1) ++ (2*\lcase,0) node[case] {#4} \dcase;
\draw (#1) ++ (0,-\lcase) node[case] {#5} \dcase;
\draw (#1) ++ (\lcase,-\lcase) node[case] {#6} \dcase;
\draw (#1) ++ (2*\lcase,-\lcase) node[case] {#7} \dcase;
}

\newcommand\WINDOW[3]{
%#1: positioin
%#2: contenu
%#3: étiquette
\WINDOWarg{#1}
{\StrChar{#2}{1}}{\StrChar{#2}{2}}{\StrChar{#2}{3}}
{\StrChar{#2}{4}}{\StrChar{#2}{5}}{\StrChar{#2}{6}}
\draw (#1) +(-\lcase*1.2,-\lcase*0.5) node {#3};
}





%%%%%%%%%%%%%%%%%%%%%%%%%%%%%%%%%%%%%%%%%%%%%%%%%%%%%%%%%%%%%%%%%%%%%%%%%%%%
%                 LISTINGS
%%%%%%%%%%%%%%%%%%%%%%%%%%%%%%%%%%%%%%%%%%%%%%%%%%%%%%%%%%%%%%%%%%%%%%%%%%%%


% LISTINGS
\RequirePackage{listings}
 \lstset{language=Java,
 mathescape=true,
 showspaces=false,
 showstringspaces=false,
 frame=single,
morekeywords={recursive,then,div,function,integer,read,out,else,par,endpar,seq,endseq,and,not,read,write,repeat,until,undef,let,to,endfor,endif,endwhile,guess,in,parallel},
% basicstyle=\ttfamily\small,
% basewidth={1.1ex},
% keywordstyle=\bf\color{blue},
% %stringstyle=\bf\ttfamily\color{green},
% commentstyle=\bfseries\color{magenta},
 literate={:=}{{$\gets$ }}1 {<=}{{$\leq$}}2  {>=}{{$\geq$}}2
 {!=}{{$\neq$}}1 
 }
% \lstset{escapechar=\%}
\let\lst\lstinline
\def\beginlisting{\begin{lstlisting}}
\def\endlisting{\end{lstlisting}}  




\newenvironment{algo}[1]{
  \lstset{escapechar=\@}
  \def\input{$#1$}
  \def\out{out}
  \def\BEGIN{\textbf{read} \input }
 \def\BEGINN{\textbf{repeat} }
  \def\END{\textbf{until out!=undef}}
  \def\ENDD{\textbf{write \out}}
  \def\ENDT{\textbf{until $\ctl$ = $q_a$ or $\ctl$ = $q_r$}}
  \def\ENDR{\textbf{until $\ctl$ = $q_a$}}
  \def\ENDP{\textbf{until $\ctl$ = $q_a$}}
}{}
\newenvironment{algon}[1]{
  \lstset{numbers=left}
  \begin{algo}{#1}
}{\end{algo}}
\newcommand\FSM[3]{FSM(#1,\lst{#2},#3)}
\newcommand\FSMi[3]{FSM(#1,{#2},#3)}




%%%%%%%%%%%%%%%%%%%%%%%%%%%%%%%%%%%%%%%%%%%%%%%%%
%
%
%                      TRADUCTION EN COURS
%
%%%%%%%%%%%%%%%%%%%%%%%%%%%%%%%%%%%%%%%%%%%%%%%%%

\newcommand\ENTREPOT[1]{\NDLR{ENTREPOT ICI: #1}}
\newcommand\scite[2][]{\cite[#1]{#2}}
\newcommand\nospellcheck[1]{{#1}}
\newcommand\nd{nd}
\newcommand\NL{
 
}

\newenvironment{FRANCAIS}{
\small
\color{orange}
}
{\normalsize
\color{black}
}


\newenvironment{SEMIFRANCAIS}{
\small
\color{violet}
}
{\normalsize
\color{black}
}

\newcommand\QFRANCAIS[1]{
\begin{FRANCAIS}
#1
\end{FRANCAIS}
}


%%%%%%%%%%%%%%%%%%%%%%%%%%%%%%%%%%%%%%%%%%%%%%%%%
%
%
%                      Wikipedia
%
%%%%%%%%%%%%%%%%%%%%%%%%%%%%%%%%%%%%%%%%%%%%%%%%%

\providecommand\tightlist{}
\providecommand\Complex{\C}
\newenvironment{myalgie}{\begin{array}{lll}}{\end{array}}

\newcommand\nl{\ \\ }




%%%%%%%%%%%%%%%%%%%%%%%%%%%%%%%%%%%%%%%%%%%%%%%%%
%
%
%                      Pour citations précises
%
%%%%%%%%%%%%%%%%%%%%%%%%%%%%%%%%%%%%%%%%%%%%%%%%%



\newcommand\citeprecis[2]{{\cite[#1]{#2}}}
\newcommand\citeprecisinv[2]{\citeprecis{#2}{#1}}

\usepackage{csquotes}



%%%%%%%%%%%%%%%%%%%%%%%%%%%%%%%%%%%%%%%%%%%%%%%%%
%
%
%                      Travail sur couleurs
%
%   toutes les réponses pointent dans le même sens, et le schéma correspond exactement 
%. au profil perceptif d’un protanope** (déficit des cônes L / longueurs d’onde longues)
%  MAIS SEVERE
%
%%%%%%%%%%%%%%%%%%%%%%%%%%%%%%%%%%%%%%%%%%%%%%%%%


% ================================
% GROUPE 1 — Verts sombres
% ================================
\definecolor{deepgreen}{RGB}{0,160,40}
\definecolor{forestgreen}{RGB}{0,150,70}
\definecolor{mossgreen}{RGB}{0,140,90}
\definecolor{greenwave}{RGB}{0,150,90}
\definecolor{forestmist}{RGB}{0,160,90}
\definecolor{darkteal}{RGB}{0,130,70}

% ================================
% GROUPE 2 — Teal / verts bleutés
% ================================
\definecolor{seagreen}{RGB}{0,130,100}
\definecolor{freshteal}{RGB}{0,140,130}
\definecolor{greenteal}{RGB}{0,140,120}
\definecolor{blueteal}{RGB}{0,150,120}      % ancien "teal" (doublon corrigé)
\definecolor{lagoon}{RGB}{0,180,120}
\definecolor{greenmist}{RGB}{0,170,110}

% ================================
% GROUPE 3 — Turquoises / Aqua
% ================================
\definecolor{springgreen}{RGB}{0,190,100}
\definecolor{mintgreen}{RGB}{10,200,150}
\definecolor{oceanfoam}{RGB}{0,200,150}
\definecolor{aquagreen}{RGB}{30,210,170}
\definecolor{coastalaqua}{RGB}{40,210,180}
\definecolor{softaqua}{RGB}{40,210,190}

% ================================
% GROUPE 4 — Cyan / Aqua clairs
% ================================
\definecolor{emeraldteal}{RGB}{0,180,140}
\definecolor{deepcyan}{RGB}{0,180,160}
\definecolor{lightcyan}{RGB}{0,200,160}
\definecolor{aquamarine}{RGB}{0,200,160}
\definecolor{lightaqua}{RGB}{60,230,200}
\definecolor{mistblue}{RGB}{80,240,210}

% ================================
% GROUPE 5 — Bleus clairs
% ================================
\definecolor{paleaqua}{RGB}{80,200,190}
\definecolor{glacier}{RGB}{90,150,200}
\definecolor{skyblueA}{RGB}{60,150,220}      % ancien "skyblue"
\definecolor{lightsky}{RGB}{90,170,230}
\definecolor{iceblue}{RGB}{120,190,240}
\definecolor{frost}{RGB}{150,210,250}

% ================================
% GROUPE 6 — Bleus moyens
% ================================
\definecolor{steelblue}{RGB}{30,100,160}
\definecolor{cadetblue}{RGB}{70,120,180}
\definecolor{coolblue}{RGB}{110,170,220}
\definecolor{morningblue}{RGB}{140,200,240}
\definecolor{blueiceA}{RGB}{40,150,210}       % ancien "blueice" (doublon corrigé)
\definecolor{blueair}{RGB}{60,170,225}

% ================================
% GROUPE 7 — Azurs / Bleu saturé
% ================================
\definecolor{tealblue}{RGB}{0,110,130}
\definecolor{deepbluecyan}{RGB}{0,110,170}
\definecolor{azur}{RGB}{0,70,180}
\definecolor{brightazure}{RGB}{20,130,200}
\definecolor{lightazure}{RGB}{80,180,235}
\definecolor{skylight}{RGB}{100,200,245}

% ================================
% GROUPE 8 — Bleus sombres / violets
% ================================
\definecolor{nightblue}{RGB}{40,0,110}
\definecolor{slateblue}{RGB}{20,40,120}
\definecolor{violetA}{RGB}{60,50,160}         % ancien "violet" (distinct de violetB ci-dessous)
\definecolor{blueviolet}{RGB}{80,60,180}
\definecolor{purpleA}{RGB}{100,70,200}        % ancien "purple"
\definecolor{lightpurple}{RGB}{130,90,210}

% ================================
% GROUPE 9 — Violets saturés / indigos
% ================================
\definecolor{lilac}{RGB}{160,110,220}
\definecolor{indigoA}{RGB}{50,0,130}          % ancien "indigo"
\definecolor{deepindigo}{RGB}{70,10,150}
\definecolor{royalpurple}{RGB}{90,20,170}

% ================================
% GROUPE 10 — Fermeture cercle perceptif
% ================================
\definecolor{coolmint}{RGB}{0,130,70}
\definecolor{deepsea}{RGB}{0,130,70}

% ================================
% PALETTE DEUTÉRANOPE — noms uniques
% ================================
\definecolor{crimson}{RGB}{200,30,60}
\definecolor{brickred}{RGB}{180,40,40}
\definecolor{coral}{RGB}{220,90,70}
\definecolor{salmon}{RGB}{240,130,120}

\definecolor{magenta}{RGB}{200,40,150}
\definecolor{fuchsia}{RGB}{220,60,170}
\definecolor{violetB}{RGB}{140,70,190}        % ancien "violet"
\definecolor{indigoB}{RGB}{90,50,180}         % ancien "indigo"

\definecolor{royalblueA}{RGB}{50,80,200}      % ancien "royalblue"
\definecolor{azure}{RGB}{30,120,230}
\definecolor{skyblueB}{RGB}{80,160,230}       % second "skyblue"
\definecolor{tealB}{RGB}{0,150,170}           % second "teal"

\definecolor{marine}{RGB}{0,80,140}




%% Alternative

% === Générateur automatique fill/text pour protanopes ===
% Usage : \ProtaColorPair{mistblue}
% Suppose que mistblue est déjà défini par \definecolor{mistblue}{RGB}{...}

\newcommand{\ProtaColorPair}[1]{%
  % fill = couleur originale
  \colorlet{#1fill}{#1}%
  % text = version assombrie idéalement lisible
  \colorlet{#1text}{#1!80!black}%
}

%Utilisation
%\ProtaColorPair{mistblue}

%\textcolor{mistbluefill}{Couleur claire} \quad
%\textcolor{mistbluetext}{Couleur pour texte}

	\foreach \i/\name in {
	  1/deepgreen,
	  2/greenwave,
	  3/lagoon,
	  4/oceanfoam,
	  5/softaqua,
	  6/mistblue,
 	 7/lightsky,
	  8/blueiceA,   % <-- corrigé
	  9/deepcyan,
	  10/azur,
	  11/blueviolet,
	  12/royalpurple}
	  {\ProtaColorPair{\name}}
	  


\newcommand\ROUEPROTANOPEDOUZE{
	% Roue chromatique protanope — 12 couleurs
	\begin{tikzpicture}[scale=3]

	\foreach \i/\name/\nametext in {
	  1/deepgreen/deepgreen,
	  2/greenwave/greenwave,
	  3/lagoon/lagoontext,
	  4/oceanfoam/oceanfoam,
	  5/softaqua/softaquatext,
	  6/mistblue/mistbluetext,
 	 7/lightsky/lightskytext,
	  8/blueiceA/blueiceA,   % <-- corrigé
	  9/deepcyan/blueiceA,
	  10/azur/azur,
	  11/blueviolet/blueviolet,
	  12/royalpurple/royalpurple}
	{
	  % secteur coloré
 	 \filldraw[fill=\name, draw=black, line width=0.2pt]
	    ({90-30*\i}:1)
	    -- ({90-30*(\i-1)}:1)
	    -- (0,0) -- cycle;

 	 % label dans la bonne couleur
	  \node at ({90-30*(\i-0.5)}:1.25)
	    {\textcolor{\nametext}{\small \nametext}};
	}

	\end{tikzpicture}
}


\newcommand\ROUEPROTANOPEVINGT{
\begin{tikzpicture}[scale=3]

\foreach \i/\name/\nametext in {
  1/deepgreen/deepgreen,
  2/greenwave/greenwave,
  3/lagoon/lagoon,
  4/oceanfoam/oceanfoam,
  5/softaqua/softaquatext,
  6/mistblue/mistbluetext,
  7/lightsky/lightsky,
  8/blueiceA/blueiceA,      % corrigé
  9/deepcyan/deepcyan,
  10/azur/azur,
  11/brightazure/brightazure,
  12/skyblueA/skyblueA,     % corrigé
  13/coolblue/coolblue,
  14/glacier/glacier,
  15/marine/marine,
  16/violetA/violetA,      % corrigé
  17/blueviolet/blueviolet,
  18/purpleA/purpleA,      % corrigé
  19/royalpurple/royalpurple,
  20/indigoA/indigoA}      % corrigé
{
  % Secteur coloré
  \filldraw[fill=\name, draw=black, line width=0.2pt]
    ({90-18*\i}:1)
    -- ({90-18*(\i-1)}:1)
    -- (0,0) -- cycle;

  % Étiquette
  \node at ({90-18*(\i-0.5)}:1.25)
    {\textcolor{\nametext}{\small \nametext}};
}

\end{tikzpicture}
}

\newcommand\ROUEDEUTERANOPEDOUZE{
% Roue chromatique deutéranope — 12 couleurs
\begin{tikzpicture}[scale=3]

\foreach \i/\name in {
  1/crimson,
  2/brickred,
  3/coral,
  4/salmon,
  5/magenta,
  6/fuchsia,
  7/violetB,       % corrigé
  8/indigoB,       % corrigé
  9/royalblueA,    % corrigé
  10/azure,
  11/skyblueB,     % corrigé
  12/tealB}        % corrigé
{
  % Secteur coloré
  \filldraw[fill=\name, draw=black, line width=0.2pt]
    ({90-30*\i}:1)
    -- ({90-30*(\i-1)}:1)
    -- (0,0) -- cycle;

  % Label dans la couleur
  \node at ({90-30*(\i-0.5)}:1.25)
    {\textcolor{\name}{\small \name}};
}

\end{tikzpicture}
}



% === Macro intelligente pour foncer une couleur de façon lisible pour protanope ===
% Usage : \protadarken{mistblue}{40}  -> crée mistblue@dark
% Le 40 signifie : 40% de mélange avec du noir

\newcommand{\DarkenForText}[2]{%
  % #1 = nom de la couleur
  % #2 = pourcentage (40 conseillé si tu ne sais pas)
  %
  % On crée une nouvelle couleur #1@dark
  \colorlet{#1text}{black!#2!#1}%
}





%%%%

% Groupe 1
\DarkenForText{deepgreen}{20}
\DarkenForText{forestgreen}{20}
\DarkenForText{mossgreen}{20}
\DarkenForText{greenwave}{20}
\DarkenForText{forestmist}{20}
\DarkenForText{darkteal}{20}

% Groupe 2
\DarkenForText{seagreen}{20}
\DarkenForText{freshteal}{20}
\DarkenForText{greenteal}{20}
\DarkenForText{blueteal}{20}
\DarkenForText{lagoon}{20}
\DarkenForText{greenmist}{20}

% Groupe 3
\DarkenForText{springgreen}{20}
\DarkenForText{mintgreen}{20}
\DarkenForText{oceanfoam}{20}
\DarkenForText{aquagreen}{20}
\DarkenForText{coastalaqua}{20}
\DarkenForText{softaqua}{20}

% Groupe 4
\DarkenForText{emeraldteal}{20}
\DarkenForText{deepcyan}{20}
\DarkenForText{lightcyan}{20}
\DarkenForText{aquamarine}{20}
\DarkenForText{lightaqua}{20}
\DarkenForText{mistblue}{30}

% Groupe 5
\DarkenForText{paleaqua}{20}
\DarkenForText{glacier}{20}
\DarkenForText{skyblueA}{20}
\DarkenForText{lightsky}{20}
\DarkenForText{iceblue}{20}
\DarkenForText{frost}{20}

% Groupe 6
\DarkenForText{steelblue}{20}
\DarkenForText{cadetblue}{20}
\DarkenForText{coolblue}{20}
\DarkenForText{morningblue}{20}
\DarkenForText{blueiceA}{20}
\DarkenForText{blueair}{20}

% Groupe 7
\DarkenForText{tealblue}{20}
\DarkenForText{deepbluecyan}{20}
\DarkenForText{azur}{20}
\DarkenForText{brightazure}{20}
\DarkenForText{lightazure}{20}
\DarkenForText{skylight}{20}

% Groupe 8
\DarkenForText{nightblue}{20}
\DarkenForText{slateblue}{20}
\DarkenForText{violetA}{20}
\DarkenForText{blueviolet}{20}
\DarkenForText{purpleA}{20}
\DarkenForText{lightpurple}{20}

% Groupe 9
\DarkenForText{lilac}{20}
\DarkenForText{indigoA}{20}
\DarkenForText{deepindigo}{20}
\DarkenForText{royalpurple}{20}

% Groupe 10
\DarkenForText{coolmint}{20}
\DarkenForText{deepsea}{20}

% Palette deutéranope
\DarkenForText{crimson}{20}
\DarkenForText{brickred}{20}
\DarkenForText{coral}{20}
\DarkenForText{salmon}{20}

\DarkenForText{magenta}{20}
\DarkenForText{fuchsia}{20}
\DarkenForText{violetB}{20}
\DarkenForText{indigoB}{20}

\DarkenForText{royalblueA}{20}
\DarkenForText{azure}{20}
\DarkenForText{skyblueB}{20}
\DarkenForText{tealB}{20}

% Couleur supplémentaire
\DarkenForText{marine}{20}

%%%%%%%%%%%%%%%%%%%%%
%
%.   apx proof
%
%%%%%%%%%%%%%%%%%%%%%



% do
\newcommand\PREUVESENAPPENDIXCOMPATIBLE{
	\newenvironment{appendixproof}{{\bf \LANGUE{Proof}{Démonstration}}:}{\hfill $\Box$}
	\newtheorem{theoremrep}[theorem]{Theorem}
	\newtheorem{lemmarep}[theorem]{Lemma}
	\newtheorem{propositionrep}[theorem]{Proposition}
	\newtheorem{corollaryrep}[theorem]{Corollary}
	\newtheorem{summaryrep}[theorem]{Summary}	
}

\providecommand\apxparameter{}

%% proof inline
%\providecommand\apxparameter{appendix=inline}
%% proof in appendix
%\providecommand\apxparameter{}


\apxproofmode{
\usepackage[\apxparameter]{apxproof}
\theoremstyle{plain}

\newtheoremrep{theorem}{Theorem}
%\newtheoremrep{theoremconjecture}[theorem]{Theorem-To-Be-Proved=Conjecture}
%\newtheoremrep{lemmaconjecture}[theorem]{Conjecture-To-Be-Proved=Conjecture}
\newtheoremrep{claim}[theorem]{Theorem}
\newtheoremrep{proposition}[theorem]{Proposition}
\newtheoremrep{lemma}[theorem]{Lemma}
\newtheoremrep{corollary}[theorem]{Corollary}
%%bug mars 26: \newtheoremrep{summary}[theorem]{Summary}
% Du coup:
\newtheoremrep{resume}[theorem]{Summary}
%bug mars 26: \newtheoremrep{openproblem}[theorem]{Open Problem}
%\newtheoremrep{claim}[theorem]{Claim}
}


%\apxproofmodesubsection{
%%apx proof, pour que compteur correct
%\makeatletter
%\AtBeginDocument{%
%\let\axp@section\subsection
%  \let\axp@oldsection\subsection
%}
%\makeatother
%\renewcommand\appendixprelim{\section{Missing proofs}}
%}


\makeatletter
\apxproofmodesubsection{
\let\axp@section\subsection
%\let\axp@oldsection\subsection
\renewcommand\appendixprelim{\section{Proofs from main documents}}
}
\makeatother


\makeatletter
\newcommand\appendixapxproofsubsection{
\appendix
\let\apx@orig@appendix\appendix
\renewcommand{\appendix}{}
}
\makeatother


%%% Pour checker/vérifier que bon


\newcommand\docheck[1]{\textcolor{check}{#1}}

\usepackage{xcolor}
%\definecolor{chose}{RGB}{0,110,50}
%\definecolor{forestgreen}{RGB}{0,150,70}
%\definecolor{ctruc}{RGB}{164,164,164}
%\definecolor{oceanfoam}{RGB}{0,200,150}
\definecolor{check}{RGB}{160,160,0}






%% defaults

\includeversion{NOTLFMI2020}
\includeversion{NOTCSE304}
\excludeversion{NOTDIFFUSIONFINALE}

%% Booleans (dont teste si existe ou pas, pour construire des macros)
% \SAFEMODE
% -> \DIFFUSIONFINALE (via entête-cours)
% -> \PAGEDEGARDE (via entête-cours)
% \POLYMODE
% \POLYCHAPTERMODE

%%%%%%%%%%%%%%%%%%%%%%%%%%%%%%%%%%%%%%%%%%%%%%%%%
%
%
%                      Hack pour DIFFUSION FINALE/ LFMI / Fin provisoire
%
%
%%%%%%%%%%%%%%%%%%%%%%%%%%%%%%%%%%%%%%%%%%%%%%%%%

%\def\DIFFUSIONFINALE{}

\providecommand\ifnotmodeDIFFUSIONFINALE[1]{
\ifdefined\DIFFUSIONFINALE
\else
#1
\fi
}

\providecommand\ifmodeDIFFUSIONFINALE[1]{
\ifdefined\DIFFUSIONFINALE
%% 2021: Comprend pas mais assure
\ifdefined\SAFEMODE
\else
#1
\fi
%#1
\fi
}



\ifnotmodeDIFFUSIONFINALE{\includeversion{NOTDIFFUSIONFINALE}}
\ifnotmodeDIFFUSIONFINALE{\newcommand\FINPROVISOIRELFMI[1]{}}

\providecommand\FINPROVISOIRELFMI{
  \section{Bibliographic notes}

  The current text is highly based (sometimes copied) from:
  
  \bibliographystyle{alpha}
  \bibliography{/Users/bournez/bibliographie/Bib-Files/bournez-avantbib2DOI,/Users/bournez/bibliographie/Bib-Files/perso}
  \end{document}
  }

%\def\POLYCHAPTERMODE{}

\providecommand\ifnotPOLYCHAPTERMODE[1]{
\ifdefined\POLYCHAPTERMODE
\else
#1
\fi
}

\providecommand\ifmodePOLYCHAPTERMODE[1]{
\ifdefined\POLYCHAPTERMODE
#1
\fi
}


%%%%%%%%%%%%%%%%%%%%%%%%%%%%%%%%%%%%%%%%%%%%%%%%%
%
%
%                      POUR INDICATIONS: ETAT
%
%
%%%%%%%%%%%%%%%%%%%%%%%%%%%%%%%%%%%%%%%%%%%%%%%%%


\newcommand\ABSTRACTINTERNAL[1]{
  \montruc[color=black!30]{
  \hrule
  \hrule
  \danger   ABSTRACT:\\[0.5cm] #1

  \vspace{0.5cm}
  \hrule
  \hrule
  }
  }

\newcommand\ABSTRACT[1]{
  \newpage
  \montruc[color=black!20]{
  \hrule
  \hrule
  \danger   ABSTRACT:\\[0.5cm] #1

  \vspace{0.5cm}
  \hrule
  \hrule
  }
  }

  
\newcommand\PARTIEFINALISEE{
  \newpage
  \montruc[color=black!15]{
  \hrule
  \hrule
  \danger   FINALISE A PARTIR ICI
  \hrule
  \hrule
  }
  }
\newcommand\PARTIESEMIFINALISEE{
  \newpage
  \montruc[color=black!15]{
  \hrule
  \hrule
  \danger   1/2 FINALISE A PARTIR ICI
  \hrule
  \hrule
  }
  }


\newcommand\PARTIENONFINALISEE{
  \newpage
  \montruc[color=black!30]{
  \hrule
  \hrule
  \danger   NON FINALISE A PARTIR ICI
  \hrule
  \hrule
  }
  }


%%%%%%%%%%%%%%%%%%%%%%%%%%%%%%%%%%%%%%%%%%%%%%%%%
%
%
%                      Affichage titre du cours: générique
%
%
%%%%%%%%%%%%%%%%%%%%%%%%%%%%%%%%%%%%%%%%%%%%%%%%%


  \newcommand\POLYWARNING[1]{
    
    \chapter*{Preliminaries}

    
\begin{center} 
  These are notes for the course: \\[1cm]
  \textbf{#1}
\end{center}

\bigskip
\fbox{\fbox{
\begin{minipage}{10cm}
\begin{center}
These notes change from days to days.  \\
Please do  not distribute. \\
\end{center}
\end{minipage}
}}

\bigskip
\begin{center}
  \textbf{
    Version of \today.
    }
\end{center}

 % This document is not stable.

\vfill
\begin{center}
Any comment (even about typos, orthography) is welcome: \\
send an email to bournez@lix.polytechnique.fr
\end{center}
}


\newcommand\MYCOURSEON[2]{ 
	%% External self-references
	\ifnotPOLYMODE{
		\ifdefined\MAINAUXFILE
			\externaldocument{\MAINAUXFILE}
		\fi
		}
	
	\ifdefined\PASLAPREMIEREFOIS
		\newpage
	\else
		\makeindex

     		\begin{document}
	\fi
	\gdef\PASLAPREMIEREFOIS{}

	\ifdefined\XPOLY
		\XPAGEDEGARDE{\TITRECOURS\ifPOLYCHAPTERMODE{\\[0.5cm] \Large \bf \LANGUE{Chapter:}{Chapitre:}  #2}}{\SUBTITRECOURS}
		\ifPOLYCHAPTERMODE{
		\renewcommand\thesection{\arabic{section}}
		\chapter*{#2}
		}
	\else
	
      \title{ \ifnotmodeDIFFUSIONFINALE{Course  #1 -}
      \ifPOLYCHAPTERMODE{\LANGUE{Chapter:}{Chapitre:} }
      #2}
      \author{
        Olivier Bournez% \inst{1}
        \\[0.5cm] %{\myaddress}
      }
      \date{Version of \today}
      % \institute{\myaddress}
      




      \ifnotPOLYCHAPTERMODE{
      	\maketitle
		}
      \ifPOLYCHAPTERMODE{\ifdefined\DOPAGEDEGARDE
      			\maketitle
		\fi
		\renewcommand\thesection{\arabic{section}}
%		\renewcommand\thesubsection{\arabic{section}.\ara}
      		\chapter*{#2}}

      % \begin{abstract}
      % \end{abstract}
      \fi
      


      \DOINDEXPRINT
      
      %% 2021: (meme si je comprends pas)
      \ifmodeDIFFUSIONFINALE{
		\excludeversion{recherche}
	}


    }



%%%%%%%%%%%%%%%%%%%%%%%%%%%%%%%%%%%%%%%%%%%%%%%%%
%
%
%                      Affichage titre du cours: variantes
%
%
%%%%%%%%%%%%%%%%%%%%%%%%%%%%%%%%%%%%%%%%%%%%%%%%%

    
\newcommand\MYCOURSE[1]{
  \MYCOURSEON{\csname numerocours#1\endcsname}{\csname nomcours#1\endcsname}
  
    \ifnotmodeDIFFUSIONFINALE{}%\listoftodos}

  }
  
  \newcommand\MYCOURSEnew[1]{
  \MYCOURSEON{???}{#1}
  
  \ifnotmodeDIFFUSIONFINALE{}%\listoftodos}

  }

%   ---------------------------
%   Pour compatilité chapter
%  ----------------------------

\newcommand\PASPREMIEREPAGE[1]{}
\providecommand\mychapter[1]{% \PASPREMIEREPAGE{\newpage} 
  \renewcommand\PASPREMIEREPAGE[1]{##1}  \hrule\hrule \section*{#1}
  \hrule\hrule \  \\[1cm]}
\providecommand\chapter[1]{\PASPREMIEREPAGE{\newpage} 
  \renewcommand\PASPREMIEREPAGE[1]{##1}  \hrule\hrule \section*{#1}
  \hrule\hrule \  \\[1cm]}

%   ---------------------------
%   if not sub
%  ----------------------------


\providecommand\IFNOTSUB[1]{#1}
\IFNOTSUB{
  
%\documentclass{article}
\usepackage{versions}
\includeversion{pasweb}
%
\providecommand\nolncs[1]{#1}	

%% Chemins par défaut (surchargables avant 
\providecommand\nolncs[1]{#1}	

%% Chemins par défaut (surchargables avant \input{macros} via \providecommand,
%% ou après via \renewcommand) :
\providecommand\BIBFILES{@@reference-biblio}
\providecommand\FIGCOMMONS{/Users/bournez/lib/LaTeX/Perso/FIG-COMMONS}

\newcommand\affiliationofficielleolivier{LIX, CNRS, École polytechnique, Institut Polytechnique de Paris, Palaiseau, France}
\newcommand\affiliationofficielleolivierslides{LIX, CNRS, École polytechnique, \\ Institut Polytechnique de Paris, Palaiseau, France}


%%%%%%%%%%%%%%%%%%%%%%%%
%
%              begin. Switch modes compatibilités
%
%%%%%%%%%%%%%%%%%%%%%%%

%% Apxproof mode
\providecommand\apxproofmode[1]{}
	% sinon: \newcommand\apxproof[1]{#1}
	
\providecommand\apxproofmodesubsection[1]{}
	% sinon: \newcommand\apxproofmodesubsection[1]{#1}
% utiliser: \appendixapxproofsubsection au lieu de \appendix dans ce cas. 
		

%% Lipics mode
\providecommand\lipicsmode{}
	% sinon: \newcommand\lipicsmode[1]{#1}
		% en vrai sert qu'à définir \nolipics
		
%%%%%%%%%%%%%%%%%%%%%%%%
%
%              end. Switch modes compatibilités
%
%%%%%%%%%%%%%%%%%%%%%%%


%% begin gestion des switchs

% en fait \nopilcs est la seule commande utilisée

\newcommand\nolipics[1]{#1}
\lipicsmode{\renewcommand\nolipics[1]{}}

%% Décorations (fourier + cadres colorés mdframed autour de definition,
%% theorem, exercice, example, remark, ...) : actives par défaut hors slides.
%% Pour les désactiver dans un document : \newcommand\NODECORATION[1]{}
%% avant \input{macros}. Retirées automatiquement en mode lipics.
\providecommand\NODECORATION[1]{\ifnotSLIDE{#1}}
\lipicsmode{\renewcommand\NODECORATION[1]{}}

\providecommand\PASSIAPXPROOF[1]{#1}
\apxproofmode{\renewcommand\PASSIAPXPROOF[1]{}}



%% end gestion des switchs

\def\MACROSPOINTEXOUVERT{}


\providecommand\ifnotTDEXAM[1]{#1}
\providecommand\ifnotSLIDE[1]{#1}


  
% pour preuves en appendix
\providecommand\PREUVESENAPPENDIX[1]{}



%%%% STYLES: APRES VOLS


\NODECORATION{
\usepackage{fourier}
}

%% Repli si fourier (et donc fourier-orns, qui définit \danger) n'est pas
%% chargé ; doit rester APRÈS le bloc ci-dessus, sinon fourier-orns échoue
%% avec « \danger already defined ».
\providecommand\danger{}

%%%% Box around environment
\usepackage[framemethod=tikz]{mdframed}

\newcommand\ENVCOLOR[1]{#1}
\newcommand\BEGINTEMPENVCOLOR[1]{\let\envcolorwas\ENVCOLOR\renewcommand\ENVCOLOR[1]{#1}}
\newcommand\ENDTEMPENVCOLOR{\let\ENVCOLOR\envcolorwas}


\NODECORATION{
\surroundwithmdframed[hidealllines=true,backgroundcolor=\ENVCOLOR{blue!10},roundcorner=8pt]{exercice}
\surroundwithmdframed[hidealllines=true,backgroundcolor=\ENVCOLOR{blue!10},roundcorner=8pt]{exercise}
\surroundwithmdframed[hidealllines=true,backgroundcolor=\ENVCOLOR{blue!10},roundcorner=8pt]{exercicee}
\surroundwithmdframed[hidealllines=true,backgroundcolor=\ENVCOLOR{blue!10},roundcorner=8pt]{exerciced}
\surroundwithmdframed[hidealllines=true,backgroundcolor=\ENVCOLOR{blue!10},roundcorner=8pt]{exercicec}
\surroundwithmdframed[hidealllines=true,backgroundcolor=\ENVCOLOR{blue!10},roundcorner=8pt]{exercicedc}
\surroundwithmdframed[hidealllines=true,backgroundcolor=\ENVCOLOR{orange!20},roundcorner=8pt]{example}
\surroundwithmdframed[hidealllines=true,backgroundcolor=\ENVCOLOR{orange!20},roundcorner=8pt]{exemple}
\surroundwithmdframed[hidealllines=true,backgroundcolor=\ENVCOLOR{orange!10},roundcorner=8pt]{remark}
\ifnotSLIDE{\surroundwithmdframed[backgroundcolor=\ENVCOLOR{green!15},roundcorner=8pt]{definition}}
\surroundwithmdframed[hidealllines=true,backgroundcolor=\ENVCOLOR{gray!35}]{proposition}
\ifnotSLIDE{\surroundwithmdframed[backgroundcolor=\ENVCOLOR{gray!35},roundcorner=8pt]{theorem}}
\surroundwithmdframed[hidealllines=true,backgroundcolor=\ENVCOLOR{gray!35}]{lemma}
\surroundwithmdframed[hidealllines=true,backgroundcolor=\ENVCOLOR{gray!35}]{corollary}


\surroundwithmdframed[backgroundcolor=red!35,roundcorner=8pt]{theoremaprouver}
\surroundwithmdframed[backgroundcolor=red!35,roundcorner=8pt]{conjectureaprouver}
}

%%% tcolorbox

\usepackage{tcolorbox}
\tcbuselibrary{skins}


  %%. Commande magique
%% \tcolorboxenvironment{⟨name⟩}{⟨options⟩}: attach options de tcolorbox à environnement
%% alternative: je cree des boites moi meme: voici des exemples


\newtcolorbox{maboiteconfiance}[2][]{colback=red!5!white, colframe=red!75!black,fonttitle=\bfseries, colbacktitle=red!85!black,enhanced,
attach boxed title to top center={yshift=-2mm},
  title={#2},#1}
  
\newtcolorbox{maboiteresume}[2][]{colback=red!5!white, colframe=red!75!black,fonttitle=\bfseries, colbacktitle=red!85!black,enhanced,
attach boxed title to top center={yshift=-2mm},
  title={#2},#1}
  \newtcolorbox{maboiteresumeperso}[2][]{colback=red!5!white, colframe=red!75!black,fonttitle=\bfseries, colbacktitle=red!85!black,enhanced,
attach boxed title to top center={yshift=-2mm},
  title={#2},#1}
  \newtcolorbox{maboitecommentaire}[2][]{colback=red!5!white, colframe=red!75!black,fonttitle=\bfseries, colbacktitle=red!85!black,enhanced,
attach boxed title to top center={yshift=-2mm},
  title={#2},#1}
\newtcolorbox{maboitetruc}[2][]{colback=red!5!white, colframe=red!75!black,fonttitle=\bfseries, colbacktitle=red!85!black,enhanced,
attach boxed title to top left={yshift=-2mm},
  title={#2},#1}

%% Boîte des mises en évidence CARE. Couleurs pensées pour la
%% protanopie : l'axe bleu-jaune reste discriminant (jaune = \IMPORTANT
%% d'Olivier, bleu = CARE), et la distinction ne repose pas que sur la
%% couleur (cadre + bandeau étiqueté, lisibles même en niveaux de gris).
\newtcolorbox{maboiteimportantcare}[2][]{colback=cyan!10!white, colframe=blue!60!black,fonttitle=\bfseries, colbacktitle=blue!70!black,enhanced,
attach boxed title to top center={yshift=-2mm},
  title={#2},#1}
  


%%%% FIN STYLES



%%%%%%%%%%%%%%%%%%%%%%%%%%%%%%%%%%%%%%%%%%%%%%%%%
%%%%%%%%%%%%%%%%%%%%%%%%%%%%%%%%%%%%%%%%%%%%%%%%%
%
%
%                      Version ENSEIGNEMENT 2020
%
%
%%%%%%%%%%%%%%%%%%%%%%%%%%%%%%%%%%%%%%%%%%%%%%%%%
%%%%%%%%%%%%%%%%%%%%%%%%%%%%%%%%%%%%%%%%%%%%%%%%%

%english, french
\ifdefined\CESTDUFRANCAIS
	\providecommand\LANGUE[2]{#2}
	\usepackage[french]{babel}
\else
	\providecommand\LANGUE[2]{#1}
\fi
\providecommand\LANGUEinv[2]{\LANGUE{#2}{#1}}

%%%%%%%%%%%%%%%%%%%%%%%%%%%%%%%%%%%%%%%%%%%%%%%%%
%
%
%                      Packages
%
%
%%%%%%%%%%%%%%%%%%%%%%%%%%%%%%%%%%%%%%%%%%%%%%%%%


%%% SURLIGNEUR
\providecommand\ifnotmodeDIFFUSIONFINALE[1]{
%% 2021: Comprend pas mais assure
\ifdefined\SAFEMODE
\else
#1
\fi
%#1
}
\providecommand\ifmodeDIFFUSIONFINALE[1]{
\ifdefined\SAFEMODE
#1
\else
\fi
}
\newcommand\SURLIGNE[1]{
\ifnotmodeDIFFUSIONFINALE{
\BEGINTEMPENVCOLOR{blue!30}
{                        \colorbox{blue!30}{

\noindent \begin{minipage}{\dimexpr\textwidth-2\fboxsep\relax}
#1
\end{minipage}
}
}
\ENDTEMPENVCOLOR}
\ifmodeDIFFUSIONFINALE{
\noindent \begin{minipage}{\textwidth}
#1
\end{minipage}
}
}


\newenvironment{confiance}{\begin{maboiteresume}[colback=brown!50]{Résumé ici}}{\end{maboiteresume}}
\newcommand\CONFIANCE[2][Confiance ici]{
\begin{maboiteconfiance}[colback=yellow!50]{#1}
\textbf{ATTENTION : #2}
\end{maboiteconfiance}
}

\newenvironment{resume}{\begin{maboiteresume}[colback=brown!50]{Résumé ici}}{\end{maboiteresume}}
\newcommand\RESUME[2][Résumé ici]{
\begin{maboiteresume}[colback=brown!50]{#1}
#2
\end{maboiteresume}
}

\newenvironment{resumeperso}{\begin{maboiteresumeperso}[colback=brown!50]{Résumé perso ici}}{\end{maboiteresumeperso}}
\newcommand\RESUMEPERSO[2][Résumé perso  ici]{
\begin{maboiteresumeperso}[colback=brown!50]{#1}
#2
\end{maboiteresumeperso}
%
%
%\todo[inline, color=brown]{résumé: #1}
%\colorbox{brown}{
%\noindent \begin{minipage}{\textwidth}
%{
%\bf 
%#2
%}
%\end{minipage}
%}
}

\newcommand\DANSLAMARGE[1]{
\ifnotmodeDIFFUSIONFINALE{
\marginpar{\textcolor{brown}{#1}}}
}

\newenvironment{commentaireperso}{\begin{maboitecommentaire}[colback=brown!30]{Commentaire perso} Commentaire perso:}{\end{maboitecommentaire}}
\newcommand\COMMENTAIREPERSO[2][Commentaire perso]{
\begin{maboitecommentaire}[colback=brown!30]{#1}
#2
\end{maboitecommentaire}
}


\newcommand\SURSURLIGNE[2][truc surligné ici]{
%\todo[inline, color=orange]{surligné: #1}
\ifnotmodeDIFFUSIONFINALE{
\BEGINTEMPENVCOLOR{blue!60}
	                 \colorbox{blue!60}{
\noindent \begin{minipage}{\dimexpr\textwidth-2\fboxsep\relax}
{
\bf 
#2
}
\end{minipage}
}
\ENDTEMPENVCOLOR
}
\ifmodeDIFFUSIONFINALE{#2}
}

\newcommand\IMPORTANT[1]{
\ifnotmodeDIFFUSIONFINALE{
\BEGINTEMPENVCOLOR{yellow!100}
		        \colorbox{yellow!100}{

\noindent \begin{minipage}{\dimexpr\textwidth-2\fboxsep\relax}
{ 
#1
}
\end{minipage}
}
\ENDTEMPENVCOLOR
}
\ifmodeDIFFUSIONFINALE{#1}
}

%% Pendant de \IMPORTANT, réservé aux mises en évidence CARE ;
%% \IMPORTANT reste réservé aux annotations d'Olivier. Comme \IMPORTANT,
%% redevient du texte simple en mode DIFFUSIONFINALE.
\newcommand\IMPORTANTCARE[2][Important (CARE)]{
\ifnotmodeDIFFUSIONFINALE{
\begin{maboiteimportantcare}{#1}
#2
\end{maboiteimportantcare}
}
\ifmodeDIFFUSIONFINALE{#2}
}


\newcommand\QUESTION[1]{
\ifnotmodeDIFFUSIONFINALE{
\BEGINTEMPENVCOLOR{oceanfoam!100}
		        \colorbox{oceanfoam!100}{

\noindent \begin{minipage}{\dimexpr\textwidth-2\fboxsep\relax}
{ 
#1
}
\end{minipage}
}
\ENDTEMPENVCOLOR
}
\ifmodeDIFFUSIONFINALE{#1}
}




\newcommand\SOUSLIGNE[1]{
\colorbox{lightgray}{

\noindent \begin{minipage}{\dimexpr\textwidth-2\fboxsep\relax}
{ \color{gray}
#1
}
\end{minipage}
}
}



%% Gestion des cross-references cross-referencing
\usepackage{xr-hyper}
\usepackage{hyperref}

%%
\usepackage{versions}
\usepackage{amsmath,amssymb,mathrsfs,amsfonts}
%\usepackage{mathabx}
\usepackage{listings}
\usepackage{url}
\usepackage{versions}
\usepackage{color}
\usepackage[colorinlistoftodos]{todonotes}
%\usepackage{todonotes}
\usepackage{proof}
\usepackage{xstring}

%% Index
\usepackage{makeidx}
\makeindex
\providecommand\DOINDEXPRINT{}
\providecommand\SIGNALEMOTNOUV[1]{}
\ifdefined\INDEXPRINT
	\ifnotSAFEMODE{
%	\usepackage{showidx}
%	\let\oldindex\index
%	\renewcommand{\index}[1]{\oldindex{#1}\marginpar{LA: \tiny#1}}
	\renewcommand\DOINDEXPRINT{
	      	\let\oldindex\index
		\renewcommand{\index}[1]{\oldindex{##1}\marginpar{IDXLA: \tiny##1}}
		\renewcommand\SIGNALEMOTNOUV[1]{\marginpar{MNV: \small\textbf{##1}}}
	}
	}
\fi

% === gestion des accents


%soit iso latin1
%\usepackage[cyr]{aeguill}
%soit utf8
\usepackage[utf8]{inputenc}
\usepackage[T1]{fontenc}


%%%  Pour citation \Cref (cref, ref)	

\usepackage{cleveref}


%%%%%%%%%%%%%%%%%%%%%%%%%%%%%%%%%%%%%%%%%%%%%%%%%
%
%
%                      Gestion des accents spéciaux.
%
%
%%%%%%%%%%%%%%%%%%%%%%%%%%%%%%%%%%%%%%%%%%%%%%%%%

\input{macros-accents}
%% Tolérant tant que macros-markdown.tex n'est pas dans le dépôt :
\InputIfFileExists{macros-markdown}{}{\typeout{*** macros-markdown.tex introuvable: ignore ***}}


%\usepackage{PDFdraftcopy}


%-- Pour état des chapitres


%\newcommand\ETAT{\draftstring{BROUILLON}}
%
%\newcommand\brouillon{\renewcommand\ETAT{\draftstring{BROUILLON}}}
%
%\newcommand\final{\renewcommand\ETAT{\draftstring{\rien}}}
%
%\newcommand\travail{\renewcommand\ETAT{\draftstring{CHANTIER}}}
%
%\newcommand\bazar{\renewcommand\ETAT{\draftstring{BAZAR TOTAL}}}
%
%\final



%%%%%%%%%%%%%%%%%%%%%%%%%%%%%%%%%%%%%%%%%%%%%%%%%
%
%
%                      TRADUCTION EN COURS
%
%%%%%%%%%%%%%%%%%%%%%%%%%%%%%%%%%%%%%%%%%%%%%%%%%


% === volé Amaury



\usepackage{stmaryrd}
\usepackage{ifthen}
\usepackage{mathtools}
\usepackage{dsfont}
\PASSIAPXPROOF{\usepackage{thmtools}}
\usepackage{longtable}
\usepackage{mathdots}
\usepackage{booktabs}


%%%%%%%%%%%%%%%%%%%%%%%%%%%%%%%%%%%%%%%%%%%%%%%%%
%
%
%                      TIKZ STUFFS
%
%
%%%%%%%%%%%%%%%%%%%%%%%%%%%%%%%%%%%%%%%%%%%%%%%%%

%%% volé pour dessin des GPAC: PAS SUR QUE BESOIN de TOUT

\usepackage{tikz}
\usetikzlibrary{calc}
\usepgflibrary{arrows}
\usetikzlibrary{fit}
\usetikzlibrary{patterns}
\usetikzlibrary{decorations.pathreplacing}
\usetikzlibrary{shapes.multipart}
\usetikzlibrary{shapes.geometric}
\usetikzlibrary{shapes.misc}
\usetikzlibrary{positioning}
\usetikzlibrary{graphs}
\usetikzlibrary{topaths}
\usetikzlibrary{arrows.meta}
\usetikzlibrary{decorations.pathreplacing}
\usetikzlibrary{decorations.markings}
\usetikzlibrary{decorations.pathmorphing}
%
 \usetikzlibrary{calc}
 \usetikzlibrary{automata}
\usetikzlibrary{chains}
\usetikzlibrary{scopes}
 \usetikzlibrary{positioning}
 \usetikzlibrary{through,backgrounds}
\usepackage{circuitikz}
% % pour accolades




%%%%%%%%%%%%%%%%%%%%%%%%%%%%%%%%%%%%%%%%%%%%%%%%%
%
%
%                      MACROS locales
%
%
%%%%%%%%%%%%%%%%%%%%%%%%%%%%%%%%%%%%%%%%%%%%%%%%%


%====================
%                      Moi
%=====================


\newcommand\myaddress{ 
{\sc LIX, CNRS, Ecole Polytechnique, Institut Polytechnique de Paris, }\\ 91128 Palaiseau Cedex, FRANCE \\
{\small Olivier.Bournez@lix.polytechnique.fr}}



%====================
%                      un peu de trucs sales
%=====================

%=  pour style lipics et ses trucs

\providecommand\ifnotPOLYCHAPTERMODE[1]{#1}
\providecommand\ifPOLYCHAPTERMODE[1]{}

\ifnotPOLYCHAPTERMODE{
\ifdefined\thechapter
\providecommand\spnewtheorem[4]{\newtheorem{#1}{#2}[chapter]}
\providecommand\spnewtheoremi[5]{\newtheorem{#1}[#5]{#2}}
\else
\providecommand\spnewtheorem[4]{\newtheorem{#1}{#2}}
\providecommand\spnewtheoremi[5]{\newtheorem{#1}[#5]{#2}}
\fi
}
\ifPOLYCHAPTERMODE{
\providecommand\spnewtheorem[4]{\newtheorem{#1}{#2}}
\providecommand\spnewtheoremi[5]{\newtheorem{#1}[#5]{#2}}
}
\providecommand\nolipics[1]{#1}




%====================
%                      un peu de trucs franchement sales
%=====================


% on l'inclus pas ici.
%\input{macros-moins-propres}

%%%%%%%%%%%%%%%%%%%%%%%%%%%%%%%%%%%%%%%%%%%%%%%%%
%
%
%                      Page de Garde Cours X
%
%
%%%%%%%%%%%%%%%%%%%%%%%%%%%%%%%%%%%%%%%%%%%%%%%%%


\includeversion{metdate}

\newcommand\XPAGEDEGARDE[2]{
\thispagestyle{empty}
\title{{#1} \\[2cm] %\\[3cm] 
{\large #2}
}
\author{
 Olivier Bournez% \inst{1}
\\[0.5cm] 
bournez@lix.polytechnique.fr%{\myaddress}
}
%\date{
%\vspace{2cm}
%\includegraphics[height=2.5cm]{/Users/bournez/lib/LaTeX/Perso/FIG-COMMONS/logo_X_new}
%}
\begin{metdate}
\date{
\vspace{2cm}
Version \LANGUE{of}{du} \today \\[1cm] 
\includegraphics[height=3.5cm]{/Users/bournez/lib/LaTeX/Perso/FIG-COMMONS/logo_X_new}
}
\end{metdate}
\maketitle
}


%%%%%%%%%%%%%%%%%%%%%%%%%%%%%%%%%%%%%%%%%%%%%%%%%
%
%
%                      MACROS UTILES
%
%
%%%%%%%%%%%%%%%%%%%%%%%%%%%%%%%%%%%%%%%%%%%%%%%%%

%%%                      Macros
 
%-------------
%--- vectors  ---
%-------------

 
\newcommand\vectorl[1]{{\mathbf#1}}
\newcommand\tu[1]{\vectorl{#1}}

\newcommand\vp{\vectorl{p}}
\newcommand\va{\vectorl{a}}
\newcommand\vb{\vectorl{b}}
\newcommand\vc{\vectorl{c}}
\newcommand\vf{\vectorl{f}}
\newcommand\vn{\vectorl{n}}
\newcommand\vx{\vectorl{x}}
\newcommand\vX{\vectorl{X}}
\newcommand\vy{\vectorl{y}}
\newcommand\vz{\vectorl{z}}
\newcommand\ve{\vectorl{e}}
\newcommand\vv{\vectorl{v}}
\newcommand\vr{\vectorl{r}}
\newcommand\vu{\vectorl{u}}
\newcommand\vA{\vectorl{A}}
\newcommand\vB{\vectorl{B}}
\newcommand\vC{\vectorl{C}}
\newcommand\vD{\vectorl{D}}

%-------------
%--- overline ---
%-------------


\newcommand\ox{\overline{x}}
\newcommand\oy{\overline{y}}
\newcommand\oc{\overline{c}}
\newcommand\oa{\overline{a}}
\newcommand\ow{\overline{w}}
\newcommand\od{\overline{d}}
\newcommand\ob{\overline{b}}
\newcommand\oP{\overline{P}}
\newcommand\oee{\overline{e}}
\newcommand\ot{\overline{t}}
\newcommand\ou{\overline{u}}
\newcommand\ov{\overline{v}}
\newcommand\oz{\overline{z}}
\newcommand\am{\overline{a}}
\newcommand\om{\overline{m}}
\newcommand\oj{\overline{j}}

%----------------
%--- ensembles ---
%----------------

\newcommand\Q{\mathbb{Q}}
\newcommand\R{\mathbb{R}}
\newcommand\Z{\mathbb{Z}}
\providecommand\C{\mathbb{C}}
\newcommand\N{\mathbb{N}}
\newcommand\K{\mathbb{K}}
\newcommand\F{\mathbb{F}}

\newcommand\dyadic{\mathbb{D}}
\newcommand\Zd{{\Z_2}}
\newcommand\Zp{{\Z_p}}
%
\newcommand\RP{\mathbb{R}^{>0}}
\newcommand\QP{\mathbb{Q}^{>0}}


%----------------
%--- lettres cal ---
%----------------

\newcommand\mK{\mathcal{K}}
\newcommand\mD{\mathcal{D}}
\newcommand\mA{\mathcal{A}}
\newcommand\mL{\mathcal{L}}
\newcommand\mX{\mathcal{X}}
\newcommand\mG{\mathcal{G}}
\newcommand\mE{\mathcal{E}}
\newcommand\mH{\mathcal{H}}
\newcommand\mR{\mathcal{R}}
\newcommand\mF{\mathcal{F}}
\newcommand\mJ{\mathcal{J}}
\newcommand\mO{\mathcal{O}}
\newcommand\mP{\mathcal{P}}
\newcommand\mN{\mathcal{N}}
\newcommand\mS{\mathcal{S}}
\newcommand\mM{\mathcal{M}}
\newcommand\mC{\mathcal{C}}
\newcommand\mV{\mathcal{V}}
\newcommand\mI{\mathcal{I}}
\newcommand\mT{\mathcal{T}}

\newcommand\fC{\mathcal{C}}

%----------------
%--- lettres frak ---
%----------------


\newcommand\kM{\mathfrak{M}}
\newcommand\kN{\mathfrak{N}}
\newcommand\kU{\mathfrak{U}}
\newcommand\kg{\mathfrak{g}}



%--------------------------
%--- Façon noter noms classes ---
%---------------------------

\newcommand{\cp}[1]{\operatorname{#1}}


%--------------------------
%--- Calculabilité  ---
%---------------------------

% -- classes
\LANGUE{
\newcommand\RE{\cp{CE}}
}
{\newcommand\RE{\cp{RE}}
}

\newcommand\Rec{\cp{D}}
\newcommand\PR{\cp{PR}}
\newcommand\Gregn{\mE_n}


%--------------------------
%--- Réductions  ---
%---------------------------

\newcommand\lem{\le_{m}}
\newcommand\lelog{\le_{L}}
\newcommand\equivlog{\equiv_{L}}



%--------------------------
%--- Complexité  ---
%---------------------------

% COMPLEXITE

% -- P
\newcommand{\Ptime}{\cp{P}}      % P
\newcommand\PTIME{\Ptime}
\newcommand\PTIMEs[1]{{\PTIME}_{#1}}
\newcommand\PTIMEsp[1]{{\PTIME}_{#1}^0}
\renewcommand\P{\PTIME}
\newcommand{\FPtime}{\cp{FPTIME}}

% -- NP
\newcommand\NP{\cp{NP}}
\newcommand{\NPtime}{\NP}
\newcommand\NPs[1]{\cp{NP}_{#1}}
\newcommand{\FNP}{\cp{FNP}}

\newcommand\NDP{\cp{NDP}}
\newcommand\coNDP{\cp{coNDP}}
\newcommand\NDPs[1]{\cp{NDP}_{#1}}

% -- coNP
\newcommand\coNP{\cp{coNP}}

% -- PSPACE
\newcommand\PSPACE{\cp{ PSPACE}}
\newcommand\NPSPACE{\cp{ NPSPACE}}
\newcommand{\Pspace}{\PSPACE}
\newcommand{\FPspace}{\cp{FPSPACE}}


% -- LOGSPACE
\newcommand\LSPACE{\cp{LOGSPACE}}
\newcommand\NLSPACE{\cp{NLOGSPACE}}
\newcommand\coNLSPACE{\cp{coNLOGSPACE}}
\newcommand\coNL{\cp{coNL}}

% -- Autres
\newcommand\EXPTIME{\cp{EXPTIME}}
\newcommand\NEXPTIME{\cp{NEXPTIME}}
\newcommand\EXPSPACE{\cp{EXPSPACE}}
\newcommand\PH{\cp{PH}}
\newcommand\NC{\cp{NC}}
\newcommand\AC{\cp{AC}}
\newcommand\PPAD{\cp{\mathbf{PPAD}}}
\newcommand\PLS{\cp{\mathbf{PLS}}}

% -- Reverse math
\newcommand{\WKL}{\ensuremath{\docheck{\mathsf{WKL}_0}}}
\newcommand{\ACA}{\ensuremath{\docheck{\mathsf{ACA}_0}}}
\newcommand{\ATR}{\ensuremath{\docheck{\mathsf{ATR}_0}}}
\newcommand{\RCAO}{\ensuremath{\docheck{\mathsf{RCA}_0}}}
\newcommand{\PiCA}{\ensuremath{\docheck{\Pi^1_1\!\text{-}\mathsf{CA}_0}}}


% -- circuits
\newcommand\CIRCUIT[2]{\cp{CIRCUIT(#1,#2)}}
\newcommand\UCIRCUIT[2]{\cp{UCIRCUIT(#1,#2)}}

% -- Classes proba
\newcommand\PP{\cp{PP}}
\newcommand\RPP{\cp{RP}}
\newcommand\ZPP{\cp{ZPP}}
\newcommand\coRP{\cp{coRP}}
\newcommand\BPP{{\cp{BPP}}}

% -- IP
\newcommand\IP{\cp{IP}}
\newcommand\PCP{\cp{PCP}}

% -- non-uniforme
\newcommand\nuP{\mathbb{P}}
\newcommand\nuPs[1]{\mathbb{P}_{#1}}
\newcommand\nuPsp[1]{\mathbb{P}_{#1}^0}
\newcommand\nuNP{\mathbb{N}\mathbb{P}}
%\newcommand\poly{/{poly}}
\newcommand\Ppoly{\PTIME_{\poly}}
\newcommand\NPpoly{\NP_{\poly}}

% -- conditions d'uniformité
\newcommand\Luniforme{\cp{L}}
\newcommand\BCuniforme{\cp{BC}}


% Classes de circuits
\newcommand\TC{\cp{TC}}
\newcommand\LFP{\cp{LFP}}
\newcommand\FAC{\cp{FAC}}
\newcommand\FACC{\cp{FACC}}
\newcommand\FTC{\cp{FTC}}
\newcommand\FNC{\cp{FNC}}

% Logiques temorelles
\newcommand\CTL{\cp{CTL}}
\newcommand\LTL{\cp{LTL}}

% -- mesure temps/espace/taille

\newcommand\DTIME{\cp{DTIME}}
\newcommand\TIME[1]{\cp{TIME(#1)}}
\newcommand\TIMEs[2]{\cp{TIME_{#2}(#1)}}
\newcommand\TIMEsp[2]{\cp{TIME^0_{#2}(#1)}}
\newcommand\NTIME[1]{\cp{NTIME(#1)}}
\newcommand\NTIMEs[2]{\cp{NTIME_{#2}(#1)}}
\newcommand\NTIMEsp[2]{\cp{NTIME_{#2}^0(#1)}}
\newcommand\DSPACE{\cp{DSPACE}}
\newcommand\SPACE[1]{\cp{SPACE(#1)}}
\newcommand\SPACEs[2]{\cp{SPACE_{#2}(#1)}}
\newcommand\NSPACE[1]{\cp{NSPACE(#1)}}
\newcommand\NSPACEs[2]{\cp{NSPACE_{#2}(#1)}}
\newcommand\coNSPACE[1]{\cp{coNSPACE(#1)}}
\newcommand\TIMEON[2]{\cp{TIME(#1,#2)}}
\newcommand\SPACEON[2]{\cp{SPACE(#1,#2)}}
\newcommand\SIZE[1]{\cp{ size(#1)}}
\newcommand\ATIME{\cp{ATIME}}
\newcommand\ASPACE{\cp{ASPACE}}


%% autre:
\newcommand\LH{\cp{LH}}
\newcommand\FO{\cp{FO}}
\newcommand\SO{\cp{SO}}
\newcommand\ESOhorn{\cp{ESO_{horn}}}
\newcommand\SOhorn{\cp{SO_{horn}}}
\newcommand\ACzero{\cp{AC_0}}
\newcommand{\LINSPACE}{\cp{LINSPACE}}
\newcommand{\mFLINSPACE}{\mF_{\cp{LINSPACE}}}
\newcommand{\LOGSPACE}{\cp{LOGSPACE}}
\newcommand\ALOGTIME{\cp{ALOGTIME}}
\newcommand\RF{\cp{RF}}


\newcommand\BOOL[1]{\cp{{BOOLEAN(#1)}}}
\newcommand\DET{\cp{ DET}}

% % -- dirty


\newcommand{\cPspace}{\cp{\sharp PSPACE}}
\newcommand{\cAPtime}{\cp{\sharp APTIME}}
%\newcommand{\FPspace}{\cp{FPSPACE}}
%\newcommand{\FPspaceN}{\cp{FPSPACE}^{+}}
%\newcommand{\FPspaceN}{\cp{FPSPACE}}



%-------------------------
%--- Problèmes de décision  ---
%--------------------------

\newcommand\decisionname[1]{\cp{#1}}


\newcommand\Luniv{\decisionname{L_{univ}}}
\newcommand\HALTINGPROBLEM{\decisionname{HALTING-PROBLEM}}
\newcommand\SAT{\decisionname{SAT}}
\newcommand\MAXtSAT{\decisionname{MAX3SAT}}
\newcommand\REACH{\decisionname{REACH}}
\newcommand\CIRCUITVALUE{\decisionname{CIRCUITVALUE}}
\newcommand\MONOTONECIRCUITVALUE{\decisionname{MONOTONECIRCUITVALUE}}
\newcommand\THEOREMS{\decisionname{THEOREMS}}
\newcommand\QCIRCUITSAT{\decisionname{QCIRCUITSAT}}
\newcommand\QSAT{\decisionname{QSAT}}
\newcommand\QBF{\decisionname{QBF}}
\newcommand\SATVALUE{\decisionname{SATVALUE}}
\newcommand\SUCCINTSAT{\decisionname{SUCCINTSAT}}
\newcommand\SUCCINTSATVALUE{\decisionname{SUCCINTSATVALUE}}
\newcommand\GEOGRAPHY{\decisionname{GEOGRAPHY}}
\newcommand\PARITY{\decisionname{PARITY}}
\newcommand\MAJ{\decisionname{MAJ}}
\newcommand\CVP{\CIRCUITVALUE}
\newcommand\BOOLEANPROGRAMVALUE{BOOLEAN~PROGRAM~VALUE}

\newcommand\CLIQUE{\cp{CLIQUE}}
\LANGUEinv{
\newcommand\RECOUVREMENTDESOMMETS{\cp{RECOUVREMENT~DE~SOMMETS}}
\newcommand\RS{\RECOUVREMENTDESOMMETS}
}{
\newcommand\RECOUVREMENTDESOMMETSeng{\cp{VERTEX~COVER}}
\newcommand\RSeng{\RECOUVREMENTDESOMMETSeng}
}

\LANGUEinv{\newcommand\COUPUREMAXIMALE{\cp{COUPURE~MAXIMALE}}}
{\newcommand\COUPUREMAXIMALEeng{\cp{MAXIMAL~CUT}}}
\newcommand\CIRCUITHAMILTONIEN{\CHAMILTONIEN}
\LANGUEinv{\newcommand\VOYAGEURDECOMMERCE{\VCOMMERCE}}
{\newcommand\VOYAGEURDECOMMERCEeng{\VCOMMERCEeng}}

\LANGUEinv{\newcommand\CIRCUITLEPLUSLONG{\cp{CIRCUIT~LE~PLUS~LONG}}}
{\newcommand\CIRCUITLEPLUSLONGeng{\cp{LONGEST~CIRCUIT}}}
\LANGUEinv{\newcommand\SOMMEDESOUSENSEMBLE{\SUBSETSUM}}{
\newcommand\SOMMEDESOUSENSEMBLEeng{\SUBSETSUMeng}}
\LANGUEinv{
\newcommand\SACADOS{\cp{SAC~A~DOS}}}
{\newcommand\SACADOSeng{\cp{KNAPSACK}}}

%français
\LANGUEinv{
\newcommand\CHAMILTONIEN{\cp{CIRCUIT~HAMILTONIEN}}}
{\newcommand\CHAMILTONIENeng{\cp{HAMILTONIAN~CIRCUIT}}}
\LANGUEinv{
\newcommand\VCOMMERCE{\cp{VOYAGEUR~DE~COMMERCE}}}
{\newcommand\VCOMMERCEeng{\cp{TRAVELING~SALESMAN}}}
\LANGUEinv{
\newcommand\kCOLORABILITE{\cp{k-COLORABILITE}}}
{\newcommand\kCOLORABILITEeng{\cp{k-COLORABILITY}}}
\LANGUEinv{
\newcommand\COLORIAGE{\cp{BON~COLORIAGE}}}
{\newcommand\COLORIAGEeng{\cp{GOOD~COLORING}}}
\LANGUEinv{
\newcommand\tCOLORABILITE{\cp{3-COLORABILITE}}}
{\newcommand\tCOLORABILITEeng{\cp{3-COLORABILITY}}}
\LANGUEinv{
\newcommand\SUBSETSUM{\cp{SOMME~DE~SOUS~ENSEMBLE}}}
{\newcommand\SUBSETSUMeng{\cp{SUBSET~SUM}}}
\newcommand\PARTITION{\cp{PARTITION}}
\LANGUEinv{
\newcommand\NOMBREPREMIER{\decisionname{NOMBRE~ PREMIER}}}
{\newcommand\NOMBREPREMIEReng{\decisionname{PRIME~NUMBER}}}
\LANGUEinv{
\newcommand\CODAGE{\decisionname{CODAGE}}}
{\newcommand\CODAGEeng{\decisionname{ENCODING}}}
\newcommand\kSAT{\cp{k-SAT}}
\newcommand\tSAT{\cp{3-SAT}}
\newcommand\NAESAT{\cp{NAESAT}}
\newcommand\NAEtSAT{\cp{NAE3SAT}}
\newcommand\NAEqSAT{\cp{NAE4SAT}}
\newcommand\STABLE{\cp{\LANGUE{INDEPENDANT~SET}{STABLE}}}

%\newcommand\CM{\cp{COUPURE MAXIMALE}}
\LANGUEinv{
\newcommand\POSTpb{\cp{Problème~de~ la~correspondance~de~Post}}}
{\newcommand\POSTpbeng{\cp{Post~correspondence~problem}}}

\LANGUEinv{
\newcommand\PARITE{\cp{PARITE}}}
{\newcommand\PARITEeng{\cp{PARITY}}}


%-------------------------
%--- Modèles parallèles  ---
%--------------------------


% PARALLELISME
\newcommand\RAM{\cp{RAM}}
%
\newcommand\PRAM{\cp{PRAM}}
\newcommand\CRCW{\cp{CRCW}}
\newcommand\CREW{\cp{CREW}}
\newcommand\EREW{\cp{EREW}}


%------------------------
%--- codages  ---
%------------------------


\newcommand\codage[1]{\langle #1 \rangle}
\newcommand\tuple[1]{\codage{#1}}
\newcommand\paire[2]{\codage{#1,#2}}
\newcommand\triplet[3]{\codage{#1,#2,#3}}
\newcommand\quadruplet[4]{\codage{#1,#2,#3,#4}}
\newcommand\cinquplet[5]{\codage{#1,#2,#3,#4,#5}}


%------------------------
%--- descriptive set theory  ---
%------------------------

\newcommand\baire{\mathcal{N}}
\newcommand\peano{\overline{\mathcal{N}}} %% A MOI
\newcommand\cantor{\mathcal{C}}
%
\newcommand\bSigma{\mathbf{\Sigma}}
\newcommand\bPi{\mathbf{\Pi}}
\newcommand\bDelta{\mathbf{\Delta}}
%


%%%%%% ABOUT NOMS QUI CHANGES

\newcommand\WF{\cp{WF}}
\newcommand\WO{\WF}
\newcommand\LO{\cp{LO}}
\newcommand\DLO{\cp{DLO}}
% 
\newcommand\rk{\cp{rank}}


%% Moins propre:

\newcommand\I{\mathbb{I}}
\renewcommand\H{\mathbb{H}}
%
\newcommand\leKB{<_{KB}}
%
\newcommand\finiomega{{<\omega}}
\newcommand\compl{^\frown}
\newcommand\prefix{ ~ prefix~  } 
\newcommand\X{{X}}
\newcommand\Y{\mathcal{Y}}
\newcommand\subsetstrict{\subset}

\renewcommand\complement[1]{#1^{c}}


%------------------------
%--- schémas de récursion ---
%------------------------

%%                      Source papier MFCS 2019


% COMPLEXITY

%% MFCS differe de Clote:
%% Version à taille restreinte

\newcommand\co{\cp{co}}

\newcommand\mFPTIME{\mF_{PTIME}}
\newcommand{\mFPSPACE}{\mF_{SPACE}}
\newcommand{\mFPH}{\mF_{PH}}

\newcommand\SigmaP{\Sigma^P}
\newcommand\DeltaP{\Delta^{P}}
\newcommand\PiP{\Pi^{P}}
\newcommand\coSigmaP{\co\SigmaP}
\newcommand\coPiP{\mF{\co\PiP}}

%% autres classes:


%% Schemas pour complexité implicite

\newcommand\COMP{\cp{ COMP}}
\newcommand\CRN{\cp{ CRN}}
\newcommand\BRN{\cp{ BRN}}
\newcommand\SBRN{\cp{ SBRN}}
\newcommand\BR{\cp{ BR}}
\newcommand\kBR{k-\cp{ BR}}
\newcommand \WBPR{\cp{ WBPR}}
\newcommand \BMIN{\cp{ BMIN}}
\newcommand \BSUM{\cp{ BSUM}}
\newcommand\SBSUM{\cp{ SBSUM}}
\newcommand \BPROD{\cp{ BPROD}}
\newcommand \SBPROD{\cp{ SBPROD}}
\newcommand \SCOMP{\cp{ SCOMP}}
\newcommand \SR{\cp{ SR}}
\newcommand \SRN{\cp{ SRN}}

%% devrait etre défini
\newcommand\WBRN{\cp{ WBRN}}



%-- godel
\newcommand\INTDIV{\cp{INTDIV}}
\newcommand\DIV{\cp{DIV}}
\newcommand\EVEN{\cp{EVEN}}
\newcommand\ODD{\cp{ODD}}
\newcommand\PRIME{\cp{PRIME}}
\newcommand\POWER{\cp{POWER}}
\newcommand\BIT{\cp{BIT}}
\newcommand\negHP{\overline{\cp{HP}}}
\newcommand\HP{\cp{HP}}
\newcommand\negLuniv{\overline{\Luniv}}
\newcommand\VALCOMP{\cp{VALCOMP}}
\newcommand\mWINDOW{\cp{LEGAL}}
\newcommand\mmWINDOW{\cp{WINDOW}}
\newcommand\LENGTH{\cp{LENGTH}}
\newcommand\DIGIT{\cp{DIGIT}}
\newcommand\MATCH{\cp{MATCH}}
\newcommand\MOVE{\cp{MOVE}}
\newcommand\START{\cp{START}}
\newcommand\CELL{\cp{CELL}}
\newcommand\HALT{\cp{HALT}}
\newcommand\SUBST{\cp{SUBST}}
\newcommand\PROOF{\cp{PROOF}}
\newcommand\PROVABLE{\cp{PROVABLE}}
\newcommand\CONSIST{\cp{CONSIST}}



%---------------------
%--- calculus (AP) ---
%---------------------

% \jacobian{f}{x}
\newcommand{\jacobian}[1]{J_{#1}}
%\newcommand{\grad}[1]{\overrightarrow{\operatorname{grad}}~{#1}}
\newcommand{\grad}[1]{\nabla{#1}}
\newcommand{\scalarprod}[2]{{#1}\cdot{#2}}
\newcommand{\bigO}[1]{\mathcal{O}\left(#1\right)}
\newcommand\smallO{o}
\newcommand{\softO}[1]{\tilde{\mathcal{O}}\left(#1\right)}
\newcommand{\taylor}[3]{{T_{#1}^{#2}#3}}
\newcommand{\transpose}[1]{{#1}^T}
\newcommand{\pastsup}[2]{{\sup}_{#1}#2}

\DeclareMathOperator\arctanh{arctanh}



%---------------
%--- norms (AP) ---
%---------------


\newcommand{\inorm}[2]{\left\lVert{#1}\right\rVert_{#2}}
\newcommand{\twonorm}[1]{\inorm{#1}{2}}
\newcommand{\infnorm}[1]{\inorm{#1}{}}
\newcommand{\normsup}[1]{\infnorm{#1}{}}
%
\newcommand\hausdorff{d_H}
%
\newcommand\norm[1]{\infnorm{#1}}


%----------------
%--- topology ---
%----------------

% \openball{pt}{radius}
 \newcommand{\openball}[2]{B_{#2}(#1)}
\newcommand{\closedball}[2]{\ovl{B}_{#2}(#1)}
\newcommand{\closure}[1]{\ovl{#1}}
% \newcommand{\interior}[1]{\mathring{#1}}
% \newcommand{\border}[1]{\partial{#1}}




%-----------------
%--- functions (AP) ---
%-----------------
 
\newcommand{\lambdafun}[2]{(#1)\mapsto{#2}}
\newcommand{\lambdafunex}[2]{#1\mapsto{#2}}
\newcommand{\argmin}{\operatornamewithlimits{argmin}}
\newcommand{\argmax}{\operatornamewithlimits{argmax}}


 
%-----------------------
%--- Machines de Turing  ---
%-----------------------

\newcommand\blanc{\vectorl{B}}



%-----------------
%--- Logique ---
%-----------------
 
\newcommand\sem[1]{{[\semspace[}#1{]\semspace]}}
\newcommand\sems[1]{{[\semspace[}#1{]\semspace]}}
\newcommand\semss[1]{\sem{#1}_{\sigma'}}
\newcommand\semspace{\kern -.25ex}
\newcommand\true{true}
\newcommand\false{false}
%extension élémentaire:
\newcommand\elem{\preceq}
\newcommand\godel[1]{\lceil #1 \rceil}

%-----------------
%--- Model theory ---
%-----------------

\newcommand\hull[1]{\langle #1 \rangle}

%-----------------
%--- Ordinaux ---
%-----------------

\newcommand\omegaunCK{\omega_1^{CK}}
\newcommand\omegaunck{\omegaunCK}

\DeclareMathOperator{\cof}{cof}

%-----------------
%--- Surréels (QG) ---
%-----------------

\newcommand{\On}{\textbf{On}}
\newcommand{\Nobf}{\textbf{No}}
\newcommand{\Nobb}{\ensuremath{\mathbbmss{No}}}
\newcommand{\Nolt}[1]{\ensuremath{\Nobf_{< #1}}}
\newcommand{\Noleq}[1]{\ensuremath{\Nobf_{\leq #1}}}
\newcommand{\Noltbb}[1]{\ensuremath{\Nobb_{< #1}}}
\providecommand\No{\Nobf}
\newcommand\Noalpha{\Nolt{\alpha}}
\newcommand\Nolambda{\Nolt{\lambda}}

\DeclareMathOperator{\length}{length}
\DeclareMathOperator{\supp}{supp} % operateur support


%----------------------
%---Maths  de base (AP)  ---
%----------------------

\DeclareMathOperator{\dom}{dom} % operateur support

\DeclareMathOperator{\size}{size}

\newcommand{\powerset}[1]{\mathcal{P}(#1)}
\newcommand{\card}[1]{{\##1}}
\newcommand\priv{ - }
\newcommand\prive{-}
\renewcommand{\restriction}{\mathord{\upharpoonright}}
\newcommand\restrict[1]{_{\restriction #1}}

% \providecommand\mod{}
%% UNDEFINE:
\let\mod\relax
\let\div\relax
\DeclareMathOperator{\mod}{mod} 
\DeclareMathOperator{\div}{div} 

% partieentier
\newcommand\entieres[1]{\lceil #1 \rceil}
\newcommand\entierei[1]{\lfloor #1 \rfloor}


%----------------------
%---Symbole danger ---
%----------------------


%\providecommand*{\TakeFourierOrnament}[1]{{%
%\fontencoding{U}\fontfamily{futs}\selectfont\char#1}}
%\providecommand*{\danger}{{\Huge \TakeFourierOrnament{66}}}



%%%%%%%%%%%%%%%%%%%%%%%%%%%%%%%%%%%%%%%%%%%%%%%%%
%
%
%                      ASMeries
%
%
%%%%%%%%%%%%%%%%%%%%%%%%%%%%%%%%%%%%%%%%%%%%%%%%%



% CIRCUITS
%-- particuliers
\newcommand\REPLACE{\cp{REPLACE}}
\newcommand\EQUAL{\cp{EQUAL}}
\newcommand\EVALUE{\cp{EVALUE}}
\newcommand\TVALUE{\cp{TVALUE}}
\newcommand\UPDATE{\cp{UPDATE}}
\newcommand\ASSIGN{\cp{ASSIGN}}
\newcommand\PAR{\cp{PAR}}
\newcommand\DO{\cp{DO}}


% ASM
\newcommand\emplacement[2]{(#1,#2)}
\newcommand\contenu[2]{\sem{(#1,#2)}}
\newcommand\affecte[3]{(#1,#2)\leftarrow#3}
\newcommand\ctl{ctl\_state}


%%%%%%%%%%%%%%%%%%%%%%%%%%%%%%%%%%%%%%%%%%%%%%%%%
%
%
%                      GPACqueries
%
%
%%%%%%%%%%%%%%%%%%%%%%%%%%%%%%%%%%%%%%%%%%%%%%%%%



%------------------------
%--- generable zoo (GPAC) ---
%------------------------

\newcommand{\myop}[1]{\operatorname{#1}}
\newcommand{\abs}{\myop{abs}}
\newcommand{\sg}{\myop{sg}}
\newcommand{\ip}[1]{\myop{ip}_{#1}}
\newcommand{\rnd}{\myop{rnd}}
\newcommand{\lil}{\myop{lil}}
\newcommand{\lxh}{\myop{lxh}}
\newcommand{\hxl}{\myop{hxl}}
\newcommand{\plil}{\myop{plil}}
\newcommand{\reach}{\myop{reach}}
%\newcommand{\mreach}{\myop{mreach}}
\newcommand{\watchdog}{\myop{watchdog}}
\newcommand{\monitor}{\myop{monitor}}
%\newcommand{\norm}{\myop{norm}}
\newcommand{\mx}{\myop{mx}}
\newcommand{\mn}{\myop{mn}}
\newcommand{\sample}{\myop{sample}}
\newcommand{\select}{\myop{select}}
\newcommand{\ssine}[1]{\sin_{#1}}
\newcommand{\clamp}{\myop{clamp}}

\newcommand{\fnsg}[3]{tanh((#1)*(#2)*(#3))}
\newcommand{\fnipone}[3]{(1+\fnsg{(#1)-1}{#2}{#3})/2.}
\newcommand{\fnipinf}[1]{#1-sin(2*pi*(#1))/2./pi}%
\newcommand{\fnlil}[5]{%
\fnipone{(#3)-(#1)+1-((#2)-(#1))/8.}{#4+log(1+4./(#2-#1)*(#5)*(#5))+log(1+4)}{8./(#2-#1)}%
*\fnipone{#2-#3+1-(#2-#1)/8.}{#4+log(1+4/(#2-#1)*(#5)*(#5))+log(1+4)}{8./(#2-#1)}%
*4./(#2-#1)*(#5)}
\newcommand{\fnlxh}[5]{\fnipone{(#3)-(#1+#2)/2.+1}{#4+log(1+(#5)*(#5))}{2./(#2-#1)}*(#5)}
\newcommand{\fnhxl}[5]{\fnipone{(#1+#2)/2.-(#3)+1}{#4+log(1+(#5)*(#5))}{2./(#2-#1)}*(#5)}
\newcommand{\fnplilf}[4]{sin(((#4)-((#1)+(#2))/2.+(#3)/4.)*2.*pi/(#3))}
\newcommand{\fnplil}[6]{(#6)*\fnlxh{\fnplilf{#1}{#2}{#3}{#1}}%
{\fnplilf{#1}{#2}{#3}{#1+((#2)-(#1))/4.}}%
{\fnplilf{#1}{#2}{#3}{#4}}%
{#5+2.+log(1+(#6)*(#6))}%
{1./4.+2./((#2)-(#1))}}
\newcommand{\fnssine}[2]{sin(#2)*(1-cos(#2)/cos(#1))}



%--------------------
%--- complexity GPAC ---
%--------------------


%\ifthenelse{\equal{#1}{}}{Empty.}{Nonempty.}
\newcommand{\myclass}[1]{\operatorname{#1}}
% \newcommand{\gcomp}[2]{\ensuremath{\myclass{GCOMP}(#1,#2)}}
% \newcommand{\ggenerable}[1]{\textcolor{red}{\ensuremath{#1-}\text{generable}}}
 \newcommand{\gval}[2][]{\ensuremath{\myclass{GVAL}_{#1}\ifthenelse{\equal{#2}{}}{}{[#2]}}}
 \newcommand{\gpval}[1][]{\ensuremath{\myclass{GPVAL}_{#1}}}
 \newcommand{\gc}[3][]{\ensuremath{\myclass{ATSC}(#2,#3)}}
 \newcommand{\gpc}[1][]{\ensuremath{\myclass{ATSP}}}
\newcommand{\cgc}[1][]{\ensuremath{\myclass{ATS}}}
\newcommand{\gexpc}[1][]{\ensuremath{\myclass{AEXP}}}
\newcommand{\gwc}[2]{\ensuremath{\myclass{AW}(#1,#2)}}
\newcommand{\gpwc}{\ensuremath{\myclass{AWP}}}
\newcommand{\cgwc}{\ensuremath{\myclass{AW}}}
\newcommand{\grc}[3]{\ensuremath{\myclass{AR}(#1,#2,#3)}}
\newcommand{\gprc}{\ensuremath{\myclass{ARP}}}
\newcommand{\gsc}[3]{\ensuremath{\myclass{AS}(#1,#2,#3)}}
\newcommand{\gpsc}{\ensuremath{\myclass{ASP}}}
\newcommand{\goc}[3]{\ensuremath{\myclass{AO}(#1,#2,#3)}}
\newcommand{\gpoc}{\ensuremath{\myclass{AOP}}}
\newcommand{\cgoc}{\ensuremath{\myclass{AO}}}
\newcommand{\guc}[4]{\ensuremath{\myclass{AX}(#1,#2,#3,#4)}}
\newcommand{\gpuc}{\ensuremath{\myclass{AXP}}}
\newcommand{\glc}[2][]{\ensuremath{\myclass{AL}(#2)}}
\newcommand{\gplc}[1][]{\ensuremath{\myclass{ALP}}}
\newcommand{\cglc}[1][]{\ensuremath{\myclass{AL}}}

%\newcommand{\intinterv}[2]{\{#1,\ldots,#2\}}
\newcommand{\intinterv}[2]{\llbracket#1,#2\rrbracket}

\newcommand{\Rgen}{\R_{G}}
\newcommand{\Rpoly}{\R_{P}}

%%%%%%%%%%%%%%%%%%%%%%%%%%%%%%%%%%%%%%%%%%%%%%%%%
%
%
%                      Mes codages obscures
%
%
%%%%%%%%%%%%%%%%%%%%%%%%%%%%%%%%%%%%%%%%%%%%%%%%%




%%% Codages

% \newcommand\gammaoc{\gamma_{[0,1]}}
% \newcommand\gammaonc{\gamma_{\N}}
\newcommand\gammaTMc{\gamma^{TM}_{[0,1]}}
\newcommand\gammaTMcatan{\gamma^{TM}_{(0,1)^2}}
\newcommand\gammaTMnc{\gamma^{TM}_{\N^3}}
\newcommand\gammaTMncd{\gamma^{TM}_{\N^2}}
% \newcommand\gammaTMe{\gamma^{TM}_{*}}
% \newcommand\gammas{\gamma_{sab}}
% \newcommand\gammaord{\gamma_{cantor}}


\newcommand\enccompact{\nu_{X}}
 \newcommand\encinfinite{\nu_\N}
 \newcommand\enc{\nu}
 
%%%%%%%%%%%%%%%%%%%%%%%%%%%%%%%%%%%%%%%%%%%%%%%%%
%
%
%                      DESSINS ET IMAGES
%
%
%%%%%%%%%%%%%%%%%%%%%%%%%%%%%%%%%%%%%%%%%%%%%%%%%


\newcommand\IMAGE[2]{ \begin{center} 
    \includegraphics[#1]{#2}
  \end{center}
  }


  
%%%% POUR FAIRE DESSIN.
\newenvironment{dessinpp}[1]{
\begin{tikzpicture}[baseline=(current bounding
      box.west),#1]
}
{\end{tikzpicture}}




\newenvironment{dessinp}[1]{
\begin{center}\begin{tikzpicture}[baseline=(current bounding
      box.north),#1]
}
{\end{tikzpicture}
\end{center}}


\newenvironment{dessin}{
\begin{center}\begin{tikzpicture}[baseline=(current bounding
      box.north)]
}
{\end{tikzpicture}
\end{center}}



\newenvironment{dessind}{
\begin{tikzpicture}[baseline=(current bounding
      box.north)]
}
{\end{tikzpicture}
}



%%%%%%%%%%%%%%%%%%%%%%%%%%%%%%%%%%%%%%%%%%%%%%%%%
%
%
%                      POUR DESSINS GPAC
%
%
%%%%%%%%%%%%%%%%%%%%%%%%%%%%%%%%%%%%%%%%%%%%%%%%%


%%%% MACROS DE BASE.


\newcommand\constant[2]
{
node (#1)  [shape=circle,draw] {#2};
\draw (#1.east) -- +(0.3,0) coordinate (#1-o) ; 
\draw (#1.west) -- +(-0.3,0) coordinate (#1-i) ; 
}

\newcommand\basicproduct[2]
{
node (#1) {#2};
\draw (#1) +(-0.4,0.4) -- +(0.4,0.4) -- +(0.8,0) -- +(0.4,-0.4) -- +(-0.4,-0.4) --
+(-0.4,0.4);
\draw (#1) ++(0.8,0) -- ++(+0.3,0) coordinate (#1-o) ;
\draw (#1) ++ (-0.4,0.2) -- +(-0.3,0) coordinate (#1-i1) ;
\draw (#1) ++(-0.4,-0.2) -- +(-0.3,0) coordinate (#1-i2) ;
}

\newcommand\product[1]{\basicproduct{#1}{$+\Pi$}}
\newcommand\negativeproduct[1]{\basicproduct{#1}{$-\Pi$}}


%%BEGIN INTERNAL
\newcommand\basicsummer[1]{
coordinate (#1) {};
\draw (#1) +(-0.5,0.6) -- +(0.5,0) -- +(-0.5,-0.6) -- +(-0.5,0.6);
\draw (#1) ++(0.5,0) -- +(+0.3,0) coordinate (#1-o) ;
}

\newcommand\basicintegrator[1]{
coordinate (#1) {};
\draw (#1) +(-0.5,0.6) -- +(0.5,0) -- +(-0.5,-0.6) -- +(-0.5,0.6);
\draw (#1) ++(0.5,0) -- +(+0.3,0) coordinate (#1-o) ;
\draw (#1) +(-0.5,-0.6) -| +(-0.8,0.6) -- +(-0.5,0.6) ;
 \draw (#1) ++(-0.65,0.6) -- +(0,+0.3) coordinate (#1-init) ;
}

\newcommand\splitfour[2]{
\draw (#2) ++ (#1,0.4) -- +(-0.3,0) coordinate (#2-i1) ;
\draw (#2) ++ (#1,0.15) -- +(-0.3,0) coordinate (#2-i2) ;
\draw (#2) ++ (#1,-0.15) -- +(-0.3,0) coordinate (#2-i3) ;
\draw (#2) ++ (#1,-0.4) -- +(-0.3,0) coordinate (#2-i4) ;
}

\newcommand\splitthree[2]{
\draw  (#2) ++(#1,0.3) -- +(-0.3,0) coordinate (#2-i1) ;
\draw  (#2) ++(#1,0) -- +(-0.3,0) coordinate (#2-i2) ;
\draw  (#2) ++(#1,-0.3) -- +(-0.3,0) coordinate (#2-i3) ;
}

\newcommand\splittwo[2]{
\draw  (#2) ++(#1,0.2) -- +(-0.3,0) coordinate (#2-i1) ;
\draw  (#2) ++(#1,-0.2) -- +(-0.3,0) coordinate (#2-i2) ;
}

\newcommand\splitone[2]{
\draw (#2) ++(#1,0) -- +(-0.3,0) coordinate (#2-i1) ;
}

%% END INTERNAL

\newcommand\summerfour[1]{
\basicsummer{#1}
\splitfour{-0.5}{#1}
}

\newcommand\summerthree[1]{
\basicsummer{#1}
\splitthree{-0.5}{#1}
}

\newcommand\summertwo[1]{
\basicsummer{#1}
\splittwo{-0.5}{#1}
}

\newcommand\summerone[1]{
\basicsummer{#1}
\splitone{-0.5}{#1}
}


\newcommand\integratorfour[1]{
\basicintegrator{#1}
\splitfour{-0.8}{#1}
<}


\newcommand\integratorthree[1]{
\basicintegrator{#1}
\splitthree{-0.8}{#1}
}


\newcommand\integratortwo[1]{
\basicintegrator{#1}
\splittwo{-0.8}{#1}
}


\newcommand\integratorone[1]{
\basicintegrator{#1}
\splitone{-0.8}{#1}
}

%%%%% FIN MACROS




%%%%%%%%%%%%%%%%%%%%%%%%%%%%%%%%%%%%%%%%%%%%%%%%%
%
%
%                      Trucs d'Amaury
%
%
%%%%%%%%%%%%%%%%%%%%%%%%%%%%%%%%%%%%%%%%%%%%%%%%%


%----------------
%--- ref (AP) ---
%----------------

% \newcommand{\defref}[1]{Definition~\ref{#1}}
\newcommand{\lemref}[1]{Lemma~\ref{#1}}
\newcommand{\thref}[1]{Theorem~\ref{#1}}
\newcommand{\amaurycorref}[1]{Corollary~\ref{#1}}
 \newcommand{\figref}[1]{Figure~\ref{#1}}
% \newcommand{\figrefmult}[1]{Figures~{#1}}
% \newcommand{\secref}[1]{Section~\ref{#1}}
% %\renewcommand{\algref}[1]{Algorithm~\ref{#1}}
 \newcommand{\propref}[1]{Proposition~\ref{#1}}
% \newcommand{\exref}[1]{Example~\ref{#1}}
\newcommand{\remref}[1]{Remark~\ref{#1}}


 %-------------------
%--- polynomials ---
%-------------------

\newcommand{\sigmap}[1]{{\Sigma{#1}}}
\newcommand{\degp}[1]{{\operatorname{deg}(#1)}}
\newcommand{\poly}{\operatorname{poly}}


 %----------------------
%--- misc functions ---
%----------------------

\newcommand{\sgn}[1]{\operatorname{sgn}(#1)}
\newcommand{\intp}{\operatorname{int}}
 \newcommand{\fracp}{\operatorname{frac}}
\newcommand{\fiter}[2]{#1^{[#2]}}
\newcommand{\frestrict}[2]{#1\restriction_{#2}}
\newcommand{\indicator}[1]{\mathds{1}_{#1}}
\newcommand{\idfun}{\operatorname{id}}
% \newcommand{\floor}[1]{\left\lfloor#1\right\rfloor}
 \newcommand{\ceil}[1]{\left\lceil#1\right\rceil}
\newcommand{\round}[1]{\left\lfloor#1\right\rceil}
\DeclareMathOperator\erf{erf}


%------------
%--- sets ---
%------------

\newcommand{\A}{\mathbb{A}}
%\newcommand{\D}{\mathbb{D}}
\newcommand{\Ns}{\mathbb{N}^*}
%\newcommand{\C}{\mathbb{C}}
%\newcommand{\K}{\mathbb{K}}
% \MAT{height}{[width]}{[field]}
\newcommand{\MAT}[3]{M_{#1\ifthenelse{\equal{#2}{}}{}{,#2}}\ifthenelse{\equal{#3}{}}{}{\left(#3\right)}}

%\newcommand{\Rcomp}{\R_{c}}
\newcommand{\Rp}{\R_{+}}
\newcommand{\Rs}{\R^{*}}
\newcommand{\Rps}{\R_{+}^{*}}

%\newcommand{\Analytic}{C^\omega}
%\newcommand{\PTIME}{\ensuremath{\operatorname{P}}}
%\newcommand{\EXPTIME}{\ensuremath{\operatorname{EXPTIME}}}
%\newcommand{\NP}{\ensuremath{\operatorname{NP}}}
%\newcommand{\NPTIME}{\NP}
\newcommand{\FP}{\ensuremath{\operatorname{FP}}}
%\newcommand{\PSPACE}{\ensuremath{\operatorname{PSPACE}}}
% \newcommand{\FPSPACE}{\ensuremath{\operatorname{FPSPACE}}}



%%%%%%%%%%%%%%%%%%%%%%%%%%%%%%%%%%%%%%%%%%%%%%%%%
%
%
%                      Trucs de Quentin
%
%
%%%%%%%%%%%%%%%%%%%%%%%%%%%%%%%%%%%%%%%%%%%%%%%%%


\newcommand{\fct}[5]{\ensuremath{#1 : \left\{\begin{array}{c c c}
#2 & \rightarrow & #3\\
#4 & \mapsto & #5
                                             \end{array}\right.}}

%Nouvelle fraction, plus jolie
\newcommand{\f}[2]{	
	\mathchoice%
		{\dfrac{#1}{#2}}
    	{\dfrac{#1}{#2}}
		{\frac{#1}{#2}}
		{\frac{#1}{#2}}
}


%%Ensembles et combinatoires

%Ensemble avec propriété à drotie \{x | blabla\}
\newcommand{\enstq}[2]{\left\{ \left.#1\ \vphantom{#2}\right|\ #2  \right\}}
%meme chose mais avec des crochets
\newcommand{\crotq}[2]{\left[ \left.#1\ \vphantom{#2}\right|\ #2  \right]}
%même chose mais avec des brackets
\newcommand{\bratq}[2]{\left\langle \left.#1\ \vphantom{#2}\right|\ #2  
  \right\rangle}

\newcommand{\intn}[2]{\ensuremath{
		\left\llbracket\, #1\ ;\ #2\,\right\rrbracket
              }}


%%%%%%%%%%%%%%%%%%%%%%%%%%%%%%%%%%%%%%%%%%%%%%%%%
%
%
%               ENVIRONNEMENTS THEOREMS & CO
%
%
%%%%%%%%%%%%%%%%%%%%%%%%%%%%%%%%%%%%%%%%%%%%%%%%%


%----------------
%--- pas dans lipics ---
%-----------------


\ifnotSLIDE{
% commente car bug Fevrier 2026: \ifnotTDEXAM{\spnewtheorem{solution}{Solution}{\bfseries}{\rmfamily}}
\spnewtheorem{openproblem}{Open Problem}{\bfseries}{\rmfamily}
\spnewtheorem{remarkolivier}{Commentaire d'Olivier: }{\bfseries}{\rmfamily}
\spnewtheorem{postulate}{Postulate}{\bfseries}{\rmfamily}
}

%----------------
%--- dans lipics ---
%-----------------
%
\newenvironment{exercice}[1]{\begin{exercise} \label{exo:#1-stt}}{\end{exercise}}
\newenvironment{exerciced}[1]{\begin{exercisee}\label{exo:#1-stt}}{\end{exercisee}}
\newenvironment{exercicec}[1]{\begin{exercise} \label{exo:#1-stt} (\LANGUE{solution on}{corrigé} page \pageref{exo:#1-cor})
}{\end{exercise}}
\newenvironment{exercicedc}[1]{\begin{exercisee} \label{exo:#1-stt} (\LANGUE{solution on}{corrigé} page \pageref{exo:#1-cor})
}{\end{exercisee}}

\newenvironment{slideproposition}{{\bf Proposition}}{}

\ifnotSLIDE{
\nolipics{
\nolncs{\spnewtheorem{theorem}{\LANGUE{Theorem}{Théorème}}{\bfseries}{\rmfamily}}
\nolncs{\spnewtheoremi{definition}{\LANGUE{Definition}{Définition}}{\bfseries}{\rmfamily}{theorem}}
\nolncs{\spnewtheoremi{lemma}{\LANGUE{Lemma}{Lemme}}{\bfseries}{\rmfamily}{theorem}}
\nolncs{\spnewtheoremi{corollary}{\LANGUE{Corollary}{Corollaire}}{\bfseries}{\rmfamily}{theorem}}
\nolncs{\spnewtheoremi{proposition}{Proposition}{\bfseries}{\rmfamily}{theorem}}
\nolncs{\spnewtheoremi{example}{\LANGUE{Example}{Exemple}}{\bfseries}{\rmfamily}{theorem}}
%\spnewtheorem{sketch}{{\bf \LANGUE{Sketch of proof}{Démonstration (principe)}}}{\bfseries}{\rmfamily}
\nolncs{\spnewtheoremi{remark}{\LANGUE{Remark}{Remarque}}{\bfseries}{\rmfamily}{theorem}}
\nolncs{\spnewtheorem{exercise}{\LANGUE{Exercise}{Exercice}}{\bfseries}{\rmfamily}}
\spnewtheorem{exercisee}{\LANGUE{*Exercise}{*Exercice}}{\bfseries}{\rmfamily}
\spnewtheorem{summary}{\LANGUE{Summary}{Résumé}}{\bfseries}{\rmfamily}
 %\spnewtheorem{exerciceenglish}{Exercise}{\bfseries}{\rmfamily}
 %
 \newenvironment{sketch}{{\bf \LANGUE{Sketch of proof}{Démonstration (principe)}}:}{\hfill $\Box$ }
\nolncs{\spnewtheorem{conjecture}{\LANGUE{Conjecture}{Conjecture}}{\bfseries}{\rmfamily}}
 %%
 \spnewtheorem{theoremaprouver}{\LANGUE{Theorem TO BE PROVED}{Théorème (A PROUVER)}}{\bfseries}{\rmfamily}
  \spnewtheorem{conjectureaprouver}{\LANGUE{Conjecture TO BE VERIFIED}{Conjecture (A VERIFIER)}}{\bfseries}{\rmfamily}
 }}
            
 \makeatletter
 \ifnotSLIDE{
  \ifnotTDEXAM{
\@ifundefined{proof}{
	\newenvironment{proof}{{\bf \LANGUE{Proof}{Démonstration}}:}{\hfill $\Box$ 
	}{}
	}}}
\makeatother

 
%%%%%%%%%%%%%%%%%%%%%%%%%%%%%%%%%%%%%%%%%%%%%%%%%
%
%
%                      Macros dont pas fier du tout
%
%
%%%%%%%%%%%%%%%%%%%%%%%%%%%%%%%%%%%%%%%%%%%%%%%%%



\newcommand\equations[1]{
\begin{array}{lll}
#1
\end{array}
}



\newcommand\systeme[1]{
\left\{
\begin{array}{lll} 
#1
\end{array}
\right.
}



%------------------
%--- shorthands (AP)---
%------------------
\newcommand{\mtt}[1]{\mathtt{#1}}
\newcommand{\ovl}[1]{\overline{#1}}
\newcommand{\panic}[1]{\textcolor{red}{#1}}
%\NewEnviron{epanic}{{\color{red}\BODY}}
\newcommand{\idest}{\emph{i.e.\thinspace}}


%------------------
%--- mes trucs  ---
%------------------


\newcommand\implique{\Rightarrow}
\newcommand\ssi{\Leftrightarrow}
\newcommand\ac[1]{\{#1\}}
\newcommand\zero{\vectorl{0}}
\newcommand\un{\vectorl{1}}

\newcommand\technique[1]{{\scriptsize #1}}


%% Preuve and Co

\newcommand\sensdirect{ Direction $\Rightarrow$. }
\newcommand\sensindirect{ Direction $\Leftarrow$. }




\newcommand\donnee{\mbox{<donnée>}}



% pour aller vite

\newcommand\gM{\mathfrak{M}}
\newcommand\gMM{(M,f_1,\cdots,f_u,r_1,\cdots,r_v)}
%
\newcommand\mygM{\mK}
\newcommand\mygMM{\left ( \K, \{op_i\}_{i\in I}, r_1 \ldots,
r_\ell, \zero,\un \right )}
\newcommand\myM{\K} 

\newcommand\gMMM{(M,f_1,\cdots,f_u,c_1,\cdots,c_k,r_1,\cdots,r_v)}
\newcommand\gMMMM{(f_1,\cdots,f_u,r_1,\cdots,r_v)}
\newcommand\gMME{(M,f_1,\cdots,f_u,f'_1,\cdots,f'_\ell,r_1,\cdots,r_v)}
\newcommand\gMMS{(f_1,\cdots,f_u,f'_1,\cdots,f'_\ell,r_1,\cdots,r_v)}

\newcommand\bool{\mathfrak{B}}
\newcommand\gbool{(\{0,1\},\zero,\un,=)}

\newcommand\osg{\overline{sg}}
\newcommand\moins{\mathrel{\dot{-}}}


%%% TRUC DE BARWISE


\newcommand\UM{\kU_{\kM}}
\newcommand\VV{\mathbb{V}}
% Hereditarily finite
\newcommand\IHF{\mathbb{I}HF}
\newcommand\IHP{\IHYP}
\newcommand\IHYP{\mathbb{I}HYP}

\providecommand\citep[1]{\cite{#1}}

%% Cofinali


% probas
% défini dans amsmath
\renewcommand\Pr{\cp{Pr}}
\newcommand\Var{\cp{Var}}


% -- mesures
%\newcommand\card{card}
\newcommand\taille{\LANGUE{size}{taille}}
\newcommand\entree{e}
\newcommand\profondeur{\LANGUE{depth}{profondeur}}


%%%%%%%%%%%%%%%%%%%%%%%%%%%%%%%%%%%%%%%%%%%%%%%%%%%%%%%%%%%%%%%%%%%%%%%%%%%%
%                 MACHINES DE TURING (graphique)
%%%%%%%%%%%%%%%%%%%%%%%%%%%%%%%%%%%%%%%%%%%%%%%%%%%%%%%%%%%%%%%%%%%%%%%%%%%%



% Configuration de machine de Turing
\newcommand\configTuring[3]{#1\textbf{#2}#3}


%%% DESSIN CROIX PORU AIDER  A DESSINER

\newcommand\croix{  +(-1,-1) -- +(1,1) +(-1,1) --
  +(1,-1) }

%% DESSIN ACCOLADE

\newcommand\ACCOLADE[3]{
%#1: coordonnées de départ
%#2: coordonnées arrivée
%#3: texte
\draw[decorate,decoration={brace}]  (#1) -- (#2)
 node[midway,above=0.1cm] {#3};
}

%%%


\def\lcase{3.5 ex}
\def\hcase{3.5ex}
\tikzstyle{case} = []
%[draw,minimum width = \lcase,minimum height = 3.6ex,baseline]
\newcommand\dcase{ +(-0.5*\lcase,-0.5*\hcase) rectangle +(0.5*\lcase,0.5*\hcase)}
%\newcommand\lettre[1]{\node [name=l,case,on chain=1] {#1};}


\newcommand\RUBANMACHINEDETURING[6]{
%#1: texte
%#2: position de la tête
%#3: nb de blancs à gauche
%#4: nb de cases totales
%#5: texte au dessus du ruban
%#6: position du coin
    \draw (#6) +(\lcase,4ex ) node {#5};
    \draw (#6) +(-0.6,0)  node  {\ldots};  
    \draw (#6) node[coordinate] (t) {};
    \foreach \x in {1,2,...,#3} {
         \draw (t) node [case]   {$\blanc$} \dcase;
         \draw (t) ++(\lcase,0) node [coordinate] (t) {};
       };
   \StrLen{#1}[\Resultat] % ATTENTION IL FAUT CETTE SYNTAXE POUR FAIRE
                         % EQUIVALENT D UN EDEF
  
\foreach \x in {1,...,\Resultat} 
             { 
        \draw (t) node [case]   {\StrChar{#1}{\x}} \dcase;
        \draw (t) ++(\lcase,0) node [coordinate] (t) {};
             }
\pgfmathtruncatemacro{\z}{#4 - \Resultat - #3}

  \foreach \x in {1,2,..., \z} {
        \draw (t) node [case]   {$\blanc$} \dcase;
        \draw (t) ++(\lcase,0) node [coordinate] (t) {};
}
%% Convention: graphiquement, la case numéro n se trouve à (n*\lcase,0)
%% en comptant les case à partir de $0$

   \draw  (t) +(0.1,0) node [name=r] {\ldots}; 
}

\newcommand\RUBANMACHINEDETURINGtexte[6]{
%#1: texte
%#2: position de la tête
%#3: nb de blancs à gauche
%#4: nb de cases totales (équivalent)
%#5: texte au dessus du ruban
%#6: position du coin
    \draw (#6) +(\lcase,4ex ) node {#5};
    \draw (#6) +(-0.6,0)  node  {\ldots};  
    \draw (#6) node[coordinate] (t) {};
   \foreach \x in {1,2,...,#3} {
        \draw (t) node [case]   {$\blanc$} \dcase;
        \draw (t) ++(\lcase,0) node [coordinate] (t) {};
      };
   \draw (#6) ++(-2*\lcase,0.5*\lcase)-- +(\lcase*#4,0);
   \draw (#6) ++(-2*\lcase,-0.5*\lcase) -- +(\lcase*#4,0);
   
    \draw (t) ++(-0.5*\lcase,0) node [case,anchor=west]  (te) {#1};
    \draw (te.east)  node [anchor=west,coordinate] (t) {};

  \draw (t) ++(0.2*\lcase,0) node [coordinate] (t) {};
   \foreach \x in {1,2,...,#3} {
        \draw (t) node [case]   {$\blanc$} \dcase;
        \draw (t) ++(\lcase,0) node [coordinate] (t) {};
      };

    \draw  (t) +(0.1,0) node [anchor=west] {\ldots}; 
}


\newcommand\RUBANMACHINEDETURINGTETE[6]{
%
% dessine la tete
%
% output: (tete)
%
\RUBANMACHINEDETURING{#1}{#2}{#3}{#4}{#5}{#6}
% TETE
     \draw (#6) + ( #2 * \lcase + #3  *\lcase,- 0.5*\lcase )
     node[coordinate] (k) {};
     \draw[->] (k) +(0,-0.5) -- (k);
     \draw (k) +(0,-0.5) node[coordinate] (tete) {};
}


\newcommand\RUBANMACHINEDETURINGTETEtexte[6]{
%
% dessine la tete
%
% output: (tete)
%
\RUBANMACHINEDETURINGtexte{#1}{#2}{#3}{#4}{#5}{#6}
% TETE
     \draw (#6) + ( #2 * \lcase + #3  *\lcase,- 0.5*\lcase )
     node[coordinate] (k) {};
     \draw[->] (k) +(0,-0.5) -- (k);
     \draw (k) +(0,-0.5) node[coordinate] (tete) {};
}


\newcommand\MACHINEDETURING[6]{
%#1: texte
%#2: position de la tête
%#3: état
%#4: nb de blancs à gauche
%#5: nb de cases totales
%#6: texte en haut du ruban
%\begin{tikzpicture}[node distance=-0.15mm]

\RUBANMACHINEDETURINGTETE{#1}{#2}{#4}{#5}{#6}{0,0}
%
%% ETAT:
     \draw (tete) node [name=etat,draw,circle,anchor=north]  {#3};
     \draw (etat) -- (tete);
%\end{tikzpicture}
}

\newcommand\MACHINEDETURINGtexte[6]{
%#1: texte
%#2: position de la tête
%#3: état
%#4: nb de blancs à gauche
%#5: nb de cases totales
%#6: texte en haut du ruban
%\begin{tikzpicture}[node distance=-0.15mm]

\RUBANMACHINEDETURINGTETEtexte{#1}{#2}{#4}{#5}{#6}{0,0}
%
%% ETAT:
     \draw (tete) node [name=etat,draw,circle,anchor=north]  {#3};
     \draw (etat) -- (tete);
%\end{tikzpicture}
}

\newcommand\MACHINEDETURINGTROISRUBANS[9]{
%#1: texte1
%#2: position de la tête 1
%#3: texte2
%#4: position de la tête 2
%#5: texte3
%#6: position de la tête 3
%#7: état
%#8: nb de blancs à gauche 1
%#9: nb de cases totales
%\begin{tikzpicture}[node distance=-0.15mm]

\RUBANMACHINEDETURINGTETE{#1}{#2}{#8}{#9}{\LANGUE{tape}{ruban} $1$}{0,0}
\draw (tete) node[coordinate] (tete1) {};
\draw (tete1) +(8,0) node[coordinate] (tete1b) {};
\draw (tete1) -| (tete1b) ;

\RUBANMACHINEDETURINGTETE{#3}{#4}{#8}{#9}{\LANGUE{tape}{ruban} $2$}{0,-2}
\draw (tete) node[coordinate] (tete2) {};
\draw (tete2) +(-8.5,0) node[coordinate] (tete2b) {};
\draw (tete2) -| (tete2b) ;

\RUBANMACHINEDETURINGTETE{#5}{#6}{#8}{#9}{\LANGUE{tape}{ruban} $3$}{0,-4}
\draw (tete) node[coordinate] (tete3) {};

%% ETAT:
     \draw (tete3) node [name=etat,draw,circle,anchor=north]  {#7};
     \draw (etat) -- (tete3);
    \draw  (etat) -|  (tete2b);
\draw  (etat) -|  (tete1b);
%\end{tikzpicture}
}

\newcommand\MACHINEDETURINGTROISRUBANSENGLISH[9]{
%#1: texte1
%#2: position de la tête 1
%#3: texte2
%#4: position de la tête 2
%#5: texte3
%#6: position de la tête 3
%#7: état
%#8: nb de blancs à gauche 1
%#9: nb de cases totales
%\begin{tikzpicture}[node distance=-0.15mm]

\RUBANMACHINEDETURINGTETE{#1}{#2}{#8}{#9}{tape $1$}{0,0}
\draw (tete) node[coordinate] (tete1) {};
\draw (tete1) +(8,0) node[coordinate] (tete1b) {};
\draw (tete1) -| (tete1b) ;

\RUBANMACHINEDETURINGTETE{#3}{#4}{#8}{#9}{tape $2$}{0,-2}
\draw (tete) node[coordinate] (tete2) {};
\draw (tete2) +(-8.5,0) node[coordinate] (tete2b) {};
\draw (tete2) -| (tete2b) ;

\RUBANMACHINEDETURINGTETE{#5}{#6}{#8}{#9}{tape $3$}{0,-4}
\draw (tete) node[coordinate] (tete3) {};

%% ETAT:
     \draw (tete3) node [name=etat,draw,circle,anchor=north]  {#7};
     \draw (etat) -- (tete3);
    \draw  (etat) -|  (tete2b);
\draw  (etat) -|  (tete1b);
%\end{tikzpicture}
}



\newcommand\MACHINEDETURINGTROISRUBANStexte[9]{
%#1: texte1
%#2: position de la tête 1
%#3: texte2
%#4: position de la tête 2
%#5: texte3
%#6: position de la tête 3
%#7: état
%#8: nb de blancs à gauche 1
%#9: nb de cases totales
% \begin{tikzpicture}[node distance=-0.15mm]

\RUBANMACHINEDETURINGTETEtexte{#1}{#2}{#8}{#9}{\LANGUE{tape}{ruban} $1$}{0,0}
\draw (tete) node[coordinate] (tete1) {};
\draw (tete1) +(8,0) node[coordinate] (tete1b) {};
\draw (tete1) -| (tete1b) ;

\RUBANMACHINEDETURINGTETEtexte{#3}{#4}{#8}{#9}{\LANGUE{tape}{ruban} $2$}{0,-2}
\draw (tete) node[coordinate] (tete2) {};
\draw (tete2) +(-8.5,0) node[coordinate] (tete2b) {};
\draw (tete2) -| (tete2b) ;

\RUBANMACHINEDETURINGTETEtexte{#5}{#6}{#8}{#9}{\LANGUE{tape}{ruban} $3$}{0,-4}
\draw (tete) node[coordinate] (tete3) {};

%% ETAT:
     \draw (tete3) node [name=etat,draw,circle,anchor=north]  {#7};
     \draw (etat) -- (tete3);
    \draw  (etat) -|  (tete2b);
\draw  (etat) -|  (tete1b);
%\end{tikzpicture}
}


\newcommand\MACHINEDETURINGDEUXRUBANStexte[9]{
%#1: texte1  
%#2: position de la tête 1 
%#3: texte2 inutilisée
%#4: position de la tête 2 inutilisée
%#5: texte3
%#6: position de la tête 3
%#7: état
%#8: nb de blancs à gauche 1
%#9: nb de cases totales
% \begin{tikzpicture}[node distance=-0.15mm]

\RUBANMACHINEDETURINGTETEtexte{#1}{#2}{#8}{#9}{\LANGUE{tape}{ruban} $1$}{0,0}
\draw (tete) node[coordinate] (tete1) {};
\draw (tete1) +(8,0) node[coordinate] (tete1b) {};
\draw (tete1) -| (tete1b) ;

% \RUBANMACHINEDETURINGTETEtexte{#3}{#4}{#8}{#9}{\LANGUE{tape}{ruban} $2$}{0,-1.5}
% \draw (tete) node[coordinate] (tete2) {};
% \draw (tete2) +(-5.5,0) node[coordinate] (tete2b) {};
% \draw (tete2) -| (tete2b) ;

\RUBANMACHINEDETURINGTETEtexte{#5}{#6}{#8}{#9}{\LANGUE{tape}{ruban} $2$}{0,-1.5}
\draw (tete) node[coordinate] (tete3) {};

%% ETAT:
     \draw (tete3) node [name=etat,draw,circle,anchor=north]  {#7};
     \draw (etat) -- (tete3);
%   \draw  (etat) -|  (tete2b);
\draw  (etat) -|  (tete1b);
%\end{tikzpicture}
}


\def\argdix{2}
\def\argonze{25}
\newcommand\MACHINEDETURINGQUATRERUBANStextep[9]{
%#1: texte1
%#2: position de la tête 1
%#3: texte2
%#4: position de la tête 2
%#5: texte3
%#6: position de la tête 3
%#7: texte4
%#8: position de la tête 3
%#9: état
%#10: nb de blancs à gauche 1
%#11: nb de cases totales
% \begin{tikzpicture}[node distance=-0.15mm]

%% BUG sur #10 devenu \argdix
%% BUG SUR #11 devenu \argonze

\RUBANMACHINEDETURINGTETEtexte{#1}{#2}{\argdix}{\argonze}{\LANGUE{tape}{ruban} $1$}{0,0}
% \draw (tete) node[coordinate] (tete1) {};
% \draw (tete1) +(8,0) node[coordinate] (tete1b) {};
% \draw (tete1) -| (tete1b) ;

\RUBANMACHINEDETURINGTETEtexte{#3}{#4}{\argdix}{\argonze}{\LANGUE{tape}{ruban} $2$}{0,-1.5}
% \draw (tete) node[coordinate] (tete2) {};
% \draw (tete2) +(-5.5,0) node[coordinate] (tete2b) {};
% \draw (tete2) -| (tete2b) ;

\RUBANMACHINEDETURINGTETEtexte{#5}{#6}{\argdix}{\argonze}{\LANGUE{tape}{ruban} $3$}{0,-3.2}
\draw (tete) node[coordinate] (tete3) {};

\RUBANMACHINEDETURINGTETEtexte{#7}{#8}{\argdix}{\argonze}{\LANGUE{tape}{ruban} $4$}{0,-4.9}
\draw (tete) node[coordinate] (tete4) {};

%% ETAT:
     \draw (tete4) node [name=etat,draw,circle,anchor=north]  {#9};
%      \draw (etat) -- (tete3);
%     \draw  (etat) -|  (tete2b);
% \draw  (etat) -|  (tete1b);
%\end{tikzpicture}
}


%%%%%%%%%%%%%%%%%%%%%%%%%%%%%%%%%%%%%%%%%%%%%%%%%%%%%%%%%%%%%%%%%%%%%%%%%%%%
%                 MACHINES DE TURING (automate)
%%%%%%%%%%%%%%%%%%%%%%%%%%%%%%%%%%%%%%%%%%%%%%%%%%%%%%%%%%%%%%%%%%%%%%%%%%%%



\newcommand\dec{1ex}
\newcommand\myACCOLADE[4]{
\ACCOLADE{#1*\lcase+#4*\lcase-0.5*\lcase,0.5*\lcase+\dec}{#1*\lcase+#4*\lcase-0.5*\lcase+#3*\lcase,0.5*\lcase+\dec}{#2}
}


\newcommand\M{\Sigma}


\newcommand\sommet[2]{\node[state] (#1) #2 {$#1$};}
\newcommand\sommetinitial[2]{\node[state,initial] (#1) #2 {$#1$};}
\newcommand\sommetfinal[2]{\node[state,accepting] (#1) #2 {$#1$};}
\newcommand\gtransition[5]{\draw (#1) edge node {#2/#4 $#5$} (#3);}
\newcommand\gtransitionba[5]{\draw (#1) [loop above] edge node {#2/#4
    $#5$} (#3);}
\newcommand\gtransitionbb[5]{\draw (#1) [loop below] edge node {#2/#4
    $#5$} (#3);}
\newcommand\gtransitionbleft[5]{\draw (#1) [loop left] edge node {#2/#4
    $#5$} (#3);}
\newcommand\gtransitionbright[5]{\draw (#1) [loop right] edge node {#2/#4
    $#5$} (#3);}
\newcommand\gtransitionbl[5]{\draw (#1) [bend left] edge node {#2/#4
    $#5$} (#3);}
\newcommand\gtransitionbr[5]{\draw (#1) [bend right] edge node {#2/#4
    $#5$} (#3);}
\newcommand\gtransitionbadouble[9]{\draw (#1) [loop above] edge node {
\begin{tabular}{l}
#2/#4 $#5$ \\
  #7/#8 $#9$ \\
\end{tabular}
} (#3);}

\newcommand\SCALE{}

\newcommand\dessinmtp[2]{
\begin{center}
\begin{tikzpicture}[->,>=stealth',shorten >=1pt,auto,node distance=2.3cm,semithick,#2] 
#1
\end{tikzpicture}
\end{center}
}

\newcommand\dessinmt[1]{\dessinmtp{#1}{scale=1}}


\newcommand\ANIMATIONMT[1]{
\begin{overprint}
\begin{dessinp}{scale=0.52}
\foreach \av/\q/\b  [count=\xi]
in  
{#1}
{ 
\only<\xi| handout 0>{ 
 \StrLen{\av}[\Resintm] % ATTENTION IL FAUT CETTE SYNTAXE POUR FAIRE
                         % EQUIVALENT D UN EDEF
 \MACHINEDETURING{\av \b}{\Resintm}{$\q$}{5}{20}{}
}}
 \end{dessinp}
\end{overprint}
}

\newcommand\ANIMATIONMTd[1]{
\begin{overprint}
\begin{dessinp}{scale=0.52}
\foreach \av/\q/\b  [count=\xi]
in  
{#1}
{ 
\only<\xi| handout 0>{ 
 \StrLen{\av}[\Resintm] % ATTENTION IL FAUT CETTE SYNTAXE POUR FAIRE
                         % EQUIVALENT D UN EDEF
 \hspace{-1cm}\MACHINEDETURING{\av \b}{\Resintm}{$\q$}{5}{33}{}
}}
 \end{dessinp}
\end{overprint}
}

\def\MAXDET{20}

\newcommand\DIAGRAMMEESPACETEMPS[1]{
\draw node[coordinate] (tt) {};
\foreach \av/\q/\b  in  
{#1}
{ 

\RUBANMACHINEDETURING{\av{$\q$}\b}{0}{4}{\MAXDET}{}{tt};
\draw (tt) +(0,-\hcase) node[coordinate] (tt) {};
}
%\foreach \x in {0,...,19} {\draw (tt) + (\x*\lcase,0) node {$\vdots$};}
}



\newcommand\progun{
/q_0/0011, % initial
X/q_1/011,X0/q_1/11,X/q_2/0Y1,/q_2/X0Y1,
X/q_0/0Y1,XX/q_1/Y1,XXY/q_1/1, XX/q_2/YY, X/q_2/XYY,
XX/q_0/YY,XXY/q_3/Y,XXYY/q_3/B,XXYYB/q_4/B
}

\newcommand\DIAGRAMMEESPACETEMPSun{
\DIAGRAMMEESPACETEMPS
%% RECOPIE progun
{/q_0/0011, % initial
X/q_1/011,X0/q_1/11,X/q_2/0Y1,/q_2/X0Y1,
X/q_0/0Y1,XX/q_1/Y1,XXY/q_1/1, XX/q_2/YY, X/q_2/XYY,
XX/q_0/YY,XXY/q_3/Y,XXYY/q_3/B,XXYYB/q_4/B}
%% FIN RECOPIE
}

\newcommand\progdeux{
/q_0/0010, X/q_1/010, X0/q_1/10, X/q_2/0Y0,/q_2/X0Y0,
X/q_0/0Y0,XX/q_1/Y0,XXY/q_1/0,XXY0/q_1/%B,
}


\newcommand\DIAGRAMMEESPACETEMPSdeux{
\DIAGRAMMEESPACETEMPS
{
%% RECOPIE progdeux
/q_0/0010, X/q_1/010, X0/q_1/10, X/q_2/0Y0,/q_2/X0Y0,
X/q_0/0Y0,XX/q_1/Y0,XXY/q_1/0,XXY0/q_1/%B,
%% FIN RECOPIE
}
}

\newcommand\DIAGRAMMEESPACETEMPSdeuxvariante{
\def\memolcase{\lcase}
\def\memoMAXDET{20}
\def\MAXDET{15}
\def\lcase{6 ex}
\DIAGRAMMEESPACETEMPS
{
%% RECOPIE progdeux
/q_00/010, X/q_10/10, X0/q_11/0, X/q_20/Y0,/q_2X/0Y0,
X/q_00/Y0,XX/q_1Y/0,XXY/q_10/,XXY0/q_1B/%B,
%% FIN RECOPIE
}
\def\lcase{\memolcase}
\def\MAXDET{\memoMAXDET}
}


\newcommand\ANIMATIONun{
\ANIMATIONMT
%% RECOPIE progun
{/q_0/0011, % initial
X/q_1/011,X0/q_1/11,X/q_2/0Y1,/q_2/X0Y1,
X/q_0/0Y1,XX/q_1/Y1,XXY/q_1/1, XX/q_2/YY, X/q_2/XYY,
XX/q_0/YY,XXY/q_3/Y,XXYY/q_3/B,XXYYB/q_4/B}
%% FIN RECOPIE
}



\newcommand\DESSINALGO[3]{
% argument 1= ou
% argument 2= texte 
% argument 3= nom du noeud ex m
\draw node[draw,#1,minimum size=1cm] (#3) {#2};

% pour fitting box
\draw (#3.west) +(-0.2cm,0) node (#3c2) {};
\draw (#3.north) +(0,+0.2cm) node (#3c3) {};
\draw (#3.south) +(0,-0.2cm) node (#3c4) {};
}

\newcommand\DESSINALGOa[3]{
\DESSINALGO{#1}{#2}{#3}
\draw (#3.east) +(-0.1cm,0.2cm) node(#3a) {};
\draw[->] (#3a) -- ++(0.5cm,0) node (#3aa) {};
\draw (#3aa) node[anchor=west] (#3aaa) {\LANGUE{Accept}{Accepte}};
\draw (#3aaa.east) node(#3aaaa) {};
}

\newcommand\DESSINALGOe[4]{
\DESSINALGO{#1}{#2}{#3}
\draw (#3.east) +(-0.1cm,0.2cm) node(#3a) {};
\draw[->] (#3a) -- ++(0.5cm,0) node (#3aa) {};
\draw (#3aa) node[anchor=west] (#3aaa) {#4};
\draw (#3aaa.east) node(#3aaaa) {};
}


\newcommand\DESSINALGOar[3]{
\DESSINALGOa{#1}{#2}{#3}
\draw (#3.east) +(-0.1cm,-0.2cm) node(#3r) {};
\draw[->] (#3r) -- ++(0.5cm,0) node (#3rr) {};
\draw (#3rr) node[anchor=west] (#3rrr) {\LANGUE{Reject}{Refuse}};
\draw (#3rrr.east) node(#3rrrr) {};
}


\newcommand\WINDOWarg[7] {
%#1: position
%#2..7: contenu du tableau
\draw (#1) node[case]  {#2} \dcase;
\draw (#1) ++ (\lcase,0) node[case] {#3} \dcase;
\draw (#1) ++ (2*\lcase,0) node[case] {#4} \dcase;
\draw (#1) ++ (0,-\lcase) node[case] {#5} \dcase;
\draw (#1) ++ (\lcase,-\lcase) node[case] {#6} \dcase;
\draw (#1) ++ (2*\lcase,-\lcase) node[case] {#7} \dcase;
}

\newcommand\WINDOW[3]{
%#1: positioin
%#2: contenu
%#3: étiquette
\WINDOWarg{#1}
{\StrChar{#2}{1}}{\StrChar{#2}{2}}{\StrChar{#2}{3}}
{\StrChar{#2}{4}}{\StrChar{#2}{5}}{\StrChar{#2}{6}}
\draw (#1) +(-\lcase*1.2,-\lcase*0.5) node {#3};
}





%%%%%%%%%%%%%%%%%%%%%%%%%%%%%%%%%%%%%%%%%%%%%%%%%%%%%%%%%%%%%%%%%%%%%%%%%%%%
%                 LISTINGS
%%%%%%%%%%%%%%%%%%%%%%%%%%%%%%%%%%%%%%%%%%%%%%%%%%%%%%%%%%%%%%%%%%%%%%%%%%%%


% LISTINGS
\RequirePackage{listings}
 \lstset{language=Java,
 mathescape=true,
 showspaces=false,
 showstringspaces=false,
 frame=single,
morekeywords={recursive,then,div,function,integer,read,out,else,par,endpar,seq,endseq,and,not,read,write,repeat,until,undef,let,to,endfor,endif,endwhile,guess,in,parallel},
% basicstyle=\ttfamily\small,
% basewidth={1.1ex},
% keywordstyle=\bf\color{blue},
% %stringstyle=\bf\ttfamily\color{green},
% commentstyle=\bfseries\color{magenta},
 literate={:=}{{$\gets$ }}1 {<=}{{$\leq$}}2  {>=}{{$\geq$}}2
 {!=}{{$\neq$}}1 
 }
% \lstset{escapechar=\%}
\let\lst\lstinline
\def\beginlisting{\begin{lstlisting}}
\def\endlisting{\end{lstlisting}}  




\newenvironment{algo}[1]{
  \lstset{escapechar=\@}
  \def\input{$#1$}
  \def\out{out}
  \def\BEGIN{\textbf{read} \input }
 \def\BEGINN{\textbf{repeat} }
  \def\END{\textbf{until out!=undef}}
  \def\ENDD{\textbf{write \out}}
  \def\ENDT{\textbf{until $\ctl$ = $q_a$ or $\ctl$ = $q_r$}}
  \def\ENDR{\textbf{until $\ctl$ = $q_a$}}
  \def\ENDP{\textbf{until $\ctl$ = $q_a$}}
}{}
\newenvironment{algon}[1]{
  \lstset{numbers=left}
  \begin{algo}{#1}
}{\end{algo}}
\newcommand\FSM[3]{FSM(#1,\lst{#2},#3)}
\newcommand\FSMi[3]{FSM(#1,{#2},#3)}




%%%%%%%%%%%%%%%%%%%%%%%%%%%%%%%%%%%%%%%%%%%%%%%%%
%
%
%                      TRADUCTION EN COURS
%
%%%%%%%%%%%%%%%%%%%%%%%%%%%%%%%%%%%%%%%%%%%%%%%%%

\newcommand\ENTREPOT[1]{\NDLR{ENTREPOT ICI: #1}}
\newcommand\scite[2][]{\cite[#1]{#2}}
\newcommand\nospellcheck[1]{{#1}}
\newcommand\nd{nd}
\newcommand\NL{
 
}

\newenvironment{FRANCAIS}{
\small
\color{orange}
}
{\normalsize
\color{black}
}


\newenvironment{SEMIFRANCAIS}{
\small
\color{violet}
}
{\normalsize
\color{black}
}

\newcommand\QFRANCAIS[1]{
\begin{FRANCAIS}
#1
\end{FRANCAIS}
}


%%%%%%%%%%%%%%%%%%%%%%%%%%%%%%%%%%%%%%%%%%%%%%%%%
%
%
%                      Wikipedia
%
%%%%%%%%%%%%%%%%%%%%%%%%%%%%%%%%%%%%%%%%%%%%%%%%%

\providecommand\tightlist{}
\providecommand\Complex{\C}
\newenvironment{myalgie}{\begin{array}{lll}}{\end{array}}

\newcommand\nl{\ \\ }




%%%%%%%%%%%%%%%%%%%%%%%%%%%%%%%%%%%%%%%%%%%%%%%%%
%
%
%                      Pour citations précises
%
%%%%%%%%%%%%%%%%%%%%%%%%%%%%%%%%%%%%%%%%%%%%%%%%%



\newcommand\citeprecis[2]{{\cite[#1]{#2}}}
\newcommand\citeprecisinv[2]{\citeprecis{#2}{#1}}

\usepackage{csquotes}



%%%%%%%%%%%%%%%%%%%%%%%%%%%%%%%%%%%%%%%%%%%%%%%%%
%
%
%                      Travail sur couleurs
%
%   toutes les réponses pointent dans le même sens, et le schéma correspond exactement 
%. au profil perceptif d’un protanope** (déficit des cônes L / longueurs d’onde longues)
%  MAIS SEVERE
%
%%%%%%%%%%%%%%%%%%%%%%%%%%%%%%%%%%%%%%%%%%%%%%%%%


% ================================
% GROUPE 1 — Verts sombres
% ================================
\definecolor{deepgreen}{RGB}{0,160,40}
\definecolor{forestgreen}{RGB}{0,150,70}
\definecolor{mossgreen}{RGB}{0,140,90}
\definecolor{greenwave}{RGB}{0,150,90}
\definecolor{forestmist}{RGB}{0,160,90}
\definecolor{darkteal}{RGB}{0,130,70}

% ================================
% GROUPE 2 — Teal / verts bleutés
% ================================
\definecolor{seagreen}{RGB}{0,130,100}
\definecolor{freshteal}{RGB}{0,140,130}
\definecolor{greenteal}{RGB}{0,140,120}
\definecolor{blueteal}{RGB}{0,150,120}      % ancien "teal" (doublon corrigé)
\definecolor{lagoon}{RGB}{0,180,120}
\definecolor{greenmist}{RGB}{0,170,110}

% ================================
% GROUPE 3 — Turquoises / Aqua
% ================================
\definecolor{springgreen}{RGB}{0,190,100}
\definecolor{mintgreen}{RGB}{10,200,150}
\definecolor{oceanfoam}{RGB}{0,200,150}
\definecolor{aquagreen}{RGB}{30,210,170}
\definecolor{coastalaqua}{RGB}{40,210,180}
\definecolor{softaqua}{RGB}{40,210,190}

% ================================
% GROUPE 4 — Cyan / Aqua clairs
% ================================
\definecolor{emeraldteal}{RGB}{0,180,140}
\definecolor{deepcyan}{RGB}{0,180,160}
\definecolor{lightcyan}{RGB}{0,200,160}
\definecolor{aquamarine}{RGB}{0,200,160}
\definecolor{lightaqua}{RGB}{60,230,200}
\definecolor{mistblue}{RGB}{80,240,210}

% ================================
% GROUPE 5 — Bleus clairs
% ================================
\definecolor{paleaqua}{RGB}{80,200,190}
\definecolor{glacier}{RGB}{90,150,200}
\definecolor{skyblueA}{RGB}{60,150,220}      % ancien "skyblue"
\definecolor{lightsky}{RGB}{90,170,230}
\definecolor{iceblue}{RGB}{120,190,240}
\definecolor{frost}{RGB}{150,210,250}

% ================================
% GROUPE 6 — Bleus moyens
% ================================
\definecolor{steelblue}{RGB}{30,100,160}
\definecolor{cadetblue}{RGB}{70,120,180}
\definecolor{coolblue}{RGB}{110,170,220}
\definecolor{morningblue}{RGB}{140,200,240}
\definecolor{blueiceA}{RGB}{40,150,210}       % ancien "blueice" (doublon corrigé)
\definecolor{blueair}{RGB}{60,170,225}

% ================================
% GROUPE 7 — Azurs / Bleu saturé
% ================================
\definecolor{tealblue}{RGB}{0,110,130}
\definecolor{deepbluecyan}{RGB}{0,110,170}
\definecolor{azur}{RGB}{0,70,180}
\definecolor{brightazure}{RGB}{20,130,200}
\definecolor{lightazure}{RGB}{80,180,235}
\definecolor{skylight}{RGB}{100,200,245}

% ================================
% GROUPE 8 — Bleus sombres / violets
% ================================
\definecolor{nightblue}{RGB}{40,0,110}
\definecolor{slateblue}{RGB}{20,40,120}
\definecolor{violetA}{RGB}{60,50,160}         % ancien "violet" (distinct de violetB ci-dessous)
\definecolor{blueviolet}{RGB}{80,60,180}
\definecolor{purpleA}{RGB}{100,70,200}        % ancien "purple"
\definecolor{lightpurple}{RGB}{130,90,210}

% ================================
% GROUPE 9 — Violets saturés / indigos
% ================================
\definecolor{lilac}{RGB}{160,110,220}
\definecolor{indigoA}{RGB}{50,0,130}          % ancien "indigo"
\definecolor{deepindigo}{RGB}{70,10,150}
\definecolor{royalpurple}{RGB}{90,20,170}

% ================================
% GROUPE 10 — Fermeture cercle perceptif
% ================================
\definecolor{coolmint}{RGB}{0,130,70}
\definecolor{deepsea}{RGB}{0,130,70}

% ================================
% PALETTE DEUTÉRANOPE — noms uniques
% ================================
\definecolor{crimson}{RGB}{200,30,60}
\definecolor{brickred}{RGB}{180,40,40}
\definecolor{coral}{RGB}{220,90,70}
\definecolor{salmon}{RGB}{240,130,120}

\definecolor{magenta}{RGB}{200,40,150}
\definecolor{fuchsia}{RGB}{220,60,170}
\definecolor{violetB}{RGB}{140,70,190}        % ancien "violet"
\definecolor{indigoB}{RGB}{90,50,180}         % ancien "indigo"

\definecolor{royalblueA}{RGB}{50,80,200}      % ancien "royalblue"
\definecolor{azure}{RGB}{30,120,230}
\definecolor{skyblueB}{RGB}{80,160,230}       % second "skyblue"
\definecolor{tealB}{RGB}{0,150,170}           % second "teal"

\definecolor{marine}{RGB}{0,80,140}




%% Alternative

% === Générateur automatique fill/text pour protanopes ===
% Usage : \ProtaColorPair{mistblue}
% Suppose que mistblue est déjà défini par \definecolor{mistblue}{RGB}{...}

\newcommand{\ProtaColorPair}[1]{%
  % fill = couleur originale
  \colorlet{#1fill}{#1}%
  % text = version assombrie idéalement lisible
  \colorlet{#1text}{#1!80!black}%
}

%Utilisation
%\ProtaColorPair{mistblue}

%\textcolor{mistbluefill}{Couleur claire} \quad
%\textcolor{mistbluetext}{Couleur pour texte}

	\foreach \i/\name in {
	  1/deepgreen,
	  2/greenwave,
	  3/lagoon,
	  4/oceanfoam,
	  5/softaqua,
	  6/mistblue,
 	 7/lightsky,
	  8/blueiceA,   % <-- corrigé
	  9/deepcyan,
	  10/azur,
	  11/blueviolet,
	  12/royalpurple}
	  {\ProtaColorPair{\name}}
	  


\newcommand\ROUEPROTANOPEDOUZE{
	% Roue chromatique protanope — 12 couleurs
	\begin{tikzpicture}[scale=3]

	\foreach \i/\name/\nametext in {
	  1/deepgreen/deepgreen,
	  2/greenwave/greenwave,
	  3/lagoon/lagoontext,
	  4/oceanfoam/oceanfoam,
	  5/softaqua/softaquatext,
	  6/mistblue/mistbluetext,
 	 7/lightsky/lightskytext,
	  8/blueiceA/blueiceA,   % <-- corrigé
	  9/deepcyan/blueiceA,
	  10/azur/azur,
	  11/blueviolet/blueviolet,
	  12/royalpurple/royalpurple}
	{
	  % secteur coloré
 	 \filldraw[fill=\name, draw=black, line width=0.2pt]
	    ({90-30*\i}:1)
	    -- ({90-30*(\i-1)}:1)
	    -- (0,0) -- cycle;

 	 % label dans la bonne couleur
	  \node at ({90-30*(\i-0.5)}:1.25)
	    {\textcolor{\nametext}{\small \nametext}};
	}

	\end{tikzpicture}
}


\newcommand\ROUEPROTANOPEVINGT{
\begin{tikzpicture}[scale=3]

\foreach \i/\name/\nametext in {
  1/deepgreen/deepgreen,
  2/greenwave/greenwave,
  3/lagoon/lagoon,
  4/oceanfoam/oceanfoam,
  5/softaqua/softaquatext,
  6/mistblue/mistbluetext,
  7/lightsky/lightsky,
  8/blueiceA/blueiceA,      % corrigé
  9/deepcyan/deepcyan,
  10/azur/azur,
  11/brightazure/brightazure,
  12/skyblueA/skyblueA,     % corrigé
  13/coolblue/coolblue,
  14/glacier/glacier,
  15/marine/marine,
  16/violetA/violetA,      % corrigé
  17/blueviolet/blueviolet,
  18/purpleA/purpleA,      % corrigé
  19/royalpurple/royalpurple,
  20/indigoA/indigoA}      % corrigé
{
  % Secteur coloré
  \filldraw[fill=\name, draw=black, line width=0.2pt]
    ({90-18*\i}:1)
    -- ({90-18*(\i-1)}:1)
    -- (0,0) -- cycle;

  % Étiquette
  \node at ({90-18*(\i-0.5)}:1.25)
    {\textcolor{\nametext}{\small \nametext}};
}

\end{tikzpicture}
}

\newcommand\ROUEDEUTERANOPEDOUZE{
% Roue chromatique deutéranope — 12 couleurs
\begin{tikzpicture}[scale=3]

\foreach \i/\name in {
  1/crimson,
  2/brickred,
  3/coral,
  4/salmon,
  5/magenta,
  6/fuchsia,
  7/violetB,       % corrigé
  8/indigoB,       % corrigé
  9/royalblueA,    % corrigé
  10/azure,
  11/skyblueB,     % corrigé
  12/tealB}        % corrigé
{
  % Secteur coloré
  \filldraw[fill=\name, draw=black, line width=0.2pt]
    ({90-30*\i}:1)
    -- ({90-30*(\i-1)}:1)
    -- (0,0) -- cycle;

  % Label dans la couleur
  \node at ({90-30*(\i-0.5)}:1.25)
    {\textcolor{\name}{\small \name}};
}

\end{tikzpicture}
}



% === Macro intelligente pour foncer une couleur de façon lisible pour protanope ===
% Usage : \protadarken{mistblue}{40}  -> crée mistblue@dark
% Le 40 signifie : 40% de mélange avec du noir

\newcommand{\DarkenForText}[2]{%
  % #1 = nom de la couleur
  % #2 = pourcentage (40 conseillé si tu ne sais pas)
  %
  % On crée une nouvelle couleur #1@dark
  \colorlet{#1text}{black!#2!#1}%
}





%%%%

% Groupe 1
\DarkenForText{deepgreen}{20}
\DarkenForText{forestgreen}{20}
\DarkenForText{mossgreen}{20}
\DarkenForText{greenwave}{20}
\DarkenForText{forestmist}{20}
\DarkenForText{darkteal}{20}

% Groupe 2
\DarkenForText{seagreen}{20}
\DarkenForText{freshteal}{20}
\DarkenForText{greenteal}{20}
\DarkenForText{blueteal}{20}
\DarkenForText{lagoon}{20}
\DarkenForText{greenmist}{20}

% Groupe 3
\DarkenForText{springgreen}{20}
\DarkenForText{mintgreen}{20}
\DarkenForText{oceanfoam}{20}
\DarkenForText{aquagreen}{20}
\DarkenForText{coastalaqua}{20}
\DarkenForText{softaqua}{20}

% Groupe 4
\DarkenForText{emeraldteal}{20}
\DarkenForText{deepcyan}{20}
\DarkenForText{lightcyan}{20}
\DarkenForText{aquamarine}{20}
\DarkenForText{lightaqua}{20}
\DarkenForText{mistblue}{30}

% Groupe 5
\DarkenForText{paleaqua}{20}
\DarkenForText{glacier}{20}
\DarkenForText{skyblueA}{20}
\DarkenForText{lightsky}{20}
\DarkenForText{iceblue}{20}
\DarkenForText{frost}{20}

% Groupe 6
\DarkenForText{steelblue}{20}
\DarkenForText{cadetblue}{20}
\DarkenForText{coolblue}{20}
\DarkenForText{morningblue}{20}
\DarkenForText{blueiceA}{20}
\DarkenForText{blueair}{20}

% Groupe 7
\DarkenForText{tealblue}{20}
\DarkenForText{deepbluecyan}{20}
\DarkenForText{azur}{20}
\DarkenForText{brightazure}{20}
\DarkenForText{lightazure}{20}
\DarkenForText{skylight}{20}

% Groupe 8
\DarkenForText{nightblue}{20}
\DarkenForText{slateblue}{20}
\DarkenForText{violetA}{20}
\DarkenForText{blueviolet}{20}
\DarkenForText{purpleA}{20}
\DarkenForText{lightpurple}{20}

% Groupe 9
\DarkenForText{lilac}{20}
\DarkenForText{indigoA}{20}
\DarkenForText{deepindigo}{20}
\DarkenForText{royalpurple}{20}

% Groupe 10
\DarkenForText{coolmint}{20}
\DarkenForText{deepsea}{20}

% Palette deutéranope
\DarkenForText{crimson}{20}
\DarkenForText{brickred}{20}
\DarkenForText{coral}{20}
\DarkenForText{salmon}{20}

\DarkenForText{magenta}{20}
\DarkenForText{fuchsia}{20}
\DarkenForText{violetB}{20}
\DarkenForText{indigoB}{20}

\DarkenForText{royalblueA}{20}
\DarkenForText{azure}{20}
\DarkenForText{skyblueB}{20}
\DarkenForText{tealB}{20}

% Couleur supplémentaire
\DarkenForText{marine}{20}

%%%%%%%%%%%%%%%%%%%%%
%
%.   apx proof
%
%%%%%%%%%%%%%%%%%%%%%



% do
\newcommand\PREUVESENAPPENDIXCOMPATIBLE{
	\newenvironment{appendixproof}{{\bf \LANGUE{Proof}{Démonstration}}:}{\hfill $\Box$}
	\newtheorem{theoremrep}[theorem]{Theorem}
	\newtheorem{lemmarep}[theorem]{Lemma}
	\newtheorem{propositionrep}[theorem]{Proposition}
	\newtheorem{corollaryrep}[theorem]{Corollary}
	\newtheorem{summaryrep}[theorem]{Summary}	
}

\providecommand\apxparameter{}

%% proof inline
%\providecommand\apxparameter{appendix=inline}
%% proof in appendix
%\providecommand\apxparameter{}


\apxproofmode{
\usepackage[\apxparameter]{apxproof}
\theoremstyle{plain}

\newtheoremrep{theorem}{Theorem}
%\newtheoremrep{theoremconjecture}[theorem]{Theorem-To-Be-Proved=Conjecture}
%\newtheoremrep{lemmaconjecture}[theorem]{Conjecture-To-Be-Proved=Conjecture}
\newtheoremrep{claim}[theorem]{Theorem}
\newtheoremrep{proposition}[theorem]{Proposition}
\newtheoremrep{lemma}[theorem]{Lemma}
\newtheoremrep{corollary}[theorem]{Corollary}
%%bug mars 26: \newtheoremrep{summary}[theorem]{Summary}
% Du coup:
\newtheoremrep{resume}[theorem]{Summary}
%bug mars 26: \newtheoremrep{openproblem}[theorem]{Open Problem}
%\newtheoremrep{claim}[theorem]{Claim}
}


%\apxproofmodesubsection{
%%apx proof, pour que compteur correct
%\makeatletter
%\AtBeginDocument{%
%\let\axp@section\subsection
%  \let\axp@oldsection\subsection
%}
%\makeatother
%\renewcommand\appendixprelim{\section{Missing proofs}}
%}


\makeatletter
\apxproofmodesubsection{
\let\axp@section\subsection
%\let\axp@oldsection\subsection
\renewcommand\appendixprelim{\section{Proofs from main documents}}
}
\makeatother


\makeatletter
\newcommand\appendixapxproofsubsection{
\appendix
\let\apx@orig@appendix\appendix
\renewcommand{\appendix}{}
}
\makeatother


%%% Pour checker/vérifier que bon


\newcommand\docheck[1]{\textcolor{check}{#1}}

\usepackage{xcolor}
%\definecolor{chose}{RGB}{0,110,50}
%\definecolor{forestgreen}{RGB}{0,150,70}
%\definecolor{ctruc}{RGB}{164,164,164}
%\definecolor{oceanfoam}{RGB}{0,200,150}
\definecolor{check}{RGB}{160,160,0}




 via \providecommand,
%% ou après via \renewcommand) :
\providecommand\BIBFILES{@@reference-biblio}
\providecommand\FIGCOMMONS{/Users/bournez/lib/LaTeX/Perso/FIG-COMMONS}

\newcommand\affiliationofficielleolivier{LIX, CNRS, École polytechnique, Institut Polytechnique de Paris, Palaiseau, France}
\newcommand\affiliationofficielleolivierslides{LIX, CNRS, École polytechnique, \\ Institut Polytechnique de Paris, Palaiseau, France}


%%%%%%%%%%%%%%%%%%%%%%%%
%
%              begin. Switch modes compatibilités
%
%%%%%%%%%%%%%%%%%%%%%%%

%% Apxproof mode
\providecommand\apxproofmode[1]{}
	% sinon: \newcommand\apxproof[1]{#1}
	
\providecommand\apxproofmodesubsection[1]{}
	% sinon: \newcommand\apxproofmodesubsection[1]{#1}
% utiliser: \appendixapxproofsubsection au lieu de \appendix dans ce cas. 
		

%% Lipics mode
\providecommand\lipicsmode{}
	% sinon: \newcommand\lipicsmode[1]{#1}
		% en vrai sert qu'à définir \nolipics
		
%%%%%%%%%%%%%%%%%%%%%%%%
%
%              end. Switch modes compatibilités
%
%%%%%%%%%%%%%%%%%%%%%%%


%% begin gestion des switchs

% en fait \nopilcs est la seule commande utilisée

\newcommand\nolipics[1]{#1}
\lipicsmode{\renewcommand\nolipics[1]{}}

%% Décorations (fourier + cadres colorés mdframed autour de definition,
%% theorem, exercice, example, remark, ...) : actives par défaut hors slides.
%% Pour les désactiver dans un document : \newcommand\NODECORATION[1]{}
%% avant 
\providecommand\nolncs[1]{#1}	

%% Chemins par défaut (surchargables avant \input{macros} via \providecommand,
%% ou après via \renewcommand) :
\providecommand\BIBFILES{@@reference-biblio}
\providecommand\FIGCOMMONS{/Users/bournez/lib/LaTeX/Perso/FIG-COMMONS}

\newcommand\affiliationofficielleolivier{LIX, CNRS, École polytechnique, Institut Polytechnique de Paris, Palaiseau, France}
\newcommand\affiliationofficielleolivierslides{LIX, CNRS, École polytechnique, \\ Institut Polytechnique de Paris, Palaiseau, France}


%%%%%%%%%%%%%%%%%%%%%%%%
%
%              begin. Switch modes compatibilités
%
%%%%%%%%%%%%%%%%%%%%%%%

%% Apxproof mode
\providecommand\apxproofmode[1]{}
	% sinon: \newcommand\apxproof[1]{#1}
	
\providecommand\apxproofmodesubsection[1]{}
	% sinon: \newcommand\apxproofmodesubsection[1]{#1}
% utiliser: \appendixapxproofsubsection au lieu de \appendix dans ce cas. 
		

%% Lipics mode
\providecommand\lipicsmode{}
	% sinon: \newcommand\lipicsmode[1]{#1}
		% en vrai sert qu'à définir \nolipics
		
%%%%%%%%%%%%%%%%%%%%%%%%
%
%              end. Switch modes compatibilités
%
%%%%%%%%%%%%%%%%%%%%%%%


%% begin gestion des switchs

% en fait \nopilcs est la seule commande utilisée

\newcommand\nolipics[1]{#1}
\lipicsmode{\renewcommand\nolipics[1]{}}

%% Décorations (fourier + cadres colorés mdframed autour de definition,
%% theorem, exercice, example, remark, ...) : actives par défaut hors slides.
%% Pour les désactiver dans un document : \newcommand\NODECORATION[1]{}
%% avant \input{macros}. Retirées automatiquement en mode lipics.
\providecommand\NODECORATION[1]{\ifnotSLIDE{#1}}
\lipicsmode{\renewcommand\NODECORATION[1]{}}

\providecommand\PASSIAPXPROOF[1]{#1}
\apxproofmode{\renewcommand\PASSIAPXPROOF[1]{}}



%% end gestion des switchs

\def\MACROSPOINTEXOUVERT{}


\providecommand\ifnotTDEXAM[1]{#1}
\providecommand\ifnotSLIDE[1]{#1}


  
% pour preuves en appendix
\providecommand\PREUVESENAPPENDIX[1]{}



%%%% STYLES: APRES VOLS


\NODECORATION{
\usepackage{fourier}
}

%% Repli si fourier (et donc fourier-orns, qui définit \danger) n'est pas
%% chargé ; doit rester APRÈS le bloc ci-dessus, sinon fourier-orns échoue
%% avec « \danger already defined ».
\providecommand\danger{}

%%%% Box around environment
\usepackage[framemethod=tikz]{mdframed}

\newcommand\ENVCOLOR[1]{#1}
\newcommand\BEGINTEMPENVCOLOR[1]{\let\envcolorwas\ENVCOLOR\renewcommand\ENVCOLOR[1]{#1}}
\newcommand\ENDTEMPENVCOLOR{\let\ENVCOLOR\envcolorwas}


\NODECORATION{
\surroundwithmdframed[hidealllines=true,backgroundcolor=\ENVCOLOR{blue!10},roundcorner=8pt]{exercice}
\surroundwithmdframed[hidealllines=true,backgroundcolor=\ENVCOLOR{blue!10},roundcorner=8pt]{exercise}
\surroundwithmdframed[hidealllines=true,backgroundcolor=\ENVCOLOR{blue!10},roundcorner=8pt]{exercicee}
\surroundwithmdframed[hidealllines=true,backgroundcolor=\ENVCOLOR{blue!10},roundcorner=8pt]{exerciced}
\surroundwithmdframed[hidealllines=true,backgroundcolor=\ENVCOLOR{blue!10},roundcorner=8pt]{exercicec}
\surroundwithmdframed[hidealllines=true,backgroundcolor=\ENVCOLOR{blue!10},roundcorner=8pt]{exercicedc}
\surroundwithmdframed[hidealllines=true,backgroundcolor=\ENVCOLOR{orange!20},roundcorner=8pt]{example}
\surroundwithmdframed[hidealllines=true,backgroundcolor=\ENVCOLOR{orange!20},roundcorner=8pt]{exemple}
\surroundwithmdframed[hidealllines=true,backgroundcolor=\ENVCOLOR{orange!10},roundcorner=8pt]{remark}
\ifnotSLIDE{\surroundwithmdframed[backgroundcolor=\ENVCOLOR{green!15},roundcorner=8pt]{definition}}
\surroundwithmdframed[hidealllines=true,backgroundcolor=\ENVCOLOR{gray!35}]{proposition}
\ifnotSLIDE{\surroundwithmdframed[backgroundcolor=\ENVCOLOR{gray!35},roundcorner=8pt]{theorem}}
\surroundwithmdframed[hidealllines=true,backgroundcolor=\ENVCOLOR{gray!35}]{lemma}
\surroundwithmdframed[hidealllines=true,backgroundcolor=\ENVCOLOR{gray!35}]{corollary}


\surroundwithmdframed[backgroundcolor=red!35,roundcorner=8pt]{theoremaprouver}
\surroundwithmdframed[backgroundcolor=red!35,roundcorner=8pt]{conjectureaprouver}
}

%%% tcolorbox

\usepackage{tcolorbox}
\tcbuselibrary{skins}


  %%. Commande magique
%% \tcolorboxenvironment{⟨name⟩}{⟨options⟩}: attach options de tcolorbox à environnement
%% alternative: je cree des boites moi meme: voici des exemples


\newtcolorbox{maboiteconfiance}[2][]{colback=red!5!white, colframe=red!75!black,fonttitle=\bfseries, colbacktitle=red!85!black,enhanced,
attach boxed title to top center={yshift=-2mm},
  title={#2},#1}
  
\newtcolorbox{maboiteresume}[2][]{colback=red!5!white, colframe=red!75!black,fonttitle=\bfseries, colbacktitle=red!85!black,enhanced,
attach boxed title to top center={yshift=-2mm},
  title={#2},#1}
  \newtcolorbox{maboiteresumeperso}[2][]{colback=red!5!white, colframe=red!75!black,fonttitle=\bfseries, colbacktitle=red!85!black,enhanced,
attach boxed title to top center={yshift=-2mm},
  title={#2},#1}
  \newtcolorbox{maboitecommentaire}[2][]{colback=red!5!white, colframe=red!75!black,fonttitle=\bfseries, colbacktitle=red!85!black,enhanced,
attach boxed title to top center={yshift=-2mm},
  title={#2},#1}
\newtcolorbox{maboitetruc}[2][]{colback=red!5!white, colframe=red!75!black,fonttitle=\bfseries, colbacktitle=red!85!black,enhanced,
attach boxed title to top left={yshift=-2mm},
  title={#2},#1}

%% Boîte des mises en évidence CARE. Couleurs pensées pour la
%% protanopie : l'axe bleu-jaune reste discriminant (jaune = \IMPORTANT
%% d'Olivier, bleu = CARE), et la distinction ne repose pas que sur la
%% couleur (cadre + bandeau étiqueté, lisibles même en niveaux de gris).
\newtcolorbox{maboiteimportantcare}[2][]{colback=cyan!10!white, colframe=blue!60!black,fonttitle=\bfseries, colbacktitle=blue!70!black,enhanced,
attach boxed title to top center={yshift=-2mm},
  title={#2},#1}
  


%%%% FIN STYLES



%%%%%%%%%%%%%%%%%%%%%%%%%%%%%%%%%%%%%%%%%%%%%%%%%
%%%%%%%%%%%%%%%%%%%%%%%%%%%%%%%%%%%%%%%%%%%%%%%%%
%
%
%                      Version ENSEIGNEMENT 2020
%
%
%%%%%%%%%%%%%%%%%%%%%%%%%%%%%%%%%%%%%%%%%%%%%%%%%
%%%%%%%%%%%%%%%%%%%%%%%%%%%%%%%%%%%%%%%%%%%%%%%%%

%english, french
\ifdefined\CESTDUFRANCAIS
	\providecommand\LANGUE[2]{#2}
	\usepackage[french]{babel}
\else
	\providecommand\LANGUE[2]{#1}
\fi
\providecommand\LANGUEinv[2]{\LANGUE{#2}{#1}}

%%%%%%%%%%%%%%%%%%%%%%%%%%%%%%%%%%%%%%%%%%%%%%%%%
%
%
%                      Packages
%
%
%%%%%%%%%%%%%%%%%%%%%%%%%%%%%%%%%%%%%%%%%%%%%%%%%


%%% SURLIGNEUR
\providecommand\ifnotmodeDIFFUSIONFINALE[1]{
%% 2021: Comprend pas mais assure
\ifdefined\SAFEMODE
\else
#1
\fi
%#1
}
\providecommand\ifmodeDIFFUSIONFINALE[1]{
\ifdefined\SAFEMODE
#1
\else
\fi
}
\newcommand\SURLIGNE[1]{
\ifnotmodeDIFFUSIONFINALE{
\BEGINTEMPENVCOLOR{blue!30}
{                        \colorbox{blue!30}{

\noindent \begin{minipage}{\dimexpr\textwidth-2\fboxsep\relax}
#1
\end{minipage}
}
}
\ENDTEMPENVCOLOR}
\ifmodeDIFFUSIONFINALE{
\noindent \begin{minipage}{\textwidth}
#1
\end{minipage}
}
}


\newenvironment{confiance}{\begin{maboiteresume}[colback=brown!50]{Résumé ici}}{\end{maboiteresume}}
\newcommand\CONFIANCE[2][Confiance ici]{
\begin{maboiteconfiance}[colback=yellow!50]{#1}
\textbf{ATTENTION : #2}
\end{maboiteconfiance}
}

\newenvironment{resume}{\begin{maboiteresume}[colback=brown!50]{Résumé ici}}{\end{maboiteresume}}
\newcommand\RESUME[2][Résumé ici]{
\begin{maboiteresume}[colback=brown!50]{#1}
#2
\end{maboiteresume}
}

\newenvironment{resumeperso}{\begin{maboiteresumeperso}[colback=brown!50]{Résumé perso ici}}{\end{maboiteresumeperso}}
\newcommand\RESUMEPERSO[2][Résumé perso  ici]{
\begin{maboiteresumeperso}[colback=brown!50]{#1}
#2
\end{maboiteresumeperso}
%
%
%\todo[inline, color=brown]{résumé: #1}
%\colorbox{brown}{
%\noindent \begin{minipage}{\textwidth}
%{
%\bf 
%#2
%}
%\end{minipage}
%}
}

\newcommand\DANSLAMARGE[1]{
\ifnotmodeDIFFUSIONFINALE{
\marginpar{\textcolor{brown}{#1}}}
}

\newenvironment{commentaireperso}{\begin{maboitecommentaire}[colback=brown!30]{Commentaire perso} Commentaire perso:}{\end{maboitecommentaire}}
\newcommand\COMMENTAIREPERSO[2][Commentaire perso]{
\begin{maboitecommentaire}[colback=brown!30]{#1}
#2
\end{maboitecommentaire}
}


\newcommand\SURSURLIGNE[2][truc surligné ici]{
%\todo[inline, color=orange]{surligné: #1}
\ifnotmodeDIFFUSIONFINALE{
\BEGINTEMPENVCOLOR{blue!60}
	                 \colorbox{blue!60}{
\noindent \begin{minipage}{\dimexpr\textwidth-2\fboxsep\relax}
{
\bf 
#2
}
\end{minipage}
}
\ENDTEMPENVCOLOR
}
\ifmodeDIFFUSIONFINALE{#2}
}

\newcommand\IMPORTANT[1]{
\ifnotmodeDIFFUSIONFINALE{
\BEGINTEMPENVCOLOR{yellow!100}
		        \colorbox{yellow!100}{

\noindent \begin{minipage}{\dimexpr\textwidth-2\fboxsep\relax}
{ 
#1
}
\end{minipage}
}
\ENDTEMPENVCOLOR
}
\ifmodeDIFFUSIONFINALE{#1}
}

%% Pendant de \IMPORTANT, réservé aux mises en évidence CARE ;
%% \IMPORTANT reste réservé aux annotations d'Olivier. Comme \IMPORTANT,
%% redevient du texte simple en mode DIFFUSIONFINALE.
\newcommand\IMPORTANTCARE[2][Important (CARE)]{
\ifnotmodeDIFFUSIONFINALE{
\begin{maboiteimportantcare}{#1}
#2
\end{maboiteimportantcare}
}
\ifmodeDIFFUSIONFINALE{#2}
}


\newcommand\QUESTION[1]{
\ifnotmodeDIFFUSIONFINALE{
\BEGINTEMPENVCOLOR{oceanfoam!100}
		        \colorbox{oceanfoam!100}{

\noindent \begin{minipage}{\dimexpr\textwidth-2\fboxsep\relax}
{ 
#1
}
\end{minipage}
}
\ENDTEMPENVCOLOR
}
\ifmodeDIFFUSIONFINALE{#1}
}




\newcommand\SOUSLIGNE[1]{
\colorbox{lightgray}{

\noindent \begin{minipage}{\dimexpr\textwidth-2\fboxsep\relax}
{ \color{gray}
#1
}
\end{minipage}
}
}



%% Gestion des cross-references cross-referencing
\usepackage{xr-hyper}
\usepackage{hyperref}

%%
\usepackage{versions}
\usepackage{amsmath,amssymb,mathrsfs,amsfonts}
%\usepackage{mathabx}
\usepackage{listings}
\usepackage{url}
\usepackage{versions}
\usepackage{color}
\usepackage[colorinlistoftodos]{todonotes}
%\usepackage{todonotes}
\usepackage{proof}
\usepackage{xstring}

%% Index
\usepackage{makeidx}
\makeindex
\providecommand\DOINDEXPRINT{}
\providecommand\SIGNALEMOTNOUV[1]{}
\ifdefined\INDEXPRINT
	\ifnotSAFEMODE{
%	\usepackage{showidx}
%	\let\oldindex\index
%	\renewcommand{\index}[1]{\oldindex{#1}\marginpar{LA: \tiny#1}}
	\renewcommand\DOINDEXPRINT{
	      	\let\oldindex\index
		\renewcommand{\index}[1]{\oldindex{##1}\marginpar{IDXLA: \tiny##1}}
		\renewcommand\SIGNALEMOTNOUV[1]{\marginpar{MNV: \small\textbf{##1}}}
	}
	}
\fi

% === gestion des accents


%soit iso latin1
%\usepackage[cyr]{aeguill}
%soit utf8
\usepackage[utf8]{inputenc}
\usepackage[T1]{fontenc}


%%%  Pour citation \Cref (cref, ref)	

\usepackage{cleveref}


%%%%%%%%%%%%%%%%%%%%%%%%%%%%%%%%%%%%%%%%%%%%%%%%%
%
%
%                      Gestion des accents spéciaux.
%
%
%%%%%%%%%%%%%%%%%%%%%%%%%%%%%%%%%%%%%%%%%%%%%%%%%

\input{macros-accents}
%% Tolérant tant que macros-markdown.tex n'est pas dans le dépôt :
\InputIfFileExists{macros-markdown}{}{\typeout{*** macros-markdown.tex introuvable: ignore ***}}


%\usepackage{PDFdraftcopy}


%-- Pour état des chapitres


%\newcommand\ETAT{\draftstring{BROUILLON}}
%
%\newcommand\brouillon{\renewcommand\ETAT{\draftstring{BROUILLON}}}
%
%\newcommand\final{\renewcommand\ETAT{\draftstring{\rien}}}
%
%\newcommand\travail{\renewcommand\ETAT{\draftstring{CHANTIER}}}
%
%\newcommand\bazar{\renewcommand\ETAT{\draftstring{BAZAR TOTAL}}}
%
%\final



%%%%%%%%%%%%%%%%%%%%%%%%%%%%%%%%%%%%%%%%%%%%%%%%%
%
%
%                      TRADUCTION EN COURS
%
%%%%%%%%%%%%%%%%%%%%%%%%%%%%%%%%%%%%%%%%%%%%%%%%%


% === volé Amaury



\usepackage{stmaryrd}
\usepackage{ifthen}
\usepackage{mathtools}
\usepackage{dsfont}
\PASSIAPXPROOF{\usepackage{thmtools}}
\usepackage{longtable}
\usepackage{mathdots}
\usepackage{booktabs}


%%%%%%%%%%%%%%%%%%%%%%%%%%%%%%%%%%%%%%%%%%%%%%%%%
%
%
%                      TIKZ STUFFS
%
%
%%%%%%%%%%%%%%%%%%%%%%%%%%%%%%%%%%%%%%%%%%%%%%%%%

%%% volé pour dessin des GPAC: PAS SUR QUE BESOIN de TOUT

\usepackage{tikz}
\usetikzlibrary{calc}
\usepgflibrary{arrows}
\usetikzlibrary{fit}
\usetikzlibrary{patterns}
\usetikzlibrary{decorations.pathreplacing}
\usetikzlibrary{shapes.multipart}
\usetikzlibrary{shapes.geometric}
\usetikzlibrary{shapes.misc}
\usetikzlibrary{positioning}
\usetikzlibrary{graphs}
\usetikzlibrary{topaths}
\usetikzlibrary{arrows.meta}
\usetikzlibrary{decorations.pathreplacing}
\usetikzlibrary{decorations.markings}
\usetikzlibrary{decorations.pathmorphing}
%
 \usetikzlibrary{calc}
 \usetikzlibrary{automata}
\usetikzlibrary{chains}
\usetikzlibrary{scopes}
 \usetikzlibrary{positioning}
 \usetikzlibrary{through,backgrounds}
\usepackage{circuitikz}
% % pour accolades




%%%%%%%%%%%%%%%%%%%%%%%%%%%%%%%%%%%%%%%%%%%%%%%%%
%
%
%                      MACROS locales
%
%
%%%%%%%%%%%%%%%%%%%%%%%%%%%%%%%%%%%%%%%%%%%%%%%%%


%====================
%                      Moi
%=====================


\newcommand\myaddress{ 
{\sc LIX, CNRS, Ecole Polytechnique, Institut Polytechnique de Paris, }\\ 91128 Palaiseau Cedex, FRANCE \\
{\small Olivier.Bournez@lix.polytechnique.fr}}



%====================
%                      un peu de trucs sales
%=====================

%=  pour style lipics et ses trucs

\providecommand\ifnotPOLYCHAPTERMODE[1]{#1}
\providecommand\ifPOLYCHAPTERMODE[1]{}

\ifnotPOLYCHAPTERMODE{
\ifdefined\thechapter
\providecommand\spnewtheorem[4]{\newtheorem{#1}{#2}[chapter]}
\providecommand\spnewtheoremi[5]{\newtheorem{#1}[#5]{#2}}
\else
\providecommand\spnewtheorem[4]{\newtheorem{#1}{#2}}
\providecommand\spnewtheoremi[5]{\newtheorem{#1}[#5]{#2}}
\fi
}
\ifPOLYCHAPTERMODE{
\providecommand\spnewtheorem[4]{\newtheorem{#1}{#2}}
\providecommand\spnewtheoremi[5]{\newtheorem{#1}[#5]{#2}}
}
\providecommand\nolipics[1]{#1}




%====================
%                      un peu de trucs franchement sales
%=====================


% on l'inclus pas ici.
%\input{macros-moins-propres}

%%%%%%%%%%%%%%%%%%%%%%%%%%%%%%%%%%%%%%%%%%%%%%%%%
%
%
%                      Page de Garde Cours X
%
%
%%%%%%%%%%%%%%%%%%%%%%%%%%%%%%%%%%%%%%%%%%%%%%%%%


\includeversion{metdate}

\newcommand\XPAGEDEGARDE[2]{
\thispagestyle{empty}
\title{{#1} \\[2cm] %\\[3cm] 
{\large #2}
}
\author{
 Olivier Bournez% \inst{1}
\\[0.5cm] 
bournez@lix.polytechnique.fr%{\myaddress}
}
%\date{
%\vspace{2cm}
%\includegraphics[height=2.5cm]{/Users/bournez/lib/LaTeX/Perso/FIG-COMMONS/logo_X_new}
%}
\begin{metdate}
\date{
\vspace{2cm}
Version \LANGUE{of}{du} \today \\[1cm] 
\includegraphics[height=3.5cm]{/Users/bournez/lib/LaTeX/Perso/FIG-COMMONS/logo_X_new}
}
\end{metdate}
\maketitle
}


%%%%%%%%%%%%%%%%%%%%%%%%%%%%%%%%%%%%%%%%%%%%%%%%%
%
%
%                      MACROS UTILES
%
%
%%%%%%%%%%%%%%%%%%%%%%%%%%%%%%%%%%%%%%%%%%%%%%%%%

%%%                      Macros
 
%-------------
%--- vectors  ---
%-------------

 
\newcommand\vectorl[1]{{\mathbf#1}}
\newcommand\tu[1]{\vectorl{#1}}

\newcommand\vp{\vectorl{p}}
\newcommand\va{\vectorl{a}}
\newcommand\vb{\vectorl{b}}
\newcommand\vc{\vectorl{c}}
\newcommand\vf{\vectorl{f}}
\newcommand\vn{\vectorl{n}}
\newcommand\vx{\vectorl{x}}
\newcommand\vX{\vectorl{X}}
\newcommand\vy{\vectorl{y}}
\newcommand\vz{\vectorl{z}}
\newcommand\ve{\vectorl{e}}
\newcommand\vv{\vectorl{v}}
\newcommand\vr{\vectorl{r}}
\newcommand\vu{\vectorl{u}}
\newcommand\vA{\vectorl{A}}
\newcommand\vB{\vectorl{B}}
\newcommand\vC{\vectorl{C}}
\newcommand\vD{\vectorl{D}}

%-------------
%--- overline ---
%-------------


\newcommand\ox{\overline{x}}
\newcommand\oy{\overline{y}}
\newcommand\oc{\overline{c}}
\newcommand\oa{\overline{a}}
\newcommand\ow{\overline{w}}
\newcommand\od{\overline{d}}
\newcommand\ob{\overline{b}}
\newcommand\oP{\overline{P}}
\newcommand\oee{\overline{e}}
\newcommand\ot{\overline{t}}
\newcommand\ou{\overline{u}}
\newcommand\ov{\overline{v}}
\newcommand\oz{\overline{z}}
\newcommand\am{\overline{a}}
\newcommand\om{\overline{m}}
\newcommand\oj{\overline{j}}

%----------------
%--- ensembles ---
%----------------

\newcommand\Q{\mathbb{Q}}
\newcommand\R{\mathbb{R}}
\newcommand\Z{\mathbb{Z}}
\providecommand\C{\mathbb{C}}
\newcommand\N{\mathbb{N}}
\newcommand\K{\mathbb{K}}
\newcommand\F{\mathbb{F}}

\newcommand\dyadic{\mathbb{D}}
\newcommand\Zd{{\Z_2}}
\newcommand\Zp{{\Z_p}}
%
\newcommand\RP{\mathbb{R}^{>0}}
\newcommand\QP{\mathbb{Q}^{>0}}


%----------------
%--- lettres cal ---
%----------------

\newcommand\mK{\mathcal{K}}
\newcommand\mD{\mathcal{D}}
\newcommand\mA{\mathcal{A}}
\newcommand\mL{\mathcal{L}}
\newcommand\mX{\mathcal{X}}
\newcommand\mG{\mathcal{G}}
\newcommand\mE{\mathcal{E}}
\newcommand\mH{\mathcal{H}}
\newcommand\mR{\mathcal{R}}
\newcommand\mF{\mathcal{F}}
\newcommand\mJ{\mathcal{J}}
\newcommand\mO{\mathcal{O}}
\newcommand\mP{\mathcal{P}}
\newcommand\mN{\mathcal{N}}
\newcommand\mS{\mathcal{S}}
\newcommand\mM{\mathcal{M}}
\newcommand\mC{\mathcal{C}}
\newcommand\mV{\mathcal{V}}
\newcommand\mI{\mathcal{I}}
\newcommand\mT{\mathcal{T}}

\newcommand\fC{\mathcal{C}}

%----------------
%--- lettres frak ---
%----------------


\newcommand\kM{\mathfrak{M}}
\newcommand\kN{\mathfrak{N}}
\newcommand\kU{\mathfrak{U}}
\newcommand\kg{\mathfrak{g}}



%--------------------------
%--- Façon noter noms classes ---
%---------------------------

\newcommand{\cp}[1]{\operatorname{#1}}


%--------------------------
%--- Calculabilité  ---
%---------------------------

% -- classes
\LANGUE{
\newcommand\RE{\cp{CE}}
}
{\newcommand\RE{\cp{RE}}
}

\newcommand\Rec{\cp{D}}
\newcommand\PR{\cp{PR}}
\newcommand\Gregn{\mE_n}


%--------------------------
%--- Réductions  ---
%---------------------------

\newcommand\lem{\le_{m}}
\newcommand\lelog{\le_{L}}
\newcommand\equivlog{\equiv_{L}}



%--------------------------
%--- Complexité  ---
%---------------------------

% COMPLEXITE

% -- P
\newcommand{\Ptime}{\cp{P}}      % P
\newcommand\PTIME{\Ptime}
\newcommand\PTIMEs[1]{{\PTIME}_{#1}}
\newcommand\PTIMEsp[1]{{\PTIME}_{#1}^0}
\renewcommand\P{\PTIME}
\newcommand{\FPtime}{\cp{FPTIME}}

% -- NP
\newcommand\NP{\cp{NP}}
\newcommand{\NPtime}{\NP}
\newcommand\NPs[1]{\cp{NP}_{#1}}
\newcommand{\FNP}{\cp{FNP}}

\newcommand\NDP{\cp{NDP}}
\newcommand\coNDP{\cp{coNDP}}
\newcommand\NDPs[1]{\cp{NDP}_{#1}}

% -- coNP
\newcommand\coNP{\cp{coNP}}

% -- PSPACE
\newcommand\PSPACE{\cp{ PSPACE}}
\newcommand\NPSPACE{\cp{ NPSPACE}}
\newcommand{\Pspace}{\PSPACE}
\newcommand{\FPspace}{\cp{FPSPACE}}


% -- LOGSPACE
\newcommand\LSPACE{\cp{LOGSPACE}}
\newcommand\NLSPACE{\cp{NLOGSPACE}}
\newcommand\coNLSPACE{\cp{coNLOGSPACE}}
\newcommand\coNL{\cp{coNL}}

% -- Autres
\newcommand\EXPTIME{\cp{EXPTIME}}
\newcommand\NEXPTIME{\cp{NEXPTIME}}
\newcommand\EXPSPACE{\cp{EXPSPACE}}
\newcommand\PH{\cp{PH}}
\newcommand\NC{\cp{NC}}
\newcommand\AC{\cp{AC}}
\newcommand\PPAD{\cp{\mathbf{PPAD}}}
\newcommand\PLS{\cp{\mathbf{PLS}}}

% -- Reverse math
\newcommand{\WKL}{\ensuremath{\docheck{\mathsf{WKL}_0}}}
\newcommand{\ACA}{\ensuremath{\docheck{\mathsf{ACA}_0}}}
\newcommand{\ATR}{\ensuremath{\docheck{\mathsf{ATR}_0}}}
\newcommand{\RCAO}{\ensuremath{\docheck{\mathsf{RCA}_0}}}
\newcommand{\PiCA}{\ensuremath{\docheck{\Pi^1_1\!\text{-}\mathsf{CA}_0}}}


% -- circuits
\newcommand\CIRCUIT[2]{\cp{CIRCUIT(#1,#2)}}
\newcommand\UCIRCUIT[2]{\cp{UCIRCUIT(#1,#2)}}

% -- Classes proba
\newcommand\PP{\cp{PP}}
\newcommand\RPP{\cp{RP}}
\newcommand\ZPP{\cp{ZPP}}
\newcommand\coRP{\cp{coRP}}
\newcommand\BPP{{\cp{BPP}}}

% -- IP
\newcommand\IP{\cp{IP}}
\newcommand\PCP{\cp{PCP}}

% -- non-uniforme
\newcommand\nuP{\mathbb{P}}
\newcommand\nuPs[1]{\mathbb{P}_{#1}}
\newcommand\nuPsp[1]{\mathbb{P}_{#1}^0}
\newcommand\nuNP{\mathbb{N}\mathbb{P}}
%\newcommand\poly{/{poly}}
\newcommand\Ppoly{\PTIME_{\poly}}
\newcommand\NPpoly{\NP_{\poly}}

% -- conditions d'uniformité
\newcommand\Luniforme{\cp{L}}
\newcommand\BCuniforme{\cp{BC}}


% Classes de circuits
\newcommand\TC{\cp{TC}}
\newcommand\LFP{\cp{LFP}}
\newcommand\FAC{\cp{FAC}}
\newcommand\FACC{\cp{FACC}}
\newcommand\FTC{\cp{FTC}}
\newcommand\FNC{\cp{FNC}}

% Logiques temorelles
\newcommand\CTL{\cp{CTL}}
\newcommand\LTL{\cp{LTL}}

% -- mesure temps/espace/taille

\newcommand\DTIME{\cp{DTIME}}
\newcommand\TIME[1]{\cp{TIME(#1)}}
\newcommand\TIMEs[2]{\cp{TIME_{#2}(#1)}}
\newcommand\TIMEsp[2]{\cp{TIME^0_{#2}(#1)}}
\newcommand\NTIME[1]{\cp{NTIME(#1)}}
\newcommand\NTIMEs[2]{\cp{NTIME_{#2}(#1)}}
\newcommand\NTIMEsp[2]{\cp{NTIME_{#2}^0(#1)}}
\newcommand\DSPACE{\cp{DSPACE}}
\newcommand\SPACE[1]{\cp{SPACE(#1)}}
\newcommand\SPACEs[2]{\cp{SPACE_{#2}(#1)}}
\newcommand\NSPACE[1]{\cp{NSPACE(#1)}}
\newcommand\NSPACEs[2]{\cp{NSPACE_{#2}(#1)}}
\newcommand\coNSPACE[1]{\cp{coNSPACE(#1)}}
\newcommand\TIMEON[2]{\cp{TIME(#1,#2)}}
\newcommand\SPACEON[2]{\cp{SPACE(#1,#2)}}
\newcommand\SIZE[1]{\cp{ size(#1)}}
\newcommand\ATIME{\cp{ATIME}}
\newcommand\ASPACE{\cp{ASPACE}}


%% autre:
\newcommand\LH{\cp{LH}}
\newcommand\FO{\cp{FO}}
\newcommand\SO{\cp{SO}}
\newcommand\ESOhorn{\cp{ESO_{horn}}}
\newcommand\SOhorn{\cp{SO_{horn}}}
\newcommand\ACzero{\cp{AC_0}}
\newcommand{\LINSPACE}{\cp{LINSPACE}}
\newcommand{\mFLINSPACE}{\mF_{\cp{LINSPACE}}}
\newcommand{\LOGSPACE}{\cp{LOGSPACE}}
\newcommand\ALOGTIME{\cp{ALOGTIME}}
\newcommand\RF{\cp{RF}}


\newcommand\BOOL[1]{\cp{{BOOLEAN(#1)}}}
\newcommand\DET{\cp{ DET}}

% % -- dirty


\newcommand{\cPspace}{\cp{\sharp PSPACE}}
\newcommand{\cAPtime}{\cp{\sharp APTIME}}
%\newcommand{\FPspace}{\cp{FPSPACE}}
%\newcommand{\FPspaceN}{\cp{FPSPACE}^{+}}
%\newcommand{\FPspaceN}{\cp{FPSPACE}}



%-------------------------
%--- Problèmes de décision  ---
%--------------------------

\newcommand\decisionname[1]{\cp{#1}}


\newcommand\Luniv{\decisionname{L_{univ}}}
\newcommand\HALTINGPROBLEM{\decisionname{HALTING-PROBLEM}}
\newcommand\SAT{\decisionname{SAT}}
\newcommand\MAXtSAT{\decisionname{MAX3SAT}}
\newcommand\REACH{\decisionname{REACH}}
\newcommand\CIRCUITVALUE{\decisionname{CIRCUITVALUE}}
\newcommand\MONOTONECIRCUITVALUE{\decisionname{MONOTONECIRCUITVALUE}}
\newcommand\THEOREMS{\decisionname{THEOREMS}}
\newcommand\QCIRCUITSAT{\decisionname{QCIRCUITSAT}}
\newcommand\QSAT{\decisionname{QSAT}}
\newcommand\QBF{\decisionname{QBF}}
\newcommand\SATVALUE{\decisionname{SATVALUE}}
\newcommand\SUCCINTSAT{\decisionname{SUCCINTSAT}}
\newcommand\SUCCINTSATVALUE{\decisionname{SUCCINTSATVALUE}}
\newcommand\GEOGRAPHY{\decisionname{GEOGRAPHY}}
\newcommand\PARITY{\decisionname{PARITY}}
\newcommand\MAJ{\decisionname{MAJ}}
\newcommand\CVP{\CIRCUITVALUE}
\newcommand\BOOLEANPROGRAMVALUE{BOOLEAN~PROGRAM~VALUE}

\newcommand\CLIQUE{\cp{CLIQUE}}
\LANGUEinv{
\newcommand\RECOUVREMENTDESOMMETS{\cp{RECOUVREMENT~DE~SOMMETS}}
\newcommand\RS{\RECOUVREMENTDESOMMETS}
}{
\newcommand\RECOUVREMENTDESOMMETSeng{\cp{VERTEX~COVER}}
\newcommand\RSeng{\RECOUVREMENTDESOMMETSeng}
}

\LANGUEinv{\newcommand\COUPUREMAXIMALE{\cp{COUPURE~MAXIMALE}}}
{\newcommand\COUPUREMAXIMALEeng{\cp{MAXIMAL~CUT}}}
\newcommand\CIRCUITHAMILTONIEN{\CHAMILTONIEN}
\LANGUEinv{\newcommand\VOYAGEURDECOMMERCE{\VCOMMERCE}}
{\newcommand\VOYAGEURDECOMMERCEeng{\VCOMMERCEeng}}

\LANGUEinv{\newcommand\CIRCUITLEPLUSLONG{\cp{CIRCUIT~LE~PLUS~LONG}}}
{\newcommand\CIRCUITLEPLUSLONGeng{\cp{LONGEST~CIRCUIT}}}
\LANGUEinv{\newcommand\SOMMEDESOUSENSEMBLE{\SUBSETSUM}}{
\newcommand\SOMMEDESOUSENSEMBLEeng{\SUBSETSUMeng}}
\LANGUEinv{
\newcommand\SACADOS{\cp{SAC~A~DOS}}}
{\newcommand\SACADOSeng{\cp{KNAPSACK}}}

%français
\LANGUEinv{
\newcommand\CHAMILTONIEN{\cp{CIRCUIT~HAMILTONIEN}}}
{\newcommand\CHAMILTONIENeng{\cp{HAMILTONIAN~CIRCUIT}}}
\LANGUEinv{
\newcommand\VCOMMERCE{\cp{VOYAGEUR~DE~COMMERCE}}}
{\newcommand\VCOMMERCEeng{\cp{TRAVELING~SALESMAN}}}
\LANGUEinv{
\newcommand\kCOLORABILITE{\cp{k-COLORABILITE}}}
{\newcommand\kCOLORABILITEeng{\cp{k-COLORABILITY}}}
\LANGUEinv{
\newcommand\COLORIAGE{\cp{BON~COLORIAGE}}}
{\newcommand\COLORIAGEeng{\cp{GOOD~COLORING}}}
\LANGUEinv{
\newcommand\tCOLORABILITE{\cp{3-COLORABILITE}}}
{\newcommand\tCOLORABILITEeng{\cp{3-COLORABILITY}}}
\LANGUEinv{
\newcommand\SUBSETSUM{\cp{SOMME~DE~SOUS~ENSEMBLE}}}
{\newcommand\SUBSETSUMeng{\cp{SUBSET~SUM}}}
\newcommand\PARTITION{\cp{PARTITION}}
\LANGUEinv{
\newcommand\NOMBREPREMIER{\decisionname{NOMBRE~ PREMIER}}}
{\newcommand\NOMBREPREMIEReng{\decisionname{PRIME~NUMBER}}}
\LANGUEinv{
\newcommand\CODAGE{\decisionname{CODAGE}}}
{\newcommand\CODAGEeng{\decisionname{ENCODING}}}
\newcommand\kSAT{\cp{k-SAT}}
\newcommand\tSAT{\cp{3-SAT}}
\newcommand\NAESAT{\cp{NAESAT}}
\newcommand\NAEtSAT{\cp{NAE3SAT}}
\newcommand\NAEqSAT{\cp{NAE4SAT}}
\newcommand\STABLE{\cp{\LANGUE{INDEPENDANT~SET}{STABLE}}}

%\newcommand\CM{\cp{COUPURE MAXIMALE}}
\LANGUEinv{
\newcommand\POSTpb{\cp{Problème~de~ la~correspondance~de~Post}}}
{\newcommand\POSTpbeng{\cp{Post~correspondence~problem}}}

\LANGUEinv{
\newcommand\PARITE{\cp{PARITE}}}
{\newcommand\PARITEeng{\cp{PARITY}}}


%-------------------------
%--- Modèles parallèles  ---
%--------------------------


% PARALLELISME
\newcommand\RAM{\cp{RAM}}
%
\newcommand\PRAM{\cp{PRAM}}
\newcommand\CRCW{\cp{CRCW}}
\newcommand\CREW{\cp{CREW}}
\newcommand\EREW{\cp{EREW}}


%------------------------
%--- codages  ---
%------------------------


\newcommand\codage[1]{\langle #1 \rangle}
\newcommand\tuple[1]{\codage{#1}}
\newcommand\paire[2]{\codage{#1,#2}}
\newcommand\triplet[3]{\codage{#1,#2,#3}}
\newcommand\quadruplet[4]{\codage{#1,#2,#3,#4}}
\newcommand\cinquplet[5]{\codage{#1,#2,#3,#4,#5}}


%------------------------
%--- descriptive set theory  ---
%------------------------

\newcommand\baire{\mathcal{N}}
\newcommand\peano{\overline{\mathcal{N}}} %% A MOI
\newcommand\cantor{\mathcal{C}}
%
\newcommand\bSigma{\mathbf{\Sigma}}
\newcommand\bPi{\mathbf{\Pi}}
\newcommand\bDelta{\mathbf{\Delta}}
%


%%%%%% ABOUT NOMS QUI CHANGES

\newcommand\WF{\cp{WF}}
\newcommand\WO{\WF}
\newcommand\LO{\cp{LO}}
\newcommand\DLO{\cp{DLO}}
% 
\newcommand\rk{\cp{rank}}


%% Moins propre:

\newcommand\I{\mathbb{I}}
\renewcommand\H{\mathbb{H}}
%
\newcommand\leKB{<_{KB}}
%
\newcommand\finiomega{{<\omega}}
\newcommand\compl{^\frown}
\newcommand\prefix{ ~ prefix~  } 
\newcommand\X{{X}}
\newcommand\Y{\mathcal{Y}}
\newcommand\subsetstrict{\subset}

\renewcommand\complement[1]{#1^{c}}


%------------------------
%--- schémas de récursion ---
%------------------------

%%                      Source papier MFCS 2019


% COMPLEXITY

%% MFCS differe de Clote:
%% Version à taille restreinte

\newcommand\co{\cp{co}}

\newcommand\mFPTIME{\mF_{PTIME}}
\newcommand{\mFPSPACE}{\mF_{SPACE}}
\newcommand{\mFPH}{\mF_{PH}}

\newcommand\SigmaP{\Sigma^P}
\newcommand\DeltaP{\Delta^{P}}
\newcommand\PiP{\Pi^{P}}
\newcommand\coSigmaP{\co\SigmaP}
\newcommand\coPiP{\mF{\co\PiP}}

%% autres classes:


%% Schemas pour complexité implicite

\newcommand\COMP{\cp{ COMP}}
\newcommand\CRN{\cp{ CRN}}
\newcommand\BRN{\cp{ BRN}}
\newcommand\SBRN{\cp{ SBRN}}
\newcommand\BR{\cp{ BR}}
\newcommand\kBR{k-\cp{ BR}}
\newcommand \WBPR{\cp{ WBPR}}
\newcommand \BMIN{\cp{ BMIN}}
\newcommand \BSUM{\cp{ BSUM}}
\newcommand\SBSUM{\cp{ SBSUM}}
\newcommand \BPROD{\cp{ BPROD}}
\newcommand \SBPROD{\cp{ SBPROD}}
\newcommand \SCOMP{\cp{ SCOMP}}
\newcommand \SR{\cp{ SR}}
\newcommand \SRN{\cp{ SRN}}

%% devrait etre défini
\newcommand\WBRN{\cp{ WBRN}}



%-- godel
\newcommand\INTDIV{\cp{INTDIV}}
\newcommand\DIV{\cp{DIV}}
\newcommand\EVEN{\cp{EVEN}}
\newcommand\ODD{\cp{ODD}}
\newcommand\PRIME{\cp{PRIME}}
\newcommand\POWER{\cp{POWER}}
\newcommand\BIT{\cp{BIT}}
\newcommand\negHP{\overline{\cp{HP}}}
\newcommand\HP{\cp{HP}}
\newcommand\negLuniv{\overline{\Luniv}}
\newcommand\VALCOMP{\cp{VALCOMP}}
\newcommand\mWINDOW{\cp{LEGAL}}
\newcommand\mmWINDOW{\cp{WINDOW}}
\newcommand\LENGTH{\cp{LENGTH}}
\newcommand\DIGIT{\cp{DIGIT}}
\newcommand\MATCH{\cp{MATCH}}
\newcommand\MOVE{\cp{MOVE}}
\newcommand\START{\cp{START}}
\newcommand\CELL{\cp{CELL}}
\newcommand\HALT{\cp{HALT}}
\newcommand\SUBST{\cp{SUBST}}
\newcommand\PROOF{\cp{PROOF}}
\newcommand\PROVABLE{\cp{PROVABLE}}
\newcommand\CONSIST{\cp{CONSIST}}



%---------------------
%--- calculus (AP) ---
%---------------------

% \jacobian{f}{x}
\newcommand{\jacobian}[1]{J_{#1}}
%\newcommand{\grad}[1]{\overrightarrow{\operatorname{grad}}~{#1}}
\newcommand{\grad}[1]{\nabla{#1}}
\newcommand{\scalarprod}[2]{{#1}\cdot{#2}}
\newcommand{\bigO}[1]{\mathcal{O}\left(#1\right)}
\newcommand\smallO{o}
\newcommand{\softO}[1]{\tilde{\mathcal{O}}\left(#1\right)}
\newcommand{\taylor}[3]{{T_{#1}^{#2}#3}}
\newcommand{\transpose}[1]{{#1}^T}
\newcommand{\pastsup}[2]{{\sup}_{#1}#2}

\DeclareMathOperator\arctanh{arctanh}



%---------------
%--- norms (AP) ---
%---------------


\newcommand{\inorm}[2]{\left\lVert{#1}\right\rVert_{#2}}
\newcommand{\twonorm}[1]{\inorm{#1}{2}}
\newcommand{\infnorm}[1]{\inorm{#1}{}}
\newcommand{\normsup}[1]{\infnorm{#1}{}}
%
\newcommand\hausdorff{d_H}
%
\newcommand\norm[1]{\infnorm{#1}}


%----------------
%--- topology ---
%----------------

% \openball{pt}{radius}
 \newcommand{\openball}[2]{B_{#2}(#1)}
\newcommand{\closedball}[2]{\ovl{B}_{#2}(#1)}
\newcommand{\closure}[1]{\ovl{#1}}
% \newcommand{\interior}[1]{\mathring{#1}}
% \newcommand{\border}[1]{\partial{#1}}




%-----------------
%--- functions (AP) ---
%-----------------
 
\newcommand{\lambdafun}[2]{(#1)\mapsto{#2}}
\newcommand{\lambdafunex}[2]{#1\mapsto{#2}}
\newcommand{\argmin}{\operatornamewithlimits{argmin}}
\newcommand{\argmax}{\operatornamewithlimits{argmax}}


 
%-----------------------
%--- Machines de Turing  ---
%-----------------------

\newcommand\blanc{\vectorl{B}}



%-----------------
%--- Logique ---
%-----------------
 
\newcommand\sem[1]{{[\semspace[}#1{]\semspace]}}
\newcommand\sems[1]{{[\semspace[}#1{]\semspace]}}
\newcommand\semss[1]{\sem{#1}_{\sigma'}}
\newcommand\semspace{\kern -.25ex}
\newcommand\true{true}
\newcommand\false{false}
%extension élémentaire:
\newcommand\elem{\preceq}
\newcommand\godel[1]{\lceil #1 \rceil}

%-----------------
%--- Model theory ---
%-----------------

\newcommand\hull[1]{\langle #1 \rangle}

%-----------------
%--- Ordinaux ---
%-----------------

\newcommand\omegaunCK{\omega_1^{CK}}
\newcommand\omegaunck{\omegaunCK}

\DeclareMathOperator{\cof}{cof}

%-----------------
%--- Surréels (QG) ---
%-----------------

\newcommand{\On}{\textbf{On}}
\newcommand{\Nobf}{\textbf{No}}
\newcommand{\Nobb}{\ensuremath{\mathbbmss{No}}}
\newcommand{\Nolt}[1]{\ensuremath{\Nobf_{< #1}}}
\newcommand{\Noleq}[1]{\ensuremath{\Nobf_{\leq #1}}}
\newcommand{\Noltbb}[1]{\ensuremath{\Nobb_{< #1}}}
\providecommand\No{\Nobf}
\newcommand\Noalpha{\Nolt{\alpha}}
\newcommand\Nolambda{\Nolt{\lambda}}

\DeclareMathOperator{\length}{length}
\DeclareMathOperator{\supp}{supp} % operateur support


%----------------------
%---Maths  de base (AP)  ---
%----------------------

\DeclareMathOperator{\dom}{dom} % operateur support

\DeclareMathOperator{\size}{size}

\newcommand{\powerset}[1]{\mathcal{P}(#1)}
\newcommand{\card}[1]{{\##1}}
\newcommand\priv{ - }
\newcommand\prive{-}
\renewcommand{\restriction}{\mathord{\upharpoonright}}
\newcommand\restrict[1]{_{\restriction #1}}

% \providecommand\mod{}
%% UNDEFINE:
\let\mod\relax
\let\div\relax
\DeclareMathOperator{\mod}{mod} 
\DeclareMathOperator{\div}{div} 

% partieentier
\newcommand\entieres[1]{\lceil #1 \rceil}
\newcommand\entierei[1]{\lfloor #1 \rfloor}


%----------------------
%---Symbole danger ---
%----------------------


%\providecommand*{\TakeFourierOrnament}[1]{{%
%\fontencoding{U}\fontfamily{futs}\selectfont\char#1}}
%\providecommand*{\danger}{{\Huge \TakeFourierOrnament{66}}}



%%%%%%%%%%%%%%%%%%%%%%%%%%%%%%%%%%%%%%%%%%%%%%%%%
%
%
%                      ASMeries
%
%
%%%%%%%%%%%%%%%%%%%%%%%%%%%%%%%%%%%%%%%%%%%%%%%%%



% CIRCUITS
%-- particuliers
\newcommand\REPLACE{\cp{REPLACE}}
\newcommand\EQUAL{\cp{EQUAL}}
\newcommand\EVALUE{\cp{EVALUE}}
\newcommand\TVALUE{\cp{TVALUE}}
\newcommand\UPDATE{\cp{UPDATE}}
\newcommand\ASSIGN{\cp{ASSIGN}}
\newcommand\PAR{\cp{PAR}}
\newcommand\DO{\cp{DO}}


% ASM
\newcommand\emplacement[2]{(#1,#2)}
\newcommand\contenu[2]{\sem{(#1,#2)}}
\newcommand\affecte[3]{(#1,#2)\leftarrow#3}
\newcommand\ctl{ctl\_state}


%%%%%%%%%%%%%%%%%%%%%%%%%%%%%%%%%%%%%%%%%%%%%%%%%
%
%
%                      GPACqueries
%
%
%%%%%%%%%%%%%%%%%%%%%%%%%%%%%%%%%%%%%%%%%%%%%%%%%



%------------------------
%--- generable zoo (GPAC) ---
%------------------------

\newcommand{\myop}[1]{\operatorname{#1}}
\newcommand{\abs}{\myop{abs}}
\newcommand{\sg}{\myop{sg}}
\newcommand{\ip}[1]{\myop{ip}_{#1}}
\newcommand{\rnd}{\myop{rnd}}
\newcommand{\lil}{\myop{lil}}
\newcommand{\lxh}{\myop{lxh}}
\newcommand{\hxl}{\myop{hxl}}
\newcommand{\plil}{\myop{plil}}
\newcommand{\reach}{\myop{reach}}
%\newcommand{\mreach}{\myop{mreach}}
\newcommand{\watchdog}{\myop{watchdog}}
\newcommand{\monitor}{\myop{monitor}}
%\newcommand{\norm}{\myop{norm}}
\newcommand{\mx}{\myop{mx}}
\newcommand{\mn}{\myop{mn}}
\newcommand{\sample}{\myop{sample}}
\newcommand{\select}{\myop{select}}
\newcommand{\ssine}[1]{\sin_{#1}}
\newcommand{\clamp}{\myop{clamp}}

\newcommand{\fnsg}[3]{tanh((#1)*(#2)*(#3))}
\newcommand{\fnipone}[3]{(1+\fnsg{(#1)-1}{#2}{#3})/2.}
\newcommand{\fnipinf}[1]{#1-sin(2*pi*(#1))/2./pi}%
\newcommand{\fnlil}[5]{%
\fnipone{(#3)-(#1)+1-((#2)-(#1))/8.}{#4+log(1+4./(#2-#1)*(#5)*(#5))+log(1+4)}{8./(#2-#1)}%
*\fnipone{#2-#3+1-(#2-#1)/8.}{#4+log(1+4/(#2-#1)*(#5)*(#5))+log(1+4)}{8./(#2-#1)}%
*4./(#2-#1)*(#5)}
\newcommand{\fnlxh}[5]{\fnipone{(#3)-(#1+#2)/2.+1}{#4+log(1+(#5)*(#5))}{2./(#2-#1)}*(#5)}
\newcommand{\fnhxl}[5]{\fnipone{(#1+#2)/2.-(#3)+1}{#4+log(1+(#5)*(#5))}{2./(#2-#1)}*(#5)}
\newcommand{\fnplilf}[4]{sin(((#4)-((#1)+(#2))/2.+(#3)/4.)*2.*pi/(#3))}
\newcommand{\fnplil}[6]{(#6)*\fnlxh{\fnplilf{#1}{#2}{#3}{#1}}%
{\fnplilf{#1}{#2}{#3}{#1+((#2)-(#1))/4.}}%
{\fnplilf{#1}{#2}{#3}{#4}}%
{#5+2.+log(1+(#6)*(#6))}%
{1./4.+2./((#2)-(#1))}}
\newcommand{\fnssine}[2]{sin(#2)*(1-cos(#2)/cos(#1))}



%--------------------
%--- complexity GPAC ---
%--------------------


%\ifthenelse{\equal{#1}{}}{Empty.}{Nonempty.}
\newcommand{\myclass}[1]{\operatorname{#1}}
% \newcommand{\gcomp}[2]{\ensuremath{\myclass{GCOMP}(#1,#2)}}
% \newcommand{\ggenerable}[1]{\textcolor{red}{\ensuremath{#1-}\text{generable}}}
 \newcommand{\gval}[2][]{\ensuremath{\myclass{GVAL}_{#1}\ifthenelse{\equal{#2}{}}{}{[#2]}}}
 \newcommand{\gpval}[1][]{\ensuremath{\myclass{GPVAL}_{#1}}}
 \newcommand{\gc}[3][]{\ensuremath{\myclass{ATSC}(#2,#3)}}
 \newcommand{\gpc}[1][]{\ensuremath{\myclass{ATSP}}}
\newcommand{\cgc}[1][]{\ensuremath{\myclass{ATS}}}
\newcommand{\gexpc}[1][]{\ensuremath{\myclass{AEXP}}}
\newcommand{\gwc}[2]{\ensuremath{\myclass{AW}(#1,#2)}}
\newcommand{\gpwc}{\ensuremath{\myclass{AWP}}}
\newcommand{\cgwc}{\ensuremath{\myclass{AW}}}
\newcommand{\grc}[3]{\ensuremath{\myclass{AR}(#1,#2,#3)}}
\newcommand{\gprc}{\ensuremath{\myclass{ARP}}}
\newcommand{\gsc}[3]{\ensuremath{\myclass{AS}(#1,#2,#3)}}
\newcommand{\gpsc}{\ensuremath{\myclass{ASP}}}
\newcommand{\goc}[3]{\ensuremath{\myclass{AO}(#1,#2,#3)}}
\newcommand{\gpoc}{\ensuremath{\myclass{AOP}}}
\newcommand{\cgoc}{\ensuremath{\myclass{AO}}}
\newcommand{\guc}[4]{\ensuremath{\myclass{AX}(#1,#2,#3,#4)}}
\newcommand{\gpuc}{\ensuremath{\myclass{AXP}}}
\newcommand{\glc}[2][]{\ensuremath{\myclass{AL}(#2)}}
\newcommand{\gplc}[1][]{\ensuremath{\myclass{ALP}}}
\newcommand{\cglc}[1][]{\ensuremath{\myclass{AL}}}

%\newcommand{\intinterv}[2]{\{#1,\ldots,#2\}}
\newcommand{\intinterv}[2]{\llbracket#1,#2\rrbracket}

\newcommand{\Rgen}{\R_{G}}
\newcommand{\Rpoly}{\R_{P}}

%%%%%%%%%%%%%%%%%%%%%%%%%%%%%%%%%%%%%%%%%%%%%%%%%
%
%
%                      Mes codages obscures
%
%
%%%%%%%%%%%%%%%%%%%%%%%%%%%%%%%%%%%%%%%%%%%%%%%%%




%%% Codages

% \newcommand\gammaoc{\gamma_{[0,1]}}
% \newcommand\gammaonc{\gamma_{\N}}
\newcommand\gammaTMc{\gamma^{TM}_{[0,1]}}
\newcommand\gammaTMcatan{\gamma^{TM}_{(0,1)^2}}
\newcommand\gammaTMnc{\gamma^{TM}_{\N^3}}
\newcommand\gammaTMncd{\gamma^{TM}_{\N^2}}
% \newcommand\gammaTMe{\gamma^{TM}_{*}}
% \newcommand\gammas{\gamma_{sab}}
% \newcommand\gammaord{\gamma_{cantor}}


\newcommand\enccompact{\nu_{X}}
 \newcommand\encinfinite{\nu_\N}
 \newcommand\enc{\nu}
 
%%%%%%%%%%%%%%%%%%%%%%%%%%%%%%%%%%%%%%%%%%%%%%%%%
%
%
%                      DESSINS ET IMAGES
%
%
%%%%%%%%%%%%%%%%%%%%%%%%%%%%%%%%%%%%%%%%%%%%%%%%%


\newcommand\IMAGE[2]{ \begin{center} 
    \includegraphics[#1]{#2}
  \end{center}
  }


  
%%%% POUR FAIRE DESSIN.
\newenvironment{dessinpp}[1]{
\begin{tikzpicture}[baseline=(current bounding
      box.west),#1]
}
{\end{tikzpicture}}




\newenvironment{dessinp}[1]{
\begin{center}\begin{tikzpicture}[baseline=(current bounding
      box.north),#1]
}
{\end{tikzpicture}
\end{center}}


\newenvironment{dessin}{
\begin{center}\begin{tikzpicture}[baseline=(current bounding
      box.north)]
}
{\end{tikzpicture}
\end{center}}



\newenvironment{dessind}{
\begin{tikzpicture}[baseline=(current bounding
      box.north)]
}
{\end{tikzpicture}
}



%%%%%%%%%%%%%%%%%%%%%%%%%%%%%%%%%%%%%%%%%%%%%%%%%
%
%
%                      POUR DESSINS GPAC
%
%
%%%%%%%%%%%%%%%%%%%%%%%%%%%%%%%%%%%%%%%%%%%%%%%%%


%%%% MACROS DE BASE.


\newcommand\constant[2]
{
node (#1)  [shape=circle,draw] {#2};
\draw (#1.east) -- +(0.3,0) coordinate (#1-o) ; 
\draw (#1.west) -- +(-0.3,0) coordinate (#1-i) ; 
}

\newcommand\basicproduct[2]
{
node (#1) {#2};
\draw (#1) +(-0.4,0.4) -- +(0.4,0.4) -- +(0.8,0) -- +(0.4,-0.4) -- +(-0.4,-0.4) --
+(-0.4,0.4);
\draw (#1) ++(0.8,0) -- ++(+0.3,0) coordinate (#1-o) ;
\draw (#1) ++ (-0.4,0.2) -- +(-0.3,0) coordinate (#1-i1) ;
\draw (#1) ++(-0.4,-0.2) -- +(-0.3,0) coordinate (#1-i2) ;
}

\newcommand\product[1]{\basicproduct{#1}{$+\Pi$}}
\newcommand\negativeproduct[1]{\basicproduct{#1}{$-\Pi$}}


%%BEGIN INTERNAL
\newcommand\basicsummer[1]{
coordinate (#1) {};
\draw (#1) +(-0.5,0.6) -- +(0.5,0) -- +(-0.5,-0.6) -- +(-0.5,0.6);
\draw (#1) ++(0.5,0) -- +(+0.3,0) coordinate (#1-o) ;
}

\newcommand\basicintegrator[1]{
coordinate (#1) {};
\draw (#1) +(-0.5,0.6) -- +(0.5,0) -- +(-0.5,-0.6) -- +(-0.5,0.6);
\draw (#1) ++(0.5,0) -- +(+0.3,0) coordinate (#1-o) ;
\draw (#1) +(-0.5,-0.6) -| +(-0.8,0.6) -- +(-0.5,0.6) ;
 \draw (#1) ++(-0.65,0.6) -- +(0,+0.3) coordinate (#1-init) ;
}

\newcommand\splitfour[2]{
\draw (#2) ++ (#1,0.4) -- +(-0.3,0) coordinate (#2-i1) ;
\draw (#2) ++ (#1,0.15) -- +(-0.3,0) coordinate (#2-i2) ;
\draw (#2) ++ (#1,-0.15) -- +(-0.3,0) coordinate (#2-i3) ;
\draw (#2) ++ (#1,-0.4) -- +(-0.3,0) coordinate (#2-i4) ;
}

\newcommand\splitthree[2]{
\draw  (#2) ++(#1,0.3) -- +(-0.3,0) coordinate (#2-i1) ;
\draw  (#2) ++(#1,0) -- +(-0.3,0) coordinate (#2-i2) ;
\draw  (#2) ++(#1,-0.3) -- +(-0.3,0) coordinate (#2-i3) ;
}

\newcommand\splittwo[2]{
\draw  (#2) ++(#1,0.2) -- +(-0.3,0) coordinate (#2-i1) ;
\draw  (#2) ++(#1,-0.2) -- +(-0.3,0) coordinate (#2-i2) ;
}

\newcommand\splitone[2]{
\draw (#2) ++(#1,0) -- +(-0.3,0) coordinate (#2-i1) ;
}

%% END INTERNAL

\newcommand\summerfour[1]{
\basicsummer{#1}
\splitfour{-0.5}{#1}
}

\newcommand\summerthree[1]{
\basicsummer{#1}
\splitthree{-0.5}{#1}
}

\newcommand\summertwo[1]{
\basicsummer{#1}
\splittwo{-0.5}{#1}
}

\newcommand\summerone[1]{
\basicsummer{#1}
\splitone{-0.5}{#1}
}


\newcommand\integratorfour[1]{
\basicintegrator{#1}
\splitfour{-0.8}{#1}
<}


\newcommand\integratorthree[1]{
\basicintegrator{#1}
\splitthree{-0.8}{#1}
}


\newcommand\integratortwo[1]{
\basicintegrator{#1}
\splittwo{-0.8}{#1}
}


\newcommand\integratorone[1]{
\basicintegrator{#1}
\splitone{-0.8}{#1}
}

%%%%% FIN MACROS




%%%%%%%%%%%%%%%%%%%%%%%%%%%%%%%%%%%%%%%%%%%%%%%%%
%
%
%                      Trucs d'Amaury
%
%
%%%%%%%%%%%%%%%%%%%%%%%%%%%%%%%%%%%%%%%%%%%%%%%%%


%----------------
%--- ref (AP) ---
%----------------

% \newcommand{\defref}[1]{Definition~\ref{#1}}
\newcommand{\lemref}[1]{Lemma~\ref{#1}}
\newcommand{\thref}[1]{Theorem~\ref{#1}}
\newcommand{\amaurycorref}[1]{Corollary~\ref{#1}}
 \newcommand{\figref}[1]{Figure~\ref{#1}}
% \newcommand{\figrefmult}[1]{Figures~{#1}}
% \newcommand{\secref}[1]{Section~\ref{#1}}
% %\renewcommand{\algref}[1]{Algorithm~\ref{#1}}
 \newcommand{\propref}[1]{Proposition~\ref{#1}}
% \newcommand{\exref}[1]{Example~\ref{#1}}
\newcommand{\remref}[1]{Remark~\ref{#1}}


 %-------------------
%--- polynomials ---
%-------------------

\newcommand{\sigmap}[1]{{\Sigma{#1}}}
\newcommand{\degp}[1]{{\operatorname{deg}(#1)}}
\newcommand{\poly}{\operatorname{poly}}


 %----------------------
%--- misc functions ---
%----------------------

\newcommand{\sgn}[1]{\operatorname{sgn}(#1)}
\newcommand{\intp}{\operatorname{int}}
 \newcommand{\fracp}{\operatorname{frac}}
\newcommand{\fiter}[2]{#1^{[#2]}}
\newcommand{\frestrict}[2]{#1\restriction_{#2}}
\newcommand{\indicator}[1]{\mathds{1}_{#1}}
\newcommand{\idfun}{\operatorname{id}}
% \newcommand{\floor}[1]{\left\lfloor#1\right\rfloor}
 \newcommand{\ceil}[1]{\left\lceil#1\right\rceil}
\newcommand{\round}[1]{\left\lfloor#1\right\rceil}
\DeclareMathOperator\erf{erf}


%------------
%--- sets ---
%------------

\newcommand{\A}{\mathbb{A}}
%\newcommand{\D}{\mathbb{D}}
\newcommand{\Ns}{\mathbb{N}^*}
%\newcommand{\C}{\mathbb{C}}
%\newcommand{\K}{\mathbb{K}}
% \MAT{height}{[width]}{[field]}
\newcommand{\MAT}[3]{M_{#1\ifthenelse{\equal{#2}{}}{}{,#2}}\ifthenelse{\equal{#3}{}}{}{\left(#3\right)}}

%\newcommand{\Rcomp}{\R_{c}}
\newcommand{\Rp}{\R_{+}}
\newcommand{\Rs}{\R^{*}}
\newcommand{\Rps}{\R_{+}^{*}}

%\newcommand{\Analytic}{C^\omega}
%\newcommand{\PTIME}{\ensuremath{\operatorname{P}}}
%\newcommand{\EXPTIME}{\ensuremath{\operatorname{EXPTIME}}}
%\newcommand{\NP}{\ensuremath{\operatorname{NP}}}
%\newcommand{\NPTIME}{\NP}
\newcommand{\FP}{\ensuremath{\operatorname{FP}}}
%\newcommand{\PSPACE}{\ensuremath{\operatorname{PSPACE}}}
% \newcommand{\FPSPACE}{\ensuremath{\operatorname{FPSPACE}}}



%%%%%%%%%%%%%%%%%%%%%%%%%%%%%%%%%%%%%%%%%%%%%%%%%
%
%
%                      Trucs de Quentin
%
%
%%%%%%%%%%%%%%%%%%%%%%%%%%%%%%%%%%%%%%%%%%%%%%%%%


\newcommand{\fct}[5]{\ensuremath{#1 : \left\{\begin{array}{c c c}
#2 & \rightarrow & #3\\
#4 & \mapsto & #5
                                             \end{array}\right.}}

%Nouvelle fraction, plus jolie
\newcommand{\f}[2]{	
	\mathchoice%
		{\dfrac{#1}{#2}}
    	{\dfrac{#1}{#2}}
		{\frac{#1}{#2}}
		{\frac{#1}{#2}}
}


%%Ensembles et combinatoires

%Ensemble avec propriété à drotie \{x | blabla\}
\newcommand{\enstq}[2]{\left\{ \left.#1\ \vphantom{#2}\right|\ #2  \right\}}
%meme chose mais avec des crochets
\newcommand{\crotq}[2]{\left[ \left.#1\ \vphantom{#2}\right|\ #2  \right]}
%même chose mais avec des brackets
\newcommand{\bratq}[2]{\left\langle \left.#1\ \vphantom{#2}\right|\ #2  
  \right\rangle}

\newcommand{\intn}[2]{\ensuremath{
		\left\llbracket\, #1\ ;\ #2\,\right\rrbracket
              }}


%%%%%%%%%%%%%%%%%%%%%%%%%%%%%%%%%%%%%%%%%%%%%%%%%
%
%
%               ENVIRONNEMENTS THEOREMS & CO
%
%
%%%%%%%%%%%%%%%%%%%%%%%%%%%%%%%%%%%%%%%%%%%%%%%%%


%----------------
%--- pas dans lipics ---
%-----------------


\ifnotSLIDE{
% commente car bug Fevrier 2026: \ifnotTDEXAM{\spnewtheorem{solution}{Solution}{\bfseries}{\rmfamily}}
\spnewtheorem{openproblem}{Open Problem}{\bfseries}{\rmfamily}
\spnewtheorem{remarkolivier}{Commentaire d'Olivier: }{\bfseries}{\rmfamily}
\spnewtheorem{postulate}{Postulate}{\bfseries}{\rmfamily}
}

%----------------
%--- dans lipics ---
%-----------------
%
\newenvironment{exercice}[1]{\begin{exercise} \label{exo:#1-stt}}{\end{exercise}}
\newenvironment{exerciced}[1]{\begin{exercisee}\label{exo:#1-stt}}{\end{exercisee}}
\newenvironment{exercicec}[1]{\begin{exercise} \label{exo:#1-stt} (\LANGUE{solution on}{corrigé} page \pageref{exo:#1-cor})
}{\end{exercise}}
\newenvironment{exercicedc}[1]{\begin{exercisee} \label{exo:#1-stt} (\LANGUE{solution on}{corrigé} page \pageref{exo:#1-cor})
}{\end{exercisee}}

\newenvironment{slideproposition}{{\bf Proposition}}{}

\ifnotSLIDE{
\nolipics{
\nolncs{\spnewtheorem{theorem}{\LANGUE{Theorem}{Théorème}}{\bfseries}{\rmfamily}}
\nolncs{\spnewtheoremi{definition}{\LANGUE{Definition}{Définition}}{\bfseries}{\rmfamily}{theorem}}
\nolncs{\spnewtheoremi{lemma}{\LANGUE{Lemma}{Lemme}}{\bfseries}{\rmfamily}{theorem}}
\nolncs{\spnewtheoremi{corollary}{\LANGUE{Corollary}{Corollaire}}{\bfseries}{\rmfamily}{theorem}}
\nolncs{\spnewtheoremi{proposition}{Proposition}{\bfseries}{\rmfamily}{theorem}}
\nolncs{\spnewtheoremi{example}{\LANGUE{Example}{Exemple}}{\bfseries}{\rmfamily}{theorem}}
%\spnewtheorem{sketch}{{\bf \LANGUE{Sketch of proof}{Démonstration (principe)}}}{\bfseries}{\rmfamily}
\nolncs{\spnewtheoremi{remark}{\LANGUE{Remark}{Remarque}}{\bfseries}{\rmfamily}{theorem}}
\nolncs{\spnewtheorem{exercise}{\LANGUE{Exercise}{Exercice}}{\bfseries}{\rmfamily}}
\spnewtheorem{exercisee}{\LANGUE{*Exercise}{*Exercice}}{\bfseries}{\rmfamily}
\spnewtheorem{summary}{\LANGUE{Summary}{Résumé}}{\bfseries}{\rmfamily}
 %\spnewtheorem{exerciceenglish}{Exercise}{\bfseries}{\rmfamily}
 %
 \newenvironment{sketch}{{\bf \LANGUE{Sketch of proof}{Démonstration (principe)}}:}{\hfill $\Box$ }
\nolncs{\spnewtheorem{conjecture}{\LANGUE{Conjecture}{Conjecture}}{\bfseries}{\rmfamily}}
 %%
 \spnewtheorem{theoremaprouver}{\LANGUE{Theorem TO BE PROVED}{Théorème (A PROUVER)}}{\bfseries}{\rmfamily}
  \spnewtheorem{conjectureaprouver}{\LANGUE{Conjecture TO BE VERIFIED}{Conjecture (A VERIFIER)}}{\bfseries}{\rmfamily}
 }}
            
 \makeatletter
 \ifnotSLIDE{
  \ifnotTDEXAM{
\@ifundefined{proof}{
	\newenvironment{proof}{{\bf \LANGUE{Proof}{Démonstration}}:}{\hfill $\Box$ 
	}{}
	}}}
\makeatother

 
%%%%%%%%%%%%%%%%%%%%%%%%%%%%%%%%%%%%%%%%%%%%%%%%%
%
%
%                      Macros dont pas fier du tout
%
%
%%%%%%%%%%%%%%%%%%%%%%%%%%%%%%%%%%%%%%%%%%%%%%%%%



\newcommand\equations[1]{
\begin{array}{lll}
#1
\end{array}
}



\newcommand\systeme[1]{
\left\{
\begin{array}{lll} 
#1
\end{array}
\right.
}



%------------------
%--- shorthands (AP)---
%------------------
\newcommand{\mtt}[1]{\mathtt{#1}}
\newcommand{\ovl}[1]{\overline{#1}}
\newcommand{\panic}[1]{\textcolor{red}{#1}}
%\NewEnviron{epanic}{{\color{red}\BODY}}
\newcommand{\idest}{\emph{i.e.\thinspace}}


%------------------
%--- mes trucs  ---
%------------------


\newcommand\implique{\Rightarrow}
\newcommand\ssi{\Leftrightarrow}
\newcommand\ac[1]{\{#1\}}
\newcommand\zero{\vectorl{0}}
\newcommand\un{\vectorl{1}}

\newcommand\technique[1]{{\scriptsize #1}}


%% Preuve and Co

\newcommand\sensdirect{ Direction $\Rightarrow$. }
\newcommand\sensindirect{ Direction $\Leftarrow$. }




\newcommand\donnee{\mbox{<donnée>}}



% pour aller vite

\newcommand\gM{\mathfrak{M}}
\newcommand\gMM{(M,f_1,\cdots,f_u,r_1,\cdots,r_v)}
%
\newcommand\mygM{\mK}
\newcommand\mygMM{\left ( \K, \{op_i\}_{i\in I}, r_1 \ldots,
r_\ell, \zero,\un \right )}
\newcommand\myM{\K} 

\newcommand\gMMM{(M,f_1,\cdots,f_u,c_1,\cdots,c_k,r_1,\cdots,r_v)}
\newcommand\gMMMM{(f_1,\cdots,f_u,r_1,\cdots,r_v)}
\newcommand\gMME{(M,f_1,\cdots,f_u,f'_1,\cdots,f'_\ell,r_1,\cdots,r_v)}
\newcommand\gMMS{(f_1,\cdots,f_u,f'_1,\cdots,f'_\ell,r_1,\cdots,r_v)}

\newcommand\bool{\mathfrak{B}}
\newcommand\gbool{(\{0,1\},\zero,\un,=)}

\newcommand\osg{\overline{sg}}
\newcommand\moins{\mathrel{\dot{-}}}


%%% TRUC DE BARWISE


\newcommand\UM{\kU_{\kM}}
\newcommand\VV{\mathbb{V}}
% Hereditarily finite
\newcommand\IHF{\mathbb{I}HF}
\newcommand\IHP{\IHYP}
\newcommand\IHYP{\mathbb{I}HYP}

\providecommand\citep[1]{\cite{#1}}

%% Cofinali


% probas
% défini dans amsmath
\renewcommand\Pr{\cp{Pr}}
\newcommand\Var{\cp{Var}}


% -- mesures
%\newcommand\card{card}
\newcommand\taille{\LANGUE{size}{taille}}
\newcommand\entree{e}
\newcommand\profondeur{\LANGUE{depth}{profondeur}}


%%%%%%%%%%%%%%%%%%%%%%%%%%%%%%%%%%%%%%%%%%%%%%%%%%%%%%%%%%%%%%%%%%%%%%%%%%%%
%                 MACHINES DE TURING (graphique)
%%%%%%%%%%%%%%%%%%%%%%%%%%%%%%%%%%%%%%%%%%%%%%%%%%%%%%%%%%%%%%%%%%%%%%%%%%%%



% Configuration de machine de Turing
\newcommand\configTuring[3]{#1\textbf{#2}#3}


%%% DESSIN CROIX PORU AIDER  A DESSINER

\newcommand\croix{  +(-1,-1) -- +(1,1) +(-1,1) --
  +(1,-1) }

%% DESSIN ACCOLADE

\newcommand\ACCOLADE[3]{
%#1: coordonnées de départ
%#2: coordonnées arrivée
%#3: texte
\draw[decorate,decoration={brace}]  (#1) -- (#2)
 node[midway,above=0.1cm] {#3};
}

%%%


\def\lcase{3.5 ex}
\def\hcase{3.5ex}
\tikzstyle{case} = []
%[draw,minimum width = \lcase,minimum height = 3.6ex,baseline]
\newcommand\dcase{ +(-0.5*\lcase,-0.5*\hcase) rectangle +(0.5*\lcase,0.5*\hcase)}
%\newcommand\lettre[1]{\node [name=l,case,on chain=1] {#1};}


\newcommand\RUBANMACHINEDETURING[6]{
%#1: texte
%#2: position de la tête
%#3: nb de blancs à gauche
%#4: nb de cases totales
%#5: texte au dessus du ruban
%#6: position du coin
    \draw (#6) +(\lcase,4ex ) node {#5};
    \draw (#6) +(-0.6,0)  node  {\ldots};  
    \draw (#6) node[coordinate] (t) {};
    \foreach \x in {1,2,...,#3} {
         \draw (t) node [case]   {$\blanc$} \dcase;
         \draw (t) ++(\lcase,0) node [coordinate] (t) {};
       };
   \StrLen{#1}[\Resultat] % ATTENTION IL FAUT CETTE SYNTAXE POUR FAIRE
                         % EQUIVALENT D UN EDEF
  
\foreach \x in {1,...,\Resultat} 
             { 
        \draw (t) node [case]   {\StrChar{#1}{\x}} \dcase;
        \draw (t) ++(\lcase,0) node [coordinate] (t) {};
             }
\pgfmathtruncatemacro{\z}{#4 - \Resultat - #3}

  \foreach \x in {1,2,..., \z} {
        \draw (t) node [case]   {$\blanc$} \dcase;
        \draw (t) ++(\lcase,0) node [coordinate] (t) {};
}
%% Convention: graphiquement, la case numéro n se trouve à (n*\lcase,0)
%% en comptant les case à partir de $0$

   \draw  (t) +(0.1,0) node [name=r] {\ldots}; 
}

\newcommand\RUBANMACHINEDETURINGtexte[6]{
%#1: texte
%#2: position de la tête
%#3: nb de blancs à gauche
%#4: nb de cases totales (équivalent)
%#5: texte au dessus du ruban
%#6: position du coin
    \draw (#6) +(\lcase,4ex ) node {#5};
    \draw (#6) +(-0.6,0)  node  {\ldots};  
    \draw (#6) node[coordinate] (t) {};
   \foreach \x in {1,2,...,#3} {
        \draw (t) node [case]   {$\blanc$} \dcase;
        \draw (t) ++(\lcase,0) node [coordinate] (t) {};
      };
   \draw (#6) ++(-2*\lcase,0.5*\lcase)-- +(\lcase*#4,0);
   \draw (#6) ++(-2*\lcase,-0.5*\lcase) -- +(\lcase*#4,0);
   
    \draw (t) ++(-0.5*\lcase,0) node [case,anchor=west]  (te) {#1};
    \draw (te.east)  node [anchor=west,coordinate] (t) {};

  \draw (t) ++(0.2*\lcase,0) node [coordinate] (t) {};
   \foreach \x in {1,2,...,#3} {
        \draw (t) node [case]   {$\blanc$} \dcase;
        \draw (t) ++(\lcase,0) node [coordinate] (t) {};
      };

    \draw  (t) +(0.1,0) node [anchor=west] {\ldots}; 
}


\newcommand\RUBANMACHINEDETURINGTETE[6]{
%
% dessine la tete
%
% output: (tete)
%
\RUBANMACHINEDETURING{#1}{#2}{#3}{#4}{#5}{#6}
% TETE
     \draw (#6) + ( #2 * \lcase + #3  *\lcase,- 0.5*\lcase )
     node[coordinate] (k) {};
     \draw[->] (k) +(0,-0.5) -- (k);
     \draw (k) +(0,-0.5) node[coordinate] (tete) {};
}


\newcommand\RUBANMACHINEDETURINGTETEtexte[6]{
%
% dessine la tete
%
% output: (tete)
%
\RUBANMACHINEDETURINGtexte{#1}{#2}{#3}{#4}{#5}{#6}
% TETE
     \draw (#6) + ( #2 * \lcase + #3  *\lcase,- 0.5*\lcase )
     node[coordinate] (k) {};
     \draw[->] (k) +(0,-0.5) -- (k);
     \draw (k) +(0,-0.5) node[coordinate] (tete) {};
}


\newcommand\MACHINEDETURING[6]{
%#1: texte
%#2: position de la tête
%#3: état
%#4: nb de blancs à gauche
%#5: nb de cases totales
%#6: texte en haut du ruban
%\begin{tikzpicture}[node distance=-0.15mm]

\RUBANMACHINEDETURINGTETE{#1}{#2}{#4}{#5}{#6}{0,0}
%
%% ETAT:
     \draw (tete) node [name=etat,draw,circle,anchor=north]  {#3};
     \draw (etat) -- (tete);
%\end{tikzpicture}
}

\newcommand\MACHINEDETURINGtexte[6]{
%#1: texte
%#2: position de la tête
%#3: état
%#4: nb de blancs à gauche
%#5: nb de cases totales
%#6: texte en haut du ruban
%\begin{tikzpicture}[node distance=-0.15mm]

\RUBANMACHINEDETURINGTETEtexte{#1}{#2}{#4}{#5}{#6}{0,0}
%
%% ETAT:
     \draw (tete) node [name=etat,draw,circle,anchor=north]  {#3};
     \draw (etat) -- (tete);
%\end{tikzpicture}
}

\newcommand\MACHINEDETURINGTROISRUBANS[9]{
%#1: texte1
%#2: position de la tête 1
%#3: texte2
%#4: position de la tête 2
%#5: texte3
%#6: position de la tête 3
%#7: état
%#8: nb de blancs à gauche 1
%#9: nb de cases totales
%\begin{tikzpicture}[node distance=-0.15mm]

\RUBANMACHINEDETURINGTETE{#1}{#2}{#8}{#9}{\LANGUE{tape}{ruban} $1$}{0,0}
\draw (tete) node[coordinate] (tete1) {};
\draw (tete1) +(8,0) node[coordinate] (tete1b) {};
\draw (tete1) -| (tete1b) ;

\RUBANMACHINEDETURINGTETE{#3}{#4}{#8}{#9}{\LANGUE{tape}{ruban} $2$}{0,-2}
\draw (tete) node[coordinate] (tete2) {};
\draw (tete2) +(-8.5,0) node[coordinate] (tete2b) {};
\draw (tete2) -| (tete2b) ;

\RUBANMACHINEDETURINGTETE{#5}{#6}{#8}{#9}{\LANGUE{tape}{ruban} $3$}{0,-4}
\draw (tete) node[coordinate] (tete3) {};

%% ETAT:
     \draw (tete3) node [name=etat,draw,circle,anchor=north]  {#7};
     \draw (etat) -- (tete3);
    \draw  (etat) -|  (tete2b);
\draw  (etat) -|  (tete1b);
%\end{tikzpicture}
}

\newcommand\MACHINEDETURINGTROISRUBANSENGLISH[9]{
%#1: texte1
%#2: position de la tête 1
%#3: texte2
%#4: position de la tête 2
%#5: texte3
%#6: position de la tête 3
%#7: état
%#8: nb de blancs à gauche 1
%#9: nb de cases totales
%\begin{tikzpicture}[node distance=-0.15mm]

\RUBANMACHINEDETURINGTETE{#1}{#2}{#8}{#9}{tape $1$}{0,0}
\draw (tete) node[coordinate] (tete1) {};
\draw (tete1) +(8,0) node[coordinate] (tete1b) {};
\draw (tete1) -| (tete1b) ;

\RUBANMACHINEDETURINGTETE{#3}{#4}{#8}{#9}{tape $2$}{0,-2}
\draw (tete) node[coordinate] (tete2) {};
\draw (tete2) +(-8.5,0) node[coordinate] (tete2b) {};
\draw (tete2) -| (tete2b) ;

\RUBANMACHINEDETURINGTETE{#5}{#6}{#8}{#9}{tape $3$}{0,-4}
\draw (tete) node[coordinate] (tete3) {};

%% ETAT:
     \draw (tete3) node [name=etat,draw,circle,anchor=north]  {#7};
     \draw (etat) -- (tete3);
    \draw  (etat) -|  (tete2b);
\draw  (etat) -|  (tete1b);
%\end{tikzpicture}
}



\newcommand\MACHINEDETURINGTROISRUBANStexte[9]{
%#1: texte1
%#2: position de la tête 1
%#3: texte2
%#4: position de la tête 2
%#5: texte3
%#6: position de la tête 3
%#7: état
%#8: nb de blancs à gauche 1
%#9: nb de cases totales
% \begin{tikzpicture}[node distance=-0.15mm]

\RUBANMACHINEDETURINGTETEtexte{#1}{#2}{#8}{#9}{\LANGUE{tape}{ruban} $1$}{0,0}
\draw (tete) node[coordinate] (tete1) {};
\draw (tete1) +(8,0) node[coordinate] (tete1b) {};
\draw (tete1) -| (tete1b) ;

\RUBANMACHINEDETURINGTETEtexte{#3}{#4}{#8}{#9}{\LANGUE{tape}{ruban} $2$}{0,-2}
\draw (tete) node[coordinate] (tete2) {};
\draw (tete2) +(-8.5,0) node[coordinate] (tete2b) {};
\draw (tete2) -| (tete2b) ;

\RUBANMACHINEDETURINGTETEtexte{#5}{#6}{#8}{#9}{\LANGUE{tape}{ruban} $3$}{0,-4}
\draw (tete) node[coordinate] (tete3) {};

%% ETAT:
     \draw (tete3) node [name=etat,draw,circle,anchor=north]  {#7};
     \draw (etat) -- (tete3);
    \draw  (etat) -|  (tete2b);
\draw  (etat) -|  (tete1b);
%\end{tikzpicture}
}


\newcommand\MACHINEDETURINGDEUXRUBANStexte[9]{
%#1: texte1  
%#2: position de la tête 1 
%#3: texte2 inutilisée
%#4: position de la tête 2 inutilisée
%#5: texte3
%#6: position de la tête 3
%#7: état
%#8: nb de blancs à gauche 1
%#9: nb de cases totales
% \begin{tikzpicture}[node distance=-0.15mm]

\RUBANMACHINEDETURINGTETEtexte{#1}{#2}{#8}{#9}{\LANGUE{tape}{ruban} $1$}{0,0}
\draw (tete) node[coordinate] (tete1) {};
\draw (tete1) +(8,0) node[coordinate] (tete1b) {};
\draw (tete1) -| (tete1b) ;

% \RUBANMACHINEDETURINGTETEtexte{#3}{#4}{#8}{#9}{\LANGUE{tape}{ruban} $2$}{0,-1.5}
% \draw (tete) node[coordinate] (tete2) {};
% \draw (tete2) +(-5.5,0) node[coordinate] (tete2b) {};
% \draw (tete2) -| (tete2b) ;

\RUBANMACHINEDETURINGTETEtexte{#5}{#6}{#8}{#9}{\LANGUE{tape}{ruban} $2$}{0,-1.5}
\draw (tete) node[coordinate] (tete3) {};

%% ETAT:
     \draw (tete3) node [name=etat,draw,circle,anchor=north]  {#7};
     \draw (etat) -- (tete3);
%   \draw  (etat) -|  (tete2b);
\draw  (etat) -|  (tete1b);
%\end{tikzpicture}
}


\def\argdix{2}
\def\argonze{25}
\newcommand\MACHINEDETURINGQUATRERUBANStextep[9]{
%#1: texte1
%#2: position de la tête 1
%#3: texte2
%#4: position de la tête 2
%#5: texte3
%#6: position de la tête 3
%#7: texte4
%#8: position de la tête 3
%#9: état
%#10: nb de blancs à gauche 1
%#11: nb de cases totales
% \begin{tikzpicture}[node distance=-0.15mm]

%% BUG sur #10 devenu \argdix
%% BUG SUR #11 devenu \argonze

\RUBANMACHINEDETURINGTETEtexte{#1}{#2}{\argdix}{\argonze}{\LANGUE{tape}{ruban} $1$}{0,0}
% \draw (tete) node[coordinate] (tete1) {};
% \draw (tete1) +(8,0) node[coordinate] (tete1b) {};
% \draw (tete1) -| (tete1b) ;

\RUBANMACHINEDETURINGTETEtexte{#3}{#4}{\argdix}{\argonze}{\LANGUE{tape}{ruban} $2$}{0,-1.5}
% \draw (tete) node[coordinate] (tete2) {};
% \draw (tete2) +(-5.5,0) node[coordinate] (tete2b) {};
% \draw (tete2) -| (tete2b) ;

\RUBANMACHINEDETURINGTETEtexte{#5}{#6}{\argdix}{\argonze}{\LANGUE{tape}{ruban} $3$}{0,-3.2}
\draw (tete) node[coordinate] (tete3) {};

\RUBANMACHINEDETURINGTETEtexte{#7}{#8}{\argdix}{\argonze}{\LANGUE{tape}{ruban} $4$}{0,-4.9}
\draw (tete) node[coordinate] (tete4) {};

%% ETAT:
     \draw (tete4) node [name=etat,draw,circle,anchor=north]  {#9};
%      \draw (etat) -- (tete3);
%     \draw  (etat) -|  (tete2b);
% \draw  (etat) -|  (tete1b);
%\end{tikzpicture}
}


%%%%%%%%%%%%%%%%%%%%%%%%%%%%%%%%%%%%%%%%%%%%%%%%%%%%%%%%%%%%%%%%%%%%%%%%%%%%
%                 MACHINES DE TURING (automate)
%%%%%%%%%%%%%%%%%%%%%%%%%%%%%%%%%%%%%%%%%%%%%%%%%%%%%%%%%%%%%%%%%%%%%%%%%%%%



\newcommand\dec{1ex}
\newcommand\myACCOLADE[4]{
\ACCOLADE{#1*\lcase+#4*\lcase-0.5*\lcase,0.5*\lcase+\dec}{#1*\lcase+#4*\lcase-0.5*\lcase+#3*\lcase,0.5*\lcase+\dec}{#2}
}


\newcommand\M{\Sigma}


\newcommand\sommet[2]{\node[state] (#1) #2 {$#1$};}
\newcommand\sommetinitial[2]{\node[state,initial] (#1) #2 {$#1$};}
\newcommand\sommetfinal[2]{\node[state,accepting] (#1) #2 {$#1$};}
\newcommand\gtransition[5]{\draw (#1) edge node {#2/#4 $#5$} (#3);}
\newcommand\gtransitionba[5]{\draw (#1) [loop above] edge node {#2/#4
    $#5$} (#3);}
\newcommand\gtransitionbb[5]{\draw (#1) [loop below] edge node {#2/#4
    $#5$} (#3);}
\newcommand\gtransitionbleft[5]{\draw (#1) [loop left] edge node {#2/#4
    $#5$} (#3);}
\newcommand\gtransitionbright[5]{\draw (#1) [loop right] edge node {#2/#4
    $#5$} (#3);}
\newcommand\gtransitionbl[5]{\draw (#1) [bend left] edge node {#2/#4
    $#5$} (#3);}
\newcommand\gtransitionbr[5]{\draw (#1) [bend right] edge node {#2/#4
    $#5$} (#3);}
\newcommand\gtransitionbadouble[9]{\draw (#1) [loop above] edge node {
\begin{tabular}{l}
#2/#4 $#5$ \\
  #7/#8 $#9$ \\
\end{tabular}
} (#3);}

\newcommand\SCALE{}

\newcommand\dessinmtp[2]{
\begin{center}
\begin{tikzpicture}[->,>=stealth',shorten >=1pt,auto,node distance=2.3cm,semithick,#2] 
#1
\end{tikzpicture}
\end{center}
}

\newcommand\dessinmt[1]{\dessinmtp{#1}{scale=1}}


\newcommand\ANIMATIONMT[1]{
\begin{overprint}
\begin{dessinp}{scale=0.52}
\foreach \av/\q/\b  [count=\xi]
in  
{#1}
{ 
\only<\xi| handout 0>{ 
 \StrLen{\av}[\Resintm] % ATTENTION IL FAUT CETTE SYNTAXE POUR FAIRE
                         % EQUIVALENT D UN EDEF
 \MACHINEDETURING{\av \b}{\Resintm}{$\q$}{5}{20}{}
}}
 \end{dessinp}
\end{overprint}
}

\newcommand\ANIMATIONMTd[1]{
\begin{overprint}
\begin{dessinp}{scale=0.52}
\foreach \av/\q/\b  [count=\xi]
in  
{#1}
{ 
\only<\xi| handout 0>{ 
 \StrLen{\av}[\Resintm] % ATTENTION IL FAUT CETTE SYNTAXE POUR FAIRE
                         % EQUIVALENT D UN EDEF
 \hspace{-1cm}\MACHINEDETURING{\av \b}{\Resintm}{$\q$}{5}{33}{}
}}
 \end{dessinp}
\end{overprint}
}

\def\MAXDET{20}

\newcommand\DIAGRAMMEESPACETEMPS[1]{
\draw node[coordinate] (tt) {};
\foreach \av/\q/\b  in  
{#1}
{ 

\RUBANMACHINEDETURING{\av{$\q$}\b}{0}{4}{\MAXDET}{}{tt};
\draw (tt) +(0,-\hcase) node[coordinate] (tt) {};
}
%\foreach \x in {0,...,19} {\draw (tt) + (\x*\lcase,0) node {$\vdots$};}
}



\newcommand\progun{
/q_0/0011, % initial
X/q_1/011,X0/q_1/11,X/q_2/0Y1,/q_2/X0Y1,
X/q_0/0Y1,XX/q_1/Y1,XXY/q_1/1, XX/q_2/YY, X/q_2/XYY,
XX/q_0/YY,XXY/q_3/Y,XXYY/q_3/B,XXYYB/q_4/B
}

\newcommand\DIAGRAMMEESPACETEMPSun{
\DIAGRAMMEESPACETEMPS
%% RECOPIE progun
{/q_0/0011, % initial
X/q_1/011,X0/q_1/11,X/q_2/0Y1,/q_2/X0Y1,
X/q_0/0Y1,XX/q_1/Y1,XXY/q_1/1, XX/q_2/YY, X/q_2/XYY,
XX/q_0/YY,XXY/q_3/Y,XXYY/q_3/B,XXYYB/q_4/B}
%% FIN RECOPIE
}

\newcommand\progdeux{
/q_0/0010, X/q_1/010, X0/q_1/10, X/q_2/0Y0,/q_2/X0Y0,
X/q_0/0Y0,XX/q_1/Y0,XXY/q_1/0,XXY0/q_1/%B,
}


\newcommand\DIAGRAMMEESPACETEMPSdeux{
\DIAGRAMMEESPACETEMPS
{
%% RECOPIE progdeux
/q_0/0010, X/q_1/010, X0/q_1/10, X/q_2/0Y0,/q_2/X0Y0,
X/q_0/0Y0,XX/q_1/Y0,XXY/q_1/0,XXY0/q_1/%B,
%% FIN RECOPIE
}
}

\newcommand\DIAGRAMMEESPACETEMPSdeuxvariante{
\def\memolcase{\lcase}
\def\memoMAXDET{20}
\def\MAXDET{15}
\def\lcase{6 ex}
\DIAGRAMMEESPACETEMPS
{
%% RECOPIE progdeux
/q_00/010, X/q_10/10, X0/q_11/0, X/q_20/Y0,/q_2X/0Y0,
X/q_00/Y0,XX/q_1Y/0,XXY/q_10/,XXY0/q_1B/%B,
%% FIN RECOPIE
}
\def\lcase{\memolcase}
\def\MAXDET{\memoMAXDET}
}


\newcommand\ANIMATIONun{
\ANIMATIONMT
%% RECOPIE progun
{/q_0/0011, % initial
X/q_1/011,X0/q_1/11,X/q_2/0Y1,/q_2/X0Y1,
X/q_0/0Y1,XX/q_1/Y1,XXY/q_1/1, XX/q_2/YY, X/q_2/XYY,
XX/q_0/YY,XXY/q_3/Y,XXYY/q_3/B,XXYYB/q_4/B}
%% FIN RECOPIE
}



\newcommand\DESSINALGO[3]{
% argument 1= ou
% argument 2= texte 
% argument 3= nom du noeud ex m
\draw node[draw,#1,minimum size=1cm] (#3) {#2};

% pour fitting box
\draw (#3.west) +(-0.2cm,0) node (#3c2) {};
\draw (#3.north) +(0,+0.2cm) node (#3c3) {};
\draw (#3.south) +(0,-0.2cm) node (#3c4) {};
}

\newcommand\DESSINALGOa[3]{
\DESSINALGO{#1}{#2}{#3}
\draw (#3.east) +(-0.1cm,0.2cm) node(#3a) {};
\draw[->] (#3a) -- ++(0.5cm,0) node (#3aa) {};
\draw (#3aa) node[anchor=west] (#3aaa) {\LANGUE{Accept}{Accepte}};
\draw (#3aaa.east) node(#3aaaa) {};
}

\newcommand\DESSINALGOe[4]{
\DESSINALGO{#1}{#2}{#3}
\draw (#3.east) +(-0.1cm,0.2cm) node(#3a) {};
\draw[->] (#3a) -- ++(0.5cm,0) node (#3aa) {};
\draw (#3aa) node[anchor=west] (#3aaa) {#4};
\draw (#3aaa.east) node(#3aaaa) {};
}


\newcommand\DESSINALGOar[3]{
\DESSINALGOa{#1}{#2}{#3}
\draw (#3.east) +(-0.1cm,-0.2cm) node(#3r) {};
\draw[->] (#3r) -- ++(0.5cm,0) node (#3rr) {};
\draw (#3rr) node[anchor=west] (#3rrr) {\LANGUE{Reject}{Refuse}};
\draw (#3rrr.east) node(#3rrrr) {};
}


\newcommand\WINDOWarg[7] {
%#1: position
%#2..7: contenu du tableau
\draw (#1) node[case]  {#2} \dcase;
\draw (#1) ++ (\lcase,0) node[case] {#3} \dcase;
\draw (#1) ++ (2*\lcase,0) node[case] {#4} \dcase;
\draw (#1) ++ (0,-\lcase) node[case] {#5} \dcase;
\draw (#1) ++ (\lcase,-\lcase) node[case] {#6} \dcase;
\draw (#1) ++ (2*\lcase,-\lcase) node[case] {#7} \dcase;
}

\newcommand\WINDOW[3]{
%#1: positioin
%#2: contenu
%#3: étiquette
\WINDOWarg{#1}
{\StrChar{#2}{1}}{\StrChar{#2}{2}}{\StrChar{#2}{3}}
{\StrChar{#2}{4}}{\StrChar{#2}{5}}{\StrChar{#2}{6}}
\draw (#1) +(-\lcase*1.2,-\lcase*0.5) node {#3};
}





%%%%%%%%%%%%%%%%%%%%%%%%%%%%%%%%%%%%%%%%%%%%%%%%%%%%%%%%%%%%%%%%%%%%%%%%%%%%
%                 LISTINGS
%%%%%%%%%%%%%%%%%%%%%%%%%%%%%%%%%%%%%%%%%%%%%%%%%%%%%%%%%%%%%%%%%%%%%%%%%%%%


% LISTINGS
\RequirePackage{listings}
 \lstset{language=Java,
 mathescape=true,
 showspaces=false,
 showstringspaces=false,
 frame=single,
morekeywords={recursive,then,div,function,integer,read,out,else,par,endpar,seq,endseq,and,not,read,write,repeat,until,undef,let,to,endfor,endif,endwhile,guess,in,parallel},
% basicstyle=\ttfamily\small,
% basewidth={1.1ex},
% keywordstyle=\bf\color{blue},
% %stringstyle=\bf\ttfamily\color{green},
% commentstyle=\bfseries\color{magenta},
 literate={:=}{{$\gets$ }}1 {<=}{{$\leq$}}2  {>=}{{$\geq$}}2
 {!=}{{$\neq$}}1 
 }
% \lstset{escapechar=\%}
\let\lst\lstinline
\def\beginlisting{\begin{lstlisting}}
\def\endlisting{\end{lstlisting}}  




\newenvironment{algo}[1]{
  \lstset{escapechar=\@}
  \def\input{$#1$}
  \def\out{out}
  \def\BEGIN{\textbf{read} \input }
 \def\BEGINN{\textbf{repeat} }
  \def\END{\textbf{until out!=undef}}
  \def\ENDD{\textbf{write \out}}
  \def\ENDT{\textbf{until $\ctl$ = $q_a$ or $\ctl$ = $q_r$}}
  \def\ENDR{\textbf{until $\ctl$ = $q_a$}}
  \def\ENDP{\textbf{until $\ctl$ = $q_a$}}
}{}
\newenvironment{algon}[1]{
  \lstset{numbers=left}
  \begin{algo}{#1}
}{\end{algo}}
\newcommand\FSM[3]{FSM(#1,\lst{#2},#3)}
\newcommand\FSMi[3]{FSM(#1,{#2},#3)}




%%%%%%%%%%%%%%%%%%%%%%%%%%%%%%%%%%%%%%%%%%%%%%%%%
%
%
%                      TRADUCTION EN COURS
%
%%%%%%%%%%%%%%%%%%%%%%%%%%%%%%%%%%%%%%%%%%%%%%%%%

\newcommand\ENTREPOT[1]{\NDLR{ENTREPOT ICI: #1}}
\newcommand\scite[2][]{\cite[#1]{#2}}
\newcommand\nospellcheck[1]{{#1}}
\newcommand\nd{nd}
\newcommand\NL{
 
}

\newenvironment{FRANCAIS}{
\small
\color{orange}
}
{\normalsize
\color{black}
}


\newenvironment{SEMIFRANCAIS}{
\small
\color{violet}
}
{\normalsize
\color{black}
}

\newcommand\QFRANCAIS[1]{
\begin{FRANCAIS}
#1
\end{FRANCAIS}
}


%%%%%%%%%%%%%%%%%%%%%%%%%%%%%%%%%%%%%%%%%%%%%%%%%
%
%
%                      Wikipedia
%
%%%%%%%%%%%%%%%%%%%%%%%%%%%%%%%%%%%%%%%%%%%%%%%%%

\providecommand\tightlist{}
\providecommand\Complex{\C}
\newenvironment{myalgie}{\begin{array}{lll}}{\end{array}}

\newcommand\nl{\ \\ }




%%%%%%%%%%%%%%%%%%%%%%%%%%%%%%%%%%%%%%%%%%%%%%%%%
%
%
%                      Pour citations précises
%
%%%%%%%%%%%%%%%%%%%%%%%%%%%%%%%%%%%%%%%%%%%%%%%%%



\newcommand\citeprecis[2]{{\cite[#1]{#2}}}
\newcommand\citeprecisinv[2]{\citeprecis{#2}{#1}}

\usepackage{csquotes}



%%%%%%%%%%%%%%%%%%%%%%%%%%%%%%%%%%%%%%%%%%%%%%%%%
%
%
%                      Travail sur couleurs
%
%   toutes les réponses pointent dans le même sens, et le schéma correspond exactement 
%. au profil perceptif d’un protanope** (déficit des cônes L / longueurs d’onde longues)
%  MAIS SEVERE
%
%%%%%%%%%%%%%%%%%%%%%%%%%%%%%%%%%%%%%%%%%%%%%%%%%


% ================================
% GROUPE 1 — Verts sombres
% ================================
\definecolor{deepgreen}{RGB}{0,160,40}
\definecolor{forestgreen}{RGB}{0,150,70}
\definecolor{mossgreen}{RGB}{0,140,90}
\definecolor{greenwave}{RGB}{0,150,90}
\definecolor{forestmist}{RGB}{0,160,90}
\definecolor{darkteal}{RGB}{0,130,70}

% ================================
% GROUPE 2 — Teal / verts bleutés
% ================================
\definecolor{seagreen}{RGB}{0,130,100}
\definecolor{freshteal}{RGB}{0,140,130}
\definecolor{greenteal}{RGB}{0,140,120}
\definecolor{blueteal}{RGB}{0,150,120}      % ancien "teal" (doublon corrigé)
\definecolor{lagoon}{RGB}{0,180,120}
\definecolor{greenmist}{RGB}{0,170,110}

% ================================
% GROUPE 3 — Turquoises / Aqua
% ================================
\definecolor{springgreen}{RGB}{0,190,100}
\definecolor{mintgreen}{RGB}{10,200,150}
\definecolor{oceanfoam}{RGB}{0,200,150}
\definecolor{aquagreen}{RGB}{30,210,170}
\definecolor{coastalaqua}{RGB}{40,210,180}
\definecolor{softaqua}{RGB}{40,210,190}

% ================================
% GROUPE 4 — Cyan / Aqua clairs
% ================================
\definecolor{emeraldteal}{RGB}{0,180,140}
\definecolor{deepcyan}{RGB}{0,180,160}
\definecolor{lightcyan}{RGB}{0,200,160}
\definecolor{aquamarine}{RGB}{0,200,160}
\definecolor{lightaqua}{RGB}{60,230,200}
\definecolor{mistblue}{RGB}{80,240,210}

% ================================
% GROUPE 5 — Bleus clairs
% ================================
\definecolor{paleaqua}{RGB}{80,200,190}
\definecolor{glacier}{RGB}{90,150,200}
\definecolor{skyblueA}{RGB}{60,150,220}      % ancien "skyblue"
\definecolor{lightsky}{RGB}{90,170,230}
\definecolor{iceblue}{RGB}{120,190,240}
\definecolor{frost}{RGB}{150,210,250}

% ================================
% GROUPE 6 — Bleus moyens
% ================================
\definecolor{steelblue}{RGB}{30,100,160}
\definecolor{cadetblue}{RGB}{70,120,180}
\definecolor{coolblue}{RGB}{110,170,220}
\definecolor{morningblue}{RGB}{140,200,240}
\definecolor{blueiceA}{RGB}{40,150,210}       % ancien "blueice" (doublon corrigé)
\definecolor{blueair}{RGB}{60,170,225}

% ================================
% GROUPE 7 — Azurs / Bleu saturé
% ================================
\definecolor{tealblue}{RGB}{0,110,130}
\definecolor{deepbluecyan}{RGB}{0,110,170}
\definecolor{azur}{RGB}{0,70,180}
\definecolor{brightazure}{RGB}{20,130,200}
\definecolor{lightazure}{RGB}{80,180,235}
\definecolor{skylight}{RGB}{100,200,245}

% ================================
% GROUPE 8 — Bleus sombres / violets
% ================================
\definecolor{nightblue}{RGB}{40,0,110}
\definecolor{slateblue}{RGB}{20,40,120}
\definecolor{violetA}{RGB}{60,50,160}         % ancien "violet" (distinct de violetB ci-dessous)
\definecolor{blueviolet}{RGB}{80,60,180}
\definecolor{purpleA}{RGB}{100,70,200}        % ancien "purple"
\definecolor{lightpurple}{RGB}{130,90,210}

% ================================
% GROUPE 9 — Violets saturés / indigos
% ================================
\definecolor{lilac}{RGB}{160,110,220}
\definecolor{indigoA}{RGB}{50,0,130}          % ancien "indigo"
\definecolor{deepindigo}{RGB}{70,10,150}
\definecolor{royalpurple}{RGB}{90,20,170}

% ================================
% GROUPE 10 — Fermeture cercle perceptif
% ================================
\definecolor{coolmint}{RGB}{0,130,70}
\definecolor{deepsea}{RGB}{0,130,70}

% ================================
% PALETTE DEUTÉRANOPE — noms uniques
% ================================
\definecolor{crimson}{RGB}{200,30,60}
\definecolor{brickred}{RGB}{180,40,40}
\definecolor{coral}{RGB}{220,90,70}
\definecolor{salmon}{RGB}{240,130,120}

\definecolor{magenta}{RGB}{200,40,150}
\definecolor{fuchsia}{RGB}{220,60,170}
\definecolor{violetB}{RGB}{140,70,190}        % ancien "violet"
\definecolor{indigoB}{RGB}{90,50,180}         % ancien "indigo"

\definecolor{royalblueA}{RGB}{50,80,200}      % ancien "royalblue"
\definecolor{azure}{RGB}{30,120,230}
\definecolor{skyblueB}{RGB}{80,160,230}       % second "skyblue"
\definecolor{tealB}{RGB}{0,150,170}           % second "teal"

\definecolor{marine}{RGB}{0,80,140}




%% Alternative

% === Générateur automatique fill/text pour protanopes ===
% Usage : \ProtaColorPair{mistblue}
% Suppose que mistblue est déjà défini par \definecolor{mistblue}{RGB}{...}

\newcommand{\ProtaColorPair}[1]{%
  % fill = couleur originale
  \colorlet{#1fill}{#1}%
  % text = version assombrie idéalement lisible
  \colorlet{#1text}{#1!80!black}%
}

%Utilisation
%\ProtaColorPair{mistblue}

%\textcolor{mistbluefill}{Couleur claire} \quad
%\textcolor{mistbluetext}{Couleur pour texte}

	\foreach \i/\name in {
	  1/deepgreen,
	  2/greenwave,
	  3/lagoon,
	  4/oceanfoam,
	  5/softaqua,
	  6/mistblue,
 	 7/lightsky,
	  8/blueiceA,   % <-- corrigé
	  9/deepcyan,
	  10/azur,
	  11/blueviolet,
	  12/royalpurple}
	  {\ProtaColorPair{\name}}
	  


\newcommand\ROUEPROTANOPEDOUZE{
	% Roue chromatique protanope — 12 couleurs
	\begin{tikzpicture}[scale=3]

	\foreach \i/\name/\nametext in {
	  1/deepgreen/deepgreen,
	  2/greenwave/greenwave,
	  3/lagoon/lagoontext,
	  4/oceanfoam/oceanfoam,
	  5/softaqua/softaquatext,
	  6/mistblue/mistbluetext,
 	 7/lightsky/lightskytext,
	  8/blueiceA/blueiceA,   % <-- corrigé
	  9/deepcyan/blueiceA,
	  10/azur/azur,
	  11/blueviolet/blueviolet,
	  12/royalpurple/royalpurple}
	{
	  % secteur coloré
 	 \filldraw[fill=\name, draw=black, line width=0.2pt]
	    ({90-30*\i}:1)
	    -- ({90-30*(\i-1)}:1)
	    -- (0,0) -- cycle;

 	 % label dans la bonne couleur
	  \node at ({90-30*(\i-0.5)}:1.25)
	    {\textcolor{\nametext}{\small \nametext}};
	}

	\end{tikzpicture}
}


\newcommand\ROUEPROTANOPEVINGT{
\begin{tikzpicture}[scale=3]

\foreach \i/\name/\nametext in {
  1/deepgreen/deepgreen,
  2/greenwave/greenwave,
  3/lagoon/lagoon,
  4/oceanfoam/oceanfoam,
  5/softaqua/softaquatext,
  6/mistblue/mistbluetext,
  7/lightsky/lightsky,
  8/blueiceA/blueiceA,      % corrigé
  9/deepcyan/deepcyan,
  10/azur/azur,
  11/brightazure/brightazure,
  12/skyblueA/skyblueA,     % corrigé
  13/coolblue/coolblue,
  14/glacier/glacier,
  15/marine/marine,
  16/violetA/violetA,      % corrigé
  17/blueviolet/blueviolet,
  18/purpleA/purpleA,      % corrigé
  19/royalpurple/royalpurple,
  20/indigoA/indigoA}      % corrigé
{
  % Secteur coloré
  \filldraw[fill=\name, draw=black, line width=0.2pt]
    ({90-18*\i}:1)
    -- ({90-18*(\i-1)}:1)
    -- (0,0) -- cycle;

  % Étiquette
  \node at ({90-18*(\i-0.5)}:1.25)
    {\textcolor{\nametext}{\small \nametext}};
}

\end{tikzpicture}
}

\newcommand\ROUEDEUTERANOPEDOUZE{
% Roue chromatique deutéranope — 12 couleurs
\begin{tikzpicture}[scale=3]

\foreach \i/\name in {
  1/crimson,
  2/brickred,
  3/coral,
  4/salmon,
  5/magenta,
  6/fuchsia,
  7/violetB,       % corrigé
  8/indigoB,       % corrigé
  9/royalblueA,    % corrigé
  10/azure,
  11/skyblueB,     % corrigé
  12/tealB}        % corrigé
{
  % Secteur coloré
  \filldraw[fill=\name, draw=black, line width=0.2pt]
    ({90-30*\i}:1)
    -- ({90-30*(\i-1)}:1)
    -- (0,0) -- cycle;

  % Label dans la couleur
  \node at ({90-30*(\i-0.5)}:1.25)
    {\textcolor{\name}{\small \name}};
}

\end{tikzpicture}
}



% === Macro intelligente pour foncer une couleur de façon lisible pour protanope ===
% Usage : \protadarken{mistblue}{40}  -> crée mistblue@dark
% Le 40 signifie : 40% de mélange avec du noir

\newcommand{\DarkenForText}[2]{%
  % #1 = nom de la couleur
  % #2 = pourcentage (40 conseillé si tu ne sais pas)
  %
  % On crée une nouvelle couleur #1@dark
  \colorlet{#1text}{black!#2!#1}%
}





%%%%

% Groupe 1
\DarkenForText{deepgreen}{20}
\DarkenForText{forestgreen}{20}
\DarkenForText{mossgreen}{20}
\DarkenForText{greenwave}{20}
\DarkenForText{forestmist}{20}
\DarkenForText{darkteal}{20}

% Groupe 2
\DarkenForText{seagreen}{20}
\DarkenForText{freshteal}{20}
\DarkenForText{greenteal}{20}
\DarkenForText{blueteal}{20}
\DarkenForText{lagoon}{20}
\DarkenForText{greenmist}{20}

% Groupe 3
\DarkenForText{springgreen}{20}
\DarkenForText{mintgreen}{20}
\DarkenForText{oceanfoam}{20}
\DarkenForText{aquagreen}{20}
\DarkenForText{coastalaqua}{20}
\DarkenForText{softaqua}{20}

% Groupe 4
\DarkenForText{emeraldteal}{20}
\DarkenForText{deepcyan}{20}
\DarkenForText{lightcyan}{20}
\DarkenForText{aquamarine}{20}
\DarkenForText{lightaqua}{20}
\DarkenForText{mistblue}{30}

% Groupe 5
\DarkenForText{paleaqua}{20}
\DarkenForText{glacier}{20}
\DarkenForText{skyblueA}{20}
\DarkenForText{lightsky}{20}
\DarkenForText{iceblue}{20}
\DarkenForText{frost}{20}

% Groupe 6
\DarkenForText{steelblue}{20}
\DarkenForText{cadetblue}{20}
\DarkenForText{coolblue}{20}
\DarkenForText{morningblue}{20}
\DarkenForText{blueiceA}{20}
\DarkenForText{blueair}{20}

% Groupe 7
\DarkenForText{tealblue}{20}
\DarkenForText{deepbluecyan}{20}
\DarkenForText{azur}{20}
\DarkenForText{brightazure}{20}
\DarkenForText{lightazure}{20}
\DarkenForText{skylight}{20}

% Groupe 8
\DarkenForText{nightblue}{20}
\DarkenForText{slateblue}{20}
\DarkenForText{violetA}{20}
\DarkenForText{blueviolet}{20}
\DarkenForText{purpleA}{20}
\DarkenForText{lightpurple}{20}

% Groupe 9
\DarkenForText{lilac}{20}
\DarkenForText{indigoA}{20}
\DarkenForText{deepindigo}{20}
\DarkenForText{royalpurple}{20}

% Groupe 10
\DarkenForText{coolmint}{20}
\DarkenForText{deepsea}{20}

% Palette deutéranope
\DarkenForText{crimson}{20}
\DarkenForText{brickred}{20}
\DarkenForText{coral}{20}
\DarkenForText{salmon}{20}

\DarkenForText{magenta}{20}
\DarkenForText{fuchsia}{20}
\DarkenForText{violetB}{20}
\DarkenForText{indigoB}{20}

\DarkenForText{royalblueA}{20}
\DarkenForText{azure}{20}
\DarkenForText{skyblueB}{20}
\DarkenForText{tealB}{20}

% Couleur supplémentaire
\DarkenForText{marine}{20}

%%%%%%%%%%%%%%%%%%%%%
%
%.   apx proof
%
%%%%%%%%%%%%%%%%%%%%%



% do
\newcommand\PREUVESENAPPENDIXCOMPATIBLE{
	\newenvironment{appendixproof}{{\bf \LANGUE{Proof}{Démonstration}}:}{\hfill $\Box$}
	\newtheorem{theoremrep}[theorem]{Theorem}
	\newtheorem{lemmarep}[theorem]{Lemma}
	\newtheorem{propositionrep}[theorem]{Proposition}
	\newtheorem{corollaryrep}[theorem]{Corollary}
	\newtheorem{summaryrep}[theorem]{Summary}	
}

\providecommand\apxparameter{}

%% proof inline
%\providecommand\apxparameter{appendix=inline}
%% proof in appendix
%\providecommand\apxparameter{}


\apxproofmode{
\usepackage[\apxparameter]{apxproof}
\theoremstyle{plain}

\newtheoremrep{theorem}{Theorem}
%\newtheoremrep{theoremconjecture}[theorem]{Theorem-To-Be-Proved=Conjecture}
%\newtheoremrep{lemmaconjecture}[theorem]{Conjecture-To-Be-Proved=Conjecture}
\newtheoremrep{claim}[theorem]{Theorem}
\newtheoremrep{proposition}[theorem]{Proposition}
\newtheoremrep{lemma}[theorem]{Lemma}
\newtheoremrep{corollary}[theorem]{Corollary}
%%bug mars 26: \newtheoremrep{summary}[theorem]{Summary}
% Du coup:
\newtheoremrep{resume}[theorem]{Summary}
%bug mars 26: \newtheoremrep{openproblem}[theorem]{Open Problem}
%\newtheoremrep{claim}[theorem]{Claim}
}


%\apxproofmodesubsection{
%%apx proof, pour que compteur correct
%\makeatletter
%\AtBeginDocument{%
%\let\axp@section\subsection
%  \let\axp@oldsection\subsection
%}
%\makeatother
%\renewcommand\appendixprelim{\section{Missing proofs}}
%}


\makeatletter
\apxproofmodesubsection{
\let\axp@section\subsection
%\let\axp@oldsection\subsection
\renewcommand\appendixprelim{\section{Proofs from main documents}}
}
\makeatother


\makeatletter
\newcommand\appendixapxproofsubsection{
\appendix
\let\apx@orig@appendix\appendix
\renewcommand{\appendix}{}
}
\makeatother


%%% Pour checker/vérifier que bon


\newcommand\docheck[1]{\textcolor{check}{#1}}

\usepackage{xcolor}
%\definecolor{chose}{RGB}{0,110,50}
%\definecolor{forestgreen}{RGB}{0,150,70}
%\definecolor{ctruc}{RGB}{164,164,164}
%\definecolor{oceanfoam}{RGB}{0,200,150}
\definecolor{check}{RGB}{160,160,0}




. Retirées automatiquement en mode lipics.
\providecommand\NODECORATION[1]{\ifnotSLIDE{#1}}
\lipicsmode{\renewcommand\NODECORATION[1]{}}

\providecommand\PASSIAPXPROOF[1]{#1}
\apxproofmode{\renewcommand\PASSIAPXPROOF[1]{}}



%% end gestion des switchs

\def\MACROSPOINTEXOUVERT{}


\providecommand\ifnotTDEXAM[1]{#1}
\providecommand\ifnotSLIDE[1]{#1}


  
% pour preuves en appendix
\providecommand\PREUVESENAPPENDIX[1]{}



%%%% STYLES: APRES VOLS


\NODECORATION{
\usepackage{fourier}
}

%% Repli si fourier (et donc fourier-orns, qui définit \danger) n'est pas
%% chargé ; doit rester APRÈS le bloc ci-dessus, sinon fourier-orns échoue
%% avec « \danger already defined ».
\providecommand\danger{}

%%%% Box around environment
\usepackage[framemethod=tikz]{mdframed}

\newcommand\ENVCOLOR[1]{#1}
\newcommand\BEGINTEMPENVCOLOR[1]{\let\envcolorwas\ENVCOLOR\renewcommand\ENVCOLOR[1]{#1}}
\newcommand\ENDTEMPENVCOLOR{\let\ENVCOLOR\envcolorwas}


\NODECORATION{
\surroundwithmdframed[hidealllines=true,backgroundcolor=\ENVCOLOR{blue!10},roundcorner=8pt]{exercice}
\surroundwithmdframed[hidealllines=true,backgroundcolor=\ENVCOLOR{blue!10},roundcorner=8pt]{exercise}
\surroundwithmdframed[hidealllines=true,backgroundcolor=\ENVCOLOR{blue!10},roundcorner=8pt]{exercicee}
\surroundwithmdframed[hidealllines=true,backgroundcolor=\ENVCOLOR{blue!10},roundcorner=8pt]{exerciced}
\surroundwithmdframed[hidealllines=true,backgroundcolor=\ENVCOLOR{blue!10},roundcorner=8pt]{exercicec}
\surroundwithmdframed[hidealllines=true,backgroundcolor=\ENVCOLOR{blue!10},roundcorner=8pt]{exercicedc}
\surroundwithmdframed[hidealllines=true,backgroundcolor=\ENVCOLOR{orange!20},roundcorner=8pt]{example}
\surroundwithmdframed[hidealllines=true,backgroundcolor=\ENVCOLOR{orange!20},roundcorner=8pt]{exemple}
\surroundwithmdframed[hidealllines=true,backgroundcolor=\ENVCOLOR{orange!10},roundcorner=8pt]{remark}
\ifnotSLIDE{\surroundwithmdframed[backgroundcolor=\ENVCOLOR{green!15},roundcorner=8pt]{definition}}
\surroundwithmdframed[hidealllines=true,backgroundcolor=\ENVCOLOR{gray!35}]{proposition}
\ifnotSLIDE{\surroundwithmdframed[backgroundcolor=\ENVCOLOR{gray!35},roundcorner=8pt]{theorem}}
\surroundwithmdframed[hidealllines=true,backgroundcolor=\ENVCOLOR{gray!35}]{lemma}
\surroundwithmdframed[hidealllines=true,backgroundcolor=\ENVCOLOR{gray!35}]{corollary}


\surroundwithmdframed[backgroundcolor=red!35,roundcorner=8pt]{theoremaprouver}
\surroundwithmdframed[backgroundcolor=red!35,roundcorner=8pt]{conjectureaprouver}
}

%%% tcolorbox

\usepackage{tcolorbox}
\tcbuselibrary{skins}


  %%. Commande magique
%% \tcolorboxenvironment{⟨name⟩}{⟨options⟩}: attach options de tcolorbox à environnement
%% alternative: je cree des boites moi meme: voici des exemples


\newtcolorbox{maboiteconfiance}[2][]{colback=red!5!white, colframe=red!75!black,fonttitle=\bfseries, colbacktitle=red!85!black,enhanced,
attach boxed title to top center={yshift=-2mm},
  title={#2},#1}
  
\newtcolorbox{maboiteresume}[2][]{colback=red!5!white, colframe=red!75!black,fonttitle=\bfseries, colbacktitle=red!85!black,enhanced,
attach boxed title to top center={yshift=-2mm},
  title={#2},#1}
  \newtcolorbox{maboiteresumeperso}[2][]{colback=red!5!white, colframe=red!75!black,fonttitle=\bfseries, colbacktitle=red!85!black,enhanced,
attach boxed title to top center={yshift=-2mm},
  title={#2},#1}
  \newtcolorbox{maboitecommentaire}[2][]{colback=red!5!white, colframe=red!75!black,fonttitle=\bfseries, colbacktitle=red!85!black,enhanced,
attach boxed title to top center={yshift=-2mm},
  title={#2},#1}
\newtcolorbox{maboitetruc}[2][]{colback=red!5!white, colframe=red!75!black,fonttitle=\bfseries, colbacktitle=red!85!black,enhanced,
attach boxed title to top left={yshift=-2mm},
  title={#2},#1}

%% Boîte des mises en évidence CARE. Couleurs pensées pour la
%% protanopie : l'axe bleu-jaune reste discriminant (jaune = \IMPORTANT
%% d'Olivier, bleu = CARE), et la distinction ne repose pas que sur la
%% couleur (cadre + bandeau étiqueté, lisibles même en niveaux de gris).
\newtcolorbox{maboiteimportantcare}[2][]{colback=cyan!10!white, colframe=blue!60!black,fonttitle=\bfseries, colbacktitle=blue!70!black,enhanced,
attach boxed title to top center={yshift=-2mm},
  title={#2},#1}
  


%%%% FIN STYLES



%%%%%%%%%%%%%%%%%%%%%%%%%%%%%%%%%%%%%%%%%%%%%%%%%
%%%%%%%%%%%%%%%%%%%%%%%%%%%%%%%%%%%%%%%%%%%%%%%%%
%
%
%                      Version ENSEIGNEMENT 2020
%
%
%%%%%%%%%%%%%%%%%%%%%%%%%%%%%%%%%%%%%%%%%%%%%%%%%
%%%%%%%%%%%%%%%%%%%%%%%%%%%%%%%%%%%%%%%%%%%%%%%%%

%english, french
\ifdefined\CESTDUFRANCAIS
	\providecommand\LANGUE[2]{#2}
	\usepackage[french]{babel}
\else
	\providecommand\LANGUE[2]{#1}
\fi
\providecommand\LANGUEinv[2]{\LANGUE{#2}{#1}}

%%%%%%%%%%%%%%%%%%%%%%%%%%%%%%%%%%%%%%%%%%%%%%%%%
%
%
%                      Packages
%
%
%%%%%%%%%%%%%%%%%%%%%%%%%%%%%%%%%%%%%%%%%%%%%%%%%


%%% SURLIGNEUR
\providecommand\ifnotmodeDIFFUSIONFINALE[1]{
%% 2021: Comprend pas mais assure
\ifdefined\SAFEMODE
\else
#1
\fi
%#1
}
\providecommand\ifmodeDIFFUSIONFINALE[1]{
\ifdefined\SAFEMODE
#1
\else
\fi
}
\newcommand\SURLIGNE[1]{
\ifnotmodeDIFFUSIONFINALE{
\BEGINTEMPENVCOLOR{blue!30}
{                        \colorbox{blue!30}{

\noindent \begin{minipage}{\dimexpr\textwidth-2\fboxsep\relax}
#1
\end{minipage}
}
}
\ENDTEMPENVCOLOR}
\ifmodeDIFFUSIONFINALE{
\noindent \begin{minipage}{\textwidth}
#1
\end{minipage}
}
}


\newenvironment{confiance}{\begin{maboiteresume}[colback=brown!50]{Résumé ici}}{\end{maboiteresume}}
\newcommand\CONFIANCE[2][Confiance ici]{
\begin{maboiteconfiance}[colback=yellow!50]{#1}
\textbf{ATTENTION : #2}
\end{maboiteconfiance}
}

\newenvironment{resume}{\begin{maboiteresume}[colback=brown!50]{Résumé ici}}{\end{maboiteresume}}
\newcommand\RESUME[2][Résumé ici]{
\begin{maboiteresume}[colback=brown!50]{#1}
#2
\end{maboiteresume}
}

\newenvironment{resumeperso}{\begin{maboiteresumeperso}[colback=brown!50]{Résumé perso ici}}{\end{maboiteresumeperso}}
\newcommand\RESUMEPERSO[2][Résumé perso  ici]{
\begin{maboiteresumeperso}[colback=brown!50]{#1}
#2
\end{maboiteresumeperso}
%
%
%\todo[inline, color=brown]{résumé: #1}
%\colorbox{brown}{
%\noindent \begin{minipage}{\textwidth}
%{
%\bf 
%#2
%}
%\end{minipage}
%}
}

\newcommand\DANSLAMARGE[1]{
\ifnotmodeDIFFUSIONFINALE{
\marginpar{\textcolor{brown}{#1}}}
}

\newenvironment{commentaireperso}{\begin{maboitecommentaire}[colback=brown!30]{Commentaire perso} Commentaire perso:}{\end{maboitecommentaire}}
\newcommand\COMMENTAIREPERSO[2][Commentaire perso]{
\begin{maboitecommentaire}[colback=brown!30]{#1}
#2
\end{maboitecommentaire}
}


\newcommand\SURSURLIGNE[2][truc surligné ici]{
%\todo[inline, color=orange]{surligné: #1}
\ifnotmodeDIFFUSIONFINALE{
\BEGINTEMPENVCOLOR{blue!60}
	                 \colorbox{blue!60}{
\noindent \begin{minipage}{\dimexpr\textwidth-2\fboxsep\relax}
{
\bf 
#2
}
\end{minipage}
}
\ENDTEMPENVCOLOR
}
\ifmodeDIFFUSIONFINALE{#2}
}

\newcommand\IMPORTANT[1]{
\ifnotmodeDIFFUSIONFINALE{
\BEGINTEMPENVCOLOR{yellow!100}
		        \colorbox{yellow!100}{

\noindent \begin{minipage}{\dimexpr\textwidth-2\fboxsep\relax}
{ 
#1
}
\end{minipage}
}
\ENDTEMPENVCOLOR
}
\ifmodeDIFFUSIONFINALE{#1}
}

%% Pendant de \IMPORTANT, réservé aux mises en évidence CARE ;
%% \IMPORTANT reste réservé aux annotations d'Olivier. Comme \IMPORTANT,
%% redevient du texte simple en mode DIFFUSIONFINALE.
\newcommand\IMPORTANTCARE[2][Important (CARE)]{
\ifnotmodeDIFFUSIONFINALE{
\begin{maboiteimportantcare}{#1}
#2
\end{maboiteimportantcare}
}
\ifmodeDIFFUSIONFINALE{#2}
}


\newcommand\QUESTION[1]{
\ifnotmodeDIFFUSIONFINALE{
\BEGINTEMPENVCOLOR{oceanfoam!100}
		        \colorbox{oceanfoam!100}{

\noindent \begin{minipage}{\dimexpr\textwidth-2\fboxsep\relax}
{ 
#1
}
\end{minipage}
}
\ENDTEMPENVCOLOR
}
\ifmodeDIFFUSIONFINALE{#1}
}




\newcommand\SOUSLIGNE[1]{
\colorbox{lightgray}{

\noindent \begin{minipage}{\dimexpr\textwidth-2\fboxsep\relax}
{ \color{gray}
#1
}
\end{minipage}
}
}



%% Gestion des cross-references cross-referencing
\usepackage{xr-hyper}
\usepackage{hyperref}

%%
\usepackage{versions}
\usepackage{amsmath,amssymb,mathrsfs,amsfonts}
%\usepackage{mathabx}
\usepackage{listings}
\usepackage{url}
\usepackage{versions}
\usepackage{color}
\usepackage[colorinlistoftodos]{todonotes}
%\usepackage{todonotes}
\usepackage{proof}
\usepackage{xstring}

%% Index
\usepackage{makeidx}
\makeindex
\providecommand\DOINDEXPRINT{}
\providecommand\SIGNALEMOTNOUV[1]{}
\ifdefined\INDEXPRINT
	\ifnotSAFEMODE{
%	\usepackage{showidx}
%	\let\oldindex\index
%	\renewcommand{\index}[1]{\oldindex{#1}\marginpar{LA: \tiny#1}}
	\renewcommand\DOINDEXPRINT{
	      	\let\oldindex\index
		\renewcommand{\index}[1]{\oldindex{##1}\marginpar{IDXLA: \tiny##1}}
		\renewcommand\SIGNALEMOTNOUV[1]{\marginpar{MNV: \small\textbf{##1}}}
	}
	}
\fi

% === gestion des accents


%soit iso latin1
%\usepackage[cyr]{aeguill}
%soit utf8
\usepackage[utf8]{inputenc}
\usepackage[T1]{fontenc}


%%%  Pour citation \Cref (cref, ref)	

\usepackage{cleveref}


%%%%%%%%%%%%%%%%%%%%%%%%%%%%%%%%%%%%%%%%%%%%%%%%%
%
%
%                      Gestion des accents spéciaux.
%
%
%%%%%%%%%%%%%%%%%%%%%%%%%%%%%%%%%%%%%%%%%%%%%%%%%


%% MOI A LA MAIN:
\usepackage{bbm}

\DeclareUnicodeCharacter{2212}{-}
\DeclareUnicodeCharacter{2297}{\ensuremath{\otimes}}       
\DeclareUnicodeCharacter{2A33}{\ensuremath{\smashtimes}}
\DeclareUnicodeCharacter{1D56F}{\ensuremath{\mathfrak{D}}}  
\DeclareUnicodeCharacter{1D40E}{\ensuremath{\mathbf{O}}}  
\DeclareUnicodeCharacter{1D427}{\ensuremath{\mathbf{n}}}  
\DeclareUnicodeCharacter{3008}{\ensuremath{\langle}} % 〈 
\DeclareUnicodeCharacter{3009}{\ensuremath{\rangle}} %  〉

% (possible de voir: http://tug.ctan.org/info/symbols/comprehensive/symbols-letter.pdf)

%% SOURCE:  https://github.com/agda/agda/blob/master/doc/user-manual/unicode-symbols-sphinx.tex.txt

%%%%%%%%%%%%%%%%%%%%%%%%%%%%%%%%%%%%%%%%%%%%%%%%%%%%%%%%%%%%%%%%%%%%%%%%%%%%%%
% Please keep the below list ordered by code points!

\DeclareUnicodeCharacter{00AC}{\ensuremath{\neg}}                     % Symbol ¬
\DeclareUnicodeCharacter{00B1}{\ensuremath{\pm}}                      % Symbol ±
\DeclareUnicodeCharacter{00B2}{\ensuremath{^2}}                       % Symbol ²
\DeclareUnicodeCharacter{00B3}{\ensuremath{^3}}                       % Symbol ³
\DeclareUnicodeCharacter{00B7}{\ensuremath{\cdot}}                    % Symbol ·
\DeclareUnicodeCharacter{00B9}{\ensuremath{^1}}                       % Symbol ¹
\DeclareUnicodeCharacter{00D7}{\ensuremath{\times}}                   % Symbol ×
\DeclareUnicodeCharacter{02B0}{\ensuremath{^h}}                       % Symbol ʰ
\DeclareUnicodeCharacter{02B3}{\ensuremath{^r}}                       % Symbol ʳ
\DeclareUnicodeCharacter{02E2}{\ensuremath{^s}}                       % Symbol ˢ
\DeclareUnicodeCharacter{0302}{\ensuremath{\hat{\phantom{x}}}}        % Symbol ̂
\DeclareUnicodeCharacter{0332}{\ensuremath{\underline{\phantom{x}}}}  % Symbol ̲
\DeclareUnicodeCharacter{0308}{\ensuremath{\ddot{\phantom{x}}}}       % Symbol ̈
\DeclareUnicodeCharacter{0393}{\ensuremath{\Gamma}}                   % Symbol Γ
\DeclareUnicodeCharacter{0394}{\ensuremath{\Delta}}                   % Symbol Δ
\DeclareUnicodeCharacter{0398}{\ensuremath{\Theta}}                   % Symbol Θ
\DeclareUnicodeCharacter{039B}{\ensuremath{\Lambda}}                  % Symbol Λ
\DeclareUnicodeCharacter{039E}{\ensuremath{\Xi}}                      % Symbol Ξ
\DeclareUnicodeCharacter{03A0}{\ensuremath{\Pi}}                      % Symbol Π
\DeclareUnicodeCharacter{03A3}{\ensuremath{\Sigma}}                   % Symbol Σ


% There is not `\Tau` command in LaTeX, so we use `\mathrm{T}`. See
% http://tex.stackexchange.com/questions/1368/why-can-i-only-use-some-capital-greek-letters-inside-my-equations.
\DeclareUnicodeCharacter{03A4}{\ensuremath{\mathrm{T}}}  % Symbol Τ

\DeclareUnicodeCharacter{03A5}{\ensuremath{\Upsilon}}   % Symbol Υ
\DeclareUnicodeCharacter{03A6}{\ensuremath{\Phi}}       % Symbol Φ
\DeclareUnicodeCharacter{03A8}{\ensuremath{\Psi}}       % Symbol Ψ
\DeclareUnicodeCharacter{03A9}{\ensuremath{\Omega}}     % Symbol Ω
\DeclareUnicodeCharacter{03B1}{\ensuremath{\alpha}}     % Symbol α
\DeclareUnicodeCharacter{03B2}{\ensuremath{\beta}}      % Symbol β
\DeclareUnicodeCharacter{03B3}{\ensuremath{\gamma}}     % Symbol γ
\DeclareUnicodeCharacter{03B4}{\ensuremath{\delta}}     % Symbol δ
\DeclareUnicodeCharacter{03B5}{\ensuremath{\epsilon}}   % Symbol ε
\DeclareUnicodeCharacter{03B6}{\ensuremath{\zeta}}      % Symbol ζ
\DeclareUnicodeCharacter{03B7}{\ensuremath{\eta}}       % Symbol η
\DeclareUnicodeCharacter{03B8}{\ensuremath{\theta}}     % Symbol θ
\DeclareUnicodeCharacter{03B9}{\ensuremath{\iota}}      % Symbol ι
\DeclareUnicodeCharacter{03BA}{\ensuremath{\kappa}}     % Symbol κ
\DeclareUnicodeCharacter{03BB}{\ensuremath{\lambda}}    % Symbol λ
\DeclareUnicodeCharacter{03BC}{\ensuremath{\mu}}        % Symbol μ
\DeclareUnicodeCharacter{03BD}{\ensuremath{\nu}}        % Symbol ν
\DeclareUnicodeCharacter{03BE}{\ensuremath{\xi}}        % Symbol ξ
\DeclareUnicodeCharacter{03C0}{\ensuremath{\pi}}        % Symbol π
\DeclareUnicodeCharacter{03C1}{\ensuremath{\rho}}       % Symbol ρ
\DeclareUnicodeCharacter{03C2}{\ensuremath{\varsigma}}  % Symbol ς
\DeclareUnicodeCharacter{03C3}{\ensuremath{\sigma}}     % Symbol σ
\DeclareUnicodeCharacter{03C4}{\ensuremath{\tau}}       % Symbol τ
\DeclareUnicodeCharacter{03C5}{\ensuremath{\upsilon}}   % Symbol υ
\DeclareUnicodeCharacter{03C6}{\ensuremath{\varphi}}    % Symbol φ
\DeclareUnicodeCharacter{03C7}{\ensuremath{\chi}}       % Symbol χ
\DeclareUnicodeCharacter{03C8}{\ensuremath{\psi}}       % Symbol ψ
\DeclareUnicodeCharacter{03C9}{\ensuremath{\omega}}     % Symbol ω
\DeclareUnicodeCharacter{03D1}{\ensuremath{\vartheta}}  % Symbol ϑ

\DeclareUnicodeCharacter{0432}{\ensuremath{\mathrm{PDF\;TODO}}}  % Symbol в
\DeclareUnicodeCharacter{0435}{\ensuremath{\mathrm{PDF\;TODO}}}  % Symbol е
\DeclareUnicodeCharacter{0438}{\ensuremath{\mathrm{PDF\;TODO}}}  % Symbol и
\DeclareUnicodeCharacter{043C}{\ensuremath{\mathrm{PDF\;TODO}}}  % Symbol м
\DeclareUnicodeCharacter{0440}{\ensuremath{\mathrm{PDF\;TODO}}}  % Symbol р
\DeclareUnicodeCharacter{0442}{\ensuremath{\mathrm{PDF\;TODO}}}  % Symbol т
\DeclareUnicodeCharacter{041F}{\ensuremath{\mathrm{PDF\;TODO}}}  % Symbol П
\DeclareUnicodeCharacter{0BA8}{\ensuremath{\mathrm{PDF\;TODO}}}  % Symbol ந
\DeclareUnicodeCharacter{0BBF}{\ensuremath{\mathrm{PDF\;TODO}}}  % Symbol  ி
\DeclareUnicodeCharacter{1100}{\ensuremath{\mathrm{PDF\;TODO}}}  % Symbol ᄀ
\DeclareUnicodeCharacter{11F9}{\ensuremath{\mathrm{PDF\;TODO}}}  % Symbol ᇹ

\DeclareUnicodeCharacter{1D48}{\ensuremath{^d}}  % Symbol ᵈ
\DeclareUnicodeCharacter{1D50}{\ensuremath{^m}}  % Symbol ᵐ
\DeclareUnicodeCharacter{1D56}{\ensuremath{^p}}  % Symbol ᵖ
\DeclareUnicodeCharacter{1D62}{\ensuremath{_i}}  % Symbol ᵢ
\DeclareUnicodeCharacter{1D63}{\ensuremath{_r}}  % Symbol ᵣ
\DeclareUnicodeCharacter{1D64}{\ensuremath{_u}}  % Symbol ᵤ
\DeclareUnicodeCharacter{1D9C}{\ensuremath{^c}}  % Symbol ᶜ

% The command `\text` requires `\usepackage{amsmath}` and the command
% `\textasciiacute` requires `\usepackage{textcomp}`.

\DeclareUnicodeCharacter{2032}{\ensuremath{\text{\textasciiacute}}}  % Symbol ′

\DeclareUnicodeCharacter{2070}{\ensuremath{^0}}  % Symbol ⁰
\DeclareUnicodeCharacter{2074}{\ensuremath{^4}}  % Symbol ⁴
\DeclareUnicodeCharacter{2075}{\ensuremath{^5}}  % Symbol ⁵
\DeclareUnicodeCharacter{2076}{\ensuremath{^6}}  % Symbol ⁶
\DeclareUnicodeCharacter{2077}{\ensuremath{^7}}  % Symbol ⁷
\DeclareUnicodeCharacter{2078}{\ensuremath{^8}}  % Symbol ⁸
\DeclareUnicodeCharacter{2079}{\ensuremath{^9}}  % Symbol ⁹
\DeclareUnicodeCharacter{207A}{\ensuremath{^+}}  % Symbol ⁺
\DeclareUnicodeCharacter{207B}{\ensuremath{^-}}  % Symbol ⁻
\DeclareUnicodeCharacter{207F}{\ensuremath{^n}}  % Symbol ⁿ
\DeclareUnicodeCharacter{2080}{\ensuremath{_0}}  % Symbol ₀
\DeclareUnicodeCharacter{2081}{\ensuremath{_1}}  % Symbol ₁
\DeclareUnicodeCharacter{2082}{\ensuremath{_2}}  % Symbol ₂
\DeclareUnicodeCharacter{2083}{\ensuremath{_3}}  % Symbol ₃
\DeclareUnicodeCharacter{2084}{\ensuremath{_4}}  % Symbol ₄
\DeclareUnicodeCharacter{2085}{\ensuremath{_5}}  % Symbol ₅
\DeclareUnicodeCharacter{2086}{\ensuremath{_6}}  % Symbol ₆
\DeclareUnicodeCharacter{2087}{\ensuremath{_7}}  % Symbol ₇
\DeclareUnicodeCharacter{2088}{\ensuremath{_8}}  % Symbol ₈
\DeclareUnicodeCharacter{2089}{\ensuremath{_9}}  % Symbol ₉
\DeclareUnicodeCharacter{208A}{\ensuremath{_+}}  % Symbol ₊
\DeclareUnicodeCharacter{208B}{\ensuremath{_-}}  % Symbol ₋
\DeclareUnicodeCharacter{2091}{\ensuremath{_e}}  % Symbol ₑ
\DeclareUnicodeCharacter{2095}{\ensuremath{_h}}  % Symbol ₕ
\DeclareUnicodeCharacter{2096}{\ensuremath{_k}}  % Symbol ₖ
\DeclareUnicodeCharacter{2097}{\ensuremath{_l}}  % Symbol ₗ
\DeclareUnicodeCharacter{2098}{\ensuremath{_m}}  % Symbol ₘ
\DeclareUnicodeCharacter{2099}{\ensuremath{_n}}  % Symbol ₙ
\DeclareUnicodeCharacter{209A}{\ensuremath{_p}}  % Symbol ₚ
\DeclareUnicodeCharacter{209B}{\ensuremath{_s}}  % Symbol ₛ
\DeclareUnicodeCharacter{209C}{\ensuremath{_t}}  % Symbol ₜ

\DeclareUnicodeCharacter{2113}{\ensuremath{\mathpzc{l}}} % Symbol ℓ

% The command `\mathbbm` requires `\usepackage{bbm}`.
\DeclareUnicodeCharacter{2115}{\ensuremath{\mathbbm{N}}}  % Symbol ℕ
\DeclareUnicodeCharacter{211A}{\ensuremath{\mathbbm{Q}}}  % Symbol ℚ
\DeclareUnicodeCharacter{211D}{\ensuremath{\mathbbm{R}}}  % Symbol ℝ
\DeclareUnicodeCharacter{2124}{\ensuremath{\mathbbm{Z}}}  % Symbol ℤ

\DeclareUnicodeCharacter{2135}{\ensuremath{\aleph}}           % Symbol ℵ
\DeclareUnicodeCharacter{2190}{\ensuremath{\leftarrow}}       % Symbol ←
\DeclareUnicodeCharacter{2192}{\ensuremath{\rightarrow}}      % Symbol →
\DeclareUnicodeCharacter{2194}{\ensuremath{\leftrightarrow}}  % Symbol ↔

% The command `\twoheadrightarrow` requires `\usepackage{amssymb}`.
\DeclareUnicodeCharacter{21A0}{\ensuremath{\twoheadrightarrow}}  % Symbol ↠

\DeclareUnicodeCharacter{21A6}{\ensuremath{\mapsto}}  % Symbol ↦
\DeclareUnicodeCharacter{21A6}{\ensuremath{\mapsto}}  % Symbol ↦

% The command `\restriction` requires `\usepackage{amssymb}`.
\DeclareUnicodeCharacter{21BE}{\ensuremath{\restriction}}  % Symbol ↾

\DeclareUnicodeCharacter{21D0}{\ensuremath{\Leftarrow}}       % Symbol ⇐
\DeclareUnicodeCharacter{21D2}{\ensuremath{\Rightarrow}}      % Symbol ⇒
\DeclareUnicodeCharacter{21D4}{\ensuremath{\Leftrightarrow}}  % Symbol ⇔

\DeclareUnicodeCharacter{2200}{\ensuremath{\forall}}      % Symbol ∀
\DeclareUnicodeCharacter{2203}{\ensuremath{\exists}}      % Symbol ∃
\DeclareUnicodeCharacter{2204}{\ensuremath{\not\exists}}  % Symbol ∃
\DeclareUnicodeCharacter{2205}{\ensuremath{\emptyset}}    % Symbol ∅
\DeclareUnicodeCharacter{2208}{\ensuremath{\in}}          % Symbol ∈
\DeclareUnicodeCharacter{2209}{\ensuremath{\not\in}}      % Symbol ∉
\DeclareUnicodeCharacter{220B}{\ensuremath{\ni}}          % Symbol ∋
\DeclareUnicodeCharacter{220C}{\ensuremath{\not\ni}}      % Symbol ∌

% The command `\blacksquare` requires `\usepackage{amssymb}`.
\DeclareUnicodeCharacter{220E}{{\tiny \ensuremath{\blacksquare}}} % Symbol ∎

\DeclareUnicodeCharacter{220F}{{\tiny \ensuremath{\prod}}}  % Symbol ∏
\DeclareUnicodeCharacter{2211}{\ensuremath{\sum}}           % Symbol ∑
\DeclareUnicodeCharacter{2218}{\ensuremath{\circ}}          % Symbol ∘
\DeclareUnicodeCharacter{2219}{\ensuremath{\bullet}}        % Symbol ∙
\DeclareUnicodeCharacter{221E}{\ensuremath{\infty}}         % Symbol ∞
\DeclareUnicodeCharacter{2223}{\ensuremath{\mid}}           % Symbol ∣
\DeclareUnicodeCharacter{2225}{\ensuremath{\parallel}}      % Symbol ∥
\DeclareUnicodeCharacter{2227}{\ensuremath{\wedge}}         % Symbol ∧
\DeclareUnicodeCharacter{2228}{\ensuremath{\vee}}           % Symbol ∨
\DeclareUnicodeCharacter{2229}{\ensuremath{\cap}}           % Symbol ∩
\DeclareUnicodeCharacter{222A}{\ensuremath{\cup}}           % Symbol ∪
\DeclareUnicodeCharacter{222B}{\ensuremath{\int}}           % Symbol ∫
\DeclareUnicodeCharacter{2234}{\ensuremath{\therefore}}     % Symbol ∴

% The command `\dblcolon` requires `\usepackage{mathtools}`.
\DeclareUnicodeCharacter{2237}{\ensuremath{\dblcolon}} % Symbol ∷

\newcommand{\dotminus}{\mathop{\mbox{$-^{\hspace{-.5em}\cdot}\,$}}}
\DeclareUnicodeCharacter{2238}{\ensuremath{\dotminus}}  % Symbol ∸

\DeclareUnicodeCharacter{223C}{\ensuremath{\sim}}       % Symbol ∼
\DeclareUnicodeCharacter{2243}{\ensuremath{\simeq}}     % Symbol ≃
\DeclareUnicodeCharacter{2245}{\ensuremath{\cong}}      % Symbol ≅
\DeclareUnicodeCharacter{2248}{\ensuremath{\approx}}    % Symbol ≈
\DeclareUnicodeCharacter{224D}{\ensuremath{\asymp}}       % Symbol ≍ : OLIVIER
\DeclareUnicodeCharacter{2250}{\ensuremath{\doteq}}     % Symbol ≐

% The command `\coloneqq` requires `\usepackage{mathtools}`.
\DeclareUnicodeCharacter{2254}{\ensuremath{\coloneqq}}  % Symbol ≔

\newcommand{\eqq}{\stackrel{\text{\tiny ?}}{=}}
\DeclareUnicodeCharacter{225F}{\ensuremath{\eqq}}  % Symbol ≟

\DeclareUnicodeCharacter{2260}{\ensuremath{\neq}}         % Symbol ≠
\DeclareUnicodeCharacter{2261}{\ensuremath{\equiv}}       % Symbol ≡
\DeclareUnicodeCharacter{2262}{\ensuremath{\not\equiv}}   % Symbol ≢
\DeclareUnicodeCharacter{2264}{\ensuremath{\le}}          % Symbol ≤
\DeclareUnicodeCharacter{2265}{\ensuremath{\ge}}          % Symbol ≥
\DeclareUnicodeCharacter{226E}{\ensuremath{\nless}}       % Symbol ≮
\DeclareUnicodeCharacter{226F}{\ensuremath{\ngtr}}        % Symbol ≯
\DeclareUnicodeCharacter{2270}{\ensuremath{\nleq}}        % Symbol ≰
\DeclareUnicodeCharacter{2271}{\ensuremath{\ngeq}}        % Symbol ≱
\DeclareUnicodeCharacter{227A}{\ensuremath{\prec}}        % Symbol ≺
\DeclareUnicodeCharacter{227C}{\ensuremath{\preceq}}      % Symbol ≼
\DeclareUnicodeCharacter{2282}{\ensuremath{\subset}}      % Symbol ⊂
\DeclareUnicodeCharacter{2284}{\ensuremath{\not\subset}} % Symbol ⊄
\DeclareUnicodeCharacter{2283}{\ensuremath{\supset}}      % Symbol ⊃
\DeclareUnicodeCharacter{2286}{\ensuremath{\subseteq}}    % Symbol ⊆
\DeclareUnicodeCharacter{228E}{\ensuremath{\uplus}}       % Symbol ⊎
\DeclareUnicodeCharacter{2291}{\ensuremath{\sqsubseteq}}  % Symbol ⊑
\DeclareUnicodeCharacter{2293}{\ensuremath{\sqcap}}       % Symbol ⊓
\DeclareUnicodeCharacter{2294}{\ensuremath{\sqcup}}       % Symbol ⊔
\DeclareUnicodeCharacter{2295}{\ensuremath{\oplus}}       % Symbol ⊕
\DeclareUnicodeCharacter{22A2}{\ensuremath{\vdash}}       % Symbol ⊢
\DeclareUnicodeCharacter{22A4}{\ensuremath{\top}}         % Symbol ⊤
\DeclareUnicodeCharacter{22A5}{\ensuremath{\bot}}         % Symbol ⊥
\DeclareUnicodeCharacter{22C0}{\ensuremath{\bigwedge}}    % Symbol ⋀
\DeclareUnicodeCharacter{22C2}{\ensuremath{\bigcap}}      % Symbol ⋂
\DeclareUnicodeCharacter{22C3}{\ensuremath{\bigcup}}      % Symbol ⋃
\DeclareUnicodeCharacter{22EE}{\ensuremath{\vdots}}       % Symbol ⋮
\DeclareUnicodeCharacter{22EF}{\ensuremath{\cdots}}       % Symbol ⋯

% The command `\iddots` requires `\usepackage{mathdots}`.
\DeclareUnicodeCharacter{22F0}{\ensuremath{\iddots}} % Symbol ⋰

\DeclareUnicodeCharacter{230A}{\ensuremath{\lfloor}}  % Symbol ⌊
\DeclareUnicodeCharacter{230B}{\ensuremath{\rfloor}}  % Symbol ⌋
\DeclareUnicodeCharacter{2308}{\ensuremath{\lceil}}   % Symbol ⌈
\DeclareUnicodeCharacter{2309}{\ensuremath{\rceil}}   % Symbol ⌉
\DeclareUnicodeCharacter{2329}{\ensuremath{\langle}}  % Symbol 〈
\DeclareUnicodeCharacter{232A}{\ensuremath{\rangle}}  % Symbol 〉

% The command `\Return` requires `\usepackage{keystroke}`.
\DeclareUnicodeCharacter{23CE}{\ensuremath{\Return}}  % Symbol ⏎

% The command `\text` requires `\usepackage{amsmath}` and the command
% `\ding` requires `\usepackage{pifont}`.
\DeclareUnicodeCharacter{2460}{\ensuremath{\text{\ding{172}}}} % Symbol ①
\DeclareUnicodeCharacter{2461}{\ensuremath{\text{\ding{173}}}} % Symbol ②
\DeclareUnicodeCharacter{2462}{\ensuremath{\text{\ding{174}}}} % Symbol ③
\DeclareUnicodeCharacter{2463}{\ensuremath{\text{\ding{175}}}} % Symbol ④
\DeclareUnicodeCharacter{2464}{\ensuremath{\text{\ding{176}}}} % Symbol ⑤
\DeclareUnicodeCharacter{2465}{\ensuremath{\text{\ding{177}}}} % Symbol ⑥
\DeclareUnicodeCharacter{2466}{\ensuremath{\text{\ding{178}}}} % Symbol ⑦
\DeclareUnicodeCharacter{2467}{\ensuremath{\text{\ding{179}}}} % Symbol ⑧
\DeclareUnicodeCharacter{2468}{\ensuremath{\text{\ding{180}}}} % Symbol ⑨

\DeclareUnicodeCharacter{25C5}{\ensuremath{\triangleleft}}  % Symbol ◅

\DeclareUnicodeCharacter{2660}{\ensuremath{\spadesuit}} % Symbol ♠
\DeclareUnicodeCharacter{266D}{\ensuremath{\flat}}      % Symbol ♭
\DeclareUnicodeCharacter{266F}{\ensuremath{\sharp}}     % Symbol ♯

\DeclareUnicodeCharacter{2713}{\ensuremath{\checkmark}}  % Symbol ✓

% The commands `\llbracket` and `\rrbracket` require
% `\usepackage{stmaryrd}`.
\DeclareUnicodeCharacter{27E6}{\ensuremath{\llbracket}}  % Symbol ⟦
\DeclareUnicodeCharacter{27E7}{\ensuremath{\rrbracket}}  % Symbol ⟧

\DeclareUnicodeCharacter{27E8}{\ensuremath{\langle}}          % Symbol ⟨
\DeclareUnicodeCharacter{27E9}{\ensuremath{\rangle}}          % Symbol 〉
\DeclareUnicodeCharacter{27F6}{\ensuremath{\longrightarrow}}  % Symbol % ⟶

\DeclareUnicodeCharacter{2983}{\ensuremath{\mathrm{PDF\;TODO}}}  % Symbol ⦃
\DeclareUnicodeCharacter{2984}{\ensuremath{\mathrm{PDF\;TODO}}}  % Symbol ⦄
\DeclareUnicodeCharacter{2985}{\ensuremath{\mathrm{PDF\;TODO}}}  % Symbol ⦅
\DeclareUnicodeCharacter{2986}{\ensuremath{\mathrm{PDF\;TODO}}}  % Symbol ⦆

% The commands `\llparenthesis` and `\rrparenthesis` require
% `\usepackage{stmaryrd}`.
\DeclareUnicodeCharacter{2987}{\ensuremath{\llparenthesis}}  % Symbol ⦇
\DeclareUnicodeCharacter{2988}{\ensuremath{\rrparenthesis}}  % Symbol ⦈

\DeclareUnicodeCharacter{29F5}{\ensuremath{\setminus}}   % Symbol ⧵

\DeclareUnicodeCharacter{2A06}{\ensuremath{\bigcup}}  % Symbol ⋃
\DeclareUnicodeCharacter{2C7C}{\ensuremath{_j}}       % Symbol ⱼ

\DeclareUnicodeCharacter{D7B0}{\ensuremath{\mathrm{PDF\;TODO}}} % Symbol ힰ

% The command `\mathbbm{1}` requires `\usepackage{bbm}`.
\DeclareUnicodeCharacter{1D7D9}{\ensuremath{\mathbbm{1}}}  % Symbol 𝟙


%%%% MOI


\DeclareUnicodeCharacter{22AC}{\ensuremath{\neg\vdash}}
\DeclareUnicodeCharacter{279C}{\ensuremath{\to}}
\DeclareUnicodeCharacter{1D442}{ }
\DeclareUnicodeCharacter{2000}{ }
\DeclareUnicodeCharacter{1D440}{M}
\DeclareUnicodeCharacter{1D442}{O}
\DeclareUnicodeCharacter{20E3}{2}
\DeclareUnicodeCharacter{1D45A}{m}
\DeclareUnicodeCharacter{1D45B}{n}

\usepackage{pifont}
\DeclareUnicodeCharacter{2714}{\ding{52}}
\DeclareUnicodeCharacter{2705}{\ding{52}}
\DeclareUnicodeCharacter{2717}{\ding{55}}
\DeclareUnicodeCharacter{274C}{\ding{53}}


%% Tolérant tant que macros-markdown.tex n'est pas dans le dépôt :
\InputIfFileExists{macros-markdown}{}{\typeout{*** macros-markdown.tex introuvable: ignore ***}}


%\usepackage{PDFdraftcopy}


%-- Pour état des chapitres


%\newcommand\ETAT{\draftstring{BROUILLON}}
%
%\newcommand\brouillon{\renewcommand\ETAT{\draftstring{BROUILLON}}}
%
%\newcommand\final{\renewcommand\ETAT{\draftstring{\rien}}}
%
%\newcommand\travail{\renewcommand\ETAT{\draftstring{CHANTIER}}}
%
%\newcommand\bazar{\renewcommand\ETAT{\draftstring{BAZAR TOTAL}}}
%
%\final



%%%%%%%%%%%%%%%%%%%%%%%%%%%%%%%%%%%%%%%%%%%%%%%%%
%
%
%                      TRADUCTION EN COURS
%
%%%%%%%%%%%%%%%%%%%%%%%%%%%%%%%%%%%%%%%%%%%%%%%%%


% === volé Amaury



\usepackage{stmaryrd}
\usepackage{ifthen}
\usepackage{mathtools}
\usepackage{dsfont}
\PASSIAPXPROOF{\usepackage{thmtools}}
\usepackage{longtable}
\usepackage{mathdots}
\usepackage{booktabs}


%%%%%%%%%%%%%%%%%%%%%%%%%%%%%%%%%%%%%%%%%%%%%%%%%
%
%
%                      TIKZ STUFFS
%
%
%%%%%%%%%%%%%%%%%%%%%%%%%%%%%%%%%%%%%%%%%%%%%%%%%

%%% volé pour dessin des GPAC: PAS SUR QUE BESOIN de TOUT

\usepackage{tikz}
\usetikzlibrary{calc}
\usepgflibrary{arrows}
\usetikzlibrary{fit}
\usetikzlibrary{patterns}
\usetikzlibrary{decorations.pathreplacing}
\usetikzlibrary{shapes.multipart}
\usetikzlibrary{shapes.geometric}
\usetikzlibrary{shapes.misc}
\usetikzlibrary{positioning}
\usetikzlibrary{graphs}
\usetikzlibrary{topaths}
\usetikzlibrary{arrows.meta}
\usetikzlibrary{decorations.pathreplacing}
\usetikzlibrary{decorations.markings}
\usetikzlibrary{decorations.pathmorphing}
%
 \usetikzlibrary{calc}
 \usetikzlibrary{automata}
\usetikzlibrary{chains}
\usetikzlibrary{scopes}
 \usetikzlibrary{positioning}
 \usetikzlibrary{through,backgrounds}
\usepackage{circuitikz}
% % pour accolades




%%%%%%%%%%%%%%%%%%%%%%%%%%%%%%%%%%%%%%%%%%%%%%%%%
%
%
%                      MACROS locales
%
%
%%%%%%%%%%%%%%%%%%%%%%%%%%%%%%%%%%%%%%%%%%%%%%%%%


%====================
%                      Moi
%=====================


\newcommand\myaddress{ 
{\sc LIX, CNRS, Ecole Polytechnique, Institut Polytechnique de Paris, }\\ 91128 Palaiseau Cedex, FRANCE \\
{\small Olivier.Bournez@lix.polytechnique.fr}}



%====================
%                      un peu de trucs sales
%=====================

%=  pour style lipics et ses trucs

\providecommand\ifnotPOLYCHAPTERMODE[1]{#1}
\providecommand\ifPOLYCHAPTERMODE[1]{}

\ifnotPOLYCHAPTERMODE{
\ifdefined\thechapter
\providecommand\spnewtheorem[4]{\newtheorem{#1}{#2}[chapter]}
\providecommand\spnewtheoremi[5]{\newtheorem{#1}[#5]{#2}}
\else
\providecommand\spnewtheorem[4]{\newtheorem{#1}{#2}}
\providecommand\spnewtheoremi[5]{\newtheorem{#1}[#5]{#2}}
\fi
}
\ifPOLYCHAPTERMODE{
\providecommand\spnewtheorem[4]{\newtheorem{#1}{#2}}
\providecommand\spnewtheoremi[5]{\newtheorem{#1}[#5]{#2}}
}
\providecommand\nolipics[1]{#1}




%====================
%                      un peu de trucs franchement sales
%=====================


% on l'inclus pas ici.
%
%% defaults

\includeversion{NOTLFMI2020}
\includeversion{NOTCSE304}
\excludeversion{NOTDIFFUSIONFINALE}

%% Booleans (dont teste si existe ou pas, pour construire des macros)
% \SAFEMODE
% -> \DIFFUSIONFINALE (via entête-cours)
% -> \PAGEDEGARDE (via entête-cours)
% \POLYMODE
% \POLYCHAPTERMODE

%%%%%%%%%%%%%%%%%%%%%%%%%%%%%%%%%%%%%%%%%%%%%%%%%
%
%
%                      Hack pour DIFFUSION FINALE/ LFMI / Fin provisoire
%
%
%%%%%%%%%%%%%%%%%%%%%%%%%%%%%%%%%%%%%%%%%%%%%%%%%

%\def\DIFFUSIONFINALE{}

\providecommand\ifnotmodeDIFFUSIONFINALE[1]{
\ifdefined\DIFFUSIONFINALE
\else
#1
\fi
}

\providecommand\ifmodeDIFFUSIONFINALE[1]{
\ifdefined\DIFFUSIONFINALE
%% 2021: Comprend pas mais assure
\ifdefined\SAFEMODE
\else
#1
\fi
%#1
\fi
}



\ifnotmodeDIFFUSIONFINALE{\includeversion{NOTDIFFUSIONFINALE}}
\ifnotmodeDIFFUSIONFINALE{\newcommand\FINPROVISOIRELFMI[1]{}}

\providecommand\FINPROVISOIRELFMI{
  \section{Bibliographic notes}

  The current text is highly based (sometimes copied) from:
  
  \bibliographystyle{alpha}
  \bibliography{/Users/bournez/bibliographie/Bib-Files/bournez-avantbib2DOI,/Users/bournez/bibliographie/Bib-Files/perso}
  \end{document}
  }

%\def\POLYCHAPTERMODE{}

\providecommand\ifnotPOLYCHAPTERMODE[1]{
\ifdefined\POLYCHAPTERMODE
\else
#1
\fi
}

\providecommand\ifmodePOLYCHAPTERMODE[1]{
\ifdefined\POLYCHAPTERMODE
#1
\fi
}


%%%%%%%%%%%%%%%%%%%%%%%%%%%%%%%%%%%%%%%%%%%%%%%%%
%
%
%                      POUR INDICATIONS: ETAT
%
%
%%%%%%%%%%%%%%%%%%%%%%%%%%%%%%%%%%%%%%%%%%%%%%%%%


\newcommand\ABSTRACTINTERNAL[1]{
  \montruc[color=black!30]{
  \hrule
  \hrule
  \danger   ABSTRACT:\\[0.5cm] #1

  \vspace{0.5cm}
  \hrule
  \hrule
  }
  }

\newcommand\ABSTRACT[1]{
  \newpage
  \montruc[color=black!20]{
  \hrule
  \hrule
  \danger   ABSTRACT:\\[0.5cm] #1

  \vspace{0.5cm}
  \hrule
  \hrule
  }
  }

  
\newcommand\PARTIEFINALISEE{
  \newpage
  \montruc[color=black!15]{
  \hrule
  \hrule
  \danger   FINALISE A PARTIR ICI
  \hrule
  \hrule
  }
  }
\newcommand\PARTIESEMIFINALISEE{
  \newpage
  \montruc[color=black!15]{
  \hrule
  \hrule
  \danger   1/2 FINALISE A PARTIR ICI
  \hrule
  \hrule
  }
  }


\newcommand\PARTIENONFINALISEE{
  \newpage
  \montruc[color=black!30]{
  \hrule
  \hrule
  \danger   NON FINALISE A PARTIR ICI
  \hrule
  \hrule
  }
  }


%%%%%%%%%%%%%%%%%%%%%%%%%%%%%%%%%%%%%%%%%%%%%%%%%
%
%
%                      Affichage titre du cours: générique
%
%
%%%%%%%%%%%%%%%%%%%%%%%%%%%%%%%%%%%%%%%%%%%%%%%%%


  \newcommand\POLYWARNING[1]{
    
    \chapter*{Preliminaries}

    
\begin{center} 
  These are notes for the course: \\[1cm]
  \textbf{#1}
\end{center}

\bigskip
\fbox{\fbox{
\begin{minipage}{10cm}
\begin{center}
These notes change from days to days.  \\
Please do  not distribute. \\
\end{center}
\end{minipage}
}}

\bigskip
\begin{center}
  \textbf{
    Version of \today.
    }
\end{center}

 % This document is not stable.

\vfill
\begin{center}
Any comment (even about typos, orthography) is welcome: \\
send an email to bournez@lix.polytechnique.fr
\end{center}
}


\newcommand\MYCOURSEON[2]{ 
	%% External self-references
	\ifnotPOLYMODE{
		\ifdefined\MAINAUXFILE
			\externaldocument{\MAINAUXFILE}
		\fi
		}
	
	\ifdefined\PASLAPREMIEREFOIS
		\newpage
	\else
		\makeindex

     		\begin{document}
	\fi
	\gdef\PASLAPREMIEREFOIS{}

	\ifdefined\XPOLY
		\XPAGEDEGARDE{\TITRECOURS\ifPOLYCHAPTERMODE{\\[0.5cm] \Large \bf \LANGUE{Chapter:}{Chapitre:}  #2}}{\SUBTITRECOURS}
		\ifPOLYCHAPTERMODE{
		\renewcommand\thesection{\arabic{section}}
		\chapter*{#2}
		}
	\else
	
      \title{ \ifnotmodeDIFFUSIONFINALE{Course  #1 -}
      \ifPOLYCHAPTERMODE{\LANGUE{Chapter:}{Chapitre:} }
      #2}
      \author{
        Olivier Bournez% \inst{1}
        \\[0.5cm] %{\myaddress}
      }
      \date{Version of \today}
      % \institute{\myaddress}
      




      \ifnotPOLYCHAPTERMODE{
      	\maketitle
		}
      \ifPOLYCHAPTERMODE{\ifdefined\DOPAGEDEGARDE
      			\maketitle
		\fi
		\renewcommand\thesection{\arabic{section}}
%		\renewcommand\thesubsection{\arabic{section}.\ara}
      		\chapter*{#2}}

      % \begin{abstract}
      % \end{abstract}
      \fi
      


      \DOINDEXPRINT
      
      %% 2021: (meme si je comprends pas)
      \ifmodeDIFFUSIONFINALE{
		\excludeversion{recherche}
	}


    }



%%%%%%%%%%%%%%%%%%%%%%%%%%%%%%%%%%%%%%%%%%%%%%%%%
%
%
%                      Affichage titre du cours: variantes
%
%
%%%%%%%%%%%%%%%%%%%%%%%%%%%%%%%%%%%%%%%%%%%%%%%%%

    
\newcommand\MYCOURSE[1]{
  \MYCOURSEON{\csname numerocours#1\endcsname}{\csname nomcours#1\endcsname}
  
    \ifnotmodeDIFFUSIONFINALE{}%\listoftodos}

  }
  
  \newcommand\MYCOURSEnew[1]{
  \MYCOURSEON{???}{#1}
  
  \ifnotmodeDIFFUSIONFINALE{}%\listoftodos}

  }

%   ---------------------------
%   Pour compatilité chapter
%  ----------------------------

\newcommand\PASPREMIEREPAGE[1]{}
\providecommand\mychapter[1]{% \PASPREMIEREPAGE{\newpage} 
  \renewcommand\PASPREMIEREPAGE[1]{##1}  \hrule\hrule \section*{#1}
  \hrule\hrule \  \\[1cm]}
\providecommand\chapter[1]{\PASPREMIEREPAGE{\newpage} 
  \renewcommand\PASPREMIEREPAGE[1]{##1}  \hrule\hrule \section*{#1}
  \hrule\hrule \  \\[1cm]}

%   ---------------------------
%   if not sub
%  ----------------------------


\providecommand\IFNOTSUB[1]{#1}
\IFNOTSUB{
  
%\documentclass{article}
\usepackage{versions}
\includeversion{pasweb}
%\input{macros}

}


 
%%%%%%%%%%%%%%%%%%%%%%%%%%%%%%%%%%%%%%%%%%%%%%%%%
%
%
%               Déclaration de cours
%
%  (l'argument 1 = numéro sert pas)
%
%%%%%%%%%%%%%%%%%%%%%%%%%%%%%%%%%%%%%%%%%%%%%%%%%


\newcommand\COURS[3]{
  \expandafter\newcommand\csname numerocours#2\endcsname{#1}
  \expandafter\newcommand\csname nomcours#2\endcsname{#3}
}



\newcommand\NANCY[3]{
  \expandafter\newcommand\csname numerocours#2\endcsname{#1}
  \expandafter\newcommand\csname nomcours#2\endcsname{#3}
}


\newcommand\INFquatrecentdouze[3]{
  \expandafter\newcommand\csname numerocours#2\endcsname{#1}
  \expandafter\newcommand\csname nomcours#2\endcsname{#3}
}

\newcommand\INFcinqsixun[3]{
  \expandafter\newcommand\csname numerocours#2\endcsname{#1}
  \expandafter\newcommand\csname nomcours#2\endcsname{#3}
}



\newcommand\MPRI[3]{
  \expandafter\newcommand\csname numerocours#2\endcsname{#1}
  \expandafter\newcommand\csname nomcours#2\endcsname{#3}
}

\newcommand\COURSEINCLUDE[2][]{
  \renewcommand\IFNOTSUB[1]{}
  \renewcommand\MYCOURSE[1]{\chapter{\ifnotmodeDIFFUSIONFINALE{#1} \csname nomcours##1\endcsname}
    \ifnotmodeDIFFUSIONFINALE{\mychapter{Fichier: #2}}
  }
   \renewcommand\MYCOURSEnew[1]{\chapter{\ifnotmodeDIFFUSIONFINALE{#1} ##1 }
    \ifnotmodeDIFFUSIONFINALE{\mychapter{Fichier: #2}}
  }
%  \renewcommand\FINSUBMPRI{}
  \input{#2}
}

\newcommand\STRESSCOURSEINCLUDE[1]{
  \COURSEINCLUDE[!STRESS!]{#1}
  }

  \newcommand\subCOURSEINCLUDE[1] {
    \NDLR{BEGIN:  J'inclus/ Je répète ICI #1 }
    \input{INCLUDE/#1}
    \NDLR{END:  J'inclus/ Je répète ICI #1 }
}

%%%%%%%%%%%%%%%%%%%%%%%%%%%%%%%%%%%%%%%%%%%%%%%%%
%
%
%                      POUR INDICATIONS: MESSAGES
%
%
%%%%%%%%%%%%%%%%%%%%%%%%%%%%%%%%%%%%%%%%%%%%%%%%%


% ancienne version
% \newcommand\montruc[1]{\begin{center} {\fbox{\fbox{{#1}
%         }}} \end{center}}




\newcommand\montruc[2][]{
  \ifnotmodeDIFFUSIONFINALE{
  \begin{maboitetruc}[#1]{\danger} #2
  \end{maboitetruc}
%%  \todo[inline,color=green!20,caption={2do}, ssss#1]{\begin{minipage}{\textwidth-4pt}
%\todo[inline,color=green!5%,caption={2do}
%]{\begin{minipage}{\textwidth-4pt}
%    #2\end{minipage}}
}
}

\newcommand\montrucdanger[3][]{
  \ifnotmodeDIFFUSIONFINALE{
                    \montruc[#1]{\danger #2: #3}
                    }
}

\newcommand\montrucdangermargin[3][]{
  \ifnotmodeDIFFUSIONFINALE{
  \montruc[#1]{\danger #2: #3}
\marginpar{\danger #2: #3}
}
 }



%%%%%%%%%%%%%%%%%%%%%%%%%%%%%%%%%%%%%%%%%%%%%%%%%
%
%
%                      GESTION DES TRUCS MARQUES VERSIONS
%
%
%%%%%%%%%%%%%%%%%%%%%%%%%%%%%%%%%%%%%%%%%%%%%%%%%

\makeatletter
\renewcommand\beginmarkversion{\montrucdanger[colback=yellow!10]{=== begin NOT}{}\@Vs@sffbox{\@currenvir$>$}}
 \renewcommand\endmarkversion{\@Vs@sffbox{$<$\@currenvir}\montrucdanger[colback=yellow!10]{end=== NOT}{}}
\makeatother 


%%%%%%%%%%%%%%%%%%%%%%%%%%%%%%%%%%%%%%%%%%%%%%%%%
%
%
%                      VARIANTES DE MESSAGES
%
%
%%%%%%%%%%%%%%%%%%%%%%%%%%%%%%%%%%%%%%%%%%%%%%%%%

 \newcommand\NDLR[1]{\montrucdanger[colback=blue!20]{NDLR}{#1}}
 \newcommand\NLDR[1]{\NDLR{#1}}
 \newcommand\VOIRAUSSI[1]{\montrucdanger[colback=yellow!20]{VOIR AUSSI}{#1}}
 \newcommand\TODO[1]{\montrucdanger[colback=green!20]{TODO}{#1}}
 \newcommand\ATTENTION[1]{\montruc[colback=brown!20]{\danger ATTENTION#1}}
 \newcommand\COMPRENDPAS[1]{\montrucdanger[colback=yellow!20]{JE
     COMPRENDS PAS CA}{#1}}
  \newcommand\POSSIBLE[1]{\montrucdanger[colback=yellow!10]{
      Serait possible}{#1}}
 \newcommand\ENLEVE[1]{\montrucdanger[colback=gray!20]{
      ENLEVE}{#1}}
 
 
  \newcommand\DEPLACE[1]{
    \vspace{0.5cm}
    \hrule
    \hrule
      \vspace{0.5cm}
    \montrucdangermargin[colback=black!40]{DEPLACE/MISSING CONCEPT
    }{#1}
      \vspace{0.5cm}
    \hrule
    \hrule
      \vspace{0.5cm}
  }

\newcommand\BESOINDE[1]{\montrucdangermargin[colback=pink!20]{ATTTENTION
     BESOIN DE}{#1}}
\newcommand\MISSINGCONCEPT[1]{\montrucdangermargin[colback=pink!20]{MISSING CONCEPT
     }{#1}}
\newcommand\SAUTE[1]{\montrucdangermargin[colback=pink!20]{SAUTE
  }{#1}}
\newcommand\SHOULDBEHERE[1]{\BEGINTEMPENVCOLOR{blue!30}\montrucdanger[colback=pink!10]{
SHOULD
    BE HERE
  }{    
    \vspace{1cm}
    
    \centerline{$\vdots$}
    
   \centerline{ #1}
    
      \centerline{$\vdots$}
          
    \vspace{1cm}
    }\ENDTEMPENVCOLOR}
\newcommand\COULDBEHERE[1]{\montrucdangermargin[colback=pink!20]{COULD
    BE HERE
  }{#1}}



%%%%%%%%%%%%%%%%%%%%%%%%%%%%%%%%%%%%%%%%%%%%%%%%%
%
%
%                      POUR INDICATIONS: SOURCES
%
%
%%%%%%%%%%%%%%%%%%%%%%%%%%%%%%%%%%%%%%%%%%%%%%%%%


\newcommand\SOURCEURL[1]{
  \ifnotmodeDIFFUSIONFINALE{
    \marginpar{ \textcolor{magenta}{Source: \url{#1}} }
    }
}

\newcommand\SOURCETXT[1]{
  \ifnotmodeDIFFUSIONFINALE{
    \marginpar{ \textcolor{magenta}{Source: {#1}} }
    }
}

\newcommand\SOURCE[1]{
  \nocite{#1}
  \ifnotmodeDIFFUSIONFINALE{
    \marginpar{ \textcolor{magenta}{Source: \cite{#1}} }
    }
}


%--------------------------------
%--- Raccourci Sources particulières  ---
%--------------------------------

\newcommand\SOURCEJECH{\SOURCE{jech2003set}}
\newcommand\SOURCEKECHRIS{\SOURCE{kechris2012classical}}
\newcommand\SOURCENOTESSETTHEORY{\SOURCE{moschovakis2006notes}}
\newcommand\SOURCEKOZENCT{\SOURCE{LivreKozenTC}}
\newcommand\SOURCEMARKER{\SOURCE{marker2002descriptive}}
\newcommand\SOURCEMOSCHO{\SOURCE{moschovakis2009descriptive}}



%%%%%%%%%%%%%%%%%%%%%%%%%%%%%%%%%%%%%%%%%%%%%%%%%
%
%
%                      POUR INDICATIONS: DELIRES/COMMENTAIRES
%
%
%%%%%%%%%%%%%%%%%%%%%%%%%%%%%%%%%%%%%%%%%%%%%%%%%


%%%%%% ABOUT PLAN

%\newcommand\PLANBUT[1]{\montruc{\begin{minipage}{0.8\textwidth} #1\end{minipage}}}


%%%%%% ABOUT RECHERCHE


\ifnotmodeDIFFUSIONFINALE{
\newenvironment{recherche}
{}{}
%{\color{blue} BEGIN-RECHERCHE:
%%\ifmodeDIFFUSIONFINALE{diffusion finale}
%%\ifnotmodeDIFFUSIONFINALE{not diffusion finale}
%%	\ifdefined\SAFEMODE
%%		SAFEMODE
%%	\else
%%		NOTSAFEMODE
%%	\fi
%  \todo[inline, color=blue]{recherche ici} }{
%  END-RECHERCHE \color{black} }
\tcolorboxenvironment{recherche}{colback=blue!5!white, colframe=red!75!black,fonttitle=\bfseries, colbacktitle=blue!85!black,enhanced,
attach boxed title to top center={yshift=-2mm},
  title={FAIRE RECHERCHE ICI}}
}


%%%%%% ABOUT DELIRE


%\ifnotmodeDIFFUSIONFINALE{
%\newenvironment{delire}{\color{violet} BEGIN-DELIRE: }{
%  END-DELIRE \color{black} }
%}


%%%%%% ABOUT PERSO


%\ifnotmodeDIFFUSIONFINALE{
%\newenvironment{perso}{\color{brown}
%  BEGIN-PERSO: }{
%  END-PERSO\color{black} }
%}



%%%%%% ABOUT COMMENTAIRE

\ifnotmodeDIFFUSIONFINALE{
\newenvironment{commentaire}{\begin{maboitecommentaire}[colback=brown!30]{Commentaire} Commentaire :}{\end{maboitecommentaire}}
\newcommand\COMMENTAIRE[2][Commentaire ]{
\begin{maboitecommentaire}[colback=brown!30]{#1}
Commentaire : #2
\end{maboitecommentaire}
}
}

%%%%%% ABOUT HORSUJET


%\ifnotmodeDIFFUSIONFINALE{
%\newenvironment{probablementhorsujet}{\color{gray}
%  BEGIN-PROBABLEMENT HORS SUJET: }{
%  END-PROBABLEMENT HORS SUJET \color{black} }
%}

%%%%%% ABOUT BONUS


\ifnotmodeDIFFUSIONFINALE{
\newenvironment{pasbesoinpourlinstant}{NOT NEEDEED FOR
  THIS SECTION: }{ }
    \tcolorboxenvironment{pasbesoinpourlinstant}{colback=green!5!white, colframe=gray!75!black,fonttitle=\bfseries, colbacktitle=green!85!black,enhanced,
attach boxed title to top left={yshift=-2mm},
  title={PAS BESOIN POUR L'INSTANT}}
}

%%%%%% ABOUT BONUS


\ifnotmodeDIFFUSIONFINALE{
\newenvironment{bonus}{BONUS: }{
  }
  \tcolorboxenvironment{bonus}{colback=gray!5!white, colframe=gray!75!black,fonttitle=\bfseries, colbacktitle=gray!85!black,enhanced,
attach boxed title to top left={yshift=-2mm},
  title={BONUS}}
}

%%%%%% ABOUT BONUS


\ifnotmodeDIFFUSIONFINALE{
\newcommand\ATTENTIONREPETE[1]{
  \ATTENTION{Repete (attention modif): #1}
  \begin{remarkolivier}
Repete (attention modif): #1
\end{remarkolivier}
}
}


%%%%%%%%%%%%%%%%%%%%%%%%%%%%%%%%%%%%%%%%%%%%%%%%%
%
%
%                      MACROS Papiers Paulin
%
%
%%%%%%%%%%%%%%%%%%%%%%%%%%%%%%%%%%%%%%%%%%%%%%%%%


\newcommand{\x}{\times}
\newcommand{\s}{\sigma}
\newcommand\sep{.}
\newcommand\bzero{\zero}
\newcommand\bun{\un}
\renewcommand{\a}{\alpha}
\def\hd{{\sf hd}}
\def\tl{{\sf tl}}
\def\cons{{\sf cons}}
\def\Cons{{\sf Cons}}
\def\Pr{{\sf Pr}}
\def\Op{{\sf Op}}
\def\Rel{{\sf Rel}}
\def\Equal{{\sf Equal}}
\def\Eval{{\sf Eval}}
\def\Gate{{\sf Gate}}
\def\Selection{{\sf Selection}}
\def\Result{{\sf Result}}
\def\Circuit{{\sf Circuit}}
\def\next{{\sf next}}
\def\comp{{\sf comp}}
\def\finalcomp{{\sf finalcomp}}
\def\Pick{{\sf Pick}}
\def\Hd{{\sf Hd}}
\def\Tl{{\sf Tl}}
\def\RevHd{{\sf RevHd}}
\def\Node{{\sf Node}}
\def\Length{{\sf Length}}
\def\Tape{{\sf Tape}}
\def\Head{{\sf Head}}
\def\Time{{\sf Time}}
\def\unpad{{\sf unpad}}
\def\concat{{\sf concat}}
\def\rg{{\sf right}}
\def\lf{{\sf left}}
\def\st{{\sf state}}
\def\tope{{\sf top}}


%%%%%%%%%%%%%%%%%%%%%%%%%%%%%%%%%%%%%%%%%%%%%%%%%
%
%
%                      MACROS POLY INF561
%
%
%%%%%%%%%%%%%%%%%%%%%%%%%%%%%%%%%%%%%%%%%%%%%%%%%

\newcommand\buglst[1]{#1}

%\newcommand\mydefref[1]{\expandafter\newcommand\csname #1\endcsname{\ref{#1}}}
%\mydefref{refthrisc}
%\mydefref{refthcinqhuit}
%\mydefref{refthcinqsix}
%\mydefref{refthfondamcaract}
%\mydefref{refdefcircuit}
%%% sinon, fait à la main
%\newcommand\refchapalgo{\ref{chapalgo}}
%\newcommand\refchapmodeles{\ref{chapmodeles}}
%\newcommand\refchappreliminaires{\ref{chappreliminaires}}
%\newcommand\refchapcircuits{\ref{chapcircuits}}
%\newcommand \refpropfondamen{\ref{propfondamen}}
%\newcommand\refthhierspace{ref{thhierspace}}
%\newcommand\refchapcomplexitetemps{\ref{chapcomplexitetemps}}
%\newcommand\refdefverifun{\ref{defverifun}}
%\newcommand\refdefverif{\ref{defverif}}
%\newcommand\refthmachinevs{ref{thmachinevs}}
%\newcommand\reflemnotationconfig{\ref{lemnotationconfig}}
%\newcommand\reflemfondam{\ref{lemfondam}}
%\newcommand\refthtroisat{\ref{thtroisat}}
%\newcommand\refthfondam{\ref{thfondam}}
%\newcommand\refthsansram{\ref{thsansram}}
%\newcommand\refdefsram{\ref{defsram}}
%!TEX encoding = UTF-8 Unicode

%% TRUCS DU POLY INF561

\ifnotSLIDE{
\LANGUE{
\spnewtheoremi{exemple}{\NDLR{ECRIT EXEMPLE AU LIEU DE
  EXAMPLE}}{\bfseries}{\rmfamily}{theorem} 
\spnewtheoremi{remarque}{\NDLR{ECRIT REMARQUE AU LIEU DE REMARK}}{\bfseries}{\rmfamily}{theorem} 
}
{
\spnewtheoremi{exemple}{Exemple}{\bfseries}{\rmfamily}{theorem} 
\spnewtheoremi{remarque}{Remarque}{\bfseries}{\rmfamily}{theorem} 
}
}
            
            \newcommand\PERSOSOURCEPOSSIBLE[1]{\PERSOSOURCE{#1}\PERSOCOMMENTAIRE{#1}}
\newcommand\PERSOPOSSIBLE[1]{\POSSIBLE{#1}}
% \newcommand\PERSOPOSSIBLEbug[1]{\NDLR{\textbf{-- possible --
%       enleve car pb compilation}}}
\newcommand\PERSOETAIT[1]{\NDLR{ETAIT: #1}}
%   \begin{minipage}{10cm}
% \NDLR{\textbf{\tiny  -- était --      #1 -- était
%     --
%   \end{minipage}
% }}}

\providecommand\nl{}

\newcommand\PERSOSOURCE[1] { \SOURCETXT{#1}}
\newcommand\PERSOCOMMENTAIRE[1]{\NDLR{#1}}
\newcommand\PERSOCOMMENTAIREd[2]{\NDLR{#1 \textbf{#2}}}


\newenvironment{dessinpb}
{%\shorthandoff{:}
%\selectlanguage{english}
\begin{center}\begin{tikzpicture}[baseline=(current bounding
      box.north)]
%babelbug with tikz
}
{\end{tikzpicture}
%\selectlanguage{frenchb}
\end{center}}



\newcommand\PERSOMANQUE[1]{\NDLR{\textbf{-- manque --
      #1 --manque--}}}

\newcommand\PERSOENLEVE[1]{\ENLEVE{#1}}






%%%%%%%%%%%%%%%%%%%%%%%%%%%%%%%%%%%%%%%%%%%%%%%%%
%
%
%                      MACROS POLY INF412
%
%
%%%%%%%%%%%%%%%%%%%%%%%%%%%%%%%%%%%%%%%%%%%%%%%%%

%\mydefref{refEEXTtruc}
%\mydefref{refEEXTchappreliminaires}
%\mydefref{refEEXTsecarbrebinaire}
%\mydefref{refEEXTsectermes}
%\mydefref{refEEXTchappredicats}
%\mydefref{refEXTchappredicats}
%\mydefref{refEXTchappreliminaires}
%\mydefref{refEEXTdefTuring}
%\mydefref{refEEXTpropnondet}
%\mydefref{refEEXTchapcompletude}
%\mydefref{refEEXTdefconfigalternative}
%\mydefref{refEEXTchapmachinesdeTuring}
%\mydefref{refEEXTpbdedecision}
%\mydefref{refEEXTchapcalc}
%\mydefref{refEEXTincompletude}
%\mydefref{refEEXTthhierspace}
%\mydefref{refEEXTchapcomplexitetemps}

\tikzstyle{nd} = [shape=circle,draw]


\newcommand\mmod{mod}

  %%% VALUATIONS
\newcommand\mmV{v}
\newcommand\mmVb{\overline{\mmV}}
\newcommand\mmVe{\mmVb}

%%% Corrections

\newcommand\exoref[1]{\ref{exo:#1-stt}}

\newenvironment{correction}[1]
{\par\medskip\noindent \label{exo:#1-cor} \hbox{\bf Exercice \ref{exo:#1-stt} 
  \   (page \pageref{exo:#1-stt}). }}
{\par\medskip}

% PROBLEMES DE DECISION

\newcommand\DECISIONint[3]{
%Le problème de décision #1
\begin{itemize}
\item[\textbf{Donnée}:] #2
\item[\textbf{Réponse}:] #3
\end{itemize}
}

\newcommand\INDECIDABLE[5]{
\DECISION{#1}{#2}{#3}

\begin{#4} \label{#5}
Le problème $#1$ % ATTENTION ENLEVE $#1$ en 2020. REMIS 2023
est indécidable.
\end{#4}
}


\newcommand\INDECIDABLEenglish[5]{
\DECISIONenglish{#1}{#2}{#3}

\begin{#4} \label{#5}
The problem $#1$
is undecidable.
\end{#4}
}

\newcommand\INDECIDABLEq[5]{
\DECISION{#1}{#2}{#3}

\begin{#4} \label{#5}
Le problème $#1$
est ?
\end{#4}
}

\newcommand\INDECIDABLEqenglish[5]{
\DECISION{#1}{#2}{#3}

\begin{#4} \label{#5}
The problem $#1$
is ?
\end{#4}
}



\newcommand\DECISIONintenglish[3]{
%Le problème de décision #1
\begin{itemize}
\item[\textbf{Input}:] #2
\item[\textbf{Answer}:] #3
\end{itemize}
}

\newcommand\DECISIONi[4]{
\item $#1$
%Le problème de décision #1
\DECISIONint{#1}{#2}{#3}
}

\newcommand\DECISIONienglish[4]{
\item $#1$
%Le problème de décision #1
\DECISIONintenglish{#1}{#2}{#3}
}


\newcommand\DECISION[3]{
\begin{definition}[$#1$]\ 

%Le problème de décision #1
\DECISIONint{#1}{#2}{#3}
\end{definition}
}


\newcommand\DECISIONenglish[3]{
\begin{definition}[$#1$]\ 
%
%Le problème de décision #1
\DECISIONintenglish{#1}{#2}{#3}
\end{definition}
}



\newcommand\NPC[4]{ 
\DECISION{#1}{#2}{#3}

\begin{#4}
Le problème $#1$
est $\NP$-complet. \LANGUE{\myindex{NP-completeness}}{\myindex{NP-completude@$\NP$-complétude}}
\end{#4}
}


\newcommand\NPCenglish[4]{ 
\DECISIONenglish{#1}{#2}{#3}
\begin{#4}
The problem $#1$
is $\NP$-complete. \LANGUE{\myindex{NP-completeness}}{\myindex{NP-completude@$\NP$-complétude}}
\end{#4}
}


\newcommand\NPCp[5]{
\DECISION{#1}{#2}{#3}

\begin{#4}[#5]
Le problème $#1$
est $\NP$-complet.
\end{#4}
}


\newcommand\NPCpenglish[5]{
\DECISIONenglish{#1}{#2}{#3}

\begin{#4}[#5]
The problem $#1$
is $\NP$-complete.
\end{#4}
}




\newcommand\BRUNOCOMMENTAIRE[1]{\NDLR{Bruno commentaire: #1}}

\ifnotSLIDE{
%\spnewtheorem{exercicec}{Exercice}{\bfseries}{\rmfamily}
%\spnewtheorem{exerciced}{Exercice}{\bfseries}{\rmfamily}
%\spnewtheorem{exercicedc}{Exercice}{\bfseries}{\rmfamily}
%\spnewtheorem{exercicecenglish}{Exercice}{\bfseries}{\rmfamily}
%\spnewtheorem{exercicedenglish}{Exercice}{\bfseries}{\rmfamily}
%\spnewtheorem{exercicedcenglish}{Exercice}{\bfseries}{\rmfamily}
\spnewtheorem{notation}{Notation}{\bfseries}{\rmfamily}
}

%%% INDEX
 \newcommand\motnouv[1]{\emph{#1}\index{#1}}%\SIGNALEMOTNOUV{#1}}%\index{#1}}
\newcommand\INTmotnouv[1]{\emph{#1}}
\newcommand\myindex[1]{\index{#1}}
%\marginpar{-#1=}

\newcommand\synonyme[2]{\kindex{#1|see{#2}}}
\newcommand\indexM[1]{\index{$#1$}}
\newcommand\idx[1]{#1\index{#1}}
\newcommand\idxm[1]{#1\index{$#1$}}
\newcommand\idxM[1]{$#1$\index{$#1$}}
%\newcommand\idxmBUG[1]{#1}
\newcommand\idxmm[2]{#1\index{#2@$#1$}}
\newcommand\idxi[2]{#1\index{#2}}
%\newcommand\motnouv[1]{\INTmotnouv{#1}\index{#1}}
 \newcommand\motnouvi[2]{\INTmotnouv{#1}\index{#2}} % indexage manuel
 \newcommand\motnouva[2]{\INTmotnouv{#1}\index{#2@#1}} % indexage: 2ième
%                                 % argumet = version sans accent
\newcommand\motnouvm[1]{\INTmotnouv{#1}} %indexage manuel fait ailleurs


\newenvironment{PERSOVERBATIM}{}{}
 \newcommand{\joNP}[3]{{ Problème \textbf{ { \motnouv{##1}}:}
            \begin{itemize}
            \item  \textbf{ Instance:} {##2} 
            \item   \textbf{ Question:} {##3}
            \end{itemize}
          }}

%%%%%%%%%%%%%%%%%%%%%%%%%%%%%%%%%%%%%%%%%%%%%%%%%
%
%
%                      POUR COMPATIBILITE AVEC SLIDES & AUTRE
%
%
%%%%%%%%%%%%%%%%%%%%%%%%%%%%%%%%%%%%%%%%%%%%%%%%%


\newcommand\defin[1]{\textbf{#1}}


%%%   COMPATIBILITE AVEC SLIDES

\newcommand\soustitre[1]{\textbf{#1}}


%% COMPATIBILITE AVEC MA THESE

\providecommand\motnouv[1]{\emph{#1}}
\providecommand\motnouvi[2]{\emph{#1}}

\newcommand\papiernorm[1]{\myop{norm}}

%%% OLIVIER
\newcommand\olivier[1]{ \ATTENTION{Olivier: #1}}


%%%%%%%%%%%%%%%%%%%%%%%%%%%%%%%%%%%%%%%%%%%%%%%%%
%
%
%                      Page de Garde Cours X
%
%
%%%%%%%%%%%%%%%%%%%%%%%%%%%%%%%%%%%%%%%%%%%%%%%%%


\includeversion{metdate}

\newcommand\XPAGEDEGARDE[2]{
\thispagestyle{empty}
\title{{#1} \\[2cm] %\\[3cm] 
{\large #2}
}
\author{
 Olivier Bournez% \inst{1}
\\[0.5cm] 
bournez@lix.polytechnique.fr%{\myaddress}
}
%\date{
%\vspace{2cm}
%\includegraphics[height=2.5cm]{/Users/bournez/lib/LaTeX/Perso/FIG-COMMONS/logo_X_new}
%}
\begin{metdate}
\date{
\vspace{2cm}
Version \LANGUE{of}{du} \today \\[1cm] 
\includegraphics[height=3.5cm]{/Users/bournez/lib/LaTeX/Perso/FIG-COMMONS/logo_X_new}
}
\end{metdate}
\maketitle
}


%%%%%%%%%%%%%%%%%%%%%%%%%%%%%%%%%%%%%%%%%%%%%%%%%
%
%
%                      MACROS UTILES
%
%
%%%%%%%%%%%%%%%%%%%%%%%%%%%%%%%%%%%%%%%%%%%%%%%%%

%%%                      Macros
 
%-------------
%--- vectors  ---
%-------------

 
\newcommand\vectorl[1]{{\mathbf#1}}
\newcommand\tu[1]{\vectorl{#1}}

\newcommand\vp{\vectorl{p}}
\newcommand\va{\vectorl{a}}
\newcommand\vb{\vectorl{b}}
\newcommand\vc{\vectorl{c}}
\newcommand\vf{\vectorl{f}}
\newcommand\vn{\vectorl{n}}
\newcommand\vx{\vectorl{x}}
\newcommand\vX{\vectorl{X}}
\newcommand\vy{\vectorl{y}}
\newcommand\vz{\vectorl{z}}
\newcommand\ve{\vectorl{e}}
\newcommand\vv{\vectorl{v}}
\newcommand\vr{\vectorl{r}}
\newcommand\vu{\vectorl{u}}
\newcommand\vA{\vectorl{A}}
\newcommand\vB{\vectorl{B}}
\newcommand\vC{\vectorl{C}}
\newcommand\vD{\vectorl{D}}

%-------------
%--- overline ---
%-------------


\newcommand\ox{\overline{x}}
\newcommand\oy{\overline{y}}
\newcommand\oc{\overline{c}}
\newcommand\oa{\overline{a}}
\newcommand\ow{\overline{w}}
\newcommand\od{\overline{d}}
\newcommand\ob{\overline{b}}
\newcommand\oP{\overline{P}}
\newcommand\oee{\overline{e}}
\newcommand\ot{\overline{t}}
\newcommand\ou{\overline{u}}
\newcommand\ov{\overline{v}}
\newcommand\oz{\overline{z}}
\newcommand\am{\overline{a}}
\newcommand\om{\overline{m}}
\newcommand\oj{\overline{j}}

%----------------
%--- ensembles ---
%----------------

\newcommand\Q{\mathbb{Q}}
\newcommand\R{\mathbb{R}}
\newcommand\Z{\mathbb{Z}}
\providecommand\C{\mathbb{C}}
\newcommand\N{\mathbb{N}}
\newcommand\K{\mathbb{K}}
\newcommand\F{\mathbb{F}}

\newcommand\dyadic{\mathbb{D}}
\newcommand\Zd{{\Z_2}}
\newcommand\Zp{{\Z_p}}
%
\newcommand\RP{\mathbb{R}^{>0}}
\newcommand\QP{\mathbb{Q}^{>0}}


%----------------
%--- lettres cal ---
%----------------

\newcommand\mK{\mathcal{K}}
\newcommand\mD{\mathcal{D}}
\newcommand\mA{\mathcal{A}}
\newcommand\mL{\mathcal{L}}
\newcommand\mX{\mathcal{X}}
\newcommand\mG{\mathcal{G}}
\newcommand\mE{\mathcal{E}}
\newcommand\mH{\mathcal{H}}
\newcommand\mR{\mathcal{R}}
\newcommand\mF{\mathcal{F}}
\newcommand\mJ{\mathcal{J}}
\newcommand\mO{\mathcal{O}}
\newcommand\mP{\mathcal{P}}
\newcommand\mN{\mathcal{N}}
\newcommand\mS{\mathcal{S}}
\newcommand\mM{\mathcal{M}}
\newcommand\mC{\mathcal{C}}
\newcommand\mV{\mathcal{V}}
\newcommand\mI{\mathcal{I}}
\newcommand\mT{\mathcal{T}}

\newcommand\fC{\mathcal{C}}

%----------------
%--- lettres frak ---
%----------------


\newcommand\kM{\mathfrak{M}}
\newcommand\kN{\mathfrak{N}}
\newcommand\kU{\mathfrak{U}}
\newcommand\kg{\mathfrak{g}}



%--------------------------
%--- Façon noter noms classes ---
%---------------------------

\newcommand{\cp}[1]{\operatorname{#1}}


%--------------------------
%--- Calculabilité  ---
%---------------------------

% -- classes
\LANGUE{
\newcommand\RE{\cp{CE}}
}
{\newcommand\RE{\cp{RE}}
}

\newcommand\Rec{\cp{D}}
\newcommand\PR{\cp{PR}}
\newcommand\Gregn{\mE_n}


%--------------------------
%--- Réductions  ---
%---------------------------

\newcommand\lem{\le_{m}}
\newcommand\lelog{\le_{L}}
\newcommand\equivlog{\equiv_{L}}



%--------------------------
%--- Complexité  ---
%---------------------------

% COMPLEXITE

% -- P
\newcommand{\Ptime}{\cp{P}}      % P
\newcommand\PTIME{\Ptime}
\newcommand\PTIMEs[1]{{\PTIME}_{#1}}
\newcommand\PTIMEsp[1]{{\PTIME}_{#1}^0}
\renewcommand\P{\PTIME}
\newcommand{\FPtime}{\cp{FPTIME}}

% -- NP
\newcommand\NP{\cp{NP}}
\newcommand{\NPtime}{\NP}
\newcommand\NPs[1]{\cp{NP}_{#1}}
\newcommand{\FNP}{\cp{FNP}}

\newcommand\NDP{\cp{NDP}}
\newcommand\coNDP{\cp{coNDP}}
\newcommand\NDPs[1]{\cp{NDP}_{#1}}

% -- coNP
\newcommand\coNP{\cp{coNP}}

% -- PSPACE
\newcommand\PSPACE{\cp{ PSPACE}}
\newcommand\NPSPACE{\cp{ NPSPACE}}
\newcommand{\Pspace}{\PSPACE}
\newcommand{\FPspace}{\cp{FPSPACE}}


% -- LOGSPACE
\newcommand\LSPACE{\cp{LOGSPACE}}
\newcommand\NLSPACE{\cp{NLOGSPACE}}
\newcommand\coNLSPACE{\cp{coNLOGSPACE}}
\newcommand\coNL{\cp{coNL}}

% -- Autres
\newcommand\EXPTIME{\cp{EXPTIME}}
\newcommand\NEXPTIME{\cp{NEXPTIME}}
\newcommand\EXPSPACE{\cp{EXPSPACE}}
\newcommand\PH{\cp{PH}}
\newcommand\NC{\cp{NC}}
\newcommand\AC{\cp{AC}}
\newcommand\PPAD{\cp{\mathbf{PPAD}}}
\newcommand\PLS{\cp{\mathbf{PLS}}}

% -- Reverse math
\newcommand{\WKL}{\ensuremath{\docheck{\mathsf{WKL}_0}}}
\newcommand{\ACA}{\ensuremath{\docheck{\mathsf{ACA}_0}}}
\newcommand{\ATR}{\ensuremath{\docheck{\mathsf{ATR}_0}}}
\newcommand{\RCAO}{\ensuremath{\docheck{\mathsf{RCA}_0}}}
\newcommand{\PiCA}{\ensuremath{\docheck{\Pi^1_1\!\text{-}\mathsf{CA}_0}}}


% -- circuits
\newcommand\CIRCUIT[2]{\cp{CIRCUIT(#1,#2)}}
\newcommand\UCIRCUIT[2]{\cp{UCIRCUIT(#1,#2)}}

% -- Classes proba
\newcommand\PP{\cp{PP}}
\newcommand\RPP{\cp{RP}}
\newcommand\ZPP{\cp{ZPP}}
\newcommand\coRP{\cp{coRP}}
\newcommand\BPP{{\cp{BPP}}}

% -- IP
\newcommand\IP{\cp{IP}}
\newcommand\PCP{\cp{PCP}}

% -- non-uniforme
\newcommand\nuP{\mathbb{P}}
\newcommand\nuPs[1]{\mathbb{P}_{#1}}
\newcommand\nuPsp[1]{\mathbb{P}_{#1}^0}
\newcommand\nuNP{\mathbb{N}\mathbb{P}}
%\newcommand\poly{/{poly}}
\newcommand\Ppoly{\PTIME_{\poly}}
\newcommand\NPpoly{\NP_{\poly}}

% -- conditions d'uniformité
\newcommand\Luniforme{\cp{L}}
\newcommand\BCuniforme{\cp{BC}}


% Classes de circuits
\newcommand\TC{\cp{TC}}
\newcommand\LFP{\cp{LFP}}
\newcommand\FAC{\cp{FAC}}
\newcommand\FACC{\cp{FACC}}
\newcommand\FTC{\cp{FTC}}
\newcommand\FNC{\cp{FNC}}

% Logiques temorelles
\newcommand\CTL{\cp{CTL}}
\newcommand\LTL{\cp{LTL}}

% -- mesure temps/espace/taille

\newcommand\DTIME{\cp{DTIME}}
\newcommand\TIME[1]{\cp{TIME(#1)}}
\newcommand\TIMEs[2]{\cp{TIME_{#2}(#1)}}
\newcommand\TIMEsp[2]{\cp{TIME^0_{#2}(#1)}}
\newcommand\NTIME[1]{\cp{NTIME(#1)}}
\newcommand\NTIMEs[2]{\cp{NTIME_{#2}(#1)}}
\newcommand\NTIMEsp[2]{\cp{NTIME_{#2}^0(#1)}}
\newcommand\DSPACE{\cp{DSPACE}}
\newcommand\SPACE[1]{\cp{SPACE(#1)}}
\newcommand\SPACEs[2]{\cp{SPACE_{#2}(#1)}}
\newcommand\NSPACE[1]{\cp{NSPACE(#1)}}
\newcommand\NSPACEs[2]{\cp{NSPACE_{#2}(#1)}}
\newcommand\coNSPACE[1]{\cp{coNSPACE(#1)}}
\newcommand\TIMEON[2]{\cp{TIME(#1,#2)}}
\newcommand\SPACEON[2]{\cp{SPACE(#1,#2)}}
\newcommand\SIZE[1]{\cp{ size(#1)}}
\newcommand\ATIME{\cp{ATIME}}
\newcommand\ASPACE{\cp{ASPACE}}


%% autre:
\newcommand\LH{\cp{LH}}
\newcommand\FO{\cp{FO}}
\newcommand\SO{\cp{SO}}
\newcommand\ESOhorn{\cp{ESO_{horn}}}
\newcommand\SOhorn{\cp{SO_{horn}}}
\newcommand\ACzero{\cp{AC_0}}
\newcommand{\LINSPACE}{\cp{LINSPACE}}
\newcommand{\mFLINSPACE}{\mF_{\cp{LINSPACE}}}
\newcommand{\LOGSPACE}{\cp{LOGSPACE}}
\newcommand\ALOGTIME{\cp{ALOGTIME}}
\newcommand\RF{\cp{RF}}


\newcommand\BOOL[1]{\cp{{BOOLEAN(#1)}}}
\newcommand\DET{\cp{ DET}}

% % -- dirty


\newcommand{\cPspace}{\cp{\sharp PSPACE}}
\newcommand{\cAPtime}{\cp{\sharp APTIME}}
%\newcommand{\FPspace}{\cp{FPSPACE}}
%\newcommand{\FPspaceN}{\cp{FPSPACE}^{+}}
%\newcommand{\FPspaceN}{\cp{FPSPACE}}



%-------------------------
%--- Problèmes de décision  ---
%--------------------------

\newcommand\decisionname[1]{\cp{#1}}


\newcommand\Luniv{\decisionname{L_{univ}}}
\newcommand\HALTINGPROBLEM{\decisionname{HALTING-PROBLEM}}
\newcommand\SAT{\decisionname{SAT}}
\newcommand\MAXtSAT{\decisionname{MAX3SAT}}
\newcommand\REACH{\decisionname{REACH}}
\newcommand\CIRCUITVALUE{\decisionname{CIRCUITVALUE}}
\newcommand\MONOTONECIRCUITVALUE{\decisionname{MONOTONECIRCUITVALUE}}
\newcommand\THEOREMS{\decisionname{THEOREMS}}
\newcommand\QCIRCUITSAT{\decisionname{QCIRCUITSAT}}
\newcommand\QSAT{\decisionname{QSAT}}
\newcommand\QBF{\decisionname{QBF}}
\newcommand\SATVALUE{\decisionname{SATVALUE}}
\newcommand\SUCCINTSAT{\decisionname{SUCCINTSAT}}
\newcommand\SUCCINTSATVALUE{\decisionname{SUCCINTSATVALUE}}
\newcommand\GEOGRAPHY{\decisionname{GEOGRAPHY}}
\newcommand\PARITY{\decisionname{PARITY}}
\newcommand\MAJ{\decisionname{MAJ}}
\newcommand\CVP{\CIRCUITVALUE}
\newcommand\BOOLEANPROGRAMVALUE{BOOLEAN~PROGRAM~VALUE}

\newcommand\CLIQUE{\cp{CLIQUE}}
\LANGUEinv{
\newcommand\RECOUVREMENTDESOMMETS{\cp{RECOUVREMENT~DE~SOMMETS}}
\newcommand\RS{\RECOUVREMENTDESOMMETS}
}{
\newcommand\RECOUVREMENTDESOMMETSeng{\cp{VERTEX~COVER}}
\newcommand\RSeng{\RECOUVREMENTDESOMMETSeng}
}

\LANGUEinv{\newcommand\COUPUREMAXIMALE{\cp{COUPURE~MAXIMALE}}}
{\newcommand\COUPUREMAXIMALEeng{\cp{MAXIMAL~CUT}}}
\newcommand\CIRCUITHAMILTONIEN{\CHAMILTONIEN}
\LANGUEinv{\newcommand\VOYAGEURDECOMMERCE{\VCOMMERCE}}
{\newcommand\VOYAGEURDECOMMERCEeng{\VCOMMERCEeng}}

\LANGUEinv{\newcommand\CIRCUITLEPLUSLONG{\cp{CIRCUIT~LE~PLUS~LONG}}}
{\newcommand\CIRCUITLEPLUSLONGeng{\cp{LONGEST~CIRCUIT}}}
\LANGUEinv{\newcommand\SOMMEDESOUSENSEMBLE{\SUBSETSUM}}{
\newcommand\SOMMEDESOUSENSEMBLEeng{\SUBSETSUMeng}}
\LANGUEinv{
\newcommand\SACADOS{\cp{SAC~A~DOS}}}
{\newcommand\SACADOSeng{\cp{KNAPSACK}}}

%français
\LANGUEinv{
\newcommand\CHAMILTONIEN{\cp{CIRCUIT~HAMILTONIEN}}}
{\newcommand\CHAMILTONIENeng{\cp{HAMILTONIAN~CIRCUIT}}}
\LANGUEinv{
\newcommand\VCOMMERCE{\cp{VOYAGEUR~DE~COMMERCE}}}
{\newcommand\VCOMMERCEeng{\cp{TRAVELING~SALESMAN}}}
\LANGUEinv{
\newcommand\kCOLORABILITE{\cp{k-COLORABILITE}}}
{\newcommand\kCOLORABILITEeng{\cp{k-COLORABILITY}}}
\LANGUEinv{
\newcommand\COLORIAGE{\cp{BON~COLORIAGE}}}
{\newcommand\COLORIAGEeng{\cp{GOOD~COLORING}}}
\LANGUEinv{
\newcommand\tCOLORABILITE{\cp{3-COLORABILITE}}}
{\newcommand\tCOLORABILITEeng{\cp{3-COLORABILITY}}}
\LANGUEinv{
\newcommand\SUBSETSUM{\cp{SOMME~DE~SOUS~ENSEMBLE}}}
{\newcommand\SUBSETSUMeng{\cp{SUBSET~SUM}}}
\newcommand\PARTITION{\cp{PARTITION}}
\LANGUEinv{
\newcommand\NOMBREPREMIER{\decisionname{NOMBRE~ PREMIER}}}
{\newcommand\NOMBREPREMIEReng{\decisionname{PRIME~NUMBER}}}
\LANGUEinv{
\newcommand\CODAGE{\decisionname{CODAGE}}}
{\newcommand\CODAGEeng{\decisionname{ENCODING}}}
\newcommand\kSAT{\cp{k-SAT}}
\newcommand\tSAT{\cp{3-SAT}}
\newcommand\NAESAT{\cp{NAESAT}}
\newcommand\NAEtSAT{\cp{NAE3SAT}}
\newcommand\NAEqSAT{\cp{NAE4SAT}}
\newcommand\STABLE{\cp{\LANGUE{INDEPENDANT~SET}{STABLE}}}

%\newcommand\CM{\cp{COUPURE MAXIMALE}}
\LANGUEinv{
\newcommand\POSTpb{\cp{Problème~de~ la~correspondance~de~Post}}}
{\newcommand\POSTpbeng{\cp{Post~correspondence~problem}}}

\LANGUEinv{
\newcommand\PARITE{\cp{PARITE}}}
{\newcommand\PARITEeng{\cp{PARITY}}}


%-------------------------
%--- Modèles parallèles  ---
%--------------------------


% PARALLELISME
\newcommand\RAM{\cp{RAM}}
%
\newcommand\PRAM{\cp{PRAM}}
\newcommand\CRCW{\cp{CRCW}}
\newcommand\CREW{\cp{CREW}}
\newcommand\EREW{\cp{EREW}}


%------------------------
%--- codages  ---
%------------------------


\newcommand\codage[1]{\langle #1 \rangle}
\newcommand\tuple[1]{\codage{#1}}
\newcommand\paire[2]{\codage{#1,#2}}
\newcommand\triplet[3]{\codage{#1,#2,#3}}
\newcommand\quadruplet[4]{\codage{#1,#2,#3,#4}}
\newcommand\cinquplet[5]{\codage{#1,#2,#3,#4,#5}}


%------------------------
%--- descriptive set theory  ---
%------------------------

\newcommand\baire{\mathcal{N}}
\newcommand\peano{\overline{\mathcal{N}}} %% A MOI
\newcommand\cantor{\mathcal{C}}
%
\newcommand\bSigma{\mathbf{\Sigma}}
\newcommand\bPi{\mathbf{\Pi}}
\newcommand\bDelta{\mathbf{\Delta}}
%


%%%%%% ABOUT NOMS QUI CHANGES

\newcommand\WF{\cp{WF}}
\newcommand\WO{\WF}
\newcommand\LO{\cp{LO}}
\newcommand\DLO{\cp{DLO}}
% 
\newcommand\rk{\cp{rank}}


%% Moins propre:

\newcommand\I{\mathbb{I}}
\renewcommand\H{\mathbb{H}}
%
\newcommand\leKB{<_{KB}}
%
\newcommand\finiomega{{<\omega}}
\newcommand\compl{^\frown}
\newcommand\prefix{ ~ prefix~  } 
\newcommand\X{{X}}
\newcommand\Y{\mathcal{Y}}
\newcommand\subsetstrict{\subset}

\renewcommand\complement[1]{#1^{c}}


%------------------------
%--- schémas de récursion ---
%------------------------

%%                      Source papier MFCS 2019


% COMPLEXITY

%% MFCS differe de Clote:
%% Version à taille restreinte

\newcommand\co{\cp{co}}

\newcommand\mFPTIME{\mF_{PTIME}}
\newcommand{\mFPSPACE}{\mF_{SPACE}}
\newcommand{\mFPH}{\mF_{PH}}

\newcommand\SigmaP{\Sigma^P}
\newcommand\DeltaP{\Delta^{P}}
\newcommand\PiP{\Pi^{P}}
\newcommand\coSigmaP{\co\SigmaP}
\newcommand\coPiP{\mF{\co\PiP}}

%% autres classes:


%% Schemas pour complexité implicite

\newcommand\COMP{\cp{ COMP}}
\newcommand\CRN{\cp{ CRN}}
\newcommand\BRN{\cp{ BRN}}
\newcommand\SBRN{\cp{ SBRN}}
\newcommand\BR{\cp{ BR}}
\newcommand\kBR{k-\cp{ BR}}
\newcommand \WBPR{\cp{ WBPR}}
\newcommand \BMIN{\cp{ BMIN}}
\newcommand \BSUM{\cp{ BSUM}}
\newcommand\SBSUM{\cp{ SBSUM}}
\newcommand \BPROD{\cp{ BPROD}}
\newcommand \SBPROD{\cp{ SBPROD}}
\newcommand \SCOMP{\cp{ SCOMP}}
\newcommand \SR{\cp{ SR}}
\newcommand \SRN{\cp{ SRN}}

%% devrait etre défini
\newcommand\WBRN{\cp{ WBRN}}



%-- godel
\newcommand\INTDIV{\cp{INTDIV}}
\newcommand\DIV{\cp{DIV}}
\newcommand\EVEN{\cp{EVEN}}
\newcommand\ODD{\cp{ODD}}
\newcommand\PRIME{\cp{PRIME}}
\newcommand\POWER{\cp{POWER}}
\newcommand\BIT{\cp{BIT}}
\newcommand\negHP{\overline{\cp{HP}}}
\newcommand\HP{\cp{HP}}
\newcommand\negLuniv{\overline{\Luniv}}
\newcommand\VALCOMP{\cp{VALCOMP}}
\newcommand\mWINDOW{\cp{LEGAL}}
\newcommand\mmWINDOW{\cp{WINDOW}}
\newcommand\LENGTH{\cp{LENGTH}}
\newcommand\DIGIT{\cp{DIGIT}}
\newcommand\MATCH{\cp{MATCH}}
\newcommand\MOVE{\cp{MOVE}}
\newcommand\START{\cp{START}}
\newcommand\CELL{\cp{CELL}}
\newcommand\HALT{\cp{HALT}}
\newcommand\SUBST{\cp{SUBST}}
\newcommand\PROOF{\cp{PROOF}}
\newcommand\PROVABLE{\cp{PROVABLE}}
\newcommand\CONSIST{\cp{CONSIST}}



%---------------------
%--- calculus (AP) ---
%---------------------

% \jacobian{f}{x}
\newcommand{\jacobian}[1]{J_{#1}}
%\newcommand{\grad}[1]{\overrightarrow{\operatorname{grad}}~{#1}}
\newcommand{\grad}[1]{\nabla{#1}}
\newcommand{\scalarprod}[2]{{#1}\cdot{#2}}
\newcommand{\bigO}[1]{\mathcal{O}\left(#1\right)}
\newcommand\smallO{o}
\newcommand{\softO}[1]{\tilde{\mathcal{O}}\left(#1\right)}
\newcommand{\taylor}[3]{{T_{#1}^{#2}#3}}
\newcommand{\transpose}[1]{{#1}^T}
\newcommand{\pastsup}[2]{{\sup}_{#1}#2}

\DeclareMathOperator\arctanh{arctanh}



%---------------
%--- norms (AP) ---
%---------------


\newcommand{\inorm}[2]{\left\lVert{#1}\right\rVert_{#2}}
\newcommand{\twonorm}[1]{\inorm{#1}{2}}
\newcommand{\infnorm}[1]{\inorm{#1}{}}
\newcommand{\normsup}[1]{\infnorm{#1}{}}
%
\newcommand\hausdorff{d_H}
%
\newcommand\norm[1]{\infnorm{#1}}


%----------------
%--- topology ---
%----------------

% \openball{pt}{radius}
 \newcommand{\openball}[2]{B_{#2}(#1)}
\newcommand{\closedball}[2]{\ovl{B}_{#2}(#1)}
\newcommand{\closure}[1]{\ovl{#1}}
% \newcommand{\interior}[1]{\mathring{#1}}
% \newcommand{\border}[1]{\partial{#1}}




%-----------------
%--- functions (AP) ---
%-----------------
 
\newcommand{\lambdafun}[2]{(#1)\mapsto{#2}}
\newcommand{\lambdafunex}[2]{#1\mapsto{#2}}
\newcommand{\argmin}{\operatornamewithlimits{argmin}}
\newcommand{\argmax}{\operatornamewithlimits{argmax}}


 
%-----------------------
%--- Machines de Turing  ---
%-----------------------

\newcommand\blanc{\vectorl{B}}



%-----------------
%--- Logique ---
%-----------------
 
\newcommand\sem[1]{{[\semspace[}#1{]\semspace]}}
\newcommand\sems[1]{{[\semspace[}#1{]\semspace]}}
\newcommand\semss[1]{\sem{#1}_{\sigma'}}
\newcommand\semspace{\kern -.25ex}
\newcommand\true{true}
\newcommand\false{false}
%extension élémentaire:
\newcommand\elem{\preceq}
\newcommand\godel[1]{\lceil #1 \rceil}

%-----------------
%--- Model theory ---
%-----------------

\newcommand\hull[1]{\langle #1 \rangle}

%-----------------
%--- Ordinaux ---
%-----------------

\newcommand\omegaunCK{\omega_1^{CK}}
\newcommand\omegaunck{\omegaunCK}

\DeclareMathOperator{\cof}{cof}

%-----------------
%--- Surréels (QG) ---
%-----------------

\newcommand{\On}{\textbf{On}}
\newcommand{\Nobf}{\textbf{No}}
\newcommand{\Nobb}{\ensuremath{\mathbbmss{No}}}
\newcommand{\Nolt}[1]{\ensuremath{\Nobf_{< #1}}}
\newcommand{\Noleq}[1]{\ensuremath{\Nobf_{\leq #1}}}
\newcommand{\Noltbb}[1]{\ensuremath{\Nobb_{< #1}}}
\providecommand\No{\Nobf}
\newcommand\Noalpha{\Nolt{\alpha}}
\newcommand\Nolambda{\Nolt{\lambda}}

\DeclareMathOperator{\length}{length}
\DeclareMathOperator{\supp}{supp} % operateur support


%----------------------
%---Maths  de base (AP)  ---
%----------------------

\DeclareMathOperator{\dom}{dom} % operateur support

\DeclareMathOperator{\size}{size}

\newcommand{\powerset}[1]{\mathcal{P}(#1)}
\newcommand{\card}[1]{{\##1}}
\newcommand\priv{ - }
\newcommand\prive{-}
\renewcommand{\restriction}{\mathord{\upharpoonright}}
\newcommand\restrict[1]{_{\restriction #1}}

% \providecommand\mod{}
%% UNDEFINE:
\let\mod\relax
\let\div\relax
\DeclareMathOperator{\mod}{mod} 
\DeclareMathOperator{\div}{div} 

% partieentier
\newcommand\entieres[1]{\lceil #1 \rceil}
\newcommand\entierei[1]{\lfloor #1 \rfloor}


%----------------------
%---Symbole danger ---
%----------------------


%\providecommand*{\TakeFourierOrnament}[1]{{%
%\fontencoding{U}\fontfamily{futs}\selectfont\char#1}}
%\providecommand*{\danger}{{\Huge \TakeFourierOrnament{66}}}



%%%%%%%%%%%%%%%%%%%%%%%%%%%%%%%%%%%%%%%%%%%%%%%%%
%
%
%                      ASMeries
%
%
%%%%%%%%%%%%%%%%%%%%%%%%%%%%%%%%%%%%%%%%%%%%%%%%%



% CIRCUITS
%-- particuliers
\newcommand\REPLACE{\cp{REPLACE}}
\newcommand\EQUAL{\cp{EQUAL}}
\newcommand\EVALUE{\cp{EVALUE}}
\newcommand\TVALUE{\cp{TVALUE}}
\newcommand\UPDATE{\cp{UPDATE}}
\newcommand\ASSIGN{\cp{ASSIGN}}
\newcommand\PAR{\cp{PAR}}
\newcommand\DO{\cp{DO}}


% ASM
\newcommand\emplacement[2]{(#1,#2)}
\newcommand\contenu[2]{\sem{(#1,#2)}}
\newcommand\affecte[3]{(#1,#2)\leftarrow#3}
\newcommand\ctl{ctl\_state}


%%%%%%%%%%%%%%%%%%%%%%%%%%%%%%%%%%%%%%%%%%%%%%%%%
%
%
%                      GPACqueries
%
%
%%%%%%%%%%%%%%%%%%%%%%%%%%%%%%%%%%%%%%%%%%%%%%%%%



%------------------------
%--- generable zoo (GPAC) ---
%------------------------

\newcommand{\myop}[1]{\operatorname{#1}}
\newcommand{\abs}{\myop{abs}}
\newcommand{\sg}{\myop{sg}}
\newcommand{\ip}[1]{\myop{ip}_{#1}}
\newcommand{\rnd}{\myop{rnd}}
\newcommand{\lil}{\myop{lil}}
\newcommand{\lxh}{\myop{lxh}}
\newcommand{\hxl}{\myop{hxl}}
\newcommand{\plil}{\myop{plil}}
\newcommand{\reach}{\myop{reach}}
%\newcommand{\mreach}{\myop{mreach}}
\newcommand{\watchdog}{\myop{watchdog}}
\newcommand{\monitor}{\myop{monitor}}
%\newcommand{\norm}{\myop{norm}}
\newcommand{\mx}{\myop{mx}}
\newcommand{\mn}{\myop{mn}}
\newcommand{\sample}{\myop{sample}}
\newcommand{\select}{\myop{select}}
\newcommand{\ssine}[1]{\sin_{#1}}
\newcommand{\clamp}{\myop{clamp}}

\newcommand{\fnsg}[3]{tanh((#1)*(#2)*(#3))}
\newcommand{\fnipone}[3]{(1+\fnsg{(#1)-1}{#2}{#3})/2.}
\newcommand{\fnipinf}[1]{#1-sin(2*pi*(#1))/2./pi}%
\newcommand{\fnlil}[5]{%
\fnipone{(#3)-(#1)+1-((#2)-(#1))/8.}{#4+log(1+4./(#2-#1)*(#5)*(#5))+log(1+4)}{8./(#2-#1)}%
*\fnipone{#2-#3+1-(#2-#1)/8.}{#4+log(1+4/(#2-#1)*(#5)*(#5))+log(1+4)}{8./(#2-#1)}%
*4./(#2-#1)*(#5)}
\newcommand{\fnlxh}[5]{\fnipone{(#3)-(#1+#2)/2.+1}{#4+log(1+(#5)*(#5))}{2./(#2-#1)}*(#5)}
\newcommand{\fnhxl}[5]{\fnipone{(#1+#2)/2.-(#3)+1}{#4+log(1+(#5)*(#5))}{2./(#2-#1)}*(#5)}
\newcommand{\fnplilf}[4]{sin(((#4)-((#1)+(#2))/2.+(#3)/4.)*2.*pi/(#3))}
\newcommand{\fnplil}[6]{(#6)*\fnlxh{\fnplilf{#1}{#2}{#3}{#1}}%
{\fnplilf{#1}{#2}{#3}{#1+((#2)-(#1))/4.}}%
{\fnplilf{#1}{#2}{#3}{#4}}%
{#5+2.+log(1+(#6)*(#6))}%
{1./4.+2./((#2)-(#1))}}
\newcommand{\fnssine}[2]{sin(#2)*(1-cos(#2)/cos(#1))}



%--------------------
%--- complexity GPAC ---
%--------------------


%\ifthenelse{\equal{#1}{}}{Empty.}{Nonempty.}
\newcommand{\myclass}[1]{\operatorname{#1}}
% \newcommand{\gcomp}[2]{\ensuremath{\myclass{GCOMP}(#1,#2)}}
% \newcommand{\ggenerable}[1]{\textcolor{red}{\ensuremath{#1-}\text{generable}}}
 \newcommand{\gval}[2][]{\ensuremath{\myclass{GVAL}_{#1}\ifthenelse{\equal{#2}{}}{}{[#2]}}}
 \newcommand{\gpval}[1][]{\ensuremath{\myclass{GPVAL}_{#1}}}
 \newcommand{\gc}[3][]{\ensuremath{\myclass{ATSC}(#2,#3)}}
 \newcommand{\gpc}[1][]{\ensuremath{\myclass{ATSP}}}
\newcommand{\cgc}[1][]{\ensuremath{\myclass{ATS}}}
\newcommand{\gexpc}[1][]{\ensuremath{\myclass{AEXP}}}
\newcommand{\gwc}[2]{\ensuremath{\myclass{AW}(#1,#2)}}
\newcommand{\gpwc}{\ensuremath{\myclass{AWP}}}
\newcommand{\cgwc}{\ensuremath{\myclass{AW}}}
\newcommand{\grc}[3]{\ensuremath{\myclass{AR}(#1,#2,#3)}}
\newcommand{\gprc}{\ensuremath{\myclass{ARP}}}
\newcommand{\gsc}[3]{\ensuremath{\myclass{AS}(#1,#2,#3)}}
\newcommand{\gpsc}{\ensuremath{\myclass{ASP}}}
\newcommand{\goc}[3]{\ensuremath{\myclass{AO}(#1,#2,#3)}}
\newcommand{\gpoc}{\ensuremath{\myclass{AOP}}}
\newcommand{\cgoc}{\ensuremath{\myclass{AO}}}
\newcommand{\guc}[4]{\ensuremath{\myclass{AX}(#1,#2,#3,#4)}}
\newcommand{\gpuc}{\ensuremath{\myclass{AXP}}}
\newcommand{\glc}[2][]{\ensuremath{\myclass{AL}(#2)}}
\newcommand{\gplc}[1][]{\ensuremath{\myclass{ALP}}}
\newcommand{\cglc}[1][]{\ensuremath{\myclass{AL}}}

%\newcommand{\intinterv}[2]{\{#1,\ldots,#2\}}
\newcommand{\intinterv}[2]{\llbracket#1,#2\rrbracket}

\newcommand{\Rgen}{\R_{G}}
\newcommand{\Rpoly}{\R_{P}}

%%%%%%%%%%%%%%%%%%%%%%%%%%%%%%%%%%%%%%%%%%%%%%%%%
%
%
%                      Mes codages obscures
%
%
%%%%%%%%%%%%%%%%%%%%%%%%%%%%%%%%%%%%%%%%%%%%%%%%%




%%% Codages

% \newcommand\gammaoc{\gamma_{[0,1]}}
% \newcommand\gammaonc{\gamma_{\N}}
\newcommand\gammaTMc{\gamma^{TM}_{[0,1]}}
\newcommand\gammaTMcatan{\gamma^{TM}_{(0,1)^2}}
\newcommand\gammaTMnc{\gamma^{TM}_{\N^3}}
\newcommand\gammaTMncd{\gamma^{TM}_{\N^2}}
% \newcommand\gammaTMe{\gamma^{TM}_{*}}
% \newcommand\gammas{\gamma_{sab}}
% \newcommand\gammaord{\gamma_{cantor}}


\newcommand\enccompact{\nu_{X}}
 \newcommand\encinfinite{\nu_\N}
 \newcommand\enc{\nu}
 
%%%%%%%%%%%%%%%%%%%%%%%%%%%%%%%%%%%%%%%%%%%%%%%%%
%
%
%                      DESSINS ET IMAGES
%
%
%%%%%%%%%%%%%%%%%%%%%%%%%%%%%%%%%%%%%%%%%%%%%%%%%


\newcommand\IMAGE[2]{ \begin{center} 
    \includegraphics[#1]{#2}
  \end{center}
  }


  
%%%% POUR FAIRE DESSIN.
\newenvironment{dessinpp}[1]{
\begin{tikzpicture}[baseline=(current bounding
      box.west),#1]
}
{\end{tikzpicture}}




\newenvironment{dessinp}[1]{
\begin{center}\begin{tikzpicture}[baseline=(current bounding
      box.north),#1]
}
{\end{tikzpicture}
\end{center}}


\newenvironment{dessin}{
\begin{center}\begin{tikzpicture}[baseline=(current bounding
      box.north)]
}
{\end{tikzpicture}
\end{center}}



\newenvironment{dessind}{
\begin{tikzpicture}[baseline=(current bounding
      box.north)]
}
{\end{tikzpicture}
}



%%%%%%%%%%%%%%%%%%%%%%%%%%%%%%%%%%%%%%%%%%%%%%%%%
%
%
%                      POUR DESSINS GPAC
%
%
%%%%%%%%%%%%%%%%%%%%%%%%%%%%%%%%%%%%%%%%%%%%%%%%%


%%%% MACROS DE BASE.


\newcommand\constant[2]
{
node (#1)  [shape=circle,draw] {#2};
\draw (#1.east) -- +(0.3,0) coordinate (#1-o) ; 
\draw (#1.west) -- +(-0.3,0) coordinate (#1-i) ; 
}

\newcommand\basicproduct[2]
{
node (#1) {#2};
\draw (#1) +(-0.4,0.4) -- +(0.4,0.4) -- +(0.8,0) -- +(0.4,-0.4) -- +(-0.4,-0.4) --
+(-0.4,0.4);
\draw (#1) ++(0.8,0) -- ++(+0.3,0) coordinate (#1-o) ;
\draw (#1) ++ (-0.4,0.2) -- +(-0.3,0) coordinate (#1-i1) ;
\draw (#1) ++(-0.4,-0.2) -- +(-0.3,0) coordinate (#1-i2) ;
}

\newcommand\product[1]{\basicproduct{#1}{$+\Pi$}}
\newcommand\negativeproduct[1]{\basicproduct{#1}{$-\Pi$}}


%%BEGIN INTERNAL
\newcommand\basicsummer[1]{
coordinate (#1) {};
\draw (#1) +(-0.5,0.6) -- +(0.5,0) -- +(-0.5,-0.6) -- +(-0.5,0.6);
\draw (#1) ++(0.5,0) -- +(+0.3,0) coordinate (#1-o) ;
}

\newcommand\basicintegrator[1]{
coordinate (#1) {};
\draw (#1) +(-0.5,0.6) -- +(0.5,0) -- +(-0.5,-0.6) -- +(-0.5,0.6);
\draw (#1) ++(0.5,0) -- +(+0.3,0) coordinate (#1-o) ;
\draw (#1) +(-0.5,-0.6) -| +(-0.8,0.6) -- +(-0.5,0.6) ;
 \draw (#1) ++(-0.65,0.6) -- +(0,+0.3) coordinate (#1-init) ;
}

\newcommand\splitfour[2]{
\draw (#2) ++ (#1,0.4) -- +(-0.3,0) coordinate (#2-i1) ;
\draw (#2) ++ (#1,0.15) -- +(-0.3,0) coordinate (#2-i2) ;
\draw (#2) ++ (#1,-0.15) -- +(-0.3,0) coordinate (#2-i3) ;
\draw (#2) ++ (#1,-0.4) -- +(-0.3,0) coordinate (#2-i4) ;
}

\newcommand\splitthree[2]{
\draw  (#2) ++(#1,0.3) -- +(-0.3,0) coordinate (#2-i1) ;
\draw  (#2) ++(#1,0) -- +(-0.3,0) coordinate (#2-i2) ;
\draw  (#2) ++(#1,-0.3) -- +(-0.3,0) coordinate (#2-i3) ;
}

\newcommand\splittwo[2]{
\draw  (#2) ++(#1,0.2) -- +(-0.3,0) coordinate (#2-i1) ;
\draw  (#2) ++(#1,-0.2) -- +(-0.3,0) coordinate (#2-i2) ;
}

\newcommand\splitone[2]{
\draw (#2) ++(#1,0) -- +(-0.3,0) coordinate (#2-i1) ;
}

%% END INTERNAL

\newcommand\summerfour[1]{
\basicsummer{#1}
\splitfour{-0.5}{#1}
}

\newcommand\summerthree[1]{
\basicsummer{#1}
\splitthree{-0.5}{#1}
}

\newcommand\summertwo[1]{
\basicsummer{#1}
\splittwo{-0.5}{#1}
}

\newcommand\summerone[1]{
\basicsummer{#1}
\splitone{-0.5}{#1}
}


\newcommand\integratorfour[1]{
\basicintegrator{#1}
\splitfour{-0.8}{#1}
<}


\newcommand\integratorthree[1]{
\basicintegrator{#1}
\splitthree{-0.8}{#1}
}


\newcommand\integratortwo[1]{
\basicintegrator{#1}
\splittwo{-0.8}{#1}
}


\newcommand\integratorone[1]{
\basicintegrator{#1}
\splitone{-0.8}{#1}
}

%%%%% FIN MACROS




%%%%%%%%%%%%%%%%%%%%%%%%%%%%%%%%%%%%%%%%%%%%%%%%%
%
%
%                      Trucs d'Amaury
%
%
%%%%%%%%%%%%%%%%%%%%%%%%%%%%%%%%%%%%%%%%%%%%%%%%%


%----------------
%--- ref (AP) ---
%----------------

% \newcommand{\defref}[1]{Definition~\ref{#1}}
\newcommand{\lemref}[1]{Lemma~\ref{#1}}
\newcommand{\thref}[1]{Theorem~\ref{#1}}
\newcommand{\amaurycorref}[1]{Corollary~\ref{#1}}
 \newcommand{\figref}[1]{Figure~\ref{#1}}
% \newcommand{\figrefmult}[1]{Figures~{#1}}
% \newcommand{\secref}[1]{Section~\ref{#1}}
% %\renewcommand{\algref}[1]{Algorithm~\ref{#1}}
 \newcommand{\propref}[1]{Proposition~\ref{#1}}
% \newcommand{\exref}[1]{Example~\ref{#1}}
\newcommand{\remref}[1]{Remark~\ref{#1}}


 %-------------------
%--- polynomials ---
%-------------------

\newcommand{\sigmap}[1]{{\Sigma{#1}}}
\newcommand{\degp}[1]{{\operatorname{deg}(#1)}}
\newcommand{\poly}{\operatorname{poly}}


 %----------------------
%--- misc functions ---
%----------------------

\newcommand{\sgn}[1]{\operatorname{sgn}(#1)}
\newcommand{\intp}{\operatorname{int}}
 \newcommand{\fracp}{\operatorname{frac}}
\newcommand{\fiter}[2]{#1^{[#2]}}
\newcommand{\frestrict}[2]{#1\restriction_{#2}}
\newcommand{\indicator}[1]{\mathds{1}_{#1}}
\newcommand{\idfun}{\operatorname{id}}
% \newcommand{\floor}[1]{\left\lfloor#1\right\rfloor}
 \newcommand{\ceil}[1]{\left\lceil#1\right\rceil}
\newcommand{\round}[1]{\left\lfloor#1\right\rceil}
\DeclareMathOperator\erf{erf}


%------------
%--- sets ---
%------------

\newcommand{\A}{\mathbb{A}}
%\newcommand{\D}{\mathbb{D}}
\newcommand{\Ns}{\mathbb{N}^*}
%\newcommand{\C}{\mathbb{C}}
%\newcommand{\K}{\mathbb{K}}
% \MAT{height}{[width]}{[field]}
\newcommand{\MAT}[3]{M_{#1\ifthenelse{\equal{#2}{}}{}{,#2}}\ifthenelse{\equal{#3}{}}{}{\left(#3\right)}}

%\newcommand{\Rcomp}{\R_{c}}
\newcommand{\Rp}{\R_{+}}
\newcommand{\Rs}{\R^{*}}
\newcommand{\Rps}{\R_{+}^{*}}

%\newcommand{\Analytic}{C^\omega}
%\newcommand{\PTIME}{\ensuremath{\operatorname{P}}}
%\newcommand{\EXPTIME}{\ensuremath{\operatorname{EXPTIME}}}
%\newcommand{\NP}{\ensuremath{\operatorname{NP}}}
%\newcommand{\NPTIME}{\NP}
\newcommand{\FP}{\ensuremath{\operatorname{FP}}}
%\newcommand{\PSPACE}{\ensuremath{\operatorname{PSPACE}}}
% \newcommand{\FPSPACE}{\ensuremath{\operatorname{FPSPACE}}}



%%%%%%%%%%%%%%%%%%%%%%%%%%%%%%%%%%%%%%%%%%%%%%%%%
%
%
%                      Trucs de Quentin
%
%
%%%%%%%%%%%%%%%%%%%%%%%%%%%%%%%%%%%%%%%%%%%%%%%%%


\newcommand{\fct}[5]{\ensuremath{#1 : \left\{\begin{array}{c c c}
#2 & \rightarrow & #3\\
#4 & \mapsto & #5
                                             \end{array}\right.}}

%Nouvelle fraction, plus jolie
\newcommand{\f}[2]{	
	\mathchoice%
		{\dfrac{#1}{#2}}
    	{\dfrac{#1}{#2}}
		{\frac{#1}{#2}}
		{\frac{#1}{#2}}
}


%%Ensembles et combinatoires

%Ensemble avec propriété à drotie \{x | blabla\}
\newcommand{\enstq}[2]{\left\{ \left.#1\ \vphantom{#2}\right|\ #2  \right\}}
%meme chose mais avec des crochets
\newcommand{\crotq}[2]{\left[ \left.#1\ \vphantom{#2}\right|\ #2  \right]}
%même chose mais avec des brackets
\newcommand{\bratq}[2]{\left\langle \left.#1\ \vphantom{#2}\right|\ #2  
  \right\rangle}

\newcommand{\intn}[2]{\ensuremath{
		\left\llbracket\, #1\ ;\ #2\,\right\rrbracket
              }}


%%%%%%%%%%%%%%%%%%%%%%%%%%%%%%%%%%%%%%%%%%%%%%%%%
%
%
%               ENVIRONNEMENTS THEOREMS & CO
%
%
%%%%%%%%%%%%%%%%%%%%%%%%%%%%%%%%%%%%%%%%%%%%%%%%%


%----------------
%--- pas dans lipics ---
%-----------------


\ifnotSLIDE{
% commente car bug Fevrier 2026: \ifnotTDEXAM{\spnewtheorem{solution}{Solution}{\bfseries}{\rmfamily}}
\spnewtheorem{openproblem}{Open Problem}{\bfseries}{\rmfamily}
\spnewtheorem{remarkolivier}{Commentaire d'Olivier: }{\bfseries}{\rmfamily}
\spnewtheorem{postulate}{Postulate}{\bfseries}{\rmfamily}
}

%----------------
%--- dans lipics ---
%-----------------
%
\newenvironment{exercice}[1]{\begin{exercise} \label{exo:#1-stt}}{\end{exercise}}
\newenvironment{exerciced}[1]{\begin{exercisee}\label{exo:#1-stt}}{\end{exercisee}}
\newenvironment{exercicec}[1]{\begin{exercise} \label{exo:#1-stt} (\LANGUE{solution on}{corrigé} page \pageref{exo:#1-cor})
}{\end{exercise}}
\newenvironment{exercicedc}[1]{\begin{exercisee} \label{exo:#1-stt} (\LANGUE{solution on}{corrigé} page \pageref{exo:#1-cor})
}{\end{exercisee}}

\newenvironment{slideproposition}{{\bf Proposition}}{}

\ifnotSLIDE{
\nolipics{
\nolncs{\spnewtheorem{theorem}{\LANGUE{Theorem}{Théorème}}{\bfseries}{\rmfamily}}
\nolncs{\spnewtheoremi{definition}{\LANGUE{Definition}{Définition}}{\bfseries}{\rmfamily}{theorem}}
\nolncs{\spnewtheoremi{lemma}{\LANGUE{Lemma}{Lemme}}{\bfseries}{\rmfamily}{theorem}}
\nolncs{\spnewtheoremi{corollary}{\LANGUE{Corollary}{Corollaire}}{\bfseries}{\rmfamily}{theorem}}
\nolncs{\spnewtheoremi{proposition}{Proposition}{\bfseries}{\rmfamily}{theorem}}
\nolncs{\spnewtheoremi{example}{\LANGUE{Example}{Exemple}}{\bfseries}{\rmfamily}{theorem}}
%\spnewtheorem{sketch}{{\bf \LANGUE{Sketch of proof}{Démonstration (principe)}}}{\bfseries}{\rmfamily}
\nolncs{\spnewtheoremi{remark}{\LANGUE{Remark}{Remarque}}{\bfseries}{\rmfamily}{theorem}}
\nolncs{\spnewtheorem{exercise}{\LANGUE{Exercise}{Exercice}}{\bfseries}{\rmfamily}}
\spnewtheorem{exercisee}{\LANGUE{*Exercise}{*Exercice}}{\bfseries}{\rmfamily}
\spnewtheorem{summary}{\LANGUE{Summary}{Résumé}}{\bfseries}{\rmfamily}
 %\spnewtheorem{exerciceenglish}{Exercise}{\bfseries}{\rmfamily}
 %
 \newenvironment{sketch}{{\bf \LANGUE{Sketch of proof}{Démonstration (principe)}}:}{\hfill $\Box$ }
\nolncs{\spnewtheorem{conjecture}{\LANGUE{Conjecture}{Conjecture}}{\bfseries}{\rmfamily}}
 %%
 \spnewtheorem{theoremaprouver}{\LANGUE{Theorem TO BE PROVED}{Théorème (A PROUVER)}}{\bfseries}{\rmfamily}
  \spnewtheorem{conjectureaprouver}{\LANGUE{Conjecture TO BE VERIFIED}{Conjecture (A VERIFIER)}}{\bfseries}{\rmfamily}
 }}
            
 \makeatletter
 \ifnotSLIDE{
  \ifnotTDEXAM{
\@ifundefined{proof}{
	\newenvironment{proof}{{\bf \LANGUE{Proof}{Démonstration}}:}{\hfill $\Box$ 
	}{}
	}}}
\makeatother

 
%%%%%%%%%%%%%%%%%%%%%%%%%%%%%%%%%%%%%%%%%%%%%%%%%
%
%
%                      Macros dont pas fier du tout
%
%
%%%%%%%%%%%%%%%%%%%%%%%%%%%%%%%%%%%%%%%%%%%%%%%%%



\newcommand\equations[1]{
\begin{array}{lll}
#1
\end{array}
}



\newcommand\systeme[1]{
\left\{
\begin{array}{lll} 
#1
\end{array}
\right.
}



%------------------
%--- shorthands (AP)---
%------------------
\newcommand{\mtt}[1]{\mathtt{#1}}
\newcommand{\ovl}[1]{\overline{#1}}
\newcommand{\panic}[1]{\textcolor{red}{#1}}
%\NewEnviron{epanic}{{\color{red}\BODY}}
\newcommand{\idest}{\emph{i.e.\thinspace}}


%------------------
%--- mes trucs  ---
%------------------


\newcommand\implique{\Rightarrow}
\newcommand\ssi{\Leftrightarrow}
\newcommand\ac[1]{\{#1\}}
\newcommand\zero{\vectorl{0}}
\newcommand\un{\vectorl{1}}

\newcommand\technique[1]{{\scriptsize #1}}


%% Preuve and Co

\newcommand\sensdirect{ Direction $\Rightarrow$. }
\newcommand\sensindirect{ Direction $\Leftarrow$. }




\newcommand\donnee{\mbox{<donnée>}}



% pour aller vite

\newcommand\gM{\mathfrak{M}}
\newcommand\gMM{(M,f_1,\cdots,f_u,r_1,\cdots,r_v)}
%
\newcommand\mygM{\mK}
\newcommand\mygMM{\left ( \K, \{op_i\}_{i\in I}, r_1 \ldots,
r_\ell, \zero,\un \right )}
\newcommand\myM{\K} 

\newcommand\gMMM{(M,f_1,\cdots,f_u,c_1,\cdots,c_k,r_1,\cdots,r_v)}
\newcommand\gMMMM{(f_1,\cdots,f_u,r_1,\cdots,r_v)}
\newcommand\gMME{(M,f_1,\cdots,f_u,f'_1,\cdots,f'_\ell,r_1,\cdots,r_v)}
\newcommand\gMMS{(f_1,\cdots,f_u,f'_1,\cdots,f'_\ell,r_1,\cdots,r_v)}

\newcommand\bool{\mathfrak{B}}
\newcommand\gbool{(\{0,1\},\zero,\un,=)}

\newcommand\osg{\overline{sg}}
\newcommand\moins{\mathrel{\dot{-}}}


%%% TRUC DE BARWISE


\newcommand\UM{\kU_{\kM}}
\newcommand\VV{\mathbb{V}}
% Hereditarily finite
\newcommand\IHF{\mathbb{I}HF}
\newcommand\IHP{\IHYP}
\newcommand\IHYP{\mathbb{I}HYP}

\providecommand\citep[1]{\cite{#1}}

%% Cofinali


% probas
% défini dans amsmath
\renewcommand\Pr{\cp{Pr}}
\newcommand\Var{\cp{Var}}


% -- mesures
%\newcommand\card{card}
\newcommand\taille{\LANGUE{size}{taille}}
\newcommand\entree{e}
\newcommand\profondeur{\LANGUE{depth}{profondeur}}


%%%%%%%%%%%%%%%%%%%%%%%%%%%%%%%%%%%%%%%%%%%%%%%%%%%%%%%%%%%%%%%%%%%%%%%%%%%%
%                 MACHINES DE TURING (graphique)
%%%%%%%%%%%%%%%%%%%%%%%%%%%%%%%%%%%%%%%%%%%%%%%%%%%%%%%%%%%%%%%%%%%%%%%%%%%%



% Configuration de machine de Turing
\newcommand\configTuring[3]{#1\textbf{#2}#3}


%%% DESSIN CROIX PORU AIDER  A DESSINER

\newcommand\croix{  +(-1,-1) -- +(1,1) +(-1,1) --
  +(1,-1) }

%% DESSIN ACCOLADE

\newcommand\ACCOLADE[3]{
%#1: coordonnées de départ
%#2: coordonnées arrivée
%#3: texte
\draw[decorate,decoration={brace}]  (#1) -- (#2)
 node[midway,above=0.1cm] {#3};
}

%%%


\def\lcase{3.5 ex}
\def\hcase{3.5ex}
\tikzstyle{case} = []
%[draw,minimum width = \lcase,minimum height = 3.6ex,baseline]
\newcommand\dcase{ +(-0.5*\lcase,-0.5*\hcase) rectangle +(0.5*\lcase,0.5*\hcase)}
%\newcommand\lettre[1]{\node [name=l,case,on chain=1] {#1};}


\newcommand\RUBANMACHINEDETURING[6]{
%#1: texte
%#2: position de la tête
%#3: nb de blancs à gauche
%#4: nb de cases totales
%#5: texte au dessus du ruban
%#6: position du coin
    \draw (#6) +(\lcase,4ex ) node {#5};
    \draw (#6) +(-0.6,0)  node  {\ldots};  
    \draw (#6) node[coordinate] (t) {};
    \foreach \x in {1,2,...,#3} {
         \draw (t) node [case]   {$\blanc$} \dcase;
         \draw (t) ++(\lcase,0) node [coordinate] (t) {};
       };
   \StrLen{#1}[\Resultat] % ATTENTION IL FAUT CETTE SYNTAXE POUR FAIRE
                         % EQUIVALENT D UN EDEF
  
\foreach \x in {1,...,\Resultat} 
             { 
        \draw (t) node [case]   {\StrChar{#1}{\x}} \dcase;
        \draw (t) ++(\lcase,0) node [coordinate] (t) {};
             }
\pgfmathtruncatemacro{\z}{#4 - \Resultat - #3}

  \foreach \x in {1,2,..., \z} {
        \draw (t) node [case]   {$\blanc$} \dcase;
        \draw (t) ++(\lcase,0) node [coordinate] (t) {};
}
%% Convention: graphiquement, la case numéro n se trouve à (n*\lcase,0)
%% en comptant les case à partir de $0$

   \draw  (t) +(0.1,0) node [name=r] {\ldots}; 
}

\newcommand\RUBANMACHINEDETURINGtexte[6]{
%#1: texte
%#2: position de la tête
%#3: nb de blancs à gauche
%#4: nb de cases totales (équivalent)
%#5: texte au dessus du ruban
%#6: position du coin
    \draw (#6) +(\lcase,4ex ) node {#5};
    \draw (#6) +(-0.6,0)  node  {\ldots};  
    \draw (#6) node[coordinate] (t) {};
   \foreach \x in {1,2,...,#3} {
        \draw (t) node [case]   {$\blanc$} \dcase;
        \draw (t) ++(\lcase,0) node [coordinate] (t) {};
      };
   \draw (#6) ++(-2*\lcase,0.5*\lcase)-- +(\lcase*#4,0);
   \draw (#6) ++(-2*\lcase,-0.5*\lcase) -- +(\lcase*#4,0);
   
    \draw (t) ++(-0.5*\lcase,0) node [case,anchor=west]  (te) {#1};
    \draw (te.east)  node [anchor=west,coordinate] (t) {};

  \draw (t) ++(0.2*\lcase,0) node [coordinate] (t) {};
   \foreach \x in {1,2,...,#3} {
        \draw (t) node [case]   {$\blanc$} \dcase;
        \draw (t) ++(\lcase,0) node [coordinate] (t) {};
      };

    \draw  (t) +(0.1,0) node [anchor=west] {\ldots}; 
}


\newcommand\RUBANMACHINEDETURINGTETE[6]{
%
% dessine la tete
%
% output: (tete)
%
\RUBANMACHINEDETURING{#1}{#2}{#3}{#4}{#5}{#6}
% TETE
     \draw (#6) + ( #2 * \lcase + #3  *\lcase,- 0.5*\lcase )
     node[coordinate] (k) {};
     \draw[->] (k) +(0,-0.5) -- (k);
     \draw (k) +(0,-0.5) node[coordinate] (tete) {};
}


\newcommand\RUBANMACHINEDETURINGTETEtexte[6]{
%
% dessine la tete
%
% output: (tete)
%
\RUBANMACHINEDETURINGtexte{#1}{#2}{#3}{#4}{#5}{#6}
% TETE
     \draw (#6) + ( #2 * \lcase + #3  *\lcase,- 0.5*\lcase )
     node[coordinate] (k) {};
     \draw[->] (k) +(0,-0.5) -- (k);
     \draw (k) +(0,-0.5) node[coordinate] (tete) {};
}


\newcommand\MACHINEDETURING[6]{
%#1: texte
%#2: position de la tête
%#3: état
%#4: nb de blancs à gauche
%#5: nb de cases totales
%#6: texte en haut du ruban
%\begin{tikzpicture}[node distance=-0.15mm]

\RUBANMACHINEDETURINGTETE{#1}{#2}{#4}{#5}{#6}{0,0}
%
%% ETAT:
     \draw (tete) node [name=etat,draw,circle,anchor=north]  {#3};
     \draw (etat) -- (tete);
%\end{tikzpicture}
}

\newcommand\MACHINEDETURINGtexte[6]{
%#1: texte
%#2: position de la tête
%#3: état
%#4: nb de blancs à gauche
%#5: nb de cases totales
%#6: texte en haut du ruban
%\begin{tikzpicture}[node distance=-0.15mm]

\RUBANMACHINEDETURINGTETEtexte{#1}{#2}{#4}{#5}{#6}{0,0}
%
%% ETAT:
     \draw (tete) node [name=etat,draw,circle,anchor=north]  {#3};
     \draw (etat) -- (tete);
%\end{tikzpicture}
}

\newcommand\MACHINEDETURINGTROISRUBANS[9]{
%#1: texte1
%#2: position de la tête 1
%#3: texte2
%#4: position de la tête 2
%#5: texte3
%#6: position de la tête 3
%#7: état
%#8: nb de blancs à gauche 1
%#9: nb de cases totales
%\begin{tikzpicture}[node distance=-0.15mm]

\RUBANMACHINEDETURINGTETE{#1}{#2}{#8}{#9}{\LANGUE{tape}{ruban} $1$}{0,0}
\draw (tete) node[coordinate] (tete1) {};
\draw (tete1) +(8,0) node[coordinate] (tete1b) {};
\draw (tete1) -| (tete1b) ;

\RUBANMACHINEDETURINGTETE{#3}{#4}{#8}{#9}{\LANGUE{tape}{ruban} $2$}{0,-2}
\draw (tete) node[coordinate] (tete2) {};
\draw (tete2) +(-8.5,0) node[coordinate] (tete2b) {};
\draw (tete2) -| (tete2b) ;

\RUBANMACHINEDETURINGTETE{#5}{#6}{#8}{#9}{\LANGUE{tape}{ruban} $3$}{0,-4}
\draw (tete) node[coordinate] (tete3) {};

%% ETAT:
     \draw (tete3) node [name=etat,draw,circle,anchor=north]  {#7};
     \draw (etat) -- (tete3);
    \draw  (etat) -|  (tete2b);
\draw  (etat) -|  (tete1b);
%\end{tikzpicture}
}

\newcommand\MACHINEDETURINGTROISRUBANSENGLISH[9]{
%#1: texte1
%#2: position de la tête 1
%#3: texte2
%#4: position de la tête 2
%#5: texte3
%#6: position de la tête 3
%#7: état
%#8: nb de blancs à gauche 1
%#9: nb de cases totales
%\begin{tikzpicture}[node distance=-0.15mm]

\RUBANMACHINEDETURINGTETE{#1}{#2}{#8}{#9}{tape $1$}{0,0}
\draw (tete) node[coordinate] (tete1) {};
\draw (tete1) +(8,0) node[coordinate] (tete1b) {};
\draw (tete1) -| (tete1b) ;

\RUBANMACHINEDETURINGTETE{#3}{#4}{#8}{#9}{tape $2$}{0,-2}
\draw (tete) node[coordinate] (tete2) {};
\draw (tete2) +(-8.5,0) node[coordinate] (tete2b) {};
\draw (tete2) -| (tete2b) ;

\RUBANMACHINEDETURINGTETE{#5}{#6}{#8}{#9}{tape $3$}{0,-4}
\draw (tete) node[coordinate] (tete3) {};

%% ETAT:
     \draw (tete3) node [name=etat,draw,circle,anchor=north]  {#7};
     \draw (etat) -- (tete3);
    \draw  (etat) -|  (tete2b);
\draw  (etat) -|  (tete1b);
%\end{tikzpicture}
}



\newcommand\MACHINEDETURINGTROISRUBANStexte[9]{
%#1: texte1
%#2: position de la tête 1
%#3: texte2
%#4: position de la tête 2
%#5: texte3
%#6: position de la tête 3
%#7: état
%#8: nb de blancs à gauche 1
%#9: nb de cases totales
% \begin{tikzpicture}[node distance=-0.15mm]

\RUBANMACHINEDETURINGTETEtexte{#1}{#2}{#8}{#9}{\LANGUE{tape}{ruban} $1$}{0,0}
\draw (tete) node[coordinate] (tete1) {};
\draw (tete1) +(8,0) node[coordinate] (tete1b) {};
\draw (tete1) -| (tete1b) ;

\RUBANMACHINEDETURINGTETEtexte{#3}{#4}{#8}{#9}{\LANGUE{tape}{ruban} $2$}{0,-2}
\draw (tete) node[coordinate] (tete2) {};
\draw (tete2) +(-8.5,0) node[coordinate] (tete2b) {};
\draw (tete2) -| (tete2b) ;

\RUBANMACHINEDETURINGTETEtexte{#5}{#6}{#8}{#9}{\LANGUE{tape}{ruban} $3$}{0,-4}
\draw (tete) node[coordinate] (tete3) {};

%% ETAT:
     \draw (tete3) node [name=etat,draw,circle,anchor=north]  {#7};
     \draw (etat) -- (tete3);
    \draw  (etat) -|  (tete2b);
\draw  (etat) -|  (tete1b);
%\end{tikzpicture}
}


\newcommand\MACHINEDETURINGDEUXRUBANStexte[9]{
%#1: texte1  
%#2: position de la tête 1 
%#3: texte2 inutilisée
%#4: position de la tête 2 inutilisée
%#5: texte3
%#6: position de la tête 3
%#7: état
%#8: nb de blancs à gauche 1
%#9: nb de cases totales
% \begin{tikzpicture}[node distance=-0.15mm]

\RUBANMACHINEDETURINGTETEtexte{#1}{#2}{#8}{#9}{\LANGUE{tape}{ruban} $1$}{0,0}
\draw (tete) node[coordinate] (tete1) {};
\draw (tete1) +(8,0) node[coordinate] (tete1b) {};
\draw (tete1) -| (tete1b) ;

% \RUBANMACHINEDETURINGTETEtexte{#3}{#4}{#8}{#9}{\LANGUE{tape}{ruban} $2$}{0,-1.5}
% \draw (tete) node[coordinate] (tete2) {};
% \draw (tete2) +(-5.5,0) node[coordinate] (tete2b) {};
% \draw (tete2) -| (tete2b) ;

\RUBANMACHINEDETURINGTETEtexte{#5}{#6}{#8}{#9}{\LANGUE{tape}{ruban} $2$}{0,-1.5}
\draw (tete) node[coordinate] (tete3) {};

%% ETAT:
     \draw (tete3) node [name=etat,draw,circle,anchor=north]  {#7};
     \draw (etat) -- (tete3);
%   \draw  (etat) -|  (tete2b);
\draw  (etat) -|  (tete1b);
%\end{tikzpicture}
}


\def\argdix{2}
\def\argonze{25}
\newcommand\MACHINEDETURINGQUATRERUBANStextep[9]{
%#1: texte1
%#2: position de la tête 1
%#3: texte2
%#4: position de la tête 2
%#5: texte3
%#6: position de la tête 3
%#7: texte4
%#8: position de la tête 3
%#9: état
%#10: nb de blancs à gauche 1
%#11: nb de cases totales
% \begin{tikzpicture}[node distance=-0.15mm]

%% BUG sur #10 devenu \argdix
%% BUG SUR #11 devenu \argonze

\RUBANMACHINEDETURINGTETEtexte{#1}{#2}{\argdix}{\argonze}{\LANGUE{tape}{ruban} $1$}{0,0}
% \draw (tete) node[coordinate] (tete1) {};
% \draw (tete1) +(8,0) node[coordinate] (tete1b) {};
% \draw (tete1) -| (tete1b) ;

\RUBANMACHINEDETURINGTETEtexte{#3}{#4}{\argdix}{\argonze}{\LANGUE{tape}{ruban} $2$}{0,-1.5}
% \draw (tete) node[coordinate] (tete2) {};
% \draw (tete2) +(-5.5,0) node[coordinate] (tete2b) {};
% \draw (tete2) -| (tete2b) ;

\RUBANMACHINEDETURINGTETEtexte{#5}{#6}{\argdix}{\argonze}{\LANGUE{tape}{ruban} $3$}{0,-3.2}
\draw (tete) node[coordinate] (tete3) {};

\RUBANMACHINEDETURINGTETEtexte{#7}{#8}{\argdix}{\argonze}{\LANGUE{tape}{ruban} $4$}{0,-4.9}
\draw (tete) node[coordinate] (tete4) {};

%% ETAT:
     \draw (tete4) node [name=etat,draw,circle,anchor=north]  {#9};
%      \draw (etat) -- (tete3);
%     \draw  (etat) -|  (tete2b);
% \draw  (etat) -|  (tete1b);
%\end{tikzpicture}
}


%%%%%%%%%%%%%%%%%%%%%%%%%%%%%%%%%%%%%%%%%%%%%%%%%%%%%%%%%%%%%%%%%%%%%%%%%%%%
%                 MACHINES DE TURING (automate)
%%%%%%%%%%%%%%%%%%%%%%%%%%%%%%%%%%%%%%%%%%%%%%%%%%%%%%%%%%%%%%%%%%%%%%%%%%%%



\newcommand\dec{1ex}
\newcommand\myACCOLADE[4]{
\ACCOLADE{#1*\lcase+#4*\lcase-0.5*\lcase,0.5*\lcase+\dec}{#1*\lcase+#4*\lcase-0.5*\lcase+#3*\lcase,0.5*\lcase+\dec}{#2}
}


\newcommand\M{\Sigma}


\newcommand\sommet[2]{\node[state] (#1) #2 {$#1$};}
\newcommand\sommetinitial[2]{\node[state,initial] (#1) #2 {$#1$};}
\newcommand\sommetfinal[2]{\node[state,accepting] (#1) #2 {$#1$};}
\newcommand\gtransition[5]{\draw (#1) edge node {#2/#4 $#5$} (#3);}
\newcommand\gtransitionba[5]{\draw (#1) [loop above] edge node {#2/#4
    $#5$} (#3);}
\newcommand\gtransitionbb[5]{\draw (#1) [loop below] edge node {#2/#4
    $#5$} (#3);}
\newcommand\gtransitionbleft[5]{\draw (#1) [loop left] edge node {#2/#4
    $#5$} (#3);}
\newcommand\gtransitionbright[5]{\draw (#1) [loop right] edge node {#2/#4
    $#5$} (#3);}
\newcommand\gtransitionbl[5]{\draw (#1) [bend left] edge node {#2/#4
    $#5$} (#3);}
\newcommand\gtransitionbr[5]{\draw (#1) [bend right] edge node {#2/#4
    $#5$} (#3);}
\newcommand\gtransitionbadouble[9]{\draw (#1) [loop above] edge node {
\begin{tabular}{l}
#2/#4 $#5$ \\
  #7/#8 $#9$ \\
\end{tabular}
} (#3);}

\newcommand\SCALE{}

\newcommand\dessinmtp[2]{
\begin{center}
\begin{tikzpicture}[->,>=stealth',shorten >=1pt,auto,node distance=2.3cm,semithick,#2] 
#1
\end{tikzpicture}
\end{center}
}

\newcommand\dessinmt[1]{\dessinmtp{#1}{scale=1}}


\newcommand\ANIMATIONMT[1]{
\begin{overprint}
\begin{dessinp}{scale=0.52}
\foreach \av/\q/\b  [count=\xi]
in  
{#1}
{ 
\only<\xi| handout 0>{ 
 \StrLen{\av}[\Resintm] % ATTENTION IL FAUT CETTE SYNTAXE POUR FAIRE
                         % EQUIVALENT D UN EDEF
 \MACHINEDETURING{\av \b}{\Resintm}{$\q$}{5}{20}{}
}}
 \end{dessinp}
\end{overprint}
}

\newcommand\ANIMATIONMTd[1]{
\begin{overprint}
\begin{dessinp}{scale=0.52}
\foreach \av/\q/\b  [count=\xi]
in  
{#1}
{ 
\only<\xi| handout 0>{ 
 \StrLen{\av}[\Resintm] % ATTENTION IL FAUT CETTE SYNTAXE POUR FAIRE
                         % EQUIVALENT D UN EDEF
 \hspace{-1cm}\MACHINEDETURING{\av \b}{\Resintm}{$\q$}{5}{33}{}
}}
 \end{dessinp}
\end{overprint}
}

\def\MAXDET{20}

\newcommand\DIAGRAMMEESPACETEMPS[1]{
\draw node[coordinate] (tt) {};
\foreach \av/\q/\b  in  
{#1}
{ 

\RUBANMACHINEDETURING{\av{$\q$}\b}{0}{4}{\MAXDET}{}{tt};
\draw (tt) +(0,-\hcase) node[coordinate] (tt) {};
}
%\foreach \x in {0,...,19} {\draw (tt) + (\x*\lcase,0) node {$\vdots$};}
}



\newcommand\progun{
/q_0/0011, % initial
X/q_1/011,X0/q_1/11,X/q_2/0Y1,/q_2/X0Y1,
X/q_0/0Y1,XX/q_1/Y1,XXY/q_1/1, XX/q_2/YY, X/q_2/XYY,
XX/q_0/YY,XXY/q_3/Y,XXYY/q_3/B,XXYYB/q_4/B
}

\newcommand\DIAGRAMMEESPACETEMPSun{
\DIAGRAMMEESPACETEMPS
%% RECOPIE progun
{/q_0/0011, % initial
X/q_1/011,X0/q_1/11,X/q_2/0Y1,/q_2/X0Y1,
X/q_0/0Y1,XX/q_1/Y1,XXY/q_1/1, XX/q_2/YY, X/q_2/XYY,
XX/q_0/YY,XXY/q_3/Y,XXYY/q_3/B,XXYYB/q_4/B}
%% FIN RECOPIE
}

\newcommand\progdeux{
/q_0/0010, X/q_1/010, X0/q_1/10, X/q_2/0Y0,/q_2/X0Y0,
X/q_0/0Y0,XX/q_1/Y0,XXY/q_1/0,XXY0/q_1/%B,
}


\newcommand\DIAGRAMMEESPACETEMPSdeux{
\DIAGRAMMEESPACETEMPS
{
%% RECOPIE progdeux
/q_0/0010, X/q_1/010, X0/q_1/10, X/q_2/0Y0,/q_2/X0Y0,
X/q_0/0Y0,XX/q_1/Y0,XXY/q_1/0,XXY0/q_1/%B,
%% FIN RECOPIE
}
}

\newcommand\DIAGRAMMEESPACETEMPSdeuxvariante{
\def\memolcase{\lcase}
\def\memoMAXDET{20}
\def\MAXDET{15}
\def\lcase{6 ex}
\DIAGRAMMEESPACETEMPS
{
%% RECOPIE progdeux
/q_00/010, X/q_10/10, X0/q_11/0, X/q_20/Y0,/q_2X/0Y0,
X/q_00/Y0,XX/q_1Y/0,XXY/q_10/,XXY0/q_1B/%B,
%% FIN RECOPIE
}
\def\lcase{\memolcase}
\def\MAXDET{\memoMAXDET}
}


\newcommand\ANIMATIONun{
\ANIMATIONMT
%% RECOPIE progun
{/q_0/0011, % initial
X/q_1/011,X0/q_1/11,X/q_2/0Y1,/q_2/X0Y1,
X/q_0/0Y1,XX/q_1/Y1,XXY/q_1/1, XX/q_2/YY, X/q_2/XYY,
XX/q_0/YY,XXY/q_3/Y,XXYY/q_3/B,XXYYB/q_4/B}
%% FIN RECOPIE
}



\newcommand\DESSINALGO[3]{
% argument 1= ou
% argument 2= texte 
% argument 3= nom du noeud ex m
\draw node[draw,#1,minimum size=1cm] (#3) {#2};

% pour fitting box
\draw (#3.west) +(-0.2cm,0) node (#3c2) {};
\draw (#3.north) +(0,+0.2cm) node (#3c3) {};
\draw (#3.south) +(0,-0.2cm) node (#3c4) {};
}

\newcommand\DESSINALGOa[3]{
\DESSINALGO{#1}{#2}{#3}
\draw (#3.east) +(-0.1cm,0.2cm) node(#3a) {};
\draw[->] (#3a) -- ++(0.5cm,0) node (#3aa) {};
\draw (#3aa) node[anchor=west] (#3aaa) {\LANGUE{Accept}{Accepte}};
\draw (#3aaa.east) node(#3aaaa) {};
}

\newcommand\DESSINALGOe[4]{
\DESSINALGO{#1}{#2}{#3}
\draw (#3.east) +(-0.1cm,0.2cm) node(#3a) {};
\draw[->] (#3a) -- ++(0.5cm,0) node (#3aa) {};
\draw (#3aa) node[anchor=west] (#3aaa) {#4};
\draw (#3aaa.east) node(#3aaaa) {};
}


\newcommand\DESSINALGOar[3]{
\DESSINALGOa{#1}{#2}{#3}
\draw (#3.east) +(-0.1cm,-0.2cm) node(#3r) {};
\draw[->] (#3r) -- ++(0.5cm,0) node (#3rr) {};
\draw (#3rr) node[anchor=west] (#3rrr) {\LANGUE{Reject}{Refuse}};
\draw (#3rrr.east) node(#3rrrr) {};
}


\newcommand\WINDOWarg[7] {
%#1: position
%#2..7: contenu du tableau
\draw (#1) node[case]  {#2} \dcase;
\draw (#1) ++ (\lcase,0) node[case] {#3} \dcase;
\draw (#1) ++ (2*\lcase,0) node[case] {#4} \dcase;
\draw (#1) ++ (0,-\lcase) node[case] {#5} \dcase;
\draw (#1) ++ (\lcase,-\lcase) node[case] {#6} \dcase;
\draw (#1) ++ (2*\lcase,-\lcase) node[case] {#7} \dcase;
}

\newcommand\WINDOW[3]{
%#1: positioin
%#2: contenu
%#3: étiquette
\WINDOWarg{#1}
{\StrChar{#2}{1}}{\StrChar{#2}{2}}{\StrChar{#2}{3}}
{\StrChar{#2}{4}}{\StrChar{#2}{5}}{\StrChar{#2}{6}}
\draw (#1) +(-\lcase*1.2,-\lcase*0.5) node {#3};
}





%%%%%%%%%%%%%%%%%%%%%%%%%%%%%%%%%%%%%%%%%%%%%%%%%%%%%%%%%%%%%%%%%%%%%%%%%%%%
%                 LISTINGS
%%%%%%%%%%%%%%%%%%%%%%%%%%%%%%%%%%%%%%%%%%%%%%%%%%%%%%%%%%%%%%%%%%%%%%%%%%%%


% LISTINGS
\RequirePackage{listings}
 \lstset{language=Java,
 mathescape=true,
 showspaces=false,
 showstringspaces=false,
 frame=single,
morekeywords={recursive,then,div,function,integer,read,out,else,par,endpar,seq,endseq,and,not,read,write,repeat,until,undef,let,to,endfor,endif,endwhile,guess,in,parallel},
% basicstyle=\ttfamily\small,
% basewidth={1.1ex},
% keywordstyle=\bf\color{blue},
% %stringstyle=\bf\ttfamily\color{green},
% commentstyle=\bfseries\color{magenta},
 literate={:=}{{$\gets$ }}1 {<=}{{$\leq$}}2  {>=}{{$\geq$}}2
 {!=}{{$\neq$}}1 
 }
% \lstset{escapechar=\%}
\let\lst\lstinline
\def\beginlisting{\begin{lstlisting}}
\def\endlisting{\end{lstlisting}}  




\newenvironment{algo}[1]{
  \lstset{escapechar=\@}
  \def\input{$#1$}
  \def\out{out}
  \def\BEGIN{\textbf{read} \input }
 \def\BEGINN{\textbf{repeat} }
  \def\END{\textbf{until out!=undef}}
  \def\ENDD{\textbf{write \out}}
  \def\ENDT{\textbf{until $\ctl$ = $q_a$ or $\ctl$ = $q_r$}}
  \def\ENDR{\textbf{until $\ctl$ = $q_a$}}
  \def\ENDP{\textbf{until $\ctl$ = $q_a$}}
}{}
\newenvironment{algon}[1]{
  \lstset{numbers=left}
  \begin{algo}{#1}
}{\end{algo}}
\newcommand\FSM[3]{FSM(#1,\lst{#2},#3)}
\newcommand\FSMi[3]{FSM(#1,{#2},#3)}




%%%%%%%%%%%%%%%%%%%%%%%%%%%%%%%%%%%%%%%%%%%%%%%%%
%
%
%                      TRADUCTION EN COURS
%
%%%%%%%%%%%%%%%%%%%%%%%%%%%%%%%%%%%%%%%%%%%%%%%%%

\newcommand\ENTREPOT[1]{\NDLR{ENTREPOT ICI: #1}}
\newcommand\scite[2][]{\cite[#1]{#2}}
\newcommand\nospellcheck[1]{{#1}}
\newcommand\nd{nd}
\newcommand\NL{
 
}

\newenvironment{FRANCAIS}{
\small
\color{orange}
}
{\normalsize
\color{black}
}


\newenvironment{SEMIFRANCAIS}{
\small
\color{violet}
}
{\normalsize
\color{black}
}

\newcommand\QFRANCAIS[1]{
\begin{FRANCAIS}
#1
\end{FRANCAIS}
}


%%%%%%%%%%%%%%%%%%%%%%%%%%%%%%%%%%%%%%%%%%%%%%%%%
%
%
%                      Wikipedia
%
%%%%%%%%%%%%%%%%%%%%%%%%%%%%%%%%%%%%%%%%%%%%%%%%%

\providecommand\tightlist{}
\providecommand\Complex{\C}
\newenvironment{myalgie}{\begin{array}{lll}}{\end{array}}

\newcommand\nl{\ \\ }




%%%%%%%%%%%%%%%%%%%%%%%%%%%%%%%%%%%%%%%%%%%%%%%%%
%
%
%                      Pour citations précises
%
%%%%%%%%%%%%%%%%%%%%%%%%%%%%%%%%%%%%%%%%%%%%%%%%%



\newcommand\citeprecis[2]{{\cite[#1]{#2}}}
\newcommand\citeprecisinv[2]{\citeprecis{#2}{#1}}

\usepackage{csquotes}



%%%%%%%%%%%%%%%%%%%%%%%%%%%%%%%%%%%%%%%%%%%%%%%%%
%
%
%                      Travail sur couleurs
%
%   toutes les réponses pointent dans le même sens, et le schéma correspond exactement 
%. au profil perceptif d’un protanope** (déficit des cônes L / longueurs d’onde longues)
%  MAIS SEVERE
%
%%%%%%%%%%%%%%%%%%%%%%%%%%%%%%%%%%%%%%%%%%%%%%%%%


% ================================
% GROUPE 1 — Verts sombres
% ================================
\definecolor{deepgreen}{RGB}{0,160,40}
\definecolor{forestgreen}{RGB}{0,150,70}
\definecolor{mossgreen}{RGB}{0,140,90}
\definecolor{greenwave}{RGB}{0,150,90}
\definecolor{forestmist}{RGB}{0,160,90}
\definecolor{darkteal}{RGB}{0,130,70}

% ================================
% GROUPE 2 — Teal / verts bleutés
% ================================
\definecolor{seagreen}{RGB}{0,130,100}
\definecolor{freshteal}{RGB}{0,140,130}
\definecolor{greenteal}{RGB}{0,140,120}
\definecolor{blueteal}{RGB}{0,150,120}      % ancien "teal" (doublon corrigé)
\definecolor{lagoon}{RGB}{0,180,120}
\definecolor{greenmist}{RGB}{0,170,110}

% ================================
% GROUPE 3 — Turquoises / Aqua
% ================================
\definecolor{springgreen}{RGB}{0,190,100}
\definecolor{mintgreen}{RGB}{10,200,150}
\definecolor{oceanfoam}{RGB}{0,200,150}
\definecolor{aquagreen}{RGB}{30,210,170}
\definecolor{coastalaqua}{RGB}{40,210,180}
\definecolor{softaqua}{RGB}{40,210,190}

% ================================
% GROUPE 4 — Cyan / Aqua clairs
% ================================
\definecolor{emeraldteal}{RGB}{0,180,140}
\definecolor{deepcyan}{RGB}{0,180,160}
\definecolor{lightcyan}{RGB}{0,200,160}
\definecolor{aquamarine}{RGB}{0,200,160}
\definecolor{lightaqua}{RGB}{60,230,200}
\definecolor{mistblue}{RGB}{80,240,210}

% ================================
% GROUPE 5 — Bleus clairs
% ================================
\definecolor{paleaqua}{RGB}{80,200,190}
\definecolor{glacier}{RGB}{90,150,200}
\definecolor{skyblueA}{RGB}{60,150,220}      % ancien "skyblue"
\definecolor{lightsky}{RGB}{90,170,230}
\definecolor{iceblue}{RGB}{120,190,240}
\definecolor{frost}{RGB}{150,210,250}

% ================================
% GROUPE 6 — Bleus moyens
% ================================
\definecolor{steelblue}{RGB}{30,100,160}
\definecolor{cadetblue}{RGB}{70,120,180}
\definecolor{coolblue}{RGB}{110,170,220}
\definecolor{morningblue}{RGB}{140,200,240}
\definecolor{blueiceA}{RGB}{40,150,210}       % ancien "blueice" (doublon corrigé)
\definecolor{blueair}{RGB}{60,170,225}

% ================================
% GROUPE 7 — Azurs / Bleu saturé
% ================================
\definecolor{tealblue}{RGB}{0,110,130}
\definecolor{deepbluecyan}{RGB}{0,110,170}
\definecolor{azur}{RGB}{0,70,180}
\definecolor{brightazure}{RGB}{20,130,200}
\definecolor{lightazure}{RGB}{80,180,235}
\definecolor{skylight}{RGB}{100,200,245}

% ================================
% GROUPE 8 — Bleus sombres / violets
% ================================
\definecolor{nightblue}{RGB}{40,0,110}
\definecolor{slateblue}{RGB}{20,40,120}
\definecolor{violetA}{RGB}{60,50,160}         % ancien "violet" (distinct de violetB ci-dessous)
\definecolor{blueviolet}{RGB}{80,60,180}
\definecolor{purpleA}{RGB}{100,70,200}        % ancien "purple"
\definecolor{lightpurple}{RGB}{130,90,210}

% ================================
% GROUPE 9 — Violets saturés / indigos
% ================================
\definecolor{lilac}{RGB}{160,110,220}
\definecolor{indigoA}{RGB}{50,0,130}          % ancien "indigo"
\definecolor{deepindigo}{RGB}{70,10,150}
\definecolor{royalpurple}{RGB}{90,20,170}

% ================================
% GROUPE 10 — Fermeture cercle perceptif
% ================================
\definecolor{coolmint}{RGB}{0,130,70}
\definecolor{deepsea}{RGB}{0,130,70}

% ================================
% PALETTE DEUTÉRANOPE — noms uniques
% ================================
\definecolor{crimson}{RGB}{200,30,60}
\definecolor{brickred}{RGB}{180,40,40}
\definecolor{coral}{RGB}{220,90,70}
\definecolor{salmon}{RGB}{240,130,120}

\definecolor{magenta}{RGB}{200,40,150}
\definecolor{fuchsia}{RGB}{220,60,170}
\definecolor{violetB}{RGB}{140,70,190}        % ancien "violet"
\definecolor{indigoB}{RGB}{90,50,180}         % ancien "indigo"

\definecolor{royalblueA}{RGB}{50,80,200}      % ancien "royalblue"
\definecolor{azure}{RGB}{30,120,230}
\definecolor{skyblueB}{RGB}{80,160,230}       % second "skyblue"
\definecolor{tealB}{RGB}{0,150,170}           % second "teal"

\definecolor{marine}{RGB}{0,80,140}




%% Alternative

% === Générateur automatique fill/text pour protanopes ===
% Usage : \ProtaColorPair{mistblue}
% Suppose que mistblue est déjà défini par \definecolor{mistblue}{RGB}{...}

\newcommand{\ProtaColorPair}[1]{%
  % fill = couleur originale
  \colorlet{#1fill}{#1}%
  % text = version assombrie idéalement lisible
  \colorlet{#1text}{#1!80!black}%
}

%Utilisation
%\ProtaColorPair{mistblue}

%\textcolor{mistbluefill}{Couleur claire} \quad
%\textcolor{mistbluetext}{Couleur pour texte}

	\foreach \i/\name in {
	  1/deepgreen,
	  2/greenwave,
	  3/lagoon,
	  4/oceanfoam,
	  5/softaqua,
	  6/mistblue,
 	 7/lightsky,
	  8/blueiceA,   % <-- corrigé
	  9/deepcyan,
	  10/azur,
	  11/blueviolet,
	  12/royalpurple}
	  {\ProtaColorPair{\name}}
	  


\newcommand\ROUEPROTANOPEDOUZE{
	% Roue chromatique protanope — 12 couleurs
	\begin{tikzpicture}[scale=3]

	\foreach \i/\name/\nametext in {
	  1/deepgreen/deepgreen,
	  2/greenwave/greenwave,
	  3/lagoon/lagoontext,
	  4/oceanfoam/oceanfoam,
	  5/softaqua/softaquatext,
	  6/mistblue/mistbluetext,
 	 7/lightsky/lightskytext,
	  8/blueiceA/blueiceA,   % <-- corrigé
	  9/deepcyan/blueiceA,
	  10/azur/azur,
	  11/blueviolet/blueviolet,
	  12/royalpurple/royalpurple}
	{
	  % secteur coloré
 	 \filldraw[fill=\name, draw=black, line width=0.2pt]
	    ({90-30*\i}:1)
	    -- ({90-30*(\i-1)}:1)
	    -- (0,0) -- cycle;

 	 % label dans la bonne couleur
	  \node at ({90-30*(\i-0.5)}:1.25)
	    {\textcolor{\nametext}{\small \nametext}};
	}

	\end{tikzpicture}
}


\newcommand\ROUEPROTANOPEVINGT{
\begin{tikzpicture}[scale=3]

\foreach \i/\name/\nametext in {
  1/deepgreen/deepgreen,
  2/greenwave/greenwave,
  3/lagoon/lagoon,
  4/oceanfoam/oceanfoam,
  5/softaqua/softaquatext,
  6/mistblue/mistbluetext,
  7/lightsky/lightsky,
  8/blueiceA/blueiceA,      % corrigé
  9/deepcyan/deepcyan,
  10/azur/azur,
  11/brightazure/brightazure,
  12/skyblueA/skyblueA,     % corrigé
  13/coolblue/coolblue,
  14/glacier/glacier,
  15/marine/marine,
  16/violetA/violetA,      % corrigé
  17/blueviolet/blueviolet,
  18/purpleA/purpleA,      % corrigé
  19/royalpurple/royalpurple,
  20/indigoA/indigoA}      % corrigé
{
  % Secteur coloré
  \filldraw[fill=\name, draw=black, line width=0.2pt]
    ({90-18*\i}:1)
    -- ({90-18*(\i-1)}:1)
    -- (0,0) -- cycle;

  % Étiquette
  \node at ({90-18*(\i-0.5)}:1.25)
    {\textcolor{\nametext}{\small \nametext}};
}

\end{tikzpicture}
}

\newcommand\ROUEDEUTERANOPEDOUZE{
% Roue chromatique deutéranope — 12 couleurs
\begin{tikzpicture}[scale=3]

\foreach \i/\name in {
  1/crimson,
  2/brickred,
  3/coral,
  4/salmon,
  5/magenta,
  6/fuchsia,
  7/violetB,       % corrigé
  8/indigoB,       % corrigé
  9/royalblueA,    % corrigé
  10/azure,
  11/skyblueB,     % corrigé
  12/tealB}        % corrigé
{
  % Secteur coloré
  \filldraw[fill=\name, draw=black, line width=0.2pt]
    ({90-30*\i}:1)
    -- ({90-30*(\i-1)}:1)
    -- (0,0) -- cycle;

  % Label dans la couleur
  \node at ({90-30*(\i-0.5)}:1.25)
    {\textcolor{\name}{\small \name}};
}

\end{tikzpicture}
}



% === Macro intelligente pour foncer une couleur de façon lisible pour protanope ===
% Usage : \protadarken{mistblue}{40}  -> crée mistblue@dark
% Le 40 signifie : 40% de mélange avec du noir

\newcommand{\DarkenForText}[2]{%
  % #1 = nom de la couleur
  % #2 = pourcentage (40 conseillé si tu ne sais pas)
  %
  % On crée une nouvelle couleur #1@dark
  \colorlet{#1text}{black!#2!#1}%
}





%%%%

% Groupe 1
\DarkenForText{deepgreen}{20}
\DarkenForText{forestgreen}{20}
\DarkenForText{mossgreen}{20}
\DarkenForText{greenwave}{20}
\DarkenForText{forestmist}{20}
\DarkenForText{darkteal}{20}

% Groupe 2
\DarkenForText{seagreen}{20}
\DarkenForText{freshteal}{20}
\DarkenForText{greenteal}{20}
\DarkenForText{blueteal}{20}
\DarkenForText{lagoon}{20}
\DarkenForText{greenmist}{20}

% Groupe 3
\DarkenForText{springgreen}{20}
\DarkenForText{mintgreen}{20}
\DarkenForText{oceanfoam}{20}
\DarkenForText{aquagreen}{20}
\DarkenForText{coastalaqua}{20}
\DarkenForText{softaqua}{20}

% Groupe 4
\DarkenForText{emeraldteal}{20}
\DarkenForText{deepcyan}{20}
\DarkenForText{lightcyan}{20}
\DarkenForText{aquamarine}{20}
\DarkenForText{lightaqua}{20}
\DarkenForText{mistblue}{30}

% Groupe 5
\DarkenForText{paleaqua}{20}
\DarkenForText{glacier}{20}
\DarkenForText{skyblueA}{20}
\DarkenForText{lightsky}{20}
\DarkenForText{iceblue}{20}
\DarkenForText{frost}{20}

% Groupe 6
\DarkenForText{steelblue}{20}
\DarkenForText{cadetblue}{20}
\DarkenForText{coolblue}{20}
\DarkenForText{morningblue}{20}
\DarkenForText{blueiceA}{20}
\DarkenForText{blueair}{20}

% Groupe 7
\DarkenForText{tealblue}{20}
\DarkenForText{deepbluecyan}{20}
\DarkenForText{azur}{20}
\DarkenForText{brightazure}{20}
\DarkenForText{lightazure}{20}
\DarkenForText{skylight}{20}

% Groupe 8
\DarkenForText{nightblue}{20}
\DarkenForText{slateblue}{20}
\DarkenForText{violetA}{20}
\DarkenForText{blueviolet}{20}
\DarkenForText{purpleA}{20}
\DarkenForText{lightpurple}{20}

% Groupe 9
\DarkenForText{lilac}{20}
\DarkenForText{indigoA}{20}
\DarkenForText{deepindigo}{20}
\DarkenForText{royalpurple}{20}

% Groupe 10
\DarkenForText{coolmint}{20}
\DarkenForText{deepsea}{20}

% Palette deutéranope
\DarkenForText{crimson}{20}
\DarkenForText{brickred}{20}
\DarkenForText{coral}{20}
\DarkenForText{salmon}{20}

\DarkenForText{magenta}{20}
\DarkenForText{fuchsia}{20}
\DarkenForText{violetB}{20}
\DarkenForText{indigoB}{20}

\DarkenForText{royalblueA}{20}
\DarkenForText{azure}{20}
\DarkenForText{skyblueB}{20}
\DarkenForText{tealB}{20}

% Couleur supplémentaire
\DarkenForText{marine}{20}

%%%%%%%%%%%%%%%%%%%%%
%
%.   apx proof
%
%%%%%%%%%%%%%%%%%%%%%



% do
\newcommand\PREUVESENAPPENDIXCOMPATIBLE{
	\newenvironment{appendixproof}{{\bf \LANGUE{Proof}{Démonstration}}:}{\hfill $\Box$}
	\newtheorem{theoremrep}[theorem]{Theorem}
	\newtheorem{lemmarep}[theorem]{Lemma}
	\newtheorem{propositionrep}[theorem]{Proposition}
	\newtheorem{corollaryrep}[theorem]{Corollary}
	\newtheorem{summaryrep}[theorem]{Summary}	
}

\providecommand\apxparameter{}

%% proof inline
%\providecommand\apxparameter{appendix=inline}
%% proof in appendix
%\providecommand\apxparameter{}


\apxproofmode{
\usepackage[\apxparameter]{apxproof}
\theoremstyle{plain}

\newtheoremrep{theorem}{Theorem}
%\newtheoremrep{theoremconjecture}[theorem]{Theorem-To-Be-Proved=Conjecture}
%\newtheoremrep{lemmaconjecture}[theorem]{Conjecture-To-Be-Proved=Conjecture}
\newtheoremrep{claim}[theorem]{Theorem}
\newtheoremrep{proposition}[theorem]{Proposition}
\newtheoremrep{lemma}[theorem]{Lemma}
\newtheoremrep{corollary}[theorem]{Corollary}
%%bug mars 26: \newtheoremrep{summary}[theorem]{Summary}
% Du coup:
\newtheoremrep{resume}[theorem]{Summary}
%bug mars 26: \newtheoremrep{openproblem}[theorem]{Open Problem}
%\newtheoremrep{claim}[theorem]{Claim}
}


%\apxproofmodesubsection{
%%apx proof, pour que compteur correct
%\makeatletter
%\AtBeginDocument{%
%\let\axp@section\subsection
%  \let\axp@oldsection\subsection
%}
%\makeatother
%\renewcommand\appendixprelim{\section{Missing proofs}}
%}


\makeatletter
\apxproofmodesubsection{
\let\axp@section\subsection
%\let\axp@oldsection\subsection
\renewcommand\appendixprelim{\section{Proofs from main documents}}
}
\makeatother


\makeatletter
\newcommand\appendixapxproofsubsection{
\appendix
\let\apx@orig@appendix\appendix
\renewcommand{\appendix}{}
}
\makeatother


%%% Pour checker/vérifier que bon


\newcommand\docheck[1]{\textcolor{check}{#1}}

\usepackage{xcolor}
%\definecolor{chose}{RGB}{0,110,50}
%\definecolor{forestgreen}{RGB}{0,150,70}
%\definecolor{ctruc}{RGB}{164,164,164}
%\definecolor{oceanfoam}{RGB}{0,200,150}
\definecolor{check}{RGB}{160,160,0}






}


 
%%%%%%%%%%%%%%%%%%%%%%%%%%%%%%%%%%%%%%%%%%%%%%%%%
%
%
%               Déclaration de cours
%
%  (l'argument 1 = numéro sert pas)
%
%%%%%%%%%%%%%%%%%%%%%%%%%%%%%%%%%%%%%%%%%%%%%%%%%


\newcommand\COURS[3]{
  \expandafter\newcommand\csname numerocours#2\endcsname{#1}
  \expandafter\newcommand\csname nomcours#2\endcsname{#3}
}



\newcommand\NANCY[3]{
  \expandafter\newcommand\csname numerocours#2\endcsname{#1}
  \expandafter\newcommand\csname nomcours#2\endcsname{#3}
}


\newcommand\INFquatrecentdouze[3]{
  \expandafter\newcommand\csname numerocours#2\endcsname{#1}
  \expandafter\newcommand\csname nomcours#2\endcsname{#3}
}

\newcommand\INFcinqsixun[3]{
  \expandafter\newcommand\csname numerocours#2\endcsname{#1}
  \expandafter\newcommand\csname nomcours#2\endcsname{#3}
}



\newcommand\MPRI[3]{
  \expandafter\newcommand\csname numerocours#2\endcsname{#1}
  \expandafter\newcommand\csname nomcours#2\endcsname{#3}
}

\newcommand\COURSEINCLUDE[2][]{
  \renewcommand\IFNOTSUB[1]{}
  \renewcommand\MYCOURSE[1]{\chapter{\ifnotmodeDIFFUSIONFINALE{#1} \csname nomcours##1\endcsname}
    \ifnotmodeDIFFUSIONFINALE{\mychapter{Fichier: #2}}
  }
   \renewcommand\MYCOURSEnew[1]{\chapter{\ifnotmodeDIFFUSIONFINALE{#1} ##1 }
    \ifnotmodeDIFFUSIONFINALE{\mychapter{Fichier: #2}}
  }
%  \renewcommand\FINSUBMPRI{}
  \input{#2}
}

\newcommand\STRESSCOURSEINCLUDE[1]{
  \COURSEINCLUDE[!STRESS!]{#1}
  }

  \newcommand\subCOURSEINCLUDE[1] {
    \NDLR{BEGIN:  J'inclus/ Je répète ICI #1 }
    \input{INCLUDE/#1}
    \NDLR{END:  J'inclus/ Je répète ICI #1 }
}

%%%%%%%%%%%%%%%%%%%%%%%%%%%%%%%%%%%%%%%%%%%%%%%%%
%
%
%                      POUR INDICATIONS: MESSAGES
%
%
%%%%%%%%%%%%%%%%%%%%%%%%%%%%%%%%%%%%%%%%%%%%%%%%%


% ancienne version
% \newcommand\montruc[1]{\begin{center} {\fbox{\fbox{{#1}
%         }}} \end{center}}




\newcommand\montruc[2][]{
  \ifnotmodeDIFFUSIONFINALE{
  \begin{maboitetruc}[#1]{\danger} #2
  \end{maboitetruc}
%%  \todo[inline,color=green!20,caption={2do}, ssss#1]{\begin{minipage}{\textwidth-4pt}
%\todo[inline,color=green!5%,caption={2do}
%]{\begin{minipage}{\textwidth-4pt}
%    #2\end{minipage}}
}
}

\newcommand\montrucdanger[3][]{
  \ifnotmodeDIFFUSIONFINALE{
                    \montruc[#1]{\danger #2: #3}
                    }
}

\newcommand\montrucdangermargin[3][]{
  \ifnotmodeDIFFUSIONFINALE{
  \montruc[#1]{\danger #2: #3}
\marginpar{\danger #2: #3}
}
 }



%%%%%%%%%%%%%%%%%%%%%%%%%%%%%%%%%%%%%%%%%%%%%%%%%
%
%
%                      GESTION DES TRUCS MARQUES VERSIONS
%
%
%%%%%%%%%%%%%%%%%%%%%%%%%%%%%%%%%%%%%%%%%%%%%%%%%

\makeatletter
\renewcommand\beginmarkversion{\montrucdanger[colback=yellow!10]{=== begin NOT}{}\@Vs@sffbox{\@currenvir$>$}}
 \renewcommand\endmarkversion{\@Vs@sffbox{$<$\@currenvir}\montrucdanger[colback=yellow!10]{end=== NOT}{}}
\makeatother 


%%%%%%%%%%%%%%%%%%%%%%%%%%%%%%%%%%%%%%%%%%%%%%%%%
%
%
%                      VARIANTES DE MESSAGES
%
%
%%%%%%%%%%%%%%%%%%%%%%%%%%%%%%%%%%%%%%%%%%%%%%%%%

 \newcommand\NDLR[1]{\montrucdanger[colback=blue!20]{NDLR}{#1}}
 \newcommand\NLDR[1]{\NDLR{#1}}
 \newcommand\VOIRAUSSI[1]{\montrucdanger[colback=yellow!20]{VOIR AUSSI}{#1}}
 \newcommand\TODO[1]{\montrucdanger[colback=green!20]{TODO}{#1}}
 \newcommand\ATTENTION[1]{\montruc[colback=brown!20]{\danger ATTENTION#1}}
 \newcommand\COMPRENDPAS[1]{\montrucdanger[colback=yellow!20]{JE
     COMPRENDS PAS CA}{#1}}
  \newcommand\POSSIBLE[1]{\montrucdanger[colback=yellow!10]{
      Serait possible}{#1}}
 \newcommand\ENLEVE[1]{\montrucdanger[colback=gray!20]{
      ENLEVE}{#1}}
 
 
  \newcommand\DEPLACE[1]{
    \vspace{0.5cm}
    \hrule
    \hrule
      \vspace{0.5cm}
    \montrucdangermargin[colback=black!40]{DEPLACE/MISSING CONCEPT
    }{#1}
      \vspace{0.5cm}
    \hrule
    \hrule
      \vspace{0.5cm}
  }

\newcommand\BESOINDE[1]{\montrucdangermargin[colback=pink!20]{ATTTENTION
     BESOIN DE}{#1}}
\newcommand\MISSINGCONCEPT[1]{\montrucdangermargin[colback=pink!20]{MISSING CONCEPT
     }{#1}}
\newcommand\SAUTE[1]{\montrucdangermargin[colback=pink!20]{SAUTE
  }{#1}}
\newcommand\SHOULDBEHERE[1]{\BEGINTEMPENVCOLOR{blue!30}\montrucdanger[colback=pink!10]{
SHOULD
    BE HERE
  }{    
    \vspace{1cm}
    
    \centerline{$\vdots$}
    
   \centerline{ #1}
    
      \centerline{$\vdots$}
          
    \vspace{1cm}
    }\ENDTEMPENVCOLOR}
\newcommand\COULDBEHERE[1]{\montrucdangermargin[colback=pink!20]{COULD
    BE HERE
  }{#1}}



%%%%%%%%%%%%%%%%%%%%%%%%%%%%%%%%%%%%%%%%%%%%%%%%%
%
%
%                      POUR INDICATIONS: SOURCES
%
%
%%%%%%%%%%%%%%%%%%%%%%%%%%%%%%%%%%%%%%%%%%%%%%%%%


\newcommand\SOURCEURL[1]{
  \ifnotmodeDIFFUSIONFINALE{
    \marginpar{ \textcolor{magenta}{Source: \url{#1}} }
    }
}

\newcommand\SOURCETXT[1]{
  \ifnotmodeDIFFUSIONFINALE{
    \marginpar{ \textcolor{magenta}{Source: {#1}} }
    }
}

\newcommand\SOURCE[1]{
  \nocite{#1}
  \ifnotmodeDIFFUSIONFINALE{
    \marginpar{ \textcolor{magenta}{Source: \cite{#1}} }
    }
}


%--------------------------------
%--- Raccourci Sources particulières  ---
%--------------------------------

\newcommand\SOURCEJECH{\SOURCE{jech2003set}}
\newcommand\SOURCEKECHRIS{\SOURCE{kechris2012classical}}
\newcommand\SOURCENOTESSETTHEORY{\SOURCE{moschovakis2006notes}}
\newcommand\SOURCEKOZENCT{\SOURCE{LivreKozenTC}}
\newcommand\SOURCEMARKER{\SOURCE{marker2002descriptive}}
\newcommand\SOURCEMOSCHO{\SOURCE{moschovakis2009descriptive}}



%%%%%%%%%%%%%%%%%%%%%%%%%%%%%%%%%%%%%%%%%%%%%%%%%
%
%
%                      POUR INDICATIONS: DELIRES/COMMENTAIRES
%
%
%%%%%%%%%%%%%%%%%%%%%%%%%%%%%%%%%%%%%%%%%%%%%%%%%


%%%%%% ABOUT PLAN

%\newcommand\PLANBUT[1]{\montruc{\begin{minipage}{0.8\textwidth} #1\end{minipage}}}


%%%%%% ABOUT RECHERCHE


\ifnotmodeDIFFUSIONFINALE{
\newenvironment{recherche}
{}{}
%{\color{blue} BEGIN-RECHERCHE:
%%\ifmodeDIFFUSIONFINALE{diffusion finale}
%%\ifnotmodeDIFFUSIONFINALE{not diffusion finale}
%%	\ifdefined\SAFEMODE
%%		SAFEMODE
%%	\else
%%		NOTSAFEMODE
%%	\fi
%  \todo[inline, color=blue]{recherche ici} }{
%  END-RECHERCHE \color{black} }
\tcolorboxenvironment{recherche}{colback=blue!5!white, colframe=red!75!black,fonttitle=\bfseries, colbacktitle=blue!85!black,enhanced,
attach boxed title to top center={yshift=-2mm},
  title={FAIRE RECHERCHE ICI}}
}


%%%%%% ABOUT DELIRE


%\ifnotmodeDIFFUSIONFINALE{
%\newenvironment{delire}{\color{violet} BEGIN-DELIRE: }{
%  END-DELIRE \color{black} }
%}


%%%%%% ABOUT PERSO


%\ifnotmodeDIFFUSIONFINALE{
%\newenvironment{perso}{\color{brown}
%  BEGIN-PERSO: }{
%  END-PERSO\color{black} }
%}



%%%%%% ABOUT COMMENTAIRE

\ifnotmodeDIFFUSIONFINALE{
\newenvironment{commentaire}{\begin{maboitecommentaire}[colback=brown!30]{Commentaire} Commentaire :}{\end{maboitecommentaire}}
\newcommand\COMMENTAIRE[2][Commentaire ]{
\begin{maboitecommentaire}[colback=brown!30]{#1}
Commentaire : #2
\end{maboitecommentaire}
}
}

%%%%%% ABOUT HORSUJET


%\ifnotmodeDIFFUSIONFINALE{
%\newenvironment{probablementhorsujet}{\color{gray}
%  BEGIN-PROBABLEMENT HORS SUJET: }{
%  END-PROBABLEMENT HORS SUJET \color{black} }
%}

%%%%%% ABOUT BONUS


\ifnotmodeDIFFUSIONFINALE{
\newenvironment{pasbesoinpourlinstant}{NOT NEEDEED FOR
  THIS SECTION: }{ }
    \tcolorboxenvironment{pasbesoinpourlinstant}{colback=green!5!white, colframe=gray!75!black,fonttitle=\bfseries, colbacktitle=green!85!black,enhanced,
attach boxed title to top left={yshift=-2mm},
  title={PAS BESOIN POUR L'INSTANT}}
}

%%%%%% ABOUT BONUS


\ifnotmodeDIFFUSIONFINALE{
\newenvironment{bonus}{BONUS: }{
  }
  \tcolorboxenvironment{bonus}{colback=gray!5!white, colframe=gray!75!black,fonttitle=\bfseries, colbacktitle=gray!85!black,enhanced,
attach boxed title to top left={yshift=-2mm},
  title={BONUS}}
}

%%%%%% ABOUT BONUS


\ifnotmodeDIFFUSIONFINALE{
\newcommand\ATTENTIONREPETE[1]{
  \ATTENTION{Repete (attention modif): #1}
  \begin{remarkolivier}
Repete (attention modif): #1
\end{remarkolivier}
}
}


%%%%%%%%%%%%%%%%%%%%%%%%%%%%%%%%%%%%%%%%%%%%%%%%%
%
%
%                      MACROS Papiers Paulin
%
%
%%%%%%%%%%%%%%%%%%%%%%%%%%%%%%%%%%%%%%%%%%%%%%%%%


\newcommand{\x}{\times}
\newcommand{\s}{\sigma}
\newcommand\sep{.}
\newcommand\bzero{\zero}
\newcommand\bun{\un}
\renewcommand{\a}{\alpha}
\def\hd{{\sf hd}}
\def\tl{{\sf tl}}
\def\cons{{\sf cons}}
\def\Cons{{\sf Cons}}
\def\Pr{{\sf Pr}}
\def\Op{{\sf Op}}
\def\Rel{{\sf Rel}}
\def\Equal{{\sf Equal}}
\def\Eval{{\sf Eval}}
\def\Gate{{\sf Gate}}
\def\Selection{{\sf Selection}}
\def\Result{{\sf Result}}
\def\Circuit{{\sf Circuit}}
\def\next{{\sf next}}
\def\comp{{\sf comp}}
\def\finalcomp{{\sf finalcomp}}
\def\Pick{{\sf Pick}}
\def\Hd{{\sf Hd}}
\def\Tl{{\sf Tl}}
\def\RevHd{{\sf RevHd}}
\def\Node{{\sf Node}}
\def\Length{{\sf Length}}
\def\Tape{{\sf Tape}}
\def\Head{{\sf Head}}
\def\Time{{\sf Time}}
\def\unpad{{\sf unpad}}
\def\concat{{\sf concat}}
\def\rg{{\sf right}}
\def\lf{{\sf left}}
\def\st{{\sf state}}
\def\tope{{\sf top}}


%%%%%%%%%%%%%%%%%%%%%%%%%%%%%%%%%%%%%%%%%%%%%%%%%
%
%
%                      MACROS POLY INF561
%
%
%%%%%%%%%%%%%%%%%%%%%%%%%%%%%%%%%%%%%%%%%%%%%%%%%

\newcommand\buglst[1]{#1}

%\newcommand\mydefref[1]{\expandafter\newcommand\csname #1\endcsname{\ref{#1}}}
%\mydefref{refthrisc}
%\mydefref{refthcinqhuit}
%\mydefref{refthcinqsix}
%\mydefref{refthfondamcaract}
%\mydefref{refdefcircuit}
%%% sinon, fait à la main
%\newcommand\refchapalgo{\ref{chapalgo}}
%\newcommand\refchapmodeles{\ref{chapmodeles}}
%\newcommand\refchappreliminaires{\ref{chappreliminaires}}
%\newcommand\refchapcircuits{\ref{chapcircuits}}
%\newcommand \refpropfondamen{\ref{propfondamen}}
%\newcommand\refthhierspace{ref{thhierspace}}
%\newcommand\refchapcomplexitetemps{\ref{chapcomplexitetemps}}
%\newcommand\refdefverifun{\ref{defverifun}}
%\newcommand\refdefverif{\ref{defverif}}
%\newcommand\refthmachinevs{ref{thmachinevs}}
%\newcommand\reflemnotationconfig{\ref{lemnotationconfig}}
%\newcommand\reflemfondam{\ref{lemfondam}}
%\newcommand\refthtroisat{\ref{thtroisat}}
%\newcommand\refthfondam{\ref{thfondam}}
%\newcommand\refthsansram{\ref{thsansram}}
%\newcommand\refdefsram{\ref{defsram}}
%!TEX encoding = UTF-8 Unicode

%% TRUCS DU POLY INF561

\ifnotSLIDE{
\LANGUE{
\spnewtheoremi{exemple}{\NDLR{ECRIT EXEMPLE AU LIEU DE
  EXAMPLE}}{\bfseries}{\rmfamily}{theorem} 
\spnewtheoremi{remarque}{\NDLR{ECRIT REMARQUE AU LIEU DE REMARK}}{\bfseries}{\rmfamily}{theorem} 
}
{
\spnewtheoremi{exemple}{Exemple}{\bfseries}{\rmfamily}{theorem} 
\spnewtheoremi{remarque}{Remarque}{\bfseries}{\rmfamily}{theorem} 
}
}
            
            \newcommand\PERSOSOURCEPOSSIBLE[1]{\PERSOSOURCE{#1}\PERSOCOMMENTAIRE{#1}}
\newcommand\PERSOPOSSIBLE[1]{\POSSIBLE{#1}}
% \newcommand\PERSOPOSSIBLEbug[1]{\NDLR{\textbf{-- possible --
%       enleve car pb compilation}}}
\newcommand\PERSOETAIT[1]{\NDLR{ETAIT: #1}}
%   \begin{minipage}{10cm}
% \NDLR{\textbf{\tiny  -- était --      #1 -- était
%     --
%   \end{minipage}
% }}}

\providecommand\nl{}

\newcommand\PERSOSOURCE[1] { \SOURCETXT{#1}}
\newcommand\PERSOCOMMENTAIRE[1]{\NDLR{#1}}
\newcommand\PERSOCOMMENTAIREd[2]{\NDLR{#1 \textbf{#2}}}


\newenvironment{dessinpb}
{%\shorthandoff{:}
%\selectlanguage{english}
\begin{center}\begin{tikzpicture}[baseline=(current bounding
      box.north)]
%babelbug with tikz
}
{\end{tikzpicture}
%\selectlanguage{frenchb}
\end{center}}



\newcommand\PERSOMANQUE[1]{\NDLR{\textbf{-- manque --
      #1 --manque--}}}

\newcommand\PERSOENLEVE[1]{\ENLEVE{#1}}






%%%%%%%%%%%%%%%%%%%%%%%%%%%%%%%%%%%%%%%%%%%%%%%%%
%
%
%                      MACROS POLY INF412
%
%
%%%%%%%%%%%%%%%%%%%%%%%%%%%%%%%%%%%%%%%%%%%%%%%%%

%\mydefref{refEEXTtruc}
%\mydefref{refEEXTchappreliminaires}
%\mydefref{refEEXTsecarbrebinaire}
%\mydefref{refEEXTsectermes}
%\mydefref{refEEXTchappredicats}
%\mydefref{refEXTchappredicats}
%\mydefref{refEXTchappreliminaires}
%\mydefref{refEEXTdefTuring}
%\mydefref{refEEXTpropnondet}
%\mydefref{refEEXTchapcompletude}
%\mydefref{refEEXTdefconfigalternative}
%\mydefref{refEEXTchapmachinesdeTuring}
%\mydefref{refEEXTpbdedecision}
%\mydefref{refEEXTchapcalc}
%\mydefref{refEEXTincompletude}
%\mydefref{refEEXTthhierspace}
%\mydefref{refEEXTchapcomplexitetemps}

\tikzstyle{nd} = [shape=circle,draw]


\newcommand\mmod{mod}

  %%% VALUATIONS
\newcommand\mmV{v}
\newcommand\mmVb{\overline{\mmV}}
\newcommand\mmVe{\mmVb}

%%% Corrections

\newcommand\exoref[1]{\ref{exo:#1-stt}}

\newenvironment{correction}[1]
{\par\medskip\noindent \label{exo:#1-cor} \hbox{\bf Exercice \ref{exo:#1-stt} 
  \   (page \pageref{exo:#1-stt}). }}
{\par\medskip}

% PROBLEMES DE DECISION

\newcommand\DECISIONint[3]{
%Le problème de décision #1
\begin{itemize}
\item[\textbf{Donnée}:] #2
\item[\textbf{Réponse}:] #3
\end{itemize}
}

\newcommand\INDECIDABLE[5]{
\DECISION{#1}{#2}{#3}

\begin{#4} \label{#5}
Le problème $#1$ % ATTENTION ENLEVE $#1$ en 2020. REMIS 2023
est indécidable.
\end{#4}
}


\newcommand\INDECIDABLEenglish[5]{
\DECISIONenglish{#1}{#2}{#3}

\begin{#4} \label{#5}
The problem $#1$
is undecidable.
\end{#4}
}

\newcommand\INDECIDABLEq[5]{
\DECISION{#1}{#2}{#3}

\begin{#4} \label{#5}
Le problème $#1$
est ?
\end{#4}
}

\newcommand\INDECIDABLEqenglish[5]{
\DECISION{#1}{#2}{#3}

\begin{#4} \label{#5}
The problem $#1$
is ?
\end{#4}
}



\newcommand\DECISIONintenglish[3]{
%Le problème de décision #1
\begin{itemize}
\item[\textbf{Input}:] #2
\item[\textbf{Answer}:] #3
\end{itemize}
}

\newcommand\DECISIONi[4]{
\item $#1$
%Le problème de décision #1
\DECISIONint{#1}{#2}{#3}
}

\newcommand\DECISIONienglish[4]{
\item $#1$
%Le problème de décision #1
\DECISIONintenglish{#1}{#2}{#3}
}


\newcommand\DECISION[3]{
\begin{definition}[$#1$]\ 

%Le problème de décision #1
\DECISIONint{#1}{#2}{#3}
\end{definition}
}


\newcommand\DECISIONenglish[3]{
\begin{definition}[$#1$]\ 
%
%Le problème de décision #1
\DECISIONintenglish{#1}{#2}{#3}
\end{definition}
}



\newcommand\NPC[4]{ 
\DECISION{#1}{#2}{#3}

\begin{#4}
Le problème $#1$
est $\NP$-complet. \LANGUE{\myindex{NP-completeness}}{\myindex{NP-completude@$\NP$-complétude}}
\end{#4}
}


\newcommand\NPCenglish[4]{ 
\DECISIONenglish{#1}{#2}{#3}
\begin{#4}
The problem $#1$
is $\NP$-complete. \LANGUE{\myindex{NP-completeness}}{\myindex{NP-completude@$\NP$-complétude}}
\end{#4}
}


\newcommand\NPCp[5]{
\DECISION{#1}{#2}{#3}

\begin{#4}[#5]
Le problème $#1$
est $\NP$-complet.
\end{#4}
}


\newcommand\NPCpenglish[5]{
\DECISIONenglish{#1}{#2}{#3}

\begin{#4}[#5]
The problem $#1$
is $\NP$-complete.
\end{#4}
}




\newcommand\BRUNOCOMMENTAIRE[1]{\NDLR{Bruno commentaire: #1}}

\ifnotSLIDE{
%\spnewtheorem{exercicec}{Exercice}{\bfseries}{\rmfamily}
%\spnewtheorem{exerciced}{Exercice}{\bfseries}{\rmfamily}
%\spnewtheorem{exercicedc}{Exercice}{\bfseries}{\rmfamily}
%\spnewtheorem{exercicecenglish}{Exercice}{\bfseries}{\rmfamily}
%\spnewtheorem{exercicedenglish}{Exercice}{\bfseries}{\rmfamily}
%\spnewtheorem{exercicedcenglish}{Exercice}{\bfseries}{\rmfamily}
\spnewtheorem{notation}{Notation}{\bfseries}{\rmfamily}
}

%%% INDEX
 \newcommand\motnouv[1]{\emph{#1}\index{#1}}%\SIGNALEMOTNOUV{#1}}%\index{#1}}
\newcommand\INTmotnouv[1]{\emph{#1}}
\newcommand\myindex[1]{\index{#1}}
%\marginpar{-#1=}

\newcommand\synonyme[2]{\kindex{#1|see{#2}}}
\newcommand\indexM[1]{\index{$#1$}}
\newcommand\idx[1]{#1\index{#1}}
\newcommand\idxm[1]{#1\index{$#1$}}
\newcommand\idxM[1]{$#1$\index{$#1$}}
%\newcommand\idxmBUG[1]{#1}
\newcommand\idxmm[2]{#1\index{#2@$#1$}}
\newcommand\idxi[2]{#1\index{#2}}
%\newcommand\motnouv[1]{\INTmotnouv{#1}\index{#1}}
 \newcommand\motnouvi[2]{\INTmotnouv{#1}\index{#2}} % indexage manuel
 \newcommand\motnouva[2]{\INTmotnouv{#1}\index{#2@#1}} % indexage: 2ième
%                                 % argumet = version sans accent
\newcommand\motnouvm[1]{\INTmotnouv{#1}} %indexage manuel fait ailleurs


\newenvironment{PERSOVERBATIM}{}{}
 \newcommand{\joNP}[3]{{ Problème \textbf{ { \motnouv{##1}}:}
            \begin{itemize}
            \item  \textbf{ Instance:} {##2} 
            \item   \textbf{ Question:} {##3}
            \end{itemize}
          }}

%%%%%%%%%%%%%%%%%%%%%%%%%%%%%%%%%%%%%%%%%%%%%%%%%
%
%
%                      POUR COMPATIBILITE AVEC SLIDES & AUTRE
%
%
%%%%%%%%%%%%%%%%%%%%%%%%%%%%%%%%%%%%%%%%%%%%%%%%%


\newcommand\defin[1]{\textbf{#1}}


%%%   COMPATIBILITE AVEC SLIDES

\newcommand\soustitre[1]{\textbf{#1}}


%% COMPATIBILITE AVEC MA THESE

\providecommand\motnouv[1]{\emph{#1}}
\providecommand\motnouvi[2]{\emph{#1}}

\newcommand\papiernorm[1]{\myop{norm}}

%%% OLIVIER
\newcommand\olivier[1]{ \ATTENTION{Olivier: #1}}


\makeatletter
\@for\@tempa:={soustitre,defin,motnouv,motnouvi,motnouvm}\do{%
  \expandafter\let\csname\@tempa\endcsname\relax}
\makeatother

\usepackage[beamer,expose,english]{olivier}   % french pour un expose en francais

\usepackage[utf8]{inputenc}
\usepackage[T1]{fontenc}
\usepackage{amsmath,amssymb,mathrsfs}
\usepackage{stmaryrd}
\usepackage{multirow}
\usepackage{mathdots}
\usepackage{xcolor}
\usepackage{tikz}
\usetikzlibrary{patterns,decorations.pathmorphing,calc,positioning}
\usepackage{pgfplots}
\pgfplotsset{compat=1.16}
\usepackage{circuitikz}

%%  PIEGE 2. plainurl.bst imprime les DOI via \path, que tikz definit aussi.
%%  Sans cette ligne, la bibliographie casse des qu'une entree a un DOI.
\usepackage{etoolbox}
\AtBeginEnvironment{thebibliography}{\renewcommand{\path}[1]{\nolinkurl{#1}}}

%%%%%%%%%%%%%%%%%%%%%%%%%%%%%%%%%%%%%%%%%%%%%%%%%%%%%%%%%%%%%%%%%%%%%%%%%%%%%%%%%
%%  PALETTE.  Daltonisme protanope: jamais rouge contre vert, et toujours un
%%  indice non colore en plus de la couleur (trait, forme, etiquette).
%%%%%%%%%%%%%%%%%%%%%%%%%%%%%%%%%%%%%%%%%%%%%%%%%%%%%%%%%%%%%%%%%%%%%%%%%%%%%%%%%
\definecolor{ocbleu}{RGB}{23,90,160}
\definecolor{ocambre}{RGB}{176,124,10}
\definecolor{ocgris}{RGB}{110,110,110}

%%%%%%%%%%%%%%%%%%%%%%%%%%%%%%%%%%%%%%%%%%%%%%%%%%%%%%%%%%%%%%%%%%%%%%%%%%%%%%%%%
%%  OUTILS D'EXPOSE
%%%%%%%%%%%%%%%%%%%%%%%%%%%%%%%%%%%%%%%%%%%%%%%%%%%%%%%%%%%%%%%%%%%%%%%%%%%%%%%%%

%%  Parentheses: un bouton dans le flot, un slide en appendice, un retour.
%%  Dans le flot :     \PAREN{Tag}{libelle du bouton}
%%  Dans l'appendice : \begin{myslide}{Titre}\RETOUR{Tag} ... \end{myslide}
%%  Le flot sans les parentheses doit tenir dans le creneau; les parentheses
%%  sont la reserve dans laquelle on pioche si la salle demande.
\newcommand\PAREN[2]{%
  {\raggedleft\hypertarget{back#1}{}\hyperlink{detail#1}{\beamerbutton{#2}}\par}}
\newcommand\RETOUR[1]{%
  \hypertarget{detail#1}{}%
  \gdef\FINPAREN{{\raggedleft\hyperlink{back#1}{\beamerreturnbutton{back}}\par}}}
\newcommand\FINPAREN{}

%%  Attribution uniforme des theoremes. A l'oral, un numero entre crochets ne
%%  dit rien: il faut des noms. Convention retenue: B. pour soi, noms complets
%%  pour les autres, et les etudiants mis en avant en premier quand c'est vrai.
%%      \begin{THEOREME}{Blanc and B., 2024}{CSL24,Icalp2024} ... \end{THEOREME}
\newenvironment{THEOREME}[2]
  {\begin{alertblock}{Theorem (#1) \cite{#2}}}{\end{alertblock}}
\newenvironment{PROPOSITION}[2]
  {\begin{alertblock}{Proposition (#1) \cite{#2}}}{\end{alertblock}}

%%  Storyboard. On ecrit d'abord tous les titres avec un \CONTENU decrivant ce
%%  qui doit y figurer, on compile, on relit le flot, et seulement ensuite on
%%  redige. Retirer la ligne quand le contenu est ecrit.
\newcommand\CONTENU[1]{{\footnotesize\itshape\color{ocgris}[#1]}}

%%  Reperes de demonstration ou de manipulation.
\newcommand\DEMO[1]{{\footnotesize\bfseries\color{ocambre}[DEMO #1]}\quad}

%%%%%%%%%%%%%%%%%%%%%%%%%%%%%%%%%%%%%%%%%%%%%%%%%%%%%%%%%%%%%%%%%%%%%%%%%%%%%%%%%
%%  TROIS PIEGES DE MISE EN PAGE, qui coutent une relecture entiere si on les
%%  laisse s'installer.
%%
%%  PIEGE 3.  \pause n'est PAS une rupture de paragraphe. Ecrire
%%      texte.
%%      \pause
%%      \medskip
%%      suite
%%  recolle les deux blocs en un seul paragraphe qui se reformate a chaque
%%  overlay. Toujours une LIGNE VIDE avant \pause, et l'espacement APRES:
%%      texte.
%%
%%      \pause
%%      \medskip
%%      suite
%%
%%  PIEGE 4.  \centerline ne coupe pas les lignes. Un texte long deborde des
%%  deux cotes du cadre sans aucun avertissement. Utiliser center:
%%      \begin{center} \soustitre{...} \end{center}
%%
%%  PIEGE 5.  \bigskip au milieu d'un paragraphe ne produit rien de visible.
%%  Il lui faut une ligne vide juste avant.
%%%%%%%%%%%%%%%%%%%%%%%%%%%%%%%%%%%%%%%%%%%%%%%%%%%%%%%%%%%%%%%%%%%%%%%%%%%%%%%%%

\begin{document}

\AtBeginSection[]{\frame{\frametitle{Where we are}%
    \tableofcontents[current,hideallsubsections]}}

\title{Titre de l'expose}
\author{\colorsoustitre{Olivier Bournez}}
\institute{\affiliationofficielleolivierslides}
\date{\colorsoustitre{Lieu et date}}

\frame{
\begin{columns}
  \begin{column}{0.25\textwidth}\raggedright
    \includegraphics[height=1.2cm]{\FIGCOMMONS/logo_CNRS}\end{column}
  \begin{column}{0.5\textwidth}\end{column}
  \begin{column}{0.25\textwidth}\raggedleft
    \includegraphics[height=1.2cm]{\FIGCOMMONS/logo_LIX}\end{column}
\end{columns}
\titlepage
\vfill
\begin{columns}
  \begin{column}{0.25\textwidth}\centerline{%
    \includegraphics[height=1.4cm]{\FIGCOMMONS/logo_IPP}}\end{column}
  \begin{column}{0.5\textwidth}\end{column}
  \begin{column}{0.25\textwidth}\centerline{%
    \includegraphics[height=2.5cm]{\FIGCOMMONS/logo_X_new}}\end{column}
\end{columns}
}

%%%%%%%%%%%%%%%%%%%%%%%%%%%%%%%%%%%%%%%%%%%%%%%%%%%%%%%%%%%%%%%%%%%%%%%%%%%%%%%%%
%%  CORPS.  Un titre de section doit dire ce que la section fait, pas nommer un
%%  concept: "Simulating a Turing machine with a GPAC" et non "How it works".
%%%%%%%%%%%%%%%%%%%%%%%%%%%%%%%%%%%%%%%%%%%%%%%%%%%%%%%%%%%%%%%%%%%%%%%%%%%%%%%%%

\section{Premiere partie}

\begin{myslide}{Un slide ordinaire}
Premier bloc de texte.

\pause
\medskip
Second bloc, apres la pause. Noter la ligne vide avant \texttt{\textbackslash pause}.

\pause
\medskip
\begin{center} \soustitre{Une chute centree, en center et non en centerline.}
\end{center}

\PAREN{Ex}{une parenthese}
\end{myslide}

\begin{myslide}{Un slide qui porte un resultat}
\begin{THEOREME}{Blanc and B., 2024}{CSL24,Icalp2024}
Enonce du theoreme.
\end{THEOREME}

\pause
\medskip
Commentaire.
\end{myslide}

\begin{myslide}{Un slide en storyboard}
\CONTENU{Ce qui doit figurer ici: la figure recurrente, la these en bloc
encadre, et le renvoi vers la section suivante.}
\end{myslide}

%%%%%%%%%%%%%%%%%%%%%%%%%%%%%%%%%%%%%%%%%%%%%%%%%%%%%%%%%%%%%%%%%%%%%%%%%%%%%%%%%
%%  APPENDICE.  Les parentheses, plus le filet de securite.
%%%%%%%%%%%%%%%%%%%%%%%%%%%%%%%%%%%%%%%%%%%%%%%%%%%%%%%%%%%%%%%%%%%%%%%%%%%%%%%%%

\appendix

\begin{myslide}{Le detail mis en parenthese}\RETOUR{Ex}
Contenu qu'on ne montre que si la salle le demande.
\FINPAREN
\end{myslide}

%%  PIEGE 6.  alphaurl et non plainurl. plainurl imprime [42], illisible pour
%%  qui ecoute; alphaurl imprime [ADK24], [BGP17], que la salle relie au nom
%%  entendu. Meme famille LIPIcs, meme gestion des DOI.
\begin{frame}[fragile,allowframebreaks]{References}
\bibliographystyle{alphaurl}
\bibliography{\BIBFILES}
\end{frame}

\end{document}

%%%%%%%%%%%%%%%%%%%%%%%%%%%%%%%%%%%%%%%%%%%%%%%%%%%%%%%%%%%%%%%%%%%%%%%%%%%%%%%%%
%%  ESTIMER LA DUREE, avant d'etre en retard le jour J.
%%
%%  Pour chaque slide, compter b = 1 + nombre de \pause, et d = nombre de mots
%%  visibles + 12 par figure + 15 par bloc de theoreme. Alors
%%      minutes ~ 0,45 + 0,28 b + 0,0125 d
%%  calibre pour qu'un slide moyen de tutoriel tombe vers deux minutes. Ajouter
%%  quinze secondes par sommaire de section, le temps reel des demonstrations, et
%%  une marge pour les questions dans le flot. La dispersion d'un orateur qui
%%  connait son sujet est de l'ordre de quinze pour cent: lire le resultat comme
%%  une fourchette, pas comme un chiffre.
%%
%%  Regle pratique: le flot SANS les parentheses doit tenir dans le creneau.
%%  Tout ce qui deborde part en parenthese plutot qu'a la poubelle.
%%
%%  NUMEROTATION.  Le numero imprime en bas a droite compte la page de titre et
%%  les sommaires de section. C'est celui-la qu'il faut citer quand on discute
%%  d'un slide a distance, sinon les decomptes divergent de deux ou trois.
%%%%%%%%%%%%%%%%%%%%%%%%%%%%%%%%%%%%%%%%%%%%%%%%%%%%%%%%%%%%%%%%%%%%%%%%%%%%%%%%%
